% Options for packages loaded elsewhere
\PassOptionsToPackage{unicode}{hyperref}
\PassOptionsToPackage{hyphens}{url}
\PassOptionsToPackage{dvipsnames,svgnames,x11names}{xcolor}
\documentclass[
  10pt,
  a4paper,
  openany]{book}
\usepackage{xcolor}
\usepackage[inner=29mm,outer=24mm,top=25mm,bottom=27mm]{geometry}
\usepackage{amsmath,amssymb}
\setcounter{secnumdepth}{5}
\usepackage{iftex}
\ifPDFTeX
  \usepackage[T1]{fontenc}
  \usepackage[utf8]{inputenc}
  \usepackage{textcomp} % provide euro and other symbols
\else % if luatex or xetex
  \usepackage{unicode-math} % this also loads fontspec
  \defaultfontfeatures{Scale=MatchLowercase}
  \defaultfontfeatures[\rmfamily]{Ligatures=TeX,Scale=1}
\fi
\usepackage{lmodern}
\ifPDFTeX\else
  % xetex/luatex font selection
\fi
% Use upquote if available, for straight quotes in verbatim environments
\IfFileExists{upquote.sty}{\usepackage{upquote}}{}
\IfFileExists{microtype.sty}{% use microtype if available
  \usepackage[]{microtype}
  \UseMicrotypeSet[protrusion]{basicmath} % disable protrusion for tt fonts
}{}
\makeatletter
\@ifundefined{KOMAClassName}{% if non-KOMA class
  \IfFileExists{parskip.sty}{%
    \usepackage{parskip}
  }{% else
    \setlength{\parindent}{0pt}
    \setlength{\parskip}{6pt plus 2pt minus 1pt}}
}{% if KOMA class
  \KOMAoptions{parskip=half}}
\makeatother
\usepackage{longtable,booktabs,array}
\usepackage{calc} % for calculating minipage widths
% Correct order of tables after \paragraph or \subparagraph
\usepackage{etoolbox}
\makeatletter
\patchcmd\longtable{\par}{\if@noskipsec\mbox{}\fi\par}{}{}
\makeatother
% Allow footnotes in longtable head/foot
\IfFileExists{footnotehyper.sty}{\usepackage{footnotehyper}}{\usepackage{footnote}}
\makesavenoteenv{longtable}
\setlength{\emergencystretch}{3em} % prevent overfull lines
\providecommand{\tightlist}{%
  \setlength{\itemsep}{0pt}\setlength{\parskip}{0pt}}
% A traditional TeX typographic palette.  Latin Modern Math is the
% OpenType continuation of Computer Modern and gives formulas the familiar
% default LaTeX appearance requested for this edition.
\IfFontExistsTF{Latin Modern Roman}{\setmainfont{Latin Modern Roman}}{}
\IfFontExistsTF{Latin Modern Sans}{\setsansfont{Latin Modern Sans}}{}
\IfFontExistsTF{Latin Modern Mono}{\setmonofont{Latin Modern Mono}}{}
\setmathfont{latinmodern-math.otf}
\usepackage{mathtools}
\usepackage{slashed}
\usepackage{booktabs}
\usepackage{longtable}
\usepackage{array}
\usepackage{enumitem}
\usepackage{fancyhdr}
\usepackage[most]{tcolorbox}
\usepackage{titlesec}
\usepackage{etoolbox}

\definecolor{chapterblue}{HTML}{173B57}
\definecolor{accent}{HTML}{8B4A32}
\definecolor{paleblue}{HTML}{F2F6F8}
\definecolor{rulegray}{HTML}{AEBBC4}
\definecolor{titlegray}{HTML}{53616B}

\titleformat{\chapter}[display]
  {\normalfont\sffamily\bfseries\color{chapterblue}}
  {\filleft\Large\chaptername\ \thechapter}
  {0.6ex}
  {\titlerule[1.1pt]\vspace{1.35ex}\huge}
  [\vspace{1.1ex}\titlerule]
\titleformat{\section}
  {\Large\sffamily\bfseries\color{chapterblue}}
  {\thesection}{0.7em}{}
\titleformat{\subsection}
  {\large\sffamily\bfseries\color{accent}}
  {\thesubsection}{0.65em}{}
\titleformat{\subsubsection}
  {\normalsize\sffamily\bfseries\color{chapterblue}}
  {\thesubsubsection}{0.6em}{}

\setlength{\parindent}{0pt}
\setlength{\parskip}{0.58em plus 0.16em minus 0.1em}
\setlist{itemsep=0.28em,topsep=0.4em,leftmargin=2em}
\linespread{1.025}
\setlength{\abovedisplayskip}{0.72em plus 0.18em minus 0.12em}
\setlength{\belowdisplayskip}{0.72em plus 0.18em minus 0.12em}
\allowdisplaybreaks[2]
\raggedbottom

\pagestyle{fancy}
\fancyhf{}
\renewcommand{\chaptermark}[1]{\markboth{\thechapter\enspace #1}{}}
\fancyhead[LE]{\footnotesize\sffamily\thepage\quad\textcolor{rulegray}{|}\quad\nouppercase{\leftmark}}
\fancyhead[RO]{\footnotesize\sffamily\nouppercase{\leftmark}\quad\textcolor{rulegray}{|}\quad\thepage}
\renewcommand{\headrulewidth}{0.35pt}
\setlength{\headheight}{14pt}
\fancypagestyle{plain}{\fancyhf{}\fancyfoot[C]{\small\sffamily\thepage}\renewcommand{\headrulewidth}{0pt}}

\AtBeginEnvironment{quote}{%
  \begin{tcolorbox}[enhanced,breakable,colback=paleblue,frame hidden,
    borderline west={1.15pt}{0pt}{chapterblue},sharp corners,
    left=3mm,right=2.5mm,top=1.8mm,bottom=1.8mm]}
\AtEndEnvironment{quote}{\end{tcolorbox}}

% A restrained title page replaces the sparse book-class default.
\makeatletter
\providecommand{\@subtitle}{}
\providecommand{\subtitle}[1]{\gdef\@subtitle{#1}}
\renewcommand{\maketitle}{%
  \begin{titlepage}
    \thispagestyle{empty}
    \centering
    \vspace*{0.14\textheight}
    {\color{chapterblue}\rule{0.78\textwidth}{1.25pt}\par}
    \vspace{2.2em}
    {\sffamily\bfseries\color{chapterblue}\fontsize{29}{35}\selectfont
      \@title\par}
    \vspace{1.4em}
    {\sffamily\color{titlegray}\fontsize{15}{21}\selectfont
      \@subtitle\par}
    \vspace{2.2em}
    {\color{chapterblue}\rule{0.38\textwidth}{0.55pt}\par}
    \vfill
    {\sffamily\large\color{titlegray}\@author\par}
    \vspace{0.7em}
    {\sffamily\normalsize\color{titlegray}\@date\par}
    \vspace*{0.08\textheight}
  \end{titlepage}}
\makeatother

\newcommand{\Tr}{\operatorname{Tr}}
\newcommand{\Str}{\operatorname{Str}}
\newcommand{\diag}{\operatorname{diag}}
\newcommand{\sign}{\operatorname{sign}}
\newcommand{\id}{\mathbf 1}
\newcommand{\ket}[1]{\lvert #1\rangle}
\newcommand{\bra}[1]{\langle #1\rvert}
\newcommand{\braket}[2]{\langle #1\mid #2\rangle}
\newcommand{\abs}[1]{\lvert #1\rvert}
\newcommand{\norm}[1]{\lVert #1\rVert}
\newcommand{\dd}{\mathrm d}

\AtBeginDocument{\hypersetup{
  pdftitle={Basic Concepts of String Theory: Reading Companion},
  pdfsubject={Consolidated reading questions and answers},
  pdfauthor={Study notebook edition},
  bookmarksopen=true,
  bookmarksnumbered=true}}
\usepackage{bookmark}
\IfFileExists{xurl.sty}{\usepackage{xurl}}{} % add URL line breaks if available
\urlstyle{same}
\hypersetup{
  pdftitle={Basic Concepts of String Theory: Reading Companion},
  pdfauthor={Study notebook edition},
  colorlinks=true,
  linkcolor={MidnightBlue},
  filecolor={Maroon},
  citecolor={Blue},
  urlcolor={BrickRed},
  pdfcreator={LaTeX via pandoc}}

\title{Basic Concepts of String Theory: Reading Companion}
\usepackage{etoolbox}
\makeatletter
\providecommand{\subtitle}[1]{% add subtitle to \maketitle
  \apptocmd{\@title}{\par {\large #1 \par}}{}{}
}
\makeatother
\subtitle{Consolidated questions, derivations, source map, and book
roadmap}
\author{Study notebook edition}
\date{August 2026}

\begin{document}
\frontmatter
\maketitle

{
\hypersetup{linkcolor=}
\setcounter{tocdepth}{1}
\tableofcontents
}
\mainmatter
\setcounter{chapter}{-1}

\chapter{Consolidated Reading
Summary}\label{consolidated-reading-summary}

This chapter collects the short conclusions from the reading
conversations. Detailed derivations and textbook locations appear in the
chapters that follow. Exact duplicate files were consolidated once; the
empty \texttt{Superstrings.md} file contributes no Q\&A material.

\section{Bosonic One-Loop Amplitudes}\label{bosonic-one-loop-amplitudes}

\begin{itemize}
\tightlist
\item
  \textbf{Question 1:} For the closed bosonic string torus amplitude,
  why does the Hamiltonian trace over a fixed torus have to be
  integrated over \(\tau\), and why is the measure \(d^2\tau/\tau_2^2\)?
  The key point is that \(\tau\) is a genuine complex structure modulus
  left after conformal gauge fixing, and \(d^2\tau/\tau_2^2\) is the
  natural Poincare/Faddeev-Popov volume form on
  \(\mathbb H/SL(2,\mathbb Z)\), not an extra physical Boltzmann weight.
\item
  \textbf{Question 2:} For the open bosonic string annulus/cylinder
  amplitude, where does the factor \(dt/(2t)\) come from? The key point
  is that the annulus has one real modulus \(t\) with logarithmic
  moduli-space measure \(dt/t\), and the standard extra factor \(1/2\)
  comes from the residual vacuum-cylinder automorphism/conformal-Killing
  normalization.
\item
  \textbf{Question 3:} A static flat D-brane with \(F=0\) has no source
  linear in the antisymmetric \(B\)-field because the trace of a
  symmetric inverse metric with \(B_{\alpha\beta}\) vanishes. Hence
  there is no one-\(B\) exchange, whereas the linear graviton and
  dilaton sources generate the NS--NS attraction.
\end{itemize}

\section{Torus Partition Functions for
RCFTs}\label{torus-partition-functions-for-rcfts}

\begin{itemize}
\tightlist
\item
  \textbf{Question 1:} The modular transformation \(S\) exchanges the
  two torus cycles. Applying it twice reverses both cycles and acts on a
  chiral sector by conjugating its representation. Therefore \(S^2=C\),
  where \(C_{ij}=\delta_{i,j^*}\) pairs each primary with its conjugate.
\end{itemize}

\section{The Classical Fermionic String and Spin
Structures}\label{the-classical-fermionic-string-and-spin-structures}

\begin{itemize}
\tightlist
\item
  \textbf{Question 1:} On a compact Riemann surface, the chiral Dirac
  operator splits into first-order operators on the two spin bundles; in
  a conformal coordinate the zero-mode equations reduce, after the
  standard Weyl rescaling, to holomorphic or anti-holomorphic
  half-differential equations. For a spin structure \(L\) with
  \(L^2=K\), the parity \(h^0(\Sigma,L)\bmod 2\) is topological: it
  equals the Arf invariant of the quadratic refinement associated to the
  spin structure. Under gluing by connected sum, the quadratic
  refinements add orthogonally, so the mod-two number of chiral zero
  modes is additive.
\item
  \textbf{Question 2:} The \(2^{2g}\) spin structures are equivalently
  quadratic refinements \(q\) of the mod-two intersection form on
  \(H_1(\Sigma_g,\mathbb Z_2)\). In a symplectic basis,
  \(\operatorname{Arf}(q)=\sum_i q(a_i)q(b_i)\), so each handle has
  three even local choices and one odd local choice. Hence the odd spin
  structures are counted by choosing an odd number of odd handles,
  giving \(2^{g-1}(2^g-1)\), while the even spin structures are counted
  by choosing an even number, giving \(2^{g-1}(2^g+1)\).
\item
  \textbf{Question 3:} The apparent
  \(\vartheta\begin{bmatrix}\alpha\\-\beta\end{bmatrix}\) versus
  \(\vartheta\begin{bmatrix}\alpha\\\beta\end{bmatrix}\) mismatch is
  only a characteristic-convention phase, and for
  \(\beta=0,\frac{1}{2}\) the spin structure is the same modulo
  integers; any remaining sign can be absorbed into the spin-structure
  phase \(\eta_{\alpha\beta}\). Also \(\vartheta_1(\tau)=0\) means
  \(\vartheta_1(0|\tau)=0\), while
  \(\vartheta_1'(\tau)=2\pi\eta^3(\tau)\) means
  \(\partial_\nu\vartheta_1(\nu|\tau)|_{\nu=0}=2\pi\eta^3(\tau)\), so
  there is no contradiction.
\item
  \textbf{Question 4:} The \(\eta=\pm 1\) in the closed-string boundary
  state is a boundary gluing label, not by itself the open-string R/NS
  sector. Under the open-closed channel exchange, the two closed-string
  boundaries have opposite induced orientations. Bosons are insensitive
  to the square root of this orientation reversal, but fermions are
  half-differentials, so the orientation reversal contributes a relative
  phase whose square is \(-1\); with the usual gluing convention this
  gives \(s_0s_l=-\eta_1\eta_2\). Hence same closed labels give the open
  NS trace, while opposite closed labels give the open R trace. In the
  NS-NS sector, \((-1)^F\) flips the sign of the left-moving fermion
  oscillators and the odd NS vacuum gives
  \((-1)^F|\mathrm{Dp},\eta\rangle_{\mathrm{NS}}=-|\mathrm{Dp},-\eta\rangle_{\mathrm{NS}}\),
  and similarly for \((-1)^{\bar F}\).
\item
  \textbf{Question 5:} The sign of a static exchange potential comes
  from the propagator numerator contracted with static source
  components. In mostly-plus signature, scalar and graviton exchange
  give \(V\sim-q_1q_2G(r)\) for positive/equal charges, hence
  attraction, while vector exchange gives \(V\sim+q_1q_2G(r)\), hence
  repulsion. A R-R \((p+1)\)-form coupled to a D\(p\)-brane behaves like
  the vector case for parallel branes: reversing orientation reverses
  the R-R charge, and equal R-R charges repel. For BPS parallel D-branes
  this R-R repulsion exactly cancels the NS-NS attraction from graviton,
  dilaton, and related fields.
\item
  \textbf{Question 6:} The type II torus vacuum amplitude vanishes
  because the GSO-summed worldsheet path integral computes a spacetime
  supertrace, not an ordinary thermal trace. In unbroken spacetime
  supersymmetry, each bosonic physical state has a fermionic partner of
  the same mass and degeneracy, so
  \(\operatorname{Str}e^{-2\pi t M^2}=0\) level by level. On the
  worldsheet this cancellation is the Jacobi abstruse identity
  \(\vartheta_3^4-\vartheta_4^4-\vartheta_2^4=0\).
\item
  \textbf{Question 7:} In tree channel, the tadpole factorization has
  the schematic form
  \(\frac{1}{2}\langle B+C|\Delta|B+C\rangle=\frac{1}{2}\langle B|\Delta|B\rangle+\langle B|\Delta|C\rangle+\frac{1}{2}\langle C|\Delta|C\rangle\).
  The annulus and Klein bottle are the two half-square terms, while the
  single Möbius strip amplitude is already the full boundary-crosscap
  cross term. Adding two copies of the Möbius amplitude would
  double-count the same unoriented worldsheet/moduli region.
\item
  \textbf{Question 8:} The orientifold projection \(\Omega\) acts on a
  BPS D\(p\) boundary state by leaving the NS-NS GSO-projected part
  invariant but multiplying the R-R GSO-projected part by
  \((-1)^{(p+3)/2}\) for odd \(p\). Therefore the R-R charged BPS branes
  invariant under type I \(\Omega\) are \(p=1,5,9\). Equivalently,
  \(\Omega\) keeps the R-R potentials \(C_2\), its magnetic dual
  \(C_6\), and the top-form \(C_{10}\), while projecting out
  \(C_0,C_4,C_8\).
\end{itemize}

\section{The Quantized Fermionic
String}\label{the-quantized-fermionic-string}

\begin{itemize}
\tightlist
\item
  \textbf{Question 1:} In old-covariant quantization the oscillator Fock
  space has indefinite metric because of timelike oscillators. The
  no-ghost theorem says that, at the critical values, the physical state
  space modulo null/spurious states is positive definite and is
  isomorphic to the transverse light-cone/DDF Fock space: for the
  bosonic string \(D=26\), \(a=1\); for the RNS superstring \(D=10\),
  with \(a_{\text{NS}}=\frac{1}{2}\) and \(a_{\text{R}}=0\). Thus
  negative-norm states are gauge artifacts or null states and decouple
  from amplitudes.
\item
  \textbf{Question 2:} The NS oscillator \(b_{-1/2}^i\) is fermionic as
  a worldsheet oscillator, but it carries an \(SO(8)\) vector index, so
  the states \(b_{-1/2}^i|0;k\rangle_{\text{NS}}\) form a vector basis
  of the massless one-string Hilbert space. The polarization \(\zeta_i\)
  is not more physical than the basis state; rather
  \(|\zeta;k\rangle=\zeta_i b_{-1/2}^i|0;k\rangle_{\text{NS}}\) is the
  physical one-particle state. In second quantization, the
  momentum-space gauge field \(A_i(k)\) is the coefficient multiplying
  this string state in the string field. Spacetime fermions instead come
  from the R sector, whose zero modes form a little-group Clifford
  algebra and produce spinor representations.
\end{itemize}

\section{Kac--Moody Algebras}\label{kacmoody-algebras}

\begin{itemize}
\tightlist
\item
  \textbf{Question 1:} A representation weight pairs integrally with
  every coroot; for simply-laced roots normalized by
  \(\vec\alpha^{\,2}=2\), this gives
  \(\vec\alpha\cdot\vec m\in\mathbb Z\). Affine unitarity then requires
  \(\vec\psi\cdot\vec m_0\leq k\), where \(\vec\psi\) and \(\vec m_0\)
  are the highest root and highest weight. At level one, the ADE
  primaries are the vacuum and the minuscule fundamental
  representations, one minimal representative for each class in \(P/Q\).
\end{itemize}

\section{Covariant Vertex Operators, BRST, and Covariant
Lattices}\label{covariant-vertex-operators-brst-and-covariant-lattices}

\begin{itemize}
\tightlist
\item
  \textbf{Question 1:} For the bosonic string, an integrated open-string
  matter vertex must have \(h=1\), while an integrated closed-string
  matter vertex must have \((h,\bar h)=(1,1)\), so that its insertion is
  invariant under the residual conformal reparametrizations.
  Equivalently, these are the \(L_0\) physical-state condition with
  intercept one and the condition for the ghost-dressed vertex
  \(cV_{\mathrm m}\) or \(c\bar cV_{\mathrm m}\) to be BRST closed.
  Weight one is the on-shell condition, not by itself a masslessness
  condition: tachyon vertices also have weight one.
\item
  \textbf{Question 2:} In the RNS superstring, the complete integrated
  operator, including its \(\beta\gamma\) superghost dressing but
  excluding the \(c\) ghost, has \(h=1\) for an open string and
  \((h,\bar h)=(1,1)\) for a closed string. The matter factor alone is
  picture-dependent: for example, the massless NS matter factor has
  \(h=1/2\) in picture \(-1\) but the superghost \(e^{-\phi}\) supplies
  the missing \(1/2\).
\item
  \textbf{Question 3:} The exponential map from the open-string strip to
  the upper half-plane sends \(\sigma=0\) to the positive real axis and
  \(\sigma=l\) to the negative real axis. Because the worldsheet
  fermions have weight \(1/2\), their square-root Jacobians have
  relative sign \(-1\) on the negative axis. Thus the general gluing is
  \(D(0)\) for \(x>0\) and \(-\eta_{\mathrm{strip}}D(l)\) for \(x<0\);
  when both ends lie on the same D-brane, this becomes \(+D\) everywhere
  in the NS sector and \(\operatorname{sign}(x)D\) in the R sector, as
  in Eq. (13.99).
\item
  \textbf{Question 4:} In BCFT, \(\bar T(\bar z)\) denotes the
  antiholomorphic stress-tensor component, not automatically the
  Hermitian adjoint of \(T(z)\). The conditions \(T(x)=\bar T(x)\) and
  \(\epsilon(x)=\bar\epsilon(x)\) express vanishing boundary flux and
  boundary-preserving conformal transformations, leaving one diagonal
  Virasoro algebra rather than two independent copies. Applying the same
  contour argument to D-brane supersymmetry currents preserves
  \(Q_L+P Q_R\); GSO chirality then requires even \(p\) in type IIA and
  odd \(p\) in type IIB.
\item
  \textbf{Question 5:} The NS/R sector of a bosonized vertex
  \(V_{\boldsymbol\lambda}\) is determined by the monodromy of a
  fundamental fermion around its insertion, not by rotating
  \(V_{\boldsymbol\lambda}\) itself. The OPE
  \(\Psi^{\pm j}(z)V_{\boldsymbol\lambda}(0)\sim z^{\pm\lambda_j}(\cdots)\)
  gives monodromy \(e^{\pm2\pi i\lambda_j}\): integer \(\lambda_j\) give
  the plane NS sector and comprise \((0)\cup(V)\), while half-integer
  \(\lambda_j\) give the plane R sector and comprise \((S)\cup(C)\).
\item
  \textbf{Question 6:} The one-loop trace is not a trace over already
  on-shell covariant states with \(L_0=1\). It is either a trace over
  the gauge-fixed worldsheet CFT state space or, in light-cone gauge,
  over the unconstrained transverse oscillator Fock space together with
  off-shell Euclidean loop momentum. The cylinder Hamiltonian is
  \(L_0-c/24\); for the 24 transverse bosons this is \(L_0-1\). In Eq.
  (13.114), \(L_0-1\) is likewise the chiral worldsheet propagation
  operator before the \(c\)-ghost dressing, while \((-1)^{F_R}\) makes
  the trace a spacetime supertrace.
\end{itemize}

\section{CFTs for Type II and Heterotic String
Vacua}\label{cfts-for-type-ii-and-heterotic-string-vacua}

\begin{itemize}
\tightlist
\item
  \textbf{Question 1:} For a toroidal orbifold \(T^6/\mathbb Z_N\), the
  untwisted Hodge classes are the \(\mathbb Z_N\)-invariant constant
  forms on \(T^6\). The twisted classes are computed from the cohomology
  of the fixed loci of each \(\theta^k\), shifted in bidegree by the age
  of the twist and projected by the full orbifold group. For case 1,
  \(T^6/\mathbb Z_3\) on \(A_2^3\), this gives
  \((h^{1,1},h^{2,1})=(9,0)+(27,0)=(36,0)\). For case 2,
  \(T^6/\mathbb Z_4\) on \(A_1^2\times B_2^2\), the \(\theta\) fixed
  points and the \(\theta^2\) fixed tori give
  \((h^{1,1},h^{2,1})=(5,1)+(16,0)+(10,6)=(31,7)\).
\item
  \textbf{Question 2:} The extra orbifold gauge bosons do not arise from
  twisted sectors. They are untwisted ten-dimensional \(E_8\)
  current-algebra states whose generators commute with the discrete
  gauge shift \(\gamma_{\mathbf V}\) used to embed the orbifold point
  group. A smooth standard embedding has full \(SU(3)\) gauge-bundle
  holonomy and leaves only its commutant \(E_6\), whereas the orbifold
  background has only a discrete subgroup generated by
  \(\gamma_{\mathbf V}\) and hence the larger centralizer
  \(E_6\times G_{\mathrm{enh}}\). For equal, two equal, or three
  distinct eigenvalues of \(\gamma_{\mathbf V}\) in \(SU(3)\),
  \(G_{\mathrm{enh}}\) is respectively \(SU(3)\), \(SU(2)\times U(1)\),
  or \(U(1)^2\). Charged twisted blow-up modes Higgs this extra factor
  when they acquire expectation values.
\item
  \textbf{Question 3:} T-duality along \(n\) coordinates conjugates the
  type-IIB orientifold generator to
  \(\mathcal T_n\Omega\mathcal T_n^{-1}=\Omega I_n[(-1)^{F_L}]^{n(n-1)/2}\).
  The reflection \(I_n\) comes from the bosonic asymmetric reflection,
  while the \((-1)^{F_L}\) factor for \(n=2,3\pmod4\) comes from the
  projective action of reflections on Ramond spinors. For
  \(n=0,1,2,3,4\) this reproduces respectively the O9, O8, O7, O6, and
  O5 rows of Table 10.3, including their type-II theory, orientifold
  action, fixed-plane multiplicity, and local Chan--Paton gauge group.
\item
  \textbf{Question 4:} Tensoring the chiral states in Table 15.3,
  imposing equal \(I_4\) parity, and then projecting with
  \(\Omega R(-1)^{F_L}\) gives the untwisted spectrum \((3,3)+11(1,1)\)
  in NS--NS, \((3,1)+(1,3)+6(1,1)\) in R--R, and \(2(2,3)+10(2,1)\) in
  the mixed sectors. The chosen extra minus sign on each twisted sector
  leaves, at each of the 16 fixed points, \(3(1,1)\) in NS--NS,
  \((1,1)\) in R--R, and \(2(2,1)\) in the mixed sectors. These states
  form one gravity, one tensor, and four hypermultiplets in the
  untwisted sector, plus sixteen twisted hypermultiplets.
\item
  \textbf{Question 5:} Because the theory is already orbifolded by
  \(I_4\), its orientifold group contains two orientation-reversing
  elements, \(\Omega R(-1)^{F_L}\) and \(\Omega RI_4(-1)^{F_L}\). The
  corresponding terms in the Klein-bottle trace identify two crosscaps
  whose geometric parts fix respectively four two-tori parallel to
  \((x^1,x^2)\) and four two-tori parallel to \((y^1,y^2)\) in the
  covering \(T^4\). Equivalently, a fixed point of the induced
  reflection on \(T^4/\mathbb Z_2\) may obey either \(Rp=p\) or
  \(Rp=I_4p\). The insertion and its zero-mode lattice sum therefore
  determine the O7-plane locus, but not as the locus of open-string
  endpoints: the transverse Klein bottle is closed-string propagation
  between crosscaps and has no worldsheet boundary. Open-string
  endpoints attach only to D-branes; annulus and Möbius amplitudes
  contain boundary states.
\item
  \textbf{Question 6:} For an \(I_4\)-invariant D7-brane wrapping
  \((x^1,x^2)\), the \(I_4\) projection retains the six-dimensional
  gauge vector while removing transverse-position and Wilson-line
  scalars. After the annulus modular transformation, the inserted R
  contribution is \(I_4\)-twisted R--R propagation, proving that the
  fractional boundary state carries twisted charge under
  \(C_6^{(i)}=\int_{E_i}C_8\). The two fractional branes differ by the
  sign of this twisted boundary-state component. Because this model
  chooses \(\Omega R(-1)^{F_L}\) to act with an extra minus sign on
  every \(I_4\)-twisted state, it leaves the untwisted component fixed
  but reverses the twisted component, giving
  \(|D7^\pm\rangle_{\mathrm{frac}}\mapsto|D7^\mp\rangle_{\mathrm{frac}}\).
  An invariant configuration therefore contains the pair, whose twisted
  charges and tadpoles cancel.
\item
  \textbf{Question 7:} Equation (15.52) follows by requiring the vertex
  in (15.51), before the \(bc\) ghosts are attached, to have conformal
  weights \((h,\bar h)=(1,1)\). The oscillator factors contribute
  \(N_L,N_{1R},N_{2R}\), the fermion-lattice exponential contributes
  \(\boldsymbol\lambda_R^2/2\), the picture factor contributes
  \(-q(q+2)/2\), the internal operator contributes
  \((h_{\mathrm{int}},\bar h_{\mathrm{int}})\), and the plane wave
  contributes \(-\alpha'm^2/4\) to each chirality. The book packages the
  two chiral weight conditions as \(m^2=m_L^2+m_R^2\), while level
  matching gives \(m_L^2=m_R^2=m^2/2\) for a physical closed-string
  state.
\item
  \textbf{Question 8:} The significant momentum factors in (15.84) are
  the extra \(\epsilon_\mu k_\nu\) and \(k_\mu\), not the universal
  plane wave \(e^{ik\cdot X}\). Antisymmetry of \(\sigma^{\mu\nu}\)
  turns \(\epsilon_\mu k_\nu\) into the graviphoton field strength
  \(F_{\mu\nu}\), while \(k_\mu C(k)\) is the Fourier transform of
  \(\partial_\mu C\). This is the four-dimensional form of the general
  R--R fact that BRST-invariant spin-field polarizations are
  field-strength bispinors. Hence perturbative local bulk couplings
  contain R--R potentials through their field strengths and do not
  produce a minimal coupling \(A_\mu J^\mu\) to perturbative charged
  matter. Wrapped D-branes can nevertheless carry R--R charge, and
  topological or Chern--Simons terms require the usual qualification.
\item
  \textbf{Question 9:} In the right-moving NS sector, BRST invariance of
  an integrated picture-zero weight-one operator requires its
  supercurrent variation to be a total derivative. If
  \(G(z)W(w)\sim A/(z-w)^2+B/(z-w)\), the \(\gamma G\) part of the BRST
  charge gives \((\partial\gamma)A+\gamma B\), which is a total
  derivative precisely when \(B=\partial A\). For a weight-\(1/2\)
  superprimary \(\hat J^a\), the \(N=1\) algebra gives
  \(J^a=G_{-1/2}\hat J^a\) and
  \(G(z)J^a(w)\sim\hat J^a/(z-w)^2+\partial\hat J^a/(z-w)\). This
  derives the superfield condition and the canonical vertex (15.91). A
  full finite-momentum picture change gives
  \(J^a+(k\cdot\psi)\hat J^a\); therefore (15.92), as printed, is exact
  at \(k=0\) or shorthand for its internal-current term, but omits the
  universal momentum-dependent completion required at \(k\neq0\).
\end{itemize}

\section{String Scattering Amplitudes and Low-Energy Effective Field
Theory}\label{string-scattering-amplitudes-and-low-energy-effective-field-theory}

\begin{itemize}
\tightlist
\item
  \textbf{Question 1:} A conventional first-quantized string calculation
  determines only the on-shell spacetime \(S\)-matrix. An ``off-shell
  amplitude'' in the low-energy field theory is more precisely an
  amputated Green function or an \(n\)-point vertex obtained by
  differentiating an effective action at momenta not obeying the free
  equations of motion. Such quantities are needed to define propagators,
  internal-line exchange, equations of motion, background dependence,
  and loop calculations, but they are not separately observable and are
  not uniquely fixed by the string \(S\)-matrix: local field
  redefinitions, gauge choices, and operators proportional to the
  equations of motion change them while leaving every on-shell amplitude
  unchanged.
\item
  \textbf{Question 2:} The three-tachyon matter factor
  \(\prod_{i<j}|z_{ij}|^{-2}\) exactly cancels the sphere ghost factor
  \(|z_{12}z_{13}z_{23}|^2\), leaving
  \(\mathcal A_3=(2\pi)^d\delta^{(d)}(\sum_i k_i)C_{S^2}g_c^3\), or
  \(C_{S^2}g_c^3\) with the delta function suppressed. The line below
  (16.24) has a factor-of-two typo: the correct condition is
  \(\frac{\alpha'}2k_i\cdot k_j=-1\), not \(-2\).
\item
  \textbf{Question 3:} On the upper half-plane, the D\(p\)-brane gluing
  matrix \(D^\mu{}_\nu=+1\) in Neumann directions and \(-1\) in
  Dirichlet directions reflects a right mover into a holomorphic image
  field in the lower half-plane. Applying the ordinary plane propagators
  to the doubled fields gives
  \(\langle X^\mu(z)\bar X^\nu(\bar w)\rangle=-\frac{\alpha'}2D^{\mu\nu}\log(z-\bar w)\)
  and
  \(\langle\psi^\mu(z)\bar\psi^\nu(\bar w)\rangle=D^{\mu\nu}/(z-\bar w)\),
  with the Dirichlet sign entirely encoded by \(D\).
\item
  \textbf{Question 4:} The sigma-model couplings \(G_{\mu\nu}(X)\),
  \(B_{\mu\nu}(X)\), and \(\Phi(X)\) are spacetime fields, and their
  Weyl-anomaly coefficients measure whether the corresponding background
  defines a consistent worldsheet CFT. The spacetime string action is
  constructed from the same worldsheet theory, and its variational
  derivatives are invertible field-space combinations of these
  coefficients. Explicitly at leading order, varying (14.8) reproduces
  (14.5), so \(\beta^G=\beta^B=\beta^\Phi=0\) is equivalent to the
  spacetime equations of motion, modulo gauge transformations, field
  redefinitions, and matched orders in \(\alpha'\) and \(g_s\).
\end{itemize}

\section{Type-II Compactifications with D-Branes and
Fluxes}\label{type-ii-compactifications-with-d-branes-and-fluxes}

\begin{itemize}
\tightlist
\item
  \textbf{Question 1:} An infinitesimal displacement of an embedded
  D6-brane cycle is kinematically a section of its normal bundle, hence
  an element of \(H^0(\Sigma,N_\Sigma)=\Gamma(N_\Sigma)\). For a special
  Lagrangian cycle, however, a displacement is a massless supersymmetric
  modulus only if it preserves both \(\omega|_\Sigma=0\) and
  \(\operatorname{Im}(e^{-i\theta}\Omega)|_\Sigma=0\). The map
  \(v\mapsto-\iota_v\omega|_\Sigma\) converts the linearized conditions
  into \(d\alpha=d^\dagger\alpha=0\), so McLean's theorem identifies the
  deformation space with \(H^1(\Sigma,\mathbb R)\) and gives
  \(b_1(\Sigma)\) real geometric moduli. A continuous Wilson-line
  fluctuation is likewise a closed one-form modulo exact gauge
  transformations, and is therefore counted by
  \(H^1(\Sigma,\mathbb R)\). The two sets of \(b_1(\Sigma)\) real modes
  combine into \(b_1(\Sigma)\) complex adjoint scalars.
\item
  \textbf{Question 2:} At tree level, the type-IIB Kähler-modulus Kähler
  potential obeys the no-scale identity
  \(K^{\alpha\bar\beta}K_\alpha K_{\bar\beta}=3\) in Planck units. Since
  the flux superpotential is independent of \(T_\alpha\), one has
  \(D_\alpha W=K_\alpha W\), and its contribution cancels the universal
  \(-3|W|^2\) term in the \(\mathcal N=1\) F-term potential. The
  remaining potential is a positive sum of the complex-structure and
  axiodilaton F-terms, so it vanishes only when those F-terms vanish. If
  also \(W=0\), all F-terms vanish and the vacuum is supersymmetric
  Minkowski; if \(W\neq0\), the vacuum can still have zero tree-level
  energy but has \(D_\alpha W\neq0\) and broken supersymmetry. The
  Kähler moduli remain flat, and \(\alpha'\) corrections, string loops,
  nonperturbative \(T_\alpha\) dependence, or nonzero D-terms generally
  spoil the zero-energy conclusion.
\end{itemize}

\section{String Dualities and
M-Theory}\label{string-dualities-and-m-theory}

\begin{itemize}
\tightlist
\item
  \textbf{Question 1:} Writing the full dilaton as \(\Phi=\Phi_0+\phi\)
  with \(g_s=e^{\Phi_0}\), the three metrics are related by
  \(G^{(E)}_{MN}=e^{-\Phi/2}G^{(S)}_{MN}\) and
  \(g^{(mE)}_{MN}=e^{-\phi/2}G^{(S)}_{MN}=\sqrt{g_s}\,G^{(E)}_{MN}\).
  The string frame is natural for the world-sheet, loop counting, and
  T-duality; the ordinary Einstein frame gives a canonical
  Einstein--Hilbert term and makes type-IIB \(SL(2)\) duality manifest;
  the modified Einstein frame retains the canonical bulk action while
  normalizing the asymptotic metric to \(\eta_{MN}\), which is why it is
  used for the brane solutions in Section 18.5 and in equation (18.26).
\item
  \textbf{Question 2:} For the warped \(p\)-brane ansatz, a single brane
  projector removes half of the spinors. After imposing it, the
  world-volume and transverse gravitino equations reduce to first-order
  BPS equations \(A'=-2\tilde n\,u/[\Delta(d-2)]\) and
  \(B'=2n\,u/[\Delta(d-2)]\) for one common flux function \(u(r)\).
  Their weighted sum vanishes, so \(nA+\tilde nB\) is constant;
  asymptotic Minkowski normalization sets it to zero. For a D\(p\)-brane
  this follows explicitly from \(B_S'=-A_S'\) and \(\phi'=(p-3)A_S'\) in
  string frame, followed by \(A=A_S-\phi/4\) and \(B=B_S-\phi/4\).
\item
  \textbf{Question 3:} In type-IIB spinor-doublet notation the D5 and
  NS5 conditions are the commuting, independent projectors
  \(K_D=\Gamma^{012345}\sigma_1\) and \(K_N=\Gamma^{012346}\sigma_3\).
  Each halves the 32 supercharges, so their intersection preserves
  \(32/4=8\). For a \((p,q)\) five-brane, the \(SL(2)/SO(2)\) scalar
  vielbein dresses the integer charge vector into the local
  central-charge matrix
  \(\mathcal Z_5=T_{\mathrm{D5}}[\operatorname{Re}(p+\tau q)\sigma_1+\operatorname{Im}(p+\tau q)\sigma_3]\);
  dividing by the BPS tension gives
  \(\widehat{\mathcal Z}_5=\cos\alpha\,\sigma_1+\sin\alpha\,\sigma_3\).
  Saturation of the BPS bound makes the supercharge anticommutator
  vanish on preserved spinors and therefore forces
  \(K_{(p,q)}=\Gamma_{(5)}\Gamma_\theta\widehat{\mathcal Z}_5\). It
  preserves the same eight supercharges precisely when
  \(\theta=\alpha=\arg(p+\tau q)\). Normalized coordinates in the real
  basis \((1,\tau)\) convert every leg to
  \(\delta\tilde x^5+i\delta\tilde x^6\propto p+iq\); this is the web
  convention \(\tau=i\), not generally an \(SL(2,\mathbb Z)\)
  equivalence of the full string background.
\item
  \textbf{Question 4:} An oriented type-II D\(p\)-brane carries the
  dimensional reduction of ten-dimensional \(\mathcal N=1\)
  super-Yang--Mills: a world-volume gauge field \(A_\alpha\), \(9-p\)
  transverse scalars \(\Phi^i\), and a GSO-projected gaugino. For a
  single type-I D1, \(\Omega\) removes \(A_{0,1}\) and one
  two-dimensional chirality of the D1--D1 Ramond ground state, leaving
  eight transverse bosons \(X^i\) and eight right-moving Green--Schwarz
  fermions \(\psi^A\). The eight mixed directions make the D1--D9 NS
  zero-point energy positive, but its Ramond ground state is massless;
  the GSO and \(\Omega\) projections leave 32 real left-moving fermions
  \(\lambda^a\). These are exactly the Green--Schwarz world-sheet fields
  of the \(\operatorname{Spin}(32)/\mathbb Z_2\) heterotic string.
  Summing over the four flat \(O(1)\cong\mathbb Z_2\) bundles on a
  toroidal D1 world-sheet supplies the heterotic GSO projection.
\item
  \textbf{Question 5:} In the eleven-to-ten-dimensional Kaluza--Klein
  reduction, the metric exponents are fixed by demanding the type-IIA
  string-frame couplings: for
  \(ds_{11}^2=e^{2a\Phi}ds_{10,S}^2+e^{2b\Phi}(dx^{11}+C_1)^2\), the
  NS--NS Ricci term requires \(8a+b=-2\), while the R--R \(F_2=dC_1\)
  term must have no \(e^{-2\Phi}\) prefactor and requires \(2a+b=0\).
  Hence \(a=-1/3\) and \(b=2/3\). Explicitly,
  \(\sqrt{-G}=e^{-8\Phi/3}\sqrt{-g}\) and
  \(R_{11}=e^{2\Phi/3}[R_{10}+\frac{14}{3}\Box\Phi-\frac{16}{3}(\partial\Phi)^2-\frac14e^{2\Phi}F_2^2]\).
  After integration by parts this gives
  \(\sqrt{-g}\{e^{-2\Phi}[R_{10}+4(\partial\Phi)^2]-\frac14F_2^2\}\), so
  the gravitational prefactor relation does not imply equality of the
  local curvature densities.
\item
  \textbf{Question 6:} For the Green--Schwarz M2-brane, the
  supersymmetric pullback makes the kinetic term strictly invariant, but
  the Wess--Zumino potential is invariant only up to an exact form. The
  resulting Noether-current improvement contributes a term proportional
  to
  \(T_{\mathrm{M2}}\epsilon^{rs}\partial_rX^M\partial_sX^N\Gamma_{MN}\Theta\)
  to the supercharge. Its Dirac bracket with the fermionic momentum
  gives the two-form extension
  \(Z^{MN}=T_{\mathrm{M2}}\int_{\Sigma_2}dX^M\wedge dX^N\), while the
  ordinary term gives \(P_M\); in the book's conventions this is exactly
  the displayed M-theory algebra. The M5 Wess--Zumino term analogously
  supplies the five-form charge.
\item
  \textbf{Question 7:} On the punctured base
  \(B^\circ=\mathbb P^1\setminus\{z_i\}\), the loops around all 24
  seven-branes obey \(\gamma_1\cdots\gamma_{24}=1\), so their ordered
  monodromies must satisfy \(K_{24}\cdots K_1=\mathbf1\). This is the
  non-Abelian \(SL(2,\mathbb Z)\) form of the integrated axion Bianchi
  identity and hence of seven-brane tadpole cancellation. Since a D7 has
  \(K_A=T\), 24 D7-branes would give \(T^{24}\neq\mathbf1\), so some
  mutually nonlocal \((p,q)\) branes are necessary. The sharper number
  18 does not follow from this product equation alone: it counts the 18
  independent complex deformations/vector multiplets of elliptic K3,
  obtained from \(9+13-3-1=18\), and therefore bounds the independent
  brane gauge rank or independently movable mutually local branes. Read
  literally as a bound on D7 constituents, the statement is too strong;
  an \(I_{19}\) fiber has monodromy \(T^{19}\) but only rank-18 algebra
  \(A_{18}\).
\item
  \textbf{Question 8:} At the \(SO(8)^4\) orientifold point, the open D7
  sector has rank \(4\times4=16\), while four closed-string vectors
  survive from \(B_{\mu a}\) and \(C_{\mu a}\), \(a=1,2\). Their
  survival follows from the combined projection: \(B_{\mu a}\) is odd
  under both \(\Omega\) and \(I_2\), while \(C_{\mu a}\) is even under
  \(\Omega\) and odd under both \((-1)^{F_L}\) and \(I_2\), so both
  products are even. Hence the full gauge group is
  \(SO(8)^4\times U(1)^4\) of rank 20. For the heterotic string on
  \(T^2\), generic Wilson lines break the ten-dimensional rank-16 group
  to \(U(1)^{16}\); reduction of \(G_{\mu a}\) and \(B_{\mu a}\)
  supplies two Kaluza--Klein and two winding vectors, giving
  \(U(1)^{20}\). Equivalently, the Narain lattice \(\Gamma^{18,2}\)
  gives \(U(1)^{18}_L\times U(1)^2_R\), with extra ADE roots becoming
  massless only at special moduli.
\item
  \textbf{Question 9:} An \(ab\) open string has one endpoint on each
  D-brane, so its mixed boundary conditions force its internal
  center-of-mass zero mode to lie on the common intersection; it can
  propagate only along the directions shared by both branes. At an
  intersection its classical stretching length vanishes, the Ramond
  ground state is massless, and supersymmetry supplies the corresponding
  NS boson. For stacks of \(N_a\) and \(N_b\) branes these states
  transform in \((\mathbf N_a,\overline{\mathbf N}_b)\), with
  multiplicity and chirality fixed by the oriented intersection number.
  In F-theory this becomes codimension-two fiber enhancement: additional
  fibral curves collapse only over the brane-intersection locus,
  producing localized matter states, while \(G_4\) flux determines the
  net four-dimensional chirality.
\item
  \textbf{Question 10:} The inequalities in the question are reversed
  relative to page 734: both descriptions require \(g_s\ll1\); the
  flat-space perturbative D-brane regime has \(g_sN\ll1\), whereas
  classical supergravity requires \(g_sN\gg1\). Since
  \(\kappa_{10}^2N\tau_p\sim g_sN\alpha'^{(7-p)/2}\), the metric
  perturbation at the string scale is controlled by \(g_sN\). For
  D3-branes, \(L^4=4\pi g_sN\alpha'^2\), so
  \(\alpha'/L^2=(4\pi g_sN)^{-1/2}\) and \(\alpha'\) corrections are
  small precisely when \(g_sN\gg1\). Combining this with \(g_s\ll1\)
  requires large \(N\) and produces a strongly warped but gently curved
  geometry.
\item
  \textbf{Question 11:} The ambient null cone in \(\mathbb R^{d,2}\) has
  one redundant scale. On the patch \(v\neq0\), fixing that scale by
  \(v=1\) gives \(u=x^2\) and the embedding
  \(Y^M=(x^\mu,(x^2-1)/2,(x^2+1)/2)\), whose induced metric is exactly
  the Minkowski metric. Fixing the scale globally with the Euclidean
  sphere gives \(S^1\times S^{d-1}\); the projective compactification is
  more precisely \((S^1\times S^{d-1})/\mathbb Z_2\), while the product
  is its double cover. Ordinary Minkowski space is the open conformal
  diamond \(v\neq0\), with \(v=0\) becoming null, timelike, and spatial
  infinity at finite locations on the Einstein cylinder.
\item
  \textbf{Question 12:} Finite-temperature \(\mathcal N=4\) \(SU(N)\)
  SYM is defined by the thermal density matrix \(e^{-\beta H}\);
  antiperiodic fermions break supersymmetry, but flat-space conformal
  invariance still implies \(\epsilon=3p\). Homogeneity and isotropy
  give the perfect-fluid form
  \(\langle T_{\mu\nu}\rangle=\operatorname{diag}(\epsilon,p,p,p)\). At
  large \(N\) and strong 't Hooft coupling, the dual AdS\(_5\) black
  brane has \(T=U_0/\pi\) and entropy density \(s=\pi^2N^2T^3/2\), so
  \(p=\pi^2N^2T^4/8\) and \(\epsilon=3\pi^2N^2T^4/8\). This is \(3/4\)
  of the large-\(N\) free-theory energy density.
\end{itemize}

\chapter{Bosonic One-Loop
Amplitudes}\label{bosonic-one-loop-amplitudes-1}

\section{Q1: Integrating over the torus
modulus}\label{q1-integrating-over-the-torus-modulus}

Consider the closed bosonic string one-loop partition function
\(A_{0}^{(g=1)}\). We denote the torus moduli as
\(\tau=\tau_{1}+i\tau_{2}\) and the fundamental region as
\(\mathcal{F}\). From the path integral formalism we can find the torus
partition function up to numerical constants \[
A_{0}^{(g=1)}\sim \frac{V_{26}}{\alpha'^{13}}\int_{\mathcal{F}} \frac{d^2\tau}{\tau_{2}^2} \tau_{2}^{12}\int [\mathcal{D'}X\mathcal{D'}c\mathcal{D'}b ]e^{-S[X,\tau]-S[b,c,\tau]}
\] Where \(\mathcal{D'}\) means that we integrate only over the non-zero
modes. After evaluating these integrals we can obtain the answer \[
A_{0}^{(g=1)}= \frac{V_{26}}{l_{s}^{26}}\int_{\mathcal{F}} \frac{d^2\tau}{4\tau_{2}^2}Z(\tau,\bar{\tau}) \ , \ Z(\tau ,\bar{\tau})=\tau_{2}^{-12}|\eta(\tau)|^{-48}
\]

While in the Hamiltonian formalism this partition function can be
written as \[
A_{0}^{(g=1)}=\int_{\mathcal{F}} \frac{d^2\tau}{4\tau_{2}^2} \mathrm{Tr}_{\mathcal{H}}(q^{L_{0}-c/24}\bar{q}^{\bar{L}_{0}-\bar{c}/24}) \ , \ q=e^{2\pi i\tau}
\] I can understand the factor in the trace as the partition function of
the closed bosonic string on a fixed torus, but how to understand the
integral over \(\tau\)? Why we have the factor
\(\frac{d^2\tau}{4\tau_{2}^2}\)? I can understand that we should
integrate over \(d^2\tau\), but why it will appear as the modular
invariant measure \(\frac{d^2\tau}{\tau_{2}^2}\)? I may regard this as
for different \(\tau\) the weight for the corresponding partition
function is different and they are weighted by \(\frac{1}{\tau_{2}^2}\),
why?

\begin{quote}
\textbf{Textbook location and related material.} Chapter 6, especially
Sections 6.2--6.3 (pp.~126--152): gauge fixing of the genus-one Polyakov
integral, Teichmüller modulus \(\tau\), modular identifications, and the
torus partition function.
\end{quote}

\subsection{Answer}\label{answer}

\subsubsection{Core point}\label{core-point}

The trace \[
\mathrm{Tr}_{\mathcal H}\left(q^{L_0-c/24}\bar q^{\bar L_0-\bar c/24}\right)
\] is the CFT partition function on one fixed complex torus of modulus
\(\tau\). In string perturbation theory, however, the worldsheet metric
is dynamical. After fixing diffeomorphisms and Weyl transformations,
most metric fluctuations are gauge, but the complex structure modulus
\(\tau\) is not gauge. Therefore the one-loop vacuum amplitude must
still integrate over inequivalent complex structures of the torus.

The quotient by large diffeomorphisms identifies \[
\tau \sim \frac{a\tau+b}{c\tau+d},
\qquad
\begin{pmatrix}
a & b\\
c & d
\end{pmatrix}
\in SL(2,\mathbb Z),
\] so the integral is over the modular fundamental region
\(\mathcal F\), not over the whole upper half-plane.

The factor \[
\frac{d^2\tau}{\tau_2^2}
\] is not an extra physical Boltzmann weight saying that some tori are
preferred. It is the natural volume form on the moduli space of flat
tori, induced by the metric on the space of worldsheet metrics after
quotienting by diffeomorphisms and Weyl transformations. Equivalently,
it is the Faddeev-Popov measure left over after conformal gauge fixing.
The numerical factor \(1/4\) in your convention is an overall
normalization convention coming from the normalization of zero modes,
residual conformal Killing symmetries, and the chosen definition of
\(d^2\tau\). The important universal dependence is the modular invariant
Poincare measure \(d^2\tau/\tau_2^2\).

\subsubsection{Derivation}\label{derivation}

Start from the Polyakov path integral for the genus-one vacuum
amplitude, \[
A_0^{(g=1)}
\sim
\int \frac{\mathcal D g\,\mathcal D X}{\mathrm{Vol}(\mathrm{Diff}\times \mathrm{Weyl})}
e^{-S_P[X,g]}.
\] For the critical bosonic string, the matter plus ghost central charge
vanishes, \[
c_X+c_{bc}=26-26=0,
\] so Weyl transformations are genuine gauge transformations at the
quantum level. Thus one can decompose a genus-one metric into \[
g_{ab}
=
e^{2\omega}
f^*\hat g_{ab}(\tau),
\] where \(f\) is a diffeomorphism, \(\omega\) is a Weyl factor, and
\(\hat g_{ab}(\tau)\) is a representative flat metric with complex
structure \(\tau\).

The crucial point is that \(\tau\) is not removable by small
diffeomorphisms or Weyl transformations. It labels a genuine modulus.
For a torus, there is one complex modulus, hence two real moduli
\(\tau_1,\tau_2\). Therefore after gauge fixing we schematically get \[
\frac{\mathcal D g}{\mathrm{Vol}(\mathrm{Diff}\times \mathrm{Weyl})}
\longrightarrow
d\mu_{\mathcal M_1}(\tau),
\] where \(d\mu_{\mathcal M_1}\) is the induced measure on the moduli
space of complex tori.

Now let us derive the \(\tau_2^{-2}\) dependence.

Use coordinates \(\sigma^1,\sigma^2\in[0,1]\) on the torus and choose
the unit-area flat metric \[
ds^2
=
\hat g_{ab}(\tau)d\sigma^a d\sigma^b
=
\frac{1}{\tau_2}\left|d\sigma^1+\tau d\sigma^2\right|^2.
\] In components, \[
\hat g_{ab}(\tau)
=
\frac{1}{\tau_2}
\begin{pmatrix}
1 & \tau_1\\
\tau_1 & |\tau|^2
\end{pmatrix}.
\] This metric has determinant \[
\det \hat g
=
\frac{1}{\tau_2^2}\left(|\tau|^2-\tau_1^2\right)
=
\frac{1}{\tau_2^2}\tau_2^2
=1,
\] so the coordinate area is normalized to one. The overall area is not
a physical modulus in the critical string because it can be removed by a
Weyl transformation. The remaining data is the shape of the torus,
namely \(\tau\).

The space of metrics has a natural ultralocal metric, often called the
DeWitt metric. For variations \(\delta g_{ab}\), take for definiteness
\[
\|\delta g\|^2
=
\int d^2\sigma\,\sqrt g\,
g^{ac}g^{bd}\delta g_{ab}\delta g_{cd}.
\] Different authors insert harmless numerical factors such as \(1/2\)
in this definition; these change only the overall normalization, not the
\(\tau_2\) dependence.

Restrict this metric on the full space of metrics to the slice of flat
unit-area representatives \(\hat g_{ab}(\tau)\). A variation of the
modulus gives \[
\delta \hat g_{ab}
=
\frac{\partial \hat g_{ab}}{\partial \tau_1}\delta\tau_1
+
\frac{\partial \hat g_{ab}}{\partial \tau_2}\delta\tau_2.
\] Because the upper half-plane metric we obtain will be
\(SL(2,\mathbb R)\) invariant, it is enough to evaluate the local form
at \(\tau_1=0\), and then rewrite the answer covariantly. At
\(\tau_1=0\), the metric is \[
\hat g_{ab}
=
\begin{pmatrix}
1/\tau_2 & 0\\
0 & \tau_2
\end{pmatrix},
\qquad
\hat g^{ab}
=
\begin{pmatrix}
\tau_2 & 0\\
0 & 1/\tau_2
\end{pmatrix}.
\] The two independent variations are \[
\frac{\partial \hat g}{\partial \tau_1}
=
\begin{pmatrix}
0 & 1/\tau_2\\
1/\tau_2 & 0
\end{pmatrix},
\] and \[
\frac{\partial \hat g}{\partial \tau_2}
=
\begin{pmatrix}
-1/\tau_2^2 & 0\\
0 & 1
\end{pmatrix}.
\] Now compute their norms using \[
G_{ij}
=
\int d^2\sigma\,\sqrt{\hat g}\,
\hat g^{ac}\hat g^{bd}
\frac{\partial \hat g_{ab}}{\partial m^i}
\frac{\partial \hat g_{cd}}{\partial m^j},
\qquad
m^i=(\tau_1,\tau_2).
\] Since \(\sqrt{\hat g}=1\) and the coordinate volume is one, the
integral only contributes a factor \(1\). For the \(\tau_1\) variation,
\[
G_{\tau_1\tau_1}
=
2\hat g^{11}\hat g^{22}
\left(\frac{1}{\tau_2}\right)^2
=
2\left(\tau_2\right)\left(\frac{1}{\tau_2}\right)\frac{1}{\tau_2^2}
=
\frac{2}{\tau_2^2}.
\] For the \(\tau_2\) variation, \[
G_{\tau_2\tau_2}
=
(\hat g^{11})^2\left(-\frac{1}{\tau_2^2}\right)^2
+
(\hat g^{22})^2(1)^2
=
\tau_2^2\frac{1}{\tau_2^4}
+
\frac{1}{\tau_2^2}
=
\frac{2}{\tau_2^2}.
\] The mixed term vanishes: \[
G_{\tau_1\tau_2}=0.
\] Therefore the induced line element on the torus moduli is \[
ds^2_{\mathrm{moduli}}
=
\frac{2}{\tau_2^2}
\left(d\tau_1^2+d\tau_2^2\right)
=
\frac{2\,d\tau d\bar\tau}{\tau_2^2}.
\] Up to the conventional factor \(2\), this is the standard Poincare
metric on the upper half-plane. Its volume element is \[
d\mu_{\mathcal T_1}
=
\sqrt{\det G}\,d\tau_1d\tau_2
=
\frac{2\,d\tau_1d\tau_2}{\tau_2^2}
\propto
\frac{d^2\tau}{\tau_2^2}.
\] Thus the factor \(\tau_2^{-2}\) is simply the Jacobian obtained by
using \(\tau_1,\tau_2\) as coordinates on the space of inequivalent flat
metrics.

This is the geometric reason. The BRST/Faddeev-Popov derivation gives
the same result in a slightly different language. In conformal gauge,
metric variations split as \[
\delta g_{ab}
=
2\delta\omega\,g_{ab}
+
\nabla_a v_b+\nabla_b v_a
+
\delta\tau\,\frac{\partial \hat g_{ab}}{\partial \tau}
+
\delta\bar\tau\,\frac{\partial \hat g_{ab}}{\partial \bar\tau}.
\] The first term is Weyl, the second is diffeomorphism, and the final
two terms are true moduli. Gauge fixing produces the \(bc\) ghost
determinant. Since the torus has one complex modulus, the ghost path
integral contains one holomorphic and one antiholomorphic antighost
insertion schematically of the form \[
(b,\mu_\tau)(\bar b,\bar\mu_{\bar\tau}),
\] where \(\mu_\tau\) is the Beltrami differential representing the
change of complex structure. Thus the genus-one amplitude has the
schematic form \[
A_0^{(g=1)}
\sim
\int_{\mathcal M_1} d^2\tau\,
\left\langle
(b,\mu_\tau)(\bar b,\bar\mu_{\bar\tau})
\right\rangle_{\tau}
Z_X(\tau,\bar\tau).
\] When one evaluates the normalization of the modulus zero modes and
ghost zero modes with the standard torus metric, this ghost/moduli
measure gives precisely the same Poincare density \[
d\mu_{\mathcal M_1}
\propto
\frac{d^2\tau}{\tau_2^2}.
\]

Now let us connect this to the Hamiltonian trace. The fixed-torus
partition function can be written as \[
Z(\tau,\bar\tau)
=
\mathrm{Tr}_{\mathcal H}
\left(
q^{L_0-c/24}
\bar q^{\bar L_0-\bar c/24}
\right),
\qquad
q=e^{2\pi i\tau}.
\] Equivalently, \[
q^{L_0-c/24}\bar q^{\bar L_0-\bar c/24}
=
e^{-2\pi\tau_2\left(L_0+\bar L_0-\frac{c+\bar c}{24}\right)}
e^{2\pi i\tau_1\left(L_0-\bar L_0-\frac{c-\bar c}{24}\right)}.
\] For the bosonic string, \(c=\bar c\), so \[
q^{L_0-c/24}\bar q^{\bar L_0-\bar c/24}
=
e^{-2\pi\tau_2\left(L_0+\bar L_0-\frac{c}{12}\right)}
e^{2\pi i\tau_1\left(L_0-\bar L_0\right)}.
\] This makes the roles of \(\tau_2\) and \(\tau_1\) transparent:

\begin{itemize}
\tightlist
\item
  \(\tau_2\) is like a Euclidean proper time for propagation around one
  cycle of the torus.
\item
  \(\tau_1\) is a twist inserted before gluing the closed string back to
  itself.
\end{itemize}

If one integrates \(\tau_1\) over a full unit interval in the strip
representation, the factor \[
e^{2\pi i\tau_1(L_0-\bar L_0)}
\] projects onto level-matched states: \[
L_0-\bar L_0=0.
\] This is the Hamiltonian version of imposing invariance under
translations around the closed string. However, the full string
amplitude is not merely a trace at one value of \(\tau\) and not merely
an integral over a strip. A single torus can be cut into a spatial
circle and Euclidean time circle in many equivalent ways. These
different choices are related by large diffeomorphisms, \[
\tau\mapsto \frac{a\tau+b}{c\tau+d}.
\] Therefore we integrate once over the quotient \[
\mathcal M_1
=
\mathbb H/SL(2,\mathbb Z),
\] which is represented by the fundamental region \(\mathcal F\).

The modular invariance of the measure follows directly. Under \[
\tau'=\frac{a\tau+b}{c\tau+d},
\] one has \[
d^2\tau'
=
\frac{d^2\tau}{|c\tau+d|^4},
\qquad
\tau_2'
=
\frac{\tau_2}{|c\tau+d|^2}.
\] Therefore \[
\frac{d^2\tau'}{(\tau_2')^2}
=
\frac{d^2\tau/|c\tau+d|^4}
{\tau_2^2/|c\tau+d|^4}
=
\frac{d^2\tau}{\tau_2^2}.
\] So the measure is invariant under large diffeomorphisms. This is
exactly what we need, because \(\tau\) and \(\tau'\) describe the same
worldsheet geometry.

The factor should therefore not be interpreted as follows: \[
\text{``The string statistically prefers small or large }\tau_2\text{ with weight }1/\tau_2^2\text{.''}
\] Rather, it is like the \(r\,dr\,d\theta\) factor in polar
coordinates. The factor \(r\) does not mean that points at larger radius
are physically preferred; it is the coordinate volume element. Here
\(1/\tau_2^2\) is the coordinate volume element on the moduli space of
torus complex structures.

Putting this together, the one-loop vacuum amplitude has the form \[
A_0^{(g=1)}
=
\frac{V_{26}}{l_s^{26}}
\int_{\mathcal F}
\frac{d^2\tau}{4\tau_2^2}
Z(\tau,\bar\tau),
\] with \[
Z(\tau,\bar\tau)
=
\tau_2^{-12}|\eta(\tau)|^{-48}.
\] Here \(Z(\tau,\bar\tau)\) is the fixed-complex-structure CFT
partition function after the nonzero oscillator modes and zero-mode
Gaussian factors have been evaluated, while \[
\frac{d^2\tau}{4\tau_2^2}
\] is the gauge-fixed measure over inequivalent torus complex
structures. The factor \(1/4\) is convention-dependent, while the
\(\tau_2^{-2}\) factor is fixed by the geometry of moduli space and by
modular invariance.

\subsubsection{Clarifications}\label{clarifications}

First, the integral over \(\tau\) is not an ordinary thermal average
over different temperatures. In the trace expression, \(\tau_2\) does
look like an inverse temperature or Schwinger proper time, and
\(\tau_1\) looks like a twist. But in string theory these are worldsheet
geometric moduli, not external thermodynamic parameters. The string path
integral sums over worldsheets, so it must integrate over their
inequivalent complex structures.

Second, \(d^2\tau/\tau_2^2\) is not chosen merely because it is modular
invariant, although modular invariance is a very efficient way to
recognize it. More fundamentally, it is the measure induced by the
metric on the space of worldsheet metrics after quotienting by gauge
transformations. Modular invariance then follows because large
diffeomorphisms are gauge redundancies.

Third, the fundamental domain \(\mathcal F\) is not an arbitrary
restriction. If we integrated over all \(\mathbb H\), we would count the
same torus infinitely many times, because \(\tau\), \(\tau+1\),
\(-1/\tau\), and more generally \((a\tau+b)/(c\tau+d)\) represent the
same torus with different choices of homology basis.

Fourth, the numerical coefficient \(1/4\) should not be
over-interpreted. Different books distribute constants differently
between the normalization of the path integral, the definition of
\(l_s\), the ghost zero-mode measure, and the volume of residual
conformal Killing symmetries. What is structurally important is \[
d\mu_{\mathcal M_1}\propto \frac{d^2\tau}{\tau_2^2}.
\]

\subsubsection{Result}\label{result}

The fixed-torus Hamiltonian trace is \[
Z(\tau,\bar\tau)
=
\mathrm{Tr}_{\mathcal H}
\left(
q^{L_0-c/24}
\bar q^{\bar L_0-\bar c/24}
\right).
\] The string one-loop vacuum amplitude integrates this fixed-torus
partition function over inequivalent complex structures: \[
A_0^{(g=1)}
=
\frac{V_{26}}{l_s^{26}}
\int_{\mathcal F}
\frac{d^2\tau}{4\tau_2^2}
Z(\tau,\bar\tau).
\] The quotient region is \[
\mathcal M_1
=
\mathbb H/SL(2,\mathbb Z),
\] because large diffeomorphisms identify \[
\tau\sim \frac{a\tau+b}{c\tau+d}.
\] The measure comes from the induced metric on the moduli space of flat
tori: \[
ds^2_{\mathrm{moduli}}
\propto
\frac{d\tau d\bar\tau}{\tau_2^2},
\qquad
d\mu_{\mathcal M_1}
\propto
\frac{d^2\tau}{\tau_2^2}.
\] Equivalently, it is the Faddeev-Popov measure left over after
conformal gauge fixing. Thus the \(\tau_2^{-2}\) factor is a
moduli-space volume element, not an additional physical weighting
assigned by hand to different fixed-torus partition functions.

\section{Q2: The annulus modulus
measure}\label{q2-the-annulus-modulus-measure}

Now consider open bosonic string one-loop partition function. In the
cylinder case there is one real modular \(t\) with \(0\leq t<\infty\);
the inequivalent cylinders are described by
\(\{ (\tau,\sigma):0\leq\sigma\leq \pi , 0\leq \tau\leq 2\pi t \}\). The
cylinder one-loop partition function can be written as \[
A_{0}^{(g=1,b=2)}=\int_{0}^\infty \frac{dt}{2t} \mathrm{Tr}_{\mathcal{H}}(q^{L_{0}-c/24}) \ , \ q=e^{-2\pi t}
\] Derive the factor \(\frac{dt}{2t}\).

\begin{quote}
\textbf{Textbook location and related material.} Chapter 6, Section 6.5
(pp.~157--163): the cylinder/annulus modulus, its open-string trace, and
the open--closed channel transformation. Section 6.1 supplies the
general vacuum-amplitude normalization.
\end{quote}

\subsection{Answer}\label{answer-1}

\subsubsection{Core point}\label{core-point-1}

For the annulus/cylinder, the trace \[
\mathrm{Tr}_{\mathcal H}\left(q^{L_0-c/24}\right),
\qquad
q=e^{-2\pi t},
\] is the open-string CFT partition function on a fixed cylinder of
modulus \(t\). The parameter \(t\) measures the shape of the cylinder:
roughly, the ratio between its periodic Euclidean-time direction and its
open-string spatial width.

After conformal gauge fixing, there is one genuine real modulus \(t\).
The natural metric on this one-dimensional moduli space is logarithmic:
\[
ds_{\mathrm{moduli}}^2\propto \frac{dt^2}{t^2}.
\] Therefore the corresponding volume element is \[
d\mu_{\mathrm{annulus}}\propto \frac{dt}{t}.
\] The conventional open-string one-loop vacuum amplitude has \[
\frac{dt}{2t}
\] rather than just \(dt/t\) because the vacuum cylinder also has a
residual conformal Killing symmetry: translations around the periodic
direction. In the standard Faddeev-Popov normalization, dividing by this
residual symmetry gives the factor \(1/(2t)\) multiplying \(dt\).
Equivalently, in the Hamiltonian sewing picture, the trace describes a
cylinder with a marked cut; dividing by the freedom to move this cut
around the loop gives the \(1/t\) factor, with the remaining \(1/2\)
being the standard vacuum-diagram/automorphism normalization.

So the factor should be thought of as \[
\frac{dt}{2t}
=
\text{gauge-fixed measure over inequivalent annuli},
\] not as an extra physical probability weight assigned by hand to
different cylinders.

\subsubsection{Derivation}\label{derivation-1}

Let us describe the annulus as a Euclidean open-string worldsheet with
coordinates \[
0\leq \sigma\leq \pi,
\qquad
\tau\sim \tau+2\pi t.
\] The boundaries are at \[
\sigma=0,\qquad \sigma=\pi.
\] The coordinate \(\tau\) is periodic, so the annulus can be viewed as
an open string propagating in Euclidean worldsheet time and then being
traced over: \[
Z_{\mathrm{open}}(t)
=
\mathrm{Tr}_{\mathcal H}
\left(e^{-2\pi t(L_0-c/24)}\right)
=
\mathrm{Tr}_{\mathcal H}
\left(q^{L_0-c/24}\right),
\qquad
q=e^{-2\pi t}.
\] This is the fixed-modulus partition function. The remaining question
is the measure over \(t\).

As in the torus case, the starting point is the Polyakov path integral
\[
A_0^{(g=1,b=2)}
\sim
\int
\frac{\mathcal D g\,\mathcal D X}
{\mathrm{Vol}(\mathrm{Diff}\times\mathrm{Weyl})}
e^{-S_P[X,g]}.
\] Here the topology is genus zero with two boundaries, equivalently an
annulus. After fixing conformal gauge, the metric can be written
schematically as \[
g_{ab}
=
e^{2\omega}
f^*\hat g_{ab}(t),
\] where \(f\) is a diffeomorphism, \(\omega\) is a Weyl factor, and
\(\hat g_{ab}(t)\) is a representative metric for the annulus with
modulus \(t\).

The annulus has one real modulus. A convenient fixed-area representative
is obtained by using coordinates \[
0\leq \sigma\leq \pi,
\qquad
0\leq \theta<2\pi,
\] and writing \[
ds^2
=
\hat g_{ab}(t)dx^a dx^b
=
\frac{1}{t}d\sigma^2+t\,d\theta^2.
\] This is conformally equivalent to the flat cylinder \[
ds^2=d\sigma^2+d\tau^2,
\qquad
\tau=t\theta,
\qquad
0\leq \tau<2\pi t.
\] The fixed-area form is useful because the Weyl factor has already
been separated off. The determinant is \[
\det \hat g
=
\frac{1}{t}\cdot t
=1,
\] so the area is independent of \(t\): \[
\int_0^\pi d\sigma\int_0^{2\pi}d\theta\,\sqrt{\hat g}
=
2\pi^2.
\] Thus \(t\) is not the area; it is the shape modulus. More explicitly,
the circumference of the periodic direction is proportional to
\(\sqrt t\), while the distance between the two boundaries is
proportional to \(1/\sqrt t\), so their ratio is proportional to \(t\).

Now restrict the natural metric on the space of worldsheet metrics to
this one-parameter family. Use the same ultralocal metric as in the
torus discussion: \[
\|\delta g\|^2
=
\int d^2x\,\sqrt g\,
g^{ac}g^{bd}\delta g_{ab}\delta g_{cd}.
\] For \[
\hat g_{\sigma\sigma}=\frac{1}{t},
\qquad
\hat g_{\theta\theta}=t,
\] we have \[
\hat g^{\sigma\sigma}=t,
\qquad
\hat g^{\theta\theta}=\frac{1}{t}.
\] The variation with respect to \(t\) is \[
\delta \hat g_{\sigma\sigma}
=
-\frac{\delta t}{t^2},
\qquad
\delta \hat g_{\theta\theta}
=
\delta t.
\] Therefore \[
\hat g^{ac}\hat g^{bd}\delta \hat g_{ab}\delta \hat g_{cd}
=
(\hat g^{\sigma\sigma})^2
(\delta \hat g_{\sigma\sigma})^2
+
(\hat g^{\theta\theta})^2
(\delta \hat g_{\theta\theta})^2.
\] Substituting the explicit expressions gives \[
(\hat g^{\sigma\sigma})^2
(\delta \hat g_{\sigma\sigma})^2
=
t^2\left(\frac{\delta t}{t^2}\right)^2
=
\frac{(\delta t)^2}{t^2},
\] and \[
(\hat g^{\theta\theta})^2
(\delta \hat g_{\theta\theta})^2
=
\frac{1}{t^2}(\delta t)^2
=
\frac{(\delta t)^2}{t^2}.
\] Thus \[
\|\delta \hat g\|^2
=
\int_0^\pi d\sigma\int_0^{2\pi}d\theta\,
\left[
\frac{(\delta t)^2}{t^2}
+
\frac{(\delta t)^2}{t^2}
\right].
\] Since \(\sqrt{\hat g}=1\), this becomes \[
\|\delta \hat g\|^2
=
2(2\pi^2)\frac{(\delta t)^2}{t^2}.
\] The overall constant depends on the normalization chosen for the
metric on the space of metrics, but the \(t\)-dependence is fixed: \[
ds^2_{\mathrm{moduli}}
\propto
\frac{dt^2}{t^2}.
\] Since the moduli space is one-dimensional, the induced volume element
is \[
d\mu_{\mathrm{Teich}}
\propto
\sqrt{G_{tt}}\,dt
\propto
\frac{dt}{t}.
\] This is the annulus analogue of the torus result \[
d\mu_{\mathcal M_1}
\propto
\frac{d^2\tau}{\tau_2^2}.
\] For the torus the modulus is complex, so the Poincare metric gives a
two-dimensional volume form. For the annulus the modulus is real, so the
logarithmic metric gives the one-dimensional volume form \(dt/t\).

Now let us explain the extra factor of \(1/2\) in \[
\frac{dt}{2t}.
\] The annulus has a residual conformal Killing symmetry generated by
translations around the periodic direction: \[
\theta\mapsto \theta+\epsilon.
\] Equivalently, in the original cylinder coordinate, \[
\tau\mapsto \tau+\epsilon.
\] This symmetry remains after choosing the conformal representative of
the annulus. In a vacuum amplitude there are no vertex operators fixing
a point on the worldsheet, so this residual symmetry must still be
divided out.

In the Hamiltonian trace picture, this has a very concrete meaning. The
expression \[
\mathrm{Tr}_{\mathcal H}
\left(e^{-2\pi t(L_0-c/24)}\right)
\] is naturally obtained by cutting the annulus open at some value of
\(\tau\), evolving the open-string state for Euclidean time \(2\pi t\),
and gluing the final state back to the initial state. But the uncut
annulus does not contain a distinguished cut. Moving the cut around the
periodic direction gives the same geometric annulus. Therefore the trace
computation overcounts by the volume of translations around the loop.

Schematically, if \[
s=2\pi t
\] is the open-string proper time around the loop, then the Schwinger
representation gives an integration over \(s\), while the unmarked
cylinder requires division by the translation volume proportional to
\(s\): \[
\frac{ds}{s}
=
\frac{d(2\pi t)}{2\pi t}
=
\frac{dt}{t}.
\] With the standard normalization of the annulus vacuum diagram, or
equivalently the standard normalization of the residual conformal
Killing group in the Faddeev-Popov procedure, this becomes \[
\frac{dt}{2t}.
\] One often phrases this as \[
d\mu_{\mathrm{annulus}}
=
\frac{dt}{2t}.
\] The factor \(1/2\) is a conventional but standard normalization
associated with the automorphism/symmetry factor of the vacuum cylinder.
The important modulus dependence is the logarithmic measure \(dt/t\).

Putting the measure and fixed-cylinder partition function together gives
\[
A_0^{(g=1,b=2)}
=
\int_0^\infty
\frac{dt}{2t}
\mathrm{Tr}_{\mathcal H}
\left(q^{L_0-c/24}\right),
\qquad
q=e^{-2\pi t}.
\]

There is a useful parallel with ordinary one-loop field theory. A
one-loop vacuum amplitude written in Schwinger proper time has the
schematic form \[
\int_0^\infty \frac{dT}{T}\,
\mathrm{Tr}\left(e^{-T H}\right).
\] The factor \(dT/T\) appears because a closed loop has no preferred
starting point: translating the starting point around the loop is a
redundancy. The open-string annulus is the stringy version of this
statement, with \[
T=2\pi t,
\qquad
H=L_0-\frac{c}{24}.
\]

\subsubsection{Clarifications}\label{clarifications-1}

First, \(t\) is not an external inverse temperature chosen by hand. It
plays a role similar to Euclidean proper time in the Hamiltonian trace,
but in the string path integral it is a modulus of the worldsheet
geometry. We integrate over it because the worldsheet metric is
dynamical.

Second, the factor \(1/t\) is not a physical preference for smaller
values of \(t\). It is the coordinate volume element on the moduli space
of inequivalent annuli. A better coordinate is \[
u=\log t.
\] Then \[
\frac{dt}{t}=du.
\] So the measure is simply flat in the logarithmic modulus \(u\).

Third, unlike the torus, the annulus modulus is real rather than
complex. Therefore the measure is not \[
\frac{d^2\tau}{\tau_2^2},
\] but its one-dimensional analogue \[
\frac{dt}{t}.
\]

Fourth, the open-channel and closed-channel descriptions are related by
a change of variables, often of the form \[
t\sim \frac{1}{\ell}.
\] This does not mean that \(t\) and \(1/t\) are being identified inside
the same integration domain in the same way that \(\tau\) and
\(-1/\tau\) are identified for the torus. Rather, it is a dual
description of the same annulus amplitude as either an open-string loop
or a closed-string tree-level exchange between two boundaries.

\subsubsection{Result}\label{result-1}

Choose a fixed-area representative metric \[
ds^2
=
\frac{1}{t}d\sigma^2+t\,d\theta^2,
\qquad
0\leq \sigma\leq \pi,
\qquad
\theta\sim\theta+2\pi.
\] Then \[
\delta g_{\sigma\sigma}
=
-\frac{\delta t}{t^2},
\qquad
\delta g_{\theta\theta}
=
\delta t,
\] and the induced metric on the modulus is \[
ds^2_{\mathrm{moduli}}
\propto
\frac{dt^2}{t^2}.
\] Therefore \[
d\mu_{\mathrm{Teich}}
\propto
\frac{dt}{t}.
\] Dividing by the residual translation symmetry of the vacuum cylinder,
with the standard annulus normalization, gives \[
d\mu_{\mathrm{annulus}}
=
\frac{dt}{2t}.
\] Thus \[
A_0^{(g=1,b=2)}
=
\int_0^\infty
\frac{dt}{2t}
\mathrm{Tr}_{\mathcal H}
\left(q^{L_0-c/24}\right),
\qquad
q=e^{-2\pi t}.
\]

\section{Q3: Massless exchange between parallel
D-branes}\label{q3-massless-exchange-between-parallel-d-branes}

Consider two parallel Dp-branes in a d-dimensional space-time. In order
to compare the string computation with a field theory computation, one
needs to study the exchange of massless modes between the two branes.
D-branes are dynamical objects which are expected to have a mass-density
(tension) and they therefore couple to gravity and possibly other
massless excitations of the closed string, such as the dilaton and the
anti-symmetric tensor. It turns out that the latter does not contribute
to lowest order because there is no field-theoretical one-particle
exchange involving this field (I do not fully understand this).

To derive the one-particle exchange amplitude between two Dp-branes, for
example the graviton, we consider the following effective space-time
action (neglected fields that do not contribute) \[
S=S_{\text{brane}}+S_{\text{bulk}}
\] Where \[
S_{\text{bulk}}=\frac{1}{2\kappa_{d}^2}\int d^dx\sqrt{ -g }\left[ R-\frac{4}{d-2}\partial_{\mu}\phi\partial^\mu \phi \right] \ , \ \Phi=\Phi_{0}+\phi
\] \[
S_{\text{brane}}=-\tau_{p}\int d^{p+1}\xi \exp\left( \frac{{4-d+2p}}{d-2}\phi \right)\sqrt{ -\Gamma }
\] We choose the normalization and conventions \[
\phi=\frac{1}{2}\kappa_{d}\sqrt{ d-2 }D \ , \ g_{\mu \nu}=\eta_{\mu \nu}+2\kappa_{d}h_{\mu \nu}
\] And the static gauge \[
X^\alpha=\xi^\alpha \ , \ X^i=\text{constant} \ , \ \alpha=0,\cdots,p \ , \ i=p+1,\cdots,d-1
\] We ignored the fluctuation of the Dp-brane and the gauge field on the
Dp-brane since they do not contribute. And we ignore the B-field for
now.

Then to the first order \[
S_{\text{brane}}=\tau_{p}\kappa_{d}\int d^{p+1}\xi\left( D^{\mu \nu}h_{\mu \nu}-\left( \frac{{4-d+2p}}{2\sqrt{ d-2 }} \right)D \right) \ , \ D^{\mu \nu}=\begin{pmatrix}
\eta^{\alpha\beta} & 0 \\
0 & 0
\end{pmatrix}
\]

\begin{quote}
\textbf{Textbook location and related material.} Chapter 6, Section 6.8
(pp.~171--174): comparison of the cylinder with massless closed-string
exchange and extraction of the D-brane tension. Related spacetime
couplings appear in Section 16.5 (pp.~630--636), especially the DBI
action.
\end{quote}

\subsection{Answer}\label{answer-2}

\subsubsection{Core point}\label{core-point-2}

For a static, flat D\(p\)-brane with vanishing world-volume gauge field,
there is no source term linear in the NS--NS two-form \(B_{mu\nu}\). A
one-particle exchange diagram needs one linear source on each brane, so
the \(B\)-field exchange vanishes at this order. The graviton and
dilaton do have linear sources and therefore contribute.

\subsubsection{\texorpdfstring{Why the \(B\)-field has no linear
source}{Why the B-field has no linear source}}\label{why-the-b-field-has-no-linear-source}

The DBI action is \[
S_{\mathrm{DBI}}=-\tau_p\int d^{p+1}\xi\,e^{-\Phi}
\sqrt{-\det\!\left(P[G+B]_{\alpha\beta}+2\pi\alpha'F_{\alpha\beta}\right)}.
\] Set \(F=0\) and expand around a flat induced metric, writing the
fluctuation matrix as
\(M_{\alpha\beta}=h_{\alpha\beta}+B_{\alpha\beta}\). Since \[
\sqrt{\det(1+M)}=1+\frac12\operatorname{tr}M+O(M^2),
\] the linear term contains \(\eta^{\alpha\beta}h_{\alpha\beta}\) but \[
\eta^{\alpha\beta}B_{\alpha\beta}=0
\] because the inverse metric is symmetric and \(B_{\alpha\beta}\) is
antisymmetric. Thus \(J_B^{\mu\nu}=0\) for the background used in the
comparison. A nonzero background flux \(F\) or a more general brane
configuration can generate \(B\) couplings through the gauge-invariant
combination \(B+2\pi\alpha'F\).

\subsubsection{Graviton and dilaton
exchange}\label{graviton-and-dilaton-exchange}

With the canonically normalized fields in the question, the source terms
are \[
S_{\mathrm{source}}
=\kappa_d\tau_p\int d^{p+1}\xi
\left(D^{\mu\nu}h_{\mu\nu}-a_pD\right),
\qquad
a_p=\frac{4-d+2p}{2\sqrt{d-2}}.
\] For two static parallel branes separated by \(r\) in the \(d-p-1\)
transverse directions, the exchange amplitude is obtained by contracting
the two sources with the massless propagators and integrating only over
transverse momentum: \[
\mathcal A_{\mathrm{NSNS}}(r)
\propto
\kappa_d^2\tau_p^2 V_{p+1}
\int\frac{d^{d-p-1}k}{(2\pi)^{d-p-1}}
\frac{e^{ik\cdot r}}{k^2}
\left(\mathcal C_h+a_p^2\right).
\] Here \(\mathcal C_h\) is the contraction of \(D^{\mu\nu}\) with the
de Donder-gauge graviton projector. Both contributions are attractive
for equal positive-tension branes. In a supersymmetric theory, R--R form
exchange supplies an equal repulsive term, yielding the no-force
condition for parallel BPS branes.

\subsubsection{Result}\label{result-2}

At leading order about the static \(F=0\) background, \[
J_h\neq0,\qquad J_D\neq0,\qquad J_B=0.
\] The absence of \(B\) exchange is therefore not a statement that
D-branes never interact with \(B\); it is the statement that this
particular background has no one-\(B\) linear vertex.

\chapter{Torus Partition Functions for
RCFTs}\label{torus-partition-functions-for-rcfts-1}

\section{\texorpdfstring{Q1: Why \(S^2=C\) in
RCFT}{Q1: Why S\^{}2=C in RCFT}}\label{q1-why-s2c-in-rcft}

In RCFT, the characters transform as

\[
\chi_i\left(-\frac{1}{\tau}\right)
=
\sum_j S_{ij}\chi_j(\tau).
\]

It is said that

\[
S^2=C,
\]

Where \(C\) is the charge conjugation matrix.

How can we derive \(S^2=C\)? What is \(C\) in the context of RCFT?

\begin{quote}
\textbf{Textbook location and related material.} Chapter 6, Section 6.4
(pp.~152--157): RCFT characters, modular \(S\) and \(T\) matrices,
modular invariants, and charge conjugation. The geometric action of
\(S^2\) is the twice-applied torus cycle exchange.
\end{quote}

\subsection{Answer}\label{answer-3}

\subsubsection{Core point}\label{core-point-3}

In an RCFT, the modular matrix \(S\) is not just a matrix describing the
scalar map \(\tau \mapsto -1/\tau\). More precisely, it is the
representation of the mapping class group of the torus on the
finite-dimensional vector space spanned by the chiral characters, or
equivalently by genus-one conformal blocks. The modular transformation

\[
s =
\begin{pmatrix}
0 & -1 \\
1 & 0
\end{pmatrix}
\]

satisfies

\[
s^2 = -I
\]

inside \(SL(2,\mathbb Z)\). Although \(-I\) acts trivially on the
complex modulus \(\tau\), it does not act trivially on the labelled
chiral block: geometrically it reverses the orientation of the torus
cycles, so a primary label \(i\) is replaced by its conjugate, or dual,
label \(i^*\). Therefore the representation of \(s^2\) is the charge
conjugation matrix:

\[
S^2 = C,
\qquad
C_{ij} = \delta_{i,j^*}.
\]

Here \(C\) is the matrix which sends each primary label, or equivalently
each basis vector in the character space, to its contragredient label:

\[
C e_i = e_{i^*}.
\]

If every primary is self-conjugate, then \(i^*=i\) for all \(i\), and
\(C=I\). In that special case \(S^2=I\). But in a general RCFT, \(C\)
can be a nontrivial permutation matrix.

\subsubsection{Derivation}\label{derivation-2}

Let the RCFT have finitely many irreducible chiral representations, or
primary labels,

\[
i,j,k \in \mathcal I.
\]

For each primary label \(i\), let \(\mathcal H_i\) be the corresponding
irreducible module of the chiral algebra. Its character is

\[
\chi_i(\tau)
=
\operatorname{Tr}_{\mathcal H_i}
q^{L_0-c/24},
\qquad
q=e^{2\pi i \tau}.
\]

The defining RCFT property is that the span of these finitely many
characters is closed under modular transformations. In particular, under

\[
s:\tau \mapsto -\frac{1}{\tau},
\]

we have

\[
\chi_i\left(-\frac{1}{\tau}\right)
=
\sum_j S_{ij}\chi_j(\tau).
\]

Equivalently, if we package the characters into a column vector

\[
\chi(\tau)
=
\begin{pmatrix}
\chi_0(\tau) \\
\chi_1(\tau) \\
\vdots
\end{pmatrix},
\]

then, up to the chosen row-versus-column convention, the modular
transformation is represented by a finite matrix \(S\).

The key point is that \(S\) represents the mapping class group action,
not merely the induced action on the upper half-plane coordinate
\(\tau\). The torus mapping class group is

\[
SL(2,\mathbb Z),
\]

with standard generators

\[
s=
\begin{pmatrix}
0 & -1 \\
1 & 0
\end{pmatrix},
\qquad
t=
\begin{pmatrix}
1 & 1 \\
0 & 1
\end{pmatrix}.
\]

The generator \(s\) acts on the modulus by

\[
\tau \mapsto -\frac{1}{\tau},
\]

while \(t\) acts by

\[
\tau \mapsto \tau+1.
\]

Now compute \(s^2\) as a matrix:

\[
s^2
=
\begin{pmatrix}
0 & -1 \\
1 & 0
\end{pmatrix}
\begin{pmatrix}
0 & -1 \\
1 & 0
\end{pmatrix}
=
\begin{pmatrix}
-1 & 0 \\
0 & -1
\end{pmatrix}
=
-I.
\]

As a fractional linear transformation,

\[
-I:
\tau \mapsto
\frac{-\tau}{-1}
=
\tau.
\]

So if we only look at the numerical point \(\tau\) in the upper
half-plane, \(s^2\) appears to do nothing. This is why one might naively
expect \(S^2=I\). But this misses the labelled conformal block data.

To see what is really happening, remember that a character can be viewed
geometrically as a genus-one chiral conformal block with a label \(i\).
Equivalently, one can think of it as a torus amplitude with an
\(i\)-labelled chiral sector propagating along one cycle of the torus.

Let \(a\) and \(b\) denote a basis of one-cycles of the torus. The
modular transformation \(s\) exchanges the two cycles, up to
orientation:

\[
s:
(a,b)
\mapsto
(b,-a),
\]

with convention-dependent signs. Applying \(s\) twice gives

\[
s^2:
(a,b)
\mapsto
(-a,-b).
\]

Thus \(s^2\) reverses the orientation of both fundamental cycles.

Now the important RCFT fact is:

\begin{quote}
Reversing the orientation of a labelled chiral line replaces the label
by its dual, or charge-conjugate, label.
\end{quote}

If a primary representation is \(\mathcal H_i\), its dual or
contragredient representation is usually denoted

\[
\mathcal H_{i^*}.
\]

Therefore, under \(s^2\), the block labelled by \(i\) is mapped to the
block labelled by \(i^*\):

\[
s^2:
i
\mapsto
i^*.
\]

Passing to the matrix representation on characters, this says

\[
\rho(s)^2
=
\rho(s^2)
=
\rho(-I)
=
C,
\]

where \(\rho(s)=S\). Hence

\[
S^2=C.
\]

This is the clean geometric derivation. The central element
\(-I \in SL(2,\mathbb Z)\) acts trivially on the complex number
\(\tau\), but nontrivially on the primary label: it sends a
representation to its conjugate representation.

There is also a useful algebraic derivation using standard properties of
the modular \(S\)-matrix. In a unitary RCFT, the \(S\)-matrix is unitary
and symmetric:

\[
S^\dagger S = I,
\qquad
S_{ij}=S_{ji}.
\]

Moreover, charge conjugation of a primary corresponds to complex
conjugation of the \(S\)-matrix row:

\[
S_{i^*j}
=
S_{ij}^*.
\]

This relation is natural because the \(S\)-matrix can be interpreted as
the invariant associated to a Hopf link whose components are labelled by
primaries. Reversing the orientation of a labelled component replaces
the label by the dual representation, and in a unitary theory this
complex-conjugates the amplitude.

Now compute the matrix square:

\[
(S^2)_{ij}
=
\sum_k S_{ik}S_{kj}.
\]

Using symmetry, \(S_{kj}=S_{jk}\), so

\[
(S^2)_{ij}
=
\sum_k S_{ik}S_{jk}.
\]

Using

\[
S_{j^*k}=S_{jk}^*,
\]

we can rewrite

\[
S_{jk}=S_{j^*k}^*.
\]

Therefore

\[
(S^2)_{ij}
=
\sum_k S_{ik}S_{j^*k}^*.
\]

But unitarity says that the rows of \(S\) are orthonormal:

\[
\sum_k S_{ik}S_{\ell k}^*
=
\delta_{i\ell}.
\]

Taking \(\ell=j^*\) gives

\[
(S^2)_{ij}
=
\delta_{i,j^*}.
\]

Thus

\[
(S^2)_{ij}
=
C_{ij},
\]

where

\[
C_{ij}
=
\delta_{i,j^*}.
\]

So again,

\[
S^2=C.
\]

This algebraic derivation is often the most efficient way to check the
statement once the standard RCFT properties of \(S\) are known. The
geometric derivation explains why the answer is charge conjugation
rather than the identity.

Let us also connect this to the torus partition function. A full
nonchiral CFT has left-moving and right-moving sectors. A common modular
invariant, often called the charge-conjugation modular invariant, is

\[
Z_C(\tau,\bar\tau)
=
\sum_{i,j} C_{ij}\chi_i(\tau)\overline{\chi_j(\tau)}
=
\sum_i \chi_i(\tau)\overline{\chi_{i^*}(\tau)}.
\]

This pairs the left-moving representation \(i\) with the right-moving
representation \(i^*\). In many familiar diagonal theories the
characters of \(i\) and \(i^*\) are identical as \(q\)-series, so this
is often written schematically as

\[
Z(\tau,\bar\tau)
=
\sum_i |\chi_i(\tau)|^2.
\]

But conceptually the pairing is really between a representation and its
conjugate. This is why the same matrix \(C\) appears both in the modular
relation

\[
S^2=C
\]

and in the canonical charge-conjugation modular invariant.

\subsubsection{Explanation of notation}\label{explanation-of-notation}

\begin{itemize}
\item
  \(\mathcal I\) is the finite set of primary labels of the RCFT.
\item
  \(\mathcal H_i\) is the irreducible chiral module associated with the
  primary label \(i\).
\item
  \(\chi_i(\tau)\) is the character of \(\mathcal H_i\):
\end{itemize}

\[
\chi_i(\tau)
=
\operatorname{Tr}_{\mathcal H_i}
q^{L_0-c/24}.
\]

\begin{itemize}
\tightlist
\item
  \(S\) is the modular matrix representing
\end{itemize}

\[
\tau \mapsto -\frac{1}{\tau}
\]

on the finite-dimensional space spanned by the RCFT characters.

\begin{itemize}
\tightlist
\item
  \(i^*\) is the conjugate, dual, or contragredient primary label. It is
  defined by
\end{itemize}

\[
\mathcal H_{i^*}
\cong
\mathcal H_i^\vee.
\]

Physically, \(i^*\) is the representation carried by the oppositely
oriented chiral line. In WZW models, it is the conjugate
finite-dimensional Lie algebra representation. For example, in an
\(SU(3)\) WZW model, the fundamental representation \(\mathbf 3\) and
antifundamental representation \(\overline{\mathbf 3}\) are conjugate.

\begin{itemize}
\tightlist
\item
  \(C\) is the charge conjugation matrix:
\end{itemize}

\[
C_{ij}
=
\delta_{i,j^*}.
\]

Equivalently,

\[
C e_j = e_{j^*},
\]

where \(e_j\) is the basis vector corresponding to the primary label
\(j\). Since conjugation is an involution,

\[
(i^*)^*=i,
\]

we have

\[
C^2=I.
\]

Also, because \(C\) is a permutation matrix,

\[
C^{-1}=C.
\]

Depending on whether one writes characters as row vectors or column
vectors, one may write the action as either

\[
(C\chi)_i=\chi_{i^*}
\]

or

\[
(\chi C)_i=\chi_{i^*}.
\]

The invariant content is the same: \(C\) permutes primary labels by
charge conjugation.

\subsubsection{Clarifications}\label{clarifications-2}

\textbf{Confusion 1: Since applying \(\tau\mapsto -1/\tau\) twice gives
\(\tau\), shouldn't \(S^2=I\)?}

Only if one forgets that the modular transformation acts on labelled
conformal blocks, not just on the complex number \(\tau\). In
\(PSL(2,\mathbb Z)\), where matrices differing by \(\pm I\) are
identified, the element \(s\) has order two. But the mapping class group
relevant for chiral conformal blocks is more naturally represented
before forgetting the central element \(-I\). That central element acts
as charge conjugation:

\[
\rho(-I)=C.
\]

So the correct RCFT statement is

\[
S^2=C,
\]

not necessarily \(S^2=I\).

\textbf{Confusion 2: Does charge conjugation mean electric charge
conjugation?}

Not necessarily. In RCFT, charge conjugation means taking the dual
representation of the chiral algebra. If the chiral algebra contains a
current algebra, this often coincides with ordinary representation
conjugation, such as

\[
\mathbf 3
\leftrightarrow
\overline{\mathbf 3}
\]

for \(SU(3)\). But the abstract definition is

\[
i \mapsto i^*,
\]

where \(i^*\) labels the contragredient module.

\textbf{Confusion 3: If \(\chi_i(\tau)=\chi_{i^*}(\tau)\) as functions,
is \(C\) still meaningful?}

Yes. The equality of the \(q\)-series can hide the distinction between
labels. The modular representation is defined on the labelled vector
space of conformal blocks. Even if two conjugate characters happen to
have the same numerical expansion, the labels \(i\) and \(i^*\) can
still be conceptually distinct.

This is why one should not derive \(S^2=I\) by treating the characters
as unlabeled scalar functions. The statement

\[
S^2=C
\]

remembers the primary-label data.

\textbf{Confusion 4: Is \(C\) always nontrivial?}

No.~In many simple RCFTs all primaries are self-conjugate. Then

\[
i^*=i
\]

for every primary label \(i\), so

\[
C=I.
\]

For example, in the Ising model all three primaries are self-conjugate,
so its modular matrix satisfies \(S^2=I\). But in theories with
genuinely complex representations, \(C\) is nontrivial.

\textbf{Confusion 5: Is \(S^2=C\) an assumption or a theorem?}

It is a theorem under the standard RCFT axioms. Geometrically, it
follows from how the torus mapping class group acts on labelled
conformal blocks. Algebraically, in a unitary RCFT it follows from
symmetry, unitarity, and the conjugation property of the modular
\(S\)-matrix:

\[
S_{ij}=S_{ji},
\qquad
S^\dagger S=I,
\qquad
S_{i^*j}=S_{ij}^*.
\]

These imply

\[
(S^2)_{ij}
=
\delta_{i,j^*}.
\]

\subsubsection{Result}\label{result-3}

The modular generator

\[
s=
\begin{pmatrix}
0 & -1 \\
1 & 0
\end{pmatrix}
\]

satisfies

\[
s^2=-I
\]

in \(SL(2,\mathbb Z)\). Although \(-I\) fixes the modulus \(\tau\), it
reverses the orientation of the torus cycles and hence sends a primary
label to its conjugate:

\[
i \mapsto i^*.
\]

Therefore, on the vector space of RCFT characters,

\[
\rho(s)^2
=
\rho(-I)
=
C.
\]

Since \(\rho(s)=S\), one obtains

\[
S^2=C.
\]

The charge conjugation matrix is

\[
C_{ij}
=
\delta_{i,j^*}.
\]

Equivalently,

\[
C e_j
=
e_{j^*}.
\]

Using symmetry and unitarity of \(S\), the same statement is

\[
(S^2)_{ij}
=
\sum_k S_{ik}S_{kj}
=
\sum_k S_{ik}S_{jk}
=
\sum_k S_{ik}S_{j^*k}^*
=
\delta_{i,j^*}
=
C_{ij}.
\]

Thus \(C\) is the matrix implementing charge conjugation on primary
labels, and \(S^2=C\) means that applying the modular
\(S\)-transformation twice sends every primary sector to its conjugate
sector.

\chapter{The Classical Fermionic String and Spin
Structures}\label{the-classical-fermionic-string-and-spin-structures-1}

\section{Q1: Dirac operators on Riemann surfaces and their zero
modes}\label{q1-dirac-operators-on-riemann-surfaces-and-their-zero-modes}

Consider the Dirac operator on compact Riemann surfaces. Denote the
Riemann surface with genus \(g\) as \(\Sigma_{g}\). In 2d, a Dirac
spinor has two complex dimensions i.e.~it can be written as \[
\psi=\begin{pmatrix}
\psi_{+} \\
\psi_{-}
\end{pmatrix}
\] And we choose the flat Dirac matrices to be \[
\gamma^0=\begin{pmatrix}
0 & -i \\
i & 0
\end{pmatrix} \ , \ \gamma^1=\begin{pmatrix}
0 & 1 \\
1 & 0
\end{pmatrix} \  , \ \gamma^2=i\gamma^0\gamma^1=\begin{pmatrix}
1 & 0 \\
0 & -1
\end{pmatrix}
\] Derive the explicit form of the Dirac operator for the two
chiralities.

On \(\Sigma_{g}\) there are \(2^{2g}\) spin structures. Prove the
following statements: 1. For a given spin structure, the number of
chiral Dirac zero modes modulo two is a topological invariant; 2. The
number of chiral Dirac zero modes modulo two is additive when we glue
together two Riemann surfaces.

\begin{quote}
\textbf{Textbook location and related material.} Chapter 7, Sections
7.3--7.5 (pp.~180--194) for the two-dimensional Dirac operator and
spinors; Chapter 9, Section 9.1 (pp.~223--232) for spin structures on
the torus. The Arf-invariant formulation used in the answer is
additional mathematical background beyond the book.
\end{quote}

\subsection{Answer}\label{answer-4}

\subsubsection{Core point}\label{core-point-4}

The question has two connected parts. First, one wants the explicit
chiral form of the Dirac operator on a compact Riemann surface. Second,
one wants to understand why, for a fixed spin structure, the number of
chiral zero modes modulo two is topological and additive under gluing.

The key geometric fact is that a spin structure on a Riemann surface is
a square root of the canonical bundle: \[
L^2\simeq K.
\] Chiral Dirac zero modes are holomorphic sections of this spin bundle,
up to the convention that the opposite chirality gives the complex
conjugate anti-holomorphic equation. Thus the number of positive or
negative chiral zero modes is controlled by \[
h^0(\Sigma,L)=\dim_{\mathbb C} H^0(\Sigma,L).
\] The ordinary integer index of the chiral Dirac operator vanishes in
two dimensions for an honest spin bundle, because \[
\deg L=g-1,
\qquad
\deg L+1-g=0.
\] So the integer index does not distinguish spin structures. The subtle
invariant is instead the mod-two index \[
h^0(\Sigma,L)\bmod 2.
\] This parity is the Arf invariant of the quadratic refinement
determined by the spin structure: \[
h^0(\Sigma,L)\equiv \operatorname{Arf}(q_L)\pmod 2.
\] Since \(q_L\) is topological, the parity of chiral zero modes is
topological. Under connected-sum gluing, the quadratic refinements
combine by orthogonal direct sum, and the Arf invariant is additive.
Therefore the parity of chiral zero modes is additive.

\subsubsection{A small correction to the chirality
matrix}\label{a-small-correction-to-the-chirality-matrix}

Using the gamma matrices in the question, \[
\gamma^0=
\begin{pmatrix}
0 & -i\\
i & 0
\end{pmatrix},
\qquad
\gamma^1=
\begin{pmatrix}
0 & 1\\
1 & 0
\end{pmatrix},
\] one computes \[
\gamma^0\gamma^1
=
\begin{pmatrix}
-i & 0\\
0 & i
\end{pmatrix},
\] and therefore \[
\gamma^2=i\gamma^0\gamma^1
=
\begin{pmatrix}
1 & 0\\
0 & -1
\end{pmatrix}.
\] So I will use \[
\gamma^2=
\begin{pmatrix}
1 & 0\\
0 & -1
\end{pmatrix}.
\] With this convention, \[
\psi=
\begin{pmatrix}
\psi_+\\
\psi_-
\end{pmatrix},
\qquad
\gamma^2\psi_\pm=\pm\psi_\pm.
\] The matrix written in the question, \[
\begin{pmatrix}
1 & 0\\
-1 & 0
\end{pmatrix},
\] is not equal to \(i\gamma^0\gamma^1\) for the displayed
\(\gamma^0,\gamma^1\), and it is not the usual chirality matrix because
it is not diagonal with eigenvalues \(\pm 1\).

\subsubsection{Local form of the Dirac
operator}\label{local-form-of-the-dirac-operator}

Let us work locally in conformal coordinates \[
z=x^0+ix^1,
\qquad
\bar z=x^0-ix^1,
\] with \[
ds^2=e^{2\omega}\left((dx^0)^2+(dx^1)^2\right)
=e^{2\omega}dz\,d\bar z.
\] Choose the orthonormal coframe \[
e^0=e^\omega dx^0,
\qquad
e^1=e^\omega dx^1.
\] The curved gamma matrices are \[
\gamma^\alpha=e^\alpha{}_a\gamma^a,
\] so in these coordinates the Dirac operator is \[
\slashed D
=
\gamma^\alpha\nabla_\alpha
=
e^{-\omega}\gamma^a\nabla_a.
\] The torsion-free spin connection satisfies \[
de^a+\omega^a{}_b\wedge e^b=0.
\] Writing \(\omega^{01}=\omega_\mu{}^{01}dx^\mu\), one finds \[
\omega^{01}
=
(\partial_1\omega)\,dx^0-(\partial_0\omega)\,dx^1.
\] The spinor covariant derivative is \[
\nabla_\mu
=
\partial_\mu+\frac{1}{4}\omega_\mu{}^{ab}\gamma_a\gamma_b.
\] Since only the \(01\) component is present, this becomes \[
\nabla_\mu
=
\partial_\mu+\frac{1}{2}\omega_\mu{}^{01}\gamma_0\gamma_1.
\] Using \[
\gamma_0\gamma_1=-i\gamma^2,
\] we get \[
\nabla_0
=
\partial_0-\frac{i}{2}(\partial_1\omega)\gamma^2,
\qquad
\nabla_1
=
\partial_1+\frac{i}{2}(\partial_0\omega)\gamma^2.
\] Therefore \[
\slashed D
=
e^{-\omega}\left[
\gamma^0\left(\partial_0-\frac{i}{2}(\partial_1\omega)\gamma^2\right)
+
\gamma^1\left(\partial_1+\frac{i}{2}(\partial_0\omega)\gamma^2\right)
\right].
\] Acting on \[
\psi=
\begin{pmatrix}
\psi_+\\
\psi_-
\end{pmatrix},
\] this gives \[
\slashed D
\begin{pmatrix}
\psi_+\\
\psi_-
\end{pmatrix}
=
e^{-\omega}
\begin{pmatrix}
-2i\left(\partial_{\bar z}+\frac{1}{2}\partial_{\bar z}\omega\right)\psi_-\\
2i\left(\partial_z+\frac{1}{2}\partial_z\omega\right)\psi_+
\end{pmatrix}.
\] Equivalently, \[
\slashed D
=
2i e^{-\omega}
\begin{pmatrix}
0 & -\left(\partial_{\bar z}+\frac{1}{2}\partial_{\bar z}\omega\right)\\
\partial_z+\frac{1}{2}\partial_z\omega & 0
\end{pmatrix}.
\] Thus the chiral pieces are \[
D_+\psi_+
=
2i e^{-\omega}
\left(\partial_z+\frac{1}{2}\partial_z\omega\right)\psi_+,
\] and \[
D_-\psi_-
=
-2i e^{-\omega}
\left(\partial_{\bar z}+\frac{1}{2}\partial_{\bar z}\omega\right)\psi_-.
\] The full Dirac equation \(\slashed D\psi=0\) is therefore equivalent
to \[
\left(\partial_z+\frac{1}{2}\partial_z\omega\right)\psi_+=0,
\qquad
\left(\partial_{\bar z}+\frac{1}{2}\partial_{\bar z}\omega\right)\psi_-=0.
\] After the Weyl rescaling \[
\widehat\psi_\pm=e^{\omega/2}\psi_\pm,
\] these become simply \[
\partial_z\widehat\psi_+=0,
\qquad
\partial_{\bar z}\widehat\psi_-=0.
\] So, with the coordinate and gamma-matrix convention used here,
\(\widehat\psi_+\) is anti-holomorphic and \(\widehat\psi_-\) is
holomorphic. If one reverses the orientation convention or swaps the
labels \(\psi_+\) and \(\psi_-\), the words ``holomorphic'' and
``anti-holomorphic'' are interchanged. The invariant statement is that
the two chiral Dirac operators are the Dolbeault operator and its
conjugate, acting on the two spin bundles.

\subsubsection{Spin bundles and chiral zero
modes}\label{spin-bundles-and-chiral-zero-modes}

Globally, one should not think of \(\psi_+\) and \(\psi_-\) as ordinary
functions. They are spinors, so they are sections of spin bundles.

On a Riemann surface, a spin structure is a holomorphic line bundle
\(L\) such that \[
L^2\simeq K,
\] where \(K=T^{*(1,0)}\Sigma\) is the canonical bundle. Such an \(L\)
is also called a theta characteristic.

Depending on the chirality convention, one chirality is represented by
sections of \(L\), while the opposite chirality is represented by
sections of the conjugate bundle. With the convention above, the
holomorphic zero-mode equation is naturally \[
\bar\partial_L \widehat\psi_-=0,
\] so negative-chirality zero modes are holomorphic sections of \(L\):
\[
\widehat\psi_-\in H^0(\Sigma,L).
\] The opposite chirality satisfies the complex conjugate equation. Thus
its zero modes are anti-holomorphic spinors, equivalently conjugates of
holomorphic sections: \[
\widehat\psi_+\in \overline{H^0(\Sigma,L)}.
\] In particular, for the untwisted Dirac operator on a compact Riemann
surface, the two chiral zero-mode spaces have the same complex
dimension: \[
\dim_{\mathbb C}\ker D_+
=
\dim_{\mathbb C}\ker D_-
=
h^0(\Sigma,L),
\] up to the assignment of which sign of chirality is called plus or
minus.

This is also consistent with the ordinary index theorem. Since
\(L^2=K\), \[
\deg L
=
\frac{1}{2}\deg K
=
\frac{1}{2}(2g-2)
=
g-1.
\] Riemann-Roch gives \[
h^0(\Sigma,L)-h^0(\Sigma,K\otimes L^{-1})
=
\deg L+1-g.
\] But \[
K\otimes L^{-1}\simeq L,
\] because \(L^2\simeq K\). Hence \[
h^0(\Sigma,L)-h^0(\Sigma,L)=0.
\] So the integer index of the chiral Dirac operator is zero: \[
\operatorname{ind}D=0.
\] The interesting invariant is not this integer index, but the parity
\[
h^0(\Sigma,L)\bmod 2.
\]

\subsubsection{\texorpdfstring{The \(2^{2g}\) spin
structures}{The 2\^{}\{2g\} spin structures}}\label{the-22g-spin-structures}

The set of spin structures on \(\Sigma_g\) is not a vector space
canonically, but it is a torsor for \[
H^1(\Sigma_g,\mathbb Z_2).
\] Since \[
\dim_{\mathbb Z_2}H^1(\Sigma_g,\mathbb Z_2)=2g,
\] there are \[
\left|H^1(\Sigma_g,\mathbb Z_2)\right|
=
2^{2g}
\] spin structures.

Equivalently, once one spin bundle \(L\) is chosen, all others are of
the form \[
L\otimes M,
\] where \(M\) is a flat line bundle of order two: \[
M^2\simeq \mathcal O.
\] The group of such \(M\) is \[
\operatorname{Jac}(\Sigma_g)[2]\simeq H^1(\Sigma_g,\mathbb Z_2),
\] which again has \(2^{2g}\) elements.

\subsubsection{Why the mod-two number of chiral zero modes is
topological}\label{why-the-mod-two-number-of-chiral-zero-modes-is-topological}

The statement is \[
h^0(\Sigma,L)\bmod 2
\] depends only on the topological spin structure, not on the chosen
metric or complex structure.

This is not obvious from the formula \(h^0(\Sigma,L)\) itself. The
actual number \(h^0(\Sigma,L)\) can jump as the complex structure
varies. What is invariant is its parity.

The clean topological explanation uses the quadratic refinement
associated to a spin structure. Let \[
V=H_1(\Sigma_g,\mathbb Z_2).
\] The mod-two intersection pairing \[
\langle\ ,\ \rangle:V\times V\to\mathbb Z_2
\] is symplectic. A spin structure determines a function \[
q_L:V\to\mathbb Z_2
\] satisfying \[
q_L(a+b)
=
q_L(a)+q_L(b)+\langle a,b\rangle.
\] Such a function is called a quadratic refinement of the intersection
form.

Geometrically, \(q_L(\gamma)\) measures whether the restriction of the
spin structure to the loop \(\gamma\) is periodic or antiperiodic, with
a convention-dependent shift chosen so that the quadratic-refinement
identity holds. The important point is that \(q_L\) depends only on the
topological spin structure.

Given a symplectic basis \[
a_1,b_1,\dots,a_g,b_g
\] of \(H_1(\Sigma_g,\mathbb Z_2)\), the Arf invariant is \[
\operatorname{Arf}(q_L)
=
\sum_{i=1}^g q_L(a_i)q_L(b_i)
\pmod 2.
\] Although this formula uses a symplectic basis, the resulting number
is independent of the chosen symplectic basis. Thus \[
\operatorname{Arf}(q_L)\in\mathbb Z_2
\] is a topological invariant of the spin structure.

The theorem relating this to zero modes is the Riemann-Mumford parity
theorem: \[
h^0(\Sigma,L)
\equiv
\operatorname{Arf}(q_L)
\pmod 2.
\] Therefore \[
h^0(\Sigma,L)\bmod 2
\] is topological.

This proves the first requested statement: \[
\boxed{
\text{For a fixed spin structure }L,\ 
\dim_{\mathbb C}\ker D_{\text{chiral}}\bmod 2
\text{ is a topological invariant.}
}
\]

The word ``fixed'' is important. If one changes the spin structure, the
parity may change. Spin structures are divided into even and odd spin
structures: \[
L\ \text{even}
\quad\Longleftrightarrow\quad
h^0(\Sigma,L)\equiv 0\pmod 2,
\] and \[
L\ \text{odd}
\quad\Longleftrightarrow\quad
h^0(\Sigma,L)\equiv 1\pmod 2.
\] For genus \(g\), the numbers of even and odd spin structures are \[
N_{\text{even}}=2^{g-1}(2^g+1),
\qquad
N_{\text{odd}}=2^{g-1}(2^g-1).
\] This is another way to see that the parity is not the same for all
spin structures; it is an invariant attached to a chosen spin structure.

\subsubsection{Additivity under gluing}\label{additivity-under-gluing}

Now consider gluing two compact Riemann surfaces by connected sum: \[
\Sigma_{g_1}\#\Sigma_{g_2}\simeq\Sigma_{g_1+g_2}.
\] Assume the spin structures \(L_1\) and \(L_2\) are glued compatibly
to produce a spin structure \[
L=L_1\# L_2
\] on \(\Sigma_{g_1+g_2}\).

Topologically, \[
H_1(\Sigma_{g_1+g_2},\mathbb Z_2)
\simeq
H_1(\Sigma_{g_1},\mathbb Z_2)
\oplus
H_1(\Sigma_{g_2},\mathbb Z_2),
\] and the mod-two intersection pairing is the orthogonal direct sum of
the two intersection pairings: \[
\langle\ ,\ \rangle
=
\langle\ ,\ \rangle_1\oplus \langle\ ,\ \rangle_2.
\] The glued spin structure gives the quadratic refinement \[
q_L=q_{L_1}\oplus q_{L_2},
\] meaning that for \[
x\in H_1(\Sigma_{g_1},\mathbb Z_2),
\qquad
y\in H_1(\Sigma_{g_2},\mathbb Z_2),
\] one has \[
q_L(x,y)=q_{L_1}(x)+q_{L_2}(y).
\] The Arf invariant is additive under orthogonal direct sums: \[
\operatorname{Arf}(q_{L_1}\oplus q_{L_2})
=
\operatorname{Arf}(q_{L_1})
+
\operatorname{Arf}(q_{L_2})
\pmod 2.
\] Using the Riemann-Mumford parity theorem on each surface gives \[
h^0(\Sigma_{g_1+g_2},L_1\# L_2)
\equiv
\operatorname{Arf}(q_L)
\pmod 2,
\] and therefore \[
h^0(\Sigma_{g_1+g_2},L_1\# L_2)
\equiv
\operatorname{Arf}(q_{L_1})
+
\operatorname{Arf}(q_{L_2})
\pmod 2.
\] Applying the same theorem again, \[
\operatorname{Arf}(q_{L_i})
\equiv
h^0(\Sigma_{g_i},L_i)
\pmod 2,
\] so \[
h^0(\Sigma_{g_1+g_2},L_1\# L_2)
\equiv
h^0(\Sigma_{g_1},L_1)
+
h^0(\Sigma_{g_2},L_2)
\pmod 2.
\] This proves the second statement: \[
\boxed{
\dim\ker D_{\text{chiral},\,L_1\# L_2}
\equiv
\dim\ker D_{\text{chiral},\,L_1}
+
\dim\ker D_{\text{chiral},\,L_2}
\pmod 2.
}
\]

The proof shows exactly why the result is true: gluing corresponds to
direct sum on \(H_1(-,\mathbb Z_2)\), and the mod-two zero-mode count is
the Arf invariant of the associated quadratic refinement.

\subsubsection{Checks}\label{checks}

For genus \(g=0\), there is only one spin structure. Since \[
K_{\mathbb P^1}\simeq\mathcal O(-2),
\] the spin bundle is \[
L\simeq\mathcal O(-1).
\] It has no holomorphic sections: \[
h^0(\mathbb P^1,\mathcal O(-1))=0.
\] Thus the spin structure is even.

For genus \(g=1\), the canonical bundle is trivial: \[
K\simeq\mathcal O.
\] There are four spin structures, corresponding to the four choices of
periodic or antiperiodic boundary conditions around the two fundamental
cycles. The trivial spin bundle has \[
h^0(T^2,\mathcal O)=1,
\] so it is odd. The other three flat square roots have no holomorphic
sections, so they are even. This matches \[
N_{\text{odd}}=2^{0}(2^1-1)=1,
\qquad
N_{\text{even}}=2^{0}(2^1+1)=3.
\]

As a gluing check, if one connected-sums two tori and chooses the odd
spin structure on both, the parity on the genus-two surface is \[
1+1\equiv 0\pmod 2.
\] So the glued spin structure is even. If one glues an odd torus spin
structure with an even torus spin structure, the resulting genus-two
spin structure is odd.

\subsubsection{Result}\label{result-4}

With \[
z=x^0+ix^1,
\qquad
ds^2=e^{2\omega}dz\,d\bar z,
\] and with \[
\gamma^2=i\gamma^0\gamma^1
=
\begin{pmatrix}
1 & 0\\
0 & -1
\end{pmatrix},
\] the Dirac operator is \[
\slashed D
=
2i e^{-\omega}
\begin{pmatrix}
0 & -\left(\partial_{\bar z}+\frac{1}{2}\partial_{\bar z}\omega\right)\\
\partial_z+\frac{1}{2}\partial_z\omega & 0
\end{pmatrix}.
\] Thus \[
\slashed D\psi=0
\quad\Longleftrightarrow\quad
\left(\partial_z+\frac{1}{2}\partial_z\omega\right)\psi_+=0,
\qquad
\left(\partial_{\bar z}+\frac{1}{2}\partial_{\bar z}\omega\right)\psi_-=0.
\] After \(\widehat\psi_\pm=e^{\omega/2}\psi_\pm\), \[
\partial_z\widehat\psi_+=0,
\qquad
\partial_{\bar z}\widehat\psi_-=0.
\] A spin structure is \[
L^2\simeq K,
\] and chiral zero modes are counted by \[
h^0(\Sigma,L)=\dim_{\mathbb C} H^0(\Sigma,L),
\] up to the convention of which chirality is called plus or minus.

The parity theorem says \[
h^0(\Sigma,L)
\equiv
\operatorname{Arf}(q_L)
\pmod 2.
\] Therefore \[
\boxed{
h^0(\Sigma,L)\bmod 2
\text{ is a topological invariant of the spin structure.}
}
\] For gluing by connected sum, \[
q_{L_1\# L_2}=q_{L_1}\oplus q_{L_2},
\] so \[
\operatorname{Arf}(q_{L_1\# L_2})
=
\operatorname{Arf}(q_{L_1})
+
\operatorname{Arf}(q_{L_2})
\pmod 2.
\] Equivalently, \[
\boxed{
h^0(\Sigma_{g_1+g_2},L_1\# L_2)
\equiv
h^0(\Sigma_{g_1},L_1)
+
h^0(\Sigma_{g_2},L_2)
\pmod 2.
}
\]

\section{Q2: Counting even and odd spin
structures}\label{q2-counting-even-and-odd-spin-structures}

We call a spin structure even if the number of zero modes of the chiral
Dirac operator is even and we call it odd otherwise. Prove: there are
\(\sum_{m\text{ odd}}\binom{g}{m}3^{g-m}=2^{g-1}(2^g-1)\) odd and
\(\sum_{m\text{ even}}\binom{g}{m}3^{g-m}=2^{g-1}(2^g+1)\) even spin
structures.

\begin{quote}
\textbf{Textbook location and related material.} Chapter 9, Section 9.1
(pp.~223--232), together with the theta-function appendix in Section 9.6
(pp.~254--262). The genus-\(g\) counting by quadratic refinements is a
generalization of the book's genus-one discussion.
\end{quote}

\subsection{Answer}\label{answer-5}

\subsubsection{Core point}\label{core-point-5}

The question asks for the number of even and odd spin structures on a
genus-\(g\) compact Riemann surface \(\Sigma_g\). By definition here, \[
L\ \text{is even}
\quad\Longleftrightarrow\quad
h^0(\Sigma_g,L)\equiv 0\pmod 2,
\] and \[
L\ \text{is odd}
\quad\Longleftrightarrow\quad
h^0(\Sigma_g,L)\equiv 1\pmod 2.
\]

The key point from Question 1 is the Riemann-Mumford parity theorem: \[
h^0(\Sigma_g,L)
\equiv
\operatorname{Arf}(q_L)
\pmod 2.
\] So the problem is equivalent to counting quadratic refinements \[
q:H_1(\Sigma_g,\mathbb Z_2)\to\mathbb Z_2
\] with Arf invariant \(0\) or \(1\).

In a symplectic basis \[
a_1,b_1,\dots,a_g,b_g
\] of \(H_1(\Sigma_g,\mathbb Z_2)\), the Arf invariant is \[
\operatorname{Arf}(q)
=
\sum_{i=1}^g q(a_i)q(b_i)
\pmod 2.
\] For each handle, there are four possible pairs \[
\left(q(a_i),q(b_i)\right)\in\mathbb Z_2^2.
\] Among these four choices, three contribute \(0\) to the Arf sum and
one contributes \(1\). Therefore an odd spin structure is obtained by
choosing an odd number of handles of the fourth type; an even spin
structure is obtained by choosing an even number.

\subsubsection{Spin structures as quadratic
refinements}\label{spin-structures-as-quadratic-refinements}

Let \[
V=H_1(\Sigma_g,\mathbb Z_2).
\] The mod-two intersection pairing \[
\langle\ ,\ \rangle:V\times V\to\mathbb Z_2
\] is a nondegenerate symplectic pairing. A spin structure determines a
quadratic refinement \[
q:V\to\mathbb Z_2
\] satisfying \[
q(x+y)
=
q(x)+q(y)+\langle x,y\rangle.
\] Conversely, every such quadratic refinement comes from a spin
structure. Thus the set of spin structures can be counted by counting
quadratic refinements of the intersection form.

Choose a symplectic basis \[
a_1,b_1,\dots,a_g,b_g
\] with \[
\langle a_i,b_j\rangle=\delta_{ij},
\qquad
\langle a_i,a_j\rangle=0,
\qquad
\langle b_i,b_j\rangle=0.
\] The values \[
q(a_i),
\qquad
q(b_i),
\qquad
i=1,\dots,g,
\] determine \(q\) on all of \(V\). Indeed, if \[
x=\sum_{i=1}^g \left(\alpha_i a_i+\beta_i b_i\right),
\qquad
\alpha_i,\beta_i\in\mathbb Z_2,
\] then the quadratic-refinement identity gives \[
q(x)
=
\sum_{i=1}^g \alpha_i q(a_i)
+
\sum_{i=1}^g \beta_i q(b_i)
+
\sum_{i=1}^g \alpha_i\beta_i
\pmod 2.
\] The final term is the correction coming from the nonzero
intersections \[
\langle a_i,b_i\rangle=1.
\] Therefore, once the \(2g\) numbers \[
q(a_1),q(b_1),\dots,q(a_g),q(b_g)
\] are chosen, the quadratic refinement is fixed.

Since each of these \(2g\) numbers can be \(0\) or \(1\), the total
number of quadratic refinements is \[
2^{2g}.
\] This agrees with the standard count of spin structures on
\(\Sigma_g\).

\subsubsection{Arf invariant handle by
handle}\label{arf-invariant-handle-by-handle}

The Arf invariant is \[
\operatorname{Arf}(q)
=
\sum_{i=1}^g q(a_i)q(b_i)
\pmod 2.
\] For a fixed handle \(i\), there are four choices: \[
\left(q(a_i),q(b_i)\right)
=
(0,0),(1,0),(0,1),(1,1).
\] Their contributions to \[
q(a_i)q(b_i)
\] are \[
\begin{array}{c|c}
\left(q(a_i),q(b_i)\right) & q(a_i)q(b_i)\\
\hline
(0,0) & 0\\
(1,0) & 0\\
(0,1) & 0\\
(1,1) & 1
\end{array}
\] So, for each handle: \[
\text{three choices contribute }0,
\qquad
\text{one choice contributes }1.
\]

Now let \(m\) be the number of handles for which \[
\left(q(a_i),q(b_i)\right)=(1,1).
\] Then \[
\operatorname{Arf}(q)
\equiv
m
\pmod 2.
\] Thus: \[
q\ \text{is odd}
\quad\Longleftrightarrow\quad
m\ \text{is odd},
\] and \[
q\ \text{is even}
\quad\Longleftrightarrow\quad
m\ \text{is even}.
\]

This already gives the desired combinatorial structure.

\subsubsection{Counting odd spin
structures}\label{counting-odd-spin-structures}

To construct an odd spin structure, choose an odd number \(m\) of
handles to have the locally odd choice \[
\left(q(a_i),q(b_i)\right)=(1,1).
\] There are \[
\binom{g}{m}
\] ways to choose those handles.

On each of these \(m\) handles, the choice is forced to be \((1,1)\). On
each of the remaining \(g-m\) handles, the contribution must be \(0\),
so there are three possible choices: \[
(0,0),
\qquad
(1,0),
\qquad
(0,1).
\] Therefore, for fixed odd \(m\), the number of spin structures with
exactly \(m\) contributing handles is \[
\binom{g}{m}3^{g-m}.
\] Summing over odd \(m\) gives \[
N_{\text{odd}}
=
\sum_{m\text{ odd}}\binom{g}{m}3^{g-m}.
\] To evaluate this sum, use the binomial identities \[
\sum_{m=0}^g \binom{g}{m}3^{g-m}
=
(3+1)^g
=
4^g,
\] and \[
\sum_{m=0}^g (-1)^m\binom{g}{m}3^{g-m}
=
(3-1)^g
=
2^g.
\] The first sum is \[
N_{\text{even}}+N_{\text{odd}},
\] while the second is \[
N_{\text{even}}-N_{\text{odd}},
\] because even \(m\) terms have sign \(+1\) and odd \(m\) terms have
sign \(-1\).

Therefore \[
N_{\text{odd}}
=
\frac{1}{2}\left(4^g-2^g\right)
=
\frac{1}{2}\left(2^{2g}-2^g\right)
=
2^{g-1}(2^g-1).
\] So \[
\boxed{
N_{\text{odd}}
=
\sum_{m\text{ odd}}\binom{g}{m}3^{g-m}
=
2^{g-1}(2^g-1).
}
\]

\subsubsection{Counting even spin
structures}\label{counting-even-spin-structures}

The same reasoning gives the even count. An even spin structure has \[
\operatorname{Arf}(q)=0,
\] so the number \(m\) of handles with \[
\left(q(a_i),q(b_i)\right)=(1,1)
\] must be even.

For fixed even \(m\), there are \[
\binom{g}{m}
\] ways to choose the contributing handles, and \[
3^{g-m}
\] ways to choose the remaining noncontributing handles. Thus \[
N_{\text{even}}
=
\sum_{m\text{ even}}\binom{g}{m}3^{g-m}.
\] Using the same binomial identities, \[
N_{\text{even}}
=
\frac{1}{2}\left(4^g+2^g\right)
=
\frac{1}{2}\left(2^{2g}+2^g\right)
=
2^{g-1}(2^g+1).
\] Therefore \[
\boxed{
N_{\text{even}}
=
\sum_{m\text{ even}}\binom{g}{m}3^{g-m}
=
2^{g-1}(2^g+1).
}
\]

As a check, \[
N_{\text{even}}+N_{\text{odd}}
=
2^{g-1}(2^g+1)+2^{g-1}(2^g-1)
=
2^{2g},
\] which is the total number of spin structures.

\subsubsection{Checks}\label{checks-1}

For \(g=1\), there are \[
N_{\text{odd}}
=
2^{0}(2^1-1)
=
1
\] odd spin structure and \[
N_{\text{even}}
=
2^{0}(2^1+1)
=
3
\] even spin structures. This is the familiar torus result: the fully
periodic spin structure is odd, while the other three are even.

For \(g=2\), \[
N_{\text{odd}}
=
2^{1}(2^2-1)
=
6,
\] and \[
N_{\text{even}}
=
2^{1}(2^2+1)
=
10.
\] Together they give \[
6+10=16=2^{4},
\] as required.

\subsubsection{Result}\label{result-5}

A spin structure \(L\) determines a quadratic refinement \[
q_L:H_1(\Sigma_g,\mathbb Z_2)\to\mathbb Z_2,
\qquad
q_L(x+y)=q_L(x)+q_L(y)+\langle x,y\rangle.
\] In a symplectic basis, \[
\operatorname{Arf}(q_L)
=
\sum_{i=1}^g q_L(a_i)q_L(b_i)
\pmod 2.
\] The parity theorem says \[
h^0(\Sigma_g,L)
\equiv
\operatorname{Arf}(q_L)
\pmod 2.
\] For each handle, the four possible pairs are \[
(0,0),(1,0),(0,1),(1,1).
\] The first three contribute \(0\) to the Arf invariant, while
\((1,1)\) contributes \(1\). Hence if \(m\) handles are of type
\((1,1)\), then \[
h^0(\Sigma_g,L)\equiv m\pmod 2.
\] Therefore \[
\boxed{
N_{\text{odd}}
=
\sum_{m\text{ odd}}\binom{g}{m}3^{g-m}
=
2^{g-1}(2^g-1)
}
\] and \[
\boxed{
N_{\text{even}}
=
\sum_{m\text{ even}}\binom{g}{m}3^{g-m}
=
2^{g-1}(2^g+1).
}
\]

\section{Q3: Fermion spin structures and theta-function
conventions}\label{q3-fermion-spin-structures-and-theta-function-conventions}

Consider the closed superstring partition function on torus. The
contribution to the partition function of the right-moving fermions with
fixed spin structure are \[
\begin{aligned}
\chi_{\text{fermion}}^{(++)}(\tau)&=\eta_{++}\mathrm{Tr}e^{2\pi i\tau H_{R}}(-1)^F\\
\chi_{\text{fermion}}^{(+-)}(\tau)&=\eta_{+ -}\mathrm{Tr}e^{2\pi i\tau H_{R}}\\
\chi_{\text{fermion}}^{(--)}(\tau)&=\eta_{--}\mathrm{Tr}e^{2\pi i\tau H_{NS}}\\
\chi_{\text{fermion}}^{(-+)}(\tau)&=\eta_{-+}\mathrm{Tr}e^{2\pi i\tau H_{NS}}(-1)^F\\
\end{aligned}
\] If we choose the complex coordinates on the torus to be \[
z=\sigma^1+i\sigma^0=2\pi(\xi^1+\tau \xi^2)
\] Then for a general twisted boundary condition \[
\psi(\xi^1+1,\xi^2)=-e^{2\pi i\alpha}\psi(\xi^1,\xi^2) \ , \ \psi(\xi^1,\xi^2+1)=-e^{2\pi i\beta}\psi(\xi^1,\xi^2)
\] The partition function can be computed to be \[
\begin{aligned}
Z\begin{bmatrix}\alpha\\\beta\end{bmatrix}(\tau)
&=\mathrm{Tr}(e^{2\pi i\beta(N-\bar N)}q^{L_0-1/24})\\
&=q^{\alpha^2/2-1/24}\prod_{n=1}^{\infty}(1+q^{n+\alpha-1/2}e^{-2\pi i\beta})
\qquad\times(1+q^{n-\alpha-1/2}e^{2\pi i\beta})\\
&=\frac{e^{2\pi i\alpha\beta}\vartheta\begin{bmatrix}\alpha\\-\beta\end{bmatrix}(\tau)}{\eta(\tau)}.
\end{aligned}
\] In the textbook, the different spin structures correspond to \[
\begin{array}{lll}
(++) & \alpha=\frac{1}{2},\beta=\frac{1}{2} & \vartheta \begin{bmatrix}
1/2 \\
1/2
\end{bmatrix}=\vartheta_{1} \\
(+ -) & \alpha=\frac{1}{2},\beta=0 & \vartheta \begin{bmatrix}
1/2 \\
0
\end{bmatrix}=\vartheta_{2} \\ 
(--) & \alpha=0,\beta=0 & \vartheta \begin{bmatrix}
0 \\
0
\end{bmatrix}=\vartheta_{3} \\
(-+) & \alpha=0,\beta=\frac{1}{2} & \vartheta \begin{bmatrix}
0 \\
1/2
\end{bmatrix}=\vartheta_{4} \\
\end{array}
\] Q1: the expression computed for twisted boundary condition is
\(\vartheta \begin{bmatrix}\alpha\\-\beta\end{bmatrix}\), while the
above are \(\vartheta \begin{bmatrix}\alpha\\\beta\end{bmatrix}\), are
these equivalent?

There are many identities among the Jacobi \(\vartheta\) functions. Some
of them are \[
\begin{aligned}
&\vartheta_{2}^4(\tau)-\vartheta_{3}^4(\tau)+\vartheta_{4}^4(\tau)=0\\
&\vartheta_{2}(\tau)\vartheta_{3}(\tau)\vartheta_{4}(\tau)=2\eta^3(\tau)\\
&\vartheta_{1}'(\tau)=2\pi \eta^3(\tau)
\end{aligned}
\] It is said that \(\vartheta_{1}(\tau)=0\). Q2: I am confused with
this: if \(\vartheta_{1}(\tau)=0\) then why
\(\vartheta_{1}'(\tau)=2\pi \eta^3(\tau)\) holds? How to compute
\(\vartheta_{1}(\tau)=0\)?

\begin{quote}
\textbf{Textbook location and related material.} Chapter 9, Sections 9.1
and 9.6 (pp.~223--232, 254--262): fermionic boundary conditions,
spin-structure traces, theta functions with characteristics, and their
identities.
\end{quote}

\subsection{Answer}\label{answer-6}

\subsubsection{Core point}\label{core-point-6}

There are two notation issues here.

First, the expression \[
Z\begin{bmatrix}
\alpha\\
\beta
\end{bmatrix}(\tau)
\propto
e^{2\pi i\alpha\beta}
\frac{
\vartheta\begin{bmatrix}
\alpha\\
-\beta
\end{bmatrix}(0|\tau)}
{\eta(\tau)}
\] does not really conflict with the table written using \[
\vartheta\begin{bmatrix}
\alpha\\
\beta
\end{bmatrix}(0|\tau).
\] For the four spin structures one only has \[
\alpha,\beta\in\left\{0,\frac{1}{2}\right\}.
\] Since characteristics are defined modulo integers, \(-\beta\) and
\(\beta\) describe the same spin structure. The theta function itself
can pick up a phase under this integer shift, but that phase is
precisely the sort of convention-dependent sign absorbed into the
spin-structure coefficient \(\eta_{\alpha\beta}\).

Second, \[
\vartheta_1(\tau)=0
\] is shorthand for \[
\vartheta_1(0|\tau)=0.
\] But \[
\vartheta_1'(\tau)=2\pi\eta^3(\tau)
\] means \[
\left.\frac{\partial}{\partial \nu}\vartheta_1(\nu|\tau)\right|_{\nu=0}
=
2\pi\eta^3(\tau).
\] So there is no contradiction: a function can vanish at a point and
still have a nonzero first derivative there.

\subsubsection{Theta functions with
characteristics}\label{theta-functions-with-characteristics}

I will use the convention \[
\vartheta\begin{bmatrix}
\alpha\\
\beta
\end{bmatrix}(\nu|\tau)
=
\sum_{n\in\mathbb Z}
\exp\left[
\pi i(n+\alpha)^2\tau
+
2\pi i(n+\alpha)(\nu+\beta)
\right].
\] At \(\nu=0\) this becomes \[
\vartheta\begin{bmatrix}
\alpha\\
\beta
\end{bmatrix}(0|\tau)
=
\sum_{n\in\mathbb Z}
q^{\frac{1}{2}(n+\alpha)^2}
e^{2\pi i(n+\alpha)\beta},
\qquad
q=e^{2\pi i\tau}.
\] The four standard Jacobi theta constants are then \[
\vartheta_1(\tau)
=
\vartheta\begin{bmatrix}
\frac{1}{2}\\
\frac{1}{2}
\end{bmatrix}(0|\tau),
\] \[
\vartheta_2(\tau)
=
\vartheta\begin{bmatrix}
\frac{1}{2}\\
0
\end{bmatrix}(0|\tau),
\] \[
\vartheta_3(\tau)
=
\vartheta\begin{bmatrix}
0\\
0
\end{bmatrix}(0|\tau),
\] and \[
\vartheta_4(\tau)
=
\vartheta\begin{bmatrix}
0\\
\frac{1}{2}
\end{bmatrix}(0|\tau).
\] Different textbooks sometimes differ by signs in the lower
characteristic or by phases in the definition of the trace. The physical
spin structure is the periodicity data; the corresponding theta constant
can differ by an overall phase depending on convention.

\subsubsection{\texorpdfstring{Why \(\beta\) and \(-\beta\) describe the
same spin
structure}{Why \textbackslash beta and -\textbackslash beta describe the same spin structure}}\label{why-beta-and--beta-describe-the-same-spin-structure}

From the definition, \[
\vartheta\begin{bmatrix}
\alpha\\
\beta+1
\end{bmatrix}(\nu|\tau)
=
e^{2\pi i\alpha}
\vartheta\begin{bmatrix}
\alpha\\
\beta
\end{bmatrix}(\nu|\tau).
\] This is because shifting \(\beta\mapsto\beta+1\) multiplies each term
in the sum by \[
e^{2\pi i(n+\alpha)}
=
e^{2\pi i\alpha}.
\] Similarly, \[
\vartheta\begin{bmatrix}
\alpha\\
\beta-1
\end{bmatrix}(\nu|\tau)
=
e^{-2\pi i\alpha}
\vartheta\begin{bmatrix}
\alpha\\
\beta
\end{bmatrix}(\nu|\tau).
\]

Now for the spin structures on the torus, \[
\beta=0
\qquad
\text{or}
\qquad
\beta=\frac{1}{2}.
\] If \(\beta=0\), then \[
-\beta=0,
\] so there is no change.

If \(\beta=\frac{1}{2}\), then \[
-\frac{1}{2}
=
\frac{1}{2}-1.
\] Therefore \[
\vartheta\begin{bmatrix}
\alpha\\
-\frac{1}{2}
\end{bmatrix}(\nu|\tau)
=
e^{-2\pi i\alpha}
\vartheta\begin{bmatrix}
\alpha\\
\frac{1}{2}
\end{bmatrix}(\nu|\tau).
\] For \(\alpha=0\), this phase is \(1\). For \(\alpha=\frac{1}{2}\),
this phase is \(-1\).

Thus \[
\vartheta\begin{bmatrix}
0\\
-\frac{1}{2}
\end{bmatrix}
=
\vartheta\begin{bmatrix}
0\\
\frac{1}{2}
\end{bmatrix},
\] while \[
\vartheta\begin{bmatrix}
\frac{1}{2}\\
-\frac{1}{2}
\end{bmatrix}
=
-
\vartheta\begin{bmatrix}
\frac{1}{2}\\
\frac{1}{2}
\end{bmatrix}.
\] At \(\nu=0\), the latter is just \[
\vartheta_1(\tau)\mapsto -\vartheta_1(\tau),
\] but \[
\vartheta_1(\tau)=0,
\] so even this sign is invisible for the theta constant itself.

More generally, even when the theta function is not zero, this sign is
not a physical change of spin structure. It is an overall phase
convention. In the fermion partition function, such phases are included
in the coefficients \[
\eta_{++},
\quad
\eta_{+-},
\quad
\eta_{--},
\quad
\eta_{-+}.
\] These are the GSO or spin-structure phases. Thus the replacement \[
\beta\longleftrightarrow -\beta
\] is harmless after one is consistent about the phase convention.

\subsubsection{Matching the four spin
structures}\label{matching-the-four-spin-structures}

The boundary conditions are \[
\psi(\xi^1+1,\xi^2)
=
-e^{2\pi i\alpha}\psi(\xi^1,\xi^2),
\] and \[
\psi(\xi^1,\xi^2+1)
=
-e^{2\pi i\beta}\psi(\xi^1,\xi^2).
\] For \[
\alpha=0,
\] the first boundary condition is antiperiodic: \[
\psi(\xi^1+1,\xi^2)=-\psi(\xi^1,\xi^2).
\] For \[
\alpha=\frac{1}{2},
\] it is periodic: \[
\psi(\xi^1+1,\xi^2)=+\psi(\xi^1,\xi^2).
\] Thus the sign convention is \[
-\quad\text{means NS,}
\qquad
+\quad\text{means R.}
\] The same interpretation applies to \(\beta\) around the other cycle.

Therefore \[
(++)
\quad\Longleftrightarrow\quad
\alpha=\frac{1}{2},
\quad
\beta=\frac{1}{2},
\] \[
(+-)
\quad\Longleftrightarrow\quad
\alpha=\frac{1}{2},
\quad
\beta=0,
\] \[
(--)
\quad\Longleftrightarrow\quad
\alpha=0,
\quad
\beta=0,
\] and \[
(-+)
\quad\Longleftrightarrow\quad
\alpha=0,
\quad
\beta=\frac{1}{2}.
\] This matches the table: \[
\begin{array}{lll}
(++) & \alpha=\frac{1}{2},\beta=\frac{1}{2}
& \vartheta\begin{bmatrix}
1/2\\
1/2
\end{bmatrix}
=\vartheta_1,\\
(+-) & \alpha=\frac{1}{2},\beta=0
& \vartheta\begin{bmatrix}
1/2\\
0
\end{bmatrix}
=\vartheta_2,\\
(--) & \alpha=0,\beta=0
& \vartheta\begin{bmatrix}
0\\
0
\end{bmatrix}
=\vartheta_3,\\
(-+) & \alpha=0,\beta=\frac{1}{2}
& \vartheta\begin{bmatrix}
0\\
1/2
\end{bmatrix}
=\vartheta_4.
\end{array}
\] If the determinant computation gives \(-\beta\) in the lower
characteristic, then the same table is recovered after using the
periodicity of the characteristic and adjusting the \(\eta\) phases.

\subsubsection{\texorpdfstring{Why
\(\vartheta_1(0|\tau)=0\)}{Why \textbackslash vartheta\_1(0\textbar\textbackslash tau)=0}}\label{why-vartheta_10tau0}

Now let us compute the vanishing of \(\vartheta_1\) directly from the
series.

By definition, \[
\vartheta_1(\nu|\tau)
=
\vartheta\begin{bmatrix}
\frac{1}{2}\\
\frac{1}{2}
\end{bmatrix}(\nu|\tau).
\] Therefore \[
\vartheta_1(\nu|\tau)
=
\sum_{n\in\mathbb Z}
\exp\left[
\pi i\left(n+\frac{1}{2}\right)^2\tau
+
2\pi i\left(n+\frac{1}{2}\right)
\left(\nu+\frac{1}{2}\right)
\right].
\] At \(\nu=0\), \[
\vartheta_1(0|\tau)
=
\sum_{n\in\mathbb Z}
\exp\left[
\pi i\left(n+\frac{1}{2}\right)^2\tau
+
\pi i\left(n+\frac{1}{2}\right)
\right].
\] Equivalently, \[
\vartheta_1(0|\tau)
=
\sum_{n\in\mathbb Z}
i(-1)^n
q^{\frac{1}{2}\left(n+\frac{1}{2}\right)^2}.
\] Now pair the term with index \(n\) with the term with index \[
n'=-n-1.
\] Then \[
n'+\frac{1}{2}
=
-n-\frac{1}{2}
=
-\left(n+\frac{1}{2}\right),
\] so the \(q\)-power is the same: \[
\left(n'+\frac{1}{2}\right)^2
=
\left(n+\frac{1}{2}\right)^2.
\] But the phase changes sign: \[
i(-1)^{n'}
=
i(-1)^{-n-1}
=
-i(-1)^n.
\] Therefore the \(n\) and \(-n-1\) terms cancel pairwise, so \[
\boxed{
\vartheta_1(0|\tau)=0.
}
\]

Equivalently, \(\vartheta_1(\nu|\tau)\) is an odd function of \(\nu\):
\[
\vartheta_1(-\nu|\tau)
=
-
\vartheta_1(\nu|\tau).
\] Every odd function vanishes at the origin, so \[
\vartheta_1(0|\tau)=0.
\]

\subsubsection{Why the derivative is
nonzero}\label{why-the-derivative-is-nonzero}

The identity \[
\vartheta_1'(\tau)=2\pi\eta^3(\tau)
\] is shorthand for \[
\vartheta_1'(0|\tau)
\equiv
\left.\frac{\partial}{\partial\nu}\vartheta_1(\nu|\tau)\right|_{\nu=0}
=
2\pi\eta^3(\tau).
\] The derivative is with respect to the elliptic variable \(\nu\), not
with respect to the modular parameter \(\tau\).

There is no contradiction with \[
\vartheta_1(0|\tau)=0.
\] For example, the elementary function \(f(\nu)=\nu\) satisfies \[
f(0)=0,
\qquad
f'(0)=1.
\] The same thing happens here: \(\vartheta_1(\nu|\tau)\) has a simple
zero at \(\nu=0\).

The cleanest derivation uses the Jacobi triple product: \[
\vartheta_1(\nu|\tau)
=
2q^{1/8}\sin(\pi\nu)
\prod_{n=1}^\infty
(1-q^n)
(1-2q^n\cos(2\pi\nu)+q^{2n}),
\] with \[
q=e^{2\pi i\tau}.
\] At \(\nu=0\), \[
\sin(\pi\nu)\big|_{\nu=0}=0,
\] so immediately \[
\vartheta_1(0|\tau)=0.
\] Now differentiate with respect to \(\nu\) and then set \(\nu=0\).
Since the zero comes from the factor \(\sin(\pi\nu)\), only the
derivative of that factor contributes at \(\nu=0\). Thus \[
\vartheta_1'(0|\tau)
=
2q^{1/8}\pi
\prod_{n=1}^\infty
(1-q^n)
(1-2q^n+q^{2n}).
\] But \[
1-2q^n+q^{2n}
=
(1-q^n)^2,
\] so \[
\vartheta_1'(0|\tau)
=
2\pi q^{1/8}
\prod_{n=1}^\infty
(1-q^n)^3.
\] Using the Dedekind eta function \[
\eta(\tau)
=
q^{1/24}
\prod_{n=1}^\infty(1-q^n),
\] we get \[
\eta^3(\tau)
=
q^{1/8}
\prod_{n=1}^\infty(1-q^n)^3.
\] Therefore \[
\boxed{
\vartheta_1'(0|\tau)
=
2\pi\eta^3(\tau).
}
\]

\subsubsection{Checks}\label{checks-2}

For the \((-+)\) spin structure, \[
\alpha=0,
\qquad
\beta=\frac{1}{2}.
\] Replacing \(\beta\) by \(-\beta\) gives \[
\vartheta\begin{bmatrix}
0\\
-1/2
\end{bmatrix}
=
\vartheta\begin{bmatrix}
0\\
1/2
\end{bmatrix},
\] because \(e^{-2\pi i\alpha}=1\). So \(\vartheta_4\) is completely
unchanged.

For the \((++)\) spin structure, \[
\alpha=\frac{1}{2},
\qquad
\beta=\frac{1}{2}.
\] Replacing \(\beta\) by \(-\beta\) gives \[
\vartheta\begin{bmatrix}
1/2\\
-1/2
\end{bmatrix}
=
-
\vartheta\begin{bmatrix}
1/2\\
1/2
\end{bmatrix}
=
-\vartheta_1.
\] At zero elliptic argument, \[
\vartheta_1(0|\tau)=0,
\] so the sign is irrelevant for the theta constant. Away from
\(\nu=0\), this sign is an overall convention that can be absorbed into
\(\eta_{++}\).

\subsubsection{Result}\label{result-6}

With \[
\vartheta\begin{bmatrix}
\alpha\\
\beta
\end{bmatrix}(\nu|\tau)
=
\sum_{n\in\mathbb Z}
\exp\left[
\pi i(n+\alpha)^2\tau
+
2\pi i(n+\alpha)(\nu+\beta)
\right],
\] one has \[
\vartheta\begin{bmatrix}
\alpha\\
\beta-1
\end{bmatrix}
=
e^{-2\pi i\alpha}
\vartheta\begin{bmatrix}
\alpha\\
\beta
\end{bmatrix}.
\] Therefore, for \(\beta=0,\frac{1}{2}\), \[
\vartheta\begin{bmatrix}
\alpha\\
-\beta
\end{bmatrix}
\] is the same spin structure as \[
\vartheta\begin{bmatrix}
\alpha\\
\beta
\end{bmatrix},
\] possibly up to a convention-dependent phase. That phase is absorbed
into the spin-structure coefficient \(\eta_{\alpha\beta}\).

The notation \[
\vartheta_1(\tau)=0
\] means \[
\vartheta_1(0|\tau)=0.
\] This follows either because \(\vartheta_1(\nu|\tau)\) is odd in
\(\nu\), or by pairwise cancellation in the series. The notation \[
\vartheta_1'(\tau)=2\pi\eta^3(\tau)
\] means \[
\left.\partial_\nu\vartheta_1(\nu|\tau)\right|_{\nu=0}
=
2\pi\eta^3(\tau).
\] Using the product formula, \[
\vartheta_1(\nu|\tau)
=
2q^{1/8}\sin(\pi\nu)
\prod_{n=1}^\infty
(1-q^n)(1-2q^n\cos(2\pi\nu)+q^{2n}),
\] one obtains \[
\vartheta_1'(0|\tau)
=
2\pi q^{1/8}\prod_{n=1}^\infty(1-q^n)^3
=
2\pi\eta^3(\tau).
\]

\section{Q4: Boundary-state gluing signs and NS-NS GSO
action}\label{q4-boundary-state-gluing-signs-and-ns-ns-gso-action}

Consider the open superstring. Use \(\sigma^{\pm}=\tau\pm\sigma\) as the
world-sheet coordinates, the fermion needs to satisfy the boundary
condition \[
\psi_{+}^\mu(0)=\pm \psi^\mu_{-}(0)\ , \ \psi_{+}^\mu(l)=\pm \psi_{-}^\mu(l)
\] Only the relative sign is relevant We separate the spacetime indices
\(\mu=(\alpha,i),\alpha=0,\cdots,p,i=p+1,\cdots,9\), and denote
\({D^\mu}_{\nu}=({\delta^\alpha}_{\beta},-{\delta^i}_{j})\), then the
boundary condition can be summarized as \[
\psi_{+}^\mu(0)={D^\mu}_{\nu}(0)\psi^\nu(0) \ , \ \psi_{+}^\mu(l)=\eta{D^\mu}_{\nu}(l)\psi^\nu(l) 
\] Where \(\eta=+1\) is called R sector and \(\eta=-1\) is the NS
sector. Only the relative sign is relevant and we can always move the
sign \(\eta\) to the \(\sigma=l\) boundary.

Now we want to construct boundary states in closed string theory via
open-closed duality. The transformation from the closed to the open
string channel via the exchange \(\tau\leftrightarrow\sigma\) or
\(\sigma^{\pm}\to\sigma'^{\pm}=\pm\sigma^{\pm}\); and
\(\psi_{+}'(\sigma'^{\pm})=\left( \frac{{\partial\sigma'^+}}{\partial\sigma^+} \right)^{-1.2}\psi_{+}(\sigma^+)\).
Then the open string boundary condition will become a gluing condition
on a boundary state: \[
(\psi_{+}'^\mu(\sigma'^{\pm})+\eta {D^\mu}_{\nu}\psi'^\nu_{-}(\sigma'^-))|B,\eta \rangle=0 \ , \ \text{at }\tau=0
\] Where \(\eta=\pm 1\). Q: Here in my understanding this \(\eta\) is
not the \(\eta\) for R/NS sector in open string. Then what is it? In the
following paragraph there may be some hints.

Through mode expansion of this gluing condition \[
(b_{r}^\mu+i\eta {D^\mu}_{\nu}\bar{b}^\nu_{-r})|B,\eta \rangle=0
\] We can solve this condition. For simplicity we only consider one
complex fermion \(\psi=\psi^1+i\psi^2\). In the NS sector \[
|B_{NS},\eta \rangle=\exp\left( \mp i\eta \sum_{i=1}^2\sum_{r>0}b_{-r}^i \bar{b}^i_{-r} \right)|0\rangle_{NS}
\] Where \(-\) sign is for \(N\) and the \(+\) sign for D boundary
condition. By computing the two NN or two DD overlaps \[
\langle B_{NS},\eta|e^{-2\pi lH_{\text{closed}}}|B_{NS},\eta \rangle=\frac{\vartheta_{3}(2il)}{\eta(2il)} \ , \ 
\langle B_{NS},\eta|e^{-2\pi lH_{\text{closed}}}|B_{NS},-\eta \rangle=\frac{\vartheta_{4}(2il)}{\eta(2il)}
\] We find that different relative signs of \(\eta\) correspond to
different spin structures. We see that \((\eta,\eta)=(+,+),(-,-)\)
overlaps yield the \(\mathrm{Tr}_{NS}q^{L_{0}-c/24}\) sector while the
\((-,+),(+,-)\) overlaps yield the \(\mathrm{Tr}_{R}q^{L_{0}-c/24}\)
sector. Here the relative sign \(-1\) will corresponds to the trace over
R sector. Also in the R sector \[
|B_{R},\eta \rangle=\exp\left( \mp i\eta \sum_{i=1}^2\sum_{r>0}b_{-r}^i \bar{b}^i_{-r} \right)|R\rangle \ , \ |R\rangle=(b_{0}^++i\eta \bar{b}^+_{0})|0\rangle_{R}
\] Where \(b_{0}^{\pm}=\frac{1}{\sqrt{ 2 }}(b_{0}^1\pm ib_{0}^2)\).
Again we can compute the NN and DD overlaps.

Now consider the boundary state corresponding to a Dp-brane. It requires
the gluing conditions \[
(\alpha^i_{n}+{D^i}_{j}\bar{\alpha}^i_{-n})|\mathrm{Dp},\eta \rangle=0 \ , \ (b^i_{r}+i\eta{D^i}_{j} \bar{b}^j_{-r})|\mathrm{Dp},\eta \rangle=0
\] Here we are in the light-cone gauge, the (0,1)-directions are
light-cone directions with N boundary conditions; and \(i=2,\cdots,9\)
are transverse coordinates i.e.~we have N for \(i=0,\cdots,p\) and D for
\(i=p+1,\cdots,9\). The solution to these conditions are \[
|{\rm Dp},\eta\rangle=\int d k_{\perp} \exp \left[-\sum_{n>0}\frac{1}{n} \alpha^{i}_{-n} D_{ij} \overline{\alpha}^{j}_{-n} -i \eta\sum_{r>0}b^{ i}_{-r} D_{ij} \overline{b}^{j}_{-r} \right] e^{i k_{\perp}\cdot x_{0}^{ \perp}} |k_{\perp}\rangle ,
\] where \(x^\perp_{0}\) denotes the position of the D-brane in the
Dirichlet directions, i.e.~transverse to the brane's world-volume. In
the NS-NS sector the ground state \(|k_{\perp}\rangle\) is simply the
unique ground state with transverse momentum \(k_{\perp}\). In the R-R
sector the degenerate ground state also has to satisfy the gluing
condition for the zero modes.

In the NS-NS sector, acting with the GSO projections on the
\(|\mathrm{Dp},\eta \rangle_{NS}\) boundary states we get \[
(-1)^{F}|\mathrm{Dp},\eta\rangle_{\mathrm{NS}}=-|\mathrm{Dp},-\eta\rangle_{ \mathrm{NS}},\qquad(-1)^{\overline{F}}|\mathrm{Dp},\eta\rangle_{\mathrm{NS}}= -|\mathrm{Dp},-\eta\rangle_{\mathrm{NS}},
\] Where the minus sign is due to the fact that the NS ground state
carries odd fermion number. Q: how to derive this equations?

\begin{quote}
\textbf{Textbook location and related material.} Chapter 9, Section 9.2
(pp.~232--235): fermionic boundary states, gluing signs, overlaps, and
GSO projection. Section 7.4 (pp.~189--192) contains the open-string NS/R
boundary conditions used in the channel exchange.
\end{quote}

\subsection{Answer}\label{answer-7}

\subsubsection{Core point}\label{core-point-7}

There are two different uses of signs here, and separating them resolves
the confusion.

In the open-string channel, the relative sign between the two ends of
the strip determines whether the open-string fermion is periodic or
antiperiodic around the doubled interval. That is what gives the
open-string R or NS sector.

In the closed-string boundary-state channel, the sign \[
\eta=\pm 1
\] in \[
\left(b_r^\mu+i\eta D^\mu{}_\nu\bar b_{-r}^\nu\right)|B,\eta\rangle=0
\] is a boundary gluing label. It says which linear combination of left-
and right-moving fermions vanishes at that boundary. A single boundary
label \(\eta\) is not itself the open-string R/NS sector. On an annulus
with two boundary states, \[
\langle B,\eta_1|e^{-2\pi lH_{\mathrm{closed}}}|B,\eta_2\rangle,
\] the relative sign \[
\eta_1\eta_2
\] does determine the spin structure, but not by the naive
identification \[
s_0s_l=\eta_1\eta_2.
\] With the conventions in the question, the open-channel endpoint sign
is instead \[
\boxed{
s_0s_l=-\eta_1\eta_2.
}
\] The extra minus sign comes from the open-closed channel exchange: the
two boundaries of the closed-string cylinder induce opposite
orientations on the open-string strip, and a worldsheet fermion is a
half-differential. Equivalently, the square-root Jacobian in the map \[
\sigma^\pm\mapsto \sigma'^\pm=\pm\sigma^\pm
\] contributes a relative phase between the two ends. This is not a new
sign in the bosonic gluing condition; it is specifically a sign/phase in
the spin lift of the coordinate transformation. This is why \[
\eta_1\eta_2=+1
\quad\Longrightarrow\quad
s_0s_l=-1
\quad\Longrightarrow\quad
\text{open NS},
\] whereas \[
\eta_1\eta_2=-1
\quad\Longrightarrow\quad
s_0s_l=+1
\quad\Longrightarrow\quad
\text{open R}.
\]

The GSO equations \[
(-1)^F|\mathrm{Dp},\eta\rangle_{\mathrm{NS}}
=
-|\mathrm{Dp},-\eta\rangle_{\mathrm{NS}},
\qquad
(-1)^{\bar F}|\mathrm{Dp},\eta\rangle_{\mathrm{NS}}
=
-|\mathrm{Dp},-\eta\rangle_{\mathrm{NS}}
\] come from two facts:

\begin{enumerate}
\def\labelenumi{\arabic{enumi}.}
\tightlist
\item
  \((-1)^F\) flips the sign of every left-moving fermion oscillator
  \(b_r^i\), while \((-1)^{\bar F}\) flips the sign of every
  right-moving fermion oscillator \(\bar b_r^i\).
\item
  In the NS sector, the ground state is conventionally assigned odd
  fermion number: \[
  (-1)^F|0\rangle_{\mathrm{NS}}
  =
  -|0\rangle_{\mathrm{NS}},
  \qquad
  (-1)^{\bar F}|0\rangle_{\mathrm{NS}}
  =
  -|0\rangle_{\mathrm{NS}}.
  \]
\end{enumerate}

\subsubsection{Open-channel relative sign versus boundary-state
label}\label{open-channel-relative-sign-versus-boundary-state-label}

For an open string on a strip, the fermion boundary conditions may be
written schematically as \[
\psi_+^\mu(0)=s_0D^\mu{}_\nu\psi_-^\nu(0),
\qquad
\psi_+^\mu(l)=s_lD^\mu{}_\nu\psi_-^\nu(l),
\] where \[
s_0,s_l=\pm 1.
\] A simultaneous flip \[
s_0\mapsto -s_0,
\qquad
s_l\mapsto -s_l
\] can be absorbed by redefining one chiral fermion, so only the
relative sign \[
s_0s_l
\] is invariant. This relative sign determines whether the open-string
fermion is periodic or antiperiodic on the doubled interval. With the
convention in the question, \[
s_0s_l=+1
\quad\text{is called R,}
\qquad
s_0s_l=-1
\quad\text{is called NS.}
\] Equivalently, one can put all the relative sign at \(\sigma=l\) and
write it as the open-channel \(\eta\).

In the closed-string channel, the cylinder is viewed as a closed string
propagating between two boundary states. Each boundary state has its own
fermionic gluing sign: \[
\left(\psi_+^\mu+i\eta_aD^\mu{}_\nu\psi_-^\nu\right)|B,\eta_a\rangle=0,
\qquad
a=1,2.
\] Here \(\eta_a\) is attached to a boundary, not to the open string as
a whole. Therefore \(\eta_a\) is not the open R/NS label. The
open-channel spin structure is recovered only after comparing the two
boundaries.

There is one further sign that matters. When we rotate the cylinder and
reinterpret the closed-string propagation as an open-string trace, the
two closed-string boundaries become the two ends of the open string with
opposite induced boundary orientations. In the convention used in the
question, the result is \[
s_0s_l=-\eta_1\eta_2.
\] Let us unpack this statement carefully.

\subsubsection{Boundary orientation and why fermions see
it}\label{boundary-orientation-and-why-fermions-see-it}

Take the closed-channel cylinder coordinates to be \[
0\leq \tau_c\leq T,
\qquad
\sigma_c\sim\sigma_c+2\pi,
\] with worldsheet orientation \[
d\tau_c\wedge d\sigma_c.
\] The two boundary circles are \[
\tau_c=0,
\qquad
\tau_c=T.
\] The induced boundary orientation is defined by the convention that \[
\text{outward normal}\wedge\text{boundary orientation}
=
\text{surface orientation}.
\] At the upper boundary \(\tau_c=T\), the outward normal is
\(+\partial_{\tau_c}\). Therefore the induced boundary tangent is
\(+\partial_{\sigma_c}\). At the lower boundary \(\tau_c=0\), the
outward normal is \(-\partial_{\tau_c}\). Therefore the induced boundary
tangent is \(-\partial_{\sigma_c}\). Thus the two boundary components
have opposite induced orientations: \[
\tau_c=T:\quad +\partial_{\sigma_c},
\qquad
\tau_c=0:\quad -\partial_{\sigma_c}.
\] This does not mean that one cannot use the same coordinate symbol
\(\sigma_c\) on both boundaries. One can. The point is that \(\sigma_c\)
increasing agrees with the induced orientation on one boundary and
disagrees with it on the other.

A bosonic coordinate \(X^\mu\) is a scalar on the worldsheet: \[
X'^\mu(\sigma')=X^\mu(\sigma).
\] So if we reverse a boundary coordinate, \[
\sigma'\,=\,-\sigma,
\] then \(X^\mu\) itself does not acquire a square-root phase. Its
derivatives transform as ordinary vectors or one-forms. Any sign from
reversing the normal derivative is already part of the usual distinction
between Neumann and Dirichlet gluing. There is no extra spin-lift
ambiguity.

By contrast, the chiral fermions are spin fields on the worldsheet.
Locally, in light-cone coordinates, one may view them as
half-differentials: \[
\psi_+\, (d\sigma^+)^{1/2},
\qquad
\psi_-\, (d\sigma^-)^{1/2}.
\] Under a change of chiral coordinate, \[
\sigma^\pm\mapsto\sigma'^\pm(\sigma^\pm),
\] the components transform as \[
\psi_\pm'(\sigma'^\pm)
=
\left(
\frac{\partial\sigma'^\pm}{\partial\sigma^\pm}
\right)^{-1/2}
\psi_\pm(\sigma^\pm).
\] If the coordinate transformation reverses the relevant chiral
coordinate, \[
\sigma'^\pm=-\sigma^\pm,
\] then \[
\frac{\partial\sigma'^\pm}{\partial\sigma^\pm}=-1,
\] and therefore the fermion picks a square-root phase: \[
\left(-1\right)^{-1/2}=\pm i.
\] The choice of sign is a choice of spin lift. Once the spin lift is
fixed consistently, the two boundaries of the cylinder get opposite
induced orientation factors. In the standard boundary-state convention
this is exactly the origin of the \(i\) in \[
\left(\psi_+^\mu+i\eta D^\mu{}_\nu\psi_-^\nu\right)|B,\eta\rangle=0,
\] and of the extra minus in \[
s_0s_l=-\eta_1\eta_2.
\] Said differently: each closed boundary gluing condition contains a
factor \(i\eta_a\) because rotating an open boundary into a closed
boundary involves a square root of \(-1\). When we compare the two
boundaries to reconstruct the open-string relative sign, the two
\(i\)-type factors multiply to \[
i^2=-1,
\] up to the convention-dependent placement of signs. This is the extra
minus between \(s_0s_l\) and \(\eta_1\eta_2\).

So the closed-channel overlaps distinguish \[
\eta_1\eta_2=+1
\qquad
\text{from}
\qquad
\eta_1\eta_2=-1.
\] In the NS-NS oscillator overlap, one finds factors of the schematic
form \[
\prod_{r>0}\left(1+\eta_1\eta_2\,q^{2r}\right).
\] Thus the same-sign overlap and opposite-sign overlap give different
theta functions: \[
\eta_1\eta_2=+1
\quad\Longrightarrow\quad
\vartheta_3,
\] whereas \[
\eta_1\eta_2=-1
\quad\Longrightarrow\quad
\vartheta_4.
\] After the modular transformation to the open-string channel, the
same-sign closed overlap gives the open NS character, while the
opposite-sign closed overlap gives the open R character. In
theta-function language, this is the familiar exchange in which the
closed-channel \(\vartheta_4\) contribution turns into the open-channel
\(\vartheta_2\) contribution. So the slogan is: \[
\boxed{
\eta_1\eta_2=+1\ \text{gives open NS, while}\ 
\eta_1\eta_2=-1\ \text{gives open R.}
}
\] This is consistent with the open-string rule \[
s_0s_l=+1\ \text{for R},
\qquad
s_0s_l=-1\ \text{for NS},
\] because the open and closed signs are related by \[
s_0s_l=-\eta_1\eta_2.
\]

\subsubsection{The NS-NS boundary state}\label{the-ns-ns-boundary-state}

Ignoring normalization and ghosts, the NS-NS fermionic part of the
D\(p\)-brane boundary state is \[
|\mathrm{Dp},\eta\rangle_{\mathrm{NS}}
=
\exp\left[
-i\eta
\sum_{r>0}
b^i_{-r}D_{ij}\bar b^j_{-r}
\right]
|0\rangle_{\mathrm{NS}}.
\] Here \[
r\in\mathbb Z+\frac{1}{2},
\qquad
r>0,
\] and \(i,j\) run over the physical transverse directions in light-cone
gauge. The matrix \(D_{ij}\) is \[
D_{ij}
=
\begin{cases}
\delta_{ij}, & \text{Neumann direction},\\
-\delta_{ij}, & \text{Dirichlet direction}.
\end{cases}
\] The precise sign convention in front of \(i\eta\) can differ between
books. What matters for the following derivation is that \(\eta\)
appears linearly in the exponent: \[
\eta\mapsto -\eta
\] is equivalent to flipping the sign of the fermion bilinear in the
exponential.

The state satisfies \[
\left(b_r^i+i\eta D^i{}_j\bar b^j_{-r}\right)
|\mathrm{Dp},\eta\rangle_{\mathrm{NS}}=0.
\] The sign \(\eta\) therefore labels the two possible fermionic gluing
conditions.

\subsubsection{\texorpdfstring{Action of \((-1)^F\) on the
oscillators}{Action of (-1)\^{}F on the oscillators}}\label{action-of--1f-on-the-oscillators}

The left-moving fermion number operator \(F\) is defined so that \[
(-1)^F b_r^i (-1)^F=-b_r^i,
\] while it leaves right-moving oscillators invariant: \[
(-1)^F\bar b_r^i(-1)^F=\bar b_r^i.
\] Equivalently, \[
(-1)^F b_{-r}^i
=
-b_{-r}^i(-1)^F,
\qquad
(-1)^F\bar b_{-r}^i
=
\bar b_{-r}^i(-1)^F.
\] Therefore, acting on the bilinear, \[
(-1)^F
\left(b^i_{-r}D_{ij}\bar b^j_{-r}\right)
(-1)^F
=
-
b^i_{-r}D_{ij}\bar b^j_{-r}.
\] Thus \((-1)^F\) changes the sign of the exponent: \[
(-1)^F
\exp\left[
-i\eta
\sum_{r>0}b^i_{-r}D_{ij}\bar b^j_{-r}
\right]
(-1)^F
=
\exp\left[
+i\eta
\sum_{r>0}b^i_{-r}D_{ij}\bar b^j_{-r}
\right].
\] But this is the same exponential that appears in \[
|\mathrm{Dp},-\eta\rangle_{\mathrm{NS}},
\] because replacing \(\eta\) by \(-\eta\) gives \[
-i(-\eta)=+i\eta.
\]

Now we must also act on the NS ground state. With the standard
superstring convention, \[
(-1)^F|0\rangle_{\mathrm{NS}}
=
-|0\rangle_{\mathrm{NS}}.
\] Therefore \[
\begin{aligned}
(-1)^F|\mathrm{Dp},\eta\rangle_{\mathrm{NS}}
&=
(-1)^F
\exp\left[
-i\eta
\sum_{r>0}b^i_{-r}D_{ij}\bar b^j_{-r}
\right]
|0\rangle_{\mathrm{NS}}\\
&=
\exp\left[
+i\eta
\sum_{r>0}b^i_{-r}D_{ij}\bar b^j_{-r}
\right]
(-1)^F|0\rangle_{\mathrm{NS}}\\
&=
-
\exp\left[
+i\eta
\sum_{r>0}b^i_{-r}D_{ij}\bar b^j_{-r}
\right]
|0\rangle_{\mathrm{NS}}\\
&=
-|\mathrm{Dp},-\eta\rangle_{\mathrm{NS}}.
\end{aligned}
\] This proves \[
\boxed{
(-1)^F|\mathrm{Dp},\eta\rangle_{\mathrm{NS}}
=
-|\mathrm{Dp},-\eta\rangle_{\mathrm{NS}}.
}
\]

\subsubsection{\texorpdfstring{Action of
\((-1)^{\bar F}\)}{Action of (-1)\^{}\{\textbackslash bar F\}}}\label{action-of--1bar-f}

The right-moving fermion number operator acts oppositely: \[
(-1)^{\bar F}\bar b_r^i(-1)^{\bar F}
=
-\bar b_r^i,
\] while leaving left-moving oscillators invariant: \[
(-1)^{\bar F}b_r^i(-1)^{\bar F}
=
b_r^i.
\] Therefore \[
(-1)^{\bar F}
\left(b^i_{-r}D_{ij}\bar b^j_{-r}\right)
(-1)^{\bar F}
=
-
b^i_{-r}D_{ij}\bar b^j_{-r}.
\] So \((-1)^{\bar F}\) also sends \[
\eta\mapsto -\eta
\] in the oscillator exponential.

The right-moving NS ground state is also odd: \[
(-1)^{\bar F}|0\rangle_{\mathrm{NS}}
=
-|0\rangle_{\mathrm{NS}}.
\] Repeating the same calculation gives \[
\begin{aligned}
(-1)^{\bar F}|\mathrm{Dp},\eta\rangle_{\mathrm{NS}}
&=
\exp\left[
+i\eta
\sum_{r>0}b^i_{-r}D_{ij}\bar b^j_{-r}
\right]
(-1)^{\bar F}|0\rangle_{\mathrm{NS}}\\
&=
-
\exp\left[
+i\eta
\sum_{r>0}b^i_{-r}D_{ij}\bar b^j_{-r}
\right]
|0\rangle_{\mathrm{NS}}\\
&=
-|\mathrm{Dp},-\eta\rangle_{\mathrm{NS}}.
\end{aligned}
\] Thus \[
\boxed{
(-1)^{\bar F}|\mathrm{Dp},\eta\rangle_{\mathrm{NS}}
=
-|\mathrm{Dp},-\eta\rangle_{\mathrm{NS}}.
}
\]

\subsubsection{Why the NS ground state is assigned odd fermion
number}\label{why-the-ns-ground-state-is-assigned-odd-fermion-number}

The minus sign in the final result is sometimes the most confusing part.
It is not produced by the oscillator exponential alone. The oscillator
exponential only gives \[
\eta\mapsto -\eta.
\] The extra overall minus sign comes from the convention \[
(-1)^F|0\rangle_{\mathrm{NS}}
=
-|0\rangle_{\mathrm{NS}}.
\] This convention is the standard one in the NS sector of the
superstring. It makes the NS tachyon odd under \((-1)^F\), so the usual
GSO projection removes it and keeps states with one fermionic
oscillator, such as \[
b_{-1/2}^i|0\rangle_{\mathrm{NS}}.
\] Indeed, \[
(-1)^F b_{-1/2}^i|0\rangle_{\mathrm{NS}}
=
\left((-1)^F b_{-1/2}^i(-1)^F\right)
(-1)^F|0\rangle_{\mathrm{NS}}
=
(-b_{-1/2}^i)(-|0\rangle_{\mathrm{NS}})
=
b_{-1/2}^i|0\rangle_{\mathrm{NS}}.
\] So the massless NS vector is even and survives the usual projection,
while the NS ground state is odd.

\subsubsection{GSO-invariant NS-NS boundary-state
combination}\label{gso-invariant-ns-ns-boundary-state-combination}

Because \[
(-1)^F|\mathrm{Dp},\eta\rangle_{\mathrm{NS}}
=
-|\mathrm{Dp},-\eta\rangle_{\mathrm{NS}},
\] neither \[
|\mathrm{Dp},+\rangle_{\mathrm{NS}}
\] nor \[
|\mathrm{Dp},-\rangle_{\mathrm{NS}}
\] is separately GSO invariant.

The invariant NS-NS combination is \[
|\mathrm{Dp}\rangle_{\mathrm{NSNS}}
\propto
|\mathrm{Dp},+\rangle_{\mathrm{NS}}
-
|\mathrm{Dp},-\rangle_{\mathrm{NS}}.
\] Check: \[
\begin{aligned}
(-1)^F
\left(
|\mathrm{Dp},+\rangle_{\mathrm{NS}}
-
|\mathrm{Dp},-\rangle_{\mathrm{NS}}
\right)
&=
-|\mathrm{Dp},-\rangle_{\mathrm{NS}}
+
|\mathrm{Dp},+\rangle_{\mathrm{NS}}\\
&=
|\mathrm{Dp},+\rangle_{\mathrm{NS}}
-
|\mathrm{Dp},-\rangle_{\mathrm{NS}}.
\end{aligned}
\] The same check works for \((-1)^{\bar F}\). This is why the physical
NS-NS part of the D-brane boundary state is built from a difference of
the two \(\eta\) gluing states.

\subsubsection{Checks}\label{checks-3}

If the NS ground state were assigned even fermion number instead, the
derivation would give \[
(-1)^F|\mathrm{Dp},\eta\rangle_{\mathrm{NS}}
=
|\mathrm{Dp},-\eta\rangle_{\mathrm{NS}}.
\] So the overall minus sign is convention-dependent in the precise
sense that it follows from the convention for the vacuum fermion number.
The standard superstring convention is the odd NS vacuum convention used
above.

Also, applying both left and right fermion parities gives \[
(-1)^{F+\bar F}|\mathrm{Dp},\eta\rangle_{\mathrm{NS}}
=
|\mathrm{Dp},\eta\rangle_{\mathrm{NS}}.
\] Indeed, the first action sends \[
|\eta\rangle\mapsto -|-\eta\rangle,
\] and the second sends \[
|-\eta\rangle\mapsto -|\eta\rangle.
\] The two minus signs cancel. This is consistent with the fact that the
NS-NS boundary state lies in a closed-string sector with matched left
and right GSO parity after projection.

\subsubsection{Result}\label{result-7}

The open-channel R/NS sign is a relative sign between the two strip
boundaries. The closed-channel boundary-state sign \[
\eta=\pm 1
\] is instead the local fermionic gluing label: \[
\left(b_r^i+i\eta D^i{}_j\bar b^j_{-r}\right)
|\mathrm{Dp},\eta\rangle=0.
\] On an annulus, the open spin structure is determined by \[
\eta_1\eta_2,
\] not by either \(\eta_1\) or \(\eta_2\) alone. More precisely, in the
conventions of the question, \[
\boxed{
s_0s_l=-\eta_1\eta_2.
}
\] Therefore \[
\eta_1\eta_2=+1
\quad\Longleftrightarrow\quad
s_0s_l=-1
\quad\Longleftrightarrow\quad
\text{open NS},
\] while \[
\eta_1\eta_2=-1
\quad\Longleftrightarrow\quad
s_0s_l=+1
\quad\Longleftrightarrow\quad
\text{open R}.
\] The extra minus sign is the spinor orientation factor from the
open-closed channel exchange. Bosonic coordinates do not acquire this
factor because they are worldsheet scalars, while fermions do because
their components transform with \[
\left(\frac{\partial\sigma'}{\partial\sigma}\right)^{-1/2}.
\]

In the NS-NS sector, \[
|\mathrm{Dp},\eta\rangle_{\mathrm{NS}}
=
\exp\left[
-i\eta
\sum_{r>0}b^i_{-r}D_{ij}\bar b^j_{-r}
\right]
|0\rangle_{\mathrm{NS}}.
\] The fermion parities act as \[
(-1)^F b_r^i(-1)^F=-b_r^i,
\qquad
(-1)^{\bar F}\bar b_r^i(-1)^{\bar F}=-\bar b_r^i,
\] and the NS vacuum convention is \[
(-1)^F|0\rangle_{\mathrm{NS}}
=
(-1)^{\bar F}|0\rangle_{\mathrm{NS}}
=
-|0\rangle_{\mathrm{NS}}.
\] Therefore \[
\boxed{
(-1)^F|\mathrm{Dp},\eta\rangle_{\mathrm{NS}}
=
-|\mathrm{Dp},-\eta\rangle_{\mathrm{NS}},
\qquad
(-1)^{\bar F}|\mathrm{Dp},\eta\rangle_{\mathrm{NS}}
=
-|\mathrm{Dp},-\eta\rangle_{\mathrm{NS}}.
}
\] Consequently the GSO-invariant NS-NS boundary state is proportional
to \[
\boxed{
|\mathrm{Dp}\rangle_{\mathrm{NSNS}}
\propto
|\mathrm{Dp},+\rangle_{\mathrm{NS}}
-
|\mathrm{Dp},-\rangle_{\mathrm{NS}}.
}
\]

\section{Q5: Force signs from even-spin, odd-spin, and R-R
exchange}\label{q5-force-signs-from-even-spin-odd-spin-and-r-r-exchange}

Consider two parallel Dp-branes.There are competing forces between these
two D-branes, namely attraction due to exchange of gravitons and
dilatons and repulsion due to the exchange of the R-R form. These two
forces precisely cancel, independent of the distance between the two
branes. Q: it is said that if the exchange particle has even spin, equal
sign charges attract and opposite charges repel. If the spin is odd, the
situation is reversed. Prove this in details. The exchange of an
antisymmetric tensor particle leads to repulsion between equal charges.

\begin{quote}
\textbf{Textbook location and related material.} Chapter 9, Section 9.3
(pp.~235--241): D-brane boundary states, NS--NS and R--R exchange, and
the BPS force cancellation. Section 6.8 gives the corresponding
long-distance field-theory comparison.
\end{quote}

\subsection{Answer}\label{answer-8}

\subsubsection{Core point}\label{core-point-8}

The question is why the sign of a force depends on the spin of the
exchanged field, and in particular why R-R form exchange between two
identical parallel D-branes is repulsive.

The clean rule for static sources is: \[
\boxed{
\text{even-spin exchange: equal charges attract, opposite charges repel,}
}
\] whereas \[
\boxed{
\text{odd-spin exchange: equal charges repel, opposite charges attract.}
}
\] For ordinary fields this means: \[
\text{scalar exchange}\quad\Rightarrow\quad\text{attraction for equal scalar charges},
\] \[
\text{vector exchange}\quad\Rightarrow\quad\text{repulsion for equal electric charges},
\] and \[
\text{graviton exchange}\quad\Rightarrow\quad\text{attraction for positive masses}.
\]

For D\(p\)-branes, the NS-NS fields couple to the tension \(T_p\) and
produce attraction between two parallel branes. The R-R potential
\(C_{p+1}\) couples to the oriented R-R charge \(\mu_p\) and produces
repulsion between two branes with equal R-R charge. For BPS branes,
these two effects exactly cancel: \[
V_{\mathrm{NSNS}}(r)+V_{\mathrm{RR}}(r)=0.
\]

\subsubsection{Static potential from integrating out a
field}\label{static-potential-from-integrating-out-a-field}

Let a field \(\Phi\) couple linearly to a source \(J\): \[
S[\Phi,J]
=
\frac{1}{2}\Phi K\Phi
+
J\Phi,
\] where \(K\) is the kinetic operator. Solving the classical equation
\[
K\Phi+J=0
\] gives \[
\Phi=-K^{-1}J.
\] Substituting back into the action gives the source-source effective
action \[
S_{\mathrm{eff}}[J]
=
-\frac{1}{2}J K^{-1}J.
\] For two separated static sources, \[
J=J_1+J_2,
\] the interaction term is \[
S_{\mathrm{int}}
=
-J_1K^{-1}J_2.
\] For a static potential, the effective action is related to the
potential by \[
S_{\mathrm{int}}
=
-\int dt\,V(r).
\] Therefore the sign of the potential is determined by the sign of the
propagator contraction between the two source currents.

This sounds abstract, so let us do the standard cases explicitly.

\subsubsection{Scalar exchange: equal charges
attract}\label{scalar-exchange-equal-charges-attract}

Take a real scalar \(\phi\) with mostly-plus signature, \[
S_\phi
=
\int d^Dx
\left[
-\frac{1}{2}\partial_\mu\phi\,\partial^\mu\phi
+
J\phi
\right].
\] For static sources, the energy functional is \[
E[\phi]
=
\int d^{D-1}x
\left[
\frac{1}{2}(\nabla\phi)^2
-
J\phi
\right].
\] The static equation is \[
-\nabla^2\phi=J.
\] Let \[
J(x)=q_1\delta(x-x_1)+q_2\delta(x-x_2).
\] The Green's function satisfies \[
-\nabla^2G(r)=\delta^{(D-1)}(r),
\qquad
G(r)>0.
\] The solution is \[
\phi(x)
=
q_1G(x-x_1)+q_2G(x-x_2).
\] Substituting the solution back into the energy gives \[
E_{\mathrm{on-shell}}
=
-\frac{1}{2}\int d^{D-1}x\,J\phi.
\] The cross term is therefore \[
V_\phi(r)
=
-q_1q_2G(r).
\] Thus for equal signs, \[
q_1q_2>0
\quad\Longrightarrow\quad
V_\phi(r)<0.
\] A negative potential that becomes more negative as the sources
approach means attraction. So scalar exchange attracts equal scalar
charges and repels opposite scalar charges.

\subsubsection{Vector exchange: equal charges
repel}\label{vector-exchange-equal-charges-repel}

Now take Maxwell theory: \[
S_A
=
\int d^Dx
\left[
-\frac{1}{4}F_{\mu\nu}F^{\mu\nu}
+
J_\mu A^\mu
\right].
\] For static electric sources, only \[
J^0=\rho
\] is nonzero. The interaction is mediated by \(A_0\). The electrostatic
equation is \[
-\nabla^2A_0=\rho.
\] Thus a point charge produces \[
A_0(x)=qG(x-x_a).
\] The potential energy of a second charge \(q_2\) in the potential
produced by the first charge is \[
V_A(r)=q_2A_0^{(1)}(x_2)
=
q_1q_2G(r).
\] Equivalently, the electric field is \[
E_i=-\partial_i A_0,
\] and the field energy \[
E_{\mathrm{field}}
=
\int d^{D-1}x
\frac{1}{2}E_iE_i
\] has a cross term \[
\int d^{D-1}x\,\nabla A_0^{(1)}\cdot\nabla A_0^{(2)}
=
q_1q_2G(r).
\] So \[
V_A(r)
=
+q_1q_2G(r).
\]

In covariant language, the same sign comes from the contraction of the
vector propagator with the time components of the current. In
mostly-plus signature, \[
J_\mu A^\mu
=
J_0A^0+\cdots
=
-J^0A^0+\cdots.
\] The propagator contraction for static sources contains \[
\eta_{00}=-1.
\] Together with the standard relation between exchange amplitude and
potential, this reverses the scalar sign. Thus \[
q_1q_2>0
\quad\Longrightarrow\quad
V_A(r)>0,
\] which means repulsion. Opposite charges have \[
q_1q_2<0,
\] so \[
V_A(r)<0,
\] which means attraction.

The vector case is the prototype of odd-spin exchange.

\subsubsection{Graviton exchange: positive masses
attract}\label{graviton-exchange-positive-masses-attract}

The linearized graviton couples to the stress tensor as \[
S_{\mathrm{int}}
=
\frac{\kappa}{2}
\int d^Dx\,h_{\mu\nu}T^{\mu\nu}.
\] For a nonrelativistic static source, \[
T^{00}\simeq m\,\delta^{(D-1)}(x-x_a),
\] and the spatial components \(T^{ij}\) are negligible. The graviton
propagator numerator in de Donder gauge is \[
P_{\mu\nu,\rho\sigma}
=
\frac{1}{2}
\left(
\eta_{\mu\rho}\eta_{\nu\sigma}
+
\eta_{\mu\sigma}\eta_{\nu\rho}
-
\frac{2}{D-2}\eta_{\mu\nu}\eta_{\rho\sigma}
\right).
\] For static sources, the relevant component is \[
P_{00,00}
=
\frac{1}{2}
\left(
2\eta_{00}^2
-
\frac{2}{D-2}\eta_{00}^2
\right)
=
1-\frac{1}{D-2}
=
\frac{D-3}{D-2}.
\] This is positive. Including the universal sign relating the exchange
amplitude to the potential, one finds \[
V_h(r)
=
-\frac{\kappa^2}{4}
\frac{D-3}{D-2}
m_1m_2G(r).
\] Therefore \[
m_1m_2>0
\quad\Longrightarrow\quad
V_h(r)<0.
\] Positive masses attract.

This is the spin-\(2\) analogue of scalar attraction. The two time
indices in \(T^{00}\) effectively produce \[
\eta_{00}^2=+1,
\] whereas the vector current had one time index and produced \[
\eta_{00}=-1.
\] This is the intuitive origin of the even-spin versus odd-spin rule.

\subsubsection{General spin sign rule for static
sources}\label{general-spin-sign-rule-for-static-sources}

For a symmetric spin-\(s\) field coupled schematically as \[
h_{\mu_1\cdots\mu_s}J^{\mu_1\cdots\mu_s},
\] a static source mainly has the component \[
J^{00\cdots 0}.
\] The propagator contraction contains one metric factor for each time
index, schematically \[
\eta_{00}^s=(-1)^s.
\] After including the standard relation between exchange amplitude and
potential, the static potential behaves as \[
V_s(r)
\sim
(-1)^{s+1}q_1q_2G(r).
\] Thus: \[
s\ \text{even}
\quad\Longrightarrow\quad
V_s(r)\sim -q_1q_2G(r),
\] so equal charges attract and opposite charges repel.

For \[
s\ \text{odd},
\] one gets \[
V_s(r)\sim +q_1q_2G(r),
\] so equal charges repel and opposite charges attract.

This is a useful rule of thumb. For gauge fields with traces,
constraints, or mixed symmetry, the detailed numerator matters, but the
static sign is still controlled by how the source indices contract with
the metric and by the sign of the kinetic term.

\subsubsection{R-R form exchange between
D-branes}\label{r-r-form-exchange-between-d-branes}

A D\(p\)-brane couples electrically to the R-R \((p+1)\)-form potential:
\[
S_{\mathrm{WZ}}
=
\mu_p
\int_{\mathrm{Dp}} C_{p+1}.
\] For a static flat D\(p\)-brane extended along \[
x^0,x^1,\dots,x^p,
\] this coupling is \[
S_{\mathrm{WZ}}
=
\mu_p
\int d^{p+1}\xi\,
C_{0 1\cdots p}.
\] So the brane is a source for the component \[
C_{0 1\cdots p}.
\] The R-R kinetic term has the conventional sign \[
S_{\mathrm{RR}}
=
-\frac{1}{2}
\int |F_{p+2}|^2
\] up to normalization. The propagator contraction between two static
parallel brane currents contains the metric factor \[
\eta_{00}\eta_{11}\cdots\eta_{pp}
=
-1.
\] There is exactly one time index and \(p\) spatial indices. Therefore
the static sign is the same vector-like sign as electromagnetic
exchange, not the scalar-like sign.

Integrating out \(C_{p+1}\) gives a source-source interaction. For two
parallel branes with the same orientation, the R-R currents have the
same sign: \[
\mu_1\mu_2>0.
\] The resulting potential has the vector-like sign \[
V_{\mathrm{RR}}(r)
=
+\mu_1\mu_2G_{9-p}(r)
\] up to a positive normalization. Thus equal R-R charges repel.

Why is this vector-like rather than scalar-like? The important physical
point is that a \(p\)-form gauge field couples to an oriented conserved
brane current, just as a vector field couples to an oriented particle
current. Reversing the brane orientation reverses the R-R charge: \[
\mu_p\mapsto-\mu_p.
\] So a brane and an anti-brane have opposite R-R charges. Consequently:
\[
\text{brane-brane:}\quad \mu_1\mu_2>0
\quad\Longrightarrow\quad
V_{\mathrm{RR}}>0
\quad\Longrightarrow\quad
\text{repulsion},
\] whereas \[
\text{brane-antibrane:}\quad \mu_1\mu_2<0
\quad\Longrightarrow\quad
V_{\mathrm{RR}}<0
\quad\Longrightarrow\quad
\text{attraction}.
\]

One sometimes loosely says that antisymmetric tensor exchange is
``odd-spin-like'' in its force sign. More precisely, a R-R form is a
gauge potential coupling to an oriented conserved current. The sign of
its static exchange between equal charges is the same as for
electromagnetic vector exchange: equal charges repel.

\subsubsection{NS-NS exchange between
D-branes}\label{ns-ns-exchange-between-d-branes}

The NS-NS sector contains the graviton \(g_{\mu\nu}\), dilaton \(\Phi\),
and Kalb-Ramond field \(B_{\mu\nu}\). For two identical static parallel
D\(p\)-branes with no worldvolume flux, the relevant universal
attractive contribution comes from the tension coupling in the DBI
action: \[
S_{\mathrm{DBI}}
=
-T_p
\int d^{p+1}\xi\,
e^{-\Phi}
\sqrt{-\det P[g+B]}.
\] Expanding around flat space gives couplings to the graviton and
dilaton proportional to \(T_p\). Since \[
T_p>0
\] for both a brane and an anti-brane, the NS-NS force is attractive for
both brane-brane and brane-antibrane pairs: \[
V_{\mathrm{NSNS}}(r)
=
-T_1T_2G_{9-p}(r)
\] up to a positive normalization and the correct tensor coefficient.

For two identical BPS D\(p\)-branes, \[
T_1=T_2=T_p,
\qquad
\mu_1=\mu_2=\mu_p,
\] and supersymmetry fixes \[
|\mu_p|=T_p
\] in matching units. The total static potential is therefore \[
V_{\mathrm{total}}(r)
=
V_{\mathrm{NSNS}}(r)+V_{\mathrm{RR}}(r)
=
-T_p^2G_{9-p}(r)
+
\mu_p^2G_{9-p}(r)
=
0.
\] This is the no-force condition for parallel BPS D-branes.

For a brane-antibrane pair, \[
\mu_1\mu_2<0,
\] so the R-R term becomes attractive: \[
V_{\mathrm{RR}}(r)<0.
\] The NS-NS term is also attractive, so the forces add rather than
cancel: \[
V_{\mathrm{brane\text{-}antibrane}}(r)<0.
\]

\subsubsection{Checks}\label{checks-4}

For ordinary electromagnetism in four dimensions, \[
G_3(r)=\frac{1}{4\pi r},
\] so \[
V_A(r)=\frac{q_1q_2}{4\pi r}.
\] Equal charges repel because \(V_A>0\) and decreases as \(r\)
increases.

For Newtonian gravity, \[
V_h(r)=-\frac{G_Nm_1m_2}{r}.
\] Positive masses attract because \(V_h<0\) and becomes more negative
as \(r\) decreases.

For D-branes, the schematic long-distance potentials are \[
V_{\mathrm{NSNS}}(r)\sim -T_p^2G_{9-p}(r),
\qquad
V_{\mathrm{RR}}(r)\sim +\mu_p^2G_{9-p}(r).
\] For BPS branes, \[
T_p^2=\mu_p^2,
\] so they cancel.

\subsubsection{Result}\label{result-8}

For static exchange of a spin-\(s\) field between sources with charges
\(q_1,q_2\), \[
V_s(r)\sim (-1)^{s+1}q_1q_2G(r).
\] Therefore: \[
s\ \text{even}
\quad\Longrightarrow\quad
V_s(r)\sim -q_1q_2G(r),
\] so equal charges attract and opposite charges repel.

For \[
s\ \text{odd},
\] one has \[
V_s(r)\sim +q_1q_2G(r),
\] so equal charges repel and opposite charges attract.

For D\(p\)-branes, \[
S_{\mathrm{DBI}}
=
-T_p\int e^{-\Phi}\sqrt{-\det P[g+B]},
\qquad
S_{\mathrm{WZ}}
=
\mu_p\int C_{p+1}.
\] The NS-NS exchange is attractive: \[
V_{\mathrm{NSNS}}(r)\sim -T_1T_2G_{9-p}(r).
\] The R-R form exchange is repulsive for equal R-R charges: \[
V_{\mathrm{RR}}(r)\sim +\mu_1\mu_2G_{9-p}(r).
\] For two identical BPS D\(p\)-branes, \[
T_p=|\mu_p|,
\] so \[
\boxed{
V_{\mathrm{NSNS}}(r)+V_{\mathrm{RR}}(r)=0.
}
\]

\section{Q6: Type II torus partition function and supersymmetric
cancellation}\label{q6-type-ii-torus-partition-function-and-supersymmetric-cancellation}

Consider the type II closed superstring one-loop partition function
\(Z(\tau,\bar{\tau})\) on a torus with moduli \(\tau\). The explicit
computation shows that \[
Z(\tau,\bar{\tau})=\frac{1}{4} \frac{1}{(\mathrm{Im}\tau)^4} \frac{1}{|\eta|^{24}}|\vartheta_{3}^4-\vartheta_{4}^4-\vartheta_{2}^4|^2=0
\] It is said that this is a consequence of a supersymmetric spectrum:
the contributions from space-time bosons and fermions cancel. Why?
Suppose we have a supersymmetric space-time theory, write down its
one-loop partition function and show this result.

\begin{quote}
\textbf{Textbook location and related material.} Chapter 9, Section 9.1
(pp.~223--232): the type-II torus amplitude, GSO sum, and Jacobi
identity. Section 8.3 supplies the supersymmetric spectrum whose
level-by-level degeneracy explains the cancellation.
\end{quote}

\subsection{Answer}\label{answer-9}

\subsubsection{Core point}\label{core-point-9}

The question is why the type II torus vacuum amplitude \[
Z(\tau,\bar\tau)
=
\frac{1}{4}
\frac{1}{(\operatorname{Im}\tau)^4}
\frac{1}{|\eta|^{24}}
\left|
\vartheta_3^4-\vartheta_4^4-\vartheta_2^4
\right|^2
\] vanishes, and how this is related to spacetime supersymmetry.

The key point is that the torus vacuum amplitude is not an ordinary
thermal partition function \[
\operatorname{Tr}e^{-\beta H}.
\] It is a one-loop vacuum functional. In spacetime field theory, a
one-loop vacuum diagram receives a plus sign from a bosonic loop and a
minus sign from a fermionic loop. Therefore it is a supertrace: \[
\operatorname{Str}e^{-tH}
=
\operatorname{Tr}_{\mathcal H_B}e^{-tH}
-
\operatorname{Tr}_{\mathcal H_F}e^{-tH}.
\] If spacetime supersymmetry is unbroken, then bosonic and fermionic
physical states pair with the same mass. Thus the supertrace cancels
level by level: \[
\operatorname{Str}e^{-tH}=0.
\] In type II string theory, the same statement is encoded
worldsheet-wise by the Jacobi abstruse identity: \[
\vartheta_3^4(\tau)-\vartheta_4^4(\tau)-\vartheta_2^4(\tau)=0.
\]

\subsubsection{Ordinary partition function versus vacuum
amplitude}\label{ordinary-partition-function-versus-vacuum-amplitude}

First separate two objects that are often both called ``partition
functions.''

An ordinary thermal partition function is \[
Z_{\mathrm{thermal}}(\beta)
=
\operatorname{Tr}_{\mathcal H}
e^{-\beta H}.
\] This does not vanish in a supersymmetric theory. Bosonic and
fermionic states both contribute positively: \[
Z_{\mathrm{thermal}}(\beta)
=
\sum_{B}e^{-\beta E_B}
+
\sum_{F}e^{-\beta E_F}.
\] Supersymmetry makes the bosonic and fermionic energies equal, but it
does not put a minus sign between them in the ordinary thermal trace.

By contrast, the one-loop vacuum amplitude is the logarithm of the
vacuum functional. In a spacetime QFT it has the schematic form \[
\Gamma^{(1)}
=
\frac{1}{2}\operatorname{Tr}_B\log \Delta_B
-
\frac{1}{2}\operatorname{Tr}_F\log \Delta_F.
\] The minus sign is the usual minus sign from a closed fermion loop,
equivalently from the Grassmann functional integral: \[
\int D\phi\,e^{-\frac{1}{2}\phi\Delta_B\phi}
\propto
(\det\Delta_B)^{-1/2},
\] whereas \[
\int D\psi\,e^{-\psi\Delta_F\psi}
\propto
\det\Delta_F.
\] Taking logarithms gives opposite signs for bosons and fermions.

Using the Schwinger representation \[
\log A
=
-
\int_0^\infty\frac{dt}{t}e^{-tA}
\] up to an irrelevant constant, the one-loop vacuum amplitude becomes
schematically \[
\Gamma^{(1)}
\sim
-
\frac{1}{2}
\int_0^\infty\frac{dt}{t}
\operatorname{Str}
e^{-t\Delta}.
\] This is why the relevant object is a supertrace, not an ordinary
trace: \[
\operatorname{Str}_{\mathcal H}(\cdots)
=
\operatorname{Tr}_{\mathcal H}
\left((-1)^{F_{\mathrm{st}}}(\cdots)\right).
\]

\subsubsection{Cancellation in a supersymmetric
spectrum}\label{cancellation-in-a-supersymmetric-spectrum}

Let the Hilbert space decompose into bosonic and fermionic states: \[
\mathcal H
=
\mathcal H_B\oplus\mathcal H_F.
\] Suppose supersymmetry is unbroken. Then for every massive bosonic
state \(|B\rangle\) there is a fermionic state \[
|F\rangle
\sim
Q|B\rangle
\] with the same energy or mass. This follows because the supercharges
commute with the Hamiltonian: \[
[H,Q]=0.
\] Thus if \[
H|B\rangle=E|B\rangle,
\] then \[
H(Q|B\rangle)
=
QH|B\rangle
=
E(Q|B\rangle).
\] So supersymmetry pairs bosons and fermions at equal energy.

Now compute the supertrace: \[
\operatorname{Str}e^{-tH}
=
\sum_{B}e^{-tE_B}
-
\sum_{F}e^{-tE_F}.
\] If the spectrum is paired with equal degeneracies at every energy, \[
d_B(E)=d_F(E),
\] then \[
\operatorname{Str}e^{-tH}
=
\sum_E
\left(d_B(E)-d_F(E)\right)e^{-tE}
=
0.
\] This is the field-theory version of the type II torus cancellation.

For a relativistic field theory in \(D\) noncompact dimensions, the
one-loop vacuum energy can be written more explicitly as \[
\Lambda^{(1)}
=
\frac{1}{2}
\sum_i
(-1)^{F_i}
d_i
\int\frac{d^Dk_E}{(2\pi)^D}
\log(k_E^2+m_i^2).
\] Using the Schwinger parameter, \[
\log(k_E^2+m_i^2)
=
-
\int_0^\infty\frac{dt}{t}
e^{-t(k_E^2+m_i^2)},
\] and integrating over \(k_E\), \[
\int\frac{d^Dk_E}{(2\pi)^D}e^{-tk_E^2}
=
\frac{1}{(4\pi t)^{D/2}},
\] one obtains \[
\Lambda^{(1)}
=
-
\frac{1}{2}
\int_0^\infty\frac{dt}{t}
\frac{1}{(4\pi t)^{D/2}}
\sum_i
(-1)^{F_i}d_i e^{-tm_i^2}.
\] The sum \[
\sum_i(-1)^{F_i}d_i e^{-tm_i^2}
\] is exactly a mass-weighted supertrace. If supersymmetry pairs all
bosonic and fermionic states at the same mass, then \[
\sum_i(-1)^{F_i}d_i e^{-tm_i^2}=0
\] for every \(t\). Hence \[
\Lambda^{(1)}=0.
\]

\subsubsection{String theory version}\label{string-theory-version}

In closed string theory, the torus amplitude is the one-loop vacuum
amplitude: \[
\Lambda_{\mathrm{string}}^{(1)}
=
\int_{\mathcal F}
\frac{d^2\tau}{(\operatorname{Im}\tau)^2}
Z(\tau,\bar\tau),
\] where \(\mathcal F\) is the fundamental domain of
\(SL(2,\mathbb Z)\).

The integrand can be written as a trace over the closed-string Hilbert
space: \[
Z(\tau,\bar\tau)
=
\operatorname{Tr}_{\mathcal H_{\mathrm{closed}}}
\left[
(-1)^{F_{\mathrm{st}}}
q^{L_0-c/24}
\bar q^{\bar L_0-\bar c/24}
\right],
\] with \[
q=e^{2\pi i\tau},
\] after the GSO sum is included. This expression is schematic because
in the worldsheet derivation one sums over spin structures rather than
inserting a single explicit spacetime \((-1)^{F_{\mathrm{st}}}\) from
the start. But after GSO projection, the result is exactly the spacetime
supertrace over physical bosonic and fermionic states.

The type II answer is \[
Z(\tau,\bar\tau)
=
\frac{1}{4}
\frac{1}{(\operatorname{Im}\tau)^4}
\frac{1}{|\eta(\tau)|^{24}}
\left|
\vartheta_3^4(\tau)
-
\vartheta_4^4(\tau)
-
\vartheta_2^4(\tau)
\right|^2.
\] The three terms have a spacetime interpretation:

\[
\vartheta_3^4
\quad\text{comes from NS states without }(-1)^F,
\] \[
\vartheta_4^4
\quad\text{comes from NS states with }(-1)^F,
\] and \[
\vartheta_2^4
\quad\text{comes from R states}.
\] The relative signs are the GSO phases. They are chosen so that the
physical spectrum is spacetime supersymmetric.

The Jacobi abstruse identity says \[
\vartheta_3^4(\tau)
-
\vartheta_4^4(\tau)
-
\vartheta_2^4(\tau)
=
0.
\] Therefore \[
Z(\tau,\bar\tau)=0
\] pointwise in \(\tau\), before even integrating over the modular
fundamental domain.

\subsubsection{Level-by-level
interpretation}\label{level-by-level-interpretation}

The identity \[
\vartheta_3^4-\vartheta_4^4-\vartheta_2^4=0
\] is stronger than saying only the total integrated amplitude vanishes.
It says that the generating function for boson-minus-fermion
degeneracies vanishes.

Schematically, one may expand the left-moving contribution as \[
\frac{1}{\eta^{12}}
\left(
\vartheta_3^4-\vartheta_4^4-\vartheta_2^4
\right)
=
\sum_N
\left(d_B(N)-d_F(N)\right)q^N.
\] The Jacobi identity says \[
d_B(N)-d_F(N)=0
\] at every oscillator level \(N\). The right-moving sector satisfies
the same identity. For type II closed strings, the full closed-string
spectrum is built from left and right movers with the level-matching
constraint \[
L_0=\bar L_0.
\] Spacetime supersymmetry then pairs every physical bosonic state with
a physical fermionic state of the same mass.

Thus the worldsheet expression \[
\vartheta_3^4-\vartheta_4^4-\vartheta_2^4
\] is the character-level implementation of \[
\operatorname{Str}e^{-tM^2}=0.
\]

\subsubsection{Why this is not a contradiction with nonzero thermal free
energy}\label{why-this-is-not-a-contradiction-with-nonzero-thermal-free-energy}

At finite temperature, spacetime fermions are antiperiodic around the
Euclidean time circle. This explicitly breaks spacetime supersymmetry.
Then the thermal partition function is \[
Z_{\mathrm{thermal}}(\beta)
=
\operatorname{Tr}e^{-\beta H},
\] not \[
\operatorname{Str}e^{-\beta H}.
\] Therefore bosonic and fermionic contributions add instead of
canceling. This is why a supersymmetric theory can have \[
\Lambda^{(1)}=0
\] for the zero-temperature vacuum amplitude but a nonzero
finite-temperature free energy.

\subsubsection{Checks}\label{checks-5}

For a simple supersymmetric harmonic spectrum with one bosonic state and
one fermionic state at the same energy \(E\), \[
Z_{\mathrm{thermal}}
=
e^{-\beta E}+e^{-\beta E}
=
2e^{-\beta E},
\] which is nonzero. But the one-loop vacuum supertrace gives \[
\operatorname{Str}e^{-\beta H}
=
e^{-\beta E}-e^{-\beta E}
=
0.
\] This is the same distinction appearing in the type II torus
amplitude.

For type II strings, the pointwise vanishing \[
Z(\tau,\bar\tau)=0
\] comes from \[
\vartheta_3^4-\vartheta_4^4-\vartheta_2^4=0.
\] If supersymmetry is broken, for example by thermal boundary
conditions or by a non-supersymmetric GSO projection, the spin-structure
sum is modified and this identity no longer forces the integrand to
vanish.

\subsubsection{Result}\label{result-9}

The ordinary thermal partition function is \[
Z_{\mathrm{thermal}}(\beta)
=
\operatorname{Tr}e^{-\beta H},
\] so bosons and fermions both contribute with positive sign.

The one-loop vacuum amplitude is instead a supertrace: \[
\Gamma^{(1)}
\sim
\int_0^\infty\frac{dt}{t}
\operatorname{Str}e^{-tH},
\qquad
\operatorname{Str}(\cdots)
=
\operatorname{Tr}\left((-1)^{F_{\mathrm{st}}}(\cdots)\right).
\] In an unbroken supersymmetric theory, \[
\mathcal H_B
\simeq
\mathcal H_F
\] at equal energy or mass, hence \[
\operatorname{Str}e^{-tH}
=
0.
\] For the type II closed superstring, \[
Z(\tau,\bar\tau)
=
\frac{1}{4}
\frac{1}{(\operatorname{Im}\tau)^4}
\frac{1}{|\eta|^{24}}
\left|
\vartheta_3^4-\vartheta_4^4-\vartheta_2^4
\right|^2.
\] The Jacobi identity \[
\vartheta_3^4-\vartheta_4^4-\vartheta_2^4=0
\] is the worldsheet character form of spacetime Bose-Fermi
cancellation. Therefore \[
\boxed{
Z(\tau,\bar\tau)=0
}
\] for type II strings with unbroken spacetime supersymmetry.

\section{Q7: Möbius amplitude and the boundary-crosscap cross
term}\label{q7-muxf6bius-amplitude-and-the-boundary-crosscap-cross-term}

Consider the type I string. The sum of the cylinder, the two Mobius and
the Klein bottle diagrams must be a perfect square such that the sum of
the disc and the crosscap tadpoles cancel. This is the requirement of no
spacetime R-R 10-form charge.

These amplitudes can be computed as follows: 1. The loop-channel annulus
amplitude for open strings stretched between a stack of \(N\) D9-branes:
\[
\mathcal{A}=N^2 \frac{V_{10}}{(4\pi^2\alpha')^5}\int_{0}^\infty \frac{dt}{2t} \frac{1}{(2t)^5} \left( \frac{{\vartheta_{3}^4-\vartheta_{4}^4-\vartheta_{2}^4}}{2\eta^{12}} \right)(it)
\] The transformation to tree-channel via \(t=\frac{1}{2l}\) gives (here
\(q^0=1\) mean 0th square) \[
\tilde{\mathcal{A}}=N^2 \frac{V_{10}}{(4\pi^2\alpha')^5}\int_{0}^\infty dl 2^{-6}(1_{\mathrm{NS}}-1_{\mathrm{R}})[16q^0+\mathcal{O}(e^{-2\pi l})]
\] 2. The O9-O9 amplitude can be viewed as the type IIB superstring
amplitude on a Klein bottle \[
\mathcal{K}=\frac{V_{10}}{(4\pi^2\alpha')^5}\int_{0}^\infty \frac{dt}{2t} \frac{1}{t^5} \left( \frac{{\vartheta_{3}^4-\vartheta_{4}^4-\vartheta_{2}^4}}{2\eta^{12}} \right)(2it)
\] Transform to tree-channel via \(t=\frac{1}{4l}\) we have \[
\tilde{\mathcal{K}}=\frac{V_{10}}{(4\pi^2\alpha')^5}\int_{0}^\infty dl 2^4\left( \frac{{\vartheta_{3}^4-\vartheta_{4}^4-\vartheta_{2}^4}}{2\eta^{12}} \right)(2il)=\frac{V_{10}}{(4\pi^2\alpha')^5}\int_{0}^\infty dl 2^4(1_{\mathrm{NS}}-1_{\mathrm{R}})[16q^0+\mathcal{O}(e^{-2\pi l})]
\] 3. The N D9-O9 amplitude can be computed in open string sector
(Chan-Paton factors supposed to be SO (N)) as Mobius strip one-loop
amplitude \[
\mathcal{M}=-N\frac{V_{10}}{(4\pi^2\alpha')^5}\int_{0}^\infty \frac{dt}{2t} \frac{1}{(2t)^5} \left( \frac{{\vartheta_{3}^4-\vartheta_{4}^4-\vartheta_{2}^4}}{2\eta^{12}} \right)\left( it+\frac{1}{2} \right)
\] Transform to tree-channel via \(t=\frac{1}{8l}\) we have \[
\tilde{\mathcal{M}}=-N \frac{V_{10}}{(4\pi^2\alpha')^5}\int_{0}^\infty dl (1_{\mathrm{NS}}-1_{\mathrm{R}})[16q^0+\mathcal{O}(e^{-2\pi l})]
\]

Combine these we find the factor \[
2^{-6}(N^2-2^6N+2^{10})=2^{-6}(N-2^5)^2
\] The cancellation requires that \(N=32\) i.e.~type I \(SO(32)\)
superstring is consistent.

Question: the square of the sum of the disc and the crosscap tadpoles
should be \[
(\text{disc}+\text{crosscap})^2=(\text{disc-disc})+2(\text{disc-crosscap})+(\text{crosscap-crosscap})
\] Why in the above computation we only sum one copy of the Mobius
amplitude, not two?

\begin{quote}
\textbf{Textbook location and related material.} Chapter 9, Section 9.4
(pp.~241--251): type-I annulus, Klein-bottle, and Möbius amplitudes and
R--R tadpole factorization. Section 6.7 (pp.~166--171) gives the general
crosscap and Möbius formalism.
\end{quote}

\subsection{Answer}\label{answer-10}

\subsubsection{Core point}\label{core-point-10}

The question is why the square \[
(\text{disc}+\text{crosscap})^2
=
(\text{disc-disc})
+
2(\text{disc-crosscap})
+
(\text{crosscap-crosscap})
\] seems to contain two disc-crosscap terms, but the one-loop type I
computation contains only one Möbius amplitude.

The answer is that the tree-channel vacuum amplitude has an overall
factor \(\frac{1}{2}\) in the closed-string exchange between tadpole
sources. Schematically, \[
\Gamma_{\mathrm{tree}}
=
\frac{1}{2}
\langle B+C|\Delta|B+C\rangle.
\] Expanding this gives \[
\Gamma_{\mathrm{tree}}
=
\frac{1}{2}\langle B|\Delta|B\rangle
+
\langle B|\Delta|C\rangle
+
\frac{1}{2}\langle C|\Delta|C\rangle.
\] Thus: \[
\tilde{\mathcal A}
=
\frac{1}{2}\langle B|\Delta|B\rangle,
\qquad
\tilde{\mathcal M}
=
\langle B|\Delta|C\rangle,
\qquad
\tilde{\mathcal K}
=
\frac{1}{2}\langle C|\Delta|C\rangle.
\] The single Möbius amplitude is already the full cross term. If one
added two copies of the Möbius amplitude, one would double-count the
boundary-crosscap exchange.

\subsubsection{Factorization of tadpoles in tree
channel}\label{factorization-of-tadpoles-in-tree-channel}

Let \(|B\rangle\) be the boundary state sourced by the D9-branes and
\(|C\rangle\) the crosscap state sourced by the O9-plane. The
tree-channel diagrams represent closed-string propagation between these
sources: \[
\tilde{\mathcal A}
\sim
\text{D9-D9 exchange},
\] \[
\tilde{\mathcal M}
\sim
\text{D9-O9 exchange},
\] and \[
\tilde{\mathcal K}
\sim
\text{O9-O9 exchange}.
\]

The total source for the massless R-R field is the sum of the boundary
and crosscap sources: \[
|J\rangle
=
|B\rangle+|C\rangle.
\] The closed-string exchange energy is quadratic in the source, but
with the standard symmetry factor: \[
\Gamma_{\mathrm{tree}}
=
\frac{1}{2}
\langle J|\Delta|J\rangle.
\] Therefore \[
\begin{aligned}
\Gamma_{\mathrm{tree}}
&=
\frac{1}{2}
\left(
\langle B|\Delta|B\rangle
+
\langle B|\Delta|C\rangle
+
\langle C|\Delta|B\rangle
+
\langle C|\Delta|C\rangle
\right)\\
&=
\frac{1}{2}\langle B|\Delta|B\rangle
+
\langle B|\Delta|C\rangle
+
\frac{1}{2}\langle C|\Delta|C\rangle.
\end{aligned}
\] In the second line we used \[
\langle B|\Delta|C\rangle
=
\langle C|\Delta|B\rangle.
\] So the algebraic factor of \(2\) in the cross term is precisely
canceled by the overall \(\frac{1}{2}\) in the quadratic exchange
functional.

This is the same basic reason that, for ordinary fields, integrating out
a field sourced by \(J=J_1+J_2\) gives \[
\frac{1}{2}(J_1+J_2)\Delta(J_1+J_2)
=
\frac{1}{2}J_1\Delta J_1
+
J_1\Delta J_2
+
\frac{1}{2}J_2\Delta J_2.
\] There is only one interaction term \(J_1\Delta J_2\), not two
separately counted ones.

\subsubsection{Relation to annulus, Möbius strip, and Klein
bottle}\label{relation-to-annulus-muxf6bius-strip-and-klein-bottle}

The three unoriented one-loop worldsheets are:

\[
\mathcal A:\quad
\text{annulus},
\] \[
\mathcal M:\quad
\text{Möbius strip},
\] and \[
\mathcal K:\quad
\text{Klein bottle}.
\] After modular transformation to tree channel, they are interpreted as
closed-string propagation: \[
\tilde{\mathcal A}
=
\frac{1}{2}\langle B|\Delta|B\rangle,
\] \[
\tilde{\mathcal M}
=
\langle B|\Delta|C\rangle,
\] \[
\tilde{\mathcal K}
=
\frac{1}{2}\langle C|\Delta|C\rangle.
\] The factors of \(\frac{1}{2}\) for the annulus and Klein bottle are
not arbitrary. They are symmetry factors for diagrams with two identical
ends:

\[
\text{annulus: two boundary ends},
\] \[
\text{Klein bottle: two crosscap ends}.
\] By contrast, the Möbius strip has one boundary and one crosscap.
These two ends are different types of objects, so there is no additional
identical-end symmetry factor. The Möbius strip is precisely the mixed
boundary-crosscap exchange.

Equivalently, in the loop-channel unoriented open-string description,
the Möbius strip is already the worldsheet with one boundary and one
crosscap. There are not two different Möbius worldsheets corresponding
to \[
\langle B|\Delta|C\rangle
\qquad
\text{and}
\qquad
\langle C|\Delta|B\rangle.
\] They are the same unoriented surface seen from opposite ends of the
closed-string propagator. Counting both would double-count the same
moduli region.

\subsubsection{Seeing the square in your
coefficients}\label{seeing-the-square-in-your-coefficients}

Your tree-channel massless coefficients are \[
\tilde{\mathcal A}
\sim
2^{-6}N^2,
\] \[
\tilde{\mathcal M}
\sim
-N,
\] and \[
\tilde{\mathcal K}
\sim
2^4.
\] Their sum is \[
2^{-6}N^2-N+2^4.
\] Rewriting everything with the same overall factor, \[
2^{-6}N^2-N+2^4
=
2^{-6}
\left(
N^2-2^6N+2^{10}
\right).
\] But \[
N^2-2^6N+2^{10}
=
(N-2^5)^2.
\] Thus \[
\tilde{\mathcal A}
+
\tilde{\mathcal M}
+
\tilde{\mathcal K}
\sim
2^{-6}(N-32)^2.
\]

To see the cross term explicitly, write \[
2^{-6}(N-32)^2
=
\left(\frac{N}{8}-4\right)^2.
\] Then \[
\left(\frac{N}{8}-4\right)^2
=
\left(\frac{N}{8}\right)^2
-
2\left(\frac{N}{8}\right)(4)
+
4^2.
\] These are \[
\left(\frac{N}{8}\right)^2
=
2^{-6}N^2
\quad
\text{from } \tilde{\mathcal A},
\] \[
-
2\left(\frac{N}{8}\right)(4)
=
-N
\quad
\text{from } \tilde{\mathcal M},
\] and \[
4^2
=
2^4
\quad
\text{from } \tilde{\mathcal K}.
\] So the single Möbius amplitude already carries the full coefficient
\[
2(\text{disc})(\text{crosscap})
\] inside the square. That is exactly why adding a second Möbius
amplitude would spoil the square.

\subsubsection{\texorpdfstring{Tadpole cancellation and
\(N=32\)}{Tadpole cancellation and N=32}}\label{tadpole-cancellation-and-n32}

The R-R tadpole is proportional to the total R-R source: \[
Q_{\mathrm{total}}
=
Q_{\mathrm{D9}}+Q_{\mathrm{O9}}.
\] The tree-channel massless R-R exchange is proportional to \[
Q_{\mathrm{total}}^2.
\] Your coefficient \[
2^{-6}(N-32)^2
\] therefore means \[
Q_{\mathrm{D9}}\propto N,
\qquad
Q_{\mathrm{O9}}\propto -32.
\] Tadpole cancellation requires \[
Q_{\mathrm{total}}=0,
\] so \[
N-32=0.
\] Hence \[
\boxed{
N=32.
}
\] This is the usual type I result: the consistent supersymmetric
ten-dimensional type I string has gauge group \(SO(32)\).

\subsubsection{Checks}\label{checks-6}

If one incorrectly added two copies of the Möbius amplitude, the
massless coefficient would become \[
2^{-6}N^2-2N+2^4.
\] This is \[
2^{-6}
\left(
N^2-2^7N+2^{10}
\right),
\] which is not \[
2^{-6}(N-32)^2.
\] It would instead correspond to the wrong cross term and would not
factorize into the square of the total R-R tadpole with the correct
O9-plane charge.

The correct counting is: \[
\boxed{
\tilde{\mathcal A}+\tilde{\mathcal M}+\tilde{\mathcal K}
=
\frac{1}{2}\langle B+C|\Delta|B+C\rangle.
}
\] The Möbius strip appears once because it is the one mixed
boundary-crosscap surface.

\subsubsection{Result}\label{result-10}

The tree-channel factorization is \[
\Gamma_{\mathrm{tree}}
=
\frac{1}{2}
\langle B+C|\Delta|B+C\rangle.
\] Expanding, \[
\Gamma_{\mathrm{tree}}
=
\frac{1}{2}\langle B|\Delta|B\rangle
+
\langle B|\Delta|C\rangle
+
\frac{1}{2}\langle C|\Delta|C\rangle.
\] Thus \[
\tilde{\mathcal A}
=
\frac{1}{2}\langle B|\Delta|B\rangle,
\qquad
\tilde{\mathcal M}
=
\langle B|\Delta|C\rangle,
\qquad
\tilde{\mathcal K}
=
\frac{1}{2}\langle C|\Delta|C\rangle.
\] So one sums only one Möbius amplitude because that one amplitude is
already the full cross term.

For the massless tadpoles, \[
\tilde{\mathcal A}+\tilde{\mathcal M}+\tilde{\mathcal K}
\sim
2^{-6}N^2-N+2^4
=
2^{-6}(N-32)^2.
\] The tadpole vanishes only for \[
\boxed{
N=32.
}
\]

\section{Q8: Orientifold action on D-brane boundary
states}\label{q8-orientifold-action-on-d-brane-boundary-states}

Previously we have talked about the Dp-brane state \[
|{\rm Dp},\eta\rangle=\int d k_{\perp} \exp \left[-\sum_{n>0}\frac{1}{n} \alpha^{i}_{-n} D_{ij} \overline{\alpha}^{j}_{-n} -i \eta\sum_{r>0}b^{ i}_{-r} D_{ij} \overline{b}^{j}_{-r} \right] e^{i k_{\perp}\cdot x_{0}^{ \perp}} |k_{\perp}\rangle ,
\] A complete Dp-brane states is \[
|\mathrm{Dp}\rangle=\frac{1}{\mathcal{N}_{\mathrm{Dp}}}(|\mathrm{Dp},\eta=1\rangle_{\mathrm{NS}}-|\mathrm{Dp},\eta=-1\rangle_{\mathrm{NS}}+i|\mathrm{Dp},\eta=1\rangle_{\mathrm{R}}+i|\mathrm{Dp},\eta=-1\rangle_{\mathrm{R}})
\] Under the orientifold projection \(\Omega\), in the resulting type I
string theory only R-R 2-form \(C_{2}\) left. Which means that only D9,
D5, and D1 branes survive. Can we see this from the action of \(\omega\)
on the Dp-brane state?

\begin{quote}
\textbf{Textbook location and related material.} Chapter 9, Sections
9.3--9.4 (pp.~235--251): GSO-projected D-brane boundary states and the
\(\Omega\) projection in type I. Table 10.3 and Section 10.6 relate the
surviving branes to toroidal orientifolds.
\end{quote}

\subsection{Answer}\label{answer-11}

\subsubsection{Core point}\label{core-point-11}

Yes. One can see the type I restriction \[
p=1,5,9
\] directly from the action of worldsheet orientation reversal
\(\Omega\) on the D\(p\) boundary state.

The BPS D\(p\) boundary state has the schematic form \[
|{\rm Dp}\rangle
=
|{\rm Dp}\rangle_{\mathrm{NSNS}}
+
|{\rm Dp}\rangle_{\mathrm{RR}},
\] where, with your convention, \[
|{\rm Dp}\rangle_{\mathrm{NSNS}}
\propto
|{\rm Dp},+\rangle_{\mathrm{NS}}
-
|{\rm Dp},-\rangle_{\mathrm{NS}},
\] and \[
|{\rm Dp}\rangle_{\mathrm{RR}}
\propto
i|{\rm Dp},+\rangle_{\mathrm{R}}
+
i|{\rm Dp},-\rangle_{\mathrm{R}}.
\] The orientifold projection keeps only \(\Omega\)-even closed-string
states and \(\Omega\)-invariant boundary-state sources. The NS-NS part
is \(\Omega\)-even: \[
\Omega |{\rm Dp}\rangle_{\mathrm{NSNS}}
=
|{\rm Dp}\rangle_{\mathrm{NSNS}}.
\] The R-R part transforms with a \(p\)-dependent sign: \[
\Omega |{\rm Dp}\rangle_{\mathrm{RR}}
=
(-1)^{(p+3)/2}
|{\rm Dp}\rangle_{\mathrm{RR}},
\qquad
p\ \text{odd}.
\] Therefore the R-R charged BPS D\(p\) boundary state survives the type
I orientifold projection only if \[
(-1)^{(p+3)/2}=+1.
\] For type IIB odd \(p\), \[
p=-1,1,3,5,7,9,
\] this gives \[
p=1,5,9.
\] These are precisely the BPS D-branes of type I string theory.

\subsubsection{\texorpdfstring{How \(\Omega\) acts on the oscillator
part}{How \textbackslash Omega acts on the oscillator part}}\label{how-omega-acts-on-the-oscillator-part}

Worldsheet parity exchanges left and right movers: \[
\Omega\,\alpha_n^\mu\,\Omega^{-1}
=
\bar\alpha_n^\mu,
\qquad
\Omega\,\bar\alpha_n^\mu\,\Omega^{-1}
=
\alpha_n^\mu,
\] and similarly for the closed-string fermion oscillators: \[
\Omega\,b_r^\mu\,\Omega^{-1}
=
\bar b_r^\mu,
\qquad
\Omega\,\bar b_r^\mu\,\Omega^{-1}
=
b_r^\mu.
\] The bosonic oscillator exponential in the D\(p\) boundary state is \[
\exp\left[
-\sum_{n>0}
\frac{1}{n}
\alpha_{-n}^iD_{ij}\bar\alpha_{-n}^j
\right].
\] Since the bosonic oscillators commute and \(D_{ij}\) is symmetric,
exchanging left and right movers leaves this part invariant: \[
\alpha_{-n}^iD_{ij}\bar\alpha_{-n}^j
\longmapsto
\bar\alpha_{-n}^iD_{ij}\alpha_{-n}^j
=
\alpha_{-n}^iD_{ij}\bar\alpha_{-n}^j.
\]

The fermionic nonzero-mode exponential is \[
\exp\left[
-i\eta
\sum_{r>0}
b_{-r}^iD_{ij}\bar b_{-r}^j
\right].
\] Under \(\Omega\), \[
b_{-r}^iD_{ij}\bar b_{-r}^j
\longmapsto
\bar b_{-r}^iD_{ij}b_{-r}^j.
\] Because fermionic oscillators anticommute, \[
\bar b_{-r}^iD_{ij}b_{-r}^j
=
-
b_{-r}^jD_{ji}\bar b_{-r}^i
=
-
b_{-r}^iD_{ij}\bar b_{-r}^j.
\] Therefore \[
-i\eta
\sum_{r>0}
b_{-r}D\bar b_{-r}
\quad
\longmapsto
\quad
+i\eta
\sum_{r>0}
b_{-r}D\bar b_{-r}.
\] This is the same as replacing \[
\eta\mapsto -\eta.
\] So the oscillator part of the boundary state naturally satisfies \[
\Omega:
|{\rm Dp},\eta\rangle
\longmapsto
\text{phase}\times
|{\rm Dp},-\eta\rangle.
\] The remaining question is the phase. In the NS-NS sector the phase is
fixed so that the GSO-projected NS-NS boundary state is invariant. In
the R-R sector the phase depends on \(p\) because of the R-R zero modes.

\subsubsection{NS-NS part}\label{ns-ns-part}

The GSO-projected NS-NS boundary state is \[
|{\rm Dp}\rangle_{\mathrm{NSNS}}
\propto
|{\rm Dp},+\rangle_{\mathrm{NS}}
-
|{\rm Dp},-\rangle_{\mathrm{NS}}.
\] With the standard ghost and superghost conventions, the action on the
unprojected NS-NS components is \[
\Omega |{\rm Dp},\eta\rangle_{\mathrm{NS}}
=
-
|{\rm Dp},-\eta\rangle_{\mathrm{NS}}.
\] Therefore \[
\begin{aligned}
\Omega |{\rm Dp}\rangle_{\mathrm{NSNS}}
&\propto
\Omega
\left(
|{\rm Dp},+\rangle_{\mathrm{NS}}
-
|{\rm Dp},-\rangle_{\mathrm{NS}}
\right)\\
&=
-
|{\rm Dp},-\rangle_{\mathrm{NS}}
+
|{\rm Dp},+\rangle_{\mathrm{NS}}\\
&=
|{\rm Dp},+\rangle_{\mathrm{NS}}
-
|{\rm Dp},-\rangle_{\mathrm{NS}}.
\end{aligned}
\] Thus \[
\boxed{
\Omega |{\rm Dp}\rangle_{\mathrm{NSNS}}
=
|{\rm Dp}\rangle_{\mathrm{NSNS}}.
}
\] This is why the tension coupling to the \(\Omega\)-even NS-NS fields,
such as \(g_{\mu\nu}\) and \(\Phi\), is not the obstruction. The
obstruction is the R-R charge.

\subsubsection{\texorpdfstring{R-R zero modes and the \(p\)-dependent
sign}{R-R zero modes and the p-dependent sign}}\label{r-r-zero-modes-and-the-p-dependent-sign}

The R-R ground states are spacetime spinors. The zero-mode part of the
R-R boundary state is usually written as a bispinor matrix involving the
gamma-matrix product along the brane worldvolume: \[
M_p(\eta)
\sim
C\Gamma^0\Gamma^1\cdots\Gamma^p
\left(1+i\eta\Gamma_{11}\right),
\] up to normalization and convention-dependent phases.

Worldsheet parity \(\Omega\) exchanges the left and right R-R
ground-state spinors. In bispinor language, this transposes the matrix
\(M_p(\eta)\). The transpose of a product of \(p+1\) gamma matrices
gives a \(p\)-dependent sign. Equivalently, this is the same sign that
appears in the \(\Omega\) transformation of R-R potentials: \[
\Omega:
C_n
\longmapsto
(-1)^{n/2+1}C_n
\qquad
\text{for even }n\text{ in type IIB}.
\] Since a D\(p\)-brane couples to \[
C_{p+1},
\] we set \[
n=p+1.
\] Hence \[
\Omega:
C_{p+1}
\longmapsto
(-1)^{(p+1)/2+1}C_{p+1}
=
(-1)^{(p+3)/2}C_{p+1}.
\] The R-R part of the D\(p\) boundary state has the same eigenvalue: \[
\boxed{
\Omega |{\rm Dp}\rangle_{\mathrm{RR}}
=
(-1)^{(p+3)/2}
|{\rm Dp}\rangle_{\mathrm{RR}}.
}
\] This is the boundary-state way of saying that the R-R potential
sourced by the brane is either kept or projected out by \(\Omega\).

\subsubsection{Which D-branes survive}\label{which-d-branes-survive}

Type IIB has BPS D\(p\)-branes with odd \(p\): \[
p=-1,1,3,5,7,9.
\] The R-R eigenvalue is \[
\lambda_p
=
(-1)^{(p+3)/2}.
\] Let us list it: \[
\begin{array}{c|c|c|c}
p & C_{p+1} & \lambda_p & \Omega\text{-even?}\\
\hline
-1 & C_0 & -1 & \text{no}\\
1 & C_2 & +1 & \text{yes}\\
3 & C_4 & -1 & \text{no}\\
5 & C_6 & +1 & \text{yes}\\
7 & C_8 & -1 & \text{no}\\
9 & C_{10} & +1 & \text{yes}
\end{array}
\] Therefore the \(\Omega\)-invariant R-R charged BPS branes are \[
\boxed{
\mathrm{D1},\quad \mathrm{D5},\quad \mathrm{D9}.
}
\]

This agrees with the spacetime field content. In type I, the propagating
R-R potential is often described as \[
C_2.
\] Its magnetic dual is \[
C_6,
\] so D5-branes are also allowed. The top-form \[
C_{10}
\] has no propagating field strength degree of freedom, but it is
crucial for the D9/O9 R-R tadpole. D9-branes couple electrically to this
top-form.

Thus the statement ``only \(C_2\) is left'' should be understood
together with its dual and the top-form tadpole sector: \[
C_2,\quad C_6,\quad C_{10}
\] are the R-R potentials relevant to \[
\mathrm{D1},\quad \mathrm{D5},\quad \mathrm{D9}.
\]

\subsubsection{Branes projected out versus branes mapped to non-BPS
objects}\label{branes-projected-out-versus-branes-mapped-to-non-bps-objects}

For \[
p=-1,3,7,
\] the R-R part is \(\Omega\)-odd: \[
\Omega |{\rm Dp}\rangle_{\mathrm{RR}}
=
-
|{\rm Dp}\rangle_{\mathrm{RR}}.
\] So the R-R charge carried by such a BPS D\(p\)-brane is projected out
in type I. The corresponding BPS R-R charged brane of type IIB therefore
does not survive as a BPS R-R charged object in type I.

This does not mean there are no possible non-BPS D-brane-like objects in
type I. There are indeed non-BPS branes in type I with more subtle
stability properties. But the simple R-R charged BPS boundary-state
argument selects precisely \[
p=1,5,9.
\]

\subsubsection{Checks}\label{checks-7}

For a D1-brane, \[
p=1,
\qquad
(-1)^{(p+3)/2}
=
(-1)^2
=
+1.
\] It couples to the \(\Omega\)-even R-R two-form \(C_2\), so it
survives.

For a D3-brane, \[
p=3,
\qquad
(-1)^{(p+3)/2}
=
(-1)^3
=
-1.
\] It couples to \(C_4\), which is \(\Omega\)-odd in type I, so the BPS
D3-brane does not survive the orientifold projection.

For a D9-brane, \[
p=9,
\qquad
(-1)^{(p+3)/2}
=
(-1)^6
=
+1.
\] It couples to \(C_{10}\) and is needed to cancel the O9-plane R-R
charge. This is the same D9-brane that led to the \(SO(32)\) tadpole
cancellation condition in Q7.

\subsubsection{Result}\label{result-11}

Worldsheet parity exchanges left and right movers: \[
\Omega b_r^\mu\Omega^{-1}=\bar b_r^\mu,
\qquad
\Omega\bar b_r^\mu\Omega^{-1}=b_r^\mu.
\] Because fermion oscillators anticommute, the fermionic gluing label
flips: \[
\Omega:
\eta\mapsto -\eta.
\] For the GSO-projected NS-NS state, \[
\Omega |{\rm Dp}\rangle_{\mathrm{NSNS}}
=
|{\rm Dp}\rangle_{\mathrm{NSNS}}.
\] For the R-R part, \[
\Omega |{\rm Dp}\rangle_{\mathrm{RR}}
=
(-1)^{(p+3)/2}
|{\rm Dp}\rangle_{\mathrm{RR}}.
\] Thus the R-R charged BPS D\(p\)-brane survives the type I orientifold
projection only when \[
(-1)^{(p+3)/2}=+1.
\] For odd type IIB \(p\), this gives \[
\boxed{
p=1,5,9.
}
\] Equivalently, type I keeps \[
C_2,\quad C_6,\quad C_{10},
\] corresponding to \[
\boxed{
\mathrm{D1},\quad \mathrm{D5},\quad \mathrm{D9}.
}
\]

\chapter{The Quantized Fermionic
String}\label{the-quantized-fermionic-string-1}

\section{Q1: The no-ghost theorem}\label{q1-the-no-ghost-theorem}

In the canonical quantization of the bosonic string and superstring, due
to the negative eigenvalue of \(\eta^{\mu \nu}\) there are negative norm
states (ghosts). One can prove a no-ghost theorem which states that the
negative norm states decouple in the critical dimension d for a
particular value of the normal ordering constant a. Explain the no-ghost
theorem for me in details.

\begin{quote}
\textbf{Textbook location and related material.} Chapter 8, Sections
8.1--8.4 (pp.~195--210), read together with Chapter 3, Sections 3.1--3.3
(pp.~35--53): covariant constraints, light-cone quantization, critical
dimension, and removal of negative-norm states. The DDF formulation is
supplementary to the book's presentation.
\end{quote}

\subsection{Answer}\label{answer-12}

\subsubsection{Core point}\label{core-point-12}

The issue is that old-covariant quantization keeps manifest Lorentz
covariance. Therefore one quantizes all target-space components \[
X^\mu(\tau,\sigma),
\qquad
\mu=0,1,\dots,D-1,
\] including the timelike component \(X^0\). Since \[
\eta^{00}<0,
\] the naive oscillator Fock space has an indefinite inner product. For
example, in the bosonic string, \[
[\alpha_m^\mu,\alpha_n^\nu]
=
m\delta_{m+n,0}\eta^{\mu\nu},
\] so \[
\|\alpha_{-1}^0|0;p\rangle\|^2
=
\langle 0;p|\alpha_1^0\alpha_{-1}^0|0;p\rangle
=
\eta^{00}
<0.
\] These negative-norm states are the covariant-quantization ghosts.

The no-ghost theorem says that these ghosts do not survive as physical
states, provided the string is quantized at the critical values of the
dimension and intercept. For the open bosonic string: \[
D=26,
\qquad
a=1.
\] For the open RNS superstring: \[
D=10,
\qquad
a_{\text{NS}}=\frac{1}{2},
\qquad
a_{\text{R}}=0.
\] More precisely, the theorem says that the physical state space modulo
null states is positive definite and is isomorphic to the transverse
light-cone Fock space. Symbolically, \[
\frac{\mathcal H_{\text{phys}}}{\mathcal H_{\text{null}}}
\cong
\mathcal H_{\text{transverse}}.
\] Thus negative-norm excitations are either absent from the physical
quotient or belong to null/spurious gauge sectors that decouple from
physical amplitudes.

\subsubsection{Why ghosts appear before imposing
constraints}\label{why-ghosts-appear-before-imposing-constraints}

For the open bosonic string, the oscillator expansion can be written
schematically as \[
X^\mu(\tau,\sigma)
=
x^\mu+2\alpha' p^\mu\tau
+i\sqrt{2\alpha'}\sum_{n\neq 0}\frac{1}{n}\alpha_n^\mu e^{-in\tau}\cos n\sigma.
\] Canonical quantization gives \[
[x^\mu,p^\nu]=i\eta^{\mu\nu},
\qquad
[\alpha_m^\mu,\alpha_n^\nu]
=
m\delta_{m+n,0}\eta^{\mu\nu}.
\] The Fock vacuum obeys \[
\alpha_n^\mu|0;p\rangle=0,
\qquad
n>0,
\] and excited states are produced by \(\alpha_{-n}^\mu\). Since the
target metric is indefinite, the Fock space inner product is also
indefinite.

At level one, a general state is \[
|\zeta;p\rangle
=
\zeta_\mu\alpha_{-1}^\mu|0;p\rangle.
\] Its norm is \[
\langle \zeta;p|\zeta;p\rangle
=
\zeta_\mu^*\zeta_\nu\eta^{\mu\nu}.
\] If \(\zeta^\mu\) is timelike, this norm is negative. So before
constraints, the covariant string Fock space definitely contains ghosts.

The purpose of the Virasoro constraints is to remove precisely the
unphysical longitudinal and timelike excitations.

\subsubsection{Physical-state conditions in the bosonic
string}\label{physical-state-conditions-in-the-bosonic-string}

Classically the Polyakov string has worldsheet diffeomorphism and Weyl
gauge symmetry. In conformal gauge, the residual constraints are \[
T_{\alpha\beta}=0.
\] In holomorphic language this becomes \[
T(z)=0,
\] or equivalently \[
L_n=0
\] on physical configurations.

Quantum mechanically, the Virasoro generators are \[
L_m
=
\frac{1}{2}\sum_{n\in\mathbb Z}:\alpha_{m-n}\cdot\alpha_n:,
\] and the zero mode is \[
L_0
=
\alpha' p^2+N
\] for the open string. Because of normal ordering, the mass-shell
condition is shifted: \[
(L_0-a)|\psi\rangle=0.
\] The old-covariant physical-state conditions are therefore \[
(L_0-a)|\psi\rangle=0,
\qquad
L_n|\psi\rangle=0
\quad
n>0.
\] The first condition gives the mass formula \[
\alpha' M^2
=
N-a,
\] where \(M^2=-p^2\).

The positive-mode constraints \(L_n|\psi\rangle=0\) with \(n>0\) are the
quantum remnants of \(T(z)=0\). They remove many timelike and
longitudinal states, but it is not obvious a priori that they remove all
negative-norm physical states. The no-ghost theorem is precisely the
nontrivial statement that they do, at the critical values.

\subsubsection{Spurious states and null
states}\label{spurious-states-and-null-states}

A key idea is the distinction between physical states, spurious states,
and null states.

A physical state \(|\phi\rangle\) obeys \[
L_n|\phi\rangle=0,
\qquad
n>0,
\] and the mass-shell condition \[
(L_0-a)|\phi\rangle=0.
\] A spurious state is a state of the form \[
|\chi\rangle_{\text{sp}}
=
\sum_{n>0}L_{-n}|\chi_n\rangle,
\] possibly with products of several negative Virasoro generators. Such
states are generated by gauge transformations on the worldsheet.

If \(|\phi\rangle\) is physical, then \[
\langle \phi|L_{-n}|\chi_n\rangle
=
\langle L_n\phi|\chi_n\rangle
=
0,
\qquad
n>0.
\] Therefore every spurious state has zero inner product with every
physical state. If a spurious state is itself physical, then it is a
null physical state: \[
|\chi\rangle_{\text{sp}}\in\mathcal H_{\text{phys}}
\quad
\Longrightarrow
\quad
\langle \phi|\chi\rangle_{\text{sp}}=0
\quad
\text{for all }|\phi\rangle\in\mathcal H_{\text{phys}}.
\] In particular, \[
\| |\chi\rangle_{\text{sp}}\|^2=0.
\] Such states decouple from amplitudes because amplitudes depend only
on inner products with physical external states, or equivalently on BRST
cohomology classes.

However, this argument alone only says that physical spurious states are
null. It does not yet prove that every negative-norm physical state is
spurious or null. That stronger statement is the content of the no-ghost
theorem.

\subsubsection{The first level as a
warm-up}\label{the-first-level-as-a-warm-up}

At level one in the open bosonic string, \[
|\zeta;p\rangle
=
\zeta_\mu\alpha_{-1}^\mu|0;p\rangle.
\] The mass-shell condition gives \[
\alpha' M^2
=
1-a.
\] At the critical intercept \(a=1\), the level-one state is massless:
\[
p^2=0.
\] The \(L_1\) constraint gives \[
0
=
L_1|\zeta;p\rangle
=
\sqrt{2\alpha'}(\zeta\cdot p)|0;p\rangle,
\] so \[
\zeta\cdot p=0.
\] Thus the polarization must be transverse to the null momentum
\(p^\mu\).

But for a null momentum, the condition \(\zeta\cdot p=0\) still permits
a polarization proportional to \(p^\mu\): \[
\zeta^\mu\sim p^\mu.
\] The corresponding state is \[
p_\mu\alpha_{-1}^\mu|0;p\rangle.
\] Using \(\alpha_0^\mu=\sqrt{2\alpha'}p^\mu\), one has \[
L_{-1}|0;p\rangle
=
\alpha_{-1}\cdot\alpha_0|0;p\rangle
=
\sqrt{2\alpha'}p_\mu\alpha_{-1}^\mu|0;p\rangle.
\] So the longitudinal polarization is a spurious state. Since
\(p^2=0\), it also has zero norm. The quotient removes it.

The remaining polarizations are transverse to both light-cone
directions. They form a vector of \(SO(D-2)\), so the level-one physical
quotient has \[
D-2
\] positive-norm states. This is exactly the light-cone answer.

This example already shows the mechanism:

\begin{enumerate}
\def\labelenumi{\arabic{enumi}.}
\tightlist
\item
  the Virasoro constraints impose transversality;
\item
  longitudinal polarizations are spurious;
\item
  null/spurious states are quotiented out;
\item
  only transverse positive-norm polarizations remain.
\end{enumerate}

The no-ghost theorem says this continues to hold at every oscillator
level.

\subsubsection{Why the critical values
matter}\label{why-the-critical-values-matter}

The Virasoro algebra of the bosonic matter oscillators is \[
[L_m,L_n]
=
(m-n)L_{m+n}
+
\frac{D}{12}(m^3-m)\delta_{m+n,0}.
\] The central term is proportional to \(D\) because each target-space
coordinate \(X^\mu\) contributes \(c=1\) to the worldsheet central
charge.

In a purely classical constraint algebra, the constraints are first
class: \[
\{L_m,L_n\}
\sim
(m-n)L_{m+n}.
\] Quantum mechanically, the central term is an anomaly. It obstructs
treating all Virasoro constraints as exact gauge constraints unless it
is canceled by the ghost contribution in BRST quantization, or
equivalently unless the old-covariant theory has the special intercept
and dimension for which Lorentz covariance and positivity are
compatible.

For the bosonic string, the reparametrization ghosts \((b,c)\) have
central charge \[
c_{bc}=-26.
\] The matter central charge is \[
c_X=D.
\] Thus BRST nilpotency requires \[
c_X+c_{bc}=0,
\] or \[
D-26=0.
\] So \[
D=26.
\] The same computation also fixes the open-string intercept to \[
a=1.
\] In old-covariant language, \(D=26\) and \(a=1\) are the values for
which the Lorentz algebra closes without anomalous terms and the
spurious physical states have the right null structure to decouple
consistently.

If \(D\) and \(a\) are not at the allowed values, the constraints no
longer remove the indefinite directions in the correct way.
Equivalently, BRST nilpotency fails or Lorentz invariance fails. Then
one cannot consistently identify physical states with a positive
transverse Hilbert space.

There is a slightly more general mathematical form of the bosonic
no-ghost theorem: for the open bosonic string, the physical state space
has no negative-norm states for \(D\leq 26\) and \(a\leq 1\) under
suitable assumptions. But the fully consistent critical bosonic string,
with Lorentz invariance and BRST nilpotency, is the special case \[
D=26,
\qquad
a=1.
\] This is the case usually meant in string theory.

\subsubsection{DDF operators and the precise structure of the
theorem}\label{ddf-operators-and-the-precise-structure-of-the-theorem}

The clean constructive proof uses DDF operators. The idea is to build
physical excitations that are manifestly transverse, while still working
in covariant quantization.

Choose a null vector \(k^\mu\) satisfying \[
k^2=0.
\] For a momentum \(p^\mu\), one chooses the normalization of \(k^\mu\)
so that the vertex operators below have the correct conformal weight. A
standard open-string choice is \[
2\alpha' p\cdot k=1.
\] Let \(i=1,\dots,D-2\) label directions transverse to both \(p^\mu\)
and \(k^\mu\). The DDF oscillator is defined by a contour integral of a
dimension-one current: \[
A_n^i
=
\sqrt{\frac{2}{\alpha'}}
\oint\frac{dz}{2\pi i}\,
\partial X^i(z)e^{in k\cdot X(z)}.
\] The current has conformal weight one because \(k^2=0\). Therefore the
contour integral commutes with the positive Virasoro constraints: \[
[L_m,A_n^i]=0,
\qquad
m>0.
\] More precisely, \(A_n^i\) maps physical states to physical states,
shifting the spacetime momentum by \(nk^\mu\) while increasing the
oscillator level.

The DDF operators obey the ordinary transverse oscillator algebra \[
[A_m^i,A_n^j]
=
m\delta^{ij}\delta_{m+n,0}.
\] This algebra has positive definite norm because the indices \(i,j\)
are transverse spacelike indices.

Starting from a suitable ground state \(|p\rangle\), one constructs
states \[
A_{-n_1}^{i_1}A_{-n_2}^{i_2}\cdots A_{-n_s}^{i_s}|p\rangle,
\qquad
n_r>0.
\] These states are physical and have the same counting and inner
products as light-cone transverse oscillator states.

The DDF form of the no-ghost theorem is: \[
\boxed{
\text{Every physical state is equivalent, modulo null states, to a DDF state.}
}
\] Equivalently, \[
\mathcal H_{\text{phys}}
=
\mathcal H_{\text{DDF}}
\oplus
\mathcal H_{\text{null}},
\] with \[
\mathcal H_{\text{DDF}}
\cong
\mathcal H_{\text{light-cone}}.
\] Because \(\mathcal H_{\text{DDF}}\) is generated only by transverse
oscillators, its inner product is positive definite. Therefore any
negative-norm contribution in \(\mathcal H_{\text{phys}}\) must lie in
the null/spurious sector and does not survive in \[
\mathcal H_{\text{phys}}/\mathcal H_{\text{null}}.
\]

\subsubsection{Relation to light-cone
quantization}\label{relation-to-light-cone-quantization}

Light-cone quantization solves the constraints before quantization.
Introduce light-cone coordinates \[
X^\pm
=
\frac{1}{\sqrt{2}}(X^0\pm X^{D-1}).
\] One chooses light-cone gauge \[
X^+(\tau,\sigma)=x^+ + 2\alpha' p^+\tau.
\] The Virasoro constraints then solve the longitudinal oscillators
\(\alpha_n^-\) in terms of the transverse oscillators \(\alpha_n^i\): \[
\alpha_n^-
=
\frac{1}{p^+}\times
\text{quadratic expression in }\alpha_m^i.
\] Thus the only independent oscillators are \[
\alpha_{-n}^i,
\qquad
i=1,\dots,D-2.
\] The light-cone Hilbert space is manifestly positive definite.

The price is that Lorentz invariance is no longer manifest. Quantum
Lorentz invariance of the light-cone theory forces \[
D=26,
\qquad
a=1
\] for the bosonic string. Thus the covariant no-ghost theorem and the
light-cone Lorentz anomaly calculation are two sides of the same fact:
\[
\text{critical covariant physical quotient}
\cong
\text{positive light-cone Hilbert space}.
\]

\subsubsection{BRST formulation of the same
result}\label{brst-formulation-of-the-same-result}

In modern language, the cleanest conceptual statement uses BRST
cohomology. After gauge fixing the worldsheet metric to conformal gauge,
one introduces \(b,c\) ghosts. The BRST charge has the schematic form \[
Q_B
=
\sum_n c_{-n}\left(L_n^{\text{m}}+\frac{1}{2}L_n^{\text{gh}}-a\delta_{n,0}\right)
-
\frac{1}{2}\sum_{m,n}(m-n):c_{-m}c_{-n}b_{m+n}:.
\] Here \(L_n^{\text{m}}\) are matter Virasoro generators and
\(L_n^{\text{gh}}\) are ghost Virasoro generators.

Consistency requires \[
Q_B^2=0.
\] For the bosonic string this condition gives \[
D=26,
\qquad
a=1.
\] Physical states are BRST cohomology classes: \[
Q_B|\psi\rangle=0,
\qquad
|\psi\rangle\sim|\psi\rangle+Q_B|\Lambda\rangle.
\] The BRST-exact states \[
Q_B|\Lambda\rangle
\] are gauge artifacts. They are the BRST version of spurious states.
The no-ghost theorem becomes the statement that the BRST cohomology is
represented by transverse excitations only: \[
H^\bullet(Q_B)
\cong
\mathcal H_{\text{transverse}}
\] at the appropriate ghost number.

This also explains the word ``decouple.'' If a state is BRST exact, then
inserting its vertex operator into a physical amplitude gives zero,
assuming the usual absence of boundary contributions in moduli space: \[
\mathcal A(Q_B\Lambda,V_2,\dots,V_n)=0.
\] So exact states, like null spurious states in old-covariant
quantization, do not affect physical S-matrix elements.

\subsubsection{Superstring: where the extra ghosts come
from}\label{superstring-where-the-extra-ghosts-come-from}

In the RNS superstring, the matter fields are \[
X^\mu(\tau,\sigma),
\qquad
\psi^\mu(\tau,\sigma).
\] The bosonic oscillators still satisfy \[
[\alpha_m^\mu,\alpha_n^\nu]
=
m\delta_{m+n,0}\eta^{\mu\nu},
\] and the fermionic modes satisfy \[
\{\psi_r^\mu,\psi_s^\nu\}
=
\eta^{\mu\nu}\delta_{r+s,0}.
\] Because \(\eta^{00}<0\), fermionic timelike excitations also have
indefinite norm. For example, \[
\|\psi_{-1/2}^0|0;p\rangle_{\text{NS}}\|^2
=
\eta^{00}
<0.
\] Thus the superstring has the same covariant ghost problem.

The new feature is local worldsheet supersymmetry. In conformal gauge,
the constraints are not only \[
T(z)=0,
\] but also \[
T_F(z)=0,
\] where \(T_F\) is the worldsheet supercurrent. In modes, the
physical-state conditions include both Virasoro and super-Virasoro
constraints.

\subsubsection{NS-sector physical-state
conditions}\label{ns-sector-physical-state-conditions}

In the NS sector, the fermions are half-integer moded: \[
r\in\mathbb Z+\frac{1}{2}.
\] The physical-state conditions are \[
(L_0-a_{\text{NS}})|\psi\rangle=0,
\] \[
L_n|\psi\rangle=0,
\qquad
n>0,
\] and \[
G_r|\psi\rangle=0,
\qquad
r>0.
\] At the critical dimension, \[
D=10,
\] the normal-ordering constant is \[
a_{\text{NS}}=\frac{1}{2}.
\] Thus the open-string NS mass formula is \[
\alpha' M^2
=
N-\frac{1}{2}.
\] Before the GSO projection, the NS ground state is tachyonic. This
tachyon is not a negative-norm ghost; it has positive norm but negative
mass squared. The no-ghost theorem concerns the sign of the inner
product, not the sign of \(M^2\).

At the first excited NS level, \[
|\zeta;p\rangle
=
\zeta_\mu\psi_{-1/2}^\mu|0;p\rangle_{\text{NS}},
\] and \(N=\frac{1}{2}\). The mass-shell condition gives \[
p^2=0.
\] The \(G_{1/2}\) constraint gives \[
p\cdot\zeta=0.
\] The polarization \(\zeta^\mu\sim p^\mu\) is a null gauge state. After
quotienting it out, the remaining states are again the \(D-2\)
transverse polarizations. For \(D=10\), this gives \[
D-2=8
\] positive physical polarizations in the open-string NS vector
multiplet.

\subsubsection{R-sector physical-state
conditions}\label{r-sector-physical-state-conditions}

In the R sector, the fermions are integer moded: \[
r\in\mathbb Z.
\] There are fermion zero modes \[
\psi_0^\mu,
\] which satisfy a spacetime Clifford algebra: \[
\{\psi_0^\mu,\psi_0^\nu\}
=
\eta^{\mu\nu}.
\] Therefore the R-sector ground state is a spacetime spinor.

The physical-state conditions are \[
(L_0-a_{\text{R}})|\psi\rangle=0,
\] \[
L_n|\psi\rangle=0,
\qquad
n>0,
\] and \[
G_n|\psi\rangle=0,
\qquad
n\geq 0.
\] The \(G_0\) condition is special. Since \[
G_0^2=L_0
\] up to the normal-ordering convention, the \(G_0\) constraint is the
spacetime Dirac equation on the R ground state. At the critical
dimension, \[
D=10,
\qquad
a_{\text{R}}=0.
\] The R ground state is massless and satisfies \[
\slashed p\,|u;p\rangle=0.
\] After imposing this Dirac equation and quotienting null gauge
components, the physical spinor states can be represented by transverse
\(SO(8)\) spinor degrees of freedom. Again, the indefinite timelike
components do not survive in the physical quotient.

\subsubsection{Super-Virasoro algebra and critical
dimension}\label{super-virasoro-algebra-and-critical-dimension}

The matter super-Virasoro algebra has the form \[
[L_m,L_n]
=
(m-n)L_{m+n}
+
\frac{c}{12}(m^3-m)\delta_{m+n,0},
\] \[
[L_m,G_r]
=
\left(\frac{m}{2}-r\right)G_{m+r},
\] \[
\{G_r,G_s\}
=
2L_{r+s}
+
\frac{c}{3}\left(r^2-\frac{1}{4}\right)\delta_{r+s,0}
\] in the NS sector. The matter central charge is \[
c_{\text{m}}
=
c_X+c_\psi
=
D+\frac{D}{2}
=
\frac{3D}{2}.
\] Gauge fixing local worldsheet supersymmetry introduces not only the
fermionic \((b,c)\) ghosts, but also the bosonic superconformal ghosts
\((\beta,\gamma)\). Their total ghost central charge is \[
c_{\text{gh}}
=
-26+11
=
-15.
\] BRST nilpotency requires \[
c_{\text{m}}+c_{\text{gh}}=0.
\] Therefore \[
\frac{3D}{2}-15=0,
\] so \[
D=10.
\] The same consistency condition gives \[
a_{\text{NS}}=\frac{1}{2},
\qquad
a_{\text{R}}=0.
\] Thus the superstring no-ghost theorem holds at the critical RNS
values \[
D=10,
\qquad
a_{\text{NS}}=\frac{1}{2},
\qquad
a_{\text{R}}=0.
\]

\subsubsection{Super-DDF form of the no-ghost
theorem}\label{super-ddf-form-of-the-no-ghost-theorem}

Just as in the bosonic string, one can build physical transverse
operators in the superstring. In addition to bosonic DDF oscillators
\(A_n^i\), there are fermionic DDF oscillators, usually denoted
schematically by \[
B_r^i
\] in the NS sector, with half-integer \(r\), and by integer-moded
analogues in the R sector.

They obey transverse oscillator relations of the form \[
[A_m^i,A_n^j]
=
m\delta^{ij}\delta_{m+n,0},
\] \[
\{B_r^i,B_s^j\}
=
\delta^{ij}\delta_{r+s,0}.
\] They commute with the positive super-Virasoro constraints in the
appropriate sense: \[
[L_n,A_m^i]\sim 0,
\qquad
[L_n,B_r^i]\sim 0,
\qquad
[G_r,A_n^i]\sim 0,
\qquad
\{G_r,B_s^i\}\sim 0
\] on physical states, with the precise relations arranged so that DDF
operators map physical states to physical states.

The superstring no-ghost theorem says: \[
\boxed{
\frac{\mathcal H_{\text{phys}}^{\text{RNS}}}
{\mathcal H_{\text{null}}^{\text{RNS}}}
\cong
\mathcal H_{\text{transverse}}^{\text{RNS}}
}
\] where the right-hand side is generated only by transverse bosonic and
fermionic oscillators: \[
A_{-n}^i,
\qquad
B_{-r}^i,
\qquad
i=1,\dots,D-2.
\] For the critical superstring, \[
D-2=8.
\] Thus the physical quotient is the positive-definite Fock space of
eight transverse bosons and eight transverse worldsheet fermions, with
the usual NS and R sector distinctions.

\subsubsection{How this explains decoupling in
amplitudes}\label{how-this-explains-decoupling-in-amplitudes}

Suppose \(|\lambda\rangle\) is a physical null state. Then \[
\langle \phi|\lambda\rangle=0
\] for every physical state \(|\phi\rangle\). In operator language, the
corresponding vertex operator is either spurious in old-covariant
quantization or BRST exact in BRST quantization: \[
V_\lambda
=
\{Q_B,\Lambda\}
\] up to picture and sector-dependent details in the superstring.

If the other external vertex operators are physical, meaning BRST
closed, \[
[Q_B,V_i]=0,
\] then a physical amplitude with one null insertion is \[
\mathcal A(V_\lambda,V_2,\dots,V_n)
=
\mathcal A(\{Q_B,\Lambda\},V_2,\dots,V_n).
\] By moving \(Q_B\) through the correlator, this becomes a total
derivative on moduli space or vanishes by BRST invariance: \[
\mathcal A(V_\lambda,V_2,\dots,V_n)=0,
\] assuming no anomalous boundary contribution. This is the precise
sense in which null and spurious states decouple.

So the phrase ``negative-norm states decouple'' should be understood
carefully. The theorem does not say that the full covariant Fock space
has no negative-norm vectors. It certainly does. It says that after
imposing the constraints and quotienting null gauge states, the physical
Hilbert space is positive definite. Negative-norm directions are gauge
artifacts of the covariant description.

\subsubsection{Closed strings}\label{closed-strings}

For closed strings, there are left- and right-moving copies: \[
L_n,
\qquad
\widetilde L_n.
\] The physical-state conditions for the bosonic closed string are \[
(L_0-a)|\psi\rangle=0,
\qquad
(\widetilde L_0-a)|\psi\rangle=0,
\] \[
L_n|\psi\rangle=0,
\qquad
\widetilde L_n|\psi\rangle=0,
\qquad
n>0.
\] Equivalently, one imposes the mass-shell condition and level
matching: \[
N-\widetilde N=0.
\] The no-ghost theorem applies separately to the left- and right-moving
sectors. The physical quotient is the tensor product of the two
transverse Hilbert spaces, subject to level matching.

For the closed superstring, the same statement holds with left and right
NS/R choices: \[
\text{NS-NS},
\qquad
\text{NS-R},
\qquad
\text{R-NS},
\qquad
\text{R-R}.
\] Each sector has a positive physical quotient at \[
D=10.
\]

\subsubsection{Checks}\label{checks-8}

For the bosonic open string at level one, the critical values give \[
D=26,
\qquad
a=1.
\] The physical massless vector has \[
D-2=24
\] positive transverse polarizations. The timelike and longitudinal
polarizations are removed by the \(L_1\) constraint and by quotienting
the null state \(L_{-1}|0;p\rangle\).

For the open superstring NS sector at the first massless level, \[
D=10,
\qquad
a_{\text{NS}}=\frac{1}{2}.
\] The physical massless vector has \[
D-2=8
\] positive transverse polarizations.

For the R sector, the ground state is a spacetime spinor. The \(G_0\)
constraint imposes the massless Dirac equation, and the physical
quotient is represented by transverse \(SO(8)\) spinor data, again with
positive norm.

\subsubsection{Result}\label{result-12}

The covariant oscillator Fock space is indefinite because \[
[\alpha_m^0,\alpha_n^0]
=
m\delta_{m+n,0}\eta^{00}
\] and, in the RNS string, \[
\{\psi_r^0,\psi_s^0\}
=
\eta^{00}\delta_{r+s,0}.
\] Physical states satisfy the old-covariant constraints. For the
bosonic open string: \[
(L_0-a)|\psi\rangle=0,
\qquad
L_n|\psi\rangle=0
\quad
n>0.
\] At the critical values \[
D=26,
\qquad
a=1,
\] the no-ghost theorem says \[
\frac{\mathcal H_{\text{phys}}}{\mathcal H_{\text{null}}}
\cong
\mathcal H_{\text{DDF}}
\cong
\mathcal H_{\text{light-cone}},
\] which is positive definite.

For the RNS superstring, the constraints are \[
(L_0-a)|\psi\rangle=0,
\qquad
L_n|\psi\rangle=0,
\qquad
G_r|\psi\rangle=0
\] with the appropriate NS or R moding. At the critical values \[
D=10,
\qquad
a_{\text{NS}}=\frac{1}{2},
\qquad
a_{\text{R}}=0,
\] the superstring no-ghost theorem says \[
\frac{\mathcal H_{\text{phys}}^{\text{RNS}}}
{\mathcal H_{\text{null}}^{\text{RNS}}}
\cong
\mathcal H_{\text{transverse}}^{\text{RNS}},
\] generated only by \[
A_{-n}^i,
\qquad
B_{-r}^i,
\qquad
i=1,\dots,D-2.
\] Therefore the final physical Hilbert space has no negative-norm
states. The ghosts of covariant quantization are removed by the Virasoro
and super-Virasoro constraints, and the null/spurious remnants decouple
from all physical amplitudes.

\section{Q2: World-sheet fermionic oscillators and spacetime
bosons}\label{q2-world-sheet-fermionic-oscillators-and-spacetime-bosons}

Consider the open superstring with NN boundary condition. In the
light-cone quantization, the spectrum is given by \[
|0\rangle \ , \ b^i_{-\frac{1}{2}}|0\rangle \ , \ \alpha_{-1}^i|0\rangle+b^i_{-\frac{1}{2}}b^j_{-\frac{1}{2}}|0\rangle\ ,\cdots
\] The first level is the tachyon. The second level which is massless
has 8 states generated by fermionic creation operator
\(b_{-\frac{1}{2}}^i\) acting on the vacuum. It is a massless particle
under the little group \(SO(8)\) i.e.~a massless vector. But I am
confused that the operator \(b\) has its anti-commuting nature, why the
states generated by them can be bosonic vector?

A general state at massless level is \[
|\zeta;k\rangle
=
\zeta_i b_{-1/2}^i|0;k\rangle_{\text{NS}}
\] If we regard \(\zeta_{i}\) as the spacetime field, then it is an
\(SO(8)\) vector. Is that in the discussion of the string theory, it is
fields such as \(\zeta_{i}\) have physical importance while the state
\(b_{-1/2}^i|0;k\rangle_{\text{NS}}\) is not such important?

\begin{quote}
\textbf{Textbook location and related material.} Chapter 8, Sections
8.2--8.3 (pp.~200--208): light-cone RNS spectrum, NS and R sectors, and
the GSO projection. The spacetime-field interpretation is connected to
the covariant vertices of Section 13.2.
\end{quote}

\subsection{Answer}\label{answer-13}

\subsubsection{Core point}\label{core-point-13}

The confusion comes from using the word ``fermionic'' in two different
senses.

The oscillator \[
b_r^i
\] is fermionic as a worldsheet oscillator. This means it comes from the
worldsheet fermion field \(\psi^i(\tau,\sigma)\) and obeys
anticommutation relations \[
\{b_r^i,b_s^j\}
=
\delta^{ij}\delta_{r+s,0}.
\] It does not mean that the one-string state created by \(b_{-1/2}^i\)
is automatically a spacetime fermion.

The massless NS state \[
|\zeta;k\rangle
=
\zeta_i b_{-1/2}^i|0;k\rangle_{\text{NS}}
\] has one transverse vector index \(i\). Under the ten-dimensional
massless little group \[
SO(8),
\] it transforms as the vector representation \[
8_v.
\] That is why it is a spacetime vector boson. Its spacetime statistics
are determined by its target-space little-group representation and by
second quantization of the string field, not by the anticommutation
algebra of the internal worldsheet oscillator.

Spacetime fermions in the RNS string come from the Ramond sector, not
from the NS oscillator \(b_{-1/2}^i\). In the R sector the fermion zero
modes form a spacetime Clifford algebra, and the ground states are
spacetime spinors.

For your follow-up question: no, the basis state \[
b_{-1/2}^i|0;k\rangle_{\text{NS}}
\] is not unimportant. It is part of the physical one-string Hilbert
space. The polarization \(\zeta_i\) and the basis state play
complementary roles. The physical state is the linear combination \[
|\zeta;k\rangle
=
\sum_{i=1}^8 \zeta_i\, b_{-1/2}^i|0;k\rangle_{\text{NS}}.
\] The \(b_{-1/2}^i|0;k\rangle_{\text{NS}}\) are the basis vectors of
the \(SO(8)\) vector representation, while the numbers \(\zeta_i\) are
the components of the polarization vector in that basis.

So the right statement is not \[
\text{only }\zeta_i\text{ is physical, while }b_{-1/2}^i|0;k\rangle_{\text{NS}}\text{ is not}.
\] Rather, the physical object is the state \[
|\zeta;k\rangle\in \mathcal H_{k,\text{massless}}^{\text{NS}},
\] or equivalently the corresponding vertex operator, or equivalently in
second quantization the spacetime field mode multiplying this state.

\subsubsection{Basis states, polarization components, and the physical
vector}\label{basis-states-polarization-components-and-the-physical-vector}

At fixed lightlike momentum \(k\), the massless NS Hilbert space is
eight-dimensional: \[
\mathcal H_{k,\text{massless}}^{\text{NS}}
=
\operatorname{span}
\left\{
b_{-1/2}^i|0;k\rangle_{\text{NS}}
\right\}_{i=1}^8.
\] This is just like an ordinary vector space. If \(e_i\) is a basis of
\(\mathbb R^8\), a vector \(v\) is written as \[
v=\sum_{i=1}^8 v_i e_i.
\] The basis vectors \(e_i\) are not ``less physical'' than the
components \(v_i\); rather, the vector is the combination of both: \[
\text{vector}
=
\text{components}\times\text{basis}.
\] In the string Hilbert space the analogous identification is \[
e_i
\longleftrightarrow
b_{-1/2}^i|0;k\rangle_{\text{NS}},
\qquad
v_i
\longleftrightarrow
\zeta_i.
\] Thus \[
|\zeta;k\rangle
=
\zeta_i b_{-1/2}^i|0;k\rangle_{\text{NS}}
\] is the actual one-particle state. The polarization \(\zeta_i\) is not
by itself the state. It is the coordinate description of the state in
the basis provided by the oscillator construction.

This is why it is meaningful to say both of the following: \[
b_{-1/2}^i|0;k\rangle_{\text{NS}}
\text{ transforms as }8_v,
\] and \[
\zeta_i
\text{ is a vector polarization}.
\] They are two sides of the same representation-theoretic statement.

\subsubsection{Active and passive transformation
viewpoints}\label{active-and-passive-transformation-viewpoints}

Let \(R\in SO(8)\). The basis states transform as \[
U(R)\,b_{-1/2}^i|0;k\rangle_{\text{NS}}
=
R^i{}_j\,b_{-1/2}^j|0;Rk\rangle_{\text{NS}}.
\] Therefore \[
U(R)|\zeta;k\rangle
=
\zeta_i R^i{}_j b_{-1/2}^j|0;Rk\rangle_{\text{NS}}.
\] If we define the transformed polarization by \[
\zeta'_j
=
\zeta_i R^i{}_j,
\] then \[
U(R)|\zeta;k\rangle
=
|\zeta';Rk\rangle.
\] This is exactly what it means for the one-particle states to furnish
the vector representation of \(SO(8)\).

Depending on convention, one may instead say that the basis is fixed and
the components transform, or that the basis transforms and the abstract
vector is the object being acted on. These are the passive and active
descriptions of the same mathematics. In neither description is the
state unimportant.

\subsubsection{First-quantized states versus spacetime
fields}\label{first-quantized-states-versus-spacetime-fields}

There are two levels of description.

In first-quantized string theory, the fundamental object is the
single-string Hilbert space. A physical particle state is a vector in
this Hilbert space: \[
|\zeta;k\rangle
=
\zeta_i b_{-1/2}^i|0;k\rangle_{\text{NS}}.
\] Equivalently, in the operator-state correspondence, it is represented
by a vertex operator. In light-cone language the oscillator state is
natural; in covariant language the corresponding massless NS vertex
operator is \[
V_A^{(-1)}
=
g_o\,\zeta_\mu \psi^\mu e^{-\phi}e^{ik\cdot X}.
\]

In second-quantized string theory, or in the low-energy effective field
theory, one expands a string field in a basis of first-quantized string
states. Schematically, for the massless open NS sector, \[
|\Psi\rangle
=
\int d^{10}k\,
A_i(k)\,b_{-1/2}^i|0;k\rangle_{\text{NS}}
+\cdots.
\] Here \[
A_i(k)
\] is the momentum-space gauge field mode. Fourier transforming gives \[
A_i(x)
=
\int d^{10}k\, e^{ik\cdot x}A_i(k)
\] up to normalization conventions.

So if by \(\zeta_i\) you mean the polarization of a single external
string state, then \(\zeta_i\) is a fixed vector labeling that
one-particle state. If by \(A_i(k)\) you mean a dynamical spacetime
field mode, then it appears as the coefficient of the string basis state
in the second-quantized string field: \[
\text{string field coefficient}
\quad
A_i(k)
\quad
\text{multiplies}
\quad
b_{-1/2}^i|0;k\rangle_{\text{NS}}.
\] This is the precise sense in which spacetime fields arise from string
states.

\subsubsection{Why the anticommuting oscillator does not make the
coefficient
Grassmann}\label{why-the-anticommuting-oscillator-does-not-make-the-coefficient-grassmann}

One might still worry: if \(b_{-1/2}^i\) is fermionic, should
\(\zeta_i\) or \(A_i(k)\) be Grassmann-odd? For the NS gauge boson, the
answer is no.

The oscillator \(b_{-1/2}^i\) is odd with respect to the worldsheet
fermion-number grading. Acting once gives a state with odd NS fermion
number: \[
(-1)^F b_{-1/2}^i|0\rangle_{\text{NS}}
=
-
b_{-1/2}^i|0\rangle_{\text{NS}},
\] up to the convention for the NS vacuum. But this \((-1)^F\) is a
worldsheet grading used in the GSO projection. It is not identical to
spacetime spin-statistics.

The spacetime statistics are determined by the spacetime little-group
representation: \[
8_v
\quad
\text{is bosonic},
\] while \[
8_s,\ 8_c
\quad
\text{are fermionic}.
\] Therefore the second-quantized coefficient \(A_i(k)\) is an ordinary
commuting bosonic field mode, even though the first-quantized oscillator
\(b_{-1/2}^i\) anticommutes with other \(b\) oscillators.

This is analogous to building a differential form with anticommuting
basis elements \(dx^i\). The components of an ordinary one-form, \[
A=A_i dx^i,
\] need not be Grassmann-odd just because the exterior basis satisfies
\[
dx^i\wedge dx^j=-dx^j\wedge dx^i.
\] The anticommutation controls the internal algebra of the basis
excitations, not the spacetime statistics of the coefficient field.

\subsubsection{The light-cone NS
spectrum}\label{the-light-cone-ns-spectrum}

For the open RNS string with Neumann-Neumann boundary conditions in all
target-space directions, the NS-sector worldsheet fermions are
half-integer moded: \[
r\in\mathbb Z+\frac{1}{2}.
\] In light-cone gauge, only the transverse directions are independent:
\[
i=1,\dots,8
\] in ten dimensions. The independent oscillators are \[
\alpha_{-n}^i,
\qquad
n=1,2,\dots,
\] and \[
b_{-r}^i,
\qquad
r=\frac{1}{2},\frac{3}{2},\dots.
\] They obey \[
[\alpha_m^i,\alpha_n^j]
=
m\delta^{ij}\delta_{m+n,0},
\] and \[
\{b_r^i,b_s^j\}
=
\delta^{ij}\delta_{r+s,0}.
\] The NS number operator is \[
N
=
\sum_{n=1}^{\infty}\alpha_{-n}^i\alpha_n^i
+
\sum_{r>0}r\, b_{-r}^i b_r^i.
\] The open-string NS mass formula is \[
\alpha' M^2
=
N-\frac{1}{2}.
\] Therefore:

\begin{enumerate}
\def\labelenumi{\arabic{enumi}.}
\item
  the NS vacuum has \(N=0\), so \[
  \alpha' M^2=-\frac{1}{2},
  \] and is tachyonic before the GSO projection;
\item
  the state \[
  b_{-1/2}^i|0;k\rangle_{\text{NS}}
  \] has \[
  N=\frac{1}{2},
  \] so \[
  M^2=0.
  \] This is the first massless NS state.
\end{enumerate}

There are eight such states because \(i=1,\dots,8\). In a covariant
description this is the familiar massless gauge boson state \[
\zeta_\mu \psi_{-1/2}^\mu|0;k\rangle_{\text{NS}},
\] with the physical constraints removing the timelike and longitudinal
polarizations. In light-cone gauge only the eight transverse
polarizations remain: \[
\zeta_i b_{-1/2}^i|0;k\rangle_{\text{NS}}.
\]

\subsubsection{What anticommuting means
here}\label{what-anticommuting-means-here}

The anticommutation relation says that the \(b\) oscillators generate an
exterior algebra. For example, \[
b_{-1/2}^i b_{-1/2}^j
=
-
b_{-1/2}^j b_{-1/2}^i,
\] and in particular \[
(b_{-1/2}^i)^2=0.
\] So each fermionic oscillator can be occupied at most once.

This affects the degeneracy and representation content of excited string
states. For example, \[
b_{-1/2}^i|0\rangle_{\text{NS}}
\] transforms as a vector of \(SO(8)\), while \[
b_{-1/2}^i b_{-1/2}^j|0\rangle_{\text{NS}}
\] is antisymmetric in \(i,j\), so it transforms as an antisymmetric
two-tensor of \(SO(8)\): \[
\wedge^2 8_v.
\] The anticommuting nature of the \(b\) oscillators tells us that
multiple \(b\) excitations are antisymmetrized. It does not tell us that
the resulting target-space particle is a spacetime fermion.

A useful analogy is the exterior algebra of differential forms. The
basis one-forms \(dx^i\) anticommute: \[
dx^i\wedge dx^j=-dx^j\wedge dx^i.
\] But a one-form is not a spacetime fermion. It is a tensorial object.
Similarly, the \(b_{-1/2}^i\) oscillators anticommute, but their index
\(i\) is a vector index, not a spinor index.

\subsubsection{Worldsheet spin versus spacetime
spin}\label{worldsheet-spin-versus-spacetime-spin}

The RNS field \(\psi^\mu\) has two kinds of transformation behavior:

\begin{enumerate}
\def\labelenumi{\arabic{enumi}.}
\tightlist
\item
  it is a fermion on the worldsheet;
\item
  it is a vector in target spacetime.
\end{enumerate}

The first statement means that \(\psi^\mu\) has conformal weight \[
h=\frac{1}{2}
\] and its modes obey anticommutation relations.

The second statement means that under spacetime Lorentz transformations,
\[
\psi^\mu
\] carries a vector index \(\mu\). In light-cone gauge the physical
components are \[
\psi^i,
\qquad
i=1,\dots,8,
\] so their modes \(b_r^i\) carry the vector representation of the
transverse rotation group \(SO(8)\).

Thus \[
b_{-1/2}^i|0;k\rangle_{\text{NS}}
\] is spacetime-vector-valued because of the index \(i\). Its worldsheet
origin as a fermionic oscillator only says that the oscillator is
Grassmann-odd in the worldsheet operator algebra.

The distinction can be summarized as \[
\boxed{
\text{worldsheet fermion}
\neq
\text{spacetime fermion}.
}
\] In the RNS string, \(\psi^\mu\) is a worldsheet fermion but a
spacetime vector.

\subsubsection{Why the NS massless state is a
vector}\label{why-the-ns-massless-state-is-a-vector}

Let \(R\in SO(8)\) be a transverse rotation. The oscillators transform
as \[
b_{-1/2}^i
\longmapsto
R^i{}_j b_{-1/2}^j.
\] Assume the NS ground state is a spacetime scalar: \[
|0;k\rangle_{\text{NS}}
\longmapsto
|0;Rk\rangle_{\text{NS}}.
\] Then the one-oscillator state transforms as \[
b_{-1/2}^i|0;k\rangle_{\text{NS}}
\longmapsto
R^i{}_j b_{-1/2}^j|0;Rk\rangle_{\text{NS}}.
\] That is exactly the transformation law of a vector.

Equivalently, a general polarization state \[
|\zeta;k\rangle
=
\zeta_i b_{-1/2}^i|0;k\rangle_{\text{NS}}
\] transforms by \[
\zeta_i
\longmapsto
R_i{}^j\zeta_j.
\] So the physical polarization \(\zeta_i\) is an \(SO(8)\) vector.
Since vectors are tensor representations, not spinor representations,
this state is a spacetime boson.

Under a \(2\pi\) spacetime rotation, a vector returns to itself: \[
8_v
\longmapsto
+8_v.
\] A spacetime spinor would pick up a minus sign under a \(2\pi\)
rotation. The NS massless state does not do that.

\subsubsection{Where spacetime fermions come from: the R
sector}\label{where-spacetime-fermions-come-from-the-r-sector}

Now compare this with the Ramond sector. In the R sector, the worldsheet
fermions are integer moded: \[
r\in\mathbb Z.
\] In particular, there are zero modes \[
b_0^i.
\] They satisfy \[
\{b_0^i,b_0^j\}
=
\delta^{ij}.
\] This is a Clifford algebra. Therefore the R-sector ground states form
a spinor representation of \(SO(8)\): \[
|A\rangle_{\text{R}},
\qquad
A=1,\dots,8
\] after the usual GSO projection, with the chirality depending on the
projection convention.

Schematically, the R ground state transforms as \[
|A\rangle_{\text{R}}
\longmapsto
S(R)^A{}_B |B\rangle_{\text{R}},
\] where \(S(R)\) is a spinor representation matrix of \(SO(8)\), not
the vector matrix \(R^i{}_j\).

Thus: \[
\text{NS sector}
\quad
\Rightarrow
\quad
\text{spacetime tensor/vector states},
\] while \[
\text{R sector}
\quad
\Rightarrow
\quad
\text{spacetime spinor states}.
\] This is why the massless open superstring spectrum after the GSO
projection contains a gauge boson from the NS sector and a gaugino from
the R sector.

\subsubsection{The role of the GSO
projection}\label{the-role-of-the-gso-projection}

Before the GSO projection, the NS sector contains the tachyonic ground
state \[
|0;k\rangle_{\text{NS}},
\] with \[
\alpha' M^2=-\frac{1}{2}.
\] The GSO projection removes this tachyon and keeps the states with the
appropriate worldsheet fermion number. In the common open-superstring
convention, it keeps the state \[
b_{-1/2}^i|0;k\rangle_{\text{NS}}
\] and removes the NS vacuum.

This sometimes adds to the confusion because the kept NS state has one
\(b\) oscillator, hence odd worldsheet fermion number. But odd
worldsheet fermion number is not the same as spacetime fermionic
statistics. It is a quantum number of the worldsheet CFT Hilbert space.

The post-GSO massless open-string states are: \[
\text{NS:}
\qquad
b_{-1/2}^i|0;k\rangle_{\text{NS}}
\quad
\in 8_v,
\] and \[
\text{R:}
\qquad
|A;k\rangle_{\text{R}}
\quad
\in 8_s
\text{ or }
8_c.
\] Together these form the field content of ten-dimensional \(N=1\)
super Yang-Mills: \[
A_\mu
\quad
\text{and}
\quad
\lambda_A.
\] In light-cone gauge, the on-shell physical polarizations are \[
A_i,
\qquad
i=1,\dots,8,
\] and an \(SO(8)\) spinor gaugino.

\subsubsection{First-quantized oscillator algebra versus
second-quantized
statistics}\label{first-quantized-oscillator-algebra-versus-second-quantized-statistics}

There is one final conceptual point. The operators \[
\alpha_{-n}^i,
\qquad
b_{-r}^i
\] are not spacetime particle creation operators in a second-quantized
quantum field theory. They are first-quantized string oscillators. They
create internal excitations of a single string.

The spacetime statistics of the corresponding particle is determined
after second quantization. A spacetime bosonic string state gives a
spacetime field whose creation and annihilation operators commute. A
spacetime fermionic string state gives a spacetime field whose creation
and annihilation operators anticommute.

For the NS massless vector, the associated spacetime field is a gauge
field: \[
A_\mu(x).
\] The one-particle state is \[
\zeta_i b_{-1/2}^i|0;k\rangle_{\text{NS}},
\] but the second-quantized creation operator creating this gauge boson
is not the same object as the worldsheet oscillator \(b_{-1/2}^i\). The
former is a spacetime particle creation operator; the latter is an
internal oscillator acting on the single-string Hilbert space.

So there is no contradiction: \[
\boxed{
b_{-1/2}^i
\text{ anticommutes as a worldsheet oscillator, but }
b_{-1/2}^i|0;k\rangle_{\text{NS}}
\text{ is a spacetime vector boson.}
}
\]

\subsubsection{Vertex-operator
viewpoint}\label{vertex-operator-viewpoint}

The same conclusion appears in the covariant vertex-operator
description. The massless NS vector vertex operator in the \(-1\)
picture is \[
V_A^{(-1)}
=
g_o\,\zeta_\mu \psi^\mu e^{-\phi}e^{ik\cdot X},
\] with \[
k^2=0,
\qquad
k\cdot \zeta=0,
\] and the gauge equivalence \[
\zeta_\mu\sim \zeta_\mu+k_\mu.
\] The polarization is \(\zeta_\mu\), a spacetime vector. The field
\(\psi^\mu\) is worldsheet-fermionic, but it carries a spacetime vector
index. Thus this vertex operator creates a spacetime vector state.

By contrast, a massless R-sector fermion vertex involves a spin field:
\[
V_\lambda^{(-1/2)}
\sim
u_A S^A e^{-\phi/2}e^{ik\cdot X}.
\] Here \(S^A\) is a spin field and \(u_A\) is a spacetime spinor
polarization. This is the vertex operator that creates a spacetime
fermion.

So the distinction is visible directly: \[
\psi^\mu
\quad
\text{carries a vector index},
\] whereas \[
S^A
\quad
\text{carries a spinor index}.
\]

\subsubsection{Checks}\label{checks-9}

For the NS massless state: \[
N=\frac{1}{2},
\qquad
\alpha' M^2=N-\frac{1}{2}=0.
\] There are eight independent light-cone states: \[
b_{-1/2}^i|0;k\rangle_{\text{NS}},
\qquad
i=1,\dots,8.
\] They transform as \[
8_v
\] of \(SO(8)\), so they are the eight transverse polarizations of a
ten-dimensional massless vector.

For the R ground state: \[
\{b_0^i,b_0^j\}
=
\delta^{ij}.
\] This Clifford algebra forces the ground state to transform as an
\(SO(8)\) spinor: \[
8_s
\quad
\text{or}
\quad
8_c.
\] Those are the spacetime fermions.

\subsubsection{Result}\label{result-13}

The NS fermionic oscillators obey \[
\{b_r^i,b_s^j\}
=
\delta^{ij}\delta_{r+s,0}.
\] This means they are worldsheet fermionic oscillators and generate
antisymmetrized multi-oscillator states.

The massless NS state is \[
|\zeta;k\rangle
=
\zeta_i b_{-1/2}^i|0;k\rangle_{\text{NS}},
\] with \[
N=\frac{1}{2},
\qquad
M^2=0.
\] At fixed \(k\), \[
\left\{
b_{-1/2}^i|0;k\rangle_{\text{NS}}
\right\}_{i=1}^8
\] is a basis of the massless one-string Hilbert space, and \(\zeta_i\)
are the components of the polarization in that basis. The physical
one-particle state is the full vector \[
|\zeta;k\rangle,
\] not \(\zeta_i\) alone and not the basis state alone.

In second quantization, the corresponding string field contains \[
|\Psi\rangle
\supset
\int d^{10}k\,
A_i(k)b_{-1/2}^i|0;k\rangle_{\text{NS}},
\] so the spacetime gauge field mode \(A_i(k)\) is the coefficient
multiplying the first-quantized string state.

Because \(b_{-1/2}^i\) carries a transverse vector index, \[
b_{-1/2}^i|0;k\rangle_{\text{NS}}
\in
8_v
\] under the massless little group \(SO(8)\). Therefore it is a
spacetime vector boson.

Spacetime fermions instead come from the R sector, where \[
\{b_0^i,b_0^j\}
=
\delta^{ij}
\] is a Clifford algebra, so the ground states transform as \[
8_s
\quad
\text{or}
\quad
8_c.
\] Thus the correct distinction is \[
\boxed{
\text{NS worldsheet fermion oscillator } b_{-1/2}^i
\text{ creates a spacetime vector state, not a spacetime fermion.}
}
\]

\chapter{Kac--Moody Algebras}\label{kacmoody-algebras-1}

\section{Q1: Integrable highest weights of simply-laced affine
algebras}\label{q1-integrable-highest-weights-of-simply-laced-affine-algebras}

Consider the chiral algebra \(\mathcal{A}=\hat{v}\ltimes\hat{g}\), where
\(\hat{v}\) is the Virasoro algebra and \(\hat{g}\) is the level \(k\)
Kac-Moody algebra of simply laced Lie algebra \(g\) whose roots are
normalized to \((\text{length})^2=2\). Consider the \(\mathfrak{su}(2)\)
subalgebra of \(\bar{g}\) generated by
\(E_{1}^{-\vec{\alpha}},E_{-1}^{\vec{\alpha}}\) and
\((k-\vec{\alpha}\cdot \vec{H}_{0})\) where \(\vec{\alpha}\) is a root,
unitary requires that the eigenvalues of
\((k-\vec{\alpha}\cdot \vec{H}_{0})\) have to be integer. Acting on a
primary state \(|\vec{m}\rangle\) with weight \(\vec{m}\). Since for
simply-laced algebra we have
\(\vec{\alpha}\cdot \vec{m}\in \mathbb{Z}\), we find
\(k\in \mathbb{Z}\). Q1: why for simply-laced algebra we have
\(\vec{\alpha}\cdot \vec{m}\in \mathbb{Z}\)?

Since \(|\vec{m}\rangle\) is primary we have
\(E_{1}^{-\vec{\alpha}}|\vec{m}\rangle=0\) and \[
||E_{-1}^{\vec{\alpha}}|\vec{m}\rangle||^2=(k-\vec{\alpha}\cdot \vec{m})|||\vec{m}\rangle||^2\geq 0
\] Gives that \(k\geq \vec{\alpha}\cdot \vec{m}\). It is said that the
right hand side of this inequality is maximized if we chose
\(\vec{\alpha}\) to be the highest root \(\vec{\psi}\) and \(\vec{m}\)
to be the highest weight \(\vec{m}_{0}\) of the given representations,
we then have \(k\geq \vec{\psi}\cdot \vec{m}_{0}\). Q2: what is a
highest root and a highest weight, why the above statement is true?
Explain with some examples.

Clearly for a given \(k\) only a finite number of highest weights
satisfies this criterion. For \(k=1\) they are those belonging to the
lowest dimensional representation in each conjugacy class. Q3: give some
examples for the k=1 case.

\begin{quote}
\textbf{Textbook location and related material.} Chapter 11, Sections
11.1, 11.5, and 11.6 (pp.~321--355): affine current algebras, unitary
highest-weight representations, the level bound, characters, and the
explicit \(\widehat{\mathfrak{su}}(2)_k\) example. Sections 11.2--11.3
provide the root/weight-lattice background.
\end{quote}

\subsection{Answer}\label{answer-14}

\subsubsection{Core point}\label{core-point-14}

There are three related facts behind the statements in the question.

First, a representation weight is not an arbitrary vector in the real
span of the roots. By definition it pairs integrally with every
\textbf{coroot}: \[
\langle \vec m,\vec\alpha^{\vee}\rangle
=\frac{2\vec\alpha\cdot\vec m}{\vec\alpha^{\,2}}
\in\mathbb Z.
\] For a simply-laced algebra in the normalization
\(\vec\alpha^{\,2}=2\), the coroot equals the root. Hence \[
\vec\alpha\cdot\vec m\in\mathbb Z.
\]

Second, after choosing positive roots, a finite-dimensional irreducible
\(g\)-representation has a dominant highest weight
\(\vec m_0=\vec\lambda\). The root system itself has a highest root
\(\vec\psi\). The strongest affine positivity condition is \[
k-\vec\psi\cdot\vec\lambda\geq0.
\] Together with integrality, this is the usual affine integrability
condition \[
\vec\psi\cdot\vec\lambda\leq k.
\]

Third, if \[
\vec\psi=\sum_{i=1}^r a_i\vec\alpha_i,
\qquad
\vec\lambda=\sum_{i=1}^r\lambda_i\vec\omega_i,
\] then, in the simply-laced normalization, \[
\vec\psi\cdot\vec\lambda=\sum_{i=1}^r a_i\lambda_i.
\] At level \(k=1\), this sum can be at most one. Thus the allowed
nonzero highest weights are precisely the fundamental weights
\(\vec\omega_i\) attached to Dynkin nodes with \(a_i=1\). For ADE these
are the minuscule fundamental weights. They give the smallest
finite-dimensional representations in the different cosets of \(P/Q\).

\subsubsection{Weights pair integrally with
coroots}\label{weights-pair-integrally-with-coroots}

Let \(\vec\alpha_1,\ldots,\vec\alpha_r\) be simple roots. The
corresponding coroots are \[
\vec\alpha_i^{\vee}
=\frac{2\vec\alpha_i}{\vec\alpha_i^{\,2}}.
\] For each simple root, the generators \[
E_0^{\vec\alpha_i},
\qquad
E_0^{-\vec\alpha_i},
\qquad
\vec\alpha_i^{\vee}\cdot\vec H_0
\] form an ordinary \(\mathfrak{su}(2)\) subalgebra of the horizontal
algebra \(g\). If \(|\vec m\rangle\) has weight \(\vec m\), then \[
(\vec\alpha_i^{\vee}\cdot\vec H_0)|\vec m\rangle
=\langle\vec m,\vec\alpha_i^{\vee}\rangle|\vec m\rangle.
\] Every finite-dimensional unitary \(\mathfrak{su}(2)\) representation
has integral eigenvalues of twice the third generator. Therefore \[
\langle\vec m,\vec\alpha_i^{\vee}\rangle
\in\mathbb Z
\] for every simple coroot. This is exactly the statement that
\(\vec m\) lies in the weight lattice \(P\).

Equivalently, introduce the fundamental weights \(\vec\omega_j\) by \[
\langle\vec\omega_j,\vec\alpha_i^{\vee}\rangle
=\delta_{ij}.
\] Every weight-lattice element can be written as \[
\vec m=\sum_{j=1}^r m_j\vec\omega_j,
\qquad
m_j\in\mathbb Z.
\] For a simply-laced algebra with \(\vec\alpha_i^{\,2}=2\), \[
\vec\alpha_i^{\vee}=\vec\alpha_i,
\] and hence \[
\vec\alpha_i\cdot\vec m=m_i\in\mathbb Z.
\]

Now let \(\vec\alpha\) be any root, not necessarily simple. In a
simply-laced root system it has an expansion \[
\vec\alpha=\sum_{i=1}^r n_i\vec\alpha_i,
\qquad
n_i\in\mathbb Z.
\] It follows that \[
\vec\alpha\cdot\vec m
=\sum_{i=1}^r n_i(\vec\alpha_i\cdot\vec m)
=\sum_{i=1}^r n_i m_i
\in\mathbb Z.
\] This proves the statement used in Q1.

The simply-laced hypothesis is needed only to replace the coroot pairing
by the ordinary inner product. For a general Lie algebra, the universal
integrality statement is \[
\frac{2\vec\alpha\cdot\vec m}{\vec\alpha^{\,2}}
=\langle\vec m,\vec\alpha^{\vee}\rangle
\in\mathbb Z,
\] not necessarily \(\vec\alpha\cdot\vec m\in\mathbb Z\) for every root
in one fixed normalization.

\subsubsection{\texorpdfstring{The affine \(\mathfrak{su}(2)\) and level
quantization}{The affine \textbackslash mathfrak\{su\}(2) and level quantization}}\label{the-affine-mathfraksu2-and-level-quantization}

For each root \(\vec\alpha\), define \[
\mathcal J_+=E_1^{-\vec\alpha},
\qquad
\mathcal J_-=E_{-1}^{\vec\alpha},
\qquad
2\mathcal J_3=k-\vec\alpha\cdot\vec H_0.
\] With the normalization in the question, the affine commutation
relations give \[
[\mathcal J_3,\mathcal J_\pm]=\pm\mathcal J_\pm,
\qquad
[\mathcal J_+,\mathcal J_-]=2\mathcal J_3.
\] Thus these operators form an \(\mathfrak{su}(2)\) algebra. On a
weight state, \[
2\mathcal J_3|\vec m\rangle
=(k-\vec\alpha\cdot\vec m)|\vec m\rangle.
\] In a unitary representation, the eigenvalues of \(2\mathcal J_3\) are
integers, so \[
k-\vec\alpha\cdot\vec m\in\mathbb Z.
\] Since the previous section gives
\(\vec\alpha\cdot\vec m\in\mathbb Z\), one obtains \[
k\in\mathbb Z.
\]

There is also a positivity condition. An affine primary multiplet is
annihilated by all positive modes, so \[
E_1^{-\vec\alpha}|\vec m\rangle=0.
\] Using \((E_{-1}^{\vec\alpha})^\dagger=E_1^{-\vec\alpha}\), \[
\begin{aligned}
\left\|E_{-1}^{\vec\alpha}|\vec m\rangle\right\|^2
&=\langle\vec m|E_1^{-\vec\alpha}E_{-1}^{\vec\alpha}|\vec m\rangle\\
&=\langle\vec m|[E_1^{-\vec\alpha},E_{-1}^{\vec\alpha}]|\vec m\rangle\\
&=(k-\vec\alpha\cdot\vec m)\langle\vec m|\vec m\rangle.
\end{aligned}
\] Therefore \[
k-\vec\alpha\cdot\vec m\geq0.
\] In fact, because \(|\vec m\rangle\) is a highest-weight vector for
this affine \(\mathfrak{su}(2)\), unitarity gives the stronger combined
statement \[
k-\vec\alpha\cdot\vec m\in\mathbb Z_{\geq0}.
\] To see why the first norm inequality strengthens to an integer
condition, set \[
\ell_{\vec\alpha}:=k-\vec\alpha\cdot\vec m.
\] Repeated use of the \(\mathfrak{su}(2)\) commutators gives \[
\left\|\mathcal J_-^N|\vec m\rangle\right\|^2
=N!\prod_{s=0}^{N-1}(\ell_{\vec\alpha}-s)
\langle\vec m|\vec m\rangle.
\] If \(\ell_{\vec\alpha}\) were positive but nonintegral, then for
sufficiently large \(N\) one factor would become negative without any
earlier factor vanishing. This would produce a negative-norm state.
Hence the lowering chain must terminate when one factor vanishes, which
requires \[
\ell_{\vec\alpha}\in\mathbb Z_{\geq0}.
\]

\subsubsection{Highest weights of finite-dimensional
representations}\label{highest-weights-of-finite-dimensional-representations}

Choose a set of positive roots \(\Delta_+\). The simple roots are the
positive roots that cannot be decomposed into sums of other positive
roots. This choice determines raising operators \(E_0^{\vec\alpha}\) for
\(\vec\alpha\in\Delta_+\).

An irreducible finite-dimensional representation has a highest-weight
state \(|\vec\lambda\rangle\) satisfying \[
E_0^{\vec\alpha_i}|\vec\lambda\rangle=0
\qquad
\text{for every simple root }\vec\alpha_i,
\] and \[
\vec H_0|\vec\lambda\rangle
=\vec\lambda|\vec\lambda\rangle.
\] The weight \(\vec\lambda\) is dominant integral: \[
\lambda_i
:=\langle\vec\lambda,\vec\alpha_i^{\vee}\rangle
\in\mathbb Z_{\geq0}.
\] It can therefore be expanded as \[
\vec\lambda=\sum_{i=1}^r\lambda_i\vec\omega_i.
\]

Every other weight \(\vec m\) in the representation is obtained by
applying lowering operators. Consequently, \[
\vec m
=\vec\lambda-\sum_{i=1}^r N_i\vec\alpha_i,
\qquad
N_i\in\mathbb Z_{\geq0}.
\] This is the precise sense in which \(\vec\lambda\) is the highest
weight. ``Highest'' depends on the chosen positive-root system; it is
not a statement about Euclidean length.

For example, for \(\mathfrak{su}(2)\) there is one simple root
\(\vec\alpha\) and one fundamental weight \[
\vec\omega=\frac{\vec\alpha}{2}.
\] The spin-\(j\) representation has highest weight \[
\vec\lambda=2j\,\vec\omega,
\] and its weights are \[
\vec\lambda,
\quad
\vec\lambda-\vec\alpha,
\quad
\ldots,
\quad
\vec\lambda-2j\vec\alpha.
\]

\subsubsection{The highest root}\label{the-highest-root}

Every positive root has a unique expansion \[
\vec\alpha=\sum_{i=1}^r n_i\vec\alpha_i,
\qquad
n_i\in\mathbb Z_{\geq0}.
\] The root partial order is defined by \[
\vec\beta\succeq\vec\alpha
\quad\Longleftrightarrow\quad
\vec\beta-\vec\alpha
=\sum_i N_i\vec\alpha_i,
\qquad
N_i\in\mathbb Z_{\geq0}.
\] For a simple Lie algebra there is a unique maximal positive root in
this order. It is the highest root, \[
\vec\psi=\sum_{i=1}^r a_i\vec\alpha_i.
\] The coefficients \(a_i\) are called marks.

Another useful characterization is that the roots are the nonzero
weights of the adjoint representation, and \(\vec\psi\) is the highest
weight of the adjoint representation. In particular, \(\vec\psi\) is
dominant: \[
\langle\vec\psi,\vec\alpha_i^{\vee}\rangle\geq0
\] for every simple root.

For \(A_1=\mathfrak{su}(2)\), \[
\vec\psi=\vec\alpha.
\] For \(A_2=\mathfrak{su}(3)\), \[
\vec\psi=\vec\alpha_1+\vec\alpha_2.
\] For \(D_r=\mathfrak{so}(2r)\), with the standard numbering in which
\(\vec\alpha_{r-1}\) and \(\vec\alpha_r\) are the two spinor-end nodes,
\[
\vec\psi
=\vec\alpha_1
+2\vec\alpha_2
+\cdots
+2\vec\alpha_{r-2}
+\vec\alpha_{r-1}
+\vec\alpha_r.
\]

\subsubsection{\texorpdfstring{Why \(\vec\psi\cdot\vec\lambda\) is the
maximal
pairing}{Why \textbackslash vec\textbackslash psi\textbackslash cdot\textbackslash vec\textbackslash lambda is the maximal pairing}}\label{why-vecpsicdotveclambda-is-the-maximal-pairing}

The statement in the question should be understood as follows: for a
fixed irreducible representation \(V_{\vec\lambda}\), maximize \[
\vec\alpha\cdot\vec m
\] over all roots \(\vec\alpha\) and all weights \(\vec m\) of
\(V_{\vec\lambda}\). For a simply-laced root system, the maximum is \[
\vec\psi\cdot\vec\lambda.
\]

Here is a direct proof.

Because the root system is irreducible and simply laced, all roots have
the same length and form a single orbit under the Weyl group \(W\).
Therefore, for any root \(\vec\alpha\), there is a Weyl transformation
\(w\) such that \[
w(\vec\alpha)=\vec\psi.
\] The Weyl group preserves the inner product, and the set of weights of
a finite-dimensional representation is Weyl invariant. Thus \[
\vec\alpha\cdot\vec m
=w(\vec\alpha)\cdot w(\vec m)
=\vec\psi\cdot\vec\mu,
\] where \(\vec\mu=w(\vec m)\) is another weight of the same
representation.

Since \(\vec\lambda\) is the highest weight, \[
\vec\lambda-\vec\mu
=\sum_i N_i\vec\alpha_i,
\qquad
N_i\in\mathbb Z_{\geq0}.
\] Since \(\vec\psi\) is dominant and the algebra is simply laced, \[
\vec\psi\cdot\vec\alpha_i
=\langle\vec\psi,\vec\alpha_i^{\vee}\rangle
\geq0.
\] It follows that \[
\begin{aligned}
\vec\psi\cdot\vec\lambda
-\vec\psi\cdot\vec\mu
&=\vec\psi\cdot(\vec\lambda-\vec\mu)\\
&=\sum_iN_i\,\vec\psi\cdot\vec\alpha_i\\
&\geq0.
\end{aligned}
\] Therefore \[
\vec\alpha\cdot\vec m
\leq\vec\psi\cdot\vec\lambda.
\] The upper bound is attained by choosing \[
\vec\alpha=\vec\psi,
\qquad
\vec m=\vec\lambda.
\] This proves the maximization statement.

Applying the norm inequality to this pair gives \[
k\geq\vec\psi\cdot\vec\lambda.
\] This is not merely one convenient necessary bound. Together with
\(\vec\lambda\) being dominant integral, it is precisely the
integrability condition for an untwisted affine highest-weight
representation of a simply-laced algebra.

\subsubsection{Dynkin-label form of the integrability
condition}\label{dynkin-label-form-of-the-integrability-condition}

Write \[
\vec\psi=\sum_i a_i\vec\alpha_i,
\qquad
\vec\lambda=\sum_j\lambda_j\vec\omega_j.
\] Since \[
\vec\alpha_i\cdot\vec\omega_j=\delta_{ij}
\] in the simply-laced normalization, \[
\vec\psi\cdot\vec\lambda
=\sum_{i,j}a_i\lambda_j
(\vec\alpha_i\cdot\vec\omega_j)
=\sum_i a_i\lambda_i.
\] Thus the allowed affine primary highest weights at level \(k\) obey
\[
\boxed{
\lambda_i\in\mathbb Z_{\geq0},
\qquad
\sum_i a_i\lambda_i\leq k.
}
\]

Equivalently, define the affine Dynkin label \[
\lambda_0
:=k-\vec\psi\cdot\vec\lambda
=k-\sum_i a_i\lambda_i.
\] Then affine integrability says that all affine Dynkin labels are
nonnegative integers: \[
\lambda_0,\lambda_1,\ldots,\lambda_r
\in\mathbb Z_{\geq0}.
\] The \(\mathfrak{su}(2)\) built from \[
E_1^{-\vec\psi},
\qquad
E_{-1}^{\vec\psi},
\qquad
k-\vec\psi\cdot\vec H_0
\] is precisely the \(\mathfrak{su}(2)\) associated with the affine
simple root \[
\vec\alpha_0=\delta-\vec\psi.
\]

\subsubsection{Examples of the bound}\label{examples-of-the-bound}

For \(A_1=\mathfrak{su}(2)\), let \[
\vec\lambda=n\vec\omega,
\qquad
n=2j\in\mathbb Z_{\geq0}.
\] Because \(\vec\psi=\vec\alpha=2\vec\omega\) and
\(\vec\alpha\cdot\vec\omega=1\), \[
\vec\psi\cdot\vec\lambda=n=2j.
\] The affine integrability condition is \[
n\leq k,
\qquad\text{or equivalently}\qquad
j\leq\frac{k}{2}.
\] At level one, only \(j=0\) and \(j=\frac12\) occur as affine primary
representations.

For \(A_2=\mathfrak{su}(3)\), write \[
\vec\lambda=p\vec\omega_1+q\vec\omega_2,
\qquad
p,q\in\mathbb Z_{\geq0}.
\] Since \[
\vec\psi=\vec\alpha_1+\vec\alpha_2,
\] one has \[
\vec\psi\cdot\vec\lambda=p+q.
\] Therefore \[
p+q\leq k.
\] At level one the possibilities are \[
(p,q)=(0,0),(1,0),(0,1),
\] corresponding to \[
\mathbf 1,
\qquad
\mathbf 3,
\qquad
\overline{\mathbf 3}.
\] For comparison, the adjoint \(\mathbf 8\) has highest weight \[
\vec\lambda=\vec\omega_1+\vec\omega_2=\vec\psi,
\] so \[
\vec\psi\cdot\vec\lambda=2.
\] It first occurs as an affine primary at level \(k=2\). At level one
the eight current states still exist, but they are level-one descendants
in the vacuum module, not a separate affine-primary module.

For \(D_r=\mathfrak{so}(2r)\), write \[
\vec\lambda=\sum_{i=1}^r\lambda_i\vec\omega_i.
\] The highest-root expansion above gives \[
\vec\psi\cdot\vec\lambda
=\lambda_1
+2\lambda_2
+\cdots
+2\lambda_{r-2}
+\lambda_{r-1}
+\lambda_r.
\] At level one, the only possibilities are \[
\vec\lambda=0,
\qquad
\vec\omega_1,
\qquad
\vec\omega_{r-1},
\qquad
\vec\omega_r.
\] These are the trivial, vector, and two chiral-spinor representations.

\subsubsection{Why only finitely many highest weights occur at fixed
level}\label{why-only-finitely-many-highest-weights-occur-at-fixed-level}

The Dynkin labels \(\lambda_i\) are nonnegative integers, and every mark
\(a_i\) is a positive integer. Therefore \[
\sum_i a_i\lambda_i\leq k
\] has only finitely many solutions. This is the elementary reason that
an affine WZW theory at positive integer level has only finitely many
integrable affine highest-weight representations.

At level one, \[
\sum_i a_i\lambda_i\leq1.
\] Because all \(a_i\geq1\) and all \(\lambda_i\geq0\) are integers,
either \[
\lambda_i=0
\qquad
\text{for every }i,
\] or there is exactly one node \(j\) such that \[
a_j=1,
\qquad
\lambda_j=1,
\] with all other Dynkin labels zero. Hence \[
\vec\lambda=0
\qquad\text{or}\qquad
\vec\lambda=\vec\omega_j
\text{ with }a_j=1.
\] For ADE root systems, these \(\vec\omega_j\) are precisely the
nonzero minuscule fundamental weights.

A dominant weight \(\vec\omega\) is minuscule when the weights of its
irreducible representation form a single Weyl orbit. Equivalently, for
every positive root \(\vec\alpha\), \[
\langle\vec\omega,\vec\alpha^{\vee}\rangle
\in\{0,1\}.
\] In particular, \[
\langle\vec\omega,\vec\psi^{\vee}\rangle=1,
\] which is exactly the nontrivial level-one condition in the
simply-laced case.

\subsubsection{Level-one ADE examples}\label{level-one-ade-examples}

The complete list is:

\begin{longtable}[]{@{}
  >{\raggedright\arraybackslash}p{(\linewidth - 6\tabcolsep) * \real{0.2308}}
  >{\raggedleft\arraybackslash}p{(\linewidth - 6\tabcolsep) * \real{0.3077}}
  >{\raggedright\arraybackslash}p{(\linewidth - 6\tabcolsep) * \real{0.2308}}
  >{\raggedright\arraybackslash}p{(\linewidth - 6\tabcolsep) * \real{0.2308}}@{}}
\toprule\noalign{}
\begin{minipage}[b]{\linewidth}\raggedright
Algebra
\end{minipage} & \begin{minipage}[b]{\linewidth}\raggedleft
\(P/Q\)
\end{minipage} & \begin{minipage}[b]{\linewidth}\raggedright
Integrable horizontal highest weights at \(k=1\)
\end{minipage} & \begin{minipage}[b]{\linewidth}\raggedright
Finite-dimensional representations
\end{minipage} \\
\midrule\noalign{}
\endhead
\bottomrule\noalign{}
\endlastfoot
\(A_r=\mathfrak{su}(r+1)\) & \(\mathbb Z_{r+1}\) &
\(0,\vec\omega_1,\ldots,\vec\omega_r\) & \(\mathbf 1\) and
\(\bigwedge^i\mathbf{(r+1)}\), \(i=1,\ldots,r\) \\
\(D_r=\mathfrak{so}(2r)\), \(r\) even & \(\mathbb Z_2\times\mathbb Z_2\)
& \(0,\vec\omega_1,\vec\omega_{r-1},\vec\omega_r\) & \(\mathbf 1\),
vector, and the two chiral spinors \\
\(D_r=\mathfrak{so}(2r)\), \(r\) odd & \(\mathbb Z_4\) &
\(0,\vec\omega_1,\vec\omega_{r-1},\vec\omega_r\) & \(\mathbf 1\),
vector, and the two chiral spinors \\
\(E_6\) & \(\mathbb Z_3\) & \(0,\vec\omega_1,\vec\omega_6\) &
\(\mathbf 1,\mathbf{27},\overline{\mathbf{27}}\) \\
\(E_7\) & \(\mathbb Z_2\) & \(0,\vec\omega_7\) &
\(\mathbf 1,\mathbf{56}\) \\
\(E_8\) & \(0\) & \(0\) & only \(\mathbf 1\) as an affine primary \\
\end{longtable}

The \(E_6\) and \(E_7\) node labels in the table use the standard
Bourbaki convention.

Some concrete special cases are:

\begin{itemize}
\tightlist
\item
  \(SU(2)_1\): \(\mathbf 1\) and \(\mathbf 2\).
\item
  \(SU(3)_1\): \(\mathbf 1\), \(\mathbf 3\), and
  \(\overline{\mathbf 3}\).
\item
  \(SU(4)_1\): \(\mathbf 1\), \(\mathbf 4\), \(\mathbf 6\), and
  \(\overline{\mathbf 4}\).
\item
  \(SO(8)_1\): \(\mathbf 1\), \(\mathbf 8_v\), \(\mathbf 8_s\), and
  \(\mathbf 8_c\).
\item
  \((E_6)_1\): \(\mathbf 1\), \(\mathbf{27}\), and
  \(\overline{\mathbf{27}}\).
\item
  \((E_7)_1\): \(\mathbf 1\) and \(\mathbf{56}\).
\item
  \((E_8)_1\): only the vacuum affine-primary module.
\end{itemize}

\subsubsection{Meaning of ``one lowest-dimensional representation in
each conjugacy
class''}\label{meaning-of-one-lowest-dimensional-representation-in-each-conjugacy-class}

The phrase ``conjugacy class'' here refers to a coset of the weight
lattice modulo the root lattice: \[
[\vec\lambda]\in P/Q.
\] Two highest weights are in the same class when \[
\vec\lambda-\vec\lambda'\in Q.
\] For the simply connected compact group \(G\) with Lie algebra \(g\),
\(P/Q\) is naturally related to the center \(Z(G)\); a weight-lattice
coset records the central character of a representation. It is not a
conjugacy class of elements of \(G\), and it is not merely the
distinction between a representation and its complex conjugate.

For simply-laced ADE algebras, each coset in \(P/Q\) has a distinguished
dominant representative of minimal size: the zero weight for the trivial
coset and a minuscule fundamental weight for every nontrivial coset.
These are exactly the level-one integrable highest weights listed above.
This explains the statement in the question.

For example, \[
P/Q\cong\mathbb Z_3
\] for \(SU(3)\). Its three cosets are represented at level one by \[
0,
\qquad
\vec\omega_1,
\qquad
\vec\omega_2,
\] corresponding to \[
\mathbf 1,
\qquad
\mathbf 3,
\qquad
\overline{\mathbf 3}.
\] Representations such as the \(\mathbf 6\) or \(\mathbf{10}\) lie in
one of the same three lattice cosets, but their highest weights have
larger \(\vec\psi\cdot\vec\lambda\) and therefore do not satisfy the
level-one bound.

For \(E_8\), one has \[
P=Q,
\qquad
P/Q=0.
\] There is only the trivial coset, whose minimal representative is
\(0\). Therefore \((E_8)_1\) has only the vacuum integrable affine
module. This does not mean that the theory has no \(E_8\) currents: the
\(248\) adjoint current states occur inside the vacuum module as
descendants.

\subsubsection{Checks}\label{checks-10}

For \(SU(2)\), the level-one inequality gives \[
2j\leq1,
\] so only \(j=0,\frac12\) are allowed. This agrees with the two classes
of \[
P/Q\cong\mathbb Z_2.
\]

For \(SU(3)\), the level-one inequality gives \[
p+q\leq1,
\] so only \((0,0),(1,0),(0,1)\) occur. This agrees with the three
classes of \[
P/Q\cong\mathbb Z_3.
\]

For \(SO(8)\), the highest-root marks are \[
(a_1,a_2,a_3,a_4)=(1,2,1,1)
\] in the standard \(D_4\) numbering. Hence the level-one weights are \[
0,\vec\omega_1,\vec\omega_3,\vec\omega_4,
\] which give \[
\mathbf 1,\mathbf 8_v,\mathbf 8_s,\mathbf 8_c.
\] This matches the four elements of \[
P/Q\cong\mathbb Z_2\times\mathbb Z_2.
\]

\subsubsection{Result}\label{result-14}

The weight-lattice integrality condition is \[
\langle\vec m,\vec\alpha^{\vee}\rangle
=\frac{2\vec\alpha\cdot\vec m}{\vec\alpha^{\,2}}
\in\mathbb Z.
\] For simply-laced \(g\) with \(\vec\alpha^{\,2}=2\), \[
\vec\alpha^{\vee}=\vec\alpha,
\qquad
\vec\alpha\cdot\vec m\in\mathbb Z.
\] The affine-root \(\mathfrak{su}(2)\) gives \[
k-\vec\alpha\cdot\vec m
\in\mathbb Z_{\geq0},
\] and therefore \[
k\in\mathbb Z_{\geq0}.
\]

For a horizontal irreducible representation with dominant highest weight
\(\vec\lambda\), the maximal root-weight pairing is \[
\max_{\vec\alpha\in\Delta,\,\vec m\in\mathrm{Wt}(V_{\vec\lambda})}
(\vec\alpha\cdot\vec m)
=\vec\psi\cdot\vec\lambda.
\] Thus affine integrability is \[
\boxed{
\vec\psi\cdot\vec\lambda
=\sum_i a_i\lambda_i
\leq k.
}
\] At level one, \[
\vec\lambda=0
\qquad\text{or}\qquad
\vec\lambda=\vec\omega_j
\text{ for a node with }a_j=1.
\] For ADE, these are the vacuum and the minuscule fundamental
representations, one minimal representative for each class in \(P/Q\).

Consider M-theory on a local CY4 given by \(\mathbb{P}^1\) with normal
bundle \(\mathcal{O}\oplus \mathcal{O}(-1)\oplus \mathcal{O}(-1)\), in
the 3d field theory, what is the susy multiplet given by a single
M2-brane wrapping this \(\mathbb{P}^1\)?

The moduli space of the \(\mathbb{P}^1\) is a Riemann surface \(\Sigma\)
extended in the direction of \(\mathcal{O}\).

\chapter{Covariant Vertex Operators, BRST, and Covariant
Lattices}\label{covariant-vertex-operators-brst-and-covariant-lattices-1}

\section{Q1: Why physical vertex operators have matter weight
one}\label{q1-why-physical-vertex-operators-have-matter-weight-one}

Why the vertex operators of physical massless states must have \(h=1\)?

\begin{quote}
\textbf{Textbook location and related material.} Chapter 13, Section
13.2 (pp.~403--411): BRST-invariant covariant vertex operators. Chapter
5, Section 5.2 (pp.~110--120) supplies the BRST cohomology, and Chapter
4 explains primary weights and integrated insertions.
\end{quote}

\subsection{Answer}\label{answer-15}

\subsubsection{Core point}\label{core-point-15}

The statement must first be phrased carefully. For the critical
\textbf{bosonic} string, it is the \textbf{matter part} of an on-shell
open-string vertex operator that has \[
h=1,
\] whereas the matter part of an on-shell closed-string vertex has \[
(h,\bar h)=(1,1).
\] The corresponding unintegrated, ghost-dressed operators are \[
cV_{\mathrm m}
\qquad\text{and}\qquad
c\bar cV_{\mathrm m},
\] and these have total conformal weight zero because \(c\) and
\(\bar c\) have weights \(-1\).

The number one arises in three equivalent descriptions of the same
constraint:

\begin{enumerate}
\def\labelenumi{\arabic{enumi}.}
\tightlist
\item
  an integrated vertex must be invariant under residual worldsheet
  conformal reparametrizations;
\item
  the state created by the vertex must obey the old covariant
  physical-state condition \((L_0-1)|\Psi\rangle=0\);
\item
  the ghost-dressed vertex must be BRST closed.
\end{enumerate}

However, \(h=1\) is not uniquely a masslessness condition. It is the
complete worldsheet on-shell condition. The bosonic tachyon vertex also
has \(h=1\). For a massless state, the oscillator or derivative part
already contributes one unit of conformal weight, while the momentum
exponential contributes zero because \(k^2=0\).

\subsubsection{Residual reparametrization invariance of an integrated
vertex}\label{residual-reparametrization-invariance-of-an-integrated-vertex}

After fixing the worldsheet metric to conformal gauge, one still has
residual conformal coordinate changes. Consider first an open-string
vertex inserted on a boundary coordinate \(x\). If \(V_{\mathrm m}(x)\)
is a boundary primary of weight \(h\), then under \(x\mapsto x'=f(x)\)
it transforms as \[
V'_{\mathrm m}(x')
=\left(\frac{dx}{dx'}\right)^h V_{\mathrm m}(x).
\] An integrated vertex appears in an amplitude as \[
\int dx\,V_{\mathrm m}(x).
\] Since \(dx'=(dx'/dx)dx\), the transformed insertion is \[
\int dx'\,V'_{\mathrm m}(x')
=\int dx\,
\left(\frac{dx'}{dx}\right)^{1-h}V_{\mathrm m}(x).
\] It is invariant for an arbitrary residual reparametrization precisely
when \[
h=1.
\] Thus the matter operator must be a marginal boundary operator.

For a closed-string insertion on a complex worldsheet coordinate
\((z,\bar z)\), \[
\int d^2z\,V_{\mathrm m}(z,\bar z),
\] the measure transforms with one holomorphic and one antiholomorphic
Jacobian. If \(V_{\mathrm m}\) has weights \((h,\bar h)\), invariance
under independent local maps \(z\mapsto f(z)\) and
\(\bar z\mapsto\bar f(\bar z)\) requires \[
h=1,
\qquad
\bar h=1.
\] Therefore an integrated closed-string vertex is a marginal bulk
operator of weights \((1,1)\).

This argument also explains the unintegrated form. Gauge fixing the
positions of some vertex insertions introduces Faddeev--Popov ghosts.
Since \[
h[c]=-1,
\qquad
\bar h[\bar c]=-1,
\] one has \[
h[cV_{\mathrm m}]=-1+1=0
\] for the open string, and \[
(h,\bar h)[c\bar cV_{\mathrm m}]=(0,0)
\] for the closed string. Hence the fixed, ghost-dressed insertion is a
worldsheet scalar, as required.

\subsubsection{The same condition in old covariant
quantization}\label{the-same-condition-in-old-covariant-quantization}

By the state--operator correspondence, a matter operator inserted at the
origin creates a matter state: \[
|\Psi\rangle=V_{\mathrm m}(0)|0\rangle.
\] If \(V_{\mathrm m}\) is a Virasoro primary of weight \(h\), then \[
L_0|\Psi\rangle=h|\Psi\rangle,
\qquad
L_n|\Psi\rangle=0
\quad(n>0).
\] In old covariant quantization, the bosonic open-string physical-state
conditions are \[
(L_n-a\delta_{n0})|\Psi\rangle=0,
\qquad n\geq0,
\] with the critical intercept \(a=1\). The \(n=0\) condition is
therefore \[
(L_0-1)|\Psi\rangle=0,
\] so the corresponding matter operator has \(h=1\). The conditions with
\(n>0\) impose primarity and remove negative-norm polarizations; they
contain information such as transversality that is not captured by the
value of \(h\) alone.

For an open-string state of momentum \(k\) and oscillator level \(N\),
\[
L_0=\alpha'k^2+N.
\] Thus \[
\alpha'k^2+N=1.
\] Using \(k^2=-m^2\) gives \[
m^2=\frac{N-1}{\alpha'}.
\] The condition \(h=1\) therefore includes every on-shell level:
\(N=0\) is the tachyon and \(N=1\) is the massless level.

For a closed-string state, \[
L_0=\frac{\alpha'k^2}{4}+N,
\qquad
\bar L_0=\frac{\alpha'k^2}{4}+\bar N.
\] The two physical-state conditions are \[
(L_0-1)|\Psi\rangle=0,
\qquad
(\bar L_0-1)|\Psi\rangle=0.
\] Equivalently, \[
\frac{\alpha'k^2}{4}+N=1,
\qquad
\frac{\alpha'k^2}{4}+\bar N=1.
\] Subtracting the two equations gives level matching \(N=\bar N\),
while either equation gives \[
m^2=\frac{4(N-1)}{\alpha'}
=\frac{4(\bar N-1)}{\alpha'}.
\] Again, \((h,\bar h)=(1,1)\) is the full on-shell condition;
masslessness is the particular level \(N=\bar N=1\).

\subsubsection{BRST closure and the role of the
ghosts}\label{brst-closure-and-the-role-of-the-ghosts}

The BRST description packages the Virasoro constraints and their gauge
equivalences into the cohomology of the BRST charge \(Q_B\). Let
\(V_{\mathrm m}\) be a matter primary of weight \(h\). The operator
product with the matter stress tensor is \[
T_{\mathrm m}(z)V_{\mathrm m}(w)
\sim
\frac{hV_{\mathrm m}(w)}{(z-w)^2}
+\frac{\partial V_{\mathrm m}(w)}{z-w}.
\] Evaluating the BRST contour integral on the ghost-number-one operator
\(cV_{\mathrm m}\) gives, up to an overall sign convention, \[
Q_B\bigl(cV_{\mathrm m}\bigr)
\propto
(1-h)c\partial c\,V_{\mathrm m}.
\] Consequently, \[
Q_B(cV_{\mathrm m})=0
\] requires \(h=1\), provided \(V_{\mathrm m}\) is a matter primary. If
the candidate operator is not primary, additional poles in its
stress-tensor OPE generate additional BRST constraints. For the massless
vector, those extra conditions give transversality.

For the closed string there are independent left- and right-moving BRST
charges. Acting on \[
c\bar cV_{\mathrm m}
\] tests the two sectors separately, so BRST closure requires \[
(h,\bar h)=(1,1)
\] together with left- and right-moving primarity. Two BRST-closed
operators differing by a BRST-exact operator describe the same physical
state: \[
\mathcal V\sim\mathcal V+Q_B\Lambda.
\] This quotient is the covariant origin of spacetime gauge equivalence.

Thus there is no contradiction between saying that the matter vertex has
weight one and saying that a physical unintegrated vertex has weight
zero. They refer to different operators: \[
V_{\mathrm m}: h=1,
\qquad
cV_{\mathrm m}: h_{\mathrm{tot}}=0.
\]

\subsubsection{Free-boson conformal weights and why weight one is not
synonymous with
massless}\label{free-boson-conformal-weights-and-why-weight-one-is-not-synonymous-with-massless}

For a free boson in the stated normalization, the exponential has
open-string boundary weight \[
h_{\mathrm{open}}\bigl[:e^{ik\cdot X}:\bigr]=\alpha'k^2,
\] whereas a closed-string exponential has \[
h_{\mathrm{closed}}\bigl[:e^{ik\cdot X}:\bigr]
=\bar h_{\mathrm{closed}}\bigl[:e^{ik\cdot X}:\bigr]
=\frac{\alpha'k^2}{4}.
\] Every holomorphic derivative \(\partial X^\mu\) contributes one unit
of \(h\), and every antiholomorphic derivative \(\bar\partial X^\nu\)
contributes one unit of \(\bar h\). More generally, oscillator
excitations contribute their oscillator level. Hence \[
h_{\mathrm{open}}=N+\alpha'k^2,
\] and \[
h_{\mathrm{closed}}=N+\frac{\alpha'k^2}{4},
\qquad
\bar h_{\mathrm{closed}}=\bar N+\frac{\alpha'k^2}{4}.
\]

For the open-string tachyon, \[
V_T=:e^{ik\cdot X}:,
\] there is no derivative contribution. The condition \(h=1\) becomes \[
\alpha'k^2=1,
\] or \[
m^2=-\frac1{\alpha'}.
\] Thus this tachyonic matter vertex has weight one when it is on shell.

Similarly, the closed-string tachyon vertex \[
V_T=:e^{ik\cdot X}:
\] has \((h,\bar h)=(1,1)\) when \[
\frac{\alpha'k^2}{4}=1,
\] which gives \[
m^2=-\frac4{\alpha'}.
\]

By contrast, a massless vertex contains a level-one oscillator or
derivative in each required chiral sector. That part supplies the unit
weight, while \(k^2=0\) makes the exponential contribute zero. This is
the precise sense in which the massless vertex has weight one.

\subsubsection{Open-string massless
vector}\label{open-string-massless-vector}

The matter vertex for an open-string vector can be written on the
boundary as \[
V_A(x)=\epsilon_\mu:\!\partial_tX^\mu e^{ik\cdot X}\!:(x),
\] where \(\partial_t\) is tangent to the boundary. Its weight is \[
h[V_A]=1+\alpha'k^2.
\] Requiring \(h=1\) gives \[
k^2=0,
\] so the state is massless. But one must also require that \(V_A\) be
primary. The potentially unwanted higher-order pole in \(T(z)V_A(x)\)
vanishes only if \[
k\cdot\epsilon=0.
\] This is the same condition obtained from \(L_1|\Psi\rangle=0\) for \[
|\Psi\rangle=\epsilon_\mu\alpha_{-1}^\mu|0;k\rangle.
\]

A longitudinal polarization is physically trivial. Indeed, \[
k_\mu\partial_tX^\mu e^{ik\cdot X}
=\frac1{i}\partial_t e^{ik\cdot X},
\] so its integrated vertex is a boundary total derivative. In BRST
language the same state is BRST exact. Therefore \[
\epsilon_\mu\sim\epsilon_\mu+\lambda k_\mu,
\] which is the momentum-space gauge equivalence of a massless vector.

\subsubsection{Closed-string massless tensor
states}\label{closed-string-massless-tensor-states}

At the first excited closed-string level the general matter vertex is \[
V_\epsilon(z,\bar z)
=\epsilon_{\mu\nu}:
\partial X^\mu\bar\partial X^\nu e^{ik\cdot X}:
\] and has weights \[
(h,\bar h)
=\left(1+\frac{\alpha'k^2}{4},
1+\frac{\alpha'k^2}{4}\right).
\] The condition \((h,\bar h)=(1,1)\) gives \[
k^2=0.
\] Left- and right-moving primarity further require \[
k^\mu\epsilon_{\mu\nu}=0,
\qquad
k^\nu\epsilon_{\mu\nu}=0.
\] BRST-exact states generate the polarization equivalence \[
\epsilon_{\mu\nu}
\sim
\epsilon_{\mu\nu}
+k_\mu\xi_\nu
+k_\nu\widetilde\xi_\mu,
\] with the usual on-shell restrictions on the gauge parameters. The
symmetric traceless, antisymmetric, and trace combinations give the
graviton, Kalb--Ramond field, and dilaton. Their spacetime gauge
redundancies are the symmetric and antisymmetric parts of the displayed
BRST equivalence.

\subsubsection{Checks}\label{checks-11}

\begin{enumerate}
\def\labelenumi{\arabic{enumi}.}
\item
  \textbf{Integrated versus unintegrated:} \(V_{\mathrm m}\) has \(h=1\)
  so that \(dx\,V_{\mathrm m}\) is invariant; \(cV_{\mathrm m}\) has
  total weight zero so that a fixed insertion is invariant. These
  statements are complementary.
\item
  \textbf{Tachyon check:} setting \(N=0\) in \(h=N+\alpha'k^2=1\) gives
  the open-string tachyon mass, not a massless state. Therefore weight
  one cannot by itself mean massless.
\item
  \textbf{Massless check:} setting \(N=1\) gives \(k^2=0\). The
  oscillator supplies one unit of weight and the exponential supplies
  zero.
\item
  \textbf{Closed-string check:} both sectors must independently have
  weight one. Requiring only \(h+\bar h=2\) would be too weak because it
  would not enforce the two Virasoro constraints or level matching
  separately.
\item
  \textbf{Superstring caveat:} in the RNS string the matter and
  superghost contributions depend on the picture. For example, in the NS
  \(-1\) picture, \(\psi^\mu e^{ik\cdot X}\) has matter weight \(1/2\)
  on shell and \(e^{-\phi}\) supplies another \(1/2\), while \(c\)
  supplies \(-1\). Thus the invariant statement is that the complete
  integrated or ghost-dressed vertex has the required total weight; the
  simple phrase ``the matter vertex has \(h=1\)'' is literally the
  bosonic-string statement, or a statement made after specifying the
  superghost picture.
\end{enumerate}

\subsubsection{Result}\label{result-15}

For the bosonic open string, \[
\boxed{
h[V_{\mathrm m}]=1,
\qquad
h[cV_{\mathrm m}]=0
}
\] because \[
\int dx\,V_{\mathrm m}
\] must be invariant, equivalently \[
(L_0-1)|\Psi\rangle=0,
\qquad
Q_B(cV_{\mathrm m})=0.
\] With \[
h=N+\alpha'k^2,
\] the on-shell condition gives \[
m^2=\frac{N-1}{\alpha'}.
\]

For the bosonic closed string, \[
\boxed{
(h,\bar h)[V_{\mathrm m}]=(1,1),
\qquad
(h,\bar h)[c\bar cV_{\mathrm m}]=(0,0)
}
\] with \[
h=N+\frac{\alpha'k^2}{4},
\qquad
\bar h=\bar N+\frac{\alpha'k^2}{4}.
\] The two conditions imply level matching and the closed-string mass
shell.

For massless states specifically, \[
N=1
\quad\text{(open)},
\qquad
N=\bar N=1
\quad\text{(closed)},
\qquad
k^2=0.
\] Therefore the derivative or oscillator part contributes the required
unit weight, while the momentum exponential contributes zero. The value
\(h=1\) expresses the worldsheet physical-state condition; masslessness
is the special level-one solution of that condition.

\section{Q2: Conformal weights and pictures in the RNS
superstring}\label{q2-conformal-weights-and-pictures-in-the-rns-superstring}

I know the case of the bosonic string theory, but what about the
superstring, do their vertex operators have h=1 and why?

\begin{quote}
\textbf{Textbook location and related material.} Chapter 13, Sections
13.1--13.2 (pp.~391--411): bosonized superghosts, picture number, BRST
invariance, and picture changing. The NS/R intercepts come from Chapter
8.
\end{quote}

\subsection{Answer}\label{answer-16}

\subsubsection{Core point}\label{core-point-16}

Yes, but in the RNS superstring one must say \textbf{which part} of the
vertex operator is being assigned weight one. The universal statement
is:

\begin{itemize}
\tightlist
\item
  an integrated open-string insertion, including its matter and
  \(\beta\gamma\) superghost factors but not a \(c\) ghost, must have
  total weight \(h=1\);
\item
  an integrated closed-string insertion must have total weights
  \((h,\bar h)=(1,1)\);
\item
  after multiplying by \(c\) for an open-string fixed insertion, or by
  \(c\bar c\) for a closed-string fixed insertion, the unintegrated
  operator has total weight zero.
\end{itemize}

What changes relative to the bosonic string is that the RNS vertex
carries a \textbf{picture number}. The matter factor alone need not have
\(h=1\), because the bosonized superghost exponential \(e^{q\phi}\) also
contributes conformal weight. For example, the massless NS matter factor
\(\epsilon\cdot\psi e^{ik\cdot X}\) has \(h=1/2\) on shell in picture
\(-1\), and \(e^{-\phi}\) contributes the remaining \(1/2\). In picture
zero the superghost exponential is absent and the massless NS matter
operator itself has \(h=1\).

Thus picture changing redistributes conformal weight between the matter
and superghost factors without changing the conformal weight or BRST
cohomology class of the complete physical insertion.

\subsubsection{Why the complete integrated insertion still has weight
one}\label{why-the-complete-integrated-insertion-still-has-weight-one}

The residual reparametrization argument is unchanged by worldsheet
supersymmetry. On an open-string boundary, an integrated insertion is \[
\int dx\,\mathcal V(x),
\] where \(\mathcal V\) now denotes the complete local RNS operator
apart from the \(c\) ghost. If \(\mathcal V\) has weight \(h\), then
under \(x\mapsto x'=f(x)\), \[
dx'\,\mathcal V'(x')
=dx\left(\frac{dx'}{dx}\right)^{1-h}\mathcal V(x).
\] Invariance for arbitrary residual conformal maps requires \[
h[\mathcal V]=1.
\]

For a closed-string bulk insertion, \[
\int d^2z\,\mathcal V(z,\bar z),
\] the two Jacobians require \[
(h,\bar h)[\mathcal V]=(1,1).
\] The reparametrization ghosts still obey \[
h[c]=-1,
\qquad
\bar h[\bar c]=-1.
\] Therefore the corresponding fixed insertions have \[
h[c\mathcal V]=0
\] for an open string and \[
(h,\bar h)[c\bar c\mathcal V]=(0,0)
\] for a closed string.

Worldsheet supersymmetry introduces additional gauge fixing, namely the
\(\beta\gamma\) superghost system. Its factors belong to \(\mathcal V\)
in the above weight count. This is why ``the matter vertex has weight
one'' is not a picture-independent superstring statement, whereas ``the
complete integrated vertex has weight one'' is.

\subsubsection{Bosonized superghost
weights}\label{bosonized-superghost-weights}

The commuting \(\beta\gamma\) system has weights \[
h[\beta]=\frac32,
\qquad
h[\gamma]=-\frac12.
\] Bosonize it as \[
\beta=e^{-\phi}\partial\xi,
\qquad
\gamma=\eta e^\phi,
\] where \[
\phi(z)\phi(w)\sim-\ln(z-w).
\] The scalar \(\phi\) has a background charge, so its stress tensor is
not that of an ordinary free scalar: \[
T_\phi(z)
=-\frac12:(\partial\phi)^2:(z)-\partial^2\phi(z).
\]

To compute the weight of \(e^{q\phi}\), first use \[
\partial\phi(z)e^{q\phi(w)}
\sim-\frac{q}{z-w}e^{q\phi(w)}.
\] The double contraction with the first term in \(T_\phi\) contributes
\[
-\frac{q^2}{2(z-w)^2}e^{q\phi(w)},
\] while the background-charge term contributes \[
-\frac{q}{(z-w)^2}e^{q\phi(w)}.
\] Hence \[
T_\phi(z)e^{q\phi(w)}
\sim
\frac{-\frac12q(q+2)}{(z-w)^2}e^{q\phi(w)}
+\frac{1}{z-w}\partial e^{q\phi(w)},
\] and therefore \[
\boxed{
h[e^{q\phi}]=-\frac12q(q+2)
}.
\] The cases needed below are \[
h[e^{-\phi}]=\frac12,
\qquad
h[e^{-\phi/2}]=\frac38,
\qquad
h[e^{0\phi}]=0.
\]

The background charge also produces a picture-number selection rule in
amplitudes. For example, tree-level amplitudes are commonly arranged to
have total picture number \(-2\) in each relevant chiral superghost
sector. This affects how pictures are distributed among the vertices,
but it does not change the local condition that every complete
integrated insertion have weight one.

\subsubsection{Massless NS vector in picture minus
one}\label{massless-ns-vector-in-picture-minus-one}

The unintegrated open-string massless vector in the NS sector and
picture \(-1\) is \[
U_A^{(-1)}(x)
=c\,e^{-\phi}\epsilon_\mu\psi^\mu e^{ik\cdot X}(x).
\] For a massless state, \(k^2=0\), so the open-string exponential
contributes \[
h[e^{ik\cdot X}]=\alpha'k^2=0.
\] The remaining weights are \[
h[c]=-1,
\qquad
h[e^{-\phi}]=\frac12,
\qquad
h[\epsilon\cdot\psi]=\frac12.
\] Thus \[
h[U_A^{(-1)}]
=-1+\frac12+\frac12+0=0.
\] Removing \(c\) gives the complete integrated local operator \[
V_A^{(-1)}
=e^{-\phi}\epsilon\cdot\psi e^{ik\cdot X},
\] whose total weight is \[
h[V_A^{(-1)}]=1.
\]

This example exhibits the main point particularly clearly: \[
\underbrace{h[\epsilon\cdot\psi e^{ik\cdot X}]}_{\text{matter}}=\frac12,
\qquad
\underbrace{h[e^{-\phi}]}_{\text{superghost}}=\frac12,
\qquad
\underbrace{h[V_A^{(-1)}]}_{\text{complete integrated operator}}=1.
\]

BRST closure imposes more than the weight condition. For this vector it
gives \[
k^2=0,
\qquad
k\cdot\epsilon=0.
\] BRST-exact vertices identify \[
\epsilon_\mu\sim\epsilon_\mu+\lambda k_\mu,
\] which is the spacetime gauge equivalence. Thus conformal weight,
transversality, and gauge redundancy are different parts of the same
BRST-cohomology problem.

\subsubsection{The same NS vector in picture
zero}\label{the-same-ns-vector-in-picture-zero}

In picture zero, the standard integrated massless vector has the
schematic form \[
V_A^{(0)}(x)
=\epsilon_\mu:
\left(
i\partial_tX^\mu
+C_{\psi}\alpha'(k\cdot\psi)\psi^\mu
\right)e^{ik\cdot X}:(x).
\] The numerical coefficient \(C_{\psi}\), as well as factors of \(i\),
depends on the normalization of \(X\), \(\psi\), and the boundary
coordinate; BRST closure fixes it once those conventions are chosen. The
weight count does not depend on that coefficient. On the mass shell, \[
h[\partial_tX^\mu]=1,
\qquad
h[(k\cdot\psi)\psi^\mu]=1,
\qquad
h[e^{ik\cdot X}]=0.
\] Therefore \[
h[V_A^{(0)}]=1.
\] There is no \(e^{q\phi}\) factor in picture zero, so in this picture
the displayed operator is entirely a matter operator and its weight is
one.

At noncanonical pictures, especially for unintegrated representatives,
picture changing can also generate additional ghost terms needed for
exact BRST closure. Therefore one should regard the simple expression
above as the standard \textbf{integrated} picture-zero representative,
rather than assume that every unintegrated representative is obtained
merely by multiplying it by \(c\) and nothing else.

\subsubsection{Picture changing preserves the physical conformal
weight}\label{picture-changing-preserves-the-physical-conformal-weight}

The picture-changing operator is \[
\mathcal X(z)=\{Q_B,\xi(z)\}.
\] It has picture number \(+1\) and conformal weight zero: \[
h[\mathcal X]=0.
\] Acting with \(\mathcal X\) on a BRST-closed picture-\(q\) vertex
produces a picture-\((q+1)\) representative of the same BRST cohomology
class, up to the usual qualifications about singular operator products
and zero-momentum exceptional states. Since \(\mathcal X\) has weight
zero, \[
h[\mathcal X\mathcal V]=h[\mathcal V].
\] In particular, picture changing maps the picture-\(-1\) NS vertex to
the picture-zero vertex, up to BRST-exact terms and total derivatives,
without changing the total weight of the complete insertion.

Picture changing therefore does not turn a nonphysical operator into a
physical one by altering its weight. It only changes how the fixed total
weight is divided among matter and superghost factors.

\subsubsection{Ramond vertex in picture minus one
half}\label{ramond-vertex-in-picture-minus-one-half}

The unintegrated massless open-string Ramond vertex in the canonical
picture \(-1/2\) is \[
U_u^{(-1/2)}(x)
=c\,e^{-\phi/2}u_AS^A e^{ik\cdot X}(x).
\] In ten dimensions the spin field has \[
h[S^A]=\frac{D}{16}=\frac58,
\] and the superghost factor has \[
h[e^{-\phi/2}]=\frac38.
\] For \(k^2=0\) the exponential again has zero weight. Hence \[
h[U_u^{(-1/2)}]
=-1+\frac38+\frac58+0=0,
\] whereas the integrated operator without \(c\) has \[
h[e^{-\phi/2}u_AS^Ae^{ik\cdot X}]
=\frac38+\frac58=1.
\]

BRST closure imposes the massless spacetime Dirac equation \[
k_\mu(\Gamma^\mu)^A{}_B u_A=0,
\] with index placement depending on the charge-conjugation and spinor
conventions. The GSO projection further selects the appropriate
spacetime chirality. Thus the equality \(3/8+5/8=1\) verifies conformal
covariance, while the Dirac equation is an additional BRST constraint.

\subsubsection{Relation to the NS and R
intercepts}\label{relation-to-the-ns-and-r-intercepts}

In old covariant quantization, the open RNS mass-shell condition is \[
(L_0-a)|\Psi\rangle=0.
\] In the NS sector, \[
a_{\mathrm{NS}}=\frac12,
\qquad
L_0=\alpha'k^2+N_{\mathrm{NS}},
\] so \[
\alpha'k^2+N_{\mathrm{NS}}=\frac12.
\] The GSO-projected massless vector has \(N_{\mathrm{NS}}=1/2\), hence
\(k^2=0\). Its plane matter operator \(\epsilon\cdot\psi e^{ik\cdot X}\)
correspondingly has weight \(1/2\), not one. The picture-\(-1\)
superghost dressing supplies the other \(1/2\) needed by the complete
integrated vertex.

In the Ramond sector, \[
a_{\mathrm R}=0,
\qquad
\alpha'k^2+N_{\mathrm R}=0.
\] The Ramond ground state has \(N_{\mathrm R}=0\) and is massless. One
must not naively identify the plane spin-field weight \(h[S^A]=5/8\)
with the cylinder intercept \(a_{\mathrm R}=0\). The spin field is a
twist operator that changes the boundary conditions of the worldsheet
fermions, and the plane-to-cylinder map includes the corresponding
sector-dependent vacuum-energy shift. The reliable local vertex-operator
statement is instead \[
h[S^A]+h[e^{-\phi/2}]
=\frac58+\frac38=1.
\] The old-covariant intercept and the plane spin-field weight are
therefore compatible, but they encode the Ramond vacuum relative to
different reference vacua.

\subsubsection{Closed NS--NS and heterotic
examples}\label{closed-nsns-and-heterotic-examples}

For a closed type-II NS--NS massless state in picture \((-1,-1)\), the
unintegrated vertex is \[
U_\epsilon^{(-1,-1)}
=c\bar c\,e^{-\phi}e^{-\bar\phi}
\epsilon_{\mu\nu}\psi^\mu\bar\psi^\nu e^{ik\cdot X}.
\] At \(k^2=0\), each chiral sector has the weight count \[
-1+\frac12+\frac12=0.
\] Thus the full unintegrated operator has weights \((0,0)\), while
removing \(c\bar c\) gives an integrated operator of weights \[
(h,\bar h)=(1,1).
\] BRST closure further yields \[
k^\mu\epsilon_{\mu\nu}=0,
\qquad
k^\nu\epsilon_{\mu\nu}=0,
\] and BRST-exact shifts generate \[
\epsilon_{\mu\nu}
\sim
\epsilon_{\mu\nu}
+k_\mu\xi_\nu
+k_\nu\widetilde\xi_\mu.
\] The symmetric and antisymmetric parts are the linearized gauge
equivalences of the graviton and Kalb--Ramond field.

The heterotic string simply combines two different chiral constructions.
In the conventional assignment, one sector is bosonic and the other is
RNS supersymmetric. A physical integrated heterotic vertex must still
have total weights \((1,1)\): the bosonic side obtains its unit weight
from a derivative, current, or oscillator factor, while the
supersymmetric side may divide its unit weight between a matter operator
and an \(e^{q\phi}\) superghost dressing. Multiplication by \(c\bar c\)
again produces an unintegrated operator of weights \((0,0)\).

\subsubsection{BRST closure is stronger than conformal weight
alone}\label{brst-closure-is-stronger-than-conformal-weight-alone}

Schematically, if \(\mathcal V\) is a complete open-string operator of
weight \(h\) before inserting \(c\), the \(cT\) part of the RNS BRST
charge gives \[
Q_B(c\mathcal V)
\supset
(1-h)c\partial c\,\mathcal V.
\] Therefore BRST closure requires \(h=1\). But the RNS BRST charge also
contains terms built from the supercurrent and the \(\beta\gamma\)
ghosts. Their operator products impose the super-Virasoro constraints
that are not visible from \(L_0\) alone. Depending on the state, these
give transversality, the Dirac equation, tracelessness conditions, or
analogous polarization constraints.

Finally, physical vertices are cohomology classes, \[
Q_B U=0,
\qquad
U\sim U+Q_B\Lambda.
\] The second relation is what removes longitudinal or otherwise
gauge-redundant polarizations. Thus having the correct conformal weight
is necessary, but it is not sufficient without BRST closure modulo
BRST-exact states.

\subsubsection{Checks}\label{checks-12}

\begin{enumerate}
\def\labelenumi{\arabic{enumi}.}
\item
  \textbf{NS picture \(-1\):} \[
  h[c]+h[e^{-\phi}]+h[\psi]+h[e^{ik\cdot X}]
  =-1+\frac12+\frac12+0=0.
  \]
\item
  \textbf{NS picture zero:} both \(\partial X\) and \((k\cdot\psi)\psi\)
  have weight one, so the integrated massless vertex has \(h=1\) without
  a superghost exponential.
\item
  \textbf{Ramond picture \(-1/2\):} \[
  h[e^{-\phi/2}]+h[S^A]
  =\frac38+\frac58=1.
  \]
\item
  \textbf{Closed NS--NS sector:} each chiral half independently has
  integrated weight one; it is not enough to require only
  \(h+\bar h=2\).
\item
  \textbf{Picture changing:} \(h[\mathcal X]=0\), so changing picture
  preserves the complete conformal weight even though it changes the
  separate matter and superghost contributions.
\item
  \textbf{Intercept check:} \(a_{\mathrm{NS}}=1/2\) explains the NS
  matter weight \(1/2\) at the massless level, but \(a_{\mathrm R}=0\)
  must not be identified directly with the plane spin-field weight
  \(5/8\).
\end{enumerate}

\subsubsection{Result}\label{result-16}

For an RNS open string, the picture-independent weight statement is \[
\boxed{
h[\mathcal V_{\mathrm{integrated}}]=1,
\qquad
h[c\mathcal V_{\mathrm{integrated}}]=0
}.
\] For an RNS closed string, \[
\boxed{
(h,\bar h)[\mathcal V_{\mathrm{integrated}}]=(1,1),
\qquad
(h,\bar h)[c\bar c\mathcal V_{\mathrm{integrated}}]=(0,0)
}.
\] The superghost exponential contributes \[
\boxed{
h[e^{q\phi}]=-\frac12q(q+2)
}.
\] Therefore the canonical massless vertices satisfy \[
\begin{aligned}
\text{NS, picture }-1:
&\qquad
h[e^{-\phi}]+h[\psi]
=\frac12+\frac12=1,\\%1
\text{NS, picture }0:
&\qquad
h[\partial X]
=h[\psi\psi]=1,\\%1
\text{R, picture }-\frac12:
&\qquad
h[e^{-\phi/2}]+h[S]
=\frac38+\frac58=1.
\end{aligned}
\] Picture changing by \[
\mathcal X=\{Q_B,\xi\},
\qquad
h[\mathcal X]=0,
\] preserves the complete conformal weight and the BRST cohomology
class. BRST closure then adds the state-specific equations such as \[
k^2=0,
\qquad
k\cdot\epsilon=0,
\qquad
k\!\!\!/\,u=0,
\] while BRST-exact shifts implement spacetime gauge equivalence. Hence
superstring vertex operators obey the same geometric weight-one
requirement as bosonic vertices, but only the \textbf{complete
picture-dressed insertion}, not the matter factor by itself, satisfies
that requirement in every picture.

\section{Q3: Deriving the upper-half-plane fermion gluing
condition}\label{q3-deriving-the-upper-half-plane-fermion-gluing-condition}

In chapter 6.3 we have the boundary conditions for open superstring \[
\begin{aligned}
\psi_{+}^\mu(0)&={D^\mu}_{\nu}(0)\psi_{-}^\nu(0)\\
\psi_{+}^\mu(l)&=\eta{D^\mu}_{\nu}(l)\psi_{-}^\nu(l)\\
\end{aligned}
\] Where \({D^\mu}_{\nu}=({\delta^\alpha}_{\beta},-{\delta^i}_{j})\) for
\(\alpha,\beta\) the NN directions and \(i,j\) the DD directions;
\(\eta=+1\) in R sector and \(\eta=-1\) in NS sector. In chapter 13.3
they use this boundary condition as \[
\psi^\mu(z)=\eta{D^\mu}_{\nu}\bar{\psi}^\nu(\bar{z})|_{\mathrm{Im}(z)=0}
\] Where \(\eta=+1\) in NS sector and
\(\eta=\mathrm{sign}(\mathrm{Re}(z))\) in R sector. Derive this equation
from the original boundary condition.

\begin{quote}
\textbf{Textbook location and related material.} Chapter 13, Section
13.3 (pp.~411--415), especially the upper-half-plane fermion gluing
condition near Eq. (13.99). The original strip boundary conditions are
in Section 7.4.
\end{quote}

\subsection{Answer}\label{answer-17}

\subsubsection{Core point}\label{core-point-17}

The two formulas are the same boundary condition written in two
coordinate systems. There are two changes between them:

\begin{enumerate}
\def\labelenumi{\arabic{enumi}.}
\tightlist
\item
  the two edges \(\sigma=0\) and \(\sigma=l\) of the open-string strip
  become the positive and negative halves of the real axis of the upper
  half-plane;
\item
  \(\psi^\mu\) is a conformal field of weight \(1/2\), so its
  transformation law contains a square root of the Jacobian. That square
  root differs by a sign between the positive and negative real axes.
\end{enumerate}

It is useful to distinguish the sign in the strip formula from the sign
in the upper-half-plane formula. Write \[
\eta_{\mathrm{strip}}
=
\begin{cases}
+1,&\mathrm R,\\
-1,&\mathrm{NS}.
\end{cases}
\] Then the strip boundary conditions are \[
\psi_+^\mu(0)
={D_0^\mu}_{\nu}\psi_-^\nu(0),
\qquad
\psi_+^\mu(l)
=\eta_{\mathrm{strip}}{D_l^\mu}_{\nu}\psi_-^\nu(l).
\] After mapping to the upper half-plane, the general result is \[
\psi^\mu(x)
=
\begin{cases}
{D_0^\mu}_{\nu}\bar\psi^\nu(x),&x>0,\\[2mm]
-\eta_{\mathrm{strip}}{D_l^\mu}_{\nu}\bar\psi^\nu(x),&x<0.
\end{cases}
\] For the single-D-brane situation used in Chapter 13, one has \[
D_0=D_l=D.
\] It follows that \[
\psi^\mu(x)
=
\begin{cases}
{D^\mu}_{\nu}\bar\psi^\nu(x),&\mathrm{NS},\\[2mm]
\operatorname{sign}(x){D^\mu}_{\nu}\bar\psi^\nu(x),&\mathrm R,
\end{cases}
\] which is precisely Eq. (13.99).

In the supplied PDF, the original strip formula is Eq. (7.62) in §7.3,
and the book itself refers back to §7.3 immediately before Eq. (13.99).
The relevant conformal map is Eq. (4.138). The chapter number in the
question may refer to a different numbering scheme, but the displayed
equations identify the passage unambiguously.

\subsubsection{What the book states and what must be
supplied}\label{what-the-book-states-and-what-must-be-supplied}

The book explicitly gives three ingredients:

\begin{itemize}
\tightlist
\item
  Eq. (7.62), the gluing conditions at \(\sigma=0\) and \(\sigma=l\);
\item
  Eq. (4.138), the exponential map from a strip to the upper half-plane;
\item
  immediately before Eq. (13.99), the statement that the worldsheet
  fermions are half-differentials and therefore acquire a nontrivial
  Jacobian.
\end{itemize}

The omitted intermediate step is the evaluation of the square-root
Jacobian on the two halves of the real axis. That evaluation is what
changes \[
\eta_{\mathrm{strip}}
=
\begin{cases}
+1,&\mathrm R,\\
-1,&\mathrm{NS}
\end{cases}
\] into \[
\eta_{\mathrm{UHP}}(x)
=
\begin{cases}
+1,&\mathrm{NS},\\
\operatorname{sign}(x),&\mathrm R.
\end{cases}
\] There is therefore no change in the definition of the R and NS
sectors. The two symbols called \(\eta\) in the two formulas encode
different quantities: the first is a relative endpoint sign on the
strip, while the second already includes the half-differential Jacobian
on the upper half-plane.

\subsubsection{Mapping the strip to the upper
half-plane}\label{mapping-the-strip-to-the-upper-half-plane}

Use Euclidean strip coordinates \[
w=\tau+i\sigma,
\qquad
0\leq\sigma\leq l,
\qquad
-\infty<\tau<\infty.
\] The exponential map \[
z=e^{\pi w/l}
=e^{\pi\tau/l}e^{i\pi\sigma/l}
\] sends the strip to the upper half-plane because \[
0\leq\arg z\leq\pi.
\] Its action on the two boundaries is \[
\begin{aligned}
\sigma=0
&\quad\Longrightarrow\quad
z=e^{\pi\tau/l}>0,\\%1
\sigma=l
&\quad\Longrightarrow\quad
z=-e^{\pi\tau/l}<0.
\end{aligned}
\] Thus \[
\boxed{
\sigma=0\longleftrightarrow x>0,
\qquad
\sigma=l\longleftrightarrow x<0,
}
\] where \(x=\operatorname{Re}z\) on the boundary.

The book's Eq. (4.138) parametrizes the strip as \(-l\leq\sigma\leq0\)
and uses \(w=\tau-i\sigma\). Replacing that \(\sigma\) by \(-\sigma\)
gives exactly the convention above, which is adapted to the endpoints
\(0\) and \(l\) used in the question.

The inverse map and its derivatives are \[
w=\frac{l}{\pi}\log z,
\qquad
\frac{dw}{dz}=\frac{l}{\pi z},
\] and \[
\bar w=\frac{l}{\pi}\log\bar z,
\qquad
\frac{d\bar w}{d\bar z}=\frac{l}{\pi\bar z}.
\]

\subsubsection{Why a fermion produces an additional
sign}\label{why-a-fermion-produces-an-additional-sign}

A holomorphic primary field of weight \(h\) transforms under
\(w\mapsto z(w)\) as \[
\Phi(z)
=
\left(\frac{dw}{dz}\right)^h\Phi(w).
\] The chiral worldsheet fermions have weight \(h=1/2\). Therefore
define \[
\psi^\mu(z)
=
\left(\frac{dw}{dz}\right)^{1/2}\psi_+^\mu(w),
\] and \[
\bar\psi^\mu(\bar z)
=
\left(\frac{d\bar w}{d\bar z}\right)^{1/2}\psi_-^\mu(\bar w).
\] Suppose that on one edge of the strip the gluing condition is \[
\psi_+^\mu(w)
=s\,{D^\mu}_{\nu}\psi_-^\nu(\bar w),
\] where \(s\) is the sign assigned to that edge. Substitution into the
transformation laws gives \[
\begin{aligned}
\psi^\mu(z)
&=
\left(\frac{dw}{dz}\right)^{1/2}
s\,{D^\mu}_{\nu}\psi_-^\nu(\bar w)\\%1
&=
s
\left[
\frac{dw/dz}{d\bar w/d\bar z}
\right]^{1/2}
{D^\mu}_{\nu}\bar\psi^\nu(\bar z).
\end{aligned}
\] Using the derivatives of the logarithm, \[
\left[
\frac{dw/dz}{d\bar w/d\bar z}
\right]^{1/2}
=
\left(\frac{\bar z}{z}\right)^{1/2}.
\] This factor is the entire origin of the extra sign.

\subsubsection{Evaluating the square root on the two halves of the real
axis}\label{evaluating-the-square-root-on-the-two-halves-of-the-real-axis}

The upper-half-plane limit fixes the branches: \[
\arg z\in(0,\pi),
\qquad
\arg\bar z\in(-\pi,0).
\]

For \(x>0\), one approaches the positive real axis with \[
\arg z\longrightarrow0,
\qquad
\arg\bar z\longrightarrow0.
\] Therefore \[
\left(\frac{\bar z}{z}\right)^{1/2}=+1,
\qquad x>0.
\]

For \(x<0\), the two boundary values approach the negative real axis
from opposite sides: \[
z=|x|e^{i\pi},
\qquad
\bar z=|x|e^{-i\pi}.
\] Consequently, \[
\left(\frac{\bar z}{z}\right)^{1/2}
=
\left(e^{-2\pi i}\right)^{1/2}
=e^{-i\pi}
=-1,
\qquad x<0.
\] Equivalently, the two square-root Jacobians are \[
\left(\frac{dw}{dz}\right)^{1/2}
=
\sqrt{\frac{l}{\pi|x|}}
\begin{cases}
1,&x>0,\\
-i,&x<0,
\end{cases}
\] and \[
\left(\frac{d\bar w}{d\bar z}\right)^{1/2}
=
\sqrt{\frac{l}{\pi|x|}}
\begin{cases}
1,&x>0,\\
+i,&x<0.
\end{cases}
\] Their ratio is \(+1\) for \(x>0\) and \(-1\) for \(x<0\).

It is useful to summarize this Jacobian as \[
J_{1/2}(x)
:=
\left(\frac{\bar z}{z}\right)^{1/2}_{\operatorname{Im}z\to0^+}
=
\operatorname{sign}(x).
\] This formula depends on the continuous branch inherited from the
upper half-plane. A simultaneous overall sign change of both square
roots is conventional, but the relative sign between the two halves of
the real axis is fixed.

\subsubsection{Applying the two original endpoint
conditions}\label{applying-the-two-original-endpoint-conditions}

At \(\sigma=0\), which maps to \(x>0\), the original condition is \[
\psi_+^\mu(0)
={D_0^\mu}_{\nu}\psi_-^\nu(0).
\] Here the strip sign is \(s=+1\) and the Jacobian is
\(J_{1/2}(x)=+1\). Therefore \[
\psi^\mu(x)
={D_0^\mu}_{\nu}\bar\psi^\nu(x),
\qquad x>0.
\]

At \(\sigma=l\), which maps to \(x<0\), the original condition is \[
\psi_+^\mu(l)
=\eta_{\mathrm{strip}}{D_l^\mu}_{\nu}\psi_-^\nu(l).
\] Here the strip sign is \(s=\eta_{\mathrm{strip}}\), while the
Jacobian is \(J_{1/2}(x)=-1\). Hence \[
\psi^\mu(x)
=-\eta_{\mathrm{strip}}{D_l^\mu}_{\nu}\bar\psi^\nu(x),
\qquad x<0.
\]

Combining the two halves gives the general upper-half-plane boundary
condition \[
\boxed{
\psi^\mu(x)
=
\begin{cases}
{D_0^\mu}_{\nu}\bar\psi^\nu(x),&x>0,\\[2mm]
-\eta_{\mathrm{strip}}{D_l^\mu}_{\nu}\bar\psi^\nu(x),&x<0.
\end{cases}
}
\]

\subsubsection{NS and R sectors for a single
D-brane}\label{ns-and-r-sectors-for-a-single-d-brane}

Equation (13.99) uses one matrix \(D\), so it assumes \[
D_0=D_l=D.
\] This is the case in which the two ends of the open string obey the
same NN/DD assignment, as for an open string ending on a single
D\(p\)-brane.

In the NS sector, \[
\eta_{\mathrm{strip}}=-1.
\] Therefore \[
-\eta_{\mathrm{strip}}=+1,
\] and the boundary condition is the same on both halves of the real
axis: \[
\psi^\mu(x)
={D^\mu}_{\nu}\bar\psi^\nu(x),
\qquad
x>0\ \text{or}\ x<0.
\] Thus \[
\eta_{\mathrm{UHP}}(x)=+1
\qquad
\text{in the NS sector}.
\]

In the R sector, \[
\eta_{\mathrm{strip}}=+1.
\] The positive real axis still has coefficient \(+1\), but the negative
real axis has coefficient \(-1\): \[
\psi^\mu(x)
=
\begin{cases}
{D^\mu}_{\nu}\bar\psi^\nu(x),&x>0,\\[1mm]
-{D^\mu}_{\nu}\bar\psi^\nu(x),&x<0.
\end{cases}
\] Since \[
\operatorname{sign}(x)
=
\begin{cases}
+1,&x>0,\\
-1,&x<0,
\end{cases}
\] this becomes \[
\psi^\mu(x)
=
\operatorname{sign}(x)
{D^\mu}_{\nu}\bar\psi^\nu(x)
\qquad
\text{in the R sector}.
\]

The two results can be written together as \[
\boxed{
\left.
\psi^\mu(z)
=
\eta_{\mathrm{UHP}}(\operatorname{Re}z)
{D^\mu}_{\nu}\bar\psi^\nu(\bar z)
\right|_{\operatorname{Im}z=0},
\qquad
\eta_{\mathrm{UHP}}(x)
=
\begin{cases}
+1,&\mathrm{NS},\\
\operatorname{sign}(x),&\mathrm R.
\end{cases}
}
\] This is Eq. (13.99).

\subsubsection{Why the R/NS signs appear to have exchanged
roles}\label{why-the-rns-signs-appear-to-have-exchanged-roles}

The strip sign and the upper-half-plane sign answer different questions.

On the strip, \(\eta_{\mathrm{strip}}\) compares the gluing condition at
\(\sigma=l\) with the one at \(\sigma=0\): \[
\eta_{\mathrm{strip}}
=
\begin{cases}
+1,&\text{same endpoint sign, R},\\
-1,&\text{opposite endpoint sign, NS}.
\end{cases}
\]

On the upper half-plane, the two endpoints have become two parts of one
real boundary, and the conformal transformation contributes \[
J_{1/2}(x)
=
\begin{cases}
+1,&x>0,\\
-1,&x<0.
\end{cases}
\] The negative-axis coefficient is therefore \[
\eta_{\mathrm{UHP}}(x<0)
=J_{1/2}(x<0)\eta_{\mathrm{strip}}
=-\eta_{\mathrm{strip}}.
\] This gives \[
\begin{array}{c|c|c}
\text{sector}
&\eta_{\mathrm{strip}}
&-\eta_{\mathrm{strip}}\ \text{on }x<0\\
\hline
\mathrm{NS}&-1&+1\\
\mathrm R&+1&-1
\end{array}
\] and explains the apparent reversal.

Geometrically, the exponential map converts translation along the
doubled strip into rotation around the origin of the \(z\)-plane. A
fermion is a spin-\(1/2\) field, so a full \(2\pi\) rotation contributes
a minus sign. The square-root Jacobian above is the local manifestation
of that spinorial sign. In the R sector the resulting real-axis gluing
changes sign at \(x=0\); this is consistent with the Ramond spin-field
branch point located at the origin, with the conjugate insertion
effectively at infinity.

\subsubsection{Different branes at the two
ends}\label{different-branes-at-the-two-ends}

The starting equation in the question allows \[
D_0\neq D_l.
\] For example, this occurs when the string stretches between different
D-branes or when the two endpoints have different NN/DD assignments. In
that case the correct upper-half-plane condition remains \[
\psi^\mu(x)
=
\begin{cases}
{D_0^\mu}_{\nu}\bar\psi^\nu(x),&x>0,\\[2mm]
-\eta_{\mathrm{strip}}{D_l^\mu}_{\nu}\bar\psi^\nu(x),&x<0.
\end{cases}
\] It cannot in general be reduced to a single matrix \(D\) multiplied
by either \(1\) or \(\operatorname{sign}(x)\). The origin then marks a
boundary-condition-changing point between the \(D_0\) and \(D_l\)
boundary conditions. Equation (13.99) is the special case relevant to
the single D\(p\)-brane discussion in §13.3.

\subsubsection{Checks}\label{checks-13}

\begin{enumerate}
\def\labelenumi{\arabic{enumi}.}
\item
  \textbf{NS sector:} On the negative axis the original endpoint sign is
  \(-1\) and the half-differential Jacobian is also \(-1\), so their
  product is \(+1\). The gluing is therefore \(D\) on the whole real
  axis.
\item
  \textbf{R sector:} On the negative axis the original endpoint sign is
  \(+1\), while the Jacobian is \(-1\). Hence the coefficient is \(-D\)
  for \(x<0\) and \(+D\) for \(x>0\), namely
  \(\operatorname{sign}(x)D\).
\item
  \textbf{Why this is special to fermions:} For an integer-weight field
  with \(h=1\), the analogous phase on the negative axis is \[
  \left(\frac{\bar z}{z}\right)^1
  =e^{-2\pi i}
  =+1.
  \] The additional minus sign arises specifically because \(h=1/2\).
\item
  \textbf{Same-brane assumption:} If \(D_0=D_l\), the book's compact
  formula follows. If \(D_0\neq D_l\), retaining both matrices is
  essential.
\item
  \textbf{Branch convention:} Changing the overall sign of a square root
  redefines one chiral fermion by an overall sign. It does not remove
  the relative sign between \(x>0\) and \(x<0\), which is the physical
  content of the derivation.
\end{enumerate}

\subsubsection{Result}\label{result-17}

The strip conditions are \[
\psi_+^\mu(0)
={D_0^\mu}_{\nu}\psi_-^\nu(0),
\qquad
\psi_+^\mu(l)
=\eta_{\mathrm{strip}}{D_l^\mu}_{\nu}\psi_-^\nu(l),
\] with \[
\eta_{\mathrm{strip}}
=
\begin{cases}
+1,&\mathrm R,\\
-1,&\mathrm{NS}.
\end{cases}
\] Under \[
z=e^{\pi(\tau+i\sigma)/l},
\] the boundaries map as \[
\sigma=0\longleftrightarrow x>0,
\qquad
\sigma=l\longleftrightarrow x<0.
\] Because \(\psi\) has conformal weight \(1/2\), \[
J_{1/2}(x)
=
\left(\frac{\bar z}{z}\right)^{1/2}_{\operatorname{Im}z\to0^+}
=
\begin{cases}
+1,&x>0,\\
-1,&x<0.
\end{cases}
\] Therefore \[
\boxed{
\psi^\mu(x)
=
\begin{cases}
{D_0^\mu}_{\nu}\bar\psi^\nu(x),&x>0,\\[2mm]
-\eta_{\mathrm{strip}}{D_l^\mu}_{\nu}\bar\psi^\nu(x),&x<0.
\end{cases}
}
\] For \(D_0=D_l=D\) this becomes \[
\boxed{
\left.
\psi^\mu(z)
=
\eta_{\mathrm{UHP}}(\operatorname{Re}z)
{D^\mu}_{\nu}\bar\psi^\nu(\bar z)
\right|_{\operatorname{Im}z=0},
\qquad
\eta_{\mathrm{UHP}}(x)
=
\begin{cases}
+1,&\mathrm{NS},\\
\operatorname{sign}(x),&\mathrm R.
\end{cases}
}
\] The apparent change of the R/NS signs is exactly the extra factor
\(-1\) acquired by a half-differential on the negative real axis.

\section{Q4: Why a D-brane preserves the diagonal
supercharge}\label{q4-why-a-d-brane-preserves-the-diagonal-supercharge}

Why it follows from this discussion that \(Q_{L}\) and \(PQ_{R}\) must
have the same chirality and that the conserved supercharges are given by
eq.(13.104)?

\begin{quote}
\textbf{Textbook location and related material.} Chapter 13, Section
13.3 (pp.~411--415), especially Eqs. (13.100)--(13.104): boundary
conformal invariance, doubled currents, and the diagonal spacetime
supercharge preserved by a D-brane.
\end{quote}

\subsection{Answer}\label{answer-18}

\subsubsection{Core point}\label{core-point-18}

There are two logically distinct statements behind Eq. (13.104):

\begin{enumerate}
\def\labelenumi{\arabic{enumi}.}
\item
  \textbf{Conservation:} a worldsheet boundary does not preserve the
  left- and right-moving supersymmetry currents independently. The
  boundary condition identifies them through the spinor gluing matrix
  \(P\). Therefore only their diagonal combination \[
  Q_\alpha
  =Q_{L\alpha}+(P Q_R)_\alpha
  \] has no current flux through the boundary and is conserved.
\item
  \textbf{Chirality:} after the GSO projection, \(Q_L\) and \(Q_R\) are
  ten-dimensional Majorana-Weyl supercharges. The two terms in the
  preserved diagonal combination must lie in the same Weyl
  representation. Therefore \(Q_L\) and \(P Q_R\) must have the same
  chirality. Since \(P\) flips chirality when the number \(9-p\) of
  Dirichlet directions is odd, this condition selects even \(p\) in type
  IIA and odd \(p\) in type IIB.
\end{enumerate}

The conservation statement is the supersymmetry-current version of the
general BCFT construction in Eqs. (4.146) and (4.147). The chirality
statement is an additional compatibility condition with the type-II GSO
projection.

In that BCFT construction, \(\bar T(\bar z)\) means the antiholomorphic
component \(T_{\bar z\bar z}\), not automatically the Hermitian adjoint
of \(T(z)=T_{zz}\). The boundary condition \(T=\bar T\) sets the normal
energy flux to zero, while the restriction
\(\varepsilon=\bar\varepsilon\) links the two conformal-transformation
parameters. The result is one \textbf{diagonal Virasoro algebra}, with
one conserved generator for every mode \(n\), rather than one isolated
numerical charge. For time translations its conserved density is
proportional to \(T+\bar T\).

\subsubsection{The two supercharges before introducing a
boundary}\label{the-two-supercharges-before-introducing-a-boundary}

For a closed type-II string without a boundary, the holomorphic and
antiholomorphic sectors are independent. The canonical \(-1/2\)-picture
supersymmetry currents are, schematically, \[
\mathcal J_R^\alpha(z)
=S_R^\alpha(z)e^{-\phi_R(z)/2},
\] and \[
\mathcal J_L^\alpha(\bar z)
=S_L^\alpha(\bar z)e^{-\phi_L(\bar z)/2}.
\] The labels \(R\) and \(L\) follow the book's convention that the
holomorphic sector is called right-moving and the antiholomorphic sector
left-moving.

Their contour integrals define independent spacetime supercharges: \[
Q_R^\alpha
=\oint\frac{dz}{2\pi i}\,
\mathcal J_R^\alpha(z),
\qquad
Q_L^\alpha
=\oint\frac{d\bar z}{2\pi i}\,
\mathcal J_L^\alpha(\bar z).
\] Thus, before a boundary is introduced, type II theory has two
independently conserved sets of supercharges: \[
Q_L
\qquad\text{and}\qquad
Q_R.
\] After the GSO projection, each is Majorana-Weyl. In type IIB they
have the same chirality, while in type IIA they have opposite
chiralities.

\subsubsection{\texorpdfstring{Does \(\bar T(\bar z)\) mean the complex
conjugate of
\(T(z)\)?}{Does \textbackslash bar T(\textbackslash bar z) mean the complex conjugate of T(z)?}}\label{does-bar-tbar-z-mean-the-complex-conjugate-of-tz}

The short answer is: \textbf{the bar primarily labels the
antiholomorphic component, not Hermitian conjugation}.

Start with the Euclidean stress tensor \(T_{ab}\) and introduce \[
z=\tau+i\sigma,
\qquad
\bar z=\tau-i\sigma.
\] Its two chiral components are denoted \[
T(z)=T_{zz}(z,\bar z),
\qquad
\bar T(\bar z)=T_{\bar z\bar z}(z,\bar z).
\] Conservation and tracelessness imply, away from operator insertions,
\[
\partial_{\bar z}T(z)=0,
\qquad
\partial_z\bar T(\bar z)=0.
\] Thus \(T\) is holomorphic and \(\bar T\) is antiholomorphic. They
generate the two independent Virasoro algebras of a bulk CFT.

If one starts from a real classical Euclidean tensor \(T_{ab}\), then
complex conjugation of tensor components exchanges \(z\) and \(\bar z\),
so on a real field configuration one has a reality relation of the
schematic form \[
\bigl(T_{zz}(z,\bar z)\bigr)^*
=T_{\bar z\bar z}(\bar z,z).
\] This is a reality property inherited from the underlying real theory.
It is not the definition of the notation \(\bar T\), and it does not
identify the two chiral operator algebras.

In radial quantization the distinction is particularly clear. Hermitian
conjugation includes inversion of the radial coordinate: \[
T(z)^\dagger
=\bar z^{-4}T\left(\frac{1}{\bar z}\right),
\] and similarly \[
\bar T(\bar z)^\dagger
=z^{-4}\bar T\left(\frac{1}{z}\right).
\] Thus Hermitian conjugation maps each chiral stress tensor to itself
at the radially reflected point; it does not simply replace \(T\) by
\(\bar T\) at the same point. A parity transformation may exchange the
two sectors, but parity is a separate operation.

Therefore the boundary equation \[
T(x)=\bar T(x),
\qquad x\in\mathbb R,
\] is an additional physical gluing condition. It is not a tautology
following from the bar notation.

\subsubsection{\texorpdfstring{Why \(T(x)=\bar T(x)\) is the no-flux
condition}{Why T(x)=\textbackslash bar T(x) is the no-flux condition}}\label{why-txbar-tx-is-the-no-flux-condition}

The relation to energy flow can be seen directly in real coordinates.
With the above choice of \(z\) and \(\bar z\), \[
T_{zz}
=\frac14
\left(
T_{\tau\tau}
-2iT_{\tau\sigma}
-T_{\sigma\sigma}
\right),
\] and \[
T_{\bar z\bar z}
=\frac14
\left(
T_{\tau\tau}
+2iT_{\tau\sigma}
-T_{\sigma\sigma}
\right).
\] Conformal invariance gives tracelessness, \[
T_{\tau\tau}+T_{\sigma\sigma}=0.
\] Consequently, \[
T_{zz}
=\frac12
\left(
T_{\tau\tau}-iT_{\tau\sigma}
\right),
\qquad
T_{\bar z\bar z}
=\frac12
\left(
T_{\tau\tau}+iT_{\tau\sigma}
\right).
\] Hence, up to the overall normalization used when the CFT fields
\(T(z)\) and \(\bar T(\bar z)\) are defined, \[
T-\bar T
=-iT_{\tau\sigma},
\] while \[
T+\bar T
=T_{\tau\tau}.
\]

For a boundary at fixed \(\sigma\), the component \(T_{\tau\sigma}\) is
the flux of \(\tau\)-momentum, namely worldsheet energy, through the
boundary. Therefore \[
T_{\tau\sigma}\big|_{\partial\Sigma}=0
\qquad\Longleftrightarrow\qquad
T(x)=\bar T(x).
\] The \textbf{difference} \(T-\bar T\) measures normal flux, whereas
the \textbf{sum} \(T+\bar T\) is the tangential energy density. This is
why the conserved time-translation charge contains the sum even though
the boundary condition sets the difference to zero.

\subsubsection{Why the boundary leaves one Ward
identity}\label{why-the-boundary-leaves-one-ward-identity}

On the full complex plane, infinitesimal conformal transformations are
\[
\delta z=\varepsilon(z),
\qquad
\delta\bar z=\bar\varepsilon(\bar z).
\] There is no boundary, so \(\varepsilon\) and \(\bar\varepsilon\) can
be chosen independently. Their Ward identities are generated separately
by \(T\) and \(\bar T\), producing two commuting Virasoro algebras: \[
\mathrm{Vir}_L\oplus\mathrm{Vir}_R.
\]

On the upper half-plane, the real axis must be mapped to itself. At a
boundary point \(z=\bar z=x\), the normal displacement is \[
\delta\sigma
=\frac{\delta z-\delta\bar z}{2i}
=\frac{\varepsilon(x)-\bar\varepsilon(x)}{2i}.
\] Preserving the boundary therefore requires \[
\boxed{
\varepsilon(x)=\bar\varepsilon(x),
\qquad x\in\mathbb R.
}
\] The holomorphic and antiholomorphic transformations are no longer
independent; they are the two boundary values of one real-analytic
transformation.

This is what the book means when it says that there is only one Ward
identity. Equation (4.141) still contains a \(T\) term and a \(\bar T\)
term, \[
\begin{aligned}
\langle\delta X\rangle
={}&
-\oint_C\frac{dz}{2\pi i}\,
\varepsilon(z)\langle T(z)X\rangle\\%1
&+
\oint_C\frac{d\bar z}{2\pi i}\,
\bar\varepsilon(\bar z)
\langle\bar T(\bar z)X\rangle,
\end{aligned}
\] but the two terms have the \textbf{same} transformation parameter at
the boundary. They therefore describe one diagonal transformation rather
than two independent transformations.

``One Ward identity'' does not mean one numerical conserved charge.
Expanding the one allowed function in modes gives one infinite family of
conserved generators: \[
\{L_n^{\mathrm{bdy}}\}_{n\in\mathbb Z}.
\] Thus the boundary reduces \[
\mathrm{Vir}_L\oplus\mathrm{Vir}_R
\quad\longrightarrow\quad
\mathrm{Vir}_{\mathrm{diag}},
\] one diagonal Virasoro algebra.

\subsubsection{The doubling trick and the diagonal Virasoro
charge}\label{the-doubling-trick-and-the-diagonal-virasoro-charge}

The book implements the preceding statement by defining a single
holomorphic stress tensor on the doubled plane: \[
\mathcal T_{\mathrm{dbl}}(z)
=
\begin{cases}
T(z),&\operatorname{Im}z>0,\\[1mm]
\bar T(z),&\operatorname{Im}z<0.
\end{cases}
\] In the lower half-plane, the argument \(z\) on the second line is the
reflected holomorphic coordinate obtained from the upper-half-plane
coordinate \(\bar z\). The boundary condition \[
T(x)=\bar T(x)
\] makes \(\mathcal T_{\mathrm{dbl}}\) continuous across the real axis.

Reflect the antiholomorphic contour in Eq. (4.141) into the lower
half-plane. Reflection reverses its orientation, which cancels the
relative sign between the two terms in the Ward identity. The two
semicircles then form one closed contour: \[
L_n^{\mathrm{bdy}}
=
\oint_C\frac{dz}{2\pi i}\,
z^{n+1}\mathcal T_{\mathrm{dbl}}(z).
\] Because this is a closed contour integral of a holomorphic current,
it can be deformed without changing its value, provided it does not
cross an insertion. This is the contour meaning of conservation.

Written before doubling, the same generator contains both chiral
contributions: \[
\begin{aligned}
L_n^{\mathrm{bdy}}
={}&
\int_{C_+}\frac{dz}{2\pi i}\,
z^{n+1}T(z)\\%1
&-
\int_{C_+}\frac{d\bar z}{2\pi i}\,
\bar z^{\,n+1}\bar T(\bar z).
\end{aligned}
\] The minus sign is an orientation sign for the antiholomorphic
semicircle. When each term is expressed using the standard orientation
used to define its chiral charge, this is the diagonal \textbf{sum} of
the two charge contributions. Thus the phrase ``the conserved charge is
\(T+\bar T\)'' is shorthand: the local fields are integrated with their
appropriate contour orientations.

For ordinary worldsheet time translations, the shorthand becomes literal
at the level of the charge density. In Lorentzian strip coordinates,
write the two chiral energy densities as \(T_{++}\) and \(T_{--}\). The
total Hamiltonian is \[
H
=
\int_0^l d\sigma\,
\left(
T_{++}+T_{--}
\right).
\] Using chirality, \[
\partial_\tau T_{++}
=\partial_\sigma T_{++},
\qquad
\partial_\tau T_{--}
=-\partial_\sigma T_{--},
\] one finds \[
\frac{dH}{d\tau}
=
\left[
T_{++}-T_{--}
\right]_{\sigma=0}^{\sigma=l}.
\] The no-flux gluing condition \(T_{++}=T_{--}\) at each endpoint gives
\[
\frac{dH}{d\tau}=0.
\] The left- and right-moving energies are not separately conserved: an
incoming left-moving excitation reflects from the boundary as a
right-moving excitation. Their sum is conserved because the boundary
only transfers energy between the two chiral sectors; it does not absorb
it.

When the same boundary is described in the closed-string channel by a
boundary state, the gluing condition is often written \[
\left(
L_n-\bar L_{-n}
\right)|B\rangle=0.
\] The minus sign and the index reversal arise from the orientation
change in passing to the closed-string channel. This is the same
diagonal symmetry, not a contradiction with saying that the physical
open-channel charge contains the sum of left- and right-moving
contributions.

\subsubsection{General BCFT reason that only a diagonal charge
survives}\label{general-bcft-reason-that-only-a-diagonal-charge-survives}

The stress-tensor discussion generalizes directly. Suppose \(W(z)\) is a
holomorphic current of weight \(h\) and \(\bar W(\bar z)\) is an
antiholomorphic current with the boundary gluing \[
\left.
W(z)
=\bar W(\bar z)
\right|_{\operatorname{Im}z=0}.
\] Equation (4.147) then combines the two semicircle integrals into a
single closed contour integral: \[
\begin{aligned}
W_n
&=
\int_{C_+}\frac{dz}{2\pi i}\,
z^{n+h-1}W(z)
-
\int_{C_+}\frac{d\bar z}{2\pi i}\,
\bar z^{\,n+h-1}\bar W(\bar z)\\%1
&=
\oint_C\frac{dz}{2\pi i}\,
z^{n+h-1}W_{\mathrm{dbl}}(z).
\end{aligned}
\] The boundary condition makes the doubled current continuous, so the
closed contour defines a conserved diagonal charge.

If the two currents fail to match on the real axis, the two semicircles
cannot be joined smoothly. The difference between the charge evaluated
at two radii is then a boundary integral of the mismatch. Physically,
the current has a nonzero normal component, so charge flows into the
boundary and that symmetry is broken.

\subsubsection{Applying the BCFT argument to the supersymmetry
currents}\label{applying-the-bcft-argument-to-the-supersymmetry-currents}

The D-brane boundary condition for the spin fields is Eq. (13.100): \[
\left.
S_\alpha(z)
=P_\alpha{}^\beta\bar S_\beta(\bar z)
\right|_{\operatorname{Im}z=0}.
\] The book's footnote 13 adds an important fact: the superconformal
ghost is continuous across the real axis. Therefore the full
supersymmetry currents, not only the matter spin fields, are glued by
the same spinor matrix \(P\).

In the spinor frame used for Eq. (13.104), this relation is \[
\left.
\mathcal J_{L\alpha}(\bar z)
=(P\mathcal J_R)_\alpha(z)
\right|_{\operatorname{Im}z=0}.
\] Whether one writes this equation with \(P\), \(P^{-1}\), or the
charge-conjugate transpose of \(P\) depends on which spinor index is
transported from the right-moving to the left-moving frame. Equation
(13.104) fixes the book's convention: \((P Q_R)_\alpha\) is the
right-moving charge expressed with a left-moving spinor index. This
convention choice does not affect the chirality or conservation
argument.

Now substitute \[
W=P\mathcal J_R,
\qquad
\bar W=\mathcal J_L
\] into the BCFT construction. The two semicircle charges combine as \[
\boxed{
Q_\alpha
=Q_{L\alpha}+(P Q_R)_\alpha.
}
\] This is Eq. (13.104).

The orthogonal combination \[
Q_{L\alpha}-(P Q_R)_\alpha
\] has the wrong sign at the boundary. Its two semicircle currents do
not join into one analytic doubled current, so that combination is not
conserved. A D-brane therefore reduces the two independent type-II
supersymmetries to a diagonal half.

\subsubsection{Why the two terms must have the same
chirality}\label{why-the-two-terms-must-have-the-same-chirality}

Before the GSO projection, the spin fields in this part of §13.3 are
organized as 32-component Majorana spinors. The GSO projection retains
one Weyl chirality in each chiral sector. Write \[
\Gamma Q_L=\chi_L Q_L,
\qquad
\Gamma Q_R=\chi_R Q_R,
\qquad
\chi_L,\chi_R=\pm1,
\] where \[
\Gamma=\Gamma^0\Gamma^1\cdots\Gamma^9
\] is the ten-dimensional chirality matrix.

The sum \[
Q_L+P Q_R
\] is a single GSO-allowed Majorana-Weyl supercharge only if its two
terms have the same chirality. Equivalently, the boundary relation
between supersymmetry parameters must map the GSO-allowed right-moving
Weyl space into the GSO-allowed left-moving Weyl space. If the
chiralities disagree, the equation relating the two parameters has only
the trivial solution, so no supersymmetry of that form is preserved.

This is the precise meaning of the book's statement that \(Q_L\) and
\(P Q_R\) must have the same chirality. It is not merely that one cannot
algebraically add two Dirac spinors of opposite chirality; one can
formally do so. Rather, the preserved generator must belong to the Weyl
subspace selected by the GSO projection and must be parametrized by one
nonzero Majorana-Weyl supersymmetry parameter.

\subsubsection{\texorpdfstring{How \(P\) acts on
chirality}{How P acts on chirality}}\label{how-p-acts-on-chirality}

The vector gluing matrix for a D\(p\)-brane is \[
D^\mu{}_\nu
=
\begin{cases}
+\delta^\mu{}_\nu,&\mu\ \text{NN},\\
-\delta^\mu{}_\nu,&\mu\ \text{DD}.
\end{cases}
\] Its spinor representative \(P\) satisfies Eq. (13.102): \[
P\left(D^\mu{}_\nu\Gamma^\nu\right)
=\Gamma^\mu P.
\] The book solves this as \[
P
=
\prod_{i\in\mathrm{DD}}
\left(\Gamma_i\Gamma\right),
\] up to an overall sign and the phase needed to preserve the Majorana
condition.

There are \[
n_D=9-p
\] Dirichlet directions. Each factor \(\Gamma_i\Gamma\) anticommutes
with the chirality matrix: \[
\Gamma(\Gamma_i\Gamma)
=-(\Gamma_i\Gamma)\Gamma.
\] Passing \(\Gamma\) through all \(n_D\) factors therefore gives \[
\boxed{
\Gamma P
=(-1)^{n_D}P\Gamma
=(-1)^{9-p}P\Gamma.
}
\]

If \(Q_R\) has chirality \(\chi_R\), then \[
\begin{aligned}
\Gamma(P Q_R)
&=
(-1)^{9-p}P\Gamma Q_R\\%1
&=
(-1)^{9-p}\chi_R\,P Q_R.
\end{aligned}
\] Thus the chirality of \(P Q_R\) is \[
\chi_{P Q_R}
=(-1)^{9-p}\chi_R.
\] The condition that \(Q_L\) and \(P Q_R\) have the same chirality is
therefore \[
\boxed{
\chi_L
=(-1)^{9-p}\chi_R.
}
\]

\subsubsection{Type IIA and type IIB
D-branes}\label{type-iia-and-type-iib-d-branes}

In type IIB, \[
\chi_L=\chi_R.
\] The chirality condition becomes \[
(-1)^{9-p}=+1.
\] Hence \(9-p\) is even, so \[
\boxed{
p\ \text{is odd in type IIB}.
}
\]

In type IIA, \[
\chi_L=-\chi_R.
\] The chirality condition becomes \[
(-1)^{9-p}=-1.
\] Hence \(9-p\) is odd, so \[
\boxed{
p\ \text{is even in type IIA}.
}
\] This is the result stated immediately after Eq. (13.104).

\subsubsection{Relation to the supersymmetry
parameters}\label{relation-to-the-supersymmetry-parameters}

The book next writes a general type-II supersymmetry generator as \[
\bar\epsilon_L Q_L+\bar\epsilon_R Q_R.
\] Without a boundary, \(\epsilon_L\) and \(\epsilon_R\) are
independent, giving \(16+16\) real supercharges. A D-brane imposes \[
\boxed{
\epsilon_L
=\pm P\epsilon_R,
}
\] which is Eq. (13.106). The plus and minus signs correspond to a brane
and an anti-brane, with the assignment depending on the orientation
convention.

This equation leaves only one independent 16-component Majorana-Weyl
parameter. Using the invariance of the Majorana spinor pairing under the
spin lift \(P\), the generator can be written as the common parameter
contracted with \[
Q_L\pm P Q_R.
\] For the sign convention chosen in Eq. (13.104), the brane preserves
\[
Q_L+P Q_R.
\] The anti-brane preserves the opposite diagonal combination \[
Q_L-P Q_R.
\] Thus the current-gluing statement and the parameter constraint are
two equivalent descriptions of the same halving of supersymmetry.

On chiral spinors, the matrix \(P\) can equivalently be written, up to
phase and sign conventions, as the product of gamma matrices tangent to
the D\(p\)-brane: \[
\epsilon_L
=\pm\Gamma^0\Gamma^1\cdots\Gamma^p\epsilon_R,
\] which is Eq. (13.107).

\subsubsection{Checks}\label{checks-14}

\begin{enumerate}
\def\labelenumi{\arabic{enumi}.}
\item
  \textbf{D9-brane:} Here \(n_D=0\), so \(P\) preserves chirality. Type
  IIB has \(\chi_L=\chi_R\), and therefore \(Q_L+P Q_R\) is a valid
  chiral conserved charge. This is the D9-brane underlying type-I string
  theory.
\item
  \textbf{D0-brane:} Here \(n_D=9\), so \(P\) reverses chirality. Type
  IIA has opposite left- and right-moving chiralities, so \(P Q_R\) has
  the same chirality as \(Q_L\). A BPS D0-brane is therefore allowed in
  type IIA.
\item
  \textbf{D3-brane:} Here \(n_D=6\) is even, so \(P\) preserves
  chirality. It is compatible with the equal chiralities of type IIB, as
  expected.
\item
  \textbf{Counting:} The boundary relation
  \(\epsilon_L=\pm P\epsilon_R\) determines one 16-component parameter
  from the other. A single flat D\(p\)-brane therefore preserves 16 of
  the original 32 real type-II supercharges.
\item
  \textbf{Broken combination:} The current associated with \(Q_L-P Q_R\)
  has a boundary mismatch for the brane sign used in Eq. (13.104). It
  cannot be represented by the closed doubled contour of Eq. (4.147), so
  it is not conserved.
\end{enumerate}

\subsubsection{Result}\label{result-18}

The bar on the antiholomorphic stress tensor labels its chirality: \[
T(z)=T_{zz},
\qquad
\bar T(\bar z)=T_{\bar z\bar z},
\] and does not by itself mean Hermitian conjugation. At the boundary,
\[
T-\bar T\propto T_{\tau\sigma},
\qquad
T+\bar T\propto T_{\tau\tau}.
\] Thus \[
T(x)=\bar T(x)
\] sets the normal energy flux to zero, while \[
\varepsilon(x)=\bar\varepsilon(x)
\] leaves one boundary-preserving conformal parameter. The doubled
stress tensor has modes \[
L_n^{\mathrm{bdy}}
=
\oint\frac{dz}{2\pi i}\,
z^{n+1}\mathcal T_{\mathrm{dbl}}(z),
\] so the surviving symmetry is one diagonal Virasoro algebra, not one
single numerical charge.

The canonical supersymmetry currents are \[
\mathcal J_R=S_R e^{-\phi_R/2},
\qquad
\mathcal J_L=S_L e^{-\phi_L/2}.
\] The spin-field boundary condition and continuity of the superghost
identify them through \(P\): \[
\left.
\mathcal J_L
=P\mathcal J_R
\right|_{\operatorname{Im}z=0}.
\] Therefore the BCFT contour construction of Eq. (4.147) preserves the
diagonal charge \[
\boxed{
Q=Q_L+P Q_R,
}
\] while the orthogonal combination is broken.

For a D\(p\)-brane, \[
n_D=9-p,
\qquad
\Gamma P=(-1)^{9-p}P\Gamma.
\] Hence \[
\Gamma(P Q_R)
=(-1)^{9-p}\chi_R(P Q_R).
\] The two terms in the preserved charge have the same GSO chirality
precisely when \[
\boxed{
\chi_L=(-1)^{9-p}\chi_R.
}
\] Consequently, \[
\boxed{
\begin{aligned}
\text{type IIA:}&\qquad p\ \text{even},\\%1
\text{type IIB:}&\qquad p\ \text{odd}.
\end{aligned}
}
\] Equivalently, the unbroken supersymmetry parameters obey \[
\boxed{
\epsilon_L=\pm P\epsilon_R,
}
\] so one flat D\(p\)-brane preserves one diagonal set of 16 real
supercharges.

\section{Q5: Reading NS and R sectors from bosonized
weights}\label{q5-reading-ns-and-r-sectors-from-bosonized-weights}

Around eq. (11.102) there is the discussion of bosonization. The vertex
operators are \[
V_{\vec{\lambda}}(z)=: e^{i\vec{\lambda}\cdot \vec{X}(z)}:c_{\vec{\lambda}}
\] Where \(c_{\vec{\lambda}}\) is the cocycle factor and
\(\vec{\lambda}\) is a weight vector in the weight lattice of \(D_{n}\).
It is said that the \(\vec{\lambda}\) of all states in the NS sector are
either in the \((0)\) or \((V)\) conjugacy class of \(D_{n}\) while that
of R sector are in \((S)\) or \((C)\). In my understanding, NS/R sector
has vertex operators of corresponding periodicity on the complex plane:
NS is periodic and R is anti-periodic. How to see the periodicity of a
vertex operator with given \(\vec{\lambda}\) so we can derive the above
statement?

\begin{quote}
\textbf{Textbook location and related material.} Chapter 11, Section
11.4 (pp.~340--343) for bosonization and \(D_n\) conjugacy classes;
Chapter 13, Sections 13.1 and 13.4 (pp.~391--403, 415--422) for
bosonized covariant vertices and the covariant lattice.
\end{quote}

\subsection{Answer}\label{answer-19}

\subsubsection{Core point}\label{core-point-19}

Your statement about the fermions on the complex plane is correct: \[
\begin{array}{c|c|c}
\text{sector}
&\text{fermion on the cylinder}
&\text{fermion on the complex plane}\\
\hline
\mathrm{NS}&\text{anti-periodic}&\text{periodic}\\
\mathrm R&\text{periodic}&\text{anti-periodic}.
\end{array}
\]

The important refinement is that one should not determine the sector
from the periodicity of \(V_{\boldsymbol\lambda}(z)\) under rotating the
vertex operator itself. Instead, \(V_{\boldsymbol\lambda}(0)\) creates a
state at the origin, and one asks for the monodromy of a
\textbf{fundamental fermion} transported around that insertion.

Using \[
\Psi^{\pm j}(z)
=:\!e^{\pm iX^j(z)}\!:c_{\pm j},
\] one finds \[
\boxed{
\Psi^{\pm j}(z)V_{\boldsymbol\lambda}(0)
\sim
z^{\pm\lambda_j}
\bigl(\text{single-valued operator and cocycle factors}\bigr).
}
\] After taking \(z\) once counterclockwise around the origin, \[
z^{\pm\lambda_j}
\longrightarrow
e^{\pm2\pi i\lambda_j}z^{\pm\lambda_j}.
\] Therefore \[
\lambda_j\in\mathbb Z
\quad\Longrightarrow\quad
\text{fermion monodromy }+1
\quad\Longrightarrow\quad
\mathrm{NS},
\] while \[
\lambda_j\in\mathbb Z+\frac12
\quad\Longrightarrow\quad
\text{fermion monodromy }-1
\quad\Longrightarrow\quad
\mathrm R.
\]

The \(D_n\) weight lattice has exactly two conjugacy classes with
integer components, \[
(0)\cup(V),
\] and two with half-integer components, \[
(S)\cup(C).
\] This gives the statement around Eq. (11.102).

The monodromy distinguishes the two unions \[
(0)\cup(V)
\qquad\text{versus}\qquad
(S)\cup(C).
\] It does not by itself distinguish \((0)\) from \((V)\) or \((S)\)
from \((C)\). Those further distinctions encode fermion-number parity
and spinor chirality, respectively.

\subsubsection{Cylinder and plane
periodicities}\label{cylinder-and-plane-periodicities}

The book begins with the complex-plane mode expansion \[
\psi^i(z)
=
\sum_r b_r^i z^{-r-1/2},
\] where \[
r\in
\begin{cases}
\mathbb Z+\frac12,&\mathrm{NS},\\
\mathbb Z,&\mathrm R.
\end{cases}
\] Under one circuit around the origin, \[
z\longrightarrow e^{2\pi i}z,
\] each term transforms by \[
z^{-r-1/2}
\longrightarrow
e^{-2\pi i(r+1/2)}
z^{-r-1/2}.
\]

For the NS sector, \(r\in\mathbb Z+\tfrac12\), so
\(r+\tfrac12\in\mathbb Z\) and \[
e^{-2\pi i(r+1/2)}=+1.
\] Hence the fermion is periodic on the plane.

For the R sector, \(r\in\mathbb Z\), so \(r+\tfrac12\) is half-integral
and \[
e^{-2\pi i(r+1/2)}=-1.
\] Hence the fermion is anti-periodic on the plane.

This is opposite to the familiar cylinder boundary conditions because
the fermion is a half-differential. Under the cylinder-to-plane map, Eq.
(4.10) gives \[
\psi_{\mathrm{plane}}(z)
=
\left(\frac{l}{2\pi}\right)^{1/2}
z^{-1/2}\psi_{\mathrm{cyl}}(w).
\] A circuit around the origin changes \[
z^{-1/2}\longrightarrow-z^{-1/2}.
\] If \[
\psi_{\mathrm{cyl}}(w+il)
=s_{\mathrm{cyl}}\psi_{\mathrm{cyl}}(w),
\] then \[
s_{\mathrm{plane}}
=-s_{\mathrm{cyl}}.
\] Thus \[
\begin{aligned}
\mathrm{NS}:&\qquad
s_{\mathrm{cyl}}=-1
\quad\Longrightarrow\quad
s_{\mathrm{plane}}=+1,\\%1
\mathrm R:&\qquad
s_{\mathrm{cyl}}=+1
\quad\Longrightarrow\quad
s_{\mathrm{plane}}=-1.
\end{aligned}
\] This is precisely the Jacobian effect mentioned immediately before
Eq. (11.102).

\subsubsection{What periodicity should be
tested}\label{what-periodicity-should-be-tested}

Let \[
|\boldsymbol\lambda\rangle
=
V_{\boldsymbol\lambda}(0)|0\rangle.
\] The question ``is this an NS or an R state?'' means:

\begin{quote}
What boundary condition does the fundamental fermion obey when radially
quantized around the state \(|\boldsymbol\lambda\rangle\)?
\end{quote}

Equivalently, one studies \[
\Psi^{\pm j}(z)|\boldsymbol\lambda\rangle
=
\Psi^{\pm j}(z)V_{\boldsymbol\lambda}(0)|0\rangle
\] and analytically continues \(z\) around the origin.

This is different from asking how \(V_{\boldsymbol\lambda}(z)\) itself
transforms under a \(2\pi\) rotation. The latter is governed by the
conformal weight \[
h_{\boldsymbol\lambda}
=\frac{\boldsymbol\lambda^2}{2}
\] and therefore by its conformal-spin phase. That phase does not
classify the NS/R sector. For example, the identity operator and a
fermion operator both belong to the NS Hilbert space, although their
conformal weights are \(0\) and \(1/2\).

\subsubsection{Deriving the monodromy from the bosonized
OPE}\label{deriving-the-monodromy-from-the-bosonized-ope}

The chiral bosons obey \[
X^i(z)X^j(w)
\sim
-\delta^{ij}\ln(z-w).
\] For normal-ordered exponentials, Wick contraction gives \[
:\!e^{i\mathbf a\cdot\mathbf X(z)}\!:
:\!e^{i\mathbf b\cdot\mathbf X(w)}\!:
\sim
(z-w)^{\mathbf a\cdot\mathbf b}
:\!e^{i(\mathbf a+\mathbf b)\cdot\mathbf X(w)}\!:.
\] This is the book's Eq. (11.92).

The bosonized complex fermions carry weights \[
\mathbf a=\pm\mathbf e_j,
\] where \(\mathbf e_j\) is the \(j\)th Euclidean unit vector. Taking \[
\mathbf b=\boldsymbol\lambda
\] therefore gives \[
\begin{aligned}
\Psi^{+j}(z)V_{\boldsymbol\lambda}(0)
&\sim
z^{\lambda_j}
V_{\boldsymbol\lambda+\mathbf e_j}(0)
\times(\text{cocycles}),\\%1
\Psi^{-j}(z)V_{\boldsymbol\lambda}(0)
&\sim
z^{-\lambda_j}
V_{\boldsymbol\lambda-\mathbf e_j}(0)
\times(\text{cocycles}).
\end{aligned}
\]

Taking \(z\) once around the origin yields \[
\begin{aligned}
\Psi^{+j}(e^{2\pi i}z)V_{\boldsymbol\lambda}(0)
&=
e^{2\pi i\lambda_j}
\Psi^{+j}(z)V_{\boldsymbol\lambda}(0),\\%1
\Psi^{-j}(e^{2\pi i}z)V_{\boldsymbol\lambda}(0)
&=
e^{-2\pi i\lambda_j}
\Psi^{-j}(z)V_{\boldsymbol\lambda}(0),
\end{aligned}
\] where only the branch monodromy is displayed.

If all \(\lambda_j\) are integers, both phases equal \(+1\) for every
\(j\). The fundamental fermions are single-valued around the insertion,
so \(V_{\boldsymbol\lambda}\) creates an NS state on the plane.

If all \(\lambda_j\) are half-integers, both phases equal \(-1\) for
every \(j\). Every fundamental fermion changes sign around the
insertion, so \(V_{\boldsymbol\lambda}\) creates an R state on the
plane.

The book assumes that all \(2n\) real fermions have the same spin
structure. If some \(\lambda_j\) were integral and others half-integral,
different complex fermions would have different boundary conditions.
Such mixed twist sectors can be considered in more general models, but
they are not the common NS/R sectors assumed in §11.4.

\subsubsection{The same conclusion from the fermion mode
indices}\label{the-same-conclusion-from-the-fermion-mode-indices}

The OPE also determines which fermion modes act on the state. In radial
quantization, \[
\Psi^{\pm j}(z)
=
\sum_r\Psi_r^{\pm j}z^{-r-1/2}.
\] Around \(V_{\boldsymbol\lambda}(0)\), the OPE powers are shifted by
\(\pm\lambda_j\): \[
\Psi^{\pm j}(z)V_{\boldsymbol\lambda}(0)
\sim
\sum_{m\in\mathbb Z}
z^{m\pm\lambda_j}(\cdots).
\] Comparing \[
-r-\frac12
=m\pm\lambda_j
\] shows that:

\begin{itemize}
\tightlist
\item
  if \(\lambda_j\in\mathbb Z\), then \(r\in\mathbb Z+\tfrac12\), which
  is the NS moding;
\item
  if \(\lambda_j\in\mathbb Z+\tfrac12\), then \(r\in\mathbb Z\), which
  is the R moding.
\end{itemize}

Thus the branch-monodromy argument and the oscillator-moding argument
are exactly equivalent.

\subsubsection{\texorpdfstring{The four \(D_n\) conjugacy
classes}{The four D\_n conjugacy classes}}\label{the-four-d_n-conjugacy-classes}

The \(D_n\) root lattice is \[
(0)
=
\left\{
(k_1,\ldots,k_n)\in\mathbb Z^n
\ \middle|\
\sum_{i=1}^n k_i=0\pmod2
\right\}.
\] The vector conjugacy class is \[
(V)
=
\left\{
(k_1,\ldots,k_n)\in\mathbb Z^n
\ \middle|\
\sum_{i=1}^n k_i=1\pmod2
\right\}.
\] Therefore \[
\boxed{
(0)\cup(V)=\mathbb Z^n.
}
\] All components are integral, so all fermions have monodromy \(+1\).
These are precisely the NS weights.

The two spinor classes are \[
(S)
=
\left\{
\left(k_1+\frac12,\ldots,k_n+\frac12\right)
\ \middle|\
k_i\in\mathbb Z,\ 
\sum_i k_i=0\pmod2
\right\},
\] and \[
(C)
=
\left\{
\left(k_1+\frac12,\ldots,k_n+\frac12\right)
\ \middle|\
k_i\in\mathbb Z,\ 
\sum_i k_i=1\pmod2
\right\}.
\] Therefore \[
\boxed{
(S)\cup(C)
=
\left(\mathbb Z+\frac12\right)^n.
}
\] All components are half-integral, so all fermions have monodromy
\(-1\). These are precisely the R weights.

Combining the two results, \[
\boxed{
\begin{aligned}
\mathrm{NS}
&\quad\Longleftrightarrow\quad
\boldsymbol\lambda\in(0)\cup(V),\\%1
\mathrm R
&\quad\Longleftrightarrow\quad
\boldsymbol\lambda\in(S)\cup(C).
\end{aligned}
}
\]

\subsubsection{\texorpdfstring{Why both \((0)\) and \((V)\) occur in the
NS
sector}{Why both (0) and (V) occur in the NS sector}}\label{why-both-0-and-v-occur-in-the-ns-sector}

The NS vacuum is a scalar of \(SO(2n)\) and corresponds to \[
\boldsymbol\lambda=\mathbf0\in(0).
\] Each fermionic creation operator transforms in the vector
representation and carries a weight \[
\pm\mathbf e_j\in(V).
\] Consequently:

\begin{itemize}
\tightlist
\item
  an even number of fermionic excitations has total weight in \((0)\);
\item
  an odd number of fermionic excitations has total weight in \((V)\).
\end{itemize}

This is also immediate from the conjugacy-class addition rule \[
(V)+(V)=(0).
\] For example, a bilinear current \[
\Psi^{\pm i}\Psi^{\pm j}
\] has weight \[
\pm\mathbf e_i\pm\mathbf e_j,
\] which is a \(D_n\) root and hence belongs to \((0)\). A single
fermion belongs to \((V)\). Both operators are local with respect to the
NS vacuum and both belong to the NS Hilbert space.

\subsubsection{\texorpdfstring{Why \((S)\) and \((C)\) occur in the R
sector}{Why (S) and (C) occur in the R sector}}\label{why-s-and-c-occur-in-the-r-sector}

The Ramond ground state is created by a spin field \[
S_{\mathbf s}(0)
=
:\!\exp\left[
\frac{i}{2}
\sum_{j=1}^n s_jX^j(0)
\right]\!:c_{\mathbf s},
\qquad
s_j=\pm1.
\] Its weight vector is \[
\boldsymbol\lambda_{\mathbf s}
=
\frac12(s_1,\ldots,s_n).
\] The OPE with a fermion is \[
\Psi^{\pm j}(z)S_{\mathbf s}(0)
\sim
z^{\pm s_j/2}(\cdots).
\] Every exponent is half-integral, so a circuit around the origin gives
\[
e^{\pm\pi i s_j}=-1.
\] Thus the spin field creates the branch cut required for an R state on
the complex plane.

The spinor weights with an even number of minus signs form \((S)\),
while those with an odd number form \((C)\). Acting with a fermion adds
a vector weight: \[
(V)+(S)=(C),
\qquad
(V)+(C)=(S).
\] Hence fermionic excitations interchange the two spinor classes, but
the components remain half-integral and the state remains in the R
sector.

Before the GSO projection, the R Hilbert space contains the two spinor
chiralities. The GSO projection selects one of \((S)\) or \((C)\)
according to the chosen convention. Similarly, in applications to the
superstring, the NS GSO projection selects a definite fermion-number
parity. The NS/R classification itself precedes this further projection.

The spin-field conformal weight is \[
h_S
=\frac{\boldsymbol\lambda_{\mathbf s}^2}{2}
=\frac12\left(\frac n4\right)
=\frac n8,
\] in agreement with the book.

\subsubsection{Role of the cocycle
factors}\label{role-of-the-cocycle-factors}

The cocycles \(c_{\boldsymbol\lambda}\) are essential because
exponentials built from different bosons would otherwise commute,
whereas different fermions must anticommute. They supply
momentum-dependent, position-independent phases that correct the
exchange algebra.

The branch power in the OPE is still \[
(z-w)^{\boldsymbol\lambda\cdot\boldsymbol\mu}.
\] Therefore the full monodromy obtained by taking one operator
completely around another contains \[
e^{2\pi i\boldsymbol\lambda\cdot\boldsymbol\mu}.
\] For a fundamental fermion, \(\boldsymbol\mu=\pm\mathbf e_j\), so this
becomes \[
e^{\pm2\pi i\lambda_j}.
\] The cocycle fixes exchange signs and operator ordering, but it does
not change whether \(\lambda_j\) is integral or half-integral and hence
does not change the NS/R assignment.

This also explains the book's locality statement. Within \((0)\cup(V)\),
the relevant mutual inner products are integral, so the NS operator
algebra contains no branch cuts. A vector weight has half-integral inner
product with a spinor weight modulo integers, \[
(V)\cdot(S)
=(V)\cdot(C)
=\frac12\pmod1,
\] so transporting a fermion around a Ramond spin field produces the
expected minus sign.

\subsubsection{Checks}\label{checks-15}

\begin{enumerate}
\def\labelenumi{\arabic{enumi}.}
\item
  \textbf{NS vacuum:} For \(\boldsymbol\lambda=\mathbf0\), \[
  \Psi^{\pm j}(z)V_{\mathbf0}(0)
  \sim z^0(\cdots),
  \] so the fermions are periodic on the plane.
\item
  \textbf{Single NS fermion:} For
  \(\boldsymbol\lambda=\mathbf e_k\in(V)\), the powers
  \(\pm\lambda_j=\pm\delta_{jk}\) are integers. The fermion remains
  single-valued around the insertion even though \(V_{\mathbf e_k}\)
  itself has \(h=1/2\).
\item
  \textbf{NS current:} For
  \(\boldsymbol\lambda=\mathbf e_i+\mathbf e_j\in(0)\), all components
  are integers and the operator is local with respect to the fermions.
\item
  \textbf{Ramond spin field:} For
  \(\boldsymbol\lambda=\tfrac12(s_1,\ldots,s_n)\), every fermion OPE
  contains a square-root power \(z^{\pm1/2}\) and therefore has
  monodromy \(-1\).
\item
  \textbf{What monodromy cannot distinguish:} Both \((0)\) and \((V)\)
  give monodromy \(+1\), while both \((S)\) and \((C)\) give monodromy
  \(-1\). The parity conditions separating the two classes within each
  pair contain representation-theoretic information beyond the NS/R spin
  structure.
\end{enumerate}

\subsubsection{Result}\label{result-19}

The bosonized fermions and a general weight vertex are \[
\Psi^{\pm j}(z)
=:\!e^{\pm iX^j(z)}\!:c_{\pm j},
\qquad
V_{\boldsymbol\lambda}(0)
=:\!e^{i\boldsymbol\lambda\cdot\mathbf X(0)}\!:
c_{\boldsymbol\lambda}.
\] Their OPE is \[
\boxed{
\Psi^{\pm j}(z)V_{\boldsymbol\lambda}(0)
\sim
z^{\pm\lambda_j}
V_{\boldsymbol\lambda\pm\mathbf e_j}(0)
\times(\text{cocycles}).
}
\] Therefore the fermion monodromy around the state created by
\(V_{\boldsymbol\lambda}\) is \[
\boxed{
M_j^\pm(\boldsymbol\lambda)
=e^{\pm2\pi i\lambda_j}.
}
\] It follows that \[
\begin{aligned}
\lambda_j\in\mathbb Z\ \forall j
&\quad\Longrightarrow\quad
M_j^\pm=+1
\quad\Longrightarrow\quad
\mathrm{NS\ on\ the\ plane},\\%1
\lambda_j\in\mathbb Z+\frac12\ \forall j
&\quad\Longrightarrow\quad
M_j^\pm=-1
\quad\Longrightarrow\quad
\mathrm{R\ on\ the\ plane}.
\end{aligned}
\] The \(D_n\) weight lattice decomposes as \[
\begin{aligned}
(0)\cup(V)
&=\mathbb Z^n,\\%1
(S)\cup(C)
&=\left(\mathbb Z+\frac12\right)^n.
\end{aligned}
\] Hence \[
\boxed{
\begin{aligned}
\mathrm{NS}
&\quad\Longleftrightarrow\quad
\boldsymbol\lambda\in(0)\cup(V),\\%1
\mathrm R
&\quad\Longleftrightarrow\quad
\boldsymbol\lambda\in(S)\cup(C).
\end{aligned}
}
\] The split \((0)\) versus \((V)\) records even versus odd
tensor/fermion-number parity, while \((S)\) versus \((C)\) records the
two spinor chiralities. The sector itself is determined by the monodromy
of the fundamental fermions around the insertion, not by the rotation
phase of \(V_{\boldsymbol\lambda}\) alone.

\section{Q6: Why the one-loop trace is not trivialized by the
physical-state
condition}\label{q6-why-the-one-loop-trace-is-not-trivialized-by-the-physical-state-condition}

In the discussion of one-loop string partition function, for example in
the open bosonic string case we will compute \[
\mathrm{Tr}_{\mathcal{H}_{\mathrm{op}}}q^{H_{\mathrm{op}}}=\mathrm{Tr}_{\mathcal{H}_{\mathrm{op}}}q^{L_{0}-c/24} \ , \ q=e^{2\pi i\tau}
\] Or in the closed bosonic string \[
\mathrm{Tr}_{\mathcal{H}_{\mathrm{cl}}}(q^{L_{0}-c/24}\bar{q}^{\bar{L}_{0}-\bar{c}/24})
\] But a physical state \(|\phi \rangle\) should satisfies
\((L_{0}-1)|\phi \rangle=0\). Does this mean that every physical state
in the Hilbert space will have \(L_{0}\) eigenvalue \(L_{0}=1\)? Then it
seems strange for the trace over \(q^{L_{0}-c/24}\).

In the fermionic string case, the holomorphic part of the one loop
partition function of the fermionic string in Hamiltonian description
has the form of eq. (13.114) \[
\chi(\tau)\sim \mathrm{Tr}[e^{2\pi i\tau(L_{0}-1)}(-1)^{F_{R}}]
\] What is the Hamiltonian and how to derive the trace over
\(q^{L_{0}-1}\) in this case?

\begin{quote}
\textbf{Textbook location and related material.} Chapter 6, Sections 6.3
and 6.5 for torus/cylinder traces, together with Chapter 13, Section
13.4 near Eq. (13.114). The relevant trace is over a gauge-fixed CFT or
transverse Fock space, not only BRST cohomology at fixed target
momentum.
\end{quote}

\subsection{Answer}\label{answer-20}

\subsubsection{Core point}\label{core-point-20}

The apparent contradiction comes from using the phrase ``string Hilbert
space'' for two related but different spaces.

The physical-state equation \[
(L_0^{\mathrm m}-1)|\phi\rangle=0
\] is an on-shell condition in the covariant bosonic string. It defines
external asymptotic states, or equivalently the BRST cohomology at fixed
on-shell target-space momentum. If one literally restricted a trace to
such states and used \(L_0^{\mathrm m}-1\) as the traced operator, then
this operator would indeed vanish.

Thus, in that precise sense, the answer to the first question is
\textbf{yes}: every on-shell covariant bosonic-string state has matter
eigenvalue \(L_0^{\mathrm m}=1\) in the open string, and
\(L_0^{\mathrm m}=\bar L_0^{\mathrm m}=1\) in the closed string. The
crucial point is that this is not the restriction imposed inside the
one-loop CFT trace.

That is not the trace appearing in the one-loop worldsheet partition
function. The latter has either of the following equivalent
descriptions:

\begin{enumerate}
\def\labelenumi{\arabic{enumi}.}
\tightlist
\item
  in covariant gauge, it is a trace over the gauge-fixed worldsheet CFT
  state space, including the appropriate ghost insertions or ghost
  cancellations;
\item
  in light-cone gauge, it is a trace over all transverse oscillator
  states and an integral over Euclidean loop momentum. The Virasoro
  constraints have already been solved, so they do not impose \(L_0=1\)
  as an additional restriction on the oscillator trace.
\end{enumerate}

The operator in the exponential is a \textbf{worldsheet cylinder
Hamiltonian}, not the target-space energy \(P^0\). For a chiral CFT it
is \[
\boxed{
H_{\mathrm{cyl}}=L_0-\frac{c}{24}.
}
\] For the 24 transverse bosons of the critical bosonic string,
\(c=24\), and hence \[
H_{\mathrm{cyl}}=L_0-1.
\]

Equation (13.114) uses the same worldsheet propagation principle. In the
covariant RNS vertex-operator description, a physical integrated chiral
vertex, including its superghost picture dressing but before multiplying
by \(c\), has weight one. Consequently its inverse propagator is
\(L_0-1\), and the holomorphic one-loop trace is schematically \[
\boxed{
\chi(\tau)\sim
\operatorname{Tr}
\left[
q_\tau^{\,L_0-1}(-1)^{F_R}
\right],
\qquad
q_\tau=e^{2\pi i\tau}.
}
\] Here \((-1)^{F_R}\) changes the ordinary trace into a spacetime
supertrace, with opposite signs for spacetime bosons and fermions.

\subsubsection{The three notions that must be
separated}\label{the-three-notions-that-must-be-separated}

It is useful to keep the following three operators or spaces distinct.

\begin{longtable}[]{@{}
  >{\raggedright\arraybackslash}p{(\linewidth - 6\tabcolsep) * \real{0.2244}}
  >{\raggedright\arraybackslash}p{(\linewidth - 6\tabcolsep) * \real{0.4098}}
  >{\raggedright\arraybackslash}p{(\linewidth - 6\tabcolsep) * \real{0.1512}}
  >{\raggedright\arraybackslash}p{(\linewidth - 6\tabcolsep) * \real{0.2146}}@{}}
\toprule\noalign{}
\begin{minipage}[b]{\linewidth}\raggedright
Object
\end{minipage} & \begin{minipage}[b]{\linewidth}\raggedright
Meaning
\end{minipage} & \begin{minipage}[b]{\linewidth}\raggedright
Condition or generator
\end{minipage} & \begin{minipage}[b]{\linewidth}\raggedright
\end{minipage} \\
\midrule\noalign{}
\endhead
\bottomrule\noalign{}
\endlastfoot
Plane CFT state & State created by a local operator at the origin &
\(L_0\) measures conformal weight & \\
On-shell string state & BRST cohomology class or old-covariant physical
state & \((L_0^{\mathrm m}-1)\) & \(\|\phi\rangle=0\) in the bosonic
open string \\
State propagating around a one-loop worldsheet & Gauge-fixed CFT state,
or a transverse light-cone state with Euclidean loop momentum &
\(H_{\mathrm{cyl}}=L_0-c/24\) & \\
\end{longtable}

The same oscillator state can appear in the second and third rows in
different ways. In the physical spectrum, its target-space momentum is
constrained to its mass shell. In a vacuum loop, the corresponding
particle species runs with arbitrary Euclidean loop momentum. Thus the
oscillator polarization is physical, while the momentum carried around
the loop is not constrained to its mass shell at every point in the
integral.

There is a second, equivalent way to state this. In covariant string
field theory, the kinetic operator is proportional to \[
L_0-1.
\] The propagator is schematically \[
\frac{b_0}{L_0-1}
=
b_0\int_0^\infty ds\,
e^{-s(L_0-1)}.
\] The on-shell condition \(L_0-1=0\) locates a pole of this propagator.
One must not impose the pole equation before constructing the
propagator; doing so would erase the propagation that the loop diagram
is meant to describe.

\subsubsection{\texorpdfstring{Deriving the cylinder Hamiltonian
\(L_0-c/24\)}{Deriving the cylinder Hamiltonian L\_0-c/24}}\label{deriving-the-cylinder-hamiltonian-l_0-c24}

Start with a chiral CFT on the complex plane and map it to a cylinder.
In dimensionless coordinates one may write \[
z=e^w.
\] The stress tensor is not a primary field. Under a holomorphic
coordinate transformation it obeys \[
T_{\mathrm{cyl}}(w)
=
\left(\frac{dz}{dw}\right)^2T_{\mathrm{plane}}(z)
+\frac{c}{12}\{z,w\},
\] where \[
\{z,w\}
=
\frac{z'''(w)}{z'(w)}
-\frac32
\left(\frac{z''(w)}{z'(w)}\right)^2
\] is the Schwarzian derivative. For \(z=e^w\), \[
\{e^w,w\}
=1-\frac32
=-\frac12.
\] Therefore \[
T_{\mathrm{cyl}}(w)
=z^2T_{\mathrm{plane}}(z)-\frac{c}{24}.
\]

The zero mode generating translations along the cylinder is consequently
\[
\boxed{
H_{\mathrm{cyl}}
=L_0-\frac{c}{24}.
}
\] The subtraction is the Casimir energy of the CFT vacuum on the
spatial circle. It is not, by itself, an extra physical-state
constraint.

One must not identify \(c/24\) with the string intercept in every
formulation. For example, the covariant bosonic matter CFT has \[
c_{\mathrm m}=26,
\] so its matter Casimir shift alone would be \(26/24\), not one. The
covariant calculation must also include the \(b,c\) ghosts and their
zero-mode insertions. After the longitudinal, timelike, and ghost
contributions cancel, the answer is equivalent to 24 transverse bosons,
for which \[
\frac{c_\perp}{24}
=
\frac{24}{24}
=1.
\] It is in this reduced light-cone description that the CFT Casimir
shift directly reproduces the bosonic intercept.

The RNS string provides an even clearer warning. The complete critical
matter-plus-ghost system has total central charge zero, but the
propagation operator before the \(c\)-ghost dressing is still written as
\(L_0-1\). Here the one is fixed by the weight-one physical-vertex
condition, while the familiar NS and R zero-point energies emerge only
after resolving the spin sector and canceling the unphysical fields.

For a non-chiral closed-string theory, the Euclidean-time and
spatial-translation generators are \[
\begin{aligned}
H_{\mathrm{cyl}}
&=
\left(L_0-\frac{c}{24}\right)
+
\left(\bar L_0-\frac{\bar c}{24}\right),\\%1
P_{\mathrm{cyl}}
&=
\left(L_0-\frac{c}{24}\right)
-
\left(\bar L_0-\frac{\bar c}{24}\right).
\end{aligned}
\] When \(c=\bar c\), the second equation reduces to \[
P_{\mathrm{cyl}}=L_0-\bar L_0.
\]

\subsubsection{Cutting and gluing the
torus}\label{cutting-and-gluing-the-torus}

Write the torus modulus as \[
\tau=\tau_1+i\tau_2.
\] Cutting the torus along a spatial cycle gives a cylinder of Euclidean
length proportional to \(\tau_2\). Gluing its two ends without a twist
gives an evolution operator \[
e^{-2\pi\tau_2H_{\mathrm{cyl}}}.
\] Before gluing, one may also shift one end relative to the other by
\(2\pi\tau_1\). This inserts \[
e^{2\pi i\tau_1P_{\mathrm{cyl}}}.
\] Taking the trace performs the gluing: \[
Z(\tau,\bar\tau)
=
\operatorname{Tr}_{\mathcal H_{\mathrm{CFT}}}
\left[
e^{-2\pi\tau_2H_{\mathrm{cyl}}}
e^{2\pi i\tau_1P_{\mathrm{cyl}}}
\right].
\] Substituting the expressions for \(H_{\mathrm{cyl}}\) and
\(P_{\mathrm{cyl}}\) gives \[
\boxed{
Z(\tau,\bar\tau)
=
\operatorname{Tr}_{\mathcal H_{\mathrm{CFT}}}
\left[
q_\tau^{\,L_0-c/24}
\bar q_\tau^{\,\bar L_0-\bar c/24}
\right],
}
\] where \[
q_\tau=e^{2\pi i\tau},
\qquad
\bar q_\tau=e^{-2\pi i\bar\tau}.
\] The bar on \(\bar q_\tau\) is important: with \(\tau_2>0\), both
factors decay exponentially.

For the open string, the spatial worldsheet is an interval. Making
Euclidean time periodic produces an annulus. There is one real modulus
\(t\), and no independent left-moving sector: \[
\boxed{
\mathcal A
=
\int_0^\infty\frac{dt}{2t}\,
\operatorname{Tr}_{\mathcal H_{\mathrm{op}}}
\left[
e^{-2\pi tH_{\mathrm{op}}}
\right],
\qquad
H_{\mathrm{op}}=L_0-\frac{c}{24}.
}
\] Equivalently, with \(\tau=it\), \[
q_\tau=e^{-2\pi t},
\qquad
\mathcal A
=
\int_0^\infty\frac{dt}{2t}\,
\operatorname{Tr}_{\mathcal H_{\mathrm{op}}}
q_\tau^{\,L_0-c/24}.
\] This is the book's Eq. (6.114).

\subsubsection{Why the bosonic physical-state condition does not make
the trace
constant}\label{why-the-bosonic-physical-state-condition-does-not-make-the-trace-constant}

For the open bosonic string, the matter zero mode is \[
L_0^{\mathrm m}
=
\alpha'p^2+N.
\] The old-covariant physical-state equation is \[
(L_0^{\mathrm m}-1)|\phi\rangle=0,
\] so, with mostly-plus target-space signature, \[
\alpha'p^2+N-1=0
\quad\Longleftrightarrow\quad
m_N^2=\frac{N-1}{\alpha'}.
\] This equation does not say that only \(N=1\) states exist. It says
that every oscillator level \(N\) has its own allowed target-space mass:
\[
\begin{array}{c|c|c}
N&\alpha'm_N^2&\text{state type}\\
\hline
0&-1&\text{tachyon}\\
1&0&\text{massless vector}\\
2&1&\text{first massive level}\\
\vdots&\vdots&\vdots
\end{array}
\] On shell, the momentum contribution changes with \(N\) so that the
complete matter eigenvalue is one.

In the annulus trace, however, the loop momentum is Wick-rotated and
integrated: \[
\operatorname{Tr}_{\mathcal H_{\mathrm{op}}}
q_\tau^{\,L_0-1}
=
\sum_N d_N
\int\frac{d^Dk_E}{(2\pi)^D}
e^{-2\pi t(\alpha'k_E^2+N-1)}.
\] Here \(d_N\) is the number of physical transverse polarizations at
level \(N\). The quantity \[
\alpha'k_E^2+N-1
=
\alpha'(k_E^2+m_N^2)
\] is not zero for a generic Euclidean loop momentum \(k_E\). It is the
inverse propagator in Schwinger proper-time form.

For the closed bosonic string, \[
L_0^{\mathrm m}
=
\frac{\alpha'p^2}{4}+N,
\qquad
\bar L_0^{\mathrm m}
=
\frac{\alpha'p^2}{4}+\bar N.
\] The two physical-state equations imply \[
N=\bar N,
\qquad
m_N^2=\frac{4}{\alpha'}(N-1).
\] Again, all levels occur. The \(\tau_1\) dependence of the torus trace
is \[
e^{2\pi i\tau_1(N-\bar N)}.
\] Integrating \(\tau_1\) over one period gives \[
\int_{-1/2}^{1/2}d\tau_1\,
e^{2\pi i\tau_1(N-\bar N)}
=
\delta_{N,\bar N},
\] so the modulus integral enforces level matching. The \(\tau_2\)
integral supplies the closed-string proper time.

\subsubsection{Proper-time derivation from the target-space one-loop
vacuum
energy}\label{proper-time-derivation-from-the-target-space-one-loop-vacuum-energy}

The same result follows without first using the cylinder map. A physical
string oscillator state \(I\) behaves in target space like a particle
species of mass \(m_I\). Its one-loop vacuum contribution is
schematically \[
\Gamma_{\mathrm{1-loop}}
\sim
\frac12
\sum_I(-1)^{F_I}
\int\frac{d^Dk_E}{(2\pi)^D}
\log(k_E^2+m_I^2).
\] Using \[
\log A
=
-\int_0^\infty\frac{ds}{s}\,e^{-sA}
+\text{constant},
\] one obtains \[
\Gamma_{\mathrm{1-loop}}
\sim
-\frac12
\int_0^\infty\frac{ds}{s}
\sum_I(-1)^{F_I}
\int\frac{d^Dk_E}{(2\pi)^D}
e^{-s(k_E^2+m_I^2)}.
\] For the open bosonic spectrum, \[
\alpha'(k_E^2+m_I^2)
=
\alpha'k_E^2+N_I-1
=
L_0-1
\] on the off-shell loop state. After setting \(s=2\pi\alpha't\), the
exponential becomes \[
e^{-2\pi t(L_0-1)}
=q_\tau^{\,L_0-1},
\qquad
q_\tau=e^{-2\pi t}.
\] Thus the annulus modulus is precisely Schwinger proper time. The
trace sums the physical particle species, while the momentum integral
supplies their off-shell propagation.

This is the logic used around Eqs. (6.89)--(6.93) of the book. It also
explains why imposing \(L_0-1=0\) inside the trace would be incorrect:
that condition identifies the poles of the propagators, whereas the
proper-time integral constructs the propagators.

\subsubsection{Covariant versus light-cone
traces}\label{covariant-versus-light-cone-traces}

There are two standard ways to obtain the same answer.

In light-cone gauge, only the \(D-2\) transverse oscillators are
present. For the critical bosonic string, \[
c_\perp=24,
\] and hence \[
H_{\mathrm{op}}
=
L_0^\perp-\frac{c_\perp}{24}
=
L_0^\perp-1.
\] All transverse oscillator Fock states are allowed in the trace. As
the book emphasizes after Eq. (6.82), the constraints leading to the
physical-state conditions have already been solved and do not impose a
further restriction on this trace.

In covariant gauge, one may include all 26 coordinate oscillators, but
one must also trace over the \(b,c\) ghosts with the correct zero-mode
insertions. The longitudinal and timelike coordinate contributions are
canceled by the ghosts. The result is again the contribution of 24
transverse oscillators: \[
Z_{\mathrm{covariant}}
=
Z_{26\,X}\,Z_{bc}
=
Z_{24\,X^\perp}.
\] This is why the covariant and light-cone calculations have the same
level-generating function.

One should therefore not interpret \[
\mathcal H_{\mathrm{op}}
\] in Eq. (6.114) as ``the set of covariant states after imposing
\(L_0=1\) at fixed momentum.'' In that equation it means the Hilbert
space appropriate to the chosen gauge-fixed formulation.

\subsubsection{The Hamiltonian in Eq.
(13.114)}\label{the-hamiltonian-in-eq.-13.114}

Equation (13.114) is a \textbf{holomorphic} or chiral partition
function. Its Hamiltonian is \[
\boxed{
H_{\mathrm{chiral}}=L_0-1.
}
\] More precisely:

\begin{itemize}
\tightlist
\item
  \(L_0\) is the plane Virasoro zero mode of the chiral RNS system used
  in the covariant vertex-operator construction;
\item
  the subtraction of one is the intercept appropriate to an integrated,
  picture-dressed chiral string vertex before insertion of the
  reparametrization ghost \(c\);
\item
  for \(\tau=it\), \(e^{2\pi i\tau(L_0-1)}=e^{-2\pi t(L_0-1)}\) is
  Euclidean worldsheet evolution;
\item
  for nonzero \(\tau_1\), the phase implements the twist used when
  gluing the cylinder into a torus;
\item
  \((-1)^{F_R}\) supplies the minus sign for right-moving spacetime
  fermions.
\end{itemize}

It is not the target-space Hamiltonian \(P^0\). It is the dimensionless
generator of translations along the chiral worldsheet cylinder,
including the intercept appropriate to the string propagation problem.

The reason the exponent is written uniformly as \(L_0-1\) in both the NS
and R sectors is that Chapter 13 includes the superghost picture
dressing. Before multiplication by \(c\), every physical chiral vertex
has total conformal weight \[
h_{\mathrm{matter+superghost}}=1.
\] Thus the uniform physical equation is \[
(L_0-1)|V\rangle=0.
\]

After removing the longitudinal, timelike, and ghost degrees of freedom,
this uniform covariant expression reduces to the familiar
sector-dependent light-cone Hamiltonians \[
\boxed{
\begin{aligned}
H_{\mathrm{NS}}^{\mathrm{lc}}
&=
\frac{\alpha'k^2}{4}
+N_{\mathrm{NS}}-\frac12,\\%1
H_{\mathrm R}^{\mathrm{lc}}
&=
\frac{\alpha'k^2}{4}
+N_{\mathrm R},
\end{aligned}
}
\] for one chiral half of a closed RNS string. The corresponding
intercepts are therefore \[
a_{\mathrm{NS}}=\frac12,
\qquad
a_{\mathrm R}=0.
\] There is no contradiction between these formulas and Eq. (13.114):
the two descriptions distribute the normal-ordering and ghost
contributions differently.

\subsubsection{Deriving the Chapter 13 lattice
trace}\label{deriving-the-chapter-13-lattice-trace}

Chapter 13 bosonizes five complex matter fermions and the
\(\beta,\gamma\) superghost system. A state is labeled by \[
\mathbf w=(\boldsymbol\lambda,q_{\mathrm{gh}})\in D_{5,1},
\] where \(\boldsymbol\lambda\) is a \(D_5\) weight and
\(q_{\mathrm{gh}}\) is the superghost picture charge. According to Eq.
(13.65), the contribution to \(L_0\) is \[
L_0
=
\frac12\boldsymbol\lambda^2
-\frac12q_{\mathrm{gh}}^2
-q_{\mathrm{gh}}
+N
+L_0^{X},
\] where \(N\) denotes oscillator excitations omitted from the simple
exponential vertex and \(L_0^X\) contains the spacetime-boson
contribution.

Substitution into Eq. (13.114) gives \[
\chi(\tau)
\sim
\sum_{\boldsymbol\lambda,q_{\mathrm{gh}},N}
e^{2\pi i\tau
\left(
\frac12\boldsymbol\lambda^2
-\frac12q_{\mathrm{gh}}^2
-q_{\mathrm{gh}}
+N+L_0^X-1
\right)}
(-1)^{F_R}.
\] The book then first isolates the nontrivial lattice zero-mode part.
Since \(q_{\mathrm{gh}}\) is integral on spacetime bosons and
half-integral on spacetime fermions, \[
\boxed{
(-1)^{F_R}=e^{-2\pi i q_{\mathrm{gh}}}.
}
\] Indeed, \[
\begin{aligned}
q_{\mathrm{gh}}\in\mathbb Z
&\quad\Longrightarrow\quad
e^{-2\pi iq_{\mathrm{gh}}}=+1,\\%1
q_{\mathrm{gh}}\in\mathbb Z+\frac12
&\quad\Longrightarrow\quad
e^{-2\pi iq_{\mathrm{gh}}}=-1.
\end{aligned}
\] This yields the lattice factor displayed in Eq. (13.115): \[
\widetilde\chi(\tau)
=
\sum_{\mathbf w=(\boldsymbol\lambda,q_{\mathrm{gh}})\in D_{5,1}}
e^{2\pi i\tau
\left(
\frac12\boldsymbol\lambda^2
-\frac12q_{\mathrm{gh}}^2
-q_{\mathrm{gh}}
\right)}
e^{-2\pi iq_{\mathrm{gh}}},
\] up to the universal oscillator and intercept factors temporarily
suppressed by the proportionality sign in Eq. (13.114). They are
restored in Eq. (13.122).

One can see the restoration of the full \(-1\) explicitly. Let \[
\mathbf e_6=(0,\ldots,0;1).
\] Because the last direction of \(D_{5,1}\) is timelike, \[
\mathbf e_6^2=-1,
\qquad
\mathbf w\cdot\mathbf e_6=-q_{\mathrm{gh}}.
\] The shifted lattice exponent in Eq. (13.122) is therefore \[
\begin{aligned}
\frac12(\mathbf w+\mathbf e_6)^2
&=
\frac12\mathbf w^2
+\mathbf w\cdot\mathbf e_6
+\frac12\mathbf e_6^2\\%1
&=
\frac12\boldsymbol\lambda^2
-\frac12q_{\mathrm{gh}}^2
-q_{\mathrm{gh}}
-\frac12.
\end{aligned}
\] Moreover, \[
\frac{1}{\eta(\tau)^{12}}
=
q_\tau^{-1/2}
\prod_{n=1}^\infty(1-q_\tau^n)^{-12}.
\] The \(-\tfrac12\) in the shifted lattice exponent and the
\(-\tfrac12\) from \(\eta^{-12}\) add to \[
-\frac12-\frac12=-1.
\] Consequently the complete expression has precisely the factor \[
q_\tau^{\,L_0-1}
\] announced in Eq. (13.114); Eq. (13.115) was displaying only the
selected zero-mode part.

The book immediately warns that the unrestricted \(D_{5,1}\) lattice sum
is not yet the physical partition function. It sums over infinitely many
equivalent ghost pictures and includes longitudinal, timelike, and ghost
degrees of freedom. The subsequent decomposition \[
D_{5,1}=D_4\oplus D_{1,1}
\] separates the transverse \(SO(8)\) degrees of freedom from the
unphysical sector. Dividing by the \(D_{1,1}\) lattice factor and fixing
the canonical-picture representative leaves the physical light-cone
partition function.

Thus Eq. (13.114) is a schematic covariant starting trace. It should not
be read as a literal trace over BRST cohomology in which \(L_0-1\) has
already been set to zero.

\subsubsection{Checks}\label{checks-16}

\begin{enumerate}
\def\labelenumi{\arabic{enumi}.}
\item
  \textbf{Open bosonic tachyon:} Omitting momentum, the \(N=0\) state
  contributes \[
  q_\tau^{-1},
  \] which is the leading term of \(\eta(\tau)^{-24}\).
\item
  \textbf{Open bosonic massless level:} The \(N=1\) oscillator states
  contribute at order \[
  q_\tau^0.
  \] They are massless, but they do not remove the other oscillator
  levels from the trace.
\item
  \textbf{On-shell pole:} When a momentum is analytically continued so
  that \[
  L_0-1=0,
  \] the propagator \((L_0-1)^{-1}\) develops a pole. This confirms that
  the physical-state equation identifies a pole rather than an operator
  to be set to zero throughout the loop calculation.
\item
  \textbf{Closed-string level matching:} The \(\tau_1\) integral removes
  terms with \[
  L_0-\bar L_0\neq0.
  \] This is why terms such as \(q_\tau^{-1}\) without the matching
  \(\bar q_\tau^{-1}\) partner occur in the unintegrated partition
  function but do not survive the full modulus integral.
\item
  \textbf{RNS intercepts:} The uniform covariant expression \(L_0-1\)
  becomes \(N_{\mathrm{NS}}-1/2\) in the light-cone NS sector and
  \(N_{\mathrm R}\) in the light-cone R sector after the superghost and
  unphysical degrees of freedom are eliminated.
\item
  \textbf{Spacetime supersymmetry:} After the GSO projection, the
  supertrace weighted by \((-1)^{F_R}\) cancels equal-mass bosonic and
  fermionic states. This is the Hamiltonian interpretation of the
  vanishing Jacobi combination in the ten-dimensional supersymmetric
  partition function.
\end{enumerate}

\subsubsection{Result}\label{result-20}

The plane-to-cylinder map gives \[
\boxed{
H_{\mathrm{cyl}}=L_0-\frac{c}{24}.
}
\] Therefore \[
\boxed{
\begin{aligned}
Z_{\mathrm{open}}(t)
&=
\operatorname{Tr}_{\mathcal H_{\mathrm{CFT}}}
e^{-2\pi t(L_0-c/24)},\\%1
Z_{\mathrm{closed}}(\tau,\bar\tau)
&=
\operatorname{Tr}_{\mathcal H_{\mathrm{CFT}}}
\left(
q_\tau^{\,L_0-c/24}
\bar q_\tau^{\,\bar L_0-\bar c/24}
\right).
\end{aligned}
}
\] For the 24 transverse bosons, \[
\frac{c_\perp}{24}=1,
\qquad
H_{\mathrm{cyl}}=L_0^\perp-1.
\] The bosonic physical-state equation \[
(L_0^{\mathrm m}-1)|\phi\rangle=0
\] means \[
\boxed{
\begin{aligned}
\text{open:}\qquad
m_N^2&=\frac{N-1}{\alpha'},\\%1
\text{closed:}\qquad
m_N^2&=\frac{4(N-1)}{\alpha'},
\qquad N=\bar N.
\end{aligned}
}
\] It fixes the mass of each oscillator level; it does not restrict the
spectrum to \(N=1\) and does not put the Euclidean loop momentum on
shell.

The proper-time identity gives \[
\boxed{
\frac{1}{L_0-1}
=
\int_0^\infty ds\,e^{-s(L_0-1)},
}
\] so the on-shell equation \(L_0-1=0\) determines the propagator poles,
while the one-loop trace constructs the propagator.

Finally, Eq. (13.114) is \[
\boxed{
\chi(\tau)
\sim
\operatorname{Tr}
\left[
q_\tau^{\,L_0-1}(-1)^{F_R}
\right].
}
\] Its Hamiltonian is the chiral worldsheet generator \[
\boxed{
H_{\mathrm{chiral}}=L_0-1,
}
\] and not the target-space energy \(P^0\). The trace is a spacetime
supertrace over the gauge-fixed chiral string state space; after
ghost-picture reduction and cancellation of unphysical degrees of
freedom, it becomes the physical light-cone partition function.

\[
A_{\perp}=\mathbb{Z}_{\geq 0}\times \mathbb{Z}_{>0}\cup \{ \perp \}
\] \[
(a,b)\cdot(c,d)=\left\{\begin{array}{ll}
(c,d-a) & d>a \\
\perp & d\leq a
\end{array}
\right.
\] \[
(a,b)\cdot \perp=\perp \ , \ \perp\cdot(a,b)=(a,b)
\]

\chapter{CFTs for Type II and Heterotic String
Vacua}\label{cfts-for-type-ii-and-heterotic-string-vacua-1}

\section{Q1: Untwisted and twisted cohomology of toroidal
orbifolds}\label{q1-untwisted-and-twisted-cohomology-of-toroidal-orbifolds}

In page 523-524, it is said that: \textgreater It is straightforward to
work out which elements of the Dolbeault cohomology of the torus survive
the orbifold projection. This is referred to as the untwisted cohomology
of the orbifold. There will be additional contributions from the twisted
sectors which are related to moduli controlling the repair of the
orbifold singularities.

The results are listed in Table 15.2.

I don't know how to obtain both the untwisted cohomology and the twisted
cohomology. Consider some examples in table 15.2, give the explicit
computation.

\begin{quote}
\textbf{Textbook location and related material.} Chapter 15, Section
15.2 (pp.~522--534), especially pp.~523--524 and Table 15.2: invariant
torus cohomology, twisted sectors, fixed loci, and blow-up modes.
Section 14.5 supplies the Calabi--Yau cohomology interpretation.
\end{quote}

\subsection{Answer}\label{answer-21}

\subsubsection{Core point}\label{core-point-21}

The entries in Table 15.2 are obtained in two logically different steps.

First, one takes the ordinary Dolbeault cohomology of the covering torus
and keeps only the forms invariant under the orbifold generator: \[
H_{\mathrm{untw}}^{p,q}(T^6/\mathbb Z_N)
=H^{p,q}(T^6)^{\mathbb Z_N}.
\] These dimensions are the numbers \(h_0^{p,q}\) written in parentheses
in the table.

Second, for every nonidentity element \(\theta^k\), one studies its
fixed locus. A twisted string is closed only up to \(\theta^k\), \[
Z^i(\sigma+2\pi)
=e^{2\pi i k v_i}Z^i(\sigma),
\] so its center of mass is confined to
\(\operatorname{Fix}(\theta^k)\). Massless twisted states are therefore
organized by the cohomology of these fixed loci. Their Hodge bidegrees
are shifted by the age of the twist, and one must then impose the
orbifold projection by the full group.

The compact mathematical statement is \[
\boxed{
H_{\mathrm{orb}}^{p,q}([T^6/G])
=
H^{p,q}(T^6)^G
\oplus
\bigoplus_{k=1}^{N-1}
H^{p-a_k,q-a_k}\!\left(\operatorname{Fix}(\theta^k)\right)^G
},
\qquad G=\mathbb Z_N.
\] When the fixed locus is disconnected, the last expression means the
cohomology of the disjoint union of all components, followed by the
\(G\)-projection. For the supersymmetric toroidal orbifolds in Table
15.2, these orbifold Hodge numbers agree with the Hodge numbers of a
crepant resolution \(\widetilde X\).

This formula will be derived and then applied explicitly to:

\begin{enumerate}
\def\labelenumi{\arabic{enumi}.}
\tightlist
\item
  case 1, \(T^6/\mathbb Z_3\) on \(A_2\times A_2\times A_2\), where the
  twisted fixed loci are isolated points;
\item
  case 2, \(T^6/\mathbb Z_4\) on \(A_1\times A_1\times B_2\times B_2\),
  where the \(\theta^2\) sector contains fixed tori and produces both
  twisted \((1,1)\) and twisted \((2,1)\) classes.
\end{enumerate}

The second example explains why the total Hodge numbers depend on the
lattice realization as well as on the twist vector.

\subsubsection{Untwisted cohomology: invariant constant
forms}\label{untwisted-cohomology-invariant-constant-forms}

On a complex torus, every Dolbeault class has a constant representative.
If \[
V^{1,0\,*}
=\operatorname{span}_{\mathbb C}\{dz^1,dz^2,dz^3\},
\] then \[
H^{p,q}(T^6)
\cong
\bigwedge^p V^{1,0\,*}
\otimes
\bigwedge^q\overline{V^{1,0\,*}}.
\] The generator acts on the one-forms by \[
\theta^*dz^i=e^{2\pi i v_i}dz^i,
\qquad
\theta^*d\bar z^{\bar i}=e^{-2\pi i v_i}d\bar z^{\bar i}.
\] Consequently, \[
\theta^*
\left(
dz^{i_1}\wedge\cdots\wedge dz^{i_p}
\wedge
d\bar z^{\bar j_1}\wedge\cdots\wedge d\bar z^{\bar j_q}
\right)
=
e^{2\pi i\left(
\sum_{a=1}^p v_{i_a}
-
\sum_{b=1}^q v_{j_b}
\right)}
\left(\cdots\right).
\] The form survives precisely when \[
\sum_{a=1}^p v_{i_a}
-
\sum_{b=1}^q v_{j_b}
\in\mathbb Z.
\]

In particular, \[
h_0^{1,1}
=
\#\left\{(i,j):v_i-v_j\in\mathbb Z\right\},
\] because a basis of \(H^{1,1}(T^6)\) is
\(dz^i\wedge d\bar z^{\bar j}\). Similarly, \[
h_0^{2,1}
=
\#\left\{(i,j;k):
i<j,\ 
v_i+v_j-v_k\in\mathbb Z
\right\}.
\]

The supersymmetry condition \[
v_1+v_2+v_3=0\pmod 1
\] implies \[
\theta^*
\left(dz^1\wedge dz^2\wedge dz^3\right)
=dz^1\wedge dz^2\wedge dz^3.
\] Thus the holomorphic volume form always survives, as stated around
equation (15.3).

\subsubsection{Twisted cohomology: fixed loci and the age
shift}\label{twisted-cohomology-fixed-loci-and-the-age-shift}

Consider the \(\theta^k\)-twisted sector. Define \[
\alpha_i^{(k)}=\{kv_i\},
\qquad
0\leq\alpha_i^{(k)}<1.
\] If \(\alpha_i^{(k)}=0\), the \(z^i\) direction is tangent to the
fixed locus. If \(\alpha_i^{(k)}\neq0\), it is a rotated normal
direction.

For a fixed component \(F\), the age is \[
a_k(F)
=
\sum_{\alpha_i^{(k)}\neq0}
\alpha_i^{(k)}.
\] In the present linear examples it is constant on all components of a
given fixed locus. A class in \(H^{r,s}(F)\) contributes after the age
shift to \[
H^{r+a_k,s+a_k}_{\mathrm{orb}}.
\]

The two cases needed below are especially simple.

For an isolated fixed point \(P\), \[
H^{0,0}(P)=\mathbb C.
\] It therefore contributes one class of bidegree \[
(a_k,a_k).
\] Thus an age-one isolated fixed point contributes to \(h^{1,1}\),
whereas an age-two fixed point contributes to \(h^{2,2}\).

For a fixed complex one-torus \(F\cong T^2\) and age \(a_k=1\), \[
\begin{array}{ccl}
H^{0,0}(F)&\longrightarrow&H_{\mathrm{orb}}^{1,1},\\
H^{1,0}(F)&\longrightarrow&H_{\mathrm{orb}}^{2,1},\\
H^{0,1}(F)&\longrightarrow&H_{\mathrm{orb}}^{1,2},\\
H^{1,1}(F)&\longrightarrow&H_{\mathrm{orb}}^{2,2}.
\end{array}
\] However, these classes survive only after projection by \(G\). If a
fixed torus is preserved by \(\theta\), one keeps only the
\(\theta\)-invariant forms on it. If several fixed tori are permuted,
one must form invariant linear combinations across the whole orbit.

The fixed components themselves are found by solving \[
\theta^k z=z+\lambda,
\qquad
\lambda\in\Lambda.
\] If there are no invariant directions, the number of solutions modulo
\(\Lambda\) is \[
\#\operatorname{Fix}(\theta^k)
=
\left|\det_{\mathbb R}(1-\theta^k)\right|
=
\prod_{i=1}^3
\left|1-e^{2\pi i k v_i}\right|^2
=
\prod_{i=1}^3
4\sin^2(\pi k v_i).
\] If \(\theta^k\) has invariant directions, this determinant vanishes
because the fixed locus is positive-dimensional. One then takes the
determinant only on the normal lattice to count components and
separately analyzes how \(G\) permutes those components.

\subsubsection{\texorpdfstring{Example 1: case 1, \(T^6/\mathbb Z_3\) on
\(A_2^3\)}{Example 1: case 1, T\^{}6/\textbackslash mathbb Z\_3 on A\_2\^{}3}}\label{example-1-case-1-t6mathbb-z_3-on-a_23}

The twist vector is \[
\mathbf v
=
\left(\frac13,\frac13,-\frac23\right).
\] Since \(-\frac23\equiv\frac13\pmod1\), the generator acts on all
three complex coordinates by the same phase: \[
\theta:(z^1,z^2,z^3)
\longmapsto
(\omega z^1,\omega z^2,\omega z^3),
\qquad
\omega=e^{2\pi i/3}.
\]

\paragraph{\texorpdfstring{Untwisted \((1,1)\)
classes}{Untwisted (1,1) classes}}\label{untwisted-11-classes}

Every basis element of \(H^{1,1}(T^6)\) is \[
\omega_{i\bar j}
=dz^i\wedge d\bar z^{\bar j},
\qquad
i,j=1,2,3.
\] Its phase is \[
\omega\bar\omega=1.
\] Equivalently, \(v_i-v_j=0\pmod1\) for every pair \((i,j)\). Therefore
all nine torus \((1,1)\) classes survive: \[
\boxed{h_0^{1,1}=3\times3=9}.
\]

\paragraph{\texorpdfstring{Untwisted \((2,1)\)
classes}{Untwisted (2,1) classes}}\label{untwisted-21-classes}

A basis element is \[
dz^i\wedge dz^j\wedge d\bar z^{\bar k},
\qquad i<j.
\] Its phase is \[
\omega^2\bar\omega
=\omega^2\omega^{-1}
=\omega\neq1.
\] Hence none of the nine \((2,1)\) classes survives: \[
\boxed{h_0^{2,1}=0}.
\] This reproduces the parenthesized entries \((9)\) and \((0)\) in the
first row of Table 15.2.

\paragraph{\texorpdfstring{Fixed points in the \(\theta\)
sector}{Fixed points in the \textbackslash theta sector}}\label{fixed-points-in-the-theta-sector}

Each \(A_2\) root lattice is a hexagonal lattice preserved by a
\(120^\circ\) rotation. On one complex torus, the number of fixed points
is \[
\left|1-\omega\right|^2
=
(1-\omega)(1-\bar\omega)
=3.
\] Equivalently, if the Coxeter action on an integral \(A_2\) basis is
represented by \[
R=
\begin{pmatrix}
0&-1\\
1&-1
\end{pmatrix},
\] then \[
\left|\det(1-R)\right|
=
\left|
\det
\begin{pmatrix}
1&1\\
-1&2
\end{pmatrix}
\right|
=3.
\] Because the six-torus is the product of three such factors, \[
\#\operatorname{Fix}(\theta)
=3^3
=27.
\]

For \(k=1\) the fractional twists are \[
\left(
\alpha_1^{(1)},\alpha_2^{(1)},\alpha_3^{(1)}
\right)
=
\left(\frac13,\frac13,\frac13\right),
\] so \[
a_1=\frac13+\frac13+\frac13=1.
\] Every fixed point therefore supplies one twisted \((1,1)\) class. The
27 fixed points are individually fixed by the full \(\mathbb Z_3\)
action, so the orbifold projection does not identify them: \[
\boxed{
\Delta_{\theta}h^{1,1}=27,
\qquad
\Delta_{\theta}h^{2,1}=0
}.
\]

Locally, each singularity is of type \[
\mathbb C^3/\mathbb Z_3,
\qquad
\frac13(1,1,1).
\] Its crepant resolution introduces an exceptional divisor. The
corresponding divisor class is the geometric realization of the age-one
twisted \((1,1)\) state, and its Kähler parameter controls the blow-up
of that singularity.

\paragraph{\texorpdfstring{The \(\theta^2\) sector and why it does not
give another 27 to
\(h^{1,1}\)}{The \textbackslash theta\^{}2 sector and why it does not give another 27 to h\^{}\{1,1\}}}\label{the-theta2-sector-and-why-it-does-not-give-another-27-to-h11}

The element \(\theta^2\) has the same 27 fixed points, but its
fractional twists are \[
\left(
\alpha_1^{(2)},\alpha_2^{(2)},\alpha_3^{(2)}
\right)
=
\left(\frac23,\frac23,\frac23\right).
\] Thus \[
a_2=\frac23+\frac23+\frac23=2.
\] The point class in \(H^{0,0}\) is now shifted to bidegree \((2,2)\),
not \((1,1)\): \[
\Delta_{\theta^2}h^{2,2}=27,
\qquad
\Delta_{\theta^2}h^{1,1}=0.
\] This is the Hodge-dual partner of the \(\theta\) contribution. The
two CFT twisted sectors therefore do not give \(54\) independent Kähler
moduli.

Combining the untwisted and age-one twisted contributions gives \[
\boxed{
h^{1,1}=9+27=36,
\qquad
h^{2,1}=0
}.
\] This is precisely case 1 of Table 15.2.

\subsubsection{\texorpdfstring{Example 2: case 2, \(T^6/\mathbb Z_4\) on
\(A_1^2\times B_2^2\)}{Example 2: case 2, T\^{}6/\textbackslash mathbb Z\_4 on A\_1\^{}2\textbackslash times B\_2\^{}2}}\label{example-2-case-2-t6mathbb-z_4-on-a_12times-b_22}

Now take \[
\mathbf v
=
\left(\frac14,\frac14,-\frac12\right).
\] Choose the first two complex planes to be the two square \(B_2\) tori
and the third to be the \(A_1\times A_1\) torus. Then \[
\theta:(z^1,z^2,z^3)
\longmapsto
(iz^1,iz^2,-z^3).
\]

\paragraph{\texorpdfstring{Untwisted \((1,1)\)
classes}{Untwisted (1,1) classes}}\label{untwisted-11-classes-1}

For \(dz^i\wedge d\bar z^{\bar j}\) to be invariant, \(z^i\) and \(z^j\)
must have the same eigenvalue. The first two coordinates both have
eigenvalue \(i\), while the third has eigenvalue \(-1\). Hence the
invariant forms are \[
\begin{aligned}
&dz^1\wedge d\bar z^{\bar1},
\qquad
dz^1\wedge d\bar z^{\bar2},
\qquad
dz^2\wedge d\bar z^{\bar1},
\qquad
dz^2\wedge d\bar z^{\bar2},\\
&dz^3\wedge d\bar z^{\bar3}.
\end{aligned}
\] Thus \[
\boxed{h_0^{1,1}=4+1=5}.
\]

\paragraph{\texorpdfstring{Untwisted \((2,1)\)
classes}{Untwisted (2,1) classes}}\label{untwisted-21-classes-1}

Consider \[
\chi
=
dz^1\wedge dz^2\wedge d\bar z^{\bar3}.
\] Its phase is \[
i\cdot i\cdot(-1)=1,
\] so it is invariant. It is the only invariant \((2,1)\) basis form.
For example, \[
dz^1\wedge dz^3\wedge d\bar z^{\bar1}
\] has phase \[
i\cdot(-1)\cdot(-i)=-1,
\] and similarly for the remaining basis elements. Therefore \[
\boxed{h_0^{2,1}=1}.
\] This reproduces the parenthesized \((5)\) and \((1)\) in the second
row of Table 15.2.

\paragraph{\texorpdfstring{The \(\theta\) and \(\theta^3\) sectors:
isolated fixed
points}{The \textbackslash theta and \textbackslash theta\^{}3 sectors: isolated fixed points}}\label{the-theta-and-theta3-sectors-isolated-fixed-points}

In the \(\theta\) sector, a quarter-turn on a square two-torus has \[
\left|1-i\right|^2=2
\] fixed points, while inversion on a two-torus has \[
\left|1-(-1)\right|^2=4
\] fixed points. Hence \[
\#\operatorname{Fix}(\theta)
=2\cdot2\cdot4
=16.
\] The fractional twists are \[
\left(\frac14,\frac14,\frac12\right),
\] and therefore \[
a_1=\frac14+\frac14+\frac12=1.
\] The \(\theta\) sector consequently contributes \[
\boxed{
\Delta_{\theta}h^{1,1}=16,
\qquad
\Delta_{\theta}h^{2,1}=0
}.
\]

The inverse element \(\theta^3\) has fractional twists \[
\left(\frac34,\frac34,\frac12\right),
\] so \[
a_3=\frac34+\frac34+\frac12=2.
\] It has the same 16 fixed points but contributes to \(h^{2,2}\) rather
than \(h^{1,1}\): \[
\Delta_{\theta^3}h^{2,2}=16.
\]

\paragraph{\texorpdfstring{The \(\theta^2\) sector: 16 fixed tori before
the orbifold
projection}{The \textbackslash theta\^{}2 sector: 16 fixed tori before the orbifold projection}}\label{the-theta2-sector-16-fixed-tori-before-the-orbifold-projection}

The square of the generator acts as \[
\theta^2:(z^1,z^2,z^3)
\longmapsto
(-z^1,-z^2,z^3).
\] The first two coordinates must be two-torsion points, while \(z^3\)
is arbitrary. Each of the first two tori has four two-torsion points, so
\[
\operatorname{Fix}(\theta^2)
=
\coprod_{a=1}^{16}F_a,
\qquad
F_a\cong T^2_{(3)}.
\] Thus there are 16 fixed two-tori before the remaining \(\mathbb Z_4\)
projection.

The fractional twists normal to a fixed torus are \[
\left(\frac12,\frac12,0\right),
\] and hence \[
a_2=\frac12+\frac12=1.
\] Before projection, \(H^{0,0}(F_a)\) could therefore give a \((1,1)\)
class and \(H^{1,0}(F_a)\) could give a \((2,1)\) class. The crucial
next step is to determine the action of \(\theta\) on the 16 components
and on their one-forms.

\paragraph{Orbifold projection on the fixed
tori}\label{orbifold-projection-on-the-fixed-tori}

Represent either square \(B_2\) torus as \[
\mathbb C/(\mathbb Z+i\mathbb Z).
\] Its four points fixed by inversion are \[
0,
\qquad
\frac12,
\qquad
\frac{i}{2},
\qquad
\frac{1+i}{2}.
\] Multiplication by \(i\) acts on them as \[
0\longmapsto0,
\qquad
\frac{1+i}{2}\longmapsto\frac{1+i}{2},
\qquad
\frac12\longleftrightarrow\frac{i}{2},
\] where equality is understood modulo the lattice. Thus, on the four
two-torsion points of each of the first two tori, \(\theta\) has two
one-element orbits and one two-element orbit.

On the Cartesian product of the two four-point sets:

\begin{itemize}
\tightlist
\item
  \(2\times2=4\) fixed-torus components are preserved individually by
  \(\theta\);
\item
  the remaining \(16-4=12\) components form \(12/2=6\) exchanged pairs.
\end{itemize}

Therefore the 16 components form \[
4+6=10
\] orbits under \(\theta\).

For \(H^{0,0}\), each orbit supplies one invariant constant combination.
After the age-one shift, this gives \[
\boxed{\Delta_{\theta^2}h^{1,1}=10}.
\]

The \(H^{1,0}\) calculation is subtler. Along every fixed torus, the
holomorphic one-form is \(dz^3\), and \[
\theta^*dz^3=-dz^3.
\]

On each of the four individually preserved components, \(dz^3\) is odd,
so it is removed by the projection.

For an exchanged pair \((F,F')\), let \(\eta_F\) and \(\eta_{F'}\)
denote the copies of \(dz^3\) supported on the two components. Then \[
\theta^*\eta_F=-\eta_{F'},
\qquad
\theta^*\eta_{F'}=-\eta_F.
\] The antisymmetric combination \[
\eta_F-\eta_{F'}
\] is invariant: \[
\theta^*(\eta_F-\eta_{F'})
=
-\eta_{F'}+\eta_F
=
\eta_F-\eta_{F'}.
\] Each of the six exchanged pairs therefore contributes one invariant
\(H^{1,0}\) class. After the age-one shift, \[
H^{1,0}(F)\longrightarrow H_{\mathrm{orb}}^{2,1},
\] so \[
\boxed{\Delta_{\theta^2}h^{2,1}=6}.
\] Complex conjugation similarly gives six \(H^{1,2}\) classes.

\paragraph{Total for case 2}\label{total-for-case-2}

Collecting the contributions by sector, \[
\begin{aligned}
h^{1,1}
&=
\underbrace{5}_{\text{untwisted}}
+
\underbrace{16}_{\theta}
+
\underbrace{10}_{\theta^2}
=31,\\
h^{2,1}
&=
\underbrace{1}_{\text{untwisted}}
+
\underbrace{0}_{\theta}
+
\underbrace{6}_{\theta^2}
=7.
\end{aligned}
\] Thus \[
\boxed{(h^{1,1},h^{2,1})=(31,7)},
\] exactly as in case 2 of Table 15.2.

\subsubsection{Why the twist vector alone does not determine the twisted
answer}\label{why-the-twist-vector-alone-does-not-determine-the-twisted-answer}

The untwisted count depends only on the eigenvalues \(e^{2\pi i v_i}\)
acting on the complex one-forms. This is why cases 2, 3, and 4, which
all use the \(\mathbb Z_4\) twist \[
\mathbf v
=
\left(\frac14,\frac14,-\frac12\right),
\] all have \[
(h_0^{1,1},h_0^{2,1})=(5,1).
\]

Their total Hodge numbers differ: \[
(31,7),
\qquad
(27,3),
\qquad
(25,1).
\] The reason is that the three cases use different six-dimensional
lattices. The number of fixed components, the way \(\theta\) permutes
the fixed tori of \(\theta^2\), and the action on the cohomology of
those tori all depend on the integral lattice automorphism. Thus the
twist vector determines the local rotation angles and age shifts, but
the lattice is also required to determine the global multiplicities.

\subsubsection{\texorpdfstring{Relation to the book's
\(T^6/\mathbb Z_3\) massless
spectrum}{Relation to the book's T\^{}6/\textbackslash mathbb Z\_3 massless spectrum}}\label{relation-to-the-books-t6mathbb-z_3-massless-spectrum}

The geometric count agrees directly with equations (15.21)--(15.22).

For type IIA on a Calabi-Yau threefold, the number of vector multiplets
is \(h^{1,1}\) and the number of hypermultiplets is \(h^{2,1}+1\), where
the extra one is the universal hypermultiplet. In the \(\mathbb Z_3\)
example, \[
h^{1,1}=9_{\mathrm{untw}}+27_{\mathrm{tw}},
\qquad
h^{2,1}=0.
\] Therefore type IIA has \[
9\ \text{untwisted vector multiplets}
+
27\ \text{twisted vector multiplets},
\] together with one universal hypermultiplet, exactly as in equation
(15.22).

For type IIB, the Kähler moduli instead occur in hypermultiplets. Thus
the same cohomology gives ten untwisted hypermultiplets---nine from
\(h_0^{1,1}=9\) plus the universal hypermultiplet---and 27 twisted
hypermultiplets, as in equation (15.21).

This also makes precise the sentence following Table 15.2: the 27
age-one twisted classes are the Kähler moduli controlling the repair of
the 27 orbifold singularities.

\subsubsection{Checks}\label{checks-17}

The age of a twist and that of its inverse obey \[
a_k+a_{N-k}
=
\operatorname{codim}_{\mathbb C}\operatorname{Fix}(\theta^k).
\] For the isolated \(\mathbb Z_3\) fixed points, \[
1+2=3,
\] so the \(\theta\) and \(\theta^2\) sectors contribute to the
Hodge-dual groups \(H^{1,1}\) and \(H^{2,2}\). For the isolated
\(\mathbb Z_4\) fixed points, the same check is \[
1+2=3.
\] For the \(\theta^2\) fixed tori, the complex codimension is two and
\(\theta^2\) is self-inverse: \[
1+1=2.
\] These relations ensure compatibility with Poincaré duality.

The resulting Euler characteristics are \[
\chi(\widetilde X_{\mathbb Z_3})
=2(h^{1,1}-h^{2,1})
=2(36-0)
=72,
\] and \[
\chi(\widetilde X_{\mathbb Z_4})
=2(31-7)
=48.
\]

\subsubsection{Result}\label{result-21}

For \(T^6/\langle\theta\rangle\) with \[
\theta:z^i\mapsto e^{2\pi i v_i}z^i,
\] the untwisted classes are selected by \[
\boxed{
\sum_{a=1}^p v_{i_a}
-
\sum_{b=1}^q v_{j_b}
\in\mathbb Z
}.
\] The twisted contribution of \(\theta^k\) is obtained from \[
\boxed{
\alpha_i^{(k)}=\{kv_i\},
\qquad
a_k=\sum_{\alpha_i^{(k)}\neq0}\alpha_i^{(k)},
\qquad
H^{r,s}(F)\mapsto H_{\mathrm{orb}}^{r+a_k,s+a_k}
},
\] followed by the full \(G\)-projection on the cohomology of the union
of fixed components.

For the two explicit examples, \[
\begin{array}{c|c|c|c|c}
\text{orbifold}
&
(h_0^{1,1},h_0^{2,1})
&
\theta\text{ contribution}
&
\theta^2\text{ contribution}
&
(h^{1,1},h^{2,1})
\\ \hline
T^6/\mathbb Z_3\text{ on }A_2^3
&
(9,0)
&
(27,0)
&
(0,0)\text{ in these bidegrees}
&
(36,0)
\\
T^6/\mathbb Z_4\text{ on }A_1^2\times B_2^2
&
(5,1)
&
(16,0)
&
(10,6)
&
(31,7)
\end{array}
\] In the \(\mathbb Z_3\) example, the \(\theta^2\) sector contributes
27 classes to \(h^{2,2}\); in the \(\mathbb Z_4\) example, the
\(\theta^3\) sector contributes 16 classes to \(h^{2,2}\). They are the
Hodge-dual partners of the age-one isolated-point contributions and are
not additional \(h^{1,1}\) moduli.

\section{Q2: Origin of the enhanced orbifold gauge
symmetry}\label{q2-origin-of-the-enhanced-orbifold-gauge-symmetry}

In page 527: \textgreater While for smooth CY compactifications this
leads to the gauge group \(E_{6}\times E_{8}\), in the orbifold limit it
is enhanced to at least \(E_{6}\times G\times E_{8}\) where \(G\) can be
\(SU(3),SU(2)\times U(1)\) or \(U(1)\times U(1)\), depending on the
shift vector. The blow-up modes, which are the elements of the
cohomology in the twisted sectors, are charged under \(G\).

I want to know where the gauge bosons of \(G\) come from i.e.~why the
symmetry in this case can be enhanced.

\begin{quote}
\textbf{Textbook location and related material.} Chapter 15, Section
15.2 around p.~527: standard embedding, orbifold gauge shift, enhanced
group \(G\), and charged blow-up modes. Sections 10.4--10.5 and
11.2--11.5 give the heterotic current-algebra and lattice machinery.
\end{quote}

\subsection{Answer}\label{answer-22}

\subsubsection{Core point}\label{core-point-22}

The gauge bosons of the extra factor \(G\) do \textbf{not} come from the
twisted sectors and do not arise from wrapped branes on the collapsing
cycles. They are ordinary perturbative heterotic gauge bosons already
present in the ten-dimensional first \(E_8\). More precisely, they are
untwisted left-moving current-algebra states that survive the orbifold
projection.

The difference between the smooth and orbifold compactifications is the
group with which a gauge generator must commute: \[
\begin{array}{ccl}
\text{smooth standard embedding}
&:&
\operatorname{Hol}(A^{\mathrm{bg}})=SU(3),
\\
\text{orbifold standard embedding}
&:&
\operatorname{Hol}_{\mathrm{orb}}
=
\langle\gamma_{\mathbf V}\rangle
\subset SU(3),
\end{array}
\] where \(\gamma_{\mathbf V}\) is a finite-order element determined by
the gauge shift \(\mathbf V\).

A generator that commutes with the full continuous group \(SU(3)\)
automatically commutes with \(\gamma_{\mathbf V}\), but the converse
need not hold. Therefore \[
C_{E_8}(SU(3))
\subseteq
C_{E_8}(\gamma_{\mathbf V}),
\] and the inclusion is generally strict. The two centralizers are \[
C_{E_8}(SU(3))=E_6,
\] and, for the standard orbifold shifts considered here, \[
C_{E_8}(\gamma_{\mathbf V})
\supseteq
E_6\times C_{SU(3)}(\gamma_{\mathbf V}).
\] The book's extra factor is \[
\boxed{
G_{\mathrm{enh}}
=
C_{SU(3)}(\gamma_{\mathbf V})
}.
\]

Thus the enhancement is a \textbf{larger-centralizer effect}: the
singular orbifold background imposes only a discrete gauge holonomy,
which projects out fewer of the original \(E_8\) currents than a smooth
connection with full \(SU(3)\) holonomy.

\subsubsection{What the book explicitly
establishes}\label{what-the-book-explicitly-establishes}

The discussion on pages 525--526 represents the heterotic gauge degrees
of freedom by a level-one current algebra. In the bosonic realization,
the generator of the orbifold point group acts on the sixteen chiral
gauge bosons as the shift \[
X_L^I
\longmapsto
X_L^I+V^I,
\] which is equation (15.7). The order-\(N\) condition is \[
N\mathbf V\in\Lambda_{E_8\times E_8}.
\]

The important source statements are:

\begin{enumerate}
\def\labelenumi{\arabic{enumi}.}
\tightlist
\item
  only invariant currents give gauge bosons in the orbifold theory;
\item
  the unbroken gauge algebra is the subalgebra commuting with the gauge
  embedding of the orbifold group;
\item
  page 528 explicitly states that gauge bosons cannot arise from twisted
  sectors.
\end{enumerate}

The last point is essential. The twisted fields mentioned in the
question are charged scalar multiplets. They break the enhanced symmetry
when they acquire expectation values, but they are not the source of its
vector multiplets.

\subsubsection{Gauge bosons as invariant current-algebra
states}\label{gauge-bosons-as-invariant-current-algebra-states}

At level one, the first \(E_8\) current algebra contains eight Cartan
currents \[
H^I(z_L)=i\partial X_L^I(z_L),
\] and, for every root \(\mathbf P\) of \(E_8\), a root current of the
schematic form \[
J_{\mathbf P}(z_L)
=
:\!e^{i\mathbf P\cdot\mathbf X_L(z_L)}\!:\, .
\] The normalization in this schematic exponential is fixed by the
gauge-lattice periodicities. In the shift-vector convention of equations
(15.8)--(15.11), the gauge shift acts as \[
\begin{aligned}
H^I
&\longmapsto H^I,\\
J_{\mathbf P}
&\longmapsto
e^{2\pi i\,\mathbf P\cdot\mathbf V}J_{\mathbf P}.
\end{aligned}
\] Hence:

\begin{itemize}
\tightlist
\item
  every Cartan current survives an Abelian shift;
\item
  a root current survives precisely if \[
  \boxed{
  \mathbf P\cdot\mathbf V\in\mathbb Z
  }.
  \]
\end{itemize}

The surviving currents close under operator products and generate the
unbroken current algebra.

A four-dimensional gauge-boson state has the schematic heterotic form \[
\widetilde\psi_{-1/2}^{\mu}|0;k\rangle_R
\otimes
J_{-1}^{T}|0\rangle_L,
\] where \(T\) is a gauge generator. Equivalently, its vertex operator
contains \[
\mathcal V_A
\sim
\epsilon_\mu\,
\widetilde\psi^\mu
e^{-\widetilde\phi}
e^{ik\cdot X}
J^T.
\] The right-moving factor supplies the spacetime vector polarization,
while the left-moving dimension-one current supplies the gauge
generator. The orbifold projection retains this vector if \[
\gamma_{\mathbf V}T\gamma_{\mathbf V}^{-1}=T.
\] Therefore the vector multiplets are labeled precisely by \[
\operatorname{Lie}C_{E_8}(\gamma_{\mathbf V}).
\]

Nothing new has to become massless in a twisted sector. At the orbifold
point, currents already present in the ten-dimensional \(E_8\) theory
are simply projected less strongly than they are by a smooth \(SU(3)\)
bundle background.

\subsubsection{\texorpdfstring{Smooth standard embedding: why only
\(E_6\)
remains}{Smooth standard embedding: why only E\_6 remains}}\label{smooth-standard-embedding-why-only-e_6-remains}

For the smooth standard embedding one identifies the gauge connection in
an \(SU(3)\subset E_8\) subgroup with the Calabi-Yau spin connection: \[
A_m^{\mathrm{bg}}=\omega_m.
\] Assuming the Calabi-Yau has holonomy exactly \(SU(3)\), the
gauge-bundle holonomy is also the full group \(SU(3)\).

Expand a possible four-dimensional vector as \[
A_\mu(x,y)
=
\sum_a A_\mu^a(x)\,s_a(y),
\] where \(s_a\) is a section of the adjoint bundle. A massless gauge
boson must have a covariantly constant internal wavefunction: \[
D_m^{\mathrm{bg}}s_a=0.
\] Parallel transport around every closed loop then requires \[
h\,s_a\,h^{-1}=s_a
\qquad
\text{for every }
h\in\operatorname{Hol}(A^{\mathrm{bg}}).
\] Thus the unbroken gauge algebra is \[
\mathfrak g_{\mathrm{4d}}
=
\operatorname{Lie}
C_{E_8}\!\left(\operatorname{Hol}(A^{\mathrm{bg}})\right).
\] For full \(SU(3)\) holonomy, \[
\boxed{
\mathfrak g_{\mathrm{smooth}}
=
\operatorname{Lie}C_{E_8}(SU(3))
=
\mathfrak e_6
}.
\]

The same conclusion follows from the familiar branching \[
\mathbf{248}
=
(\mathbf8,\mathbf1)
\oplus
(\mathbf1,\mathbf{78})
\oplus
(\mathbf3,\mathbf{27})
\oplus
(\bar{\mathbf3},\overline{\mathbf{27}})
\] under \(SU(3)\times E_6\).

Only \((\mathbf1,\mathbf{78})\) is invariant under the full \(SU(3)\)
structure group. The \(\mathbf8\) contains no singlet under the adjoint
action of the entire \(SU(3)\): equivalently, the Lie algebra
\(\mathfrak{su}(3)\) has no nonzero element commuting with all of
\(\mathfrak{su}(3)\). Therefore the smooth background has \(E_6\), not
\(E_6\times SU(3)\), as its four-dimensional gauge group.

This conclusion can also be seen from the ten-dimensional Yang--Mills
action. For a mode \(A_\mu(x)T\) independent of the internal
coordinates, \[
F_{\mu m}
=
-\left[A_m^{\mathrm{bg}},A_\mu T\right].
\] The term \[
\int_X
\operatorname{tr}(F_{\mu m}F^{\mu m})
\] becomes a four-dimensional mass term proportional to \[
\int_X
\operatorname{tr}
\left(
[A_m^{\mathrm{bg}},T]
[A^{\mathrm{bg}\,m},T]
\right).
\] It vanishes exactly for generators commuting with the background
holonomy. A generic smooth \(SU(3)\) connection therefore removes the
would-be \(SU(3)\) gauge vectors.

\subsubsection{Orbifold standard embedding: only a discrete element must
commute}\label{orbifold-standard-embedding-only-a-discrete-element-must-commute}

In the orbifold CFT, the standard embedding is \[
\mathbf V
=
(v_1,v_2,v_3,0,\ldots,0),
\qquad
\mathbf V'=0.
\] The geometric point-group generator \(\theta\) is embedded as the
finite gauge transformation \[
\gamma_{\mathbf V}
=
\exp(2\pi i\,\mathbf V\cdot\mathbf H)
\in SU(3)\subset E_8.
\] Away from the orbifold singularities, the corresponding idealized
connection is locally flat. Its nontrivial information is the finite
monodromy \(\gamma_{\mathbf V}\) around the singular locus. Thus a gauge
generator need commute only with this one finite-order element, not with
every element of \(SU(3)\): \[
\gamma_{\mathbf V}T\gamma_{\mathbf V}^{-1}=T.
\]

Using the \(SU(3)\times E_6\) branching:

\begin{itemize}
\tightlist
\item
  all \((\mathbf1,\mathbf{78})\) currents survive, giving \(E_6\);
\item
  the invariant part of \((\mathbf8,\mathbf1)\) is \[
  \operatorname{Lie}C_{SU(3)}(\gamma_{\mathbf V}),
  \] giving \(G_{\mathrm{enh}}\);
\item
  the hidden second \(E_8\) is unaffected because \(\mathbf V'=0\);
\item
  for the four-dimensional supersymmetric twists in Table 15.1, all
  \(v_i\) are nonzero modulo one, so the mixed
  \((\mathbf3,\mathbf{27})\) and
  \((\bar{\mathbf3},\overline{\mathbf{27}})\) currents are projected out
  as gauge bosons.
\end{itemize}

This yields \[
\boxed{
E_6\times G_{\mathrm{enh}}\times E_8
\subseteq
C_{E_8\times E_8}(\gamma_{\mathbf V})
}.
\] The book says ``at least'' because in more general embeddings or for
special shifts there can be further invariant roots. For the minimal
four-dimensional standard-embedding examples with no Wilson lines, the
displayed group is the relevant result.

\subsubsection{\texorpdfstring{Computing \(G_{\mathrm{enh}}\) from the
three
eigenvalues}{Computing G\_\{\textbackslash mathrm\{enh\}\} from the three eigenvalues}}\label{computing-g_mathrmenh-from-the-three-eigenvalues}

Inside the \(SU(3)\) used in the standard embedding, the shift element
may be written in the defining representation as \[
\gamma_{\mathbf V}
=
\operatorname{diag}
\left(
e^{2\pi i v_1},
e^{2\pi i v_2},
e^{2\pi i v_3}
\right),
\] with \[
v_1+v_2+v_3=0\pmod1.
\] The off-diagonal \(SU(3)\) generator \(E_{ij}\) transforms as \[
\gamma_{\mathbf V}E_{ij}\gamma_{\mathbf V}^{-1}
=
e^{2\pi i(v_i-v_j)}E_{ij}.
\] Therefore \[
E_{ij}\text{ survives}
\quad\Longleftrightarrow\quad
v_i-v_j\in\mathbb Z.
\] The two Cartan generators always survive. The answer is then
determined by the degeneracy pattern of the three eigenvalues:

\[
\boxed{
\begin{array}{c|c}
\text{eigenvalues of }\gamma_{\mathbf V}
&
G_{\mathrm{enh}}
\\ \hline
\lambda_1=\lambda_2=\lambda_3
&
SU(3)
\\
\text{exactly two }\lambda_i\text{ are equal}
&
SU(2)\times U(1)
\\
\lambda_1,\lambda_2,\lambda_3\text{ all distinct}
&
U(1)\times U(1)
\end{array}
}
\]

This is the origin of the three possibilities listed on page 527.

\subsubsection{\texorpdfstring{Explicit \(\mathbb Z_3\) example: the
eight \(SU(3)\) gauge
bosons}{Explicit \textbackslash mathbb Z\_3 example: the eight SU(3) gauge bosons}}\label{explicit-mathbb-z_3-example-the-eight-su3-gauge-bosons}

For the \(T^6/\mathbb Z_3\) model discussed later in the chapter, \[
\mathbf v
=
\left(
\frac13,\frac13,-\frac23
\right),
\qquad
\mathbf V
=
\left(
\frac13,\frac13,-\frac23,0,\ldots,0
\right).
\] Modulo integers, all three phases are equal: \[
e^{2\pi i v_1}
=
e^{2\pi i v_2}
=
e^{2\pi i v_3}
=
\omega.
\] Thus \[
\gamma_{\mathbf V}
=
\omega\,\mathbf1_3.
\] This is a central element of \(SU(3)\), so it commutes with the
entire \(SU(3)\): \[
C_{SU(3)}(\gamma_{\mathbf V})=SU(3).
\]

In current-algebra language, take the six \(SU(3)\) roots \[
\mathbf P_{ij}=\mathbf e_i-\mathbf e_j,
\qquad
i\neq j,
\qquad
i,j\in\{1,2,3\}.
\] Their shift phases are controlled by \[
\mathbf P_{ij}\cdot\mathbf V
=
v_i-v_j.
\] Explicitly, \[
v_1-v_2=0,
\qquad
v_1-v_3=1,
\qquad
v_2-v_3=1,
\] and the negatives are also integers. Hence all six root currents
survive. Together with the two Cartan currents, they give \[
6+2=8
\] massless vector multiplets forming the adjoint \(\mathbf8\) of
\(SU(3)\).

The first \(E_8\) gauge group at this orbifold point is therefore \[
\boxed{
E_6\times SU(3)
},
\] and including the unaffected hidden factor gives \[
\boxed{
E_6\times SU(3)\times E_8
}.
\] These are precisely the untwisted gauge multiplets listed in equation
(15.24).

Notice that the central element \(\omega\mathbf1_3\) acts trivially on
the adjoint \(\mathbf8\) but nontrivially on the fundamental
\(\mathbf3\): \[
\mathbf8\longmapsto\mathbf8,
\qquad
\mathbf3\longmapsto\omega\,\mathbf3.
\] This is why the \(SU(3)\) gauge currents survive while the mixed
\((\mathbf3,\mathbf{27})\) currents do not become additional gauge
bosons.

\subsubsection{\texorpdfstring{Explicit \(\mathbb Z_4\) example:
\(SU(2)\times U(1)\)}{Explicit \textbackslash mathbb Z\_4 example: SU(2)\textbackslash times U(1)}}\label{explicit-mathbb-z_4-example-su2times-u1}

For \[
\mathbf v
=
\left(
\frac14,\frac14,-\frac12
\right),
\] the three eigenvalues are \[
(i,i,-1).
\] Only the first two coincide. Equivalently, \[
v_1-v_2=0,
\qquad
v_1-v_3=\frac34,
\qquad
v_2-v_3=\frac34.
\] Therefore the root currents \(J_{\mathbf e_1-\mathbf e_2}\) and
\(J_{\mathbf e_2-\mathbf e_1}\) survive, while those connecting the
third direction to the first two are projected out. The two Cartan
currents remain.

The surviving algebra consists of:

\begin{itemize}
\tightlist
\item
  one \(SU(2)\) Cartan generator and its two root generators;
\item
  one additional commuting \(U(1)\) Cartan generator.
\end{itemize}

Thus \[
\boxed{
G_{\mathrm{enh}}
=
SU(2)\times U(1)
}
\] up to a finite global quotient.

If all three phases are distinct, none of the six off-diagonal \(SU(3)\)
currents survives, but the two Cartan currents do. This gives \[
\boxed{
G_{\mathrm{enh}}=U(1)^2
}.
\]

\subsubsection{Why resolving the orbifold Higgses the extra
group}\label{why-resolving-the-orbifold-higgses-the-extra-group}

The twisted sectors contain massless scalar fields localized at the
orbifold fixed loci. Under the gauge shift, these states generally carry
nonzero weights of \(G_{\mathrm{enh}}\). The book therefore identifies
the blow-up modes with charged chiral multiplets.

Let \(b^A\) denote such fields. Their kinetic terms contain \[
\sum_A
\left|D_\mu b^A\right|^2,
\qquad
D_\mu b^A
=
\partial_\mu b^A
-igA_\mu^a(T_a b)^A.
\] If the resolution requires \[
\langle b^A\rangle\neq0,
\] then \[
\sum_A|D_\mu b^A|^2
\supset
g^2 A_\mu^a A^{\mu b}
\sum_A
(T_a\langle b\rangle)^A
(T_b\langle b\rangle)^{A*}.
\] The gauge bosons for generators satisfying \[
T_a\langle b\rangle\neq0
\] become massive. The unbroken group is the stabilizer of all blow-up
expectation values.

In the \(T^6/\mathbb Z_3\) example, equation (15.24) exhibits twisted
fields charged as \(\mathbf3\) under the enhanced \(SU(3)\). Moving
along suitable D-flat blow-up directions breaks this \(SU(3)\)
completely. As the book explains after equation (15.24), eight
chiral-multiplet directions combine with the eight \(SU(3)\) vector
multiplets to form massive vector multiplets.

Geometrically, the singular orbifold with finite gauge monodromy is
replaced by a smooth Calabi-Yau together with a smooth gauge connection
whose holonomy fills out \(SU(3)\). In the worldsheet description this
is a deformation by twisted operators; in the four-dimensional
description it is a Higgs effect. These are two descriptions of the same
transition: \[
\begin{gathered}
\boxed{\begin{array}{c}
\text{orbifold point: discrete }\langle\gamma_{\mathbf V}\rangle\\
E_6\times G_{\mathrm{enh}}\times E_8
\end{array}}\\[4pt]
\underset{\langle b_{\mathrm{blow\text{-}up}}\rangle\neq0}{\Downarrow}\\[4pt]
\boxed{\begin{array}{c}
\text{smooth standard embedding: full }SU(3)\text{ holonomy}\\
E_6\times E_8
\end{array}}
\end{gathered}.
\]

The resolved Kähler moduli are neutral under the final
\(E_6\times E_8\). There is no contradiction with their being charged
under \(G_{\mathrm{enh}}\) at the orbifold point: the smooth moduli are
gauge-invariant coordinates on the D-flat Higgs quotient, not the
individual charged fields used as coordinates at the enhanced-symmetry
point.

\subsubsection{Checks}\label{checks-18}

The rank count is consistent. An Abelian shift lies in a Cartan torus
and leaves all Cartan currents invariant. Hence the first \(E_8\)
retains rank eight: \[
\operatorname{rank}(E_6)
+
\operatorname{rank}(G_{\mathrm{enh}})
=6+2=8.
\] The smooth standard embedding has only rank-six \(E_6\) because the
continuous \(SU(3)\) background does not leave its own Cartan generators
covariantly constant.

For the \(\mathbb Z_3\) example, the dimensions also agree: \[
\dim(E_6\times SU(3))
=78+8
=86.
\] At the level of roots, this consists of 72 \(E_6\) roots, six
\(SU(3)\) roots, and eight Cartan generators: \[
72+6+8=86.
\]

Finally, the location of the states agrees with page 528: \[
\boxed{
\text{enhanced gauge bosons: untwisted sector},
\qquad
\text{charged blow-up fields: twisted sectors}
}.
\]

\subsubsection{Result}\label{result-22}

For a heterotic orbifold gauge shift \(\mathbf V\), the root current
\(J_{\mathbf P}\) survives when \[
\boxed{
\mathbf P\cdot\mathbf V\in\mathbb Z
},
\] so the orbifold gauge algebra is \[
\boxed{
\mathfrak g_{\mathrm{orb}}
=
\operatorname{Lie}
C_{E_8\times E_8}(\gamma_{\mathbf V})
}.
\] For the standard embedding in the first \(E_8\), \[
\mathbf V=(v_1,v_2,v_3,0,\ldots,0),
\] this contains \[
\boxed{
\mathfrak e_6
\oplus
\operatorname{Lie}C_{SU(3)}(\gamma_{\mathbf V})
\oplus
\mathfrak e_8
}.
\] The extra factor is determined by \[
\boxed{
G_{\mathrm{enh}}
=
\begin{cases}
SU(3),&
v_1=v_2=v_3\pmod1,\\
SU(2)\times U(1),&
\text{exactly two }v_i\text{ are equal mod }1,\\
U(1)^2,&
\text{all three }v_i\text{ are distinct mod }1.
\end{cases}
}
\]

The smooth standard embedding instead requires commutation with the full
gauge-bundle holonomy: \[
\boxed{
\mathfrak g_{\mathrm{smooth}}
=
\operatorname{Lie}C_{E_8}(SU(3))
=
\mathfrak e_6
}.
\] Therefore the symmetry is enhanced at the orbifold point because \[
\boxed{
\langle\gamma_{\mathbf V}\rangle
\subset SU(3)
\quad\Longrightarrow\quad
C_{E_8}(SU(3))
\subseteq
C_{E_8}(\gamma_{\mathbf V})
}.
\] The extra vectors are untwisted \(E_8\) current-algebra states.
Charged twisted blow-up fields acquire expectation values upon
resolution and Higgs \(G_{\mathrm{enh}}\), returning the gauge group to
\(E_6\times E_8\).

\section{Q3: T-duality orbit of the toroidal Type-II
orientifolds}\label{q3-t-duality-orbit-of-the-toroidal-type-ii-orientifolds}

In page 536, it is said that one can check that the orientifolds in
Table 10.3 on page 317 are related by T-dualities. For instance one can
verifies that \(T_{2}\Omega T_{2}^{-1}=\Omega I_{2}(-1)^{F_{L}}\).

Verify the T-duality relations in Table 10.3.

\begin{quote}
\textbf{Textbook location and related material.} Chapter 15, Section
15.3 around p.~536, together with Table 10.3 on p.~317 and Section 10.6:
the T-duality orbit of type-II orientifolds and the actions \(I_n\) and
\((-1)^{F_L}\).
\end{quote}

\subsection{Answer}\label{answer-23}

\subsubsection{Core point}\label{core-point-23}

Start with the type-I orientifold \[
\frac{\text{type IIB}}{\Omega}
\] on \(T^6\). Let \(\mathcal T_A\) denote T-duality along a set \(A\)
of \(n\) torus directions, and let \(I_A\) reflect those same
target-space coordinates. With the spin-field phases used by the book,
the complete operator relation on the physical closed-string Hilbert
space is \[
\boxed{
\mathcal T_A\Omega\mathcal T_A^{-1}
=
\Omega I_A
\left[(-1)^{F_L}\right]^{\frac{n(n-1)}2}
},
\qquad n=|A|.
\] Only the parity of \(\frac{n(n-1)}2\) matters. Equivalently, \[
\mathcal T_n\Omega\mathcal T_n^{-1}
=
\begin{cases}
\Omega I_n,
&
n=0,1\pmod4,\\
\Omega I_n(-1)^{F_L},
&
n=2,3\pmod4.
\end{cases}
\]

For \(n=0,1,2,3,4\), this gives exactly the five orientifold actions in
Table 10.3: \[
\begin{array}{c|c|c|c}
n
&
\text{type-II theory}
&
\mathcal T_n\Omega\mathcal T_n^{-1}
&
\text{O-plane}
\\ \hline
0&IIB&\Omega&O9\\
1&IIA&\Omega I_1&O8\\
2&IIB&\Omega I_2(-1)^{F_L}&O7\\
3&IIA&\Omega I_3(-1)^{F_L}&O6\\
4&IIB&\Omega I_4&O5
\end{array}
\]

There are three ingredients to verify:

\begin{enumerate}
\def\labelenumi{\arabic{enumi}.}
\tightlist
\item
  the bosonic worldsheet fields give \(\Omega I_n\);
\item
  the lift of \(I_n\) to Ramond spinors produces the additional
  \((-1)^{F_L}\) for \(n=2,3\pmod4\);
\item
  the fixed loci and the T-dual D-brane distribution reproduce the
  O-plane and Chan--Paton columns.
\end{enumerate}

\subsubsection{Bosonic conjugation: why a target-space reflection
appears}\label{bosonic-conjugation-why-a-target-space-reflection-appears}

For one dualized coordinate \(X^a\), decompose \[
X^a(\tau,\sigma)
=
X_L^a(\tau+\sigma)
+
X_R^a(\tau-\sigma).
\] In the convention used around equation (10.211), T-duality is the
asymmetric reflection \[
\mathcal T_a:
\qquad
X_L^a\longmapsto X_L^a,
\qquad
X_R^a\longmapsto-X_R^a.
\] Worldsheet parity exchanges the chiral halves: \[
\Omega:
\qquad
X_L^a\longleftrightarrow X_R^a.
\]

Now conjugate \(\Omega\) by \(\mathcal T_a\). Acting on \(X_L^a\) gives
\[
\begin{aligned}
X_L^a
&\xrightarrow{\ \mathcal T_a^{-1}\ }
X_L^a
\xrightarrow{\ \Omega\ }
X_R^a
\xrightarrow{\ \mathcal T_a\ }
-X_R^a,
\end{aligned}
\] while acting on \(X_R^a\) gives \[
\begin{aligned}
X_R^a
&\xrightarrow{\ \mathcal T_a^{-1}\ }
-X_R^a
\xrightarrow{\ \Omega\ }
-X_L^a
\xrightarrow{\ \mathcal T_a\ }
-X_L^a.
\end{aligned}
\] Therefore \[
\mathcal T_a\Omega\mathcal T_a^{-1}:
\qquad
(X_L^a,X_R^a)
\longmapsto
(-X_R^a,-X_L^a).
\] This is precisely worldsheet parity followed by the spacetime
reflection \[
I_a:
\qquad
(X_L^a,X_R^a)
\longmapsto
(-X_L^a,-X_R^a).
\] Hence, on bosons, \[
\boxed{
\mathcal T_a\Omega\mathcal T_a^{-1}
=
\Omega I_a
},
\] which is equation (10.211).

The calculation factorizes over different coordinates. For a set \(A\)
of \(n\) dualized coordinates, \[
\boxed{
\mathcal T_A\Omega\mathcal T_A^{-1}
=
\Omega I_A
\qquad
\text{on the bosonic and NS sectors}
}.
\] The worldsheet fermions in an NS sector transform in the vector
representation and obey the same calculation. The subtlety is entirely
in the Ramond ground states, because reflections act projectively on
spinors.

\subsubsection{Spinor lift of a
reflection}\label{spinor-lift-of-a-reflection}

Page 535 derives the action of a target-space reflection on spin fields.
If the directions in \(A\) are reflected, then \[
I_A:
\qquad
S_\alpha
\longmapsto
(P_A)_\alpha{}^\beta S_\beta,
\] where equation (15.26) writes the lift in terms of gamma matrices.
Its overall phase is fixed by the Majorana condition. It is useful to
express the same lift as \[
P_A
=
\rho_{a_1}\rho_{a_2}\cdots\rho_{a_n},
\] where each \(\rho_a\) is the properly phased lift of a
single-coordinate reflection.

The phases can be chosen, as in the book's light-cone discussion, so
that \[
\rho_a^2=1.
\] Distinct reflection matrices anticommute: \[
\rho_a\rho_b
=
-\rho_b\rho_a,
\qquad
a\neq b.
\] The anticommutation is the spinorial statement that orthogonal
reflections are represented by Clifford generators.

Squaring \(P_A\) requires moving the second copy of every \(\rho_a\)
through the matrices to its left. The number of pairwise interchanges is
\[
(n-1)+(n-2)+\cdots+1
=
\frac{n(n-1)}2.
\] Therefore \[
\boxed{
P_A^2
=
(-1)^{\frac{n(n-1)}2}
}.
\] In particular, \[
\begin{array}{c|ccccc}
n&0&1&2&3&4\\ \hline
P_A^2&+1&+1&-1&-1&+1
\end{array}
\] This reproduces the book's statement that reflection in two
directions is a rotation by \(\pi\) and squares to \(-1\) on a spinor.

Since \(P_A^2=\pm1\), \[
P_A^{-1}
=
(-1)^{\frac{n(n-1)}2}P_A.
\] This inverse is the origin of the extra fermion-number operator in
the conjugation of \(\Omega\).

\subsubsection{\texorpdfstring{Conjugating \(\Omega\) on Ramond
states}{Conjugating \textbackslash Omega on Ramond states}}\label{conjugating-omega-on-ramond-states}

Equation (15.27) says that T-duality acts by \(P_A\) on the Ramond spin
field of one chiral sector and trivially on the other. Worldsheet parity
exchanges the two chiral sectors. Consequently, in \[
\mathcal T_A\Omega\mathcal T_A^{-1},
\] one of the two spin-field transformations appears as \(P_A^{-1}\),
whereas the geometric operation \(\Omega I_A\) contains \(P_A\).

Using \[
P_A^{-1}
=
(-1)^{\frac{n(n-1)}2}P_A,
\] the two actions differ by \[
(-1)^{\frac{n(n-1)}2}
\] on the affected Ramond factor and agree on its NS factor. Tracking
the exchange of the two chiral sectors with the book's convention for
\(\Omega\), and writing the final operator in the ordering used in Table
10.3, assigns this sign to the left-moving Ramond number. It is
therefore precisely \[
\left[(-1)^{F_L}\right]^{\frac{n(n-1)}2}.
\] The label can be checked independently from \[
\Omega(-1)^{F_L}
=
(-1)^{F_R}\Omega,
\] which is the identity used by the book on page 318. Thus moving a
fermion-number operator through \(\Omega\) exchanges \(F_L\) and
\(F_R\); its placement and label must be kept together.

Thus the full operator relation is \[
\boxed{
\mathcal T_A\Omega\mathcal T_A^{-1}
=
\Omega I_A
\left[(-1)^{F_L}\right]^{\frac{n(n-1)}2}
}.
\]

The apparently innocuous bosonic identity
\(\mathcal T_A\Omega\mathcal T_A^{-1}=\Omega I_A\) therefore has to be
corrected whenever the spin lift satisfies \(P_A^2=-1\). This occurs for
\(n=2,3\pmod4\).

\subsubsection{Independent involution
check}\label{independent-involution-check}

The same factor can be derived without tracking the detailed order of
the spin matrices. On the complete closed-string Hilbert space, \[
I_A^2
=
\left[
(-1)^{F_L+F_R}
\right]^{\frac{n(n-1)}2}.
\] Indeed, \(I_A^2\) is \(P_A^2\) on each Ramond factor and is \(+1\) on
an NS factor.

Worldsheet parity exchanges left and right fermion numbers: \[
\Omega(-1)^{F_L}\Omega^{-1}
=
(-1)^{F_R}.
\] Let \[
\mathcal O_n
=
\Omega I_n
\left[(-1)^{F_L}\right]^{a_n},
\qquad
a_n=\frac{n(n-1)}2\pmod2.
\] Because \(\Omega\) commutes with the symmetric target reflection
\(I_n\), \[
\begin{aligned}
\mathcal O_n^2
&=
(\Omega I_n)^2
\left[(-1)^{F_L+F_R}\right]^{a_n}\\
&=
I_n^2
\left[(-1)^{F_L+F_R}\right]^{a_n}\\
&=
\left[(-1)^{F_L+F_R}\right]^{2a_n}\\
&=1.
\end{aligned}
\] Thus every orientifold action in the table is a genuine involution on
physical states.

For \(n=2\), this reduces exactly to the argument on page 318: \[
I_2^2=(-1)^{F_L+F_R},
\] so \[
(\Omega I_2)^2
=
(-1)^{F_L+F_R}
\neq1
\] on spacetime fermions. Inserting \((-1)^{F_L}\) gives \[
\boxed{
\left(
\Omega I_2(-1)^{F_L}
\right)^2
=1
}.
\] The \(n=3\) case has the same sign because \(P_3^2=-1\). For \(n=1\)
and \(n=4\), \(P_n^2=+1\), so no fermion-number insertion is required.

\subsubsection{Type IIA versus type IIB}\label{type-iia-versus-type-iib}

A single T-duality reverses the chirality of the Ramond ground state in
the chiral sector on which it acts. Therefore: \[
\begin{array}{c|c}
n\text{ even}
&
\text{type IIB}\longleftrightarrow\text{type IIB},
\\
n\text{ odd}
&
\text{type IIB}\longleftrightarrow\text{type IIA}.
\end{array}
\] This is equation (10.25) and the observation following equation
(15.27). It immediately explains the distribution of the rows between
the type-IIB and type-IIA blocks of Table 10.3.

\subsubsection{Verification row by row}\label{verification-row-by-row}

Choose, for definiteness, consecutive internal coordinates
\(X^9,X^8,X^7,X^6\). The cumulative T-duality chain is \[
\begin{aligned}
\frac{IIB}{\Omega}
&\xleftrightarrow{\ \mathcal T_9\ }
\frac{IIA}{\Omega I_9}\\
&\xleftrightarrow{\ \mathcal T_8\ }
\frac{IIB}{\Omega I_{89}(-1)^{F_L}}\\
&\xleftrightarrow{\ \mathcal T_7\ }
\frac{IIA}{\Omega I_{789}(-1)^{F_L}}\\
&\xleftrightarrow{\ \mathcal T_6\ }
\frac{IIB}{\Omega I_{6789}}.
\end{aligned}
\] Here each displayed quotient is understood as the theory obtained
from the original type-IIB orientifold after the cumulative set of
T-dualities. Relabeling \(I_9,I_{89},I_{789},I_{6789}\) as
\(I_1,I_2,I_3,I_4\) gives the notation in Table 10.3.

\paragraph{\texorpdfstring{Row 1: type IIB with
\(\Omega\)}{Row 1: type IIB with \textbackslash Omega}}\label{row-1-type-iib-with-omega}

For \(n=0\), \[
\frac{n(n-1)}2=0.
\] There is no geometric reflection and no fermion-number factor: \[
\mathcal O_0=\Omega.
\] The fixed locus fills all ten spacetime dimensions, so it is an
O9-plane. This is the type-I theory.

\paragraph{\texorpdfstring{Row 2: type IIA with
\(\Omega I_1\)}{Row 2: type IIA with \textbackslash Omega I\_1}}\label{row-2-type-iia-with-omega-i_1}

For \(n=1\), \[
\frac{n(n-1)}2=0,
\qquad
P_1^2=+1.
\] Thus \[
\boxed{
\mathcal T_1\Omega\mathcal T_1^{-1}
=
\Omega I_1
}.
\] The odd number of T-dualities converts IIB to IIA. The reflection of
one circle has two fixed points, so the O9-plane becomes two O8-planes.

In light-cone gauge, the book chooses the spin lift of the single
reflection with the Majorana-compatible phase, for example \[
I_1=i\Gamma^9\Gamma
\] on spinors. With this phase, \[
(\Omega I_1)^2=1,
\] as stated on page 536.

\paragraph{\texorpdfstring{Row 3: type IIB with
\(\Omega I_2(-1)^{F_L}\)}{Row 3: type IIB with \textbackslash Omega I\_2(-1)\^{}\{F\_L\}}}\label{row-3-type-iib-with-omega-i_2-1f_l}

For \(n=2\), \[
\frac{n(n-1)}2=1,
\qquad
P_2^2=-1.
\] Therefore \[
\boxed{
\mathcal T_2\Omega\mathcal T_2^{-1}
=
\Omega I_2(-1)^{F_L}
}.
\] For example, with the directions \(8,9\), \[
\boxed{
\mathcal T_{89}\Omega\mathcal T_{89}^{-1}
=
\Omega I_{89}(-1)^{F_L}
},
\] which is the relation quoted on pages 318 and 536.

The even number of dualities leaves the theory type IIB. Reflection of
two circles has \[
2^2=4
\] fixed loci, so there are four O7-planes.

\paragraph{\texorpdfstring{Row 4: type IIA with
\(\Omega I_3(-1)^{F_L}\)}{Row 4: type IIA with \textbackslash Omega I\_3(-1)\^{}\{F\_L\}}}\label{row-4-type-iia-with-omega-i_3-1f_l}

For \(n=3\), \[
\frac{n(n-1)}2=3,
\qquad
P_3^2=-1.
\] Hence \[
\boxed{
\mathcal T_3\Omega\mathcal T_3^{-1}
=
\Omega I_3(-1)^{F_L}
}.
\] Because \(n\) is odd, the dual theory is type IIA. The reflection has
\[
2^3=8
\] fixed loci, giving eight O6-planes.

\paragraph{\texorpdfstring{Row 5: type IIB with
\(\Omega I_4\)}{Row 5: type IIB with \textbackslash Omega I\_4}}\label{row-5-type-iib-with-omega-i_4}

For \(n=4\), \[
\frac{n(n-1)}2=6,
\qquad
P_4^2=+1.
\] The fermion-number factor occurs to an even power and drops out: \[
\boxed{
\mathcal T_4\Omega\mathcal T_4^{-1}
=
\Omega I_4
}.
\] Four T-dualities return to type IIB. The reflection has \[
2^4=16
\] fixed loci, so there are sixteen O5-planes.

\subsubsection{O-plane dimensions and
multiplicities}\label{o-plane-dimensions-and-multiplicities}

T-duality along a direction parallel to a D\(p\)-brane or O\(p\)-plane
lowers its dimension: \[
D_p\longmapsto D_{p-1},
\qquad
O_p\longmapsto O_{p-1}.
\] Starting from the type-I O9-plane and dualizing \(n\) internal
directions therefore gives \[
p=9-n.
\]

On each reflected circle, the equation \[
x^a=-x^a\pmod{2\pi R_a}
\] has the two solutions \[
x^a=0,
\qquad
x^a=\pi R_a.
\] Thus \(I_n\) has \[
\boxed{
2^n
}
\] fixed components, each supporting an \(O_{9-n}\)-plane. This gives \[
\begin{array}{c|ccccc}
n&0&1&2&3&4\\ \hline
\text{number of O-planes}&1&2&4&8&16\\
\text{type}&O9&O8&O7&O6&O5
\end{array}
\] and reproduces the third column of Table 10.3.

\subsubsection{Chan--Paton groups and local tadpole
cancellation}\label{chanpaton-groups-and-local-tadpole-cancellation}

T-duality preserves the total number of D-branes in the upstairs count:
the 32 type-I D9-branes become 32 D\((9-n)\)-branes. Equation (10.212)
gives the charge of a standard orientifold plane: \[
\mu_{O_p}
=
-2^{p-4}\mu_{D_p}.
\] For \(p=9-n\), \[
\mu_{O_{9-n}}
=
-2^{5-n}\mu_{D_{9-n}}.
\] Therefore local RR-charge cancellation at each fixed plane requires
\[
\boxed{
N_{\mathrm{local}}
=
2^{5-n}
=
\frac{32}{2^n}
}
\] coincident D\((9-n)\)-branes on each O\((9-n)\)-plane.

The T-duals of the type-I \(O9^-\) plane impose the orthogonal
Chan--Paton projection. A stack of \(N_{\mathrm{local}}\) branes on each
such plane gives \(SO(N_{\mathrm{local}})\). Since there are \(2^n\)
fixed planes, \[
\boxed{
G_{\mathrm{CP}}
=
SO\!\left(2^{5-n}\right)^{2^n}
}.
\] Evaluating this formula gives \[
\begin{array}{c|c|c|c}
n
&
\text{O-planes}
&
\text{D-branes per O-plane}
&
G_{\mathrm{CP}}
\\ \hline
0&1\times O9&32&SO(32)\\
1&2\times O8&16&SO(16)^2\\
2&4\times O7&8&SO(8)^4\\
3&8\times O6&4&SO(4)^8\\
4&16\times O5&2&SO(2)^{16}
\end{array}
\] This is the final column of Table 10.3.

As a charge check, \[
2^n\mu_{O_{9-n}}
=
2^n
\left(
-2^{5-n}\mu_{D_{9-n}}
\right)
=
-32\mu_{D_{9-n}},
\] which is canceled by the 32 T-dual D-branes for every \(n\).

\subsubsection{Check on the massless closed-string
sectors}\label{check-on-the-massless-closed-string-sectors}

The conjugation can also be seen directly on massless states.

An NS--NS state has the schematic form \[
\psi_{-1/2}^M
\widetilde\psi_{-1/2}^N
|0;k\rangle_{\mathrm{NS\text{-}NS}}.
\] Worldsheet parity exchanges \(M\) and \(N\). Conjugation by
\(\mathcal T_A\) supplies a minus sign whenever an index lies in a
dualized direction, so the polarization transforms as \[
E_{MN}
\longmapsto
(R_A)_M{}^P
(R_A)_N{}^Q
E_{QP},
\] which is exactly \(\Omega I_A\). Since \((-1)^{F_L}=+1\) in the left
NS sector, the extra factor is invisible here.

For R--R and mixed R--NS/NS--R states, \(\mathcal T_A\) acts on one spin
field by \(P_A\). Exchanging the two chiral sectors with \(\Omega\)
reverses the order of the spin lift and produces \(P_A^{-1}\). The
relation \[
P_A^{-1}
=
(-1)^{\frac{n(n-1)}2}P_A
\] then gives precisely the factor recorded by \((-1)^{F_L}\). This
ensures that the transformed orientifold action has invariant fermions
and preserves the same sixteen supercharges as type I.

\subsubsection{Checks}\label{checks-19}

The alternation between IIA and IIB is correct: \[
n=0,2,4
\quad\Longrightarrow\quad
IIB,
\qquad
n=1,3
\quad\Longrightarrow\quad
IIA.
\]

The fermion-number insertion has the independent pattern \[
\frac{n(n-1)}2
=
0,0,1,3,6
\quad
\text{for }
n=0,1,2,3,4.
\] Thus it occurs for \(n=2,3\), not simply for odd or even \(n\). These
are exactly the two rows in Table 10.3 containing \((-1)^{F_L}\).

The O-plane dimension and number obey \[
p=9-n,
\qquad
N_{O_p}=2^n.
\] The total O-plane charge is independent of \(n\): \[
N_{O_p}\mu_{O_p}
=
-32\mu_{D_p},
\] as required by T-duality.

\subsubsection{Result}\label{result-23}

For T-duality along a set \(A\) of \(n\) compact coordinates, \[
\boxed{
\mathcal T_A:
\quad
X_L^a\mapsto X_L^a,
\qquad
X_R^a\mapsto-X_R^a
},
\] and therefore the bosonic conjugation is \[
\boxed{
\mathcal T_A\Omega\mathcal T_A^{-1}
=
\Omega I_A
\quad\text{on NS states}.
}
\]

The spin lift satisfies \[
\boxed{
P_A^2
=
(-1)^{\frac{n(n-1)}2}
},
\] so the complete orientifold action is \[
\boxed{
\mathcal T_A\Omega\mathcal T_A^{-1}
=
\Omega I_A
\left[(-1)^{F_L}\right]^{\frac{n(n-1)}2}
}.
\]

Consequently, \[
\boxed{
\begin{array}{c|c|c|c|c}
n
&
\text{theory}
&
\text{orientifold}
&
\text{O-planes}
&
\text{local CP group}
\\ \hline
0&IIB&\Omega&O9&SO(32)\\
1&IIA&\Omega I_1&2\,O8&SO(16)^2\\
2&IIB&\Omega I_2(-1)^{F_L}&4\,O7&SO(8)^4\\
3&IIA&\Omega I_3(-1)^{F_L}&8\,O6&SO(4)^8\\
4&IIB&\Omega I_4&16\,O5&SO(2)^{16}
\end{array}
}.
\] This verifies every T-duality relation and every associated entry in
Table 10.3.

\section{\texorpdfstring{Q4: Closed-string spectrum of the
\(T^4/\mathbb Z_2\)
orientifold}{Q4: Closed-string spectrum of the T\^{}4/\textbackslash mathbb Z\_2 orientifold}}\label{q4-closed-string-spectrum-of-the-t4mathbb-z_2-orientifold}

In page 536-537, they obtained the massless closed string states for
this \(T^4/\mathbb{Z}_{2}\) orientifold construction. The detailed
computation is omitted. Following their method, compute the resulting
\(SU(2)\times SU(2)\) representations.

\begin{quote}
\textbf{Textbook location and related material.} Chapter 15, Section
15.3 (pp.~536--537) and Table 15.3: the \(T^4/\mathbb Z_2\) chiral
building blocks, orbifold/orientifold projections, and assembly into
six-dimensional multiplets.
\end{quote}

\subsection{Answer}\label{answer-24}

\subsubsection{Core point}\label{core-point-24}

The computation consists of two successive projections.

First construct the type-IIB \(T^4/\mathbb Z_2\) orbifold spectrum by
tensoring a left-moving state and a right-moving state from Table 15.3
and retaining only equal \(I_4\) parities: \[
q_{I_4}^{(L)}q_{I_4}^{(R)}=+1.
\]

Then project the surviving closed-string states with \[
\mathcal O
=
\Omega R(-1)^{F_L},
\] where \[
R:
\qquad
y^1\mapsto-y^1,
\qquad
y^2\mapsto-y^2.
\] In the twisted sectors, the model additionally assigns the twist
field the sign \[
\epsilon_{\mathrm{tw}}=-1.
\]

Carrying out these steps gives \[
\boxed{
\begin{array}{c|c|l}
&\text{sector}&SU(2)_L\times SU(2)_R\text{ content}\\ \hline
\text{untwisted}
&NS\text{--}NS
&(\mathbf3,\mathbf3)+11(\mathbf1,\mathbf1)\\
&
R\text{--}R
&(\mathbf3,\mathbf1)+(\mathbf1,\mathbf3)+6(\mathbf1,\mathbf1)\\
&
NS\text{--}R+R\text{--}NS
&2(\mathbf2,\mathbf3)+10(\mathbf2,\mathbf1)\\ \hline
\text{twisted}
&NS\text{--}NS
&48(\mathbf1,\mathbf1)\\
&
R\text{--}R
&16(\mathbf1,\mathbf1)\\
&
NS\text{--}R+R\text{--}NS
&32(\mathbf2,\mathbf1)
\end{array}
}.
\]

The nontrivial multiplicities follow from how \(R\) acts on the internal
degeneracy spaces and from the fact that \(\Omega\) exchanges the two
chiral factors. We now derive them one by one.

\subsubsection{Chiral building blocks from Table
15.3}\label{chiral-building-blocks-from-table-15.3}

For one chiral half of the untwisted NS sector, Table 15.3 gives \[
\begin{aligned}
V&=(\mathbf2,\mathbf2)_{+},\\
Y&=4(\mathbf1,\mathbf1)_{-},
\end{aligned}
\] where the subscript is the \(I_4\) eigenvalue.

The block \(V\) comes from an oscillator in one of the four transverse
noncompact directions. The block \(Y\) comes from the four internal
oscillators, one for each real coordinate \[
x^1,\ y^1,\ x^2,\ y^2.
\] Under \(R\), \[
\begin{array}{c|cccc}
\text{internal oscillator}&x^1&y^1&x^2&y^2\\ \hline
R\text{ eigenvalue}&+1&-1&+1&-1
\end{array}
\] so \(Y\) has two \(R\)-even and two \(R\)-odd basis states.

For one chiral half of the untwisted Ramond sector, define \[
\begin{aligned}
A&=2(\mathbf2,\mathbf1)_{-},\\
B&=2(\mathbf1,\mathbf2)_{+}.
\end{aligned}
\] The factor \(2\) in each block is an internal two-state degeneracy.

Equation (15.29) gives \[
R:
\quad
|s_1,s_2;s_3,s_4\rangle
\longmapsto
(-1)^{\frac12+s_4}
|s_1,s_2;-s_3,-s_4\rangle.
\] For either \(A\) or \(B\), \(R\) exchanges the two internally
degenerate states with a relative minus sign. In a suitable basis, \[
J
=
\begin{pmatrix}
0&1\\
-1&0
\end{pmatrix},
\qquad
J^2=-\mathbf1_2.
\] The precise sign of \(J\) can be changed by rephasing the basis;
\(J^2=-1\) is invariant and is what enters the projection.

At each of the 16 \(I_4\) fixed points, the twisted chiral blocks are \[
\begin{aligned}
C&=2(\mathbf1,\mathbf1)_{+}
\qquad &&\text{in the NS sector},\\
D&=(\mathbf2,\mathbf1)_{+}
\qquad &&\text{in the Ramond sector}.
\end{aligned}
\] The reflection acts on the two-dimensional degeneracy of \(C\) by
another matrix \(J\) with \(J^2=-1\). The internal part of \(D\) is
nondegenerate. Its geometric \(R\) phase and the phase assigned to the
twist field enter the orientifold action only through their product;
this model fixes that product by choosing the additional twisted-sector
sign \(\epsilon_{\mathrm{tw}}=-1\) on page 536.

\subsubsection{\texorpdfstring{Useful \(SU(2)\) tensor
products}{Useful SU(2) tensor products}}\label{useful-su2-tensor-products}

The only nontrivial little-group identity needed is \[
\mathbf2\otimes\mathbf2
=
\operatorname{Sym}^2\mathbf2
\oplus
\bigwedge^2\mathbf2
=
\mathbf3\oplus\mathbf1.
\] Consequently, \[
(\mathbf2,\mathbf2)\otimes(\mathbf2,\mathbf2)
=
(\mathbf3,\mathbf3)
\oplus
(\mathbf3,\mathbf1)
\oplus
(\mathbf1,\mathbf3)
\oplus
(\mathbf1,\mathbf1).
\] The symmetric and antisymmetric squares are \[
\begin{aligned}
\operatorname{Sym}^2(\mathbf2,\mathbf2)
&=
(\mathbf3,\mathbf3)
\oplus
(\mathbf1,\mathbf1),\\
\bigwedge^2(\mathbf2,\mathbf2)
&=
(\mathbf3,\mathbf1)
\oplus
(\mathbf1,\mathbf3).
\end{aligned}
\]

We will also repeatedly use the exchange operator \[
\tau(u\otimes v)=v\otimes u.
\] For a finite-dimensional space \(W\) and an operator \(M\) on \(W\),
\[
\operatorname{Tr}_{W\otimes W}
\left[
\tau(M\otimes M)
\right]
=
\operatorname{Tr}_W(M^2).
\] If \(\tau(M\otimes M)\) squares to one, the dimensions of its two
eigenspaces are therefore \[
d_\pm
=
\frac12
\left[
(\dim W)^2
\pm
\operatorname{Tr}(M^2)
\right].
\]

\subsubsection{Signs in the orientifold
projection}\label{signs-in-the-orientifold-projection}

The book uses the standard action of \(\Omega\):

\begin{itemize}
\tightlist
\item
  in NS--NS, \(\Omega\) symmetrizes the two chiral factors;
\item
  in R--R, \(\Omega\) antisymmetrizes them.
\end{itemize}

Meanwhile, \[
(-1)^{F_L}
=
\begin{cases}
+1,&\text{left NS},\\
-1,&\text{left R}.
\end{cases}
\] Therefore, in the untwisted diagonal sectors, the two minus signs in
R--R cancel: \[
\begin{array}{c|c}
\text{sector}&\mathcal O\text{ on identical chiral blocks}\\ \hline
NS\text{--}NS&+\tau(R\otimes R)\\
R\text{--}R&(-\tau)(R\otimes R)(-1)
=+\tau(R\otimes R)
\end{array}
\] Thus both untwisted diagonal sectors require the \(+1\) eigenspace of
\[
K_R=\tau(R\otimes R).
\]

In a twisted sector the chosen extra sign reverses the projection: \[
\mathcal O_{\mathrm{tw}}
=
-K_R
\] in both NS--NS and R--R. Hence an \(\mathcal O_{\mathrm{tw}}\)-even
state is a \(K_R\)-odd state.

In the mixed NS--R and R--NS sectors, \(\Omega\) exchanges two different
sectors. It pairs every state in NS--R with one state in R--NS, and
precisely one linear combination survives. Therefore the orientifold
projection divides the doubled mixed-sector count by two but does not
otherwise change its little-group representation content.

\subsubsection{Untwisted NS--NS sector}\label{untwisted-nsns-sector}

The tensor product of the two chiral NS spectra is \[
(V\oplus Y)_L
\otimes
(V\oplus Y)_R.
\] The \(I_4\) projection retains equal parities: \[
\boxed{
\mathcal H_{\mathrm{untw}}^{NS\text{--}NS}
=
(V\otimes V)
\oplus
(Y\otimes Y)
}.
\] The cross terms \(V\otimes Y\) and \(Y\otimes V\) have total \(I_4\)
eigenvalue \(-1\) and are removed.

\paragraph{\texorpdfstring{The \(V\otimes V\)
contribution}{The V\textbackslash otimes V contribution}}\label{the-votimes-v-contribution}

The reflection \(R\) acts trivially on \(V\), because \(V\) consists of
noncompact oscillators. Hence the orientifold projection is ordinary
symmetrization: \[
(V\otimes V)_{\mathcal O=+1}
=
\operatorname{Sym}^2V.
\] Using the tensor-product identity above, \[
\boxed{
\operatorname{Sym}^2(\mathbf2,\mathbf2)
=
(\mathbf3,\mathbf3)
\oplus
(\mathbf1,\mathbf1)
}.
\]

\paragraph{\texorpdfstring{The \(Y\otimes Y\)
contribution}{The Y\textbackslash otimes Y contribution}}\label{the-yotimes-y-contribution}

The space \(Y\) has dimension four and \(R_Y^2=1\). The orientifold
operator on \(Y\otimes Y\) is \[
K_Y
=
\tau(R_Y\otimes R_Y).
\] Its trace is \[
\operatorname{Tr}K_Y
=
\operatorname{Tr}(R_Y^2)
=4.
\] Therefore its invariant subspace has dimension \[
\dim(Y\otimes Y)_{K_Y=+1}
=
\frac12(4^2+4)
=10.
\] All these states are little-group singlets: \[
(Y\otimes Y)_{\mathcal O=+1}
=
10(\mathbf1,\mathbf1).
\]

One can display the ten states explicitly. If \(e_i\) is an \(R\)
eigenbasis with eigenvalues \(r_i=\pm1\), the invariant tensors are \[
e_i\otimes e_i,
\qquad
e_i\otimes e_j
+
r_ir_j\,e_j\otimes e_i,
\qquad
i<j.
\] There are four diagonal tensors and six off-diagonal combinations.

Combining \(V\otimes V\) and \(Y\otimes Y\) gives \[
\boxed{
\mathcal H_{\mathrm{untw}}^{NS\text{--}NS}
=
(\mathbf3,\mathbf3)
+
11(\mathbf1,\mathbf1)
}.
\]

\subsubsection{Untwisted R--R sector}\label{untwisted-rr-sector}

The chiral R spectrum is \(A\oplus B\), with opposite \(I_4\) parities.
Thus \[
\boxed{
\mathcal H_{\mathrm{untw}}^{R\text{--}R}
=
(A\otimes A)
\oplus
(B\otimes B)
}.
\] The cross terms \(A\otimes B\) and \(B\otimes A\) are orbifold odd.

Consider \(A\otimes A\). Write \[
A
=
(\mathbf2,\mathbf1)\otimes U_A,
\qquad
\dim U_A=2,
\] where \(R\) acts as \(J\) on \(U_A\). On \(U_A\otimes U_A\), define
\[
K_U
=
\tau_U(J\otimes J).
\] Since \(J^2=-1\), \[
\operatorname{Tr}K_U
=
\operatorname{Tr}(J^2)
=-2.
\] The \(K_U\) eigenspaces therefore have dimensions \[
\dim U_{+}
=
\frac12(4-2)
=1,
\qquad
\dim U_{-}
=
\frac12(4+2)
=3.
\]

The exchange of the complete \(A\) factors is \[
\tau_A
=
\tau_{\mathbf2}\otimes\tau_U.
\] Hence the orientifold operator factorizes as \[
\tau_A(R\otimes R)
=
\tau_{\mathbf2}\otimes K_U.
\] An invariant state can arise in two ways:

\begin{enumerate}
\def\labelenumi{\arabic{enumi}.}
\tightlist
\item
  the external doublet product is symmetric and the internal factor is
  \(K_U\)-even;
\item
  the external doublet product is antisymmetric and the internal factor
  is \(K_U\)-odd.
\end{enumerate}

Since \[
\operatorname{Sym}^2\mathbf2=\mathbf3,
\qquad
\bigwedge^2\mathbf2=\mathbf1,
\] the first possibility gives one triplet, while the second gives three
singlets: \[
\boxed{
(A\otimes A)_{\mathcal O=+1}
=
(\mathbf3,\mathbf1)
+
3(\mathbf1,\mathbf1)
}.
\]

Exactly the same calculation for \[
B
=
(\mathbf1,\mathbf2)\otimes U_B
\] gives \[
\boxed{
(B\otimes B)_{\mathcal O=+1}
=
(\mathbf1,\mathbf3)
+
3(\mathbf1,\mathbf1)
}.
\]

Therefore \[
\boxed{
\mathcal H_{\mathrm{untw}}^{R\text{--}R}
=
(\mathbf3,\mathbf1)
+
(\mathbf1,\mathbf3)
+
6(\mathbf1,\mathbf1)
}.
\]

\subsubsection{Untwisted mixed sectors}\label{untwisted-mixed-sectors}

Consider first NS--R. Equal \(I_4\) parity allows \[
V\otimes B
\qquad\text{and}\qquad
Y\otimes A.
\]

For the first term, \[
\begin{aligned}
V\otimes B
&=
(\mathbf2,\mathbf2)
\otimes
2(\mathbf1,\mathbf2)\\
&=
2(\mathbf2,\mathbf2\otimes\mathbf2)\\
&=
2(\mathbf2,\mathbf3)
+
2(\mathbf2,\mathbf1).
\end{aligned}
\] For the second, \[
\begin{aligned}
Y\otimes A
&=
4(\mathbf1,\mathbf1)
\otimes
2(\mathbf2,\mathbf1)\\
&=
8(\mathbf2,\mathbf1).
\end{aligned}
\] Thus one orientation has \[
\mathcal H_{\mathrm{untw}}^{NS\text{--}R}
=
2(\mathbf2,\mathbf3)
+
10(\mathbf2,\mathbf1).
\]

The R--NS sector has an isomorphic set of states. The orientifold
exchanges NS--R with R--NS and retains one linear combination from each
pair. It therefore removes the overall doubling: \[
\boxed{
\mathcal H_{\mathrm{untw}}^{NS\text{--}R+R\text{--}NS}
=
2(\mathbf2,\mathbf3)
+
10(\mathbf2,\mathbf1)
}.
\]

\subsubsection{Twisted sectors at one fixed
point}\label{twisted-sectors-at-one-fixed-point}

There are 16 \(I_4\) fixed points. For the factorized rectangular torus
used by the book, \(R\) maps each of them to itself, because each fixed
coordinate is either \(0\) or a half-period. We may therefore compute
the spectrum at one fixed point and multiply by 16.

The chiral twisted blocks are \[
C=2(\mathbf1,\mathbf1),
\qquad
D=(\mathbf2,\mathbf1).
\] They are both \(I_4\) even, so all diagonal and mixed tensor products
relevant below survive the orbifold projection.

\subsubsection{Twisted NS--NS sector}\label{twisted-nsns-sector}

Before the orientifold projection, \[
C\otimes C
=
4(\mathbf1,\mathbf1).
\] On the two-dimensional degeneracy space of \(C\), \(R\) acts by \(J\)
with \(J^2=-1\). Thus \[
K_C
=
\tau(J\otimes J)
\] has \[
\operatorname{Tr}K_C
=
\operatorname{Tr}J^2
=-2.
\] Its \(+1\) and \(-1\) eigenspaces have dimensions one and three: \[
\dim(K_C=+1)=1,
\qquad
\dim(K_C=-1)=3.
\]

Because the chosen twisted-sector sign is \(\epsilon_{\mathrm{tw}}=-1\),
the full orientifold operator is \[
\mathcal O_{\mathrm{tw}}
=
-K_C.
\] An orientifold-even state must therefore lie in the \(K_C=-1\)
eigenspace. Hence one fixed point contributes \[
\boxed{
3(\mathbf1,\mathbf1)
}
\] in NS--NS. Multiplying by 16 gives \[
\boxed{
\mathcal H_{\mathrm{tw}}^{NS\text{--}NS}
=
48(\mathbf1,\mathbf1)
}.
\]

\subsubsection{Twisted R--R sector}\label{twisted-rr-sector}

At one fixed point, \[
D\otimes D
=
(\mathbf2,\mathbf1)\otimes(\mathbf2,\mathbf1)
=
(\mathbf3,\mathbf1)
\oplus
(\mathbf1,\mathbf1).
\] The symmetric square is \((\mathbf3,\mathbf1)\) and the antisymmetric
square is \((\mathbf1,\mathbf1)\).

After including the standard R--R sign of \(\Omega\), the \((-1)^{F_L}\)
sign, and the additional twisted-sector sign, the effective projection
is \[
\mathcal O_{\mathrm{tw}}
=
-\tau
\] on \(D\otimes D\). Thus the symmetric triplet is odd and the
antisymmetric singlet is even: \[
\boxed{
(D\otimes D)_{\mathcal O=+1}
=
(\mathbf1,\mathbf1)
}.
\] With 16 fixed points, \[
\boxed{
\mathcal H_{\mathrm{tw}}^{R\text{--}R}
=
16(\mathbf1,\mathbf1)
}.
\]

\subsubsection{Twisted mixed sectors}\label{twisted-mixed-sectors}

For one orientation and one fixed point, \[
C\otimes D
=
2(\mathbf1,\mathbf1)
\otimes
(\mathbf2,\mathbf1)
=
2(\mathbf2,\mathbf1).
\] There is an isomorphic \(D\otimes C\) contribution in the opposite
R--NS sector. The orientifold exchanges the two orientations and leaves
one linear combination, so one fixed point contributes \[
2(\mathbf2,\mathbf1).
\] Multiplying by 16, \[
\boxed{
\mathcal H_{\mathrm{tw}}^{NS\text{--}R+R\text{--}NS}
=
32(\mathbf2,\mathbf1)
}.
\]

\subsubsection{\texorpdfstring{Assembly into six-dimensional
\(\mathcal N=(1,0)\)
multiplets}{Assembly into six-dimensional \textbackslash mathcal N=(1,0) multiplets}}\label{assembly-into-six-dimensional-mathcal-n10-multiplets}

The total untwisted bosonic content is \[
\begin{aligned}
&\left[
(\mathbf3,\mathbf3)
+
11(\mathbf1,\mathbf1)
\right]_{NS\text{--}NS}\\
&\qquad
\oplus
\left[
(\mathbf3,\mathbf1)
+
(\mathbf1,\mathbf3)
+
6(\mathbf1,\mathbf1)
\right]_{R\text{--}R}\\
&=
(\mathbf3,\mathbf3)
+
(\mathbf3,\mathbf1)
+
(\mathbf1,\mathbf3)
+
17(\mathbf1,\mathbf1).
\end{aligned}
\] The untwisted fermions are \[
2(\mathbf2,\mathbf3)
+
10(\mathbf2,\mathbf1).
\]

Using Table 15.9:

\begin{itemize}
\tightlist
\item
  one gravity multiplet contains \[
  (\mathbf3,\mathbf3)
  +2(\mathbf2,\mathbf3)
  +(\mathbf1,\mathbf3);
  \]
\item
  one tensor multiplet contains \[
  (\mathbf3,\mathbf1)
  +2(\mathbf2,\mathbf1)
  +(\mathbf1,\mathbf1).
  \]
\end{itemize}

After subtracting these, the remaining untwisted states are \[
16(\mathbf1,\mathbf1)
+
8(\mathbf2,\mathbf1).
\] One hypermultiplet contains four real singlet bosonic degrees of
freedom and two \((\mathbf2,\mathbf1)\) fermionic degrees of freedom.
The remainder is therefore four hypermultiplets. Thus \[
\boxed{
\text{untwisted sector}
=
\text{gravity}
+
\text{tensor}
+
4\ \text{hypermultiplets}
}.
\]

At each fixed point the twisted states are \[
\underbrace{
3(\mathbf1,\mathbf1)
+
(\mathbf1,\mathbf1)
}_{4\text{ bosonic singlets}}
+
\underbrace{
2(\mathbf2,\mathbf1)
}_{\text{fermions}},
\] which is exactly one hypermultiplet. Therefore \[
\boxed{
\text{twisted sector}
=
16\ \text{hypermultiplets}
}.
\] This reproduces equation (15.30).

\subsubsection{Checks}\label{checks-20}

The untwisted bosonic degree count is \[
\begin{aligned}
N_B^{\mathrm{untw}}
&=
\left(9+11\right)
+
\left(3+3+6\right)\\
&=20+12\\
&=32.
\end{aligned}
\] The fermionic count is \[
N_F^{\mathrm{untw}}
=
2(2\cdot3)
+
10(2\cdot1)
=12+20
=32.
\]

In the twisted sector, \[
N_B^{\mathrm{tw}}
=48+16
=64,
\] and \[
N_F^{\mathrm{tw}}
=32(2\cdot1)
=64.
\] Thus bosonic and fermionic physical degrees of freedom match
separately in the untwisted and twisted sectors.

The role of the twisted-sector sign can also be checked. If one chose \[
\epsilon_{\mathrm{tw}}=+1
\] instead, the twisted NS--NS sector would retain the one-dimensional
\(K_C=+1\) space, while twisted R--R would retain the symmetric
\((\mathbf3,\mathbf1)\). At each fixed point the result would be \[
(\mathbf3,\mathbf1)
+
2(\mathbf2,\mathbf1)
+
(\mathbf1,\mathbf1),
\] which is a tensor multiplet rather than a hypermultiplet, precisely
as stated by the book.

\subsubsection{Result}\label{result-24}

The chiral blocks and orbifold parities are \[
\begin{array}{c|cc}
&\text{block}&I_4\\ \hline
\text{untwisted NS}
&V=(\mathbf2,\mathbf2),\quad Y=4(\mathbf1,\mathbf1)
&+,\ -\\
\text{untwisted R}
&A=2(\mathbf2,\mathbf1),\quad B=2(\mathbf1,\mathbf2)
&-,\ +\\
\text{twisted NS}
&C=2(\mathbf1,\mathbf1)
&+\\
\text{twisted R}
&D=(\mathbf2,\mathbf1)
&+
\end{array}
\] Equal \(I_4\) parities are tensored, after which one imposes \[
\mathcal O
=
\Omega R(-1)^{F_L},
\qquad
\epsilon_{\mathrm{tw}}=-1.
\] The result is \[
\boxed{
\begin{aligned}
\text{untwisted }NS\text{--}NS
&:
(\mathbf3,\mathbf3)
+
11(\mathbf1,\mathbf1),\\
\text{untwisted }R\text{--}R
&:
(\mathbf3,\mathbf1)
+
(\mathbf1,\mathbf3)
+
6(\mathbf1,\mathbf1),\\
\text{untwisted mixed}
&:
2(\mathbf2,\mathbf3)
+
10(\mathbf2,\mathbf1),\\
\text{twisted }NS\text{--}NS
&:
48(\mathbf1,\mathbf1),\\
\text{twisted }R\text{--}R
&:
16(\mathbf1,\mathbf1),\\
\text{twisted mixed}
&:
32(\mathbf2,\mathbf1).
\end{aligned}
}
\] These representations assemble as \[
\boxed{
\text{untwisted: gravity}+\text{tensor}+4\text{ hypers},
\qquad
\text{twisted: }16\text{ hypers}.
}
\]

\section{\texorpdfstring{Q5: Why the orientifold contains both \(x\)-
and \(y\)-oriented
O7-planes}{Q5: Why the orientifold contains both x- and y-oriented O7-planes}}\label{q5-why-the-orientifold-contains-both-x--and-y-oriented-o7-planes}

In page 538-539, the prefactor in front of the massless tadpole
indicates that it corresponds to an O7-plane, one occupies the six
noncompact dimensions and two x-axes, the other occupies the six
noncompact dimensions and two y-axes. I agree with this conclusion, but
I also wander why the position of the O7-brane is like this. We have the
orientifold action \(\Omega R(-1)^{F_{L}}\) with
\(R:y_{i}\to-y_{i},i=1,2\); so naively the orientifold plane may locate
at the fixed locus of \(R\) i.e.~it occupies the six noncompact
dimensions and two x-axes? What about the other O7-plane? Or how to
formally define the position of the O-plane?

\begin{quote}
\textbf{Textbook location and related material.} Chapter 15, Section
15.3 (pp.~538--539): Klein-bottle lattice sums and the two O7-plane
families. Sections 6.7 and 10.6 provide the general crosscap and
orientifold fixed-locus interpretation.
\end{quote}

\subsection{Answer}\label{answer-25}

\subsubsection{Core point}\label{core-point-25}

Your fixed-locus argument is correct for the single group element \[
\mathcal O=\Omega R(-1)^{F_L}.
\] Its geometric part is \(R\), and \(\operatorname{Fix}(R)\) consists
of loci extended along \(x^1,x^2\) and localized in \(y^1,y^2\). These
are the \(O7_x\)-planes.

The missing point is that the background has already been divided by the
orbifold generator \(I_4\). The orientifold projection must therefore be
imposed on the \(T^4/\mathbb Z_2\) theory, not on the covering torus by
itself. Closure of the full orientifold group gives two
orientation-reversing elements: \[
\mathcal O
=
\Omega R(-1)^{F_L},
\qquad
\mathcal O I_4
=
\Omega RI_4(-1)^{F_L}.
\] The geometric part of the second element is \(RI_4\). Since \[
R:(x^i,y^i)\mapsto(x^i,-y^i),
\qquad
I_4:(x^i,y^i)\mapsto(-x^i,-y^i),
\] their product acts as \[
RI_4:
\qquad
(x^i,y^i)\mapsto(-x^i,y^i).
\] Its fixed locus is localized in \(x^1,x^2\) and extended along
\(y^1,y^2\). These are the \(O7_y\)-planes.

Thus the second family is not an extra object inserted by hand. It is
required by the \(I_4\) term in the orbifold projector in equation
(15.31).

\subsubsection{The full orientifold group and the two
crosscaps}\label{the-full-orientifold-group-and-the-two-crosscaps}

The full group is \[
\Gamma_{\mathrm{or}}
=
\left\{
1,\,
I_4,\,
\Omega R(-1)^{F_L},\,
\Omega RI_4(-1)^{F_L}
\right\}.
\] The first two elements preserve worldsheet orientation. The last two
reverse it and therefore produce crosscap sectors.

Indeed, equation (15.31) contains \[
\begin{aligned}
\mathcal K
&=
\int_0^\infty\frac{dt}{2t}\,
\operatorname{Tr}
\left[
\Omega R(-1)^{F_L}
\frac{1+I_4}{2}
P_{\mathrm{GSO}}
e^{-4\pi t(L_0-c/24)}
\right]\\
&=
\frac12\mathcal K_{\mathcal O}
+
\frac12\mathcal K_{\mathcal O I_4}.
\end{aligned}
\] The first term is a trace with insertion \(\mathcal O\), whereas the
second is a trace with insertion \(\mathcal OI_4\). After transforming
the Klein bottle to the tree channel, these become closed-string
propagation between the corresponding crosscap states: \[
\widetilde{\mathcal K}_{\mathcal O}
\sim
\langle C_R|e^{-2\pi lH_{\mathrm{cl}}}|C_R\rangle,
\qquad
\widetilde{\mathcal K}_{\mathcal OI_4}
\sim
\langle C_{RI_4}|e^{-2\pi lH_{\mathrm{cl}}}|C_{RI_4}\rangle.
\] The target-space loci of these crosscaps are what are called the
orientifold planes.

The factors \(\Omega\) and \((-1)^{F_L}\) do not provide additional
target-space coordinate transformations. The former creates the crosscap
identification, while the latter affects supersymmetry and the relative
NS--NS/R--R crosscap signs. The geometric support is determined by \(R\)
for the first crosscap and by \(RI_4\) for the second.

\subsubsection{Formal definition of the O-plane
locus}\label{formal-definition-of-the-o-plane-locus}

For an orientation-reversing element \(\Omega g\) whose target-space
part is the involution \(g\), the crosscap condition has the schematic
form \[
X(\tau=0,\sigma_{\mathrm{ws}}+\pi)
=
g\!\left(X(\tau=0,\sigma_{\mathrm{ws}})\right).
\] In a coordinate reflected by \(g\), the zero mode must obey \[
q=-q
\pmod{2\pi R}.
\] It is therefore localized at \[
q=0
\qquad\text{or}\qquad
q=\pi R.
\] In a coordinate left invariant by \(g\), the crosscap is not
localized. Consequently, the associated O-plane worldvolume is the
noncompact spacetime times \(\operatorname{Fix}(g)\) in the compact
target.

This is also the book's definition in Chapter 9: after transforming the
Klein bottle to the tree channel, the closed string couples to a
crosscap source, and the target-space locus of that source is called an
orientifold plane. Strictly speaking, it is an O7-plane rather than an
``O7-brane'': it is not a dynamical D-brane, has no transverse position
modulus, and does not support open-string endpoints.

For an orbifold target, one must include every orientation-reversing
group element: \[
\boxed{
\mathcal L_O
=
\bigcup_{\substack{\gamma\in\Gamma_{\mathrm{or}}\\
\gamma\ {\rm reverses\ worldsheet\ orientation}}}
\operatorname{Fix}(g_\gamma)
\Big/
\langle I_4\rangle
}.
\] Here \(g_\gamma\) means the geometric part of \(\gamma\). In the
present model this gives \[
\mathcal L_O
=
\left[
\operatorname{Fix}(R)
\cup
\operatorname{Fix}(RI_4)
\right]
\Big/
\langle I_4\rangle.
\]

This also gives a useful formulation directly on the quotient \[
X=T^4/\langle I_4\rangle.
\] The reflection induced by \(R\) acts on an orbifold point as \[
\widehat R:[p]\mapsto[Rp].
\] Such a point is fixed when \[
[Rp]=[p].
\] But the orbit represented by \([p]\) contains both \(p\) and
\(I_4p\). Hence there are two ways for this condition to hold: \[
Rp=p
\pmod{\Lambda},
\qquad\text{or}\qquad
Rp=I_4p
\pmod{\Lambda}.
\] The first condition is the \(R\) branch. The second can be rewritten
as \[
RI_4p=p
\pmod{\Lambda},
\] and is the \(RI_4\) branch. Therefore \[
\boxed{
\operatorname{Fix}_X(\widehat R)
=
\left[
\operatorname{Fix}_{T^4}(R)
\cup
\operatorname{Fix}_{T^4}(RI_4)
\right]
\Big/
\langle I_4\rangle
}.
\] This is the precise quotient-space version of the answer: even the
fixed locus of the induced reflection on \(T^4/\mathbb Z_2\) has both
the \(x\)-oriented and the \(y\)-oriented branches.

\subsubsection{Explicit fixed loci and their
multiplicities}\label{explicit-fixed-loci-and-their-multiplicities}

For the rectangular torus, \[
x^i\sim x^i+2\pi R_x,
\qquad
y^i\sim y^i+2\pi R_y.
\]

For \(R\), the fixed-point equations are \[
y^i=-y^i
\pmod{2\pi R_y},
\] or \[
y^i\in\{0,\pi R_y\},
\qquad i=1,2,
\] while \(x^1,x^2\) are arbitrary. Thus, in the covering geometry, \[
\boxed{
\operatorname{Fix}(R)
=
\bigcup_{\epsilon_1,\epsilon_2\in\{0,1\}}
\left\{
y^1=\epsilon_1\pi R_y,\,
y^2=\epsilon_2\pi R_y
\right\}
\cong
4T_x^2
}.
\] Each component has worldvolume \[
\mathbb R^{1,5}\times T_x^2,
\] so it is an O7-plane.

For \(RI_4\), the fixed-point equations are instead \[
x^i=-x^i
\pmod{2\pi R_x},
\] or \[
x^i\in\{0,\pi R_x\},
\qquad i=1,2,
\] while \(y^1,y^2\) are arbitrary. Hence \[
\boxed{
\operatorname{Fix}(RI_4)
=
\bigcup_{\epsilon_1,\epsilon_2\in\{0,1\}}
\left\{
x^1=\epsilon_1\pi R_x,\,
x^2=\epsilon_2\pi R_x
\right\}
\cong
4T_y^2
}.
\] Each component has worldvolume \[
\mathbb R^{1,5}\times T_y^2,
\] and is the second kind of O7-plane.

Thus the upstairs covering geometry contains four \(O7_x\) components
and four \(O7_y\) components. The book's singular phrase ``an O7-plane
along the two \(x\)-axes'' or ``an O7-plane along the two \(y\)-axes''
refers to the corresponding crosscap source or family, not to only one
connected component.

The two families intersect when both \(x^i\) and \(y^i\) are
half-periods. These are precisely the 16 fixed points of \(I_4\).

\subsubsection{The two loci from the loop-channel lattice
sums}\label{the-two-loci-from-the-loop-channel-lattice-sums}

The same result is encoded directly in equation (15.33). For one compact
circle, \[
\Omega|m,n\rangle
=
|m,-n\rangle,
\qquad
I|m,n\rangle
=
|-m,-n\rangle,
\] where \(m\) and \(n\) are momentum and winding numbers.

Consider first the insertion \(\mathcal O=\Omega R(-1)^{F_L}\).

\begin{itemize}
\tightlist
\item
  Along each \(x\) direction, \(R\) acts trivially. The insertion
  therefore acts as \(\Omega\), and a state contributes to the trace
  only when \(n_x=0\). The surviving sum is a momentum sum \(P_x\).
\item
  Along each \(y\) direction, \(R\) is a reflection. The combination
  \(\Omega R\) maps \[
  |m_y,n_y\rangle
  \longmapsto
  |-m_y,n_y\rangle,
  \] so the trace requires \(m_y=0\). The surviving sum is a winding sum
  \(W_y\).
\end{itemize}

There are two coordinates of each kind, so this insertion produces \[
\mathcal K_{\mathcal O}^{(0)}
\propto
P_x(t)^2W_y(t)^2.
\] Momentum in \(x\) and winding in \(y\) are the closed-string
zero-mode signature of Neumann-like crosscap conditions along \(x\) and
Dirichlet-like localization along \(y\). Thus this is the \(O7_x\)
contribution.

Now consider \(\mathcal OI_4\). Its geometric action \(RI_4\) reflects
\(x\) and leaves \(y\) unchanged.

\begin{itemize}
\tightlist
\item
  Along \(x\), the trace requires \(m_x=0\) and leaves a winding sum
  \(W_x\).
\item
  Along \(y\), the trace requires \(n_y=0\) and leaves a momentum sum
  \(P_y\).
\end{itemize}

Therefore \[
\mathcal K_{\mathcal OI_4}^{(0)}
\propto
P_y(t)^2W_x(t)^2,
\] which is the \(O7_y\) contribution. Equation (15.33) is exactly \[
\boxed{
\mathcal K^{(0)}
\propto
P_x^2W_y^2
+
P_y^2W_x^2
}.
\] The notation ``\((x\leftrightarrow y)\)'' in equation (15.35) is the
tree-channel remnant of these two distinct insertions.

\subsubsection{What can be read directly from the Klein-bottle
trace}\label{what-can-be-read-directly-from-the-klein-bottle-trace}

Your interpretation of the two terms is correct, with a small refinement
in terminology. Expanding the orbifold projector gives \[
\begin{aligned}
\mathcal K
\supset
\frac12\operatorname{Tr}
\left(
\mathcal O
P_{\mathrm{GSO}}
e^{-4\pi t(L_0-c/24)}
\right)
+
\frac12\operatorname{Tr}
\left(
\mathcal O I_4
P_{\mathrm{GSO}}
e^{-4\pi t(L_0-c/24)}
\right).
\end{aligned}
\] Writing a general orientation-reversing element as \[
\gamma
=
\Omega g\,\Xi,
\] where \(g\) is its geometric target-space action and \(\Xi\) is any
additional internal operator such as \((-1)^{F_L}\), the rule is \[
\boxed{
\operatorname{Tr}_{\mathcal H}
\left(
\gamma\,q^{L_0-c/24}\bar q^{\,\bar L_0-\bar c/24}
\right)
\quad\longleftrightarrow\quad
\text{crosscap associated with }\gamma,
\qquad
\mathcal L_{O,\gamma}
=
\operatorname{Fix}(g).
}
\] Thus, in this model, \[
\begin{array}{c|c|c}
\text{trace insertion}
&\text{geometric part}
&\text{O-plane locus}\\ \hline
\mathcal O=\Omega R(-1)^{F_L}
&R
&\operatorname{Fix}(R)=4T_x^2\\
\mathcal OI_4=\Omega RI_4(-1)^{F_L}
&RI_4
&\operatorname{Fix}(RI_4)=4T_y^2
\end{array}
\]

It is therefore common, but slightly imprecise, to say that the planes
are the ``fixed loci of \(\mathcal O\) and \(\mathcal OI_4\).'' Since
\(\Omega\) reverses the worldsheet and \((-1)^{F_L}\) acts on the string
Hilbert space, neither is an ordinary map of target-space points. The
precise statement is that the planes are fixed loci of the geometric
parts \(R\) and \(RI_4\) of those two elements.

There are two complementary levels at which the trace supplies
information:

\begin{enumerate}
\def\labelenumi{\arabic{enumi}.}
\tightlist
\item
  \textbf{The insertion label identifies the geometric involution.} From
  \(\mathcal O\) one extracts \(R\), and from \(\mathcal OI_4\) one
  extracts \(RI_4\). Solving their fixed-point equations gives the
  actual loci.
\item
  \textbf{The zero-mode part verifies the wrapped and transverse
  directions.} A momentum sum in a coordinate signals that the crosscap
  is extended along that direction, whereas a winding sum signals
  localization in that direction. Hence \[
  P_x^2W_y^2
  \quad\Longrightarrow\quad
  O7_x,
  \qquad
  P_y^2W_x^2
  \quad\Longrightarrow\quad
  O7_y.
  \] After Poisson resummation, the corresponding volume factors provide
  the same information.
\end{enumerate}

So in this unshifted geometric orientifold, one can indeed read off the
O-plane locus from each term in the Klein-bottle trace by identifying
the geometric part of its insertion and then examining its zero modes.

There is, however, a useful limitation to this statement. A volume
prefactor by itself generally determines only how many directions are
wrapped and which radii are parallel or transverse; it need not
determine every absolute position or every local O-plane charge.
Translated reflections, freely acting shifts, discrete torsion, and
signs in the crosscap state can insert phases into the lattice sums or
alter the distribution and type of O-planes. In such cases one must
retain the complete insertion, including its lattice phases and crosscap
signs, rather than infer the locus from the massless volume scaling
alone. No such ambiguity occurs here because \(R\) and \(RI_4\) are
explicit unshifted reflections on the rectangular torus.

\subsubsection{Why the Klein bottle does not describe open-string
endpoints}\label{why-the-klein-bottle-does-not-describe-open-string-endpoints}

The proposed endpoint interpretation is not correct for the Klein
bottle. Although one often speaks broadly of ``open--closed channel
duality,'' the modular transformation of a Klein bottle does not turn it
into an open-string diagram. A Klein bottle has no worldsheet boundary
in either channel.

The relevant one-loop surfaces are distinguished as follows:

\begin{longtable}[]{@{}
  >{\raggedright\arraybackslash}p{(\linewidth - 4\tabcolsep) * \real{0.3333}}
  >{\raggedright\arraybackslash}p{(\linewidth - 4\tabcolsep) * \real{0.3333}}
  >{\raggedright\arraybackslash}p{(\linewidth - 4\tabcolsep) * \real{0.3333}}@{}}
\toprule\noalign{}
\begin{minipage}[b]{\linewidth}\raggedright
Worldsheet
\end{minipage} & \begin{minipage}[b]{\linewidth}\raggedright
Loop-channel description
\end{minipage} & \begin{minipage}[b]{\linewidth}\raggedright
Transverse-channel description
\end{minipage} \\
\midrule\noalign{}
\endhead
\bottomrule\noalign{}
\endlastfoot
Klein bottle & Closed-string trace with an orientation-reversing
insertion & Closed string propagating between two crosscaps \\
Annulus & Open-string trace & Closed string propagating between two
D-brane boundaries \\
Möbius strip & Open-string trace with an orientation-reversing insertion
& Closed string propagating between a D-brane boundary and a crosscap \\
\end{longtable}

A modular transformation changes the parametrization and propagation
direction on the same worldsheet; it does not change its topology. In
particular, it cannot turn either crosscap of a Klein bottle into a
boundary. Although all three transverse-channel diagrams can be drawn as
tubes, the states placed at the tube ends are different.

Schematically, \[
\begin{aligned}
\mathcal K_g
&=
\operatorname{Tr}_{\mathcal H_{\mathrm{cl}}}
\left(
\Omega g\,\Xi\,q^{H_{\mathrm{cl}}}
\right)
\quad
\xleftrightarrow{\text{modular}}
\quad
\langle C_g|
e^{-2\pi lH_{\mathrm{cl}}}
|C_g\rangle,\\
\mathcal A_{ab}
&=
\operatorname{Tr}_{\mathcal H_{ab}^{\mathrm{op}}}
\left(
q^{H_{\mathrm{op}}}
\right)
\quad
\xleftrightarrow{\text{modular}}
\quad
\langle B_a|
e^{-2\pi lH_{\mathrm{cl}}}
|B_b\rangle,\\
\mathcal M_{a,g}
&=
\operatorname{Tr}_{\mathcal H_{a,g(a)}^{\mathrm{op}}}
\left(
\Omega g\,\Xi\,q^{H_{\mathrm{op}}}
\right)
\quad
\xleftrightarrow{\text{modular}}
\quad
\langle B_a|
e^{-2\pi lH_{\mathrm{cl}}}
|C_g\rangle.
\end{aligned}
\] Numerical factors in the modular parameters have been suppressed here
because the important point is the nature of the states and the ends of
the transverse tube.

\paragraph{What the operator means in the Klein-bottle
trace}\label{what-the-operator-means-in-the-klein-bottle-trace}

In \[
\operatorname{Tr}_{\mathcal H_{\mathrm{cl}}}
\left(
\Omega g\,\Xi\,
e^{-4\pi t(L_0-c/24)}
\right),
\] the trace is over the closed-string Hilbert space. The insertion
means that, after propagating once around the loop-channel Euclidean
time direction, the closed-string configuration is identified with its
\(\Omega g\,\Xi\) image. Because \(\Omega\) reverses worldsheet
orientation, this identification produces a nonorientable closed
worldsheet.

After the modular transformation, the same surface is a tube with a
crosscap at each end. There is still no open string and no open-string
endpoint. The target-space fixed locus appears because the crosscap
state satisfies \[
X(\sigma_{\mathrm{ws}}+\pi)
=
g\!\left(X(\sigma_{\mathrm{ws}})\right).
\] For a reflected zero mode this reduces to \[
q=g(q)=-q
\pmod{2\pi R},
\] which localizes the crosscap source at \(\operatorname{Fix}(g)\).
Thus the trace detects the O-plane through a crosscap condition, not
through an endpoint condition.

For the two terms in equation (15.31), \[
\begin{aligned}
\operatorname{Tr}_{\mathrm{cl}}
\left(
\mathcal O\,P_{\mathrm{GSO}}e^{-4\pi t(L_0-c/24)}
\right)
&\longleftrightarrow
|C_R\rangle
&&\longrightarrow
\operatorname{Fix}(R),\\
\operatorname{Tr}_{\mathrm{cl}}
\left(
\mathcal OI_4\,P_{\mathrm{GSO}}e^{-4\pi t(L_0-c/24)}
\right)
&\longleftrightarrow
|C_{RI_4}\rangle
&&\longrightarrow
\operatorname{Fix}(RI_4).
\end{aligned}
\] Neither line contains a D-brane boundary state.

\paragraph{Where open strings actually
enter}\label{where-open-strings-actually-enter}

Open strings enter only after D7-branes are introduced to cancel the
R--R tadpoles. Their endpoints satisfy boundary conditions defined by
the D7-brane boundary states \(|B_a\rangle\). The annulus is a loop of
such open strings, and its transverse channel is closed-string exchange
between two D-branes.

The Möbius strip is the surface on which the orientifold operator acts
inside an open-string trace. If \(a\) denotes a D-brane, then
\(\Omega g\) reverses the open-string orientation and maps the brane
data to its image under \(g\). An open-string sector can contribute to
the Möbius trace only when it is mapped back to an equivalent sector.
For an invariant brane this may require \[
g(a)=a.
\] For example, a D7-brane parallel to the \(x\) directions and
localized at \((y^1,y^2)\) is invariant under \(R\) only when \[
(y^1,y^2)
\in
\left\{
(0,0),\,
(\pi R_y,0),\,
(0,\pi R_y),\,
(\pi R_y,\pi R_y)
\right\}.
\] Such a D7-brane lies on an \(O7_x\) component. Nevertheless, an
open-string endpoint there is attached to the D7-brane, not to the
O7-plane. The coincidence of their loci does not turn the O-plane into
an object supporting open-string endpoints.

A D-brane need not always be individually fixed. A brane away from the
fixed locus can occur together with its image \(g(a)\); the orientifold
then identifies the two brane sectors. This is why O-planes determine
image rules and Chan--Paton projections without themselves carrying
Chan--Paton gauge fields.

The logical chain is therefore \[
\boxed{
\begin{aligned}
\text{Klein-bottle insertion }\Omega g\,\Xi
&\Longrightarrow
\text{crosscap }|C_g\rangle
\Longrightarrow
\text{O-plane on }\operatorname{Fix}(g),\\
\text{D-brane boundary }|B_a\rangle
&\Longrightarrow
\text{open-string endpoints on }a,\\
\text{Möbius insertion }\Omega g\,\Xi
&\Longrightarrow
\text{orientifold identification of }a\text{ with }g(a).
\end{aligned}
}
\]

\subsubsection{Why the volume prefactors identify the wrapped
directions}\label{why-the-volume-prefactors-identify-the-wrapped-directions}

Poisson resummation gives \[
\sum_{m\in\mathbb Z}
e^{-\pi t(m/\rho)^2}
=
\frac{\rho}{\sqrt t}
\sum_{\widetilde n\in\mathbb Z}
e^{-\pi\rho^2\widetilde n^{\,2}/t},
\] and \[
\sum_{n\in\mathbb Z}
e^{-\pi t(\rho n)^2}
=
\frac{1}{\rho\sqrt t}
\sum_{\widetilde m\in\mathbb Z}
e^{-\pi\widetilde m^{\,2}/(\rho^2t)}.
\] After the Klein-bottle modular transformation \(t=1/(4l)\), the first
lattice product therefore carries \[
P_x^2W_y^2
\longrightarrow
\left(\frac{\rho_x}{\rho_y}\right)^2
\times
\text{tree-channel lattice sums},
\] while the second carries \[
P_y^2W_x^2
\longrightarrow
\left(\frac{\rho_y}{\rho_x}\right)^2
\times
\text{tree-channel lattice sums}.
\] In the massless limit \(l\to\infty\), all nonzero dual momentum and
winding modes are exponentially suppressed. The remaining coefficients
are \[
v_6\left(\frac{\rho_x}{\rho_y}\right)^2,
\qquad
v_6\left(\frac{\rho_y}{\rho_x}\right)^2.
\]

The first coefficient scales as internal wrapped volume divided by
transverse volume, \[
\frac{V_x}{V_y}
=
\frac{R_x^2}{R_y^2},
\] so it describes a source extended along \(x^1,x^2\) and localized in
\(y^1,y^2\). The second has the reciprocal scaling and describes a
source extended along \(y^1,y^2\) and localized in \(x^1,x^2\). This is
why the tadpole calculation independently recovers the two fixed-locus
geometries.

\subsubsection{Checks}\label{checks-21}

If the \(I_4\) orbifold projector were absent, equation (15.31) would
contain only \(\mathcal O\). Then only \(\operatorname{Fix}(R)\) and the
\(O7_x\) family would occur. The \(O7_y\) family is therefore entirely
tied to the fact that the orientifold is performed after the \(I_4\)
quotient.

The two fixed sets have real dimension two inside \(T^4\). Adding the
six noncompact dimensions gives an eight-dimensional worldvolume, which
is precisely the worldvolume dimension of an O7-plane.

Finally, \((-1)^{F_L}\) does not change either fixed locus. It changes
the action on spacetime fermions and the R--R part of the crosscap
state. Likewise, the extra twisted-sector sign chosen on page 536
changes twisted-sector charges and the resulting multiplet projection,
but not the geometric equations defining the two O7 loci.

\subsubsection{Result}\label{result-25}

The orientifold group contains \[
\boxed{
\Gamma_{\mathrm{or}}
=
\left\{
1,\,
I_4,\,
\Omega R(-1)^{F_L},\,
\Omega RI_4(-1)^{F_L}
\right\}.
}
\] Its two orientation-reversing elements have geometric parts \[
\begin{aligned}
R:(x^i,y^i)&\mapsto(x^i,-y^i),\\
RI_4:(x^i,y^i)&\mapsto(-x^i,y^i).
\end{aligned}
\] Consequently, \[
\boxed{
\begin{aligned}
\operatorname{Fix}(R)
&=
4T_x^2
&&\Longrightarrow&
4\ O7_x\text{ components},\\
\operatorname{Fix}(RI_4)
&=
4T_y^2
&&\Longrightarrow&
4\ O7_y\text{ components}.
\end{aligned}
}
\] Equivalently, on the quotient, \[
\boxed{
[Rp]=[p]
\quad\Longleftrightarrow\quad
Rp=p
\ \text{or}\
Rp=I_4p,
}
\] which produces the same two branches.

The loop-channel zero modes and tree-channel tadpole prefactors are \[
\boxed{
\begin{aligned}
\Omega R(-1)^{F_L}
&:
P_x^2W_y^2
\longrightarrow
v_6\left(\frac{\rho_x}{\rho_y}\right)^2
&&\Longrightarrow O7_x,\\
\Omega RI_4(-1)^{F_L}
&:
P_y^2W_x^2
\longrightarrow
v_6\left(\frac{\rho_y}{\rho_x}\right)^2
&&\Longrightarrow O7_y.
\end{aligned}
}
\] More generally, the trace-to-locus rule is \[
\boxed{
\gamma=\Omega g\,\Xi
\quad\Longrightarrow\quad
\operatorname{Tr}(\gamma\cdots)
\text{ identifies the }\gamma\text{-crosscap},
\qquad
\mathcal L_{O,\gamma}=\operatorname{Fix}(g),
}
\] with the zero-mode lattice sums determining the parallel and
transverse directions. This is a crosscap statement, not an open-string
endpoint statement: \[
\boxed{
\mathcal K:\ C_g\text{--}C_g,
\qquad
\mathcal A:\ B_a\text{--}B_b,
\qquad
\mathcal M:\ B_a\text{--}C_g.
}
\] Thus your naive fixed-locus argument correctly finds the first
family; the second comes from the equally mandatory
orientation-reversing element obtained by multiplying the orientifold
generator by the orbifold element \(I_4\). Open-string endpoints appear
only after D7-brane boundary states are introduced.

\section{Q6: Derivation and massless interpretation of the fractional-D7
annulus}\label{q6-derivation-and-massless-interpretation-of-the-fractional-d7-annulus}

In page 539, they compute the amplitude for open string starting and
ending on the same D7-brane in (15.36) and the result is (15.37).
\textgreater Since \(\vartheta_{4}\) and \(\vartheta_{3}\) differ by the
sign for the first excited level, this term implies that in the massless
sector only the excitations \(\psi_{-\frac{1}{2}}^\mu\) in the six
non-compact directions survive the \(I_{4}\) projection, whereas the
excitations \(\psi_{-\frac{1}{2}}^i\) in the internal directions
\(i=5,\cdots,8\) are projected out. The D7-brane therefore supports a
six-dimensional gauge field and cannot be moved in the internal
directions, i.e.~it is frozen at \((y_{1},y_{2})=(0,0)\) and has no
continuous Wilson-line in the x-directions.

I am not such familiar with the computation of this partition function
and its interpretation. Do the explicit computation for me and explain
the above results from the computation.

\begin{quote}
\textbf{Textbook location and related material.} Chapter 15, Section
15.3 around Eqs. (15.36)--(15.41), pp.~539--544: the fractional-D7
annulus, \(I_4\) projection, tree-channel twisted exchange, and
orientifold images.
\end{quote}

\subsection{Answer}\label{answer-26}

\subsubsection{Core point}\label{core-point-26}

The annulus in equation (15.36) is an open-string one-loop trace with
two independent projections: \[
\mathcal A
=
\int_0^\infty\frac{dt}{2t}\,
\operatorname{Tr}_{(D7,D7)}
\left[
\frac{1+I_4}{2}\,
P_{\mathrm{GSO}}\,
e^{-2\pi t(L_0-c/24)}
\right].
\] For a D7-brane parallel to the two \(x\)-axes, the trace has the
following structure:

\begin{enumerate}
\def\labelenumi{\arabic{enumi}.}
\tightlist
\item
  the six noncompact zero modes give the measure \(v_6\,dt/t^4\);
\item
  the two Neumann \(x\) directions give Kaluza-Klein momentum sums;
\item
  the two Dirichlet \(y\) directions give winding sums;
\item
  the identity part of \((1+I_4)/2\) gives the ordinary D7--D7
  oscillator character;
\item
  the \(I_4\)-inserted part has no nonzero compact zero modes and
  changes the signs of all internal oscillators.
\end{enumerate}

Combining these pieces gives exactly \[
\boxed{
\begin{aligned}
\mathcal A_{xx}
&=
v_6\int_0^\infty\frac{dt}{t^4}\,
\frac{1}{2^6}
(1_{\mathrm{NS}}-1_{\mathrm R})\\
&\quad\times
\left[
\frac{\vartheta_4^4}{\eta^{12}}(2it)\,
\left(
\sum_{m\in\mathbb Z}
e^{-2\pi t(m/\rho_x)^2}
\right)^2
\left(
\sum_{n\in\mathbb Z}
e^{-2\pi t(\rho_y n)^2}
\right)^2
+
4\,
\frac{\vartheta_4^2\vartheta_3^2}
{\eta^6\vartheta_2^2}(2it)
\right].
\end{aligned}
}
\] The first term in brackets is the trace without \(I_4\); the second
is the trace with \(I_4\).

At the massless NS level, before the orbifold projection, the physical
states are \[
\psi_{-1/2}^{M}|0;k\rangle_{\mathrm{NS}},
\qquad
M=\mu\ \text{or}\ i.
\] Their \(I_4\) eigenvalues are \[
I_4\psi_{-1/2}^{\mu}|0\rangle_{\mathrm{NS}}
=
+\psi_{-1/2}^{\mu}|0\rangle_{\mathrm{NS}},
\qquad
I_4\psi_{-1/2}^{i}|0\rangle_{\mathrm{NS}}
=
-\psi_{-1/2}^{i}|0\rangle_{\mathrm{NS}}.
\] Therefore \[
\frac{1+I_4}{2}
\psi_{-1/2}^{\mu}|0\rangle_{\mathrm{NS}}
=
\psi_{-1/2}^{\mu}|0\rangle_{\mathrm{NS}},
\qquad
\frac{1+I_4}{2}
\psi_{-1/2}^{i}|0\rangle_{\mathrm{NS}}
=
0.
\] The surviving state is the six-dimensional gauge vector. The four
removed internal states are precisely two transverse-position scalars
and two Wilson-line scalars. We now derive every ingredient of the
amplitude and this interpretation.

\subsubsection{Geometry and open-string boundary
conditions}\label{geometry-and-open-string-boundary-conditions}

Take \[
D7_x:
\qquad
\mathbb R^{1,5}\times T_x^2,
\] localized at a point \((y^1,y^2)\) of the transverse \(T_y^2\). Its
boundary conditions are \[
\begin{array}{c|c|c}
\text{directions}
&\text{boundary condition}
&\text{open-string zero modes}\\ \hline
\mathbb R^{1,5}
&NN
&k_\mu\\
x^1,x^2
&NN
&m_a/R_x\\
y^1,y^2
&DD
&n_aR_y/\alpha'
\end{array}
\] with \(m_a,n_a\in\mathbb Z\).

The Dirichlet winding deserves emphasis. Both endpoints lie on the same
point of the compact \(y\)-circle, but the string can stretch from that
point to one of its images in the covering space: \[
Y(\pi)-Y(0)
=
2\pi nR_y.
\] Its classical stretching energy is therefore proportional to \[
\left(\frac{nR_y}{\alpha'}\right)^2.
\] This produces the two types of lattice factor \[
P_x(t)
=
\sum_{m\in\mathbb Z}
e^{-2\pi t(m/\rho_x)^2},
\qquad
W_y(t)
=
\sum_{n\in\mathbb Z}
e^{-2\pi t(\rho_y n)^2}.
\] Because there are two \(x\) and two \(y\) directions, the complete
compact zero-mode factor is \[
P_x(t)^2W_y(t)^2.
\]

\subsubsection{\texorpdfstring{Why the D7-brane must be at an \(I_4\)
fixed
position}{Why the D7-brane must be at an I\_4 fixed position}}\label{why-the-d7-brane-must-be-at-an-i_4-fixed-position}

The orbifold inversion maps \[
D7_x(y^1,y^2)
\longmapsto
D7_x(-y^1,-y^2).
\] The operator \(I_4\) therefore maps the Hilbert space \[
\mathcal H_{(D7(y),D7(y))}
\longrightarrow
\mathcal H_{(D7(-y),D7(-y))}.
\] A trace with an \(I_4\) insertion is defined within the same
open-string Hilbert space only if the brane is mapped to itself: \[
(y^1,y^2)
=
(-y^1,-y^2)
\pmod{2\pi R_y}.
\] Hence \[
\boxed{
(y^1,y^2)
\in
\left\{
(0,0),\,
(\pi R_y,0),\,
(0,\pi R_y),\,
(\pi R_y,\pi R_y)
\right\}.
}
\] The book chooses \((0,0)\) as a representative. Thus ``frozen at
\((0,0)\)'' means frozen at the particular fixed point selected in the
example; the same calculation applies at any of the four fixed
positions.

If the brane is away from these points, a consistent orbifold
configuration must also contain its image at \((-y^1,-y^2)\). The
individual \((D7(y),D7(y))\) trace with \(I_4\) then vanishes because
\(I_4\) maps it to a distinct sector.

\subsubsection{\texorpdfstring{Zero modes in the identity and \(I_4\)
traces}{Zero modes in the identity and I\_4 traces}}\label{zero-modes-in-the-identity-and-i_4-traces}

Expanding the orbifold projector, \[
\mathcal A
=
\frac12\mathcal A[1]
+
\frac12\mathcal A[I_4].
\]

For the identity trace, every momentum and winding state contributes: \[
\mathcal A[1]_{\mathrm{zero}}
=
P_x(t)^2W_y(t)^2.
\]

For the inserted trace, \(I_4\) reverses all four compact coordinates.
It therefore acts on the zero-mode labels as \[
I_4:
\qquad
|m_1,m_2;n_1,n_2\rangle
\longmapsto
|-m_1,-m_2;-n_1,-n_2\rangle.
\] The diagonal matrix element entering a trace is \[
\begin{aligned}
&\langle m_1,m_2;n_1,n_2|
I_4
|m_1,m_2;n_1,n_2\rangle\\
&\qquad=
\delta_{m_1,-m_1}
\delta_{m_2,-m_2}
\delta_{n_1,-n_1}
\delta_{n_2,-n_2}.
\end{aligned}
\] For integer labels, all four conditions require \[
m_1=m_2=n_1=n_2=0.
\] Consequently, \[
\boxed{
\mathcal A[I_4]_{\mathrm{zero}}=1.
}
\] This explains why the last term of equation (15.37) has no momentum
or winding sums.

\subsubsection{\texorpdfstring{The six-dimensional measure and the
factor
\(2^{-6}\)}{The six-dimensional measure and the factor 2\^{}\{-6\}}}\label{the-six-dimensional-measure-and-the-factor-2-6}

The Gaussian trace over the six noncompact momenta is \[
\begin{aligned}
V_6\int\frac{d^6k}{(2\pi)^6}
e^{-2\pi t\alpha' k^2}
&=
\frac{V_6}{(8\pi^2\alpha't)^3}\\
&=
\frac{v_6}{2^3t^3},
\end{aligned}
\] where \[
v_6
=
\frac{V_6}{(4\pi^2\alpha')^3}.
\] Combining this with the annulus measure \(dt/(2t)\) gives \[
v_6\frac{dt}{t^4}\frac1{2^4}.
\] The orbifold projector contributes another factor \(1/2\), and the
GSO projector another factor \(1/2\). Hence \[
\frac1{2^4}\times\frac12\times\frac12
=
\boxed{\frac1{2^6}},
\] which is the normalization displayed in equation (15.37) for a single
Chan--Paton label.

\subsubsection{\texorpdfstring{Oscillator trace without
\(I_4\)}{Oscillator trace without I\_4}}\label{oscillator-trace-without-i_4}

In light-cone gauge there are eight physical bosonic and eight physical
fermionic oscillators. The eight unreflected bosons give \[
Z_B[1]
=
\frac1{\eta^8}.
\] After the open-string GSO sum, the NS and R characters are equal by
the Jacobi identity and occur with opposite spacetime-statistics signs.
In the book's annulus convention this is written \[
Z_F[1]
=
(1_{\mathrm{NS}}-1_{\mathrm R})
\frac{\vartheta_4^4}{\eta^4}.
\] Therefore \[
\boxed{
Z_{\mathrm{osc}}[1]
=
(1_{\mathrm{NS}}-1_{\mathrm R})
\frac{\vartheta_4^4}{\eta^{12}}.
}
\] Multiplication by \(P_x^2W_y^2\) produces the first term in the
square brackets of equation (15.37).

The symbol \((1_{\mathrm{NS}}-1_{\mathrm R})\) records the cancellation
of the complete vacuum amplitude between spacetime bosons and fermions.
It should not be used to discard the amplitude before reading the
spectrum: one first examines the NS and R characters separately.

\subsubsection{\texorpdfstring{Oscillator trace with
\(I_4\)}{Oscillator trace with I\_4}}\label{oscillator-trace-with-i_4}

The inversion \(I_4\) acts trivially on the four physical noncompact
transverse oscillators and as \(-1\) on the four internal oscillators.
The inserted trace therefore splits into four even and four odd real
fields.

For an \(I_4\)-even bosonic oscillator, \[
\operatorname{Tr}
\left(
q^{nN}
\right)
=
\sum_{N=0}^{\infty}q^{nN}
=
\frac1{1-q^n}.
\] For an \(I_4\)-odd bosonic oscillator, \[
\operatorname{Tr}
\left(
I_4q^{nN}
\right)
=
\sum_{N=0}^{\infty}(-1)^Nq^{nN}
=
\frac1{1+q^n}.
\] The four internal real bosons form two complex bosons. Using the
product formula for \(\vartheta_2\), \[
Z_{B,\mathrm{internal}}[I_4]
=
\left(\frac{2\eta}{\vartheta_2}\right)^2.
\] The four noncompact real bosons give \(\eta^{-4}\), so \[
\boxed{
Z_B[I_4]
=
\frac1{\eta^4}
\left(\frac{2\eta}{\vartheta_2}\right)^2
=
\frac{4}{\eta^2\vartheta_2^2}.
}
\] This is the origin of the explicit factor \(4\) in equation (15.37).
Equivalently, \(\vartheta_2\) has a leading factor \(2\), so the
explicit \(4\) removes the factor contained in \(\vartheta_2^2\).

For an NS fermionic oscillator, the occupation number is \(N=0,1\). An
even oscillator contributes \[
\sum_{N=0}^{1}q^{rN}
=
1+q^r,
\] whereas an odd oscillator contributes \[
\sum_{N=0}^{1}(-1)^Nq^{rN}
=
1-q^r.
\] The product formulas identify these two signs with \(\vartheta_3\)
and \(\vartheta_4\). Since two complex fermions are even and two are
odd, the inserted fermion character is \[
Z_F[I_4]
=
(1_{\mathrm{NS}}-1_{\mathrm R})
\frac{\vartheta_4^2\vartheta_3^2}{\eta^4}.
\] Therefore \[
\boxed{
Z_{\mathrm{osc}}[I_4]
=
(1_{\mathrm{NS}}-1_{\mathrm R})
4\frac{\vartheta_4^2\vartheta_3^2}
{\eta^6\vartheta_2^2}.
}
\] Because the inserted trace has no compact zero-mode sum, this is
exactly the second term in equation (15.37).

\subsubsection{\texorpdfstring{Reading the first excited level from
\(\vartheta_3\) and
\(\vartheta_4\)}{Reading the first excited level from \textbackslash vartheta\_3 and \textbackslash vartheta\_4}}\label{reading-the-first-excited-level-from-vartheta_3-and-vartheta_4}

For vanishing compact momentum and winding, the NS mass formula is \[
\alpha'M^2
=
N_{\mathrm{NS}}-\frac12.
\] The GSO projection removes the \(N_{\mathrm{NS}}=0\) tachyon. The
first allowed states have \[
N_{\mathrm{NS}}=\frac12,
\] and are created by one fermionic oscillator: \[
\psi_{-1/2}^{M}|0\rangle_{\mathrm{NS}}.
\] These are the massless states whose \(I_4\) eigenvalues must be
extracted from the first excited terms of the fermion characters.

Let \(Q=e^{\pi i\tau}\). The relevant leading terms are \[
\vartheta_3(\tau)
=
1+2Q+\mathcal O(Q^4),
\qquad
\vartheta_4(\tau)
=
1-2Q+\mathcal O(Q^4).
\] The opposite signs of the first excited contribution are precisely
the insertion eigenvalues. Four real fermions, or two complex fermions,
contribute \[
\vartheta_3^2
=
1+4Q+\cdots,
\qquad
\vartheta_4^2
=
1-4Q+\cdots.
\] Thus the \(I_4\)-inserted character weights the first excited states
in the two four-dimensional oscillator blocks with opposite signs. Up to
the conventional overall phase assigned to the NS vacuum, this is \[
\operatorname{Tr}_{\mathrm{massless,NS}}I_4
=
4-4
=
0.
\] The four positive contributions are the noncompact excitations and
the four negative contributions are the internal ones.

The trace without insertion counts all eight massless NS states: \[
\operatorname{Tr}_{\mathrm{massless,NS}}1
=
4+4
=
8.
\] Applying the projector gives the total number \[
\operatorname{Tr}_{\mathrm{massless,NS}}
\frac{1+I_4}{2}
=
\frac12(8+0)
=
4.
\] More importantly, applying it state by state gives \[
\begin{array}{c|c|c}
\text{state}
&I_4
&\frac12(1+I_4)\\ \hline
\psi_{-1/2}^{\mu}|0\rangle_{\mathrm{NS}}
&+1
&\text{kept}\\
\psi_{-1/2}^{i}|0\rangle_{\mathrm{NS}}
&-1
&\text{removed}
\end{array}
\] which is the statement made below equation (15.37).

This projection assumes that both endpoints lie on the same
fractional-brane type. The Chan--Paton action is then \[
\lambda
\longmapsto
\gamma_{I_4}\lambda\gamma_{I_4}^{-1}
=
\lambda,
\] so the total \(I_4\) parity is the oscillator parity displayed above.
For an open string stretched between fractional branes with opposite
twisted charges, the Chan--Paton factor contributes an additional minus
sign. The projection is then reversed: the vector is removed while
internal scalar states survive, in agreement with the discussion
following equation (15.40).

\subsubsection{From the surviving states to six-dimensional
fields}\label{from-the-surviving-states-to-six-dimensional-fields}

Before the orbifold projection, a D7-brane supports an eight-dimensional
gauge field. Reducing it on the wrapped \(T_x^2\) and viewing the result
in six dimensions gives \[
A_M
\longrightarrow
\left(
A_\mu,\,
A_{x^1},\,
A_{x^2}
\right),
\] together with the two transverse scalars \[
\Phi_{y^1},\,
\Phi_{y^2}.
\] The correspondence with NS oscillators is \[
\begin{array}{c|c|c}
\text{oscillator}
&\text{six-dimensional field}
&\text{geometric meaning}\\ \hline
\psi_{-1/2}^{\mu}
&A_\mu
&\text{gauge vector}\\
\psi_{-1/2}^{x^1},\psi_{-1/2}^{x^2}
&A_{x^1},A_{x^2}
&\text{Wilson-line scalars}\\
\psi_{-1/2}^{y^1},\psi_{-1/2}^{y^2}
&\Phi_{y^1},\Phi_{y^2}
&\text{transverse-position scalars}
\end{array}
\]

Since \(I_4\) leaves the noncompact coordinates invariant, \[
A_\mu\mapsto+A_\mu.
\] Since it reverses all four internal coordinates, a constant internal
vector component and a transverse displacement are both odd: \[
A_{x^a}\mapsto-A_{x^a},
\qquad
\Phi_{y^a}\mapsto-\Phi_{y^a}.
\] Therefore only \(A_\mu\) survives.

The four states \(\psi_{-1/2}^{\mu}|0\rangle\) are the four physical
polarizations of a massless vector in six dimensions: \[
d-2
=
6-2
=
4.
\] The R-sector projection retains the corresponding gaugino states, so
the result is a six-dimensional supersymmetric vector multiplet.

The absence of \(\Phi_{y^a}\) means that an individual invariant brane
has no infinitesimal transverse displacement. Moving it by \(\delta y\)
would produce a brane at \(\delta y\) and an image at \(-\delta y\),
which is not a deformation of a single invariant fractional brane.

The absence of \(A_{x^a}\) means there is no continuous Wilson-line
modulus. This does not rule out isolated discrete Wilson-line choices
fixed by \(A_x\sim-A_x\) modulo large gauge transformations; it says
that there is no massless scalar generating a continuous Wilson-line
deformation.

\subsubsection{Tree-channel check and the meaning of
``fractional''}\label{tree-channel-check-and-the-meaning-of-fractional}

Using \[
t=\frac1{2l},
\] Poisson resummation of the lattice sums, and the modular
transformations of \(\eta\) and \(\vartheta_a\), equation (15.37)
becomes \[
\boxed{
\begin{aligned}
\widetilde{\mathcal A}_{xx}
&=
v_6\int_0^\infty dl\,
\frac1{2^7}
(1_{\mathrm{NS}}-1_{\mathrm R})\\
&\quad\times
\left[
\left(\frac{\rho_x}{\rho_y}\right)^2
\frac{\vartheta_4^4}{\eta^{12}}(2il)
\left(
\sum_{n\in\mathbb Z}
e^{-\pi l(\rho_xn)^2}
\right)^2
\left(
\sum_{m\in\mathbb Z}
e^{-\pi l(m/\rho_y)^2}
\right)^2\right.\\
&\qquad\left.
+
4^2
\frac{\vartheta_2^2\vartheta_3^2}
{\eta^6\vartheta_4^2}(2il)
\right].
\end{aligned}
}
\] This is equation (15.38).

The first term is closed-string exchange from the \(I_4\)-untwisted
sector between two copies of the D7 boundary state. Its factor \[
\left(\frac{\rho_x}{\rho_y}\right)^2
\] has the correct scaling for a D7-brane extended in \(x^1,x^2\) and
localized in \(y^1,y^2\).

The second term has no momentum or winding sums and transforms into a
twisted closed-string character. It is exchange of \(I_4\)-twisted
closed-string states. Hence the boundary state has both untwisted and
twisted components. With the book's normalization, equation (15.39) is
\[
|D7_x^\pm\rangle_{\mathrm{frac}}
=
\frac{1}{2\mathcal N_{D7}}
\left[
\left(\frac{\rho_x}{\rho_y}\right)
|D7_x\rangle_{\mathrm{untw}}
\pm
2\sum_{a=1}^{4}
|D7_x,a\rangle_{\mathrm{tw}}
\right].
\] The four terms correspond to the four \(I_4\) fixed points through
which the \(x\)-wrapping D7-brane passes.

A twisted R--R field is localized at the orbifold fixed set. A brane
carrying its charge cannot move continuously away from that set without
violating charge conservation. This is the closed-string channel
explanation of the same freezing already seen in the open-string
spectrum.

Two fractional branes with opposite twisted charges have \[
\left(
|D7_x\rangle_{\mathrm{untw}}
+
|D7_x\rangle_{\mathrm{tw}}
\right)
+
\left(
|D7_x\rangle_{\mathrm{untw}}
-
|D7_x\rangle_{\mathrm{tw}}
\right)
=
2|D7_x\rangle_{\mathrm{untw}}.
\] Their twisted charges cancel, and together they can form an ordinary
bulk D7-brane with movable brane-image degrees of freedom. Thus the
statement that the brane is frozen applies to an individual fractional
D7-brane, not to every possible D7-brane combination in the orbifold.

\subsubsection{Why the annulus proves that the D7-brane carries twisted
R--R
charge}\label{why-the-annulus-proves-that-the-d7-brane-carries-twisted-rr-charge}

There are three steps:

\begin{enumerate}
\def\labelenumi{\arabic{enumi}.}
\tightlist
\item
  the \(I_4\) insertion in the open loop channel becomes an \(I_4\)
  twist in the closed transverse channel;
\item
  the R contribution to that closed-string exchange is therefore an
  \(I_4\)-twisted R--R state;
\item
  a nonzero boundary-state overlap with an R--R state is an R--R source
  coupling and hence an R--R charge.
\end{enumerate}

\paragraph{From an inserted open trace to a twisted closed
string}\label{from-an-inserted-open-trace-to-a-twisted-closed-string}

The \(I_4\) contribution to equation (15.36) is \[
\mathcal A[I_4]
=
\operatorname{Tr}_{\mathcal H_{\mathrm{op}}}
\left[
I_4P_{\mathrm{GSO}}
e^{-2\pi t(L_0-c/24)}
\right].
\] In the open loop channel, the insertion means that fields are
identified after one trip around Euclidean time by \[
X(\sigma,\tau+2\pi t)
=
I_4X(\sigma,\tau).
\] The annulus modular transformation exchanges its spatial and
Euclidean-time cycles. In the transverse channel, the propagating closed
string consequently obeys \[
X(\sigma+2\pi,\tau)
=
I_4X(\sigma,\tau).
\] This is the defining boundary condition of the \(I_4\)-twisted
closed-string Hilbert space: \[
\boxed{
\operatorname{Tr}_{\mathrm{open}}
\left(
I_4q^{H_{\mathrm{op}}}
\right)
\quad\xleftrightarrow{S}\quad
\text{closed-string propagation in }
\mathcal H_{\mathrm{cl}}^{(I_4\text{-tw})}.
}
\]

The theta functions in equations (15.37) and (15.38) implement this
exchange: \[
4\,
\frac{\vartheta_4^2\vartheta_3^2}
{\eta^6\vartheta_2^2}(2it)
\quad\xrightarrow{\,t=1/(2l)\,}\quad
4^2\,
\frac{\vartheta_2^2\vartheta_3^2}
{\eta^6\vartheta_4^2}(2il).
\] The expression on the right is the twisted closed-string character.
Its lack of compact momentum and winding sums is another diagnostic: an
\(I_4\)-twisted closed string is localized at an orbifold fixed point
and has no translational zero mode in the four inverted directions.

\paragraph{Why the R term means R--R
charge}\label{why-the-r-term-means-rr-charge}

The shorthand \[
(1_{\mathrm{NS}}-1_{\mathrm R})
\] in equation (15.38) contains a nonzero NS--NS exchange and a nonzero
R--R exchange. They cancel only after the supersymmetric vacuum
amplitude is summed; to determine charges, the two sectors must be kept
separate.

In the transverse channel the annulus factorizes as \[
\widetilde{\mathcal A}
=
\langle B|
e^{-2\pi lH_{\mathrm{cl}}}
|B\rangle.
\] Inserting a complete set of closed-string states gives \[
\widetilde{\mathcal A}
=
\sum_{\alpha}
\langle B|\alpha\rangle
e^{-2\pi lE_\alpha}
\langle\alpha|B\rangle.
\] For a long tube, \(l\to\infty\), only the lightest closed-string
states remain. If \(\alpha\) is a massless twisted R--R state, its
contribution is proportional to \[
\left|
\langle\alpha_{\mathrm{tw,R}}|B\rangle
\right|^2.
\] The overlap \[
\langle\alpha_{\mathrm{tw,R}}|B\rangle
\] is the disk one-point function of that R--R field in the presence of
the D-brane. A nonzero R--R one-point function is the string-theory
definition of a source coupling. Therefore \[
\boxed{
\widetilde{\mathcal A}[I_4]_{\mathrm R}\neq0
\quad\Longrightarrow\quad
\langle\alpha_{\mathrm{tw,R}}|D7_x\rangle\neq0
\quad\Longrightarrow\quad
D7_x\text{ carries twisted R--R charge}.
}
\]

This is visible directly in the boundary state. Suppressing
normalization, \[
|D7_x^\pm\rangle_{\mathrm{frac}}
=
|B\rangle_{\mathrm{untw}}
\pm
|B\rangle_{\mathrm{tw}}.
\] In a self-overlap, the twisted part occurs quadratically: \[
\langle D7_x^\pm|
e^{-2\pi lH}
|D7_x^\pm\rangle
\supset
\langle B_{\mathrm{tw}}|
e^{-2\pi lH}
|B_{\mathrm{tw}}\rangle,
\] so the sign of the twisted charge is invisible. To measure its sign,
compare opposite fractional branes: \[
\langle D7_x^+|
e^{-2\pi lH}
|D7_x^-\rangle
\supset
-
\langle B_{\mathrm{tw}}|
e^{-2\pi lH}
|B_{\mathrm{tw}}\rangle.
\] This is why the book obtains equation (15.37) with the sign of the
\(I_4\) term reversed for branes of opposite twisted charge.

\paragraph{\texorpdfstring{Why the twisted field is
\(\int_{E_i}C_8\)}{Why the twisted field is \textbackslash int\_\{E\_i\}C\_8}}\label{why-the-twisted-field-is-int_e_ic_8}

The partition function establishes that the exchanged state is
localized, twisted, and R--R. Identifying the spacetime differential
form requires the geometric dictionary for the resolution of \[
T^4/\mathbb Z_2.
\] Near each fixed point, the singularity is locally \[
\mathbb C^2/\mathbb Z_2.
\] Resolving it introduces an exceptional cycle \[
E_i\simeq S^2
\] whose area tends to zero at the orbifold point. Associated with
\(E_i\) is a localized harmonic two-form \(\omega_i\) satisfying \[
\int_{E_j}\omega_i
=
\delta_{ij}.
\]

In the democratic formulation of type-IIB supergravity, expand the R--R
eight-form as \[
C_8(x,y)
=
\sum_{i=1}^{16}
C_6^{(i)}(x)\wedge\omega_i(y)
+
\cdots.
\] Integrating over an exceptional cycle gives \[
\boxed{
C_6^{(i)}(x)
=
\int_{E_i}C_8(x,y).
}
\] Because \(\omega_i\) is localized near the \(i\)th exceptional cycle,
\(C_6^{(i)}\) is the spacetime field associated with the \(i\)th twisted
R--R sector. Although \(E_i\) has zero area at the orbifold point, it
remains a vanishing homology cycle; its twisted state and R--R flux
integral do not disappear.

The D7 Wess--Zumino coupling is \[
S_{\mathrm{WZ}}
=
\mu_7\int_{W_8}C_8
+
\cdots.
\] For a fractional D7 whose wrapped cycle contains an exceptional
component at \(E_i\), the corresponding part of its worldvolume is \[
W_8
\supset
\mathbb R^{1,5}\times E_i.
\] Substituting the expansion of \(C_8\) gives \[
\begin{aligned}
\mu_7\int_{\mathbb R^{1,5}\times E_i}C_8
&=
\mu_7\int_{\mathbb R^{1,5}}
\left(
\int_{E_i}C_8
\right)\\
&=
\mu_7\int_{\mathbb R^{1,5}}C_6^{(i)}.
\end{aligned}
\] Thus the exceptional component of the fractional brane is
electrically charged under \[
\boxed{
C_6^{(i)}
=
\int_{E_i}C_8.
}
\]

The \(D7_x\) wraps the \((x^1,x^2)\) torus and passes through four
\(I_4\) fixed points. This is why equation (15.39) contains \[
\sum_{i=1}^{4}|D7_x,i\rangle_{\mathrm{tw}}
\] and why the brane couples to the four corresponding fields \[
C_6^{(i)}
=
\int_{E_i}C_8.
\] The two choices \(D7_x^\pm\) have opposite coefficients for all these
twisted components. Their sum has no exceptional charge and can become
an ordinary movable bulk D7-brane.

\paragraph{What is and is not determined by the partition
function}\label{what-is-and-is-not-determined-by-the-partition-function}

From equations (15.37)--(15.38) alone one can read:

\begin{itemize}
\tightlist
\item
  the presence of an \(I_4\)-twisted closed-string exchange;
\item
  its NS--NS and R--R parts;
\item
  the absence of internal zero modes, hence localization at fixed
  points;
\item
  a nonzero twisted R--R boundary-state coupling;
\item
  the relative sign of the twisted charge by comparing \(D7_x^+\) and
  \(D7_x^-\) amplitudes.
\end{itemize}

The further identification \[
\text{twisted R--R state}
\quad\longleftrightarrow\quad
\int_{E_i}C_8
\] uses the geometric dictionary between orbifold twisted sectors and
exceptional cycles, together with the D7 Wess--Zumino coupling. Thus the
partition function proves that the charge is twisted R--R charge, while
the resolved-orbifold geometry identifies the R--R potential that
represents it.

\subsubsection{Derivation of the orientifold image in equation
(15.41)}\label{derivation-of-the-orientifold-image-in-equation-15.41}

Equation (15.41) states \[
\Omega R:
\qquad
|D7^\pm\rangle_{\mathrm{frac}}
\longmapsto
|D7^\mp\rangle_{\mathrm{frac}}.
\] Strictly, this is the action induced by the full orientifold
generator \[
\mathcal O
=
\Omega R(-1)^{F_L}
\] on the GSO-complete boundary state; the book suppresses
\((-1)^{F_L}\) in the label on the left-hand side of equation (15.41).

The result is not determined by the geometry of \(R\) alone. It follows
from combining:

\begin{enumerate}
\def\labelenumi{\arabic{enumi}.}
\tightlist
\item
  geometric invariance of the D7-brane under \(R\);
\item
  the decomposition of a fractional boundary state into untwisted and
  twisted pieces;
\item
  the extra sign chosen for the orientifold action on \(I_4\)-twisted
  states on page 536.
\end{enumerate}

\paragraph{Action on the untwisted
component}\label{action-on-the-untwisted-component}

For a \(D7_x\) at one of the fixed positions \[
(y^1,y^2)
=
(-y^1,-y^2)
\pmod{2\pi R_y},
\] the reflection \(R\) maps the geometric brane to itself. Worldsheet
parity exchanges the two ends of the open string, but a D7--D7 boundary
condition is unchanged. The factor \((-1)^{F_L}\) supplies the action
required for this supersymmetric type-IIB orientifold and does not
convert this D7 into an anti-D7. With the common overall boundary-state
phase fixed conventionally, \[
\boxed{
\mathcal O
|D7_x\rangle_{\mathrm{untw}}
=
+|D7_x\rangle_{\mathrm{untw}}.
}
\]

If a brane were not at an \(R\)-fixed position, the same equation would
instead contain the boundary state of its geometrically reflected image.
Equation (15.41) concerns the invariant fractional branes used in the
construction.

\paragraph{Action on the twisted
component}\label{action-on-the-twisted-component}

On page 536 the model makes the discrete choice that \(\mathcal O\) acts
with an additional minus sign on every \(I_4\)-twisted-sector state. In
the notation used earlier, \[
\epsilon_{\mathrm{tw}}
=
-1.
\] The twisted boundary state \(|D7_x,i\rangle_{\mathrm{tw}}\) in
equation (15.40) is a coherent state built on the localized twisted
ground state \(|t_i\rangle\). The oscillator part is carried into the
corresponding boundary-state oscillator part by \(\Omega R\). The
additional twist-field sign acts on its ground state and gives \[
\mathcal O
|D7_x,i\rangle_{\mathrm{tw}}
=
\epsilon_{\mathrm{tw}}\,
|D7_x,R(i)\rangle_{\mathrm{tw}}.
\] For the four fixed points lying on an invariant \(D7_x\), \(R(i)=i\)
modulo the torus lattice, so \[
\boxed{
\mathcal O
|D7_x,i\rangle_{\mathrm{tw}}
=
-|D7_x,i\rangle_{\mathrm{tw}}.
}
\]

More generally, if \(R\) permuted fixed points, the orientifold would
apply the same minus sign while permuting the corresponding twisted
boundary states. Here the relevant four fixed points are individually
fixed.

\paragraph{Acting on the full fractional boundary
state}\label{acting-on-the-full-fractional-boundary-state}

Write equation (15.39) as \[
|D7_x^\pm\rangle_{\mathrm{frac}}
=
A\,|D7_x\rangle_{\mathrm{untw}}
\pm
B\sum_{i=1}^{4}|D7_x,i\rangle_{\mathrm{tw}},
\] where \[
A
=
\frac{1}{2\mathcal N_{D7}}
\frac{\rho_x}{\rho_y},
\qquad
B
=
\frac{1}{\mathcal N_{D7}}.
\] Using the two transformation laws above, \[
\begin{aligned}
\mathcal O
|D7_x^\pm\rangle_{\mathrm{frac}}
&=
A\,\mathcal O|D7_x\rangle_{\mathrm{untw}}
\pm
B\sum_{i=1}^{4}
\mathcal O|D7_x,i\rangle_{\mathrm{tw}}\\
&=
A\,|D7_x\rangle_{\mathrm{untw}}
\mp
B\sum_{i=1}^{4}
|D7_x,i\rangle_{\mathrm{tw}}\\
&=
|D7_x^\mp\rangle_{\mathrm{frac}}.
\end{aligned}
\] Therefore \[
\boxed{
\mathcal O:
\qquad
|D7_x^+\rangle_{\mathrm{frac}}
\longleftrightarrow
|D7_x^-\rangle_{\mathrm{frac}}.
}
\] The identical reasoning applies to the \(D7_y^\pm\) pair.

\paragraph{The same result from the twisted R--R
coupling}\label{the-same-result-from-the-twisted-rr-coupling}

The relevant Wess--Zumino couplings have the schematic form \[
\begin{aligned}
S_{\mathrm{WZ}}^\pm
&=
\mu_{\mathrm{ut}}
\int_{\mathbb R^{1,5}\times\Pi_x}
C_8^{\mathrm{untw}}\\
&\quad
\pm
\mu_{\mathrm{tw}}
\sum_{i=1}^{4}
\int_{\mathbb R^{1,5}}
C_6^{(i)},
\qquad
C_6^{(i)}
=
\int_{E_i}C_8.
\end{aligned}
\] The chosen orientifold parity leaves the relevant untwisted source
component invariant but reverses the twisted one: \[
\mathcal O:
\qquad
C_8^{\mathrm{untw}}
\mapsto
+C_8^{\mathrm{untw}},
\qquad
C_6^{(i)}
\mapsto
-C_6^{(i)}.
\] A source with \(q_{\mathrm{tw}}=+1\) must therefore be mapped to a
source with \(q_{\mathrm{tw}}=-1\), and conversely. This reproduces the
boundary-state derivation.

This also clarifies the terminology: \(D7^+\) and \(D7^-\) are parallel
and have the same untwisted D7 charge and tension. They differ only in
the sign of their localized twisted charge.

\paragraph{Why the pair is required and why its twisted tadpole
cancels}\label{why-the-pair-is-required-and-why-its-twisted-tadpole-cancels}

An orientifold-invariant D-brane configuration must contain complete
orbits of \(\mathcal O\). Since neither fractional brane is invariant by
itself, the minimal invariant combination is \[
|D7_x^+\rangle_{\mathrm{frac}}
+
|D7_x^-\rangle_{\mathrm{frac}}.
\] Its twisted component cancels: \[
\begin{aligned}
|D7_x^+\rangle_{\mathrm{frac}}
+
|D7_x^-\rangle_{\mathrm{frac}}
&=
2A\,|D7_x\rangle_{\mathrm{untw}}\\
&\quad
+
B\sum_i|D7_x,i\rangle_{\mathrm{tw}}
-
B\sum_i|D7_x,i\rangle_{\mathrm{tw}}\\
&=
2A\,|D7_x\rangle_{\mathrm{untw}}.
\end{aligned}
\] Consequently, \[
q_{\mathrm{tw}}^{(+)}
+
q_{\mathrm{tw}}^{(-)}
=
+1-1
=
0.
\] This is why the book says that introducing the fractional branes in
orientifold-image pairs automatically cancels their twisted-sector
tadpoles.

\paragraph{Dependence on the discrete orientifold
choice}\label{dependence-on-the-discrete-orientifold-choice}

The exchange in equation (15.41) is not universal for every orientifold
of \(T^4/\mathbb Z_2\). If one chooses the opposite twisted-sector
action, \[
\epsilon_{\mathrm{tw}}
=
+1,
\] then \[
\mathcal O
|D7_x,i\rangle_{\mathrm{tw}}
=
+|D7_x,i\rangle_{\mathrm{tw}},
\] and the same calculation gives \[
\mathcal O
|D7_x^\pm\rangle_{\mathrm{frac}}
=
|D7_x^\pm\rangle_{\mathrm{frac}}.
\] Each fractional brane is then separately invariant. This is precisely
the alternative model mentioned below equation (15.41), whose closed
twisted sectors contain tensor multiplets instead of hypermultiplets.

Thus equation (15.41) cannot be inferred from the annulus amplitude
alone. The annulus establishes the existence of the two signs of twisted
charge. Determining how \(\Omega R(-1)^{F_L}\) maps those signs
additionally requires the orientifold action on the twisted
closed-string ground states, and the model fixed that action by the
extra minus sign on page 536.

\subsubsection{Checks}\label{checks-22}

At the massless NS level, \[
N_{\mathrm{before}}
=
4_{\mathrm{vector}}
+
2_{\mathrm{Wilson}}
+
2_{\mathrm{position}}
=
8,
\] while after the \(I_4\) projection, \[
N_{\mathrm{after}}
=
4_{\mathrm{vector}}.
\] This is the correct bosonic content of a six-dimensional massless
vector multiplet once the R-sector gauginos are included.

The trace with \(I_4\) has no compact lattice sums because only \[
m_a=n_a=0
\] is fixed by sign reversal. Under the annulus modular transformation,
precisely this inserted trace becomes the twisted closed-string
exchange. This matches the general orbifold rule that an insertion in
the open loop channel becomes propagation in the corresponding twisted
closed sector in the transverse channel.

Finally, changing the fixed position from \((0,0)\) to any of \[
(\pi R_y,0),
\qquad
(0,\pi R_y),
\qquad
(\pi R_y,\pi R_y)
\] does not change the local oscillator projection. It only changes
which fixed locus carries the fractional-brane boundary state.

\subsubsection{Result}\label{result-26}

The trace decomposes as \[
\boxed{
\mathcal A
=
\frac12\mathcal A[1]
+
\frac12\mathcal A[I_4],
}
\] with \[
\boxed{
\begin{aligned}
\mathcal A[1]
&:
\quad
P_x^2W_y^2\,
(1_{\mathrm{NS}}-1_{\mathrm R})
\frac{\vartheta_4^4}{\eta^{12}},\\
\mathcal A[I_4]
&:
\quad
(1_{\mathrm{NS}}-1_{\mathrm R})
4\frac{\vartheta_4^2\vartheta_3^2}
{\eta^6\vartheta_2^2}.
\end{aligned}
}
\] The \(I_4\) trace is possible only for \[
(y^1,y^2)
=
(-y^1,-y^2)
\pmod{2\pi R_y},
\] and its compact zero modes obey \[
m_1=m_2=n_1=n_2=0.
\]

At the massless NS level, \[
\boxed{
\begin{aligned}
\frac{1+I_4}{2}
\psi_{-1/2}^{\mu}|0\rangle_{\mathrm{NS}}
&=
\psi_{-1/2}^{\mu}|0\rangle_{\mathrm{NS}}
&&\Longrightarrow A_\mu,\\
\frac{1+I_4}{2}
\psi_{-1/2}^{x^a}|0\rangle_{\mathrm{NS}}
&=
0
&&\Longrightarrow\text{no continuous Wilson line},\\
\frac{1+I_4}{2}
\psi_{-1/2}^{y^a}|0\rangle_{\mathrm{NS}}
&=
0
&&\Longrightarrow\text{no transverse-position modulus}.
\end{aligned}
}
\] The transverse-channel interpretation is \[
\boxed{
\mathcal A[1]
\longleftrightarrow
\text{untwisted closed-string exchange},
\qquad
\mathcal A[I_4]
\longleftrightarrow
\text{twisted closed-string exchange}.
}
\] Separating the R contribution and factorizing the long tube gives \[
\boxed{
\widetilde{\mathcal A}[I_4]_{\mathrm R}
\sim
\sum_i
\left|
\langle C_6^{(i)}|D7_x\rangle
\right|^2,
\qquad
C_6^{(i)}
=
\int_{E_i}C_8.
}
\] Finally, the model's twisted-sector sign gives \[
\boxed{
\begin{aligned}
\mathcal O|B\rangle_{\mathrm{untw}}
&=
+|B\rangle_{\mathrm{untw}},\\
\mathcal O|B\rangle_{\mathrm{tw}}
&=
-|B\rangle_{\mathrm{tw}},
\end{aligned}
\qquad
\Longrightarrow
\qquad
\mathcal O|D7^\pm\rangle_{\mathrm{frac}}
=
|D7^\mp\rangle_{\mathrm{frac}}.
}
\] Hence an invariant pair has \[
\boxed{
q_{\mathrm{tw}}^{(+)}
+
q_{\mathrm{tw}}^{(-)}
=
0,
}
\] which is equation (15.41) and its twisted-tadpole consequence.

Therefore the surviving object supports a six-dimensional gauge
multiplet and is an \(I_4\)-charged fractional D7-brane pinned at the
chosen orbifold fixed position.

\section{Q7: Deriving the vertex-operator mass formula from conformal
weights}\label{q7-deriving-the-vertex-operator-mass-formula-from-conformal-weights}

How to derive the mass formula for a vertex operator? For example, in
the chapter15.4, how to derive eq (15.52) for the vertex operator given
in (15.51)?

\begin{quote}
\textbf{Textbook location and related material.} Chapter 15, Section
15.4 around Eqs. (15.51)--(15.52): conformal weights of the general
compactified vertex. Section 13.2 supplies the BRST and picture
conventions.
\end{quote}

\subsection{Answer}\label{answer-27}

\subsubsection{Core point}\label{core-point-27}

The mass formula is the conformal-weight condition for a physical string
vertex operator. The vertex written in equation (15.51) contains the
matter and superghost factors but does not display the reparametrization
ghosts \(c\bar c\). Therefore it must have weights \[
(h,\bar h)=(1,1).
\] Equivalently, the unintegrated BRST vertex \(c\bar c\,V\) has total
weights \((0,0)\), because \(c\) and \(\bar c\) each have weight \(-1\).

In the convention used on pages 545--546, the \(z\) sector is called
\textbf{right-moving} and the \(\bar z\) sector \textbf{left-moving}.
This is why \(h_{\mathrm{int}}\) appears in the \(m_R^2\) equation,
while \(\bar h_{\mathrm{int}}\) appears in the \(m_L^2\) equation.

Suppressing normal-ordering symbols, equation (15.51) is \[
\begin{aligned}
V(z,\bar z)
={}&
\bigl(\bar\partial^{n_L}X^\mu(\bar z)\bigr)^{m_L}
\bigl(\partial^{n_{1R}}X^\nu(z)\bigr)^{m_{1R}}
\bigl(\partial^{n_{2R}}\phi^i(z)\bigr)^{m_{2R}}
\\
&\times
e^{i\boldsymbol\lambda_R\cdot\boldsymbol\phi(z)}
e^{q\phi(z)}
V_{\mathrm{int}}(z,\bar z)
e^{ik\cdot X(z,\bar z)}.
\end{aligned}
\] Here the indexed fields \(\phi^i\) bosonize the four external
right-moving fermions, whereas the unindexed \(\phi\) bosonizes the
\(\beta\gamma\) superghost. The derivation consists of adding the
weights of all these factors and setting the result equal to \((1,1)\).

\subsubsection{The plane-wave
contribution}\label{the-plane-wave-contribution}

With the free-boson OPE \[
X^\mu(z,\bar z)X^\nu(w,\bar w)
\sim
-\frac{\alpha'}2\eta^{\mu\nu}
\bigl[\log(z-w)+\log(\bar z-\bar w)\bigr],
\] the stress tensor gives \[
T_X(z)e^{ik\cdot X(w,\bar w)}
\sim
\frac{\alpha'k^2/4}{(z-w)^2}e^{ik\cdot X(w,\bar w)}
+\frac1{z-w}\partial e^{ik\cdot X(w,\bar w)}.
\] Thus \[
h\bigl(e^{ik\cdot X}\bigr)
=
\bar h\bigl(e^{ik\cdot X}\bigr)
=
\frac{\alpha'k^2}{4}.
\] For a state of spacetime mass \(m\), the mostly-plus convention gives
\[
k^2=-m^2,
\] and hence the plane wave contributes \[
\boxed{
h_k=\bar h_k=-\frac{\alpha'm^2}{4}.
}
\] This negative contribution is not pathological: the time component of
the Lorentzian free boson has negative norm in target-space signature,
and the physical-state condition balances it against positive oscillator
and internal-CFT weights.

\subsubsection{Oscillator and internal-CFT
contributions}\label{oscillator-and-internal-cft-contributions}

A chiral free boson has \(h(\partial X)=1\). More generally, \[
h(\partial^nX)=n,
\qquad
\bar h(\bar\partial^nX)=n.
\] Therefore the three oscillator blocks in (15.51) contribute \[
N_L=m_Ln_L,
\qquad
N_{1R}=m_{1R}n_{1R},
\qquad
N_{2R}=m_{2R}n_{2R}
\] for the displayed monomials. For a general product of oscillators,
these become sums, for example \[
N_L=\sum_a m_{L,a}n_{L,a}.
\] The internal operator is defined to have weights \[
V_{\mathrm{int}}(z,\bar z):\qquad
(h,\bar h)=(h_{\mathrm{int}},\bar h_{\mathrm{int}}).
\] It therefore contributes \(h_{\mathrm{int}}\) to the right-moving
\(z\) equation and \(\bar h_{\mathrm{int}}\) to the left-moving
\(\bar z\) equation.

\subsubsection{The bosonized-fermion
exponential}\label{the-bosonized-fermion-exponential}

For the canonically normalized bosons representing the external
right-moving fermions, \[
\phi^i(z)\phi^j(w)
\sim
-\delta^{ij}\log(z-w).
\] Their stress tensor is \[
T_{\boldsymbol\phi}(z)
=
-\frac12:\!\partial\boldsymbol\phi\cdot
\partial\boldsymbol\phi\!:(z).
\] The standard exponential OPE then gives \[
T_{\boldsymbol\phi}(z)
e^{i\boldsymbol\lambda_R\cdot\boldsymbol\phi(w)}
\sim
\frac{\boldsymbol\lambda_R^2/2}{(z-w)^2}
e^{i\boldsymbol\lambda_R\cdot\boldsymbol\phi(w)}
+\frac1{z-w}\partial
e^{i\boldsymbol\lambda_R\cdot\boldsymbol\phi(w)}.
\] Consequently, \[
\boxed{
h\left(e^{i\boldsymbol\lambda_R\cdot\boldsymbol\phi}\right)
=\frac12\boldsymbol\lambda_R^2.
}
\] This factor encodes whether the external right-moving state lies in
an NS conjugacy class or a Ramond spinor conjugacy class; together with
the picture charge \(q\), \(\boldsymbol\lambda_R\) forms the
covariant-lattice vector \((\boldsymbol\lambda_R,q)\in(D_{2,1})_R\).

\subsubsection{The superghost-picture
contribution}\label{the-superghost-picture-contribution}

The bosonized superghost is not an ordinary free boson: it has a
background charge. With \[
\beta=e^{-\phi}\partial\xi,
\qquad
\gamma=\eta e^\phi,
\] its scalar stress tensor is \[
T_\phi(z)
=
-\frac12:\!(\partial\phi)^2\!:(z)
-\partial^2\phi(z).
\] Since \[
\phi(z)\phi(w)\sim-\log(z-w),
\] the double contraction with the first term contributes \(-q^2/2\),
while the background-charge term contributes \(-q\). Hence \[
T_\phi(z)e^{q\phi(w)}
\sim
\frac{-\frac12q^2-q}{(z-w)^2}e^{q\phi(w)}
+\frac1{z-w}\partial e^{q\phi(w)},
\] so that \[
\boxed{
h(e^{q\phi})
=
-\frac12q(q+2)
=
-\frac12q^2-q.
}
\] For example, \[
h(e^{-\phi})=\frac12,
\qquad
h(e^{-\phi/2})=\frac38,
\] as required in the conventional NS picture \(q=-1\) and Ramond
picture \(q=-\frac12\).

\subsubsection{The left-moving weight
condition}\label{the-left-moving-weight-condition}

Only three pieces of (15.51) contribute nontrivially to the
antiholomorphic weight:

\begin{itemize}
\tightlist
\item
  the left-moving bosonic oscillators, with weight \(N_L\);
\item
  the internal operator, with weight \(\bar h_{\mathrm{int}}\);
\item
  the plane wave, with weight \(-\alpha'm^2/4\).
\end{itemize}

Therefore \[
\bar h(V)
=
N_L+\bar h_{\mathrm{int}}
-\frac{\alpha'm^2}{4}.
\] The physical condition \(\bar h(V)=1\) gives \[
\frac{\alpha'm^2}{4}
=
N_L+\bar h_{\mathrm{int}}-1.
\] The book defines the left-moving contribution \(m_L^2\) by \[
\frac{\alpha'}2m_L^2
\equiv
N_L+\bar h_{\mathrm{int}}-1.
\] Thus the first line of (15.52) is \[
\boxed{
\frac{\alpha'}2m_L^2
=N_L+\bar h_{\mathrm{int}}-1.
}
\] The constant \(-1\) is simply the result of moving the required
vertex weight \(1\) to the other side. In oscillator language it is the
left-moving bosonic intercept.

\subsubsection{The right-moving weight
condition}\label{the-right-moving-weight-condition}

Adding all holomorphic weights gives \[
\begin{aligned}
h(V)
={}&
N_{1R}+N_{2R}
+\frac12\boldsymbol\lambda_R^2
-\frac12q^2-q
+h_{\mathrm{int}}
-\frac{\alpha'm^2}{4}.
\end{aligned}
\] Imposing \(h(V)=1\) yields \[
\frac{\alpha'm^2}{4}
=
\frac12\boldsymbol\lambda_R^2
-\frac12q^2-q
+N_{1R}+N_{2R}
+h_{\mathrm{int}}-1.
\] Defining the right-moving contribution \(m_R^2\) by \[
\frac{\alpha'}2m_R^2
\equiv
\frac12\boldsymbol\lambda_R^2
-\frac12q^2-q
+N_{1R}+N_{2R}
+h_{\mathrm{int}}-1,
\] we obtain the second line of (15.52): \[
\boxed{
\frac{\alpha'}2m_R^2
=
\frac12\boldsymbol\lambda_R^2
-\frac12q^2-q
+N_{1R}+N_{2R}
+h_{\mathrm{int}}-1.
}
\] The familiar NS and R zero-point shifts are encoded here in a
picture-covariant way: part of the right-moving conformal weight resides
in the lattice exponential and part in \(e^{q\phi}\).

\subsubsection{\texorpdfstring{Why the book writes
\(m^2=m_L^2+m_R^2\)}{Why the book writes m\^{}2=m\_L\^{}2+m\_R\^{}2}}\label{why-the-book-writes-m2m_l2m_r2}

There is only one spacetime momentum \(k^\mu\) and one physical mass
\(m\). The symbols \(m_L^2\) and \(m_R^2\) are therefore not two
independent masses. They are a bookkeeping convention that assigns half
of the total mass formula to each chiral sector.

Indeed, either physical weight condition gives \[
N_L+\bar h_{\mathrm{int}}-1
=
\frac{\alpha'm^2}{4},
\] and \[
\frac12\boldsymbol\lambda_R^2
-\frac12q^2-q
+N_{1R}+N_{2R}
+h_{\mathrm{int}}-1
=
\frac{\alpha'm^2}{4}.
\] Comparing these equations with the definitions of \(m_L^2\) and
\(m_R^2\) gives \[
m_L^2=m_R^2=\frac{m^2}{2}.
\] Consequently, \[
\boxed{m^2=m_L^2+m_R^2,}
\] which is the third line of (15.52).

The equality of the two chiral expressions is also the level-matching
condition: \[
\boxed{
N_L+\bar h_{\mathrm{int}}
=
\frac12\boldsymbol\lambda_R^2
-\frac12q^2-q
+N_{1R}+N_{2R}
+h_{\mathrm{int}}.
}
\] Thus one should not first compute two unrelated masses and then add
them. One computes the two chiral conformal weights, imposes level
matching, and uses the book's definitions to package the common result
as \(m^2=m_L^2+m_R^2\).

\subsubsection{Checks}\label{checks-23}

For a massless NS state in the canonical picture \(q=-1\), a
right-moving fermion \(\psi^\mu\) is represented by a lattice
exponential with \[
\frac12\boldsymbol\lambda_R^2=\frac12.
\] The superghost supplies \[
-\frac12q^2-q
=
-\frac12+1
=
\frac12.
\] With \(N_{1R}=N_{2R}=h_{\mathrm{int}}=0\), the right-moving
non-plane-wave weight is therefore \(1\), and \[
m_R^2=0.
\] On the left, a factor \(\bar\partial X^\mu\) has \(N_L=1\) and
\(\bar h_{\mathrm{int}}=0\), giving \(m_L^2=0\). This is the standard
massless heterotic vertex.

For a four-dimensional Ramond state in picture \(q=-\frac12\), the
external spin field has \[
\boldsymbol\lambda_R
=
\left(\pm\frac12,\pm\frac12\right),
\qquad
\frac12\boldsymbol\lambda_R^2=\frac14.
\] The superghost contributes \(3/8\). If the internal operator is an
internal Ramond ground state of the \(c_{\mathrm{int}}=9\) SCFT, then \[
h_{\mathrm{int}}=\frac{c_{\mathrm{int}}}{24}=\frac38.
\] Hence \[
\frac14+\frac38+\frac38=1,
\] again giving \(m_R^2=0\).

Picture changing alters \(q\) and simultaneously alters the other
right-moving factors, but it preserves their total conformal weight. The
mass formula is therefore picture independent even though the separate
terms in the second line of (15.52) are picture dependent.

\subsubsection{Result}\label{result-27}

The conformal weights of the factors in (15.51) are \[
\begin{array}{c|c|c}
\text{factor}&h&\bar h\\ \hline
e^{ik\cdot X}&-\alpha'm^2/4&-\alpha'm^2/4\\
(\bar\partial^{n_L}X)^{m_L}&0&N_L\\
(\partial^{n_{1R}}X)^{m_{1R}}&N_{1R}&0\\
(\partial^{n_{2R}}\phi^i)^{m_{2R}}&N_{2R}&0\\
e^{i\boldsymbol\lambda_R\cdot\boldsymbol\phi}&\boldsymbol\lambda_R^2/2&0\\
e^{q\phi}&-q^2/2-q&0\\
V_{\mathrm{int}}&h_{\mathrm{int}}&\bar h_{\mathrm{int}}
\end{array}
\] and the physical-state condition is \[
h(V)=\bar h(V)=1.
\] It follows that \[
\boxed{
\begin{aligned}
\frac{\alpha'}2m_L^2
&=N_L+\bar h_{\mathrm{int}}-1,\\
\frac{\alpha'}2m_R^2
&=\frac12\boldsymbol\lambda_R^2
-\frac12q^2-q
+N_{1R}+N_{2R}
+h_{\mathrm{int}}-1,\\
m^2&=m_L^2+m_R^2.
\end{aligned}
}
\] For a physical level-matched closed-string vertex, \[
\boxed{
m_L^2=m_R^2=\frac{m^2}{2}.
}
\]

\section{Q8: Momentum dependence of R--R vertices and derivative
couplings}\label{q8-momentum-dependence-of-rr-vertices-and-derivative-couplings}

In page553-554 the discussion of type IIB (R, R) vertex operators. It is
said that \textgreater The explicit k-dependence of the (R, R) vertex
operators indicates that the corresponding massless fields have only
derivative couplings in the lowenergy effective action. In other words
they only appear through their field strength which implies that there
are e.g.~no fields charged with respect to the graviphoton.

Why the k-dependence of these vertex operators indicates this?

\begin{quote}
\textbf{Textbook location and related material.} Chapter 15, Section
15.4 around pp.~553--554 and Eq. (15.84): type-IIB R--R vertices and
their explicit momentum factors. Section 13.2 gives the ten-dimensional
R--R field-strength bispinor construction.
\end{quote}

\subsection{Answer}\label{answer-28}

\subsubsection{Core point}\label{core-point-28}

The relevant \(k\)-dependence is \textbf{not} the plane wave
\(e^{ik\cdot X}\), since every momentum-eigenstate vertex operator
contains that factor. The important feature of equation (15.84) is the
additional momentum multiplying the polarization: \[
V^A\propto\epsilon_\mu k_\nu(\text{spin fields})\,e^{ik\cdot X},
\qquad
V^C\propto k_\mu(\text{spin fields})\,e^{ik\cdot X},
\qquad
V^{\bar C}\propto k_\mu(\text{spin fields})\,e^{ik\cdot X}.
\] In Fourier space, multiplication by \(ik_\mu\) is differentiation in
spacetime: \[
\partial_\mu C(x)
=
\int\frac{d^4k}{(2\pi)^4}\,
ik_\mu C(k)e^{ik\cdot x}.
\] Consequently, \(k_\mu C(k)\) represents the one-form field strength
\(dC\), while the antisymmetric combination of \(\epsilon_\mu k_\nu\)
represents the graviphoton field strength
\(F_{\mu\nu}=2\partial_{[\mu}A_{\nu]}\).

This is not an accidental feature of the chosen vertices. In the RNS
formulation, a BRST-invariant R--R vertex is naturally labeled by an
R--R \textbf{field-strength bispinor}, not directly by an R--R gauge
potential. Equation (15.84) is the four-dimensional decomposition of
that general statement.

\subsubsection{The extra momentum in equation
(15.84)}\label{the-extra-momentum-in-equation-15.84}

Suppressing the common superghost and plane-wave factors, the type-IIB
vertices in (15.84) have the structure \[
\begin{aligned}
V^A
={}&
\epsilon_\mu k_\nu
\left[
\bar S^\alpha\bar\Sigma\,
\sigma^{\mu\nu}_{\alpha}{}^\beta
S_\beta\Sigma
+
\bar S_{\dot\alpha}\bar\Sigma^\dagger\,
\bar\sigma^{\mu\nu\dot\alpha}{}_{\dot\beta}
S^{\dot\beta}\Sigma^\dagger
\right]
e^{-\frac12(\phi+\bar\phi)}e^{ik\cdot X},
\\
V^C
={}&
k_\mu\,
\bar S_{\dot\alpha}\bar\Sigma^\dagger
\bar\sigma^{\mu\dot\alpha\beta}
S_\beta\Sigma\,
e^{-\frac12(\phi+\bar\phi)}e^{ik\cdot X},
\\
V^{\bar C}
={}&
k_\mu\,
\bar S^\alpha\bar\Sigma
\sigma^\mu_{\alpha\dot\beta}
S^{\dot\beta}\Sigma^\dagger\,
e^{-\frac12(\phi+\bar\phi)}e^{ik\cdot X}.
\end{aligned}
\] The spin fields and the internal spectral-flow operators determine
the Lorentz representation and the internal R--R state. The factors
\(\epsilon_\mu k_\nu\) and \(k_\mu\) determine which gauge-invariant
spacetime tensor serves as the physical polarization.

For comparison, writing only \[
\epsilon_\mu(\text{spin fields})\,e^{ik\cdot X}
\] would describe the potential polarization itself. The book emphasizes
that such a potential-level expression is not what appears in the
BRST-invariant R--R vertex.

\subsubsection{The graviphoton vertex is a field-strength
vertex}\label{the-graviphoton-vertex-is-a-field-strength-vertex}

The Lorentz generators appearing in the spin-field bilinears are
antisymmetric: \[
\sigma^{\mu\nu}=-\sigma^{\nu\mu},
\qquad
\bar\sigma^{\mu\nu}=-\bar\sigma^{\nu\mu}.
\] Therefore only the antisymmetric part of \(\epsilon_\mu k_\nu\)
contributes: \[
\epsilon_\mu k_\nu\sigma^{\mu\nu}
=
\frac12
\left(
\epsilon_\mu k_\nu-\epsilon_\nu k_\mu
\right)\sigma^{\mu\nu}.
\] With \[
F_{\mu\nu}(k)
=
i\left(k_\mu\epsilon_\nu-k_\nu\epsilon_\mu\right),
\] we obtain, up to the conventional Fourier factor, \[
\boxed{
\epsilon_\mu k_\nu\sigma^{\mu\nu}
\ \propto\
F_{\mu\nu}(k)\sigma^{\mu\nu},
\qquad
\epsilon_\mu k_\nu\bar\sigma^{\mu\nu}
\ \propto\
F_{\mu\nu}(k)\bar\sigma^{\mu\nu}.
}
\] Thus \(V^A\) is really a vertex for the two chiral spinor components
of \(F_{\mu\nu}\), not for a bare \(A_\mu\).

In four-dimensional Lorentzian signature, \(\sigma^{\mu\nu}\) and
\(\bar\sigma^{\mu\nu}\) select the two duality components \[
F^\pm_{\mu\nu}
=
\frac12\left(F_{\mu\nu}\pm i\,{*F}_{\mu\nu}\right).
\] This is the reason for the two terms in \(V^A\): as the book states,
they encode the self-dual and anti-self-dual parts of the graviphoton
field strength.

The vertex is also manifestly invariant under the linearized gauge
transformation \[
\epsilon_\mu\longrightarrow\epsilon_\mu+a\,k_\mu,
\] because \[
k_\mu k_\nu\sigma^{\mu\nu}
=
k_\mu k_\nu\bar\sigma^{\mu\nu}
=0.
\] This is exactly the momentum-space version of \[
A_\mu\longrightarrow A_\mu+\partial_\mu\Lambda,
\qquad
F_{\mu\nu}\longrightarrow F_{\mu\nu}.
\]

\subsubsection{\texorpdfstring{The scalar vertices contain \(dC\) and
\(d\bar C\)}{The scalar vertices contain dC and d\textbackslash bar C}}\label{the-scalar-vertices-contain-dc-and-dbar-c}

Equation (15.84) suppresses the scalar wavefunction, effectively setting
its polarization to one. Restoring the Fourier coefficient \(C(k)\), the
prefactor of the vertex is \(k_\mu C(k)\).

For a scalar mode, \[
C(x)
=
\int\frac{d^4k}{(2\pi)^4}\,
C(k)e^{ik\cdot x},
\] so \[
k_\mu C(k)e^{ik\cdot x}
=
\frac1i\,\partial_\mu
\left(C(k)e^{ik\cdot x}\right).
\] The tensor polarizations in the two scalar vertices are therefore \[
\mathcal F_{\mu}^{C}(k)\propto k_\mu C(k),
\qquad
\mathcal F_{\mu}^{\bar C}(k)\propto k_\mu\bar C(k).
\] In position space, \[
\boxed{
\mathcal F_1^C=dC,
\qquad
\mathcal F_1^{\bar C}=d\bar C.
}
\] A naive spin-field operator with a coefficient \(C(k)\) but no
\(k_\mu\) would represent the R--R potential rather than its field
strength. As the book says, that naive expression is not BRST invariant.

There is a useful direct check. The scalar bispinor polarization has the
form \[
\mathcal F_{\dot\alpha\beta}
\propto
k_\mu\bar\sigma^\mu_{\dot\alpha\beta}C(k).
\] Acting once more with the momentum slash gives \[
k_\nu\sigma^\nu
\mathcal F
\propto
k_\nu\sigma^\nu k_\mu\bar\sigma^\mu C(k)
\propto
k^2C(k).
\] It vanishes for the massless state, \(k^2=0\), which is precisely the
form of the R--R BRST constraint.

\subsubsection{Why BRST invariance gives a field-strength
bispinor}\label{why-brst-invariance-gives-a-field-strength-bispinor}

The general ten-dimensional R--R vertex in the symmetric
\((-\frac12,-\frac12)\) picture has the schematic form \[
V_{\mathrm{R\text{--}R}}
=
c\bar c\,
e^{-\phi/2}S^A\,
e^{-\bar\phi/2}\bar S^B\,
\mathcal F_{AB}(k)\,
e^{ik\cdot X}.
\] Here \(\mathcal F_{AB}\) is a bispinor. To see the origin of its BRST
constraint, recall that the BRST charge contains the
worldsheet-supercurrent term \[
Q_{\mathrm{BRST}}
\supset
\oint\frac{dz}{2\pi i}\,
\gamma(z)T_F(z),
\qquad
T_F(z)\supset\psi_\mu\partial X^\mu(z).
\] The two OPEs responsible for the momentum constraint are \[
\partial X^\mu(z)e^{ik\cdot X(w)}
\sim
-\frac{i\alpha'}2
\frac{k^\mu}{z-w}e^{ik\cdot X(w)},
\] and \[
\psi^\mu(z)S^A(w)
\sim
\frac{1}{(z-w)^{1/2}}
(\Gamma^\mu)^A{}_{C}S^C(w).
\] Together with the superghost OPE, these produce an unwanted BRST
singularity proportional to \[
k_\mu(\Gamma^\mu)^A{}_{C}\mathcal F_{AB}.
\] Its vanishing gives the first Dirac equation. Repeating the same
calculation in the other chiral sector gives the second: \[
\boxed{
k_\mu\Gamma^\mu\mathcal F=0,
\qquad
\mathcal F\,k_\mu\Gamma^\mu=0.
}
\] Expand the bispinor in antisymmetrized gamma matrices: \[
\mathcal F
=
\sum_n\frac1{n!}
F^{(n)}_{\mu_1\cdots\mu_n}
\Gamma^{\mu_1\cdots\mu_n}.
\] For type IIB the sum contains the odd-degree R--R field strengths
\(F_1,F_3,F_5,\ldots\), with the appropriate duality relations.
Gamma-matrix identities turn the two bispinor Dirac equations into \[
k\wedge F^{(n)}=0,
\qquad
k\mathbin{\lrcorner}F^{(n)}=0.
\] In position space these are \[
dF^{(n)}=0,
\qquad
d{*F^{(n)}}=0,
\] namely the linearized Bianchi identity and equation of motion.
Locally, the first equation implies \[
F^{(n)}=dC^{(n-1)}.
\] Equivalently, in momentum space, \[
\boxed{
F^{(n)}(k)
=
i\,k\wedge C^{(n-1)}(k).
}
\] The explicit \(k\) in (15.84) is therefore the compactified remnant
of the universal relation \(F=dC\).

\subsubsection{From a momentum factor to a derivative
coupling}\label{from-a-momentum-factor-to-a-derivative-coupling}

Consider a local term in the low-energy effective action that is linear
in an R--R scalar: \[
S_{\mathrm{int}}
=
\int d^4x\,
\partial_\mu C(x)\,\mathcal O^\mu(x).
\] After Fourier transformation, its vertex contains \[
ik_\mu C(k)\,\mathcal O^\mu.
\] Conversely, if every local amplitude with one external \(C\) state
contains an explicit \(k_\mu C(k)\), then the corresponding local
interaction contains a derivative on \(C\). A constant mode, \[
C(x)=C_0,
\] has \(dC=0\) and drops out of such perturbative local couplings. This
is the perturbative shift symmetry \[
C\longrightarrow C+C_0.
\]

For the vector, the analogous local interactions are built from \[
F_{\mu\nu}=2\partial_{[\mu}A_{\nu]},
\] for example the kinetic term \[
S_A
=
-\frac1{4g_A^2}
\int d^4x\,F_{\mu\nu}F^{\mu\nu}.
\] An external \(A_\mu\) leg extracted from an interaction expressed
through \(F_{\mu\nu}\) necessarily brings down a factor \[
i(k_\mu\epsilon_\nu-k_\nu\epsilon_\mu).
\] This is exactly the structure exhibited by \(V^A\).

There is a small S-matrix qualification. The argument applies directly
to \textbf{local contact terms}, or equivalently to the local
one-particle-irreducible low-energy action. A complete amplitude can
contain a massless exchange pole proportional to \(1/k\) that
compensates a soft momentum. Such a nonlocal pole describes propagation
through intermediate states, not a local nonderivative R--R coupling.

\subsubsection{Why this excludes perturbative matter charged under the
graviphoton}\label{why-this-excludes-perturbative-matter-charged-under-the-graviphoton}

If a field \(\Psi\) had charge \(Q\) under the graviphoton, its kinetic
term would contain the covariant derivative \[
D_\mu\Psi
=
\left(\partial_\mu-iQ A_\mu\right)\Psi.
\] Expanding \(|D_\mu\Psi|^2\) gives a minimal coupling \[
S_{\mathrm{min}}
\supset
\int d^4x\,A_\mu J^\mu,
\] where, for a charged complex scalar, \[
J^\mu
=
iQ\left(
\Psi^*\partial^\mu\Psi
-\Psi\partial^\mu\Psi^*
\right).
\] The corresponding three-point vertex contains the graviphoton
polarization \(\epsilon_\mu\) itself: \[
\mathcal A_3\propto\epsilon_\mu J^\mu.
\] It is not forced to carry the graviphoton momentum in the
antisymmetric combination \(k_{[\mu}\epsilon_{\nu]}\). Therefore a
minimal coupling requires the potential \(A_\mu\), rather than only
\(F_{\mu\nu}\).

The perturbative closed-string R--R vertex, however, supplies only
\(F_{\mu\nu}\). Accordingly, the perturbative bulk spectrum has no
matter field with a minimal electric coupling to this R--R graviphoton.
This is the precise content of the book's statement that there are no
fields charged with respect to the graviphoton.

This conclusion is stronger than the ordinary Ward identity. Gauge
invariance alone permits \[
\int A_\mu J^\mu
\] when \(\partial_\mu J^\mu=0\). What excludes the coupling here is the
string-theoretic fact that the perturbative R--R state is represented by
a field-strength vertex and that perturbative fundamental closed strings
carry no R--R charge.

Neutral fields may still interact with the graviphoton through
field-strength couplings, for example Pauli-type terms schematically of
the form \[
F_{\mu\nu}\bar\psi\Gamma^{\mu\nu}\psi.
\] Such an interaction does not make \(\psi\) electrically charged: the
distinction is between an \(F_{\mu\nu}\) coupling and a minimal
\(A_\mu J^\mu\) coupling.

\subsubsection{Important scope of the
statement}\label{important-scope-of-the-statement}

The statement is about the perturbative closed-string bulk theory being
discussed in this section. It does not mean that type-II string theory
has no R--R-charged objects.

D-branes carry R--R charge and have Wess--Zumino couplings \[
S_{\mathrm{WZ}}
=
\mu_p\int_{D_p}
\mathbf C_{\mathrm{RR}}\wedge
e^{2\pi\alpha'F_{\mathrm{wv}}+B}.
\] After compactification, a wrapped D-brane can become a particle
charged under a four-dimensional R--R vector, including a graviphoton in
an appropriate compactification. Such states are nonperturbative from
the viewpoint of the closed-string genus expansion. Their
potential-level coupling is therefore not contradicted by the
perturbative sphere-vertex argument. In the RNS description of D-brane
couplings, one must also treat superghost pictures and boundaries
carefully; asymmetric-picture R--R vertices can display the potential
rather than only the symmetric-picture field strength.

Likewise, supergravity may contain gauge-invariant Chern--Simons terms
or modified field strengths, schematically \[
\widetilde F=dC+\text{terms involving other form potentials}.
\] An action in a particular duality frame can then contain an
undifferentiated potential in a topological term while remaining gauge
invariant up to a total derivative. Nonperturbative effects can also
break a continuous R--R axion shift symmetry to a discrete subgroup.
Thus ``only derivative couplings'' should be read as the linearized
perturbative bulk statement encoded by the physical R--R vertex, not as
a prohibition of every potential-level or nonperturbative coupling.

\subsubsection{Checks}\label{checks-24}

Under a graviphoton gauge transformation, \[
\epsilon_\mu\mapsto\epsilon_\mu+a k_\mu,
\] the vertex changes by \[
\delta V^A
\propto
a\,k_\mu k_\nu\sigma^{\mu\nu}
+a\,k_\mu k_\nu\bar\sigma^{\mu\nu}
=0.
\] Thus the explicit vertex depends only on the gauge-equivalence class
of \(A_\mu\).

In the soft limit, \[
k_\mu\longrightarrow0,
\] the local vertices \(V^A,V^C,V^{\bar C}\) vanish at least linearly in
\(k\). A nonzero constant graviphoton potential is locally pure gauge,
and a constant R--R scalar has zero perturbative field strength. Both
conclusions agree with the spacetime interpretation.

Finally, the two terms in \(V^A\) contain \(\sigma^{\mu\nu}\) and
\(\bar\sigma^{\mu\nu}\). They therefore reconstruct the two duality
components \(F^+\) and \(F^-\) of one real graviphoton field strength,
rather than describing two independent vector potentials.

\subsubsection{Result}\label{result-28}

The extra momentum dependence in (15.84) has the direct spacetime
meaning \[
\boxed{
\begin{aligned}
\epsilon_\mu k_\nu\sigma^{\mu\nu}
&\propto
F^+_{\mu\nu}\sigma^{\mu\nu},
\\
\epsilon_\mu k_\nu\bar\sigma^{\mu\nu}
&\propto
F^-_{\mu\nu}\bar\sigma^{\mu\nu},
\\
k_\mu C(k)
&\longleftrightarrow
\frac1i\,\partial_\mu C(x),
\\
k_\mu\bar C(k)
&\longleftrightarrow
\frac1i\,\partial_\mu\bar C(x).
\end{aligned}
}
\] The general R--R spin-field polarization obeys \[
\boxed{
k\!\!\!/\ \mathcal F=0,
\qquad
\mathcal F\,k\!\!\!/=0,
\qquad
\mathcal F=\sum_n\frac1{n!}F^{(n)}\Gamma_{(n)},
\qquad
F^{(n)}=dC^{(n-1)}.
}
\] Therefore the physical R--R vertices describe \[
\boxed{
\text{R--R field strengths, not bare R--R potentials}.
}
\] For the perturbative local bulk action this means \[
\boxed{
C,\bar C\text{ enter through }dC,d\bar C,
\qquad
A_\mu\text{ enters through }F_{\mu\nu},
}
\] so there is no perturbative minimal coupling \(A_\mu J^\mu\) and
hence no perturbative matter charged under the R--R graviphoton.
R--R-charged D-branes and gauge-invariant topological terms lie outside
this restricted statement.

\section{Q9: BRST invariance, current superfields, and picture
changing}\label{q9-brst-invariance-current-superfields-and-picture-changing}

In page 556: \textgreater In the latter case, BRST invariance demands
that \(J^a(z)\) is the upper component of a two-dimensional superfield,
i.e.~it can be obtained via picture changing from a dimension
\(\frac{1}{2}\) (fermionic) field \(\hat{J}^a(z)\).

Then they derive the vertex operator (15.91) and (15.92).

I have no idea how to derive this from the BRST invariance. Show the
explicit computation.

\begin{quote}
\textbf{Textbook location and related material.} Chapter 15, Section
15.4 around p.~556 and Eqs. (15.91)--(15.92): heterotic gauge-current
vertices, superfield components, BRST invariance, and picture changing.
See also Sections 12.1 and 13.2.
\end{quote}

\subsection{Answer}\label{answer-29}

\subsubsection{Core point}\label{core-point-29}

The statement follows from the \(\gamma G\) term in the RNS BRST charge.
For an integrated picture-zero operator \(W(z)\) of weight one, BRST
invariance does not require \(G(z)W(w)\) to vanish. It requires the BRST
variation to be a total derivative. If \[
G(z)W(w)
\sim
\frac{A(w)}{(z-w)^2}
+
\frac{B(w)}{z-w}
+\text{finite},
\] then the \(\gamma G\) term in \(Q_{\mathrm B}\) produces \[
(\partial\gamma)A+\gamma B.
\] This is a total derivative precisely when \[
B=\partial A,
\] because then \[
(\partial\gamma)A+\gamma\partial A
=
\partial(\gamma A).
\]

The \(N=1\) superconformal algebra supplies exactly such a pair: \[
A=\hat J^a,
\qquad
W=J^a=G_{-1/2}\hat J^a,
\] where \(\hat J^a\) is a weight-\(1/2\) superprimary. Indeed, \[
G(z)J^a(w)
\sim
\frac{\hat J^a(w)}{(z-w)^2}
+
\frac{\partial\hat J^a(w)}{z-w}.
\] Therefore the \(\gamma G\) contribution to the BRST variation is \[
\partial\!\left(\gamma\hat J^a\right),
\] which integrates to zero. This is the explicit reason that the
weight-one current must be the upper component of a weight-\(1/2\)
superfield.

There is also an important finite-momentum subtlety. Acting with the
\textbf{full} picture-changing operator on (15.91) produces not only
\(J^a\), but also a term proportional to \[
(k\cdot\psi)\hat J^a.
\] Thus the complete finite-momentum picture-zero matter vertex is \[
\epsilon_\mu\bar\partial X^\mu
\left[
J^a+c_X(k\cdot\psi)\hat J^a
\right]e^{ik\cdot X},
\] where \(c_X\) is convention dependent. Equation (15.92) displays only
the first term. It is exact at \(k=0\), and it can be read as shorthand
for the internal-current part, but at nonzero momentum the omitted term
is required by a literal BRST/picture-changing computation.

\subsubsection{The relevant part of the heterotic BRST
charge}\label{the-relevant-part-of-the-heterotic-brst-charge}

The right-moving, holomorphic half of the heterotic string is an RNS
superstring. Bosonizing \[
\beta=e^{-\phi}\partial\xi,
\qquad
\gamma=\eta e^\phi,
\] the BRST charge contains \[
Q_{\mathrm B}
=
\oint\frac{dz}{2\pi i}
\left[
c\left(T_{\mathrm m}+T_{\mathrm{gh}}\right)
+\gamma G_{\mathrm m}
-\frac14b\gamma^2
+\cdots
\right].
\] Only two parts are needed for the present argument:

\begin{itemize}
\tightlist
\item
  \(cT\) imposes the conformal-weight and Virasoro-primary conditions;
\item
  \(\gamma G_{\mathrm m}\) imposes the worldsheet-supersymmetry
  constraint.
\end{itemize}

The matter supercurrent splits into external and internal pieces, \[
G_{\mathrm m}
=
G_{\mathrm{st}}+G_{\mathrm{int}},
\] with \[
G_{\mathrm{st}}
\propto
\psi_\mu\partial X^\mu.
\] The antiholomorphic half of the heterotic string is bosonic. Its BRST
conditions give the usual spacetime mass-shell and transversality
constraints \[
k^2=0,
\qquad
k^\mu\epsilon_\mu=0.
\] The new condition discussed around (15.91) comes from the holomorphic
\(\gamma G_{\mathrm m}\) term.

\subsubsection{\texorpdfstring{The weight-\(1/2\) lower
component}{The weight-1/2 lower component}}\label{the-weight-12-lower-component}

Let \(\hat J^a\) be an internal NS field satisfying \[
h(\hat J^a)=\frac12.
\] For it to be an \(N=1\) superconformal primary, its state must obey
\[
L_n|\hat J^a\rangle=0
\quad(n>0),
\qquad
G_r|\hat J^a\rangle=0
\quad(r>0).
\] Its upper component is defined by \[
\boxed{
|J^a\rangle
=
G_{-1/2}^{\mathrm{int}}|\hat J^a\rangle.
}
\] The commutator \[
[L_0,G_{-1/2}]
=
\frac12G_{-1/2}
\] immediately gives \[
h(J^a)
=
h(\hat J^a)+\frac12
=1.
\] Thus \[
\mathscr J^a(z,\theta)
=
\hat J^a(z)+\theta J^a(z)
\] is a fermionic \(N=1\) current superfield with lower component of
weight \(1/2\) and upper component of weight one.

\subsubsection{Deriving the two supercurrent
OPEs}\label{deriving-the-two-supercurrent-opes}

The NS \(N=1\) superconformal algebra is \[
\{G_r,G_s\}
=
2L_{r+s}
+
\frac{c}{3}
\left(r^2-\frac14\right)\delta_{r+s,0}.
\] The definition \(J^a=G_{-1/2}\hat J^a\) gives the first component OPE
\[
\boxed{
G(z)\hat J^a(w)
\sim
\frac{J^a(w)}{z-w}.
}
\]

To obtain the OPE with the upper component, act with \(G_{1/2}\): \[
\begin{aligned}
G_{1/2}|J^a\rangle
&=
G_{1/2}G_{-1/2}|\hat J^a\rangle
\\
&=
\{G_{1/2},G_{-1/2}\}|\hat J^a\rangle
-
G_{-1/2}G_{1/2}|\hat J^a\rangle
\\
&=
2L_0|\hat J^a\rangle
\\
&=
|\hat J^a\rangle.
\end{aligned}
\] Here \(G_{1/2}|\hat J^a\rangle=0\) and \(2h(\hat J^a)=1\). Similarly,
\[
G_{-1/2}|J^a\rangle
=
G_{-1/2}^2|\hat J^a\rangle
=
L_{-1}|\hat J^a\rangle.
\] Through the state-operator correspondence, these two relations give
\[
\boxed{
G(z)J^a(w)
\sim
\frac{\hat J^a(w)}{(z-w)^2}
+
\frac{\partial\hat J^a(w)}{z-w}.
}
\] This is the component transformation law of an \(N=1\) superfield
whose lower component has weight \(1/2\).

\subsubsection{Explicit BRST variation of a picture-zero
current}\label{explicit-brst-variation-of-a-picture-zero-current}

Consider first the internal operator \(J^a(w)\) at zero spacetime
momentum. The supercurrent part of its BRST variation is \[
\left[
\oint_w\frac{dz}{2\pi i}\,
\gamma(z)G(z),
J^a(w)
\right].
\] Expand \[
\gamma(z)
=
\gamma(w)+(z-w)\partial\gamma(w)+\cdots
\] and insert the OPE just derived: \[
\gamma(z)G(z)J^a(w)
\sim
\gamma(z)
\left[
\frac{\hat J^a(w)}{(z-w)^2}
+
\frac{\partial\hat J^a(w)}{z-w}
\right].
\] The contour residue is \[
\begin{aligned}
\delta_{\gamma G}J^a
&=
(\partial\gamma)\hat J^a
+
\gamma\partial\hat J^a
\\
&=
\boxed{
\partial\!\left(\gamma\hat J^a\right).
}
\end{aligned}
\] The \(cT\) term likewise gives a total derivative for a weight-one
primary: \[
\delta_{cT}J^a
=
c\,\partial J^a+(\partial c)J^a
=
\partial(cJ^a).
\] Consequently, \[
[Q_{\mathrm B},J^a]
=
\partial\!\left(cJ^a+\gamma\hat J^a\right)
\] up to terms that vanish in the small Hilbert space or belong to the
standard ghost completion. Therefore \[
\boxed{
\left[
Q_{\mathrm B},
\int dz\,J^a(z)
\right]
=0.
}
\] This is the promised explicit BRST computation.

\subsubsection{Why an arbitrary weight-one current
fails}\label{why-an-arbitrary-weight-one-current-fails}

Now suppose that \(K^a\) is merely a bosonic weight-one field, without
assuming it belongs to a superfield. Write its most general relevant
supercurrent OPE as \[
G(z)K^a(w)
\sim
\frac{A^a(w)}{(z-w)^2}
+
\frac{B^a(w)}{z-w}
+\cdots.
\] The same contour calculation gives \[
\delta_{\gamma G}K^a
=
(\partial\gamma)A^a+\gamma B^a.
\] This is a total derivative only if \[
B^a=\partial A^a.
\] Dimensional analysis gives \[
h(A^a)=\frac12.
\] Calling \(A^a=\hat J^a\), the required OPE becomes \[
G(z)K^a(w)
\sim
\frac{\hat J^a(w)}{(z-w)^2}
+
\frac{\partial\hat J^a(w)}{z-w}.
\] By the \(N=1\) algebra, this is precisely the OPE of the upper
component \[
K^a=G_{-1/2}\hat J^a,
\] up to possible null or BRST-trivial operators. Thus a generic
Kac--Moody current of weight one is not sufficient: it must extend to a
super-Kac--Moody multiplet.

\subsubsection{Deriving the canonical vertex
(15.91)}\label{deriving-the-canonical-vertex-15.91}

The complete integrated canonical-picture matter vertex is \[
\boxed{
V_{-1}^a
=
\epsilon_\mu\,
\bar\partial X^\mu(\bar z)\,
\hat J^a(z)\,
e^{-\phi(z)}
e^{ik\cdot X(z,\bar z)}.
}
\] Its conformal weights are easily checked. In the right-moving sector,
\[
h(\hat J^a)=\frac12,
\qquad
h(e^{-\phi})=\frac12,
\qquad
h(e^{ik\cdot X})=0
\quad(k^2=0),
\] so \[
h(V_{-1}^a)=1.
\] In the left-moving sector, \[
\bar h(\bar\partial X^\mu)=1,
\qquad
\bar h(e^{ik\cdot X})=0,
\] so \[
\bar h(V_{-1}^a)=1.
\] The corresponding unintegrated representative is \[
U_{-1}^a
=
c\bar c\,V_{-1}^a.
\] The \(cT\) part of BRST closure gives \(k^2=0\) and the left-moving
vector condition \(k\cdot\epsilon=0\). On the right, the factor
\(e^{-\phi}\) supplies half a unit of weight, while the superprimary
\(\hat J^a\) supplies the other half. Because \(G(z)\hat J^a(w)\) has at
most the allowed simple pole, the \(\gamma G\) term produces no
uncancelled BRST residue in picture \(-1\).

This derives equation (15.91).

\subsubsection{Picture changing to picture
zero}\label{picture-changing-to-picture-zero}

The picture-changing operator is \[
\mathcal X(z)
=
\{Q_{\mathrm B},\xi(z)\}
=
e^\phi G_{\mathrm m}
+\text{ghost terms}.
\] For the integrated matter vertex, the displayed term supplies the
matter part of the picture-zero representative; the omitted ghost terms
contribute BRST-exact pieces or the ghost completion of an unintegrated
operator.

The superghost OPE is \[
e^{\phi(z)}e^{-\phi(w)}
\sim
(z-w).
\] The internal supercurrent gives \[
G_{\mathrm{int}}(z)\hat J^a(w)
\sim
\frac{J^a(w)}{z-w}.
\] Their product has a finite coincident limit: \[
e^\phi G_{\mathrm{int}}\,
\hat J^a e^{-\phi}
\longrightarrow
J^a.
\] This is the internal term shown in equation (15.92).

However, the full supercurrent also acts on the plane wave. Using \[
\partial X^\mu(z)e^{ik\cdot X(w,\bar w)}
\sim
-\frac{i\alpha'}2
\frac{k^\mu}{z-w}
e^{ik\cdot X(w,\bar w)},
\] one finds \[
G_{\mathrm{st}}(z)e^{ik\cdot X(w,\bar w)}
\sim
\frac{c_X\,k\cdot\psi(w)}{z-w}
e^{ik\cdot X(w,\bar w)},
\] where \(c_X\) contains the convention-dependent factor of
\(i\alpha'\) and the normalization of \(G_{\mathrm{st}}\). Hence \[
e^\phi G_{\mathrm{st}}\,
\hat J^a e^{-\phi}e^{ik\cdot X}
\longrightarrow
c_X(k\cdot\psi)\hat J^a e^{ik\cdot X}.
\]

The complete matter result of picture changing is therefore \[
\boxed{
V_0^a
=
\epsilon_\mu\bar\partial X^\mu
\left[
J^a+c_X(k\cdot\psi)\hat J^a
\right]
e^{ik\cdot X}.
}
\] Equivalently, if \[
\mathcal U^a
=
\hat J^a e^{ik\cdot X},
\] then the bracket is the full upper component \[
\boxed{
\mathcal W^a
=
G_{-1/2}^{\mathrm{tot}}\mathcal U^a
=
J^a e^{ik\cdot X}
+
c_X(k\cdot\psi)\hat J^a e^{ik\cdot X}.
}
\] The superconformal algebra then guarantees \[
G(z)\mathcal W^a(w)
\sim
\frac{\mathcal U^a(w)}{(z-w)^2}
+
\frac{\partial\mathcal U^a(w)}{z-w},
\] because \(h(\mathcal U^a)=1/2\) on shell. Its BRST variation is
therefore \[
\delta_{\gamma G}\mathcal W^a
=
\partial(\gamma\mathcal U^a).
\]

\subsubsection{How equation (15.92) should be
read}\label{how-equation-15.92-should-be-read}

Equation (15.92) prints \[
V_{\mathrm{book},0}^a
=
\epsilon_\mu\bar\partial X^\mu
J^a e^{ik\cdot X}.
\] At zero momentum this agrees exactly with the full result: \[
k^\mu=0
\quad\Longrightarrow\quad
V_0^a
=
\epsilon_\mu\bar\partial X^\mu J^a.
\] This is the limit in which \(J^a\) itself is the conserved worldsheet
current generating the spacetime gauge symmetry.

At nonzero momentum, a literal application of the full picture-changing
operator produces the additional term \[
c_X\,
\epsilon_\mu\bar\partial X^\mu
(k\cdot\psi)\hat J^a e^{ik\cdot X}.
\] It is not removed merely by \(k\cdot\epsilon=0\), because \(k\)
contracts with \(\psi^\nu\) whereas \(\epsilon_\mu\) contracts with the
independent left-moving oscillator \(\bar\partial X^\mu\).

This is also visible by comparing with equation (15.87). There the book
picture-changes a canonical vertex with lower component \(\Lambda^i\)
and obtains \[
\left[
\Lambda^i k\cdot\psi
+2\hat\Lambda^i
\right]e^{ik\cdot X}.
\] The same calculation with \[
\Lambda^i\longrightarrow\hat J^a,
\qquad
2\hat\Lambda^i\longrightarrow J^a
\] necessarily gives \[
\left[
(k\cdot\psi)\hat J^a
+J^a
\right]e^{ik\cdot X},
\] up to normalization.

Therefore the precise conclusion is:

\begin{enumerate}
\def\labelenumi{\arabic{enumi}.}
\tightlist
\item
  equation (15.91) is the complete canonical \(-1\)-picture matter
  vertex;
\item
  the \(J^a\) term in (15.92) is derived from the internal part of
  picture changing;
\item
  as printed, (15.92) is exact at \(k=0\) or is shorthand for the
  internal-current contribution;
\item
  at finite \(k\), BRST invariance requires the full upper component of
  \(\hat J^a e^{ik\cdot X}\), including \((k\cdot\psi)\hat J^a\).
\end{enumerate}

\subsubsection{Checks}\label{checks-25}

The lower component has the correct weight for picture \(-1\): \[
h(\hat J^a)+h(e^{-\phi})
=
\frac12+\frac12
=1.
\] The two terms in the full picture-zero bracket both have weight one:
\[
h(J^a)=1,
\qquad
h\bigl((k\cdot\psi)\hat J^a\bigr)
=
\frac12+\frac12
=1.
\]

At \(k=0\), the momentum-dependent completion vanishes and the BRST
computation reduces to \[
[Q_{\mathrm B},J^a]
=
\partial(cJ^a+\gamma\hat J^a),
\] so the integrated current is BRST invariant.

If \(J^a\) were not an upper component, its OPE with \(G\) would give \[
G(z)J^a(w)
\sim
\frac{A^a}{(z-w)^2}
+
\frac{B^a}{z-w},
\qquad
B^a\neq\partial A^a,
\] and the uncancelled variation \[
(\partial\gamma)A^a+\gamma B^a
\] would not be a total derivative. The corresponding integrated
operator would therefore fail BRST invariance, as the book states.

\subsubsection{Result}\label{result-29}

For a weight-\(1/2\) \(N=1\) superprimary \(\hat J^a\), \[
\boxed{
\begin{aligned}
J^a&=G_{-1/2}^{\mathrm{int}}\hat J^a,
\\
G(z)\hat J^a(w)
&\sim
\frac{J^a(w)}{z-w},
\\
G(z)J^a(w)
&\sim
\frac{\hat J^a(w)}{(z-w)^2}
+
\frac{\partial\hat J^a(w)}{z-w}.
\end{aligned}
}
\] The last OPE implies \[
\boxed{
\delta_{\gamma G}J^a
=
(\partial\gamma)\hat J^a
+\gamma\partial\hat J^a
=
\partial(\gamma\hat J^a),
}
\] which is the explicit BRST derivation of the superfield condition.

The canonical vertex is \[
\boxed{
V_{-1}^a
=
\epsilon_\mu\bar\partial X^\mu
\hat J^a e^{-\phi}e^{ik\cdot X},
\qquad
k^2=k\cdot\epsilon=0.
}
\] Full picture changing gives \[
\boxed{
V_0^a
=
\epsilon_\mu\bar\partial X^\mu
\left[
J^a+c_X(k\cdot\psi)\hat J^a
\right]e^{ik\cdot X}.
}
\] Thus the book's displayed expression \[
V_{\mathrm{book},0}^a
=
\epsilon_\mu\bar\partial X^\mu J^a e^{ik\cdot X}
\] is the \(k=0\) form, or the internal-current part, of the complete
picture-zero representative. The requirement that \(J^a\) be the upper
component of \(\hat J^a\) is exactly what makes the \(\gamma G\) BRST
variation a total derivative.

\chapter{String Scattering Amplitudes and Low-Energy Effective Field
Theory}\label{string-scattering-amplitudes-and-low-energy-effective-field-theory-1}

\section{Q1: On-shell string amplitudes and off-shell effective
vertices}\label{q1-on-shell-string-amplitudes-and-off-shell-effective-vertices}

When computing string scattering amplitudes by computing the correlation
functions of vertex operators, we can only compute on-shell scattering
amplitudes since BRST-invariance forces the vertex operators to be
on-shell. In the resulting low-energy effective action, what is an
off-shell amplitude? Why these off-shell amplitudes are also important
that we cannot neglect them?

\begin{quote}
\textbf{Textbook location and related material.} Chapter 16, Sections
16.1 and 16.3 (pp.~585--588, 615--622): on-shell string amplitudes,
low-energy expansions, field redefinitions, and construction of an
effective action. BRST on-shell conditions come from Section 13.2.
\end{quote}

\subsection{Answer}\label{answer-30}

\subsubsection{Core point}\label{core-point-30}

The phrase \textbf{off-shell amplitude} is somewhat loose. An amplitude
in the strict physical sense is an \(S\)-matrix element, and its
external states are on shell. What the low-energy field theory possesses
away from the mass shell are instead \textbf{off-shell Green functions}
and \textbf{off-shell vertices}. They are obtained from the action
without imposing the external equations of motion. Only after one
computes connected Green functions, amputates their external
propagators, takes the external momenta to their poles, and contracts
with physical polarizations does one obtain the on-shell \(S\)-matrix.

Thus there is no contradiction:

\[
\boxed{
\text{first-quantized string theory computes the on-shell }S\text{-matrix}
}
\]

whereas

\[
\boxed{
S_{\mathrm{eff}}[\Phi]\text{ is an off-shell functional whose on-shell }S\text{-matrix agrees with string theory.}
}
\]

The second box does \textbf{not} mean that the string \(S\)-matrix
uniquely determines every term in \(S_{\mathrm{eff}}\). It determines
the effective action only up to transformations that leave the on-shell
\(S\)-matrix invariant, most importantly local field redefinitions,
gauge-fixing choices, total derivatives, and operators proportional to
the lower-order equations of motion.

This is precisely the ambiguity mentioned by the book in §16.1. The book
explicitly states that BRST invariance restricts the scattered string
states to be on shell and that this produces ambiguities in
reconstructing the low-energy action. Section 16.3 then gives the
practical matching prescription: construct
\(\mathcal L_{2\mathrm{pt}}\), match the three-point amplitudes with
\(\mathcal L_{3\mathrm{pt}}\), and at four points subtract the massless
exchange graphs already generated by \(\mathcal L_{3\mathrm{pt}}\); the
remaining analytic terms determine local interactions in
\(\mathcal L_{4\mathrm{pt}}\). The action so obtained is one convenient
off-shell representative of an entire on-shell equivalence class.

Off-shell information is nevertheless indispensable for using a field
theory. It supplies the equations of motion, propagators, interaction
vertices attached to internal lines, response to sources and
backgrounds, and the starting point for loop calculations. The precise
statement is therefore not that every off-shell vertex is an independent
physical datum that must be retained. Rather:

\[
\boxed{
\begin{gathered}
\text{a conventional local effective-action description requires off-shell data,}\\
\text{but one must not assign invariant meaning to a particular off-shell continuation.}
\end{gathered}
}
\]

\subsubsection{From an action to off-shell
vertices}\label{from-an-action-to-off-shell-vertices}

Let \(\Phi^I(x)\) be the massless spacetime fields. Around a chosen
background, write the effective action schematically as

\[
S_{\mathrm{eff}}[\Phi]
=S_2[\Phi]+S_3[\Phi]+S_4[\Phi]+\cdots,
\]

where \(S_n\) contains terms of order \(\Phi^n\). Its quadratic part has
the momentum-space form

\[
S_2
=\frac12\int\frac{d^dk}{(2\pi)^d}\,
\Phi^I(-k)K_{IJ}(k)\Phi^J(k).
\]

After gauge fixing, the inverse of \(K_{IJ}\) is the propagator. For a
scalar,

\[
K(k)\sim k^2+m^2,
\qquad
G(k)\sim\frac{i}{k^2+m^2-i\epsilon}.
\]

The physical mass shell is the zero locus of the inverse propagator,

\[
K(k)=0.
\]

For tensor or gauge fields this equation must be supplemented by the
physical-polarization conditions. For example, a massless vector
satisfies

\[
k^2=0,
\qquad
k\cdot\epsilon=0,
\qquad
\epsilon_\mu\sim\epsilon_\mu+k_\mu\lambda.
\]

The \(n\)-point 1PI vertex is obtained by differentiating the effective
action:

\[
\Gamma^{(n)}_{I_1\cdots I_n}(x_1,\ldots,x_n)
=
\left.
\frac{\delta^n\Gamma}
{\delta\Phi^{I_1}(x_1)\cdots\delta\Phi^{I_n}(x_n)}
\right|_{\Phi=\Phi_0}.
\]

After Fourier transformation, translation invariance gives

\[
\Gamma^{(n)}(k_1,\ldots,k_n)
=(2\pi)^d\delta^{(d)}\!\left(\sum_{i=1}^n k_i\right)
\widehat\Gamma^{(n)}(k_1,\ldots,k_n).
\]

Momentum conservation is already imposed here, but the individual
momenta need not satisfy

\[
k_i^2=-m_i^2.
\]

The function \(\widehat\Gamma^{(n)}\) at such general momenta is what is
often called an off-shell amplitude. More accurately, it is an amputated
off-shell 1PI vertex. It depends on the chosen fields and, in gauge
theories, on gauge fixing.

The physical scattering amplitude is obtained only after restricting the
external momenta to the poles of the exact propagators and contracting
the external indices with physical wave functions. Schematically, the
LSZ relation has the form

\[
\mathcal A_n
\sim
\lim_{K_i(k_i)\to0}
\prod_{i=1}^n
\left[Z_i^{-1/2}K_i(k_i)\right]
G^{(n)}_{\mathrm{conn}}(k_1,\ldots,k_n),
\]

where \(G^{(n)}_{\mathrm{conn}}\) is the connected Green function before
amputation. The factors \(K_i\) remove the external propagators, and the
final restriction puts the external particles on shell. For gauge
fields, only the BRST cohomology classes of external polarizations
survive.

This field-theory distinction mirrors the worldsheet statement. A
BRST-invariant string vertex operator satisfies

\[
Q_{\mathrm B}V_i=0,
\]

which includes the target-space mass-shell and polarization constraints.
A BRST-exact change,

\[
V_i\longrightarrow V_i+Q_{\mathrm B}\Lambda_i,
\]

decouples from the physical amplitude. Consequently, the ordinary
worldsheet correlator directly produces the analogue of the final,
LSZ-reduced quantity, not a unique function \(\widehat\Gamma^{(n)}\) at
arbitrary \(k_i^2\).

\subsubsection{Why internal lines are off
shell}\label{why-internal-lines-are-off-shell}

Even when all external particles are on shell, the momenta carried by
internal lines are generally off shell. Consider a four-point process
with external momenta \(k_1,k_2,k_3,k_4\) and an \(s\)-channel internal
momentum

\[
q=k_1+k_2,
\qquad
s=-q^2
\]

in the mostly-plus convention.

The external conditions \(k_i^2=-m_i^2\) do not imply

\[
q^2=-m_{\mathrm{int}}^2.
\]

The field-theory exchange contribution therefore uses cubic vertices
evaluated with one off-shell leg:

\[
\mathcal A_4^{(s)\text{-exchange}}
=
V_3(k_1,k_2,-q)
\frac{i}{K(q)}
V_3(q,k_3,k_4).
\]

There are analogous \(t\)- and \(u\)-channel terms. The full tree-level
field-theory amplitude is

\[
\mathcal A_4^{\mathrm{EFT}}
=
\mathcal A_4^{(s)\text{-exchange}}
+\mathcal A_4^{(t)\text{-exchange}}
+\mathcal A_4^{(u)\text{-exchange}}
+\mathcal A_4^{\mathrm{contact}}.
\]

Only this complete sum, with its external legs put on shell, is required
to agree with the string amplitude.

This is the content of the matching procedure described in §16.3. The
poles of the string four-point amplitude due to massless states must be
reproduced by exchange diagrams constructed from the lower-point
effective interactions. Once those poles are subtracted, the part
analytic in the external momenta is expanded as

\[
\mathcal A_{4,\mathrm{analytic}}
=c_0+c_2\alpha' k^2+c_4(\alpha' k^2)^2+\cdots,
\]

with the kinematic tensors understood. These terms are represented by
local contact operators in \(\mathcal L_{4\mathrm{pt}}\) of the
schematic form

\[
D^nF^m,
\qquad m\leq4,
\]

with the appropriate powers of \(\alpha'\) fixed by dimensions. As the
book notes below equation (16.130), terms with \(m>4\) require
amplitudes with more than four external gauge fields.

Equivalently, the full string amplitude contains massive-string poles
such as

\[
\frac{1}{q^2+M_N^2},
\qquad
M_N^2\sim\frac{N}{\alpha'}.
\]

At \(|q^2|\ll M_N^2\), a massive propagator admits the expansion

\[
\frac{1}{q^2+M_N^2}
=
\frac{1}{M_N^2}
\left(
1-\frac{q^2}{M_N^2}
+\frac{q^4}{M_N^4}
-\cdots
\right).
\]

Each term becomes a local higher-derivative interaction among the light
fields. This is how an on-shell string amplitude determines the
derivative expansion of an off-shell low-energy action.

\subsubsection{Field redefinitions and the non-uniqueness of off-shell
vertices}\label{field-redefinitions-and-the-non-uniqueness-of-off-shell-vertices}

The central subtlety can be seen with a scalar field. Let

\[
S_2[\phi]
=\frac12\int\frac{d^dk}{(2\pi)^d}\,
\phi(-k)K(k)\phi(k),
\qquad
K(k)=k^2+m^2,
\]

and add the local cubic operator

\[
\Delta S
=\frac{a}{\Lambda^p}\int d^dx\,
\phi^2 K\phi.
\]

Here \(K\phi=0\) is the free equation of motion. Upon symmetrizing the
three identical external legs, this operator changes the cubic off-shell
vertex by a term proportional to

\[
\Delta V_3(k_1,k_2,k_3)
\propto
\frac{a}{\Lambda^p}
\left[K(k_1)+K(k_2)+K(k_3)\right].
\]

If all three external legs are on shell, then \(K(k_i)=0\), and hence

\[
\left.\Delta V_3\right|_{\mathrm{on\ shell}}=0.
\]

The string three-point \(S\)-matrix therefore cannot determine the
coefficient \(a\). Nevertheless, if the third leg is an internal line
carrying momentum \(q\), then

\[
\Delta V_3(k_1,k_2,q)
\propto K(q)
\]

for on-shell \(k_1\) and \(k_2\). In an exchange graph this factor can
cancel the propagator:

\[
\Delta V_3(k_1,k_2,-q)
\frac{i}{K(q)}
V_3(q,k_3,k_4)
\propto
iV_3(q,k_3,k_4).
\]

The resulting contribution is local and has the form of a contact
interaction. Thus changing an off-shell cubic vertex changes the
division of a four-point amplitude into

\[
\text{exchange part}+\text{contact part}.
\]

But it need not change their on-shell sum. Indeed, perform the local
field redefinition

\[
\phi
\longrightarrow
\phi-\frac{a}{\Lambda^p}\phi^2.
\]

To first order in \(a\), its change of the quadratic action is

\[
\delta S_2
=
\int d^dx\,
\frac{\delta S_2}{\delta\phi}\,\delta\phi
=
-\frac{a}{\Lambda^p}
\int d^dx\,
(K\phi)\phi^2,
\]

which removes \(\Delta S\). At the same time, this redefinition changes
the quartic and higher interactions already present in the action. Those
induced changes precisely compensate the altered exchange graphs. The
equivalence theorem therefore gives

\[
S_{\mathrm{eff}}[\phi]
\quad\text{and}\quad
S'_{\mathrm{eff}}[\phi']
\qquad\Longrightarrow\qquad
\mathcal S[S_{\mathrm{eff}}]
=
\mathcal S[S'_{\mathrm{eff}}]
\]

for local, perturbatively invertible field redefinitions, with the usual
consistent treatment of wave-function normalization and quantum
Jacobians.

More generally, if

\[
E_I(\Phi)=\frac{\delta S_{\mathrm{eff}}}{\delta\Phi^I},
\]

then an infinitesimal field redefinition

\[
\Phi^I\longrightarrow\Phi^I-F^I(\Phi,\partial\Phi,\ldots)
\]

changes the action by

\[
\Delta S_{\mathrm{eff}}
=
-\int d^dx\,F^I E_I
+O(F^2).
\]

Therefore operators proportional to the equations of motion are
redundant in the on-shell operator basis. They can change off-shell
vertices while leaving the physical \(S\)-matrix unchanged.

This also explains the book's discussion around equation (16.116). A
Weyl rescaling of the metric together with field redefinitions converts
the Einstein-frame action into the string-frame action. The two actions
have different off-shell couplings---for example, the dilaton is
distributed differently among the kinetic terms---but describe the same
perturbative on-shell physics when fields and external states are
translated consistently.

\subsubsection{What on-shell string amplitudes do and do not
determine}\label{what-on-shell-string-amplitudes-do-and-do-not-determine}

Suppose two local low-energy actions satisfy

\[
S'_{\mathrm{eff}}
=S_{\mathrm{eff}}
+\int d^dx\,F^I E_I
+\int d^dx\,\partial_\mu J^\mu
+O(F^2).
\]

They may have different off-shell vertices, but they give the same
on-shell scattering amplitudes. Consequently, matching string amplitudes
determines

\[
\boxed{
S_{\mathrm{eff}}
\quad\text{modulo field redefinitions, total derivatives, and gauge-fixing choices.}
}
\]

At a fixed order in \(\alpha'\), this is usually implemented by choosing
a basis of local gauge- and diffeomorphism-invariant operators and
eliminating redundant operators with integrations by parts, algebraic
identities, and the lower-order equations of motion. On-shell string
amplitudes then determine the coefficients of the remaining independent
operators.

For example, a curvature-squared correction may be written in several
bases involving

\[
R_{\mu\nu\rho\sigma}R^{\mu\nu\rho\sigma},
\qquad
R_{\mu\nu}R^{\mu\nu},
\qquad
R^2.
\]

A local metric redefinition of the schematic form

\[
g_{\mu\nu}
\longrightarrow
g_{\mu\nu}
+\alpha'\left(aR_{\mu\nu}+b g_{\mu\nu}R\right)
\]

shifts the coefficients of terms involving the leading Einstein
equation. Hence the on-shell graviton \(S\)-matrix cannot determine all
three coefficients separately in an arbitrary off-shell basis. It
determines the field-redefinition-invariant combinations. This is why
the book can say that higher-derivative terms are inferred from
amplitudes while also warning that the resulting action is ambiguous.

The same qualification applies to gauge dependence. An off-shell
gauge-field Green function contains longitudinal components and
generally depends on the gauge-fixing parameter. A physical external
state is transverse and belongs to BRST cohomology, so these
gauge-dependent components cancel or decouple in the on-shell
\(S\)-matrix.

\subsubsection{Why an off-shell action is still
necessary}\label{why-an-off-shell-action-is-still-necessary}

Although its detailed form is not unique, an off-shell action is needed
for several distinct reasons.

\paragraph{It generates internal propagation and complete higher-point
amplitudes}\label{it-generates-internal-propagation-and-complete-higher-point-amplitudes}

An \(n\)-point \(S\)-matrix element is not generally one \(n\)-field
contact vertex. It is a sum of contact and exchange graphs. Internal
momenta are off shell, so one needs the quadratic operator, its inverse
propagator, and vertices defined away from the external mass shell. The
book's four-point matching prescription is an explicit use of precisely
this off-shell field-theory machinery.

One must retain the \textbf{complete action in a chosen field basis}. It
would be inconsistent to delete a term merely because its contribution
vanishes when every leg attached to that vertex is put on shell, while
leaving all other vertices unchanged. Such a term may contribute when
one of those legs is internal. It may be removed only through a
consistent field redefinition, including all induced changes in the
remaining interactions.

\paragraph{The equations of motion come from variations away from a
solution}\label{the-equations-of-motion-come-from-variations-away-from-a-solution}

A classical background \(\Phi_0\) is determined by

\[
\left.
\frac{\delta S_{\mathrm{eff}}}{\delta\Phi^I(x)}
\right|_{\Phi_0}=0.
\]

This derivative compares the action at configurations infinitesimally
displaced from \(\Phi_0\). If the action existed only on solutions, its
first variation---and therefore its equations of motion---could not be
defined.

The same observation applies to tadpoles. A one-point function is the
first derivative of the effective action. A nonzero tadpole means that
the chosen background is not a stationary point. This is not an ordinary
scattering amplitude between asymptotic particles, but it is essential
for deciding whether perturbation theory has been set up around a valid
vacuum.

\paragraph{Potentials and vacuum structure are intrinsically off
shell}\label{potentials-and-vacuum-structure-are-intrinsically-off-shell}

For scalar fields that are constant in spacetime,

\[
\Gamma[\phi_{\mathrm{const}}]
=-\operatorname{Vol}_d\,V_{\mathrm{eff}}(\phi).
\]

A generic value of \(\phi\) is not a solution because

\[
\frac{dV_{\mathrm{eff}}}{d\phi}\neq0.
\]

Nevertheless, the entire function \(V_{\mathrm{eff}}\) is needed to
locate vacua, test stability, calculate masses from its Hessian, and
study rolling or tunnelling solutions. On-shell scattering about one
stationary point determines only local information accessible around
that vacuum and may not determine a global potential without additional
input.

\paragraph{Quantum calculations use off-shell loop
momenta}\label{quantum-calculations-use-off-shell-loop-momenta}

In a loop integral, an internal momentum is integrated over all values:

\[
\int\frac{d^dq}{(2\pi)^d}
\frac{N(q,k_i)}{K_1(q)K_2(q+k_1)\cdots}.
\]

The internal lines are not restricted to the zeros of their inverse
propagators. Off-shell Feynman rules are therefore necessary for
ordinary perturbative calculations in the effective theory. The final
on-shell observable remains invariant under consistent local field
redefinitions and gauge choices, even though intermediate integrands and
off-shell Green functions change.

\paragraph{Symmetries are most usefully encoded off
shell}\label{symmetries-are-most-usefully-encoded-off-shell}

Gauge invariance and diffeomorphism invariance relate vertices of
different orders. As §16.3 emphasizes, a covariant kinetic term already
contains interactions with gauge bosons and arbitrarily many gravitons
when expanded in fluctuations. Writing a gauge- and
diffeomorphism-invariant action makes these relations manifest and
guarantees the Ward or Slavnov--Taylor identities needed for
cancellations among diagrams.

Supersymmetry may close only modulo equations of motion if no
auxiliary-field formulation is available. This is another sense in which
an on-shell symmetry algebra can be sufficient for the physical spectrum
while an off-shell action remains the practical object used to organize
interactions.

\subsubsection{How string theory can discuss genuinely off-shell
configurations}\label{how-string-theory-can-discuss-genuinely-off-shell-configurations}

The limitation belongs to the ordinary first-quantized vertex-operator
prescription, not to every possible formulation of string theory.

In covariant string field theory, the string field \(\Psi\) is not
restricted to BRST cohomology. For example, the schematic cubic
open-string-field-theory action begins as

\[
S_{\mathrm{SFT}}[\Psi]
=\frac12\langle\Psi,Q_{\mathrm B}\Psi\rangle
+\frac{g_s}{3}\langle\Psi,\Psi*\Psi\rangle
+\cdots.
\]

The equation \(Q_{\mathrm B}\Psi=0\) emerges from the linearized
equation of motion instead of being imposed before the field is
inserted. String field theory therefore supplies off-shell vertices and
propagators, although these depend on choices such as gauge fixing,
local coordinates on punctured surfaces, and the definition of the
string field. Its on-shell \(S\)-matrix is independent of those choices.

There is also the worldsheet sigma-model route discussed earlier in the
book. One allows general spacetime-dependent background couplings and
computes their beta functions. The conditions

\[
\beta^I=0
\]

are spacetime equations of motion in a particular field or
renormalization scheme. Away from \(\beta^I=0\), the worldsheet theory
is not an exact conformal background, so this is not an ordinary
correlator of physical BRST-invariant scattering vertices. Scheme
changes of the beta functions correspond closely to spacetime field
redefinitions, again reflecting the non-uniqueness of off-shell data.

\subsubsection{Checks}\label{checks-26}

\begin{enumerate}
\def\labelenumi{\arabic{enumi}.}
\item
  \textbf{A free external leg.} For a scalar, LSZ multiplies the
  connected Green function by \(K(k)=k^2+m^2\) and then takes
  \(K(k)\to0\). Terms in an amputated vertex proportional to \(K(k)\)
  vanish for that external leg. This is the elementary origin of the
  equation-of-motion ambiguity.
\item
  \textbf{An internal leg.} For \(q=k_1+k_2\), there is no reason for
  \(K(q)\) to vanish. A factor \(K(q)\) in a vertex can cancel
  \(1/K(q)\) from the propagator, converting what looks like exchange
  into a local contact contribution. This confirms that contact and
  exchange pieces separately depend on the off-shell convention.
\item
  \textbf{A field redefinition.} The identity
\end{enumerate}

\[
S[\Phi+\delta\Phi]
=S[\Phi]
+\int d^dx\,
\frac{\delta S}{\delta\Phi^I}\delta\Phi^I
+O(\delta\Phi^2)
\]

shows directly that equation-of-motion operators are generated or
removed by changing field coordinates. Their coefficients cannot be
on-shell observables.

\begin{enumerate}
\def\labelenumi{\arabic{enumi}.}
\setcounter{enumi}{3}
\tightlist
\item
  \textbf{Einstein frame versus string frame.} The two frames have
  visibly different dilaton and metric couplings off shell, yet they are
  related by invertible field redefinitions. Physical amplitudes agree
  after the external states and normalizations are translated. This is a
  concrete example already present in §16.3.
\end{enumerate}

\subsubsection{Result}\label{result-30}

An off-shell \(n\)-point vertex is

\[
\Gamma^{(n)}_{I_1\cdots I_n}
=
\left.
\frac{\delta^n\Gamma}
{\delta\Phi^{I_1}\cdots\delta\Phi^{I_n}}
\right|_{\Phi_0},
\qquad
\sum_i k_i=0,
\]

with no requirement that

\[
K_i(k_i)=0.
\]

The physical amplitude is the pole residue with external legs on shell:

\[
\mathcal A_n
\sim
\lim_{K_i(k_i)\to0}
\prod_i Z_i^{-1/2}K_i(k_i)\,
G^{(n)}_{\mathrm{conn}}
\big|_{\mathrm{physical\ polarizations}}.
\]

At four points,

\[
\mathcal A_4^{\mathrm{EFT}}
=
\sum_{s,t,u}
V_3\frac{i}{K}V_3
+V_4,
\]

so internal off-shell vertices and propagators are required even though
the external states are on shell.

Under a local field redefinition,

\[
\Phi^I\to\Phi^I-F^I,
\qquad
S_{\mathrm{eff}}\to
S_{\mathrm{eff}}-
\int d^dx\,F^I\frac{\delta S_{\mathrm{eff}}}{\delta\Phi^I}
+O(F^2),
\]

the off-shell vertices change but the on-shell \(S\)-matrix does not.
Therefore the string amplitudes determine

\[
\boxed{
S_{\mathrm{eff}}
\text{ only up to field redefinitions and other on-shell-trivial terms.}
}
\]

The correct conclusion is:

\[
\boxed{
\begin{aligned}
&\text{Off-shell Green functions are necessary ingredients of the conventional action formalism,}\\
&\text{but they are not separately physical or uniquely fixed by first-quantized string amplitudes.}
\end{aligned}
}
\]

\section{Q2: Explicit three-tachyon sphere
amplitude}\label{q2-explicit-three-tachyon-sphere-amplitude}

I want to see some explicit computation for the correlation function of
vertex operators. Compute the simplest three-point amplitude of the
closed bosonic string three-tachyon eq(16.24).

\begin{quote}
\textbf{Textbook location and related material.} Chapter 16, Section
16.2 around Eq. (16.24): the closed-bosonic three-tachyon sphere
correlator, conformal Killing group fixing, and ghost insertions.
\end{quote}

\subsection{Answer}\label{answer-31}

\subsubsection{Core point}\label{core-point-31}

Equation (16.24) is the simplest closed-string tree amplitude because
the worldsheet is the sphere and the conformal Killing group
\(SL(2,\mathbb C)\) fixes all three vertex positions. There is therefore
no modulus and no remaining position integral.

The complete gauge-fixed correlator is

\[
\mathcal A_3
=C_{S^2}g_c^3
\left\langle
c\bar c\,V_T(k_1;z_1,\bar z_1)
c\bar c\,V_T(k_2;z_2,\bar z_2)
c\bar c\,V_T(k_3;z_3,\bar z_3)
\right\rangle_{S^2},
\]

where

\[
V_T(k_i;z_i,\bar z_i)
=:e^{ik_i\cdot X(z_i,\bar z_i)}:,
\qquad
k_i^2=\frac{4}{\alpha'}.
\]

The answer factorizes into a ghost correlator, a nonzero-mode matter
correlator, and the integral over the spacetime zero mode of \(X^\mu\):

\[
\mathcal A_3
=C_{S^2}g_c^3
\underbrace{\langle c_1c_2c_3\rangle
\langle\bar c_1\bar c_2\bar c_3\rangle}_{\text{ghosts}}
\underbrace{\left\langle\prod_{i=1}^3:e^{ik_i\cdot X_i}:\right\rangle'}_{\text{matter nonzero modes}}
\underbrace{\int d^dx_0\,e^{ix_0\cdot\sum_i k_i}}_{\text{matter zero mode}}.
\]

The three factors are

\[
\langle c_1c_2c_3\rangle
\langle\bar c_1\bar c_2\bar c_3\rangle
=|z_{12}z_{13}z_{23}|^2,
\]

\[
\left\langle\prod_{i=1}^3:e^{ik_i\cdot X_i}:\right\rangle'
=\prod_{i<j}|z_{ij}|^{\alpha'k_i\cdot k_j},
\]

and

\[
\int d^dx_0\,e^{ix_0\cdot\sum_i k_i}
=(2\pi)^d\delta^{(d)}\!\left(\sum_i k_i\right).
\]

For three on-shell tachyons, momentum conservation implies

\[
\alpha'k_i\cdot k_j=-2,
\qquad i\neq j.
\]

There is a factor-of-two typo in the sentence immediately below equation
(16.24): the book prints \(\frac{\alpha'}2k_i\cdot k_j=-2\). The
derivation below shows that the correct condition is

\[
\boxed{
\frac{\alpha'}2k_i\cdot k_j=-1,
\qquad i\neq j.
}
\]

The corrected value is also required for the matter and ghost position
dependences to cancel. Equation (16.24) itself is correct.

Consequently,

\[
\prod_{i<j}|z_{ij}|^{\alpha'k_i\cdot k_j}
=|z_{12}z_{13}z_{23}|^{-2},
\]

which exactly cancels the ghost factor. Restoring the delta function
that the book suppresses,

\[
\boxed{
\mathcal A_3
=(2\pi)^d\delta^{(d)}(k_1+k_2+k_3)
C_{S^2}g_c^3.
}
\]

After adopting the book's convention of omitting the momentum-conserving
delta function, this is precisely

\[
\boxed{
\mathcal A_3=C_{S^2}g_c^3,
}
\]

which is equation (16.24).

\subsubsection{\texorpdfstring{Why three \(c\bar c\) insertions
appear}{Why three c\textbackslash bar c insertions appear}}\label{why-three-cbar-c-insertions-appear}

Before gauge fixing, the formal closed-string sphere amplitude is

\[
\mathcal A_3
\sim
g_c^3 C_{S^2}
\int
\frac{d^2z_1\,d^2z_2\,d^2z_3}
{\operatorname{Vol}SL(2,\mathbb C)}
\left\langle
V_T(k_1;z_1,\bar z_1)
V_T(k_2;z_2,\bar z_2)
V_T(k_3;z_3,\bar z_3)
\right\rangle.
\]

The Möbius group acts by

\[
z\longmapsto\frac{az+b}{cz+d},
\qquad
ad-bc=1,
\]

and has three complex parameters. On a sphere with three punctures,
these parameters can move any three distinct points to any other three
distinct points. Thus all three \(z_i\) are gauge, rather than moduli.

The Faddeev-Popov determinant for fixing \(z_1,z_2,z_3\) is represented
by three holomorphic \(c\) ghosts and three antiholomorphic \(\bar c\)
ghosts. Equivalently, the sphere has three holomorphic \(c\) zero modes,
corresponding to the conformal Killing vectors

\[
1,
\qquad
z,
\qquad
z^2,
\]

and the correlator vanishes unless all three are saturated. Therefore
every one of the three fixed matter vertices becomes an unintegrated
vertex,

\[
U_T(k_i;z_i,\bar z_i)
=c(z_i)\bar c(\bar z_i)V_T(k_i;z_i,\bar z_i).
\]

The tachyon matter vertex has weights

\[
h_T=\bar h_T=\frac{\alpha'k_i^2}{4}.
\]

BRST invariance requires \(h_T=\bar h_T=1\), so

\[
k_i^2=\frac4{\alpha'}.
\]

Since \(c\) and \(\bar c\) have weights \(-1\) and \(-1\), respectively,
the unintegrated vertex \(c\bar cV_T\) has total weights \((0,0)\). This
is why its correlator can be independent of the three fixed positions.

\subsubsection{Matter correlator from Wick's
theorem}\label{matter-correlator-from-wicks-theorem}

We now calculate the matter part explicitly. The book's free-field
convention is

\[
\langle X^\mu(z)X^\nu(w)\rangle
=-\frac{\alpha'}2\eta^{\mu\nu}\log(z-w),
\]

together with

\[
\langle\bar X^\mu(\bar z)\bar X^\nu(\bar w)\rangle
=-\frac{\alpha'}2\eta^{\mu\nu}\log(\bar z-\bar w).
\]

Combining the two chiralities gives

\[
\langle X^\mu(z,\bar z)X^\nu(w,\bar w)\rangle
=-\frac{\alpha'}2\eta^{\mu\nu}\log|z-w|^2.
\]

For normal-ordered exponentials, self-contractions within a single
vertex are omitted. Wick's theorem retains only contractions between
distinct vertices. If

\[
A_i=ik_i\cdot X(z_i,\bar z_i),
\]

then

\[
\left\langle\prod_{i=1}^3:e^{A_i}:\right\rangle'
=
\exp\left(\sum_{i<j}\langle A_iA_j\rangle\right).
\]

For \(i\neq j\),

\[
\begin{aligned}
\langle A_iA_j\rangle
&=i^2 k_{i\mu}k_{j\nu}
\langle X^\mu(z_i,\bar z_i)X^\nu(z_j,\bar z_j)\rangle\\
&=(-1)k_{i\mu}k_{j\nu}
\left[-\frac{\alpha'}2\eta^{\mu\nu}\log|z_{ij}|^2\right]\\
&=\frac{\alpha'}2 k_i\cdot k_j\log|z_{ij}|^2.
\end{aligned}
\]

Exponentiating this result gives

\[
\begin{aligned}
\left\langle
\prod_{i=1}^3:e^{ik_i\cdot X(z_i,\bar z_i)}:
\right\rangle'
&=\exp\left[
\frac{\alpha'}2
\sum_{i<j}k_i\cdot k_j\log|z_{ij}|^2
\right]\\
&=\prod_{i<j}|z_{ij}|^{\alpha'k_i\cdot k_j}\\
&=|z_{12}|^{\alpha'k_1\cdot k_2}
|z_{13}|^{\alpha'k_1\cdot k_3}
|z_{23}|^{\alpha'k_2\cdot k_3}.
\end{aligned}
\]

This is the three-point specialization of the Koba-Nielsen factor.

\subsubsection{The zero mode and momentum
conservation}\label{the-zero-mode-and-momentum-conservation}

The preceding Wick contraction concerns the oscillating part of
\(X^\mu\). To see the remaining factor, split

\[
X^\mu(z,\bar z)=x_0^\mu+X'^\mu(z,\bar z),
\]

where \(x_0^\mu\) is constant on the sphere. The product of the three
exponentials contains

\[
\prod_{i=1}^3e^{ik_i\cdot x_0}
=e^{ix_0\cdot(k_1+k_2+k_3)}.
\]

The path integral over \(x_0\) is an ordinary Fourier integral:

\[
\int d^dx_0\,
e^{ix_0\cdot(k_1+k_2+k_3)}
=(2\pi)^d\delta^{(d)}(k_1+k_2+k_3).
\]

Thus the full matter correlator is

\[
\left\langle
\prod_{i=1}^3:e^{ik_i\cdot X_i}:
\right\rangle
=(2\pi)^d\delta^{(d)}\!\left(\sum_i k_i\right)
\prod_{i<j}|z_{ij}|^{\alpha'k_i\cdot k_j}.
\]

The paragraph immediately preceding equation (16.24) states that the
book assumes momentum conservation and drops this delta function. It is
therefore absent from the printed final answer, not from the underlying
path integral.

\subsubsection{Deriving the pairwise on-shell
products}\label{deriving-the-pairwise-on-shell-products}

For the closed bosonic tachyon,

\[
k_i^2=-m_T^2=\frac4{\alpha'}.
\]

Momentum conservation is

\[
k_1+k_2+k_3=0.
\]

For example,

\[
k_1+k_2=-k_3.
\]

Squaring both sides yields

\[
k_1^2+k_2^2+2k_1\cdot k_2=k_3^2.
\]

Substituting the three mass-shell conditions,

\[
\frac4{\alpha'}+\frac4{\alpha'}
+2k_1\cdot k_2
=\frac4{\alpha'},
\]

so

\[
k_1\cdot k_2=-\frac2{\alpha'}.
\]

Cyclically,

\[
\boxed{
k_1\cdot k_2
=k_1\cdot k_3
=k_2\cdot k_3
=-\frac2{\alpha'}.
}
\]

This gives the matter factor

\[
\left\langle\prod_{i=1}^3:e^{ik_i\cdot X_i}:\right\rangle'
=|z_{12}z_{13}z_{23}|^{-2}.
\]

Thus the correct condition to use in the step below equation (16.24) is

\[
\frac{\alpha'}2 k_i\cdot k_j=-1,
\qquad i\neq j.
\]

The printed right-hand side \(-2\) would instead give
\(\alpha'k_i\cdot k_j=-4\) and hence the matter factor
\(|z_{12}z_{13}z_{23}|^{-4}\). Multiplication by the ghost factor would
then leave \(|z_{12}z_{13}z_{23}|^{-2}\), contradicting the required
\(SL(2,\mathbb C)\) independence. This provides an independent check
that the printed value is a typo.

\subsubsection{Ghost correlator}\label{ghost-correlator}

The \(bc\) system has the basic OPE

\[
b(z)c(w)\sim\frac1{z-w}.
\]

On the sphere, the normalized correlator saturating the three
holomorphic \(c\) zero modes is

\[
\langle c(z_1)c(z_2)c(z_3)\rangle
=z_{12}z_{13}z_{23}.
\]

This formula can be seen directly from the mode expansion. Only the
three zero modes can contribute:

\[
c(z)=c_1+zc_0+z^2c_{-1}+\text{nonzero-mode terms}.
\]

With the standard sphere normalization

\[
\langle0|c_{-1}c_0c_1|0\rangle=1,
\]

the coefficient of \(c_{-1}c_0c_1\) in the product of the three
\(c(z_i)\) is the Vandermonde determinant

\[
\begin{aligned}
\langle c(z_1)c(z_2)c(z_3)\rangle
&=
\det
\begin{pmatrix}
z_1^2&z_1&1\\
z_2^2&z_2&1\\
z_3^2&z_3&1
\end{pmatrix}\\
&=(z_1-z_2)(z_1-z_3)(z_2-z_3)\\
&=z_{12}z_{13}z_{23}.
\end{aligned}
\]

The antiholomorphic sector similarly gives

\[
\langle\bar c(\bar z_1)\bar c(\bar z_2)\bar c(\bar z_3)\rangle
=\bar z_{12}\bar z_{13}\bar z_{23}.
\]

Therefore

\[
\begin{aligned}
\langle c_1c_2c_3\rangle
\langle\bar c_1\bar c_2\bar c_3\rangle
&=(z_{12}z_{13}z_{23})
(\bar z_{12}\bar z_{13}\bar z_{23})\\
&=|z_{12}z_{13}z_{23}|^2.
\end{aligned}
\]

This is exactly the Faddeev-Popov factor displayed in equations
(16.5)-(16.6).

\subsubsection{Cancellation of the insertion-point
dependence}\label{cancellation-of-the-insertion-point-dependence}

Multiplying the ghost and matter nonzero-mode factors gives

\[
\begin{aligned}
&\langle c_1c_2c_3\rangle
\langle\bar c_1\bar c_2\bar c_3\rangle
\left\langle\prod_{i=1}^3:e^{ik_i\cdot X_i}:\right\rangle'\\
&\qquad=
|z_{12}z_{13}z_{23}|^2
|z_{12}z_{13}z_{23}|^{-2}\\
&\qquad=1.
\end{aligned}
\]

This cancellation is not accidental. Each matter tachyon vertex is a
\((1,1)\) primary, while \(c\bar c\) has weights \((-1,-1)\). Hence
every unintegrated operator \(c\bar cV_T\) has weights \((0,0)\), and
the gauge-fixed three-point correlator cannot depend on which three
distinct points were chosen.

It follows that

\[
\begin{aligned}
\mathcal A_3
&=C_{S^2}g_c^3
(2\pi)^d\delta^{(d)}(k_1+k_2+k_3)\\
&\qquad\times
|z_{12}z_{13}z_{23}|^2
|z_{12}z_{13}z_{23}|^{-2}\\
&=(2\pi)^d\delta^{(d)}(k_1+k_2+k_3)
C_{S^2}g_c^3.
\end{aligned}
\]

Dropping the delta function gives equation (16.24).

\subsubsection{\texorpdfstring{Explicit check with \(z_1=0\), \(z_2=1\),
\(z_3=\infty\)}{Explicit check with z\_1=0, z\_2=1, z\_3=\textbackslash infty}}\label{explicit-check-with-z_10-z_21-z_3infty}

The usual choice is

\[
z_1=0,
\qquad
z_2=1,
\qquad
z_3=\infty.
\]

To treat the point at infinity carefully, first set \(z_3=R\) and take
\(|R|\to\infty\). Then

\[
z_{12}=-1,
\qquad
z_{13}=-R,
\qquad
z_{23}=1-R.
\]

The ghost factor behaves as

\[
|z_{12}z_{13}z_{23}|^2
=|R|^2|1-R|^2
\sim|R|^4.
\]

The matter factor behaves as

\[
|z_{12}z_{13}z_{23}|^{-2}
=|R|^{-2}|1-R|^{-2}
\sim|R|^{-4}.
\]

Their product is one before and after the limit:

\[
\lim_{|R|\to\infty}
\left[|R|^2|1-R|^2\right]
\left[|R|^{-2}|1-R|^{-2}\right]
=1.
\]

Equivalently, since \(U_T=c\bar cV_T\) has weights \((0,0)\), its
insertion at infinity is simply defined by the conformal limit of the
complete operator. Separating the matter and ghost pieces makes each
piece singular, but their product is finite.

\subsubsection{Checks}\label{checks-27}

\begin{enumerate}
\def\labelenumi{\arabic{enumi}.}
\tightlist
\item
  \textbf{Closed-string coupling counting.} Each of the three vertices
  contributes \(g_c\), while the sphere normalization satisfies
  \(C_{S^2}\propto g_c^{-2}\). Therefore
\end{enumerate}

\[
\mathcal A_3\propto g_c^{3-2}=g_c,
\]

in agreement with the general closed-string tree-level rule
\(\mathcal A_N\propto g_s^{N-2}\).

\begin{enumerate}
\def\labelenumi{\arabic{enumi}.}
\setcounter{enumi}{1}
\item
  \textbf{No modulus integral.} The complex dimension of the moduli
  space of a sphere with \(N\) marked points is \(N-3\). For \(N=3\)
  this is zero. The three-point answer is consequently a constant rather
  than a beta or gamma function of Mandelstam invariants.
\item
  \textbf{Worldsheet scaling.} The matter vertex has weight \(+1\) at
  each holomorphic and antiholomorphic insertion; the attached
  \(c\bar c\) factor has weight \(-1\) in each chirality. The total
  unintegrated vertex has weight zero, explaining the exact position
  independence.
\item
  \textbf{Spacetime interpretation.} After removing the universal delta
  function, the amplitude is momentum independent. At this order it
  therefore matches a local cubic tachyon interaction in spacetime. Its
  precise coefficient and phase depend on the normalization conventions
  collected in \(g_c\) and \(C_{S^2}\); the book postpones fixing
  \(C_{S^2}\) until factorization and unitarity are imposed.
\item
  \textbf{Formal nature of the tachyon amplitude.} The bosonic-string
  tachyon signals that the perturbative vacuum is unstable. The
  correlator is nevertheless a perfectly well-defined on-shell CFT
  amplitude, understood with the usual analytic continuation of momenta
  when necessary, and is the cleanest example for learning the
  worldsheet calculation.
\end{enumerate}

\subsubsection{Result}\label{result-31}

The unintegrated tachyon vertex is

\[
U_T(k_i;z_i,\bar z_i)
=c(z_i)\bar c(\bar z_i):e^{ik_i\cdot X(z_i,\bar z_i)}:,
\qquad
k_i^2=\frac4{\alpha'}.
\]

The matter correlator is

\[
\left\langle\prod_{i=1}^3:e^{ik_i\cdot X_i}:\right\rangle
=(2\pi)^d\delta^{(d)}\!\left(\sum_i k_i\right)
\prod_{i<j}|z_{ij}|^{\alpha'k_i\cdot k_j}.
\]

Using

\[
k_1+k_2+k_3=0,
\qquad
k_i^2=\frac4{\alpha'},
\]

one obtains

\[
\alpha'k_i\cdot k_j=-2,
\qquad
\left\langle\prod_i:e^{ik_i\cdot X_i}:\right\rangle'
=|z_{12}z_{13}z_{23}|^{-2}.
\]

The ghost correlator is

\[
\langle c_1c_2c_3\rangle
\langle\bar c_1\bar c_2\bar c_3\rangle
=|z_{12}z_{13}z_{23}|^2.
\]

Therefore

\[
\boxed{
\mathcal A_3
=(2\pi)^d\delta^{(d)}(k_1+k_2+k_3)
C_{S^2}g_c^3,
}
\]

or, suppressing the delta function as in the book,

\[
\boxed{
\mathcal A_3=C_{S^2}g_c^3.
}
\]

\section{Q3: Mixed left-right propagators on the upper
half-plane}\label{q3-mixed-left-right-propagators-on-the-upper-half-plane}

In page603, due to the presence of the boundary along the real axis,
there are non-trivial two-point functions between left- and right-moving
fields. Compute the
\(\langle X^\mu(z)\overline{X}^\nu(\overline{w})\rangle\) and
\(\langle\psi^\mu(z)\overline{\psi}(\overline{w})\rangle\) correlation
functions in eq (16.56) in details.

\begin{quote}
\textbf{Textbook location and related material.} Chapter 16, Section
16.2 around p.~603 and Eqs. (16.56)--(16.58): the upper-half-plane
doubling trick and mixed left--right propagators in the presence of a
D-brane.
\end{quote}

\subsection{Answer}\label{answer-32}

\subsubsection{Core point}\label{core-point-32}

On the full complex plane, the holomorphic and antiholomorphic fields
are independent, so mixed contractions such as

\[
\langle X^\mu(z)\bar X^\nu(\bar w)\rangle_{\mathbb C}
\]

vanish. On the upper half-plane, however, the real boundary glues left
and right movers. For a D\(p\)-brane this gluing is encoded by

\[
D^\mu{}_\nu
=\operatorname{diag}(+1_{\mathrm N},-1_{\mathrm D}),
\]

so a right mover can be replaced by a holomorphic image field at the
reflected point \(\bar w\) in the lower half-plane. The ordinary plane
propagators then immediately give

\[
\boxed{
\left\langle X^\mu(z)\bar X^\nu(\bar w)\right\rangle_{\mathbb H_+}
=-\frac{\alpha'}2D^{\mu\nu}\log(z-\bar w)
}
\]

and

\[
\boxed{
\left\langle\psi^\mu(z)\bar\psi^\nu(\bar w)\right\rangle_{\mathbb H_+}
=\frac{D^{\mu\nu}}{z-\bar w}.
}
\]

These are the first two mixed correlators in equation (16.56). The sign
difference between Neumann and Dirichlet directions is entirely
contained in \(D^{\mu\nu}\).

I interpret the second correlator in the question as
\(\langle\psi^\mu(z)\bar\psi^\nu(\bar w)\rangle\), with the missing
target-space index \(\nu\) restored as in equation (16.56).

The derivation has three steps:

\begin{enumerate}
\def\labelenumi{\arabic{enumi}.}
\tightlist
\item
  derive the left-right gluing conditions from Neumann and Dirichlet
  boundary conditions;
\item
  extend, or double, the fields from \(\mathbb H_+\) to the full plane;
\item
  use the standard holomorphic plane two-point functions between a point
  and its image.
\end{enumerate}

\subsubsection{From Neumann and Dirichlet conditions to the gluing
matrix}\label{from-neumann-and-dirichlet-conditions-to-the-gluing-matrix}

Write the Euclidean worldsheet coordinate as

\[
z=x+iy,
\qquad
\bar z=x-iy,
\qquad
y\geq0.
\]

The boundary is \(y=0\). With

\[
\partial=\frac12(\partial_x-i\partial_y),
\qquad
\bar\partial=\frac12(\partial_x+i\partial_y),
\]

the tangential and normal derivatives are, up to the orientation
convention for the outward normal,

\[
\partial_t=\partial_x=\partial+\bar\partial,
\qquad
\partial_n=\partial_y=i(\partial-\bar\partial).
\]

Split the closed-string coordinate into chiral pieces:

\[
X^\mu(z,\bar z)
=X^\mu(z)+\bar X^\mu(\bar z).
\]

For a world-volume direction \(X^\alpha\), the Neumann condition is

\[
\left.\partial_nX^\alpha\right|_{y=0}=0.
\]

Because the left-moving piece depends only on \(z\) and the right-moving
piece only on \(\bar z\), this gives

\[
\left.\partial X^\alpha(z)\right|_{z=x}
=
\left.\bar\partial\bar X^\alpha(\bar z)\right|_{\bar z=x}.
\]

After integration along the boundary,

\[
X^\alpha(x)=+\bar X^\alpha(x)+\text{constant}.
\]

The additive constant is irrelevant for contractions and can be
suppressed.

For a transverse direction \(X^i\), the D-brane fixes the endpoint
position:

\[
\left.X^i(z,\bar z)\right|_{y=0}=y_0^i.
\]

The boundary value is constant along the real axis, so

\[
\left.\partial_tX^i\right|_{y=0}=0.
\]

This implies

\[
\left.\partial X^i(z)\right|_{z=x}
=-\left.\bar\partial\bar X^i(\bar z)\right|_{\bar z=x},
\]

and hence

\[
X^i(x)=-\bar X^i(x)+y_0^i.
\]

Suppressing the position constant, the two cases combine into

\[
\boxed{
X^\mu(x)=D^\mu{}_\nu\bar X^\nu(x),
}
\]

where

\[
D^\alpha{}_\beta=\delta^\alpha{}_\beta,
\qquad
D^i{}_j=-\delta^i{}_j,
\qquad
D^\alpha{}_i=D^i{}_\alpha=0.
\]

The fermionic boundary condition follows either from the boundary term
in the variation of the NSR fermion action or by requiring that
worldsheet supersymmetry preserve the bosonic boundary conditions. In
the sign convention chosen in equation (16.54), it is

\[
\boxed{
\psi^\mu(x)=D^\mu{}_\nu\bar\psi^\nu(x).
}
\]

Thus

\[
\psi^\alpha(x)=+\bar\psi^\alpha(x)
\qquad(\mathrm N),
\]

whereas

\[
\psi^i(x)=-\bar\psi^i(x)
\qquad(\mathrm D).
\]

There can be a common sign multiplying all fermionic gluing conditions,
associated with the choice of spin structure or boundary-state sign.
Equation (16.54) has chosen this common sign to be \(+1\), so it will
not be displayed below.

The matrix \(D\) is a target-space reflection. In particular,

\[
D^\mu{}_\rho D^\rho{}_\nu=\delta^\mu{}_\nu,
\qquad
D^T\eta D=\eta.
\]

It is therefore legitimate to use \(D^{-1}=D\) when undoing the doubling
transformation.

\subsubsection{The doubling trick}\label{the-doubling-trick}

Define a single holomorphic bosonic field on the full complex plane by

\[
\mathbb X^\mu(\zeta)
=
\begin{cases}
X^\mu(\zeta),&\operatorname{Im}\zeta>0,\\[2mm]
D^\mu{}_\rho\bar X^\rho(\zeta),&\operatorname{Im}\zeta<0.
\end{cases}
\]

Here \(\bar X^\rho(\zeta)\) in the second line means the analytic
right-moving field continued in its own argument into the lower
half-plane. At the boundary \(\zeta=x\), the two definitions agree
because of the gluing condition.

Similarly define the doubled fermion

\[
\Psi^\mu(\zeta)
=
\begin{cases}
\psi^\mu(\zeta),&\operatorname{Im}\zeta>0,\\[2mm]
D^\mu{}_\rho\bar\psi^\rho(\zeta),&\operatorname{Im}\zeta<0.
\end{cases}
\]

The doubled fields are ordinary holomorphic free fields on the full
plane. Their two-point functions are therefore the plane correlators
from equation (16.2):

\[
\left\langle\mathbb X^\mu(\zeta)\mathbb X^\nu(\xi)\right\rangle_{\mathbb C}
=-\frac{\alpha'}2\eta^{\mu\nu}\log(\zeta-\xi),
\]

\[
\left\langle\Psi^\mu(\zeta)\Psi^\nu(\xi)\right\rangle_{\mathbb C}
=\frac{\eta^{\mu\nu}}{\zeta-\xi}.
\]

The entire effect of the boundary has now been moved into the definition
of the fields below the real axis.

\subsubsection{Bosonic mixed correlator}\label{bosonic-mixed-correlator}

Let \(z,w\in\mathbb H_+\). The reflected point \(\bar w\) lies in
\(\mathbb H_-\). From the lower-half-plane definition,

\[
\mathbb X^\lambda(\bar w)
=D^\lambda{}_\rho\bar X^\rho(\bar w).
\]

Multiplying by \(D\) and using \(D^2=1\) gives

\[
\bar X^\nu(\bar w)
=D^\nu{}_\lambda\mathbb X^\lambda(\bar w).
\]

The desired mixed correlator is therefore

\[
\begin{aligned}
\left\langle X^\mu(z)\bar X^\nu(\bar w)\right\rangle_{\mathbb H_+}
&=D^\nu{}_\lambda
\left\langle\mathbb X^\mu(z)\mathbb X^\lambda(\bar w)\right\rangle_{\mathbb C}\\
&=D^\nu{}_\lambda
\left[-\frac{\alpha'}2\eta^{\mu\lambda}\log(z-\bar w)\right]\\
&=-\frac{\alpha'}2
\eta^{\mu\lambda}D^\nu{}_\lambda
\log(z-\bar w).
\end{aligned}
\]

Define

\[
D^{\mu\nu}
=\eta^{\mu\lambda}D^\nu{}_\lambda.
\]

For the diagonal reflection matrix used here, \(D^{\mu\nu}=D^{\nu\mu}\).
Hence

\[
\boxed{
\left\langle X^\mu(z)\bar X^\nu(\bar w)\right\rangle_{\mathbb H_+}
=-\frac{\alpha'}2D^{\mu\nu}\log(z-\bar w).
}
\]

This correlator is holomorphic in \(z\) and antiholomorphic in \(w\), as
it must be. The logarithm has its singularity at \(z=\bar w\), the image
of \(w\) across the boundary.

\subsubsection{Fermionic mixed
correlator}\label{fermionic-mixed-correlator}

The same argument applies to the worldsheet fermion. For the reflected
point,

\[
\bar\psi^\nu(\bar w)
=D^\nu{}_\lambda\Psi^\lambda(\bar w).
\]

Therefore

\[
\begin{aligned}
\left\langle\psi^\mu(z)\bar\psi^\nu(\bar w)\right\rangle_{\mathbb H_+}
&=D^\nu{}_\lambda
\left\langle\Psi^\mu(z)\Psi^\lambda(\bar w)\right\rangle_{\mathbb C}\\
&=D^\nu{}_\lambda
\frac{\eta^{\mu\lambda}}{z-\bar w}\\
&=\frac{D^{\mu\nu}}{z-\bar w}.
\end{aligned}
\]

Thus

\[
\boxed{
\left\langle\psi^\mu(z)\bar\psi^\nu(\bar w)\right\rangle_{\mathbb H_+}
=\frac{D^{\mu\nu}}{z-\bar w}.
}
\]

No additional minus sign arises from the doubling calculation itself.
The relative Dirichlet sign mentioned below equation (16.56) is
precisely the \(-1\) already present in \(D^i{}_j\).

\subsubsection{Neumann and Dirichlet
components}\label{neumann-and-dirichlet-components}

It is useful to write the formulas separately. Along Neumann directions,

\[
D^{\alpha\beta}=\eta^{\alpha\beta},
\]

so

\[
\boxed{
\left\langle X^\alpha(z)\bar X^\beta(\bar w)\right\rangle
=-\frac{\alpha'}2\eta^{\alpha\beta}\log(z-\bar w),
}
\]

\[
\boxed{
\left\langle\psi^\alpha(z)\bar\psi^\beta(\bar w)\right\rangle
=\frac{\eta^{\alpha\beta}}{z-\bar w}.
}
\]

Along spatial Dirichlet directions,

\[
D^{ij}=-\delta^{ij},
\]

and hence

\[
\boxed{
\left\langle X^i(z)\bar X^j(\bar w)\right\rangle
=+\frac{\alpha'}2\delta^{ij}\log(z-\bar w),
}
\]

\[
\boxed{
\left\langle\psi^i(z)\bar\psi^j(\bar w)\right\rangle
=-\frac{\delta^{ij}}{z-\bar w}.
}
\]

Mixed Neumann-Dirichlet components vanish because

\[
D^{\alpha i}=D^{i\alpha}=0.
\]

\subsubsection{Equivalent method-of-images Green
function}\label{equivalent-method-of-images-green-function}

The doubling result can be repackaged as the Green function of the full
nonchiral coordinate

\[
X^\mu(z,\bar z)=X^\mu(z)+\bar X^\mu(\bar z).
\]

Expanding its two-point function gives four chiral contractions:

\[
\begin{aligned}
&\left\langle X^\mu(z,\bar z)X^\nu(w,\bar w)\right\rangle_{\mathbb H_+}\\
&\quad=
\langle X^\mu(z)X^\nu(w)\rangle
+\langle\bar X^\mu(\bar z)\bar X^\nu(\bar w)\rangle\\
&\qquad
+\langle X^\mu(z)\bar X^\nu(\bar w)\rangle
+\langle\bar X^\mu(\bar z)X^\nu(w)\rangle.
\end{aligned}
\]

Using the same-chirality plane propagators and the mixed propagators
derived above,

\[
\boxed{
\begin{aligned}
&\left\langle X^\mu(z,\bar z)X^\nu(w,\bar w)\right\rangle_{\mathbb H_+}\\
&\qquad=
-\frac{\alpha'}2\eta^{\mu\nu}\log|z-w|^2
-\frac{\alpha'}2D^{\mu\nu}\log|z-\bar w|^2.
\end{aligned}
}
\]

The first term is the direct source at \(w\); the second is its image at
\(\bar w\). For a Neumann direction the image has the same sign, whereas
for a Dirichlet direction it has the opposite sign.

This immediately verifies the boundary conditions. For Neumann
directions,

\[
G_{\mathrm N}(z,w)
=-\frac{\alpha'}2
\left[\log|z-w|^2+\log|z-\bar w|^2\right],
\]

and the normal derivative vanishes at \(z=x\in\mathbb R\):

\[
\left.\partial_nG_{\mathrm N}(z,w)\right|_{z=x}=0.
\]

For Dirichlet directions,

\[
G_{\mathrm D}(z,w)
=-\frac{\alpha'}2
\left[\log|z-w|^2-\log|z-\bar w|^2\right].
\]

At \(z=x\in\mathbb R\), the two distances are equal,

\[
|x-w|=|x-\bar w|,
\]

so

\[
\left.G_{\mathrm D}(z,w)\right|_{z=x}=0,
\]

as appropriate for fluctuations of a coordinate fixed on the boundary.

\subsubsection{Boundary-limit checks and equation
(16.58)}\label{boundary-limit-checks-and-equation-16.58}

For a Neumann coordinate, the gluing condition gives

\[
X^\alpha(x,\bar x)
=X^\alpha(x)+\bar X^\alpha(x)
=2X^\alpha(x).
\]

Consequently,

\[
\begin{aligned}
\left\langle X^\alpha(x)X^\beta(x')\right\rangle_{\mathrm{boundary}}
&=4\left[-\frac{\alpha'}2\eta^{\alpha\beta}\log|x-x'|\right]\\
&=-2\alpha'\eta^{\alpha\beta}\log|x-x'|,
\end{aligned}
\]

which is the first line of equation (16.58).

For a Dirichlet coordinate the boundary value itself does not fluctuate,
but its normal derivative does. At the boundary,

\[
\partial_nX^i
=i(\partial-\bar\partial)X^i
=2i\partial X^i,
\]

where the last equality uses \(\partial X^i=-\bar\partial\bar X^i\).
From

\[
\left\langle\partial X^i(x)\partial X^j(x')\right\rangle
=-\frac{\alpha'}2
\frac{\delta^{ij}}{(x-x')^2},
\]

we obtain

\[
\begin{aligned}
\left\langle\partial_nX^i(x)\partial_nX^j(x')\right\rangle
&=(2i)^2
\left[-\frac{\alpha'}2
\frac{\delta^{ij}}{(x-x')^2}\right]\\
&=2\alpha'
\frac{\delta^{ij}}{(x-x')^2},
\end{aligned}
\]

which is the second line of equation (16.58), up to the immaterial
choice of orientation for each normal derivative.

For fermions, moving both points to the boundary and using the chiral
field gives

\[
\left\langle\psi^\mu(x)\psi^\nu(x')\right\rangle
=\frac{\eta^{\mu\nu}}{x-x'},
\]

for both Neumann and Dirichlet directions. The differing left-right sign
has already been absorbed into the gluing relation \(\bar\psi=D\psi\).
This reproduces the third line of (16.58).

\subsubsection{Why there was no mixed correlator on the
sphere}\label{why-there-was-no-mixed-correlator-on-the-sphere}

On the sphere or full plane, the path integral separates into
independent holomorphic and antiholomorphic sectors:

\[
X^\mu(z,\bar z)=X^\mu(z)+\bar X^\mu(\bar z),
\]

with no boundary condition relating the two. One therefore has

\[
\left\langle X^\mu(z)\bar X^\nu(\bar w)\right\rangle_{\mathbb C}=0,
\qquad
\left\langle\psi^\mu(z)\bar\psi^\nu(\bar w)\right\rangle_{\mathbb C}=0.
\]

The disk boundary reduces the independent left and right conformal
symmetries to a diagonal subgroup and imposes the gluing conditions. A
right mover is then not an independent field: it is the reflected image
of a left mover. This is the conceptual origin of equation (16.56).

\subsubsection{Checks}\label{checks-28}

\begin{enumerate}
\def\labelenumi{\arabic{enumi}.}
\item
  \textbf{Image-point singularity.} The denominators depend on
  \(z-\bar w\), not \(z-w\). A mixed contraction becomes singular when
  the left-moving insertion approaches the mirror of the right-moving
  insertion. For two points strictly inside \(\mathbb H_+\),
  \(z=\bar w\) cannot occur; the singularity is reached when the points
  approach the boundary together.
\item
  \textbf{Neumann sign.} For \(D=+1\), the image source has the same
  sign as the original source. This makes the normal derivative of the
  scalar Green function vanish at the real axis.
\item
  \textbf{Dirichlet sign.} For \(D=-1\), the image source has the
  opposite sign. This makes the boundary value of the fluctuating scalar
  Green function vanish.
\item
  \textbf{Fermionic sign.} The standard plane OPE has a positive
  numerator, \(\eta^{\mu\nu}/(z-w)\). Replacing the right mover by its
  image multiplies it by \(D\), giving a plus sign for Neumann
  components and a minus sign for Dirichlet components.
\item
  \textbf{Upper half-plane versus unit disk.} On \(\mathbb H_+\),
  reflection is simply \(w\mapsto\bar w\), so the doubled fermion
  acquires no conformal Jacobian. On the unit disk, the image map is
  \(w\mapsto1/\bar w\); because \(\psi\) has weight \(1/2\), its
  transformation includes the square-root Jacobian mentioned in footnote
  10 on page 603.
\end{enumerate}

\subsubsection{Result}\label{result-32}

The boundary gluing conditions are

\[
X^\mu(x)=D^\mu{}_\nu\bar X^\nu(x),
\qquad
\psi^\mu(x)=D^\mu{}_\nu\bar\psi^\nu(x),
\]

with

\[
D^\mu{}_\nu
=\operatorname{diag}(+1_{\mathrm N},-1_{\mathrm D}),
\qquad
D^2=1.
\]

The doubling trick replaces a right mover at \(\bar w\in\mathbb H_-\) by
a holomorphic image multiplied by \(D\). Using

\[
\langle X^\mu(z)X^\nu(w)\rangle_{\mathbb C}
=-\frac{\alpha'}2\eta^{\mu\nu}\log(z-w),
\]

\[
\langle\psi^\mu(z)\psi^\nu(w)\rangle_{\mathbb C}
=\frac{\eta^{\mu\nu}}{z-w},
\]

one obtains

\[
\boxed{
\left\langle X^\mu(z)\bar X^\nu(\bar w)\right\rangle_{\mathbb H_+}
=-\frac{\alpha'}2D^{\mu\nu}\log(z-\bar w),
}
\]

\[
\boxed{
\left\langle\psi^\mu(z)\bar\psi^\nu(\bar w)\right\rangle_{\mathbb H_+}
=\frac{D^{\mu\nu}}{z-\bar w}.
}
\]

Equivalently, the full bosonic image-charge Green function is

\[
\boxed{
\left\langle X^\mu(z,\bar z)X^\nu(w,\bar w)\right\rangle_{\mathbb H_+}
=-\frac{\alpha'}2\eta^{\mu\nu}\log|z-w|^2
-\frac{\alpha'}2D^{\mu\nu}\log|z-\bar w|^2.
}
\]

\section{Q4: Why vanishing sigma-model beta functions are spacetime
equations of
motion}\label{q4-why-vanishing-sigma-model-beta-functions-are-spacetime-equations-of-motion}

Why the vanishing of the sigma-model \(\beta\) function gives string
equation of motion i.e.~the equation of motion derived by string
effective field theory action?

In my understanding, sigma-model (Polyakov) action shows how string
couples to background fields which are string excitations (and we only
consider the massless part). The vanishing of \(\beta\) function is a
somehow consistency condition required by the definition of the string
theory and this constraint the massless fields. But why these
constraints are exactly equivalent to the equation of motion from the
effective action of these massless fields?

\begin{quote}
\textbf{Textbook location and related material.} Chapter 14, Section
14.1 (pp.~427--432): sigma-model beta functions and background
equations; Chapter 16, Sections 16.3--16.4 (pp.~615--630): the matching
spacetime effective actions and supergravities.
\end{quote}

\subsection{Answer}\label{answer-33}

\subsubsection{Core point}\label{core-point-33}

Your understanding is essentially correct. The missing link is that the
worldsheet sigma model and the spacetime effective action are not two
unrelated theories that happen to impose the same constraints. They are
two descriptions of the dynamics of the same string degrees of freedom:

\begin{itemize}
\tightlist
\item
  in the worldsheet description, \(G_{\mu\nu}(X)\), \(B_{\mu\nu}(X)\),
  and \(\Phi(X)\) are position-dependent couplings multiplying
  worldsheet operators;
\item
  in the spacetime description, those same coupling functions are
  dynamical fields;
\item
  failure of the worldsheet theory to be Weyl invariant is a tadpole, or
  failure of stationarity, for the corresponding spacetime string field.
\end{itemize}

The precise relation is not generally

\[
\beta^I=\frac{\delta S_{\mathrm{eff}}}{\delta\lambda^I}.
\]

Instead it has the schematic form

\[
\boxed{
\frac{\delta S_{\mathrm{eff}}}{\delta\lambda^I(x)}
=\mathcal K_{IJ}(\lambda,\partial\lambda,\ldots)\,
\beta^J(x),
}
\]

up to target-space gauge transformations and higher-order corrections.
The kernel \(\mathcal K_{IJ}\) is a field-space metric or, more
generally, an invertible differential operator. It also permits mixing
among the metric and dilaton equations. Therefore the invariant
statement is

\[
\boxed{
\beta^I=0
\quad\Longleftrightarrow\quad
\frac{\delta S_{\mathrm{eff}}}{\delta\lambda^I}=0,
}
\]

provided \(\mathcal K\) is nondegenerate after gauge redundancies have
been removed and both sides are compared in the same renormalization
scheme and at the same order in \(\alpha'\) and \(g_s\).

Equations (14.5) and (14.8) give an explicit leading-order
demonstration. Varying the string-frame action (14.8) gives equations
that are invertible linear combinations of the three Weyl-anomaly
coefficients in (14.5). Thus their common zero locus is exactly the
same, even though an individual raw Euler-Lagrange derivative is not
identical to an individual beta function.

\subsubsection{Background fields are worldsheet
couplings}\label{background-fields-are-worldsheet-couplings}

The bosonic part of the NS-NS sigma model in equation (14.3) is

\[
S_\sigma
=-\frac{1}{4\pi\alpha'}
\int_\Sigma d^2\sigma\,\sqrt{-h}
\left[
h^{ab}\partial_aX^\mu\partial_bX^\nu G_{\mu\nu}(X)
+\epsilon^{ab}\partial_aX^\mu\partial_bX^\nu B_{\mu\nu}(X)
+\alpha'\Phi(X)R^{(2)}
\right].
\]

From the two-dimensional viewpoint, this is a quantum field theory with
infinitely many couplings: every value of \(G_{\mu\nu}(x)\),
\(B_{\mu\nu}(x)\), and \(\Phi(x)\) is part of the coupling data.
Schematically,

\[
S_\sigma
=S_*+\sum_I\int_\Sigma d^2z\,
\lambda^I(X)\mathcal O_I(z,\bar z).
\]

For example,

\[
\mathcal O^G_{\mu\nu}
\sim \partial X^\mu\bar\partial X^\nu,
\qquad
\mathcal O^B_{\mu\nu}
\sim \partial X^{[\mu}\bar\partial X^{\nu]},
\qquad
\mathcal O^\Phi
\sim R^{(2)}.
\]

Expanding a background around flat spacetime makes the relation to
string states explicit:

\[
G_{\mu\nu}(X)
=\eta_{\mu\nu}+h_{\mu\nu}(X).
\]

Fourier expanding

\[
h_{\mu\nu}(X)
=\int\frac{d^dk}{(2\pi)^d}\,
h_{\mu\nu}(k)e^{ik\cdot X}
\]

turns the perturbation of the sigma-model action into a coherent sum of
integrated graviton vertex operators,

\[
\delta S_\sigma
\sim
\int\frac{d^dk}{(2\pi)^d}\,
h_{\mu\nu}(k)
\int_\Sigma d^2z\,
\partial X^\mu\bar\partial X^\nu e^{ik\cdot X}.
\]

Thus a spacetime field profile is literally a collection of worldsheet
couplings to string vertex operators. This is why a condition on
sigma-model couplings can become a differential equation for spacetime
fields.

\subsubsection{Why Weyl invariance requires the beta functions to
vanish}\label{why-weyl-invariance-requires-the-beta-functions-to-vanish}

After fixing the worldsheet metric to conformal gauge,

\[
h_{ab}=e^{2\omega}\hat h_{ab},
\]

the Weyl factor \(\omega\) is a gauge degree of freedom. It must not
become a physical field merely because of a quantum anomaly.
Equivalently, the trace of the renormalized worldsheet stress tensor
must vanish, apart from the central-charge term whose matter and ghost
contributions must cancel.

For spacetime-dependent couplings, the local Weyl anomaly has the form
displayed schematically in equation (14.4):

\[
T^a{}_a
=
\beta^G_{\mu\nu}\,
h^{ab}\partial_aX^\mu\partial_bX^\nu
+\beta^B_{\mu\nu}\,
\epsilon^{ab}\partial_aX^\mu\partial_bX^\nu
+\beta^\Phi R^{(2)}
+\text{total derivatives}.
\]

The operators multiplying the independent coefficients are independent
local worldsheet operators. Consequently, quantum Weyl invariance
requires the corresponding Weyl-anomaly coefficients to vanish, modulo
redundant operators generated by target-space gauge transformations:

\[
\beta^G_{\mu\nu}=0,
\qquad
\beta^B_{\mu\nu}=0,
\qquad
\beta^\Phi=0.
\]

At linear order around flat space these are just the BRST physical-state
conditions on the graviton, antisymmetric tensor, and dilaton. At
nonlinear order they include the interactions among those fields. The
sigma-model beta functions are therefore nonlinear completions of the
free mass-shell and polarization constraints.

\subsubsection{The leading beta functions in the
book}\label{the-leading-beta-functions-in-the-book}

For the type-II NS-NS fields, equation (14.5) gives

\[
\beta^G_{\mu\nu}
=\alpha'
\left(
R_{\mu\nu}
-\frac14H_{\mu\rho\sigma}H_\nu{}^{\rho\sigma}
+2\nabla_\mu\nabla_\nu\Phi
\right)
+O(\alpha'^2),
\]

\[
\beta^B_{\mu\nu}
=\alpha'
\left(
-\frac12\nabla^\rho H_{\rho\mu\nu}
+H_{\rho\mu\nu}\nabla^\rho\Phi
\right)
+O(\alpha'^2),
\]

and

\[
\beta^\Phi
=\frac14\Delta
+\alpha'
\left[
(\nabla\Phi)^2
-\frac12\nabla^2\Phi
-\frac1{24}H^2
\right]
+O(\alpha'^2),
\qquad
\Delta=d-d_{\mathrm{crit}}.
\]

The first two equations say that the metric and \(B\)-field couplings
cease to run only when the spacetime fields satisfy particular
differential equations. The constant term in \(\beta^\Phi\) is the
central-charge deficit, including the ghost contribution.

\subsubsection{The string-frame action}\label{the-string-frame-action}

The spacetime action displayed in equation (14.8) is, up to an overall
normalization,

\[
\boxed{
S_{\mathrm{eff}}
=
\int d^dx\,\sqrt{-G}\,e^{-2\Phi}
\left[
R
+4(\nabla\Phi)^2
-\frac1{12}H^2
-\frac{\Delta}{\alpha'}
\right].
}
\]

This is the tree-level NS-NS action in string frame to the derivative
order shown in (14.5). The factor \(e^{-2\Phi}\) reflects the sphere
Euler characteristic and hence the closed-string tree-level factor
\(g_s^{-2}\).

We now vary this action and compare the resulting equations directly
with the beta functions.

\subsubsection{\texorpdfstring{Variation with respect to the
\(B\)-field}{Variation with respect to the B-field}}\label{variation-with-respect-to-the-b-field}

Since

\[
H_{\mu\nu\rho}=3\partial_{[\mu}B_{\nu\rho]},
\]

its variation is

\[
\delta H_{\mu\nu\rho}
=3\nabla_{[\mu}\delta B_{\nu\rho]}.
\]

Only the \(H^2\) term varies:

\[
\delta_BS_{\mathrm{eff}}
=-\frac16
\int d^dx\,\sqrt{-G}\,e^{-2\Phi}
H^{\mu\nu\rho}\delta H_{\mu\nu\rho}.
\]

Using antisymmetry,

\[
H^{\mu\nu\rho}\delta H_{\mu\nu\rho}
=3H^{\mu\nu\rho}\nabla_\mu\delta B_{\nu\rho},
\]

and integrating by parts gives

\[
\delta_BS_{\mathrm{eff}}
=\frac12
\int d^dx\,
\delta B_{\nu\rho}\,
\nabla_\mu
\left(
\sqrt{-G}\,e^{-2\Phi}H^{\mu\nu\rho}
\right).
\]

Hence the \(B\)-field equation is

\[
\nabla_\mu\left(e^{-2\Phi}H^{\mu\nu\rho}\right)=0,
\]

or

\[
\nabla_\mu H^{\mu\nu\rho}
-2(\nabla_\mu\Phi)H^{\mu\nu\rho}=0.
\]

Multiplying by \(-\alpha'/2\) and lowering indices gives exactly

\[
\boxed{
\beta^B_{\nu\rho}=0.
}
\]

Thus the \(B\)-field beta function and Euler-Lagrange equation agree up
to a nonzero normalization.

\subsubsection{Variation with respect to the
dilaton}\label{variation-with-respect-to-the-dilaton}

Define

\[
\mathcal L
=R+4(\nabla\Phi)^2-\frac1{12}H^2-\frac{\Delta}{\alpha'}.
\]

Both the prefactor \(e^{-2\Phi}\) and the explicit kinetic term depend
on \(\Phi\). The variation is

\[
\delta_\Phi S_{\mathrm{eff}}
=\int d^dx\,\sqrt{-G}\,e^{-2\Phi}
\left[
-2\mathcal L\,\delta\Phi
+8\nabla_\mu\Phi\,\nabla^\mu\delta\Phi
\right].
\]

Integrating the second term by parts with the weighted measure gives

\[
\int\sqrt{-G}\,e^{-2\Phi}
8\nabla_\mu\Phi\,\nabla^\mu\delta\Phi
=
-8\int\sqrt{-G}\,e^{-2\Phi}
\left[
\nabla^2\Phi-2(\nabla\Phi)^2
\right]\delta\Phi.
\]

Therefore the dilaton equation is

\[
\boxed{
\mathcal E_\Phi
\equiv
R
+4\nabla^2\Phi
-4(\nabla\Phi)^2
-\frac1{12}H^2
-\frac{\Delta}{\alpha'}
=0.
}
\]

This is not equal to \(\beta^\Phi=0\) by itself. Instead, taking the
trace of \(\beta^G_{\mu\nu}\) gives

\[
\frac1{\alpha'}G^{\mu\nu}\beta^G_{\mu\nu}
=
R-\frac14H^2+2\nabla^2\Phi
+O(\alpha').
\]

Combining this with \(\beta^\Phi\) yields

\[
\boxed{
\mathcal E_\Phi
=
\frac1{\alpha'}
\left(
G^{\mu\nu}\beta^G_{\mu\nu}
-4\beta^\Phi
\right)
+O(\alpha').
}
\]

This is the first explicit instance of field-space mixing: the dilaton
Euler-Lagrange equation is a linear combination of the metric and
dilaton Weyl-anomaly coefficients.

\subsubsection{Variation with respect to the
metric}\label{variation-with-respect-to-the-metric}

It is useful to do the metric variation without first imposing the
dilaton equation. Varying with respect to the inverse metric, the
relevant identity is

\[
\delta\!\left(\sqrt{-G}\,fR\right)
=\sqrt{-G}
\left[
f\left(R_{\mu\nu}-\frac12G_{\mu\nu}R\right)
+\left(G_{\mu\nu}\nabla^2-\nabla_\mu\nabla_\nu\right)f
\right]\delta G^{\mu\nu}
\]

up to a boundary term. Here \(f=e^{-2\Phi}\), so

\[
\frac{\nabla_\mu\nabla_\nu f}{f}
=-2\nabla_\mu\nabla_\nu\Phi
+4\nabla_\mu\Phi\nabla_\nu\Phi,
\]

\[
\frac{\nabla^2 f}{f}
=-2\nabla^2\Phi+4(\nabla\Phi)^2.
\]

The other metric variations are

\[
\delta\sqrt{-G}
=-\frac12\sqrt{-G}\,G_{\mu\nu}\delta G^{\mu\nu},
\]

\[
\delta\bigl((\nabla\Phi)^2\bigr)
=\nabla_\mu\Phi\nabla_\nu\Phi\,\delta G^{\mu\nu},
\]

and, because \(H^2\) contains three inverse metrics,

\[
\delta H^2
=3H_{\mu\rho\sigma}H_\nu{}^{\rho\sigma}\delta G^{\mu\nu}.
\]

After collecting these terms, the coefficient of \(\delta G^{\mu\nu}\)
can be written particularly transparently in terms of the dilaton
equation already found:

\[
\delta_GS_{\mathrm{eff}}
=\int d^dx\,\sqrt{-G}\,e^{-2\Phi}
\left[
R_{\mu\nu}
+2\nabla_\mu\nabla_\nu\Phi
-\frac14H_{\mu\rho\sigma}H_\nu{}^{\rho\sigma}
-\frac12G_{\mu\nu}\mathcal E_\Phi
\right]\delta G^{\mu\nu}.
\]

Thus a convenient representative of the metric Euler-Lagrange equation
is

\[
\boxed{
\mathcal E^G_{\mu\nu}
\equiv
R_{\mu\nu}
+2\nabla_\mu\nabla_\nu\Phi
-\frac14H_{\mu\rho\sigma}H_\nu{}^{\rho\sigma}
-\frac12G_{\mu\nu}\mathcal E_\Phi
=0.
}
\]

Comparison with \(\beta^G_{\mu\nu}\) gives

\[
\boxed{
\mathcal E^G_{\mu\nu}
=\frac1{\alpha'}\beta^G_{\mu\nu}
-\frac12G_{\mu\nu}\mathcal E_\Phi
+O(\alpha').
}
\]

The term proportional to \(G_{\mu\nu}\mathcal E_\Phi\) is why the raw
metric variation is not literally \(\beta^G_{\mu\nu}/\alpha'\). Once the
dilaton equation is imposed, however, the metric equation reduces
precisely to \(\beta^G_{\mu\nu}=0\) at this order.

\subsubsection{Equivalence of the complete sets of
equations}\label{equivalence-of-the-complete-sets-of-equations}

The three variations can now be compared with the three beta functions
without any hand-waving. To the order displayed in the book,

\[
\mathcal E^B_{\mu\nu}=0
\quad\Longleftrightarrow\quad
\beta^B_{\mu\nu}=0,
\]

\[
\mathcal E_\Phi
=\frac1{\alpha'}
\left(G^{\mu\nu}\beta^G_{\mu\nu}-4\beta^\Phi\right)
+O(\alpha'),
\]

\[
\mathcal E^G_{\mu\nu}
=\frac1{\alpha'}\beta^G_{\mu\nu}
-\frac12G_{\mu\nu}\mathcal E_\Phi
+O(\alpha').
\]

Therefore:

\begin{enumerate}
\def\labelenumi{\arabic{enumi}.}
\tightlist
\item
  If all beta functions vanish, then \(\mathcal E_\Phi=0\),
  \(\mathcal E^G_{\mu\nu}=0\), and \(\mathcal E^B_{\mu\nu}=0\).
\item
  Conversely, if all Euler-Lagrange equations vanish, the metric
  relation gives \(\beta^G_{\mu\nu}=0\), the dilaton relation then gives
  \(\beta^\Phi=0\), and the \(B\)-field equation gives
  \(\beta^B_{\mu\nu}=0\).
\end{enumerate}

Hence

\[
\boxed{
\left\{
\frac{\delta S_{\mathrm{eff}}}{\delta G}=0,
\frac{\delta S_{\mathrm{eff}}}{\delta B}=0,
\frac{\delta S_{\mathrm{eff}}}{\delta\Phi}=0
\right\}
\quad\Longleftrightarrow\quad
\left\{
\beta^G=0,
\beta^B=0,
\beta^\Phi=0
\right\}
}
\]

to the common order in \(\alpha'\) and \(g_s\). The assertion is
equality of the two solution sets, not component-by-component identity
between every variational derivative and the beta function carrying the
same label.

More generally, if all couplings are denoted by \(\lambda^I\), the
relation takes the schematic form

\[
\frac{1}{\sqrt{-G}}\frac{\delta S_{\mathrm{eff}}}{\delta\lambda^I(x)}
=e^{-2\Phi}\mathcal K_{IJ}(\lambda)\beta^J(x),
\]

where \(\mathcal K_{IJ}\) is a local field-space kernel. After
quotienting by redundant directions associated with diffeomorphisms and
gauge transformations, \(\mathcal K_{IJ}\) is nondegenerate
perturbatively. It may mix fields, exactly as the explicit
\(G\)-\(\Phi\) relations above do, but it does not change the common
zero locus.

\subsubsection{Why the worldsheet and spacetime computations encode the
same
dynamics}\label{why-the-worldsheet-and-spacetime-computations-encode-the-same-dynamics}

This agreement is structural rather than accidental. A variation of a
background field changes the worldsheet action by an integrated vertex
operator:

\[
\delta S_\sigma
=\delta\lambda^I\int d^2z\,\mathcal O_I(z,\bar z).
\]

Thus the background coupling \(\lambda^I\) is the spacetime field whose
fluctuation is created by \(\mathcal O_I\). At linear order, requiring
\(\mathcal O_I\) to be a marginal, BRST-invariant vertex operator gives
the free spacetime equation of motion. At nonlinear order, operator
products generate the beta functions; their vanishing is the condition
that the deformed worldsheet theory remain conformal.

The spacetime effective action packages the same string amplitudes in
target-space language. Its first variation is the one-point function, or
tadpole, of the corresponding string state in the chosen background. A
nonzero beta function is likewise a failure of the background to solve
the string consistency conditions. Both computations are therefore
measuring the same obstruction: the background is off shell.

String field theory makes this relation especially direct. The string
field contains the coefficients of vertex operators, including
\(G_{\mu\nu}\), \(B_{\mu\nu}\), and \(\Phi\) among its massless
components. Its classical equation of motion is the nonlinear BRST
consistency condition. Expanding that equation in components gives the
spacetime field equations, while interpreting the same deformation as a
two-dimensional QFT gives the conformal-invariance conditions. The
low-energy effective action is the massless, derivative-expanded form of
this string-field-theoretic dynamics.

This also explains why reconstructing the action from on-shell
scattering amplitudes is not a third, unrelated mechanism. The
amplitudes determine interaction vertices modulo field redefinitions and
terms proportional to lower-order equations of motion. The sigma-model
beta functions determine the same equivalence class by demanding quantum
Weyl invariance.

\subsubsection{What ``exactly'' means and what it does not
mean}\label{what-exactly-means-and-what-it-does-not-mean}

Several qualifications are essential.

First, one should use the Weyl-anomaly coefficients. Ordinary
renormalization-group beta functions can differ from them by
target-space diffeomorphisms or \(B\)-field gauge transformations. Those
differences correspond to redundant worldsheet operators and do not
alter the physical conformal background.

Second, the comparison must be made in a common renormalization scheme.
A local field redefinition,

\[
\lambda'^I=\lambda'^I(\lambda),
\]

changes the beta functions as a vector field on the space of couplings,

\[
\beta'^I
=\frac{\partial\lambda'^I}{\partial\lambda^J}\beta^J,
\]

and rewrites the spacetime action in the new variables. Their component
formulas can change, but the condition \(\beta^I=0\) and the physical
content of the Euler-Lagrange equations do not.

Third, orders must be matched. A one-loop sigma-model calculation
produces the leading two-derivative spacetime equations. Higher
sigma-model loops produce higher powers of \(\alpha'\) and must be
compared with higher-derivative terms such as curvature-squared or
curvature-quartic corrections in the effective action. Separately,
worldsheet genus counts powers of \(g_s\): equation (14.8) is a
sphere-level, or closed-string tree-level, action, so quantum spacetime
corrections require higher-genus contributions. The sigma-model loop
expansion in \(\alpha'\) and the genus expansion in \(g_s\) are distinct
expansions.

Fourth, equation (14.8) is a truncation to the NS-NS massless fields. If
additional massless fields, Ramond-Ramond backgrounds, open-string
fields, or massive string modes are sourced, their equations and beta
functions must be included. Setting omitted fields to zero is valid only
when that is a consistent truncation.

Finally, the effective action is not unique as a formula. Total
derivatives, local field redefinitions, and terms proportional to
lower-order equations of motion can change its appearance without
changing on-shell physics. The invariant statement is the equivalence
class of equations and their solution space.

\subsubsection{Checks}\label{checks-29}

\textbf{Critical, constant-dilaton, fluxless limit.} Take \(\Delta=0\),
\(\nabla\Phi=0\), and \(H=0\). The beta functions give

\[
\beta^G_{\mu\nu}=\alpha'R_{\mu\nu},
\qquad
\beta^\Phi=0,
\]

so conformal invariance requires \(R_{\mu\nu}=0\). The action reduces to
the Einstein-Hilbert action. Its raw metric equation is
\(R_{\mu\nu}-\tfrac12G_{\mu\nu}R=0\); tracing it for \(d\neq2\) gives
\(R=0\), and hence \(R_{\mu\nu}=0\). This simple example already shows
why identical solution sets do not require identical component
expressions.

\textbf{Derivative counting.} Each leading beta function contains two
target-space derivatives: \(R_{\mu\nu}\) has two derivatives of \(G\),
\(\nabla H\) has two derivatives of \(B\), and \(\nabla^2\Phi\) has two
derivatives of \(\Phi\). This matches the two-derivative action (14.8).

\textbf{Gauge invariance.} The action depends on \(B\) only through
\(H=dB\), so it is invariant under \(B\mapsto B+d\Lambda\).
Correspondingly, the \(B\) equation is a divergence equation, and its
beta function is defined modulo the same redundant gauge direction.

\textbf{Linearized spectrum.} Expanding around
\(G_{\mu\nu}=\eta_{\mu\nu}\), constant \(\Phi\), and \(B=0\), the
linearized beta-function equations reproduce the massless wave equations
and transversality conditions after gauge fixing. These are the same
conditions obtained from BRST invariance of the massless vertex
operators.

\subsubsection{Result}\label{result-33}

The sigma-model background fields are worldsheet couplings,

\[
S_\sigma=S_\sigma[\lambda^I],
\qquad
\lambda^I=\{G_{\mu\nu},B_{\mu\nu},\Phi,\ldots\},
\]

and quantum Weyl invariance requires

\[
\boxed{\beta^I=0.}
\]

At leading order, the same equations follow from

\[
\boxed{
S_{\mathrm{eff}}
=\int d^dx\,\sqrt{-G}\,e^{-2\Phi}
\left[
R+4(\nabla\Phi)^2-\frac1{12}H^2-\frac{d-d_{\mathrm{crit}}}{\alpha'}
\right].
}
\]

Explicitly,

\[
\mathcal E^B=0\Longleftrightarrow\beta^B=0,
\]

\[
\mathcal E_\Phi
=\frac1{\alpha'}\left(G^{\mu\nu}\beta^G_{\mu\nu}-4\beta^\Phi\right)
+O(\alpha'),
\]

\[
\mathcal E^G_{\mu\nu}
=\frac1{\alpha'}\beta^G_{\mu\nu}
-\frac12G_{\mu\nu}\mathcal E_\Phi
+O(\alpha').
\]

Consequently,

\[
\boxed{
\frac{\delta S_{\mathrm{eff}}}{\delta\lambda^I}=0
\quad\Longleftrightarrow\quad
\beta^I=0,
}
\]

provided the same fields, \(\alpha'\) order, string-loop order,
renormalization scheme, and gauge-equivalence class are used. The
variational derivatives and beta functions need not be literally
identical components; they are related by an invertible field-space
transformation and have the same physical zero locus.

\chapter{Type-II Compactifications with D-Branes and
Fluxes}\label{type-ii-compactifications-with-d-branes-and-fluxes-1}

\section{Q1: Special-Lagrangian deformations and Wilson-line
cohomology}\label{q1-special-lagrangian-deformations-and-wilson-line-cohomology}

In page650, it is said that: \textgreater Additional non-chiral matter
transforming in the adjoint representation of \(U(N_{a})\) arises from
open strings stretched between branes in the same stack, i.e.~branes
lying on top of each other. Geometrically, these correspond to
deformations of the three-cycle inside the Calabi-Yau or to continuous
Wilson lines along non-trivial one-cycles inside the three-cycles. They
are counted by \(H^0(\Sigma,N_{\Sigma})\) and \(H^1(\Sigma)\)
respectively (\(\Sigma\) is a special Lagrangian inside the Calabi-Yau).

Why the number of these states are counted by these cohomologies?

\begin{quote}
\textbf{Textbook location and related material.} Chapter 17, Section
17.2 around p.~650: same-stack open strings, geometric brane
deformations, Wilson lines, and adjoint matter. Sections 14.4--14.5 give
the normal-bundle and special-Lagrangian geometry.
\end{quote}

\subsection{Answer}\label{answer-34}

\subsubsection{Core point}\label{core-point-34}

There are three related but logically distinct statements in the quoted
passage:

\begin{enumerate}
\def\labelenumi{\arabic{enumi}.}
\tightlist
\item
  an arbitrary infinitesimal deformation of an embedded submanifold is a
  section of its normal bundle;
\item
  a deformation of a \textbf{special Lagrangian} submanifold is much
  more constrained, and its massless deformation space is isomorphic to
  \(H^1(\Sigma,\mathbb R)\);
\item
  an infinitesimal continuous Wilson line is a closed one-form modulo an
  exact one-form, so it is also represented by
  \(H^1(\Sigma,\mathbb R)\).
\end{enumerate}

The first statement explains the appearance of \(H^0(\Sigma,N_\Sigma)\).
The second and third explain why the actual finite multiplicity of
massless adjoint matter is \(b_1(\Sigma)\).

More precisely, the book's \(H^0(\Sigma,N_\Sigma)\) should not be read
as saying that the multiplicity is literally \[
\dim H^0(\Sigma,N_\Sigma)
\] with \(H^0\) the space of all smooth normal sections. That space is
infinite-dimensional. It is the kinematic space in which the normal
deformation field takes values. The massless supersymmetric modes are
the sections satisfying the linearized special-Lagrangian equations.
McLean's theorem gives \[
\boxed{
T_{[\Sigma]}\mathcal M_{\mathrm{sLag}}
\cong
\mathcal H^1(\Sigma)
\cong
H^1_{\mathrm{dR}}(\Sigma,\mathbb R)
},
\] and therefore \[
\dim_{\mathbb R}\mathcal M_{\mathrm{sLag}}=b_1(\Sigma).
\]

The Wilson-line moduli form another real space of dimension
\(b_1(\Sigma)\). Four-dimensional \(\mathcal N=1\) supersymmetry pairs
each real geometric modulus with one real Wilson-line modulus. Thus a
stack with gauge group \(U(N_a)\) has \[
\boxed{b_1(\Sigma)\text{ complex scalars in }\operatorname{Adj}U(N_a)}.
\]

Here ``non-chiral matter'' means that the adjoint representation is real
and therefore produces no net four-dimensional gauge chirality. The
fields can still be packaged into \(\mathcal N=1\) chiral
supermultiplets.

\subsubsection{Why same-stack strings give adjoint-valued
fields}\label{why-same-stack-strings-give-adjoint-valued-fields}

Consider \(N_a\) coincident D6-branes wrapping \[
\mathbb R^{1,3}\times\Sigma.
\] An open string whose two endpoints lie on this stack carries
Chan--Paton labels \((r,s)\), with \(r,s=1,\ldots,N_a\). Its state is
therefore accompanied by an \(N_a\times N_a\) matrix \(\lambda\). Under
a gauge transformation \(U\in U(N_a)\), \[
\lambda\longmapsto U\lambda U^{-1}.
\] Consequently the corresponding worldvolume fields take values in \[
\operatorname{End}(\mathbb C^{N_a})\cong\mathfrak u(N_a),
\] which is the adjoint representation.

Before compactification, the massless bosonic fields on a D6-brane are a
seven-dimensional gauge field together with three transverse scalar
fields. After wrapping the D6-brane on \(\Sigma\):

\begin{itemize}
\tightlist
\item
  the components \(A_\mu\), \(\mu=0,\ldots,3\), give the
  four-dimensional gauge field;
\item
  the internal components \(A_m\), \(m=1,2,3\) along \(\Sigma\), give
  possible Wilson-line scalars;
\item
  the three transverse scalar fields combine geometrically into a
  section of the rank-three real normal bundle \(N_\Sigma\).
\end{itemize}

The question is which internal profiles of these seven-dimensional
fields have zero Kaluza--Klein mass. Those are precisely the zero modes
of the appropriate internal differential operators.

\subsubsection{\texorpdfstring{Arbitrary embedded deformations and
\(H^0(\Sigma,N_\Sigma)\)}{Arbitrary embedded deformations and H\^{}0(\textbackslash Sigma,N\_\textbackslash Sigma)}}\label{arbitrary-embedded-deformations-and-h0sigman_sigma}

Let \(i:\Sigma\hookrightarrow X\) be the embedding. An infinitesimal
variation of the embedding is a vector field along \(\Sigma\), \[
\delta i\in\Gamma(i^*TX).
\] Using the induced metric, it decomposes into tangent and normal
parts, \[
\delta i=(\delta i)^\parallel+(\delta i)^\perp.
\] The tangent part changes only the parametrization of \(\Sigma\); it
is generated by an infinitesimal diffeomorphism of the worldvolume and
does not change the embedded submanifold. The physical infinitesimal
displacement is therefore the normal part, \[
v=(\delta i)^\perp\in\Gamma(N_\Sigma).
\] Equivalently, \[
T_{[i]}\frac{\operatorname{Emb}(\Sigma,X)}{\operatorname{Diff}(\Sigma)}
\cong
\Gamma(N_\Sigma)
=H^0(\Sigma,N_\Sigma).
\]

This is the precise kinematic meaning of the first cohomology group in
the quoted sentence. A zeroth cohomology group records global sections:
one must choose the normal displacement consistently over all of
\(\Sigma\), not independently in unrelated coordinate patches.

There is an important qualification. Since \(\Sigma\) is a real
three-manifold and \(N_\Sigma\) is a real smooth vector bundle,
\(\Gamma(N_\Sigma)\) is generically infinite-dimensional. Most of its
sections do not preserve the Lagrangian or special condition. They
change the D6-brane energy and are not massless moduli. The BPS and
zero-mode equations reduce this kinematic space to a finite-dimensional
subspace.

This differs from the familiar statement that deformations of a
\textbf{complex} submanifold are counted by holomorphic sections of a
holomorphic normal bundle. In that B-brane setting, a Dolbeault group
such as \(H^0(\Sigma,N^{1,0}_\Sigma)\) can already be
finite-dimensional. A special Lagrangian D6-brane is an A-brane, and the
relevant finite-dimensional theorem is instead McLean's theorem.

\subsubsection{The Lagrangian condition identifies normal vectors with
one-forms}\label{the-lagrangian-condition-identifies-normal-vectors-with-one-forms}

Because \(\Sigma\) is Lagrangian, \[
i^*\omega=0.
\] This implies that \(J\) exchanges tangent and normal directions: \[
J(T\Sigma)=N_\Sigma,
\qquad
J(N_\Sigma)=T\Sigma.
\] Indeed, if \(u,w\in T\Sigma\), then \[
g(Ju,w)=\omega(u,w)=0,
\] so \(Ju\) is orthogonal to \(T\Sigma\) and hence normal. Since both
bundles have real rank three, this is an isomorphism.

The metric then identifies a normal vector field \(v\) with a one-form
on \(\Sigma\). A convenient sign convention is \[
\alpha_v=-i^*(\iota_v\omega),
\qquad
\alpha_v(w)=-\omega(v,w),\quad w\in T\Sigma.
\] The map \[
N_\Sigma\longrightarrow T^*\Sigma,
\qquad
v\longmapsto\alpha_v,
\] is a pointwise bundle isomorphism. Thus the three normal components
of the D6-brane deformation may equivalently be treated as the
components of a one-form on \(\Sigma\).

To see this explicitly, choose at a point of \(\Sigma\) an oriented
orthonormal tangent basis \(e_1,e_2,e_3\). Then \(Je_1,Je_2,Je_3\) is a
normal basis. If \[
v=v^IJe_I,
\] then, in the corresponding dual tangent coframe \(e^I\), \[
\alpha_v=v^Ie^I.
\] No information about \(v\) is lost in passing to \(\alpha_v\).

\subsubsection{\texorpdfstring{Linearizing the Lagrangian condition
gives
\(d\alpha_v=0\)}{Linearizing the Lagrangian condition gives d\textbackslash alpha\_v=0}}\label{linearizing-the-lagrangian-condition-gives-dalpha_v0}

Let \(i_t:\Sigma\hookrightarrow X\) be a one-parameter family of
embeddings with \(i_0=i\) and normal velocity \(v\). For every
differential form \(\rho\) on \(X\), \[
\left.\frac{d}{dt}\right|_{t=0}i_t^*\rho
=i^*(\mathcal L_v\rho).
\] Using Cartan's formula, \[
\mathcal L_v\rho=d(\iota_v\rho)+\iota_v(d\rho).
\] The Calabi--Yau Kähler form is closed, \(d\omega=0\). Therefore \[
\left.\frac{d}{dt}\right|_{t=0}i_t^*\omega
=i^*d(\iota_v\omega)
=-d\alpha_v.
\] For the deformed cycle to remain Lagrangian to first order, this
variation must vanish. Hence \[
\boxed{d\alpha_v=0}.
\]

Thus infinitesimal Lagrangian deformations correspond locally to closed
one-forms. If one imposed only the Lagrangian condition, exact one-forms
would also occur; the special condition supplies the second equation
needed for harmonicity.

\subsubsection{\texorpdfstring{Linearizing the special condition gives
\(d^\dagger\alpha_v=0\)}{Linearizing the special condition gives d\^{}\textbackslash dagger\textbackslash alpha\_v=0}}\label{linearizing-the-special-condition-gives-ddaggeralpha_v0}

For the phase \(\theta\) selected by the orientifold and the preserved
\(\mathcal N=1\) supersymmetry, the second special-Lagrangian condition
is \[
i^*\operatorname{Im}(e^{-i\theta}\Omega)=0.
\] Since \(\Omega\) is closed, \[
\left.\frac{d}{dt}\right|_{t=0}
i_t^*\operatorname{Im}(e^{-i\theta}\Omega)
=d\,i^*\!\left(
\iota_v\operatorname{Im}(e^{-i\theta}\Omega)
\right).
\]

On a special Lagrangian three-cycle there is a pointwise identity \[
i^*\!\left(
\iota_v\operatorname{Im}(e^{-i\theta}\Omega)
\right)
=*_\Sigma\alpha_v.
\] With a different sign convention for \(\alpha_v\) the right-hand side
acquires an overall minus sign, which does not affect the resulting
equation.

The identity can be checked in an adapted orthonormal frame. Let \(f^I\)
be the coframe dual to the normal vectors \(Je_I\). At the chosen point
one may write \[
\omega=\sum_{I=1}^3 e^I\wedge f^I,
\qquad
e^{-i\theta}\Omega
=(e^1+if^1)\wedge(e^2+if^2)\wedge(e^3+if^3).
\] For \(v=v^IJe_I\), contracting the terms in
\(\operatorname{Im}(e^{-i\theta}\Omega)\) containing one normal coframe
gives \[
i^*(\iota_v\operatorname{Im}(e^{-i\theta}\Omega))
=v^1e^2\wedge e^3-v^2e^1\wedge e^3+v^3e^1\wedge e^2
=*_\Sigma(v^Ie^I)
=*_\Sigma\alpha_v.
\]

The linearized special condition is consequently \[
d(*_\Sigma\alpha_v)=0.
\] On a three-manifold, for a one-form \(\alpha\), \[
d^\dagger\alpha=-*_\Sigma d *_\Sigma\alpha.
\] Therefore \[
\boxed{d^\dagger\alpha_v=0}.
\]

Combining the two linearized equations gives \[
d\alpha_v=0,
\qquad
d^\dagger\alpha_v=0,
\] or equivalently \[
\Delta_1\alpha_v=0.
\] Hence \(\alpha_v\) is a harmonic one-form.

\subsubsection{McLean's theorem and the finite deformation
count}\label{mcleans-theorem-and-the-finite-deformation-count}

The preceding calculation determines the tangent space to the
special-Lagrangian deformation problem: \[
\left\{
v\in\Gamma(N_\Sigma):
\delta_v(i^*\omega)=0,
\quad
\delta_v i^*\operatorname{Im}(e^{-i\theta}\Omega)=0
\right\}
\cong
\mathcal H^1(\Sigma).
\] Hodge theory gives a unique harmonic representative of every de Rham
cohomology class, so \[
\mathcal H^1(\Sigma)\cong H^1_{\mathrm{dR}}(\Sigma,\mathbb R).
\]

McLean's theorem says more than this linear identification: for a
compact special Lagrangian, these infinitesimal deformations are
unobstructed. Locally, the moduli space of special Lagrangian
deformations is a smooth manifold whose tangent space is
\(H^1(\Sigma,\mathbb R)\). Therefore \[
\boxed{
\#\{\text{real geometric moduli}\}
=b_1(\Sigma)
}.
\]

This is the precise content behind the book's statement on the following
page that ``for special Lagrangian cycles the number of these two kinds
of moduli are equal.'' The normal-bundle notation describes where the
deformation field lives; the special-Lagrangian equations and McLean's
theorem determine its finite set of massless zero modes.

The same result also has a direct Kaluza--Klein interpretation. Under
the bundle isomorphism \(N_\Sigma\cong T^*\Sigma\), the quadratic
operator governing supersymmetric normal fluctuations reduces to the
one-form Hodge Laplacian. If \[
\Delta_1\eta_n=\lambda_n\eta_n,
\] then the corresponding four-dimensional field has a Kaluza--Klein
mass proportional to \(\lambda_n\). The massless modes have
\(\lambda_n=0\), so they are exactly the harmonic one-forms.

\subsubsection{\texorpdfstring{Why continuous Wilson lines are counted
by
\(H^1(\Sigma)\)}{Why continuous Wilson lines are counted by H\^{}1(\textbackslash Sigma)}}\label{why-continuous-wilson-lines-are-counted-by-h1sigma}

Now consider the internal part \(a\) of the D6-brane gauge field. First
take a single brane, or the Abelian part of a coincident stack.
Infinitesimally, \[
a\in\Omega^1(\Sigma).
\] A Wilson-line modulus is a deformation of a flat connection, so its
linearized field strength must vanish: \[
da=0.
\] An infinitesimal gauge transformation with parameter
\(\lambda\in\Omega^0(\Sigma)\) acts by \[
a\longmapsto a+d\lambda.
\] The space of inequivalent infinitesimal flat connections is therefore
\[
\frac{\ker(d:\Omega^1\to\Omega^2)}
{\operatorname{im}(d:\Omega^0\to\Omega^1)}
=H^1_{\mathrm{dR}}(\Sigma,\mathbb R).
\]

Equivalently, in a gauge such as \(d^\dagger a=0\), the zero-mode
equation is \[
\Delta_1 a=0.
\] Thus one may expand \[
a(x,y)=\sum_{I=1}^{b_1(\Sigma)}\xi^I(x)\eta_I(y)
+\text{massive modes},
\] where the \(\eta_I\) form a basis of harmonic one-forms. Each
coefficient \(\xi^I(x)\) is a real massless scalar in four dimensions.

The terminology ``Wilson line'' comes from its holonomy around a closed
one-cycle \(\gamma\subset\Sigma\), \[
W_\gamma
=\exp\!\left(i\oint_\gamma A\right).
\] If \(a=d\lambda\) is pure gauge, its integral around every closed
cycle vanishes. A non-exact closed one-form can have nonzero periods and
changes these holonomies. Globally, for a \(U(1)\) connection the
continuous moduli space is the torus \[
H^1(\Sigma,\mathbb R)\big/2\pi H^1(\Sigma,\mathbb Z)
\cong
\operatorname{Hom}(H_1(\Sigma,\mathbb Z),U(1))_{\mathrm{continuous}},
\] whose real dimension is \(b_1(\Sigma)\). Torsion in
\(H_1(\Sigma,\mathbb Z)\) can produce discrete Wilson lines, but
discrete choices do not give continuous four-dimensional scalar fields
and are not counted by \(b_1\).

For a \(U(N_a)\) stack with a background flat bundle \(E\), the correct
infinitesimal space is \[
H^1(\Sigma,\operatorname{ad}E).
\] For the trivial flat connection used implicitly in the quoted
counting, \[
H^1(\Sigma,\operatorname{ad}E)
\cong
H^1(\Sigma,\mathbb R)\otimes\mathfrak u(N_a).
\] This makes the adjoint Chan--Paton representation explicit.

\subsubsection{Pairing into complex adjoint
scalars}\label{pairing-into-complex-adjoint-scalars}

Choose a harmonic basis \(\eta_I\), \(I=1,\ldots,b_1(\Sigma)\). Through
the special-Lagrangian isomorphism, a normal deformation can be written
as \[
\alpha_v(x,y)
=\sum_{I=1}^{b_1(\Sigma)}s^I(x)\eta_I(y),
\] while the Wilson-line fluctuation is \[
a(x,y)
=\sum_{I=1}^{b_1(\Sigma)}\xi^I(x)\eta_I(y).
\] Thus for each \(I\) there are two real four-dimensional scalar
fields: \[
(s^I,\xi^I).
\] They are paired by the supersymmetry preserved by the special
Lagrangian D6-brane. After a convention-dependent normalization, one can
form complex coordinates \[
\Phi^I=\xi^I+i\,c_I s^I,
\] where the factors \(c_I\) depend on the normalization of the harmonic
forms, the D6-brane action, and the chosen four-dimensional fields. The
invariant statement is the counting: \[
\underbrace{b_1(\Sigma)}_{\text{real sLag deformations}}
+
\underbrace{b_1(\Sigma)}_{\text{real Wilson lines}}
=
\underbrace{b_1(\Sigma)}_{\text{complex scalars}}.
\]

Because both types of open-string zero modes carry the same endpoint
Chan--Paton matrices, every \(\Phi^I\) transforms in
\(\operatorname{Adj}U(N_a)\).

\subsubsection{Checks and
qualifications}\label{checks-and-qualifications}

For a factorized toroidal cycle \(\Sigma=T^3\), \[
b_1(T^3)=3.
\] There are three independent harmonic one-forms \(dx^1,dx^2,dx^3\),
hence three real Wilson-line moduli and three real special-Lagrangian
position/deformation moduli. They combine into three complex adjoint
scalars, in agreement with the familiar toroidal D6-brane spectrum.

For \(\Sigma=S^3\), \[
b_1(S^3)=0.
\] The special Lagrangian is locally rigid and there are no continuous
Wilson lines, so there are no adjoint chiral multiplets of this type.
Notice that \(\Gamma(N_{S^3})\) is nevertheless infinite-dimensional.
This is a simple demonstration that the book's \(H^0(\Sigma,N_\Sigma)\)
cannot, by itself, be the literal finite multiplicity of massless
fields.

The counting above assumes a compact smooth special Lagrangian without
boundary, a fixed ambient Calabi--Yau structure, no worldvolume flux,
and the trivial flat Chan--Paton bundle. A nontrivial flat bundle
replaces ordinary cohomology by \(H^1(\Sigma,\operatorname{ad}E)\).
Orientifold projections can retain only particular parity eigenspaces,
and fluxes or other effects can lift some modes. For a generic
non-invariant brane and its orientifold image giving \(U(N_a)\), the
book is using the unprojected \(b_1(\Sigma)\) count described above.

\subsubsection{Result}\label{result-34}

The hierarchy of spaces is \[
\boxed{
\begin{aligned}
\text{all infinitesimal embedded deformations:}
&\quad
v\in H^0(\Sigma,N_\Sigma)=\Gamma(N_\Sigma),\\[2mm]
\alpha_v:
&\quad
v\longmapsto-i^*(\iota_v\omega)
\in\Omega^1(\Sigma),\\[2mm]
\text{linearized sLag conditions:}
&\quad
d\alpha_v=0,
\qquad
d^\dagger\alpha_v=0,\\[2mm]
\text{massless geometric moduli:}
&\quad
T_{[\Sigma]}\mathcal M_{\mathrm{sLag}}
\cong\mathcal H^1(\Sigma)
\cong H^1(\Sigma,\mathbb R),\\[2mm]
\text{continuous Wilson lines:}
&\quad
\frac{\{a\in\Omega^1:da=0\}}
{\{a\sim a+d\lambda\}}
\cong H^1(\Sigma,\mathbb R).
\end{aligned}
}
\]

Therefore, \[
\boxed{
b_1(\Sigma)\text{ real geometric modes}
+b_1(\Sigma)\text{ real Wilson-line modes}
=b_1(\Sigma)\text{ complex adjoint scalars}
}.
\]

\section{Q2: No-scale structure and the tree-level vacuum
energy}\label{q2-no-scale-structure-and-the-tree-level-vacuum-energy}

In page661, they derived the holomorphic superpotential for type IIB on
Calabi-Yau threefold with non-trivial three-form flux \[
W=\frac{1}{\tilde{\kappa}_{10}^2} \int_{X}\Omega \wedge G_{3}
\] The holomorphic three-form depends on the complex structure moduli,
\(G_{3}\) on the complex scalar but \(W\) is independent of the Kähler
moduli. The Kähler potential, which is also needed to compute the scalar
potential, depends on the complex-structure moduli, on the Kähler-moduli
and on \(\tau\). The dependence on the Kähler moduli features the
so-called no-scale-structure, which means that supersymmetric minima
have vanishing vacuum energy.

I am interested with this no-scale-structure. What is its concrete
statement, and to what extend it leads to the vanishing of vacuum
energy?

\begin{quote}
\textbf{Textbook location and related material.} Chapter 17, Section
17.3 around pp.~657--663, especially p.~661: the Gukov--Vafa--Witten
superpotential, Kähler potential, F-term potential, and no-scale
structure.
\end{quote}

\subsection{Answer}\label{answer-35}

\subsubsection{Core point}\label{core-point-35}

The concrete no-scale statement is an identity satisfied by the Kähler
potential of the Kähler moduli: \[
\boxed{
K^{\alpha\bar\beta}K_\alpha K_{\bar\beta}=3
}
\] in four-dimensional Planck units. Together with \[
\partial_{T_\alpha}W=0,
\] this implies \[
D_\alpha W=K_\alpha W
\] and hence \[
K^{\alpha\bar\beta}D_\alpha W
D_{\bar\beta}\overline W
=3|W|^2.
\] This term cancels the universal negative term \(-3|W|^2\) in the
\(\mathcal N=1\) F-term potential.

The cancellation does \textbf{not} imply that the scalar potential
vanishes identically. It removes only the Kähler-modulus part. At tree
level the result is \[
\boxed{
V_F
=e^K K^{A\bar B}D_AW
D_{\bar B}\overline W
\geq0,
}
\] where \(A,B\) run over the complex-structure moduli and the
axiodilaton. Therefore:

\begin{itemize}
\tightlist
\item
  \(V_F\) is zero when \(D_AW=0\) for all complex-structure and
  axiodilaton directions;
\item
  this zero-energy point need not be supersymmetric, because
  \(D_\alpha W=K_\alpha W\) is nonzero when \(W\neq0\);
\item
  a fully supersymmetric vacuum must also obey \(D_\alpha W=0\), which
  in a no-scale model forces \(W=0\) and therefore gives a
  supersymmetric Minkowski vacuum;
\item
  the Kähler moduli are not stabilized at this order;
\item
  corrections that change the Kähler potential or introduce \(T_\alpha\)
  dependence in \(W\) generally break the cancellation.
\end{itemize}

Thus there are two distinct conclusions: \[
\boxed{
\begin{aligned}
D_AW=0,\quad W\neq0
&\quad\Longrightarrow\quad
V_F=0,\quad D_\alpha W\neq0,
\quad\text{nonsupersymmetric no-scale Minkowski},\\[1mm]
D_AW=0,\quad W=0
&\quad\Longrightarrow\quad
V_F=0,\quad D_IW=0,
\quad\text{supersymmetric Minkowski}.
\end{aligned}
}
\]

The sentence in the book that ``supersymmetric minima have vanishing
vacuum energy'' refers to the second line. The same no-scale
cancellation actually allows the larger first class of zero-energy but
nonsupersymmetric minima, as the book explains immediately below
equation (17.58).

\subsubsection{\texorpdfstring{The generic \(\mathcal N=1\) F-term
potential}{The generic \textbackslash mathcal N=1 F-term potential}}\label{the-generic-mathcal-n1-f-term-potential}

For four-dimensional \(\mathcal N=1\) supergravity, the F-term scalar
potential is \[
V_F
=e^K\left(
K^{I\bar J}D_IW D_{\bar J}\overline W
-3|W|^2
\right),
\] where \[
D_IW=\partial_IW+K_IW.
\] The index \(I\) runs over every chiral multiplet. We use
\(M_{\mathrm P}=1\) here. The book retains \(\kappa_4\) and writes \[
V_F
=e^{\kappa_4^2K}\left(
K^{I\bar J}D_IW D_{\bar J}\overline W
-3\kappa_4^2|W|^2
\right),
\qquad
D_IW=\partial_IW+\kappa_4^2K_IW.
\]

The first term is nonnegative because the Kähler metric is positive
definite. The second term is negative. Consequently, in a generic
\(\mathcal N=1\) theory, a supersymmetric point \[
D_IW=0\quad\text{for every }I
\] has \[
V_{\mathrm{susy}}=-3e^K|W|^2.
\] It is therefore: \[
\begin{array}{c|c}
W=0 & \text{supersymmetric Minkowski},\\
W\neq0 & \text{supersymmetric anti-de Sitter}.
\end{array}
\]

The special feature of the flux compactification is that the Kähler
moduli occur in a sector whose positive F-term contribution is forced to
reproduce exactly \(3|W|^2\). That is the no-scale structure.

\subsubsection{Tree-level Kähler potential in type
IIB}\label{tree-level-kuxe4hler-potential-in-type-iib}

For the Calabi--Yau orientifold compactification considered around
equation (17.58), the tree-level Kähler potential has the schematic
separated form \[
K
=K_{\mathrm{cs}}(z,\bar z)
+K_\tau(\tau,\bar\tau)
+K_{\mathrm K}(T,\bar T),
\] with \[
K_{\mathrm{cs}}
=-\log\left(
i\int_X\Omega(z)\wedge\overline{\Omega(z)}
\right),
\] \[
K_\tau
=-\log[-i(\tau-\bar\tau)],
\] and \[
\boxed{
K_{\mathrm K}=-2\log\mathcal V(T+\bar T).
}
\] Additive constants and convention-dependent powers of
\(\tilde\kappa_{10}\) do not affect the no-scale identity.

The flux superpotential is \[
W(z,\tau)
=\frac{1}{\tilde\kappa_{10}^2}
\int_X\Omega(z)\wedge
\bigl(F_3-\tau H_3\bigr).
\] It depends on \(z^i\) through \(\Omega\) and on \(\tau\) through
\(G_3\), but at string tree level \[
\boxed{\partial_{T_\alpha}W=0.}
\]

The absence of \(T_\alpha\) from \(W\) is only one half of the no-scale
mechanism. The other half is the precise homogeneity of \(\mathcal V\).

\subsubsection{Deriving the no-scale identity from volume
homogeneity}\label{deriving-the-no-scale-identity-from-volume-homogeneity}

Let \[
x^\alpha=T_\alpha+\overline T_\alpha.
\] The real parts of the \(T_\alpha\) are four-cycle volumes: \[
\tau_\alpha
=\frac{\partial\mathcal V}{\partial v^\alpha}
=\frac12\kappa_{\alpha\beta\gamma}
v^\beta v^\gamma,
\] while \[
\mathcal V
=\frac16\kappa_{\alpha\beta\gamma}
v^\alpha v^\beta v^\gamma.
\] Under a rescaling \(v^\alpha\mapsto\lambda v^\alpha\), \[
\mathcal V\mapsto\lambda^3\mathcal V,
\qquad
\tau_\alpha\mapsto\lambda^2\tau_\alpha.
\] Therefore, when regarded as a function of the four-cycle volumes,
\(\mathcal V\) is homogeneous of degree \(3/2\): \[
\mathcal V(\lambda x)
=\lambda^{3/2}\mathcal V(x).
\] Euler's theorem for homogeneous functions gives \[
x^\alpha\partial_{x^\alpha}\mathcal V
=\frac32\mathcal V.
\]

Because \[
K_{\mathrm K}=-2\log\mathcal V,
\] we obtain \[
x^\alpha K_\alpha
=-2\,\frac{x^\alpha\partial_{x^\alpha}\mathcal V}
{\mathcal V}
=-3.
\] Here \(\partial/\partial T_\alpha\) acts as
\(\partial/\partial x^\alpha\) on a function that depends only on
\(T+\bar T\).

Differentiate the last identity with respect to \(x^\beta\): \[
K_{\bar\beta}
+x^\alpha K_{\alpha\bar\beta}
=0.
\] Multiplying by the inverse Kähler metric yields \[
K^{\gamma\bar\beta}K_{\bar\beta}
=-x^\gamma.
\] Contracting once more with \(K_\gamma\) gives \[
K^{\gamma\bar\beta}K_\gamma K_{\bar\beta}
=-x^\gamma K_\gamma
=3.
\] Therefore \[
\boxed{
K^{\alpha\bar\beta}K_\alpha K_{\bar\beta}=3.
}
\]

This is not an accidental one-modulus relation. It follows from the
degree-\(3/2\) homogeneity of the Calabi--Yau volume as a function of
divisor volumes and therefore holds for any number of Kähler moduli at
the classical level.

In the dimensionful convention of equation (17.56), the corresponding
formulas are \[
K^{\alpha\bar\beta}K_\alpha K_{\bar\beta}
=\frac{3}{\kappa_4^2},
\qquad
D_\alpha W=\kappa_4^2K_\alpha W.
\] Consequently, \[
K^{\alpha\bar\beta}D_\alpha W
D_{\bar\beta}\overline W
=3\kappa_4^2|W|^2,
\] which cancels precisely the \(-3\kappa_4^2|W|^2\) term printed in
equation (17.56).

\subsubsection{Explicit one-modulus
calculation}\label{explicit-one-modulus-calculation}

The example in footnote 21 has one Kähler modulus and \[
K
=-3\log(T+\overline T)
+\widetilde K(\phi_A,\overline\phi_{\bar A}),
\qquad
W=W(\phi_A).
\] Then \[
K_T=-\frac{3}{T+\overline T},
\qquad
K_{T\overline T}
=\frac{3}{(T+\overline T)^2},
\qquad
K^{T\overline T}
=\frac{(T+\overline T)^2}{3}.
\] Hence \[
K^{T\overline T}K_TK_{\overline T}
=\frac{(T+\overline T)^2}{3}
\frac{9}{(T+\overline T)^2}
=3.
\]

Since \(\partial_TW=0\), \[
D_TW=K_TW.
\] It follows that \[
K^{T\overline T}|D_TW|^2
=K^{T\overline T}K_TK_{\overline T}|W|^2
=3|W|^2.
\] Substitution into the full potential gives \[
\begin{aligned}
V_F
&=e^K\left[
K^{A\bar B}D_AW D_{\bar B}\overline W
+K^{T\overline T}|D_TW|^2
-3|W|^2
\right]\\
&=e^K K^{A\bar B}
D_AW D_{\bar B}\overline W.
\end{aligned}
\] This is the formula displayed in compressed form in footnote 21.

\subsubsection{What remains after the
cancellation}\label{what-remains-after-the-cancellation}

For several Kähler moduli, the same calculation gives \[
\begin{aligned}
V_F
&=e^K\bigl[
K^{A\bar B}D_AW D_{\bar B}\overline W
+K^{\alpha\bar\beta}D_\alpha W
D_{\bar\beta}\overline W
-3|W|^2
\bigr]\\
&=e^K K^{A\bar B}D_AW
D_{\bar B}\overline W.
\end{aligned}
\] Thus \(V_F\) is positive semidefinite, not identically zero: \[
V_F\geq0.
\] It vanishes precisely when \[
D_AW=0
\] for all complex-structure and axiodilaton fields, assuming the
corresponding Kähler metric is nonsingular and positive definite.

At such a point the fluxes have fixed \(z^i\) and \(\tau\), but the
potential has no dependence that fixes \(T_\alpha\). After inserting the
solution for \(z^i\) and \(\tau\), one obtains a constant \[
W_0=W(z_\star,\tau_\star),
\] while \[
V_F(T,\bar T)=0
\] along the entire Kähler-moduli space at this order. This degenerate
family of flat directions is the origin of the name ``no-scale.''

\subsubsection{\texorpdfstring{Relation to ISD flux and the Hodge
decomposition of
\(G_3\)}{Relation to ISD flux and the Hodge decomposition of G\_3}}\label{relation-to-isd-flux-and-the-hodge-decomposition-of-g_3}

The four-dimensional result is the effective-supergravity version of the
ten-dimensional positive-semidefinite potential in equations
(17.53)--(17.55). For harmonic three-form flux on a Calabi--Yau
threefold, the imaginary self-dual and imaginary anti-self-dual
components are \[
*_6G_3^{\mathrm{ISD}}
=iG_3^{\mathrm{ISD}},
\qquad
*_6G_3^{\mathrm{IASD}}
=-iG_3^{\mathrm{IASD}}.
\] The flux potential, after the tadpole and source contribution has
been included, is proportional to the norm of the IASD component: \[
V_{\mathrm{flux}}
\sim
\frac{1}{\operatorname{Im}\tau\,\mathcal V^2}
\int_X
G_3^{\mathrm{IASD}}
\wedge *_6\overline{G_3^{\mathrm{IASD}}}
\geq0.
\] Therefore its zero-energy minima obey \[
G_3^{\mathrm{IASD}}=0,
\] so \(G_3\) is ISD.

In the Calabi--Yau Hodge decomposition, and under the
harmonic/primitivity assumptions used in this discussion, \[
\begin{array}{c|c}
\text{ISD} & G_3^{(2,1)}+G_3^{(0,3)},\\
\text{IASD} & G_3^{(1,2)}+G_3^{(3,0)}.
\end{array}
\] The complex-structure equations \(D_iW=0\) remove the \((1,2)\) part,
while the axiodilaton equation \(D_\tau W=0\) removes the \((3,0)\)
part. These equations therefore impose the ISD condition but can leave a
\((0,3)\) component.

The value of the superpotential detects precisely that remaining
component: \[
W
=\frac{1}{\tilde\kappa_{10}^2}
\int_X\Omega^{(3,0)}
\wedge G_3^{(0,3)}.
\] Thus:

\begin{itemize}
\tightlist
\item
  primitive \((2,1)\) flux has \(W=0\) and preserves \(\mathcal N=1\)
  supersymmetry;
\item
  a \((0,3)\) component gives \(W\neq0\), is still ISD, and can have
  zero no-scale potential, but it breaks supersymmetry.
\end{itemize}

This is why the condition ``ISD'' is sufficient for a tree-level
zero-energy minimum but is not, by itself, sufficient for a
supersymmetric minimum.

\subsubsection{\texorpdfstring{Why a fully supersymmetric no-scale
vacuum forces
\(W=0\)}{Why a fully supersymmetric no-scale vacuum forces W=0}}\label{why-a-fully-supersymmetric-no-scale-vacuum-forces-w0}

Unbroken four-dimensional \(\mathcal N=1\) supersymmetry requires \[
D_IW=0
\] for \textbf{all} chiral multiplets, including the Kähler moduli.
Since \(W\) is independent of them, \[
D_\alpha W=K_\alpha W.
\] The no-scale Euler identity gives \[
x^\alpha K_\alpha=-3.
\] If every \(D_\alpha W\) vanishes, then \[
0
=x^\alpha D_\alpha W
=x^\alpha K_\alpha W
=-3W.
\] Therefore \[
\boxed{
\text{full supersymmetry}\quad\Longrightarrow\quad W=0.
}
\] At such a point, \[
D_IW=0,
\qquad
W=0,
\] and the generic supergravity expression gives \[
\boxed{V_{\mathrm{susy}}=0.}
\] In flux language, \(W=0\) removes \(G_3^{(0,3)}\), and the
supersymmetric flux is primitive of type \((2,1)\), exactly as stated
below equation (17.58).

This explains the book's statement more precisely. A generic
supersymmetric \(\mathcal N=1\) vacuum may be AdS when \(W\neq0\), but
in this tree-level no-scale compactification the Kähler-modulus F-term
equations themselves force \(W=0\). Hence the supersymmetric flux vacua
in this approximation are Minkowski.

\subsubsection{Zero vacuum energy does not imply
supersymmetry}\label{zero-vacuum-energy-does-not-imply-supersymmetry}

Suppose instead that \[
D_AW=0,
\qquad
W_0\neq0.
\] The no-scale potential still vanishes: \[
V_F=0.
\] However, \[
D_\alpha W=K_\alpha W_0\neq0.
\] Equivalently, the Kähler-modulus auxiliary fields are nonzero: \[
F^\alpha
=-e^{K/2}K^{\alpha\bar\beta}
D_{\bar\beta}\overline W
\neq0.
\] The gravitino mass is also nonzero, \[
m_{3/2}=e^{K/2}|W_0|.
\] Supersymmetry is therefore spontaneously broken even though the
classical vacuum energy vanishes.

The equality \(V_F=0\) is a cancellation between the positive
supersymmetry-breaking contribution \[
K^{\alpha\bar\beta}D_\alpha W
D_{\bar\beta}\overline W
=3|W_0|^2
\] and the negative supergravity term \(-3|W_0|^2\). It is not a
statement that every auxiliary field vanishes.

This distinction is particularly important: \[
\boxed{
V=0\ \not\Longrightarrow\ \text{unbroken supersymmetry}.
}
\] The ISD flux \(G_3^{(0,3)}\) provides the standard string-theory
realization of this nonsupersymmetric no-scale Minkowski situation.

\subsubsection{Scope and limitations of the vanishing-energy
result}\label{scope-and-limitations-of-the-vanishing-energy-result}

The cancellation is powerful but limited. The tree-level conclusion
requires all of the following.

\paragraph{The remaining F-terms must
vanish}\label{the-remaining-f-terms-must-vanish}

No-scale structure cancels only the Kähler-modulus block and the
universal negative term. If \[
D_iW\neq0
\quad\text{or}\quad
D_\tau W\neq0,
\] then \[
V_F=e^K K^{A\bar B}D_AW
D_{\bar B}\overline W>0.
\] Thus no-scale structure does not make an arbitrary flux configuration
a Minkowski vacuum.

\paragraph{D-terms must vanish as
well}\label{d-terms-must-vanish-as-well}

The full scalar potential is \[
V=V_F+V_D.
\] Gauge fluxes, charged fields, or D7-branes can generate a
positive-semidefinite D-term potential. A zero-energy vacuum also
requires \[
V_D=0.
\] The F-term no-scale identity does not cancel a nonzero D-term.

\paragraph{The Kähler moduli are not
stabilized}\label{the-kuxe4hler-moduli-are-not-stabilized}

At the no-scale minimum, \[
V_F=0
\] for a continuous family of \(T_\alpha\). The tree-level flux
potential fixes the complex structure and axiodilaton but leaves the
overall volume and other Kähler moduli massless. A vacuum with exactly
flat moduli is not yet a fully stabilized phenomenological vacuum.

\paragraph{\texorpdfstring{Corrections break the classical identity or
the \(T\)-independence of
\(W\)}{Corrections break the classical identity or the T-independence of W}}\label{corrections-break-the-classical-identity-or-the-t-independence-of-w}

The relation \[
K^{\alpha\bar\beta}K_\alpha K_{\bar\beta}=3
\] uses the classical homogeneous Kähler potential
\(K_{\mathrm K}=-2\log\mathcal V\). Higher-derivative \(\alpha'\)
corrections and string-loop corrections change \(K\) and generally
change the left-hand side away from \(3\).

Nonperturbative effects such as Euclidean D3-brane instantons or gaugino
condensation can generate \[
W
=W_0+\sum_\alpha A_\alpha e^{-a_\alpha T_\alpha}.
\] Then \[
\partial_{T_\alpha}W\neq0,
\] so \(D_\alpha W\) is no longer simply \(K_\alpha W\) and the
cancellation no longer follows. These effects are precisely what are
used in constructions such as KKLT or the Large Volume Scenario to
stabilize Kähler moduli. The resulting vacuum energy can be negative,
zero, or positive depending on the complete stabilization and uplift
mechanism; no-scale structure alone no longer fixes it.

\paragraph{Warping and backreaction require the appropriate effective
theory}\label{warping-and-backreaction-require-the-appropriate-effective-theory}

The book states equation (17.58) while initially neglecting backreaction
and then explains that the consistent solution is a warped Calabi--Yau
geometry. The leading Giddings--Kachru--Polchinski effective theory
retains the relevant no-scale structure, but strong warping, localized
sources, and corrections must be incorporated consistently. The simple
unwarped formulas should therefore be understood as the leading
two-derivative approximation.

Consequently, the no-scale cancellation is not a solution of the
cosmological-constant problem. It is a classical structural cancellation
in a controlled approximation. Once the effects required to stabilize
the Kähler moduli are included, the vacuum energy must be computed
again.

\subsubsection{Result}\label{result-35}

At type-IIB tree level, \[
\begin{aligned}
K
&=K_{\mathrm{cs}}(z,\bar z)
+K_\tau(\tau,\bar\tau)
-2\log\mathcal V(T+\bar T),\\
W
&=\frac{1}{\tilde\kappa_{10}^2}
\int_X\Omega(z)\wedge(F_3-\tau H_3),
\qquad
\partial_{T_\alpha}W=0.
\end{aligned}
\] Homogeneity of \(\mathcal V\) gives \[
\boxed{
K^{\alpha\bar\beta}K_\alpha K_{\bar\beta}=3,
\qquad
D_\alpha W=K_\alpha W.
}
\] Therefore \[
\boxed{
\begin{aligned}
V_F
&=e^K\left[
K^{A\bar B}D_AW D_{\bar B}\overline W
+K^{\alpha\bar\beta}D_\alpha W
D_{\bar\beta}\overline W
-3|W|^2
\right]\\
&=e^K K^{A\bar B}D_AW
D_{\bar B}\overline W
\geq0.
\end{aligned}
}
\] The vacuum classes are \[
\boxed{
\begin{array}{c|c|c|c}
\text{conditions}
&V_F
&\text{supersymmetry}
&\text{Kähler moduli}\\ \hline
D_AW=0,\ W=0
&0
&\text{unbroken}
&\text{unfixed}\\
D_AW=0,\ W\neq0
&0
&\text{broken by }D_\alpha W
&\text{unfixed}\\
D_AW\neq0
&>0
&\text{broken}
&\text{generically still unstabilized at tree level}
\end{array}
}
\] Thus no-scale structure ensures that fully supersymmetric flux vacua
are Minkowski and also permits nonsupersymmetric tree-level Minkowski
vacua. It does not make all vacua have zero energy, does not stabilize
the Kähler moduli, and does not protect the zero energy once the
no-scale structure is broken by corrections or additional effects.

\chapter{String Dualities and
M-Theory}\label{string-dualities-and-m-theory-1}

\section{Q1: String, Einstein, and modified Einstein
frames}\label{q1-string-einstein-and-modified-einstein-frames}

In page690 eq(18.26), they used the so-called modified Einstein frame:
\textgreater The frame we have used in (18.26) is called modified
Einstein frame. It is related to the Einstein frame through a rescaling
of the metric by the constant (asymptotic) value of the dilaton.

There are at least three frames used in this book: string frame,
Einstein frame and this modified Einstein frame. List the low-energy
effective action in each frame and show how they are related, and the
special properties of each frame, and in which discussion which frame is
used.

\begin{quote}
\textbf{Textbook location and related material.} Chapter 18, Sections
18.3 and 18.5 (pp.~680--684, 688--705), especially Eq. (18.26): string,
Einstein, and asymptotically normalized modified Einstein frames for
brane solutions.
\end{quote}

\subsection{Answer}\label{answer-36}

\subsubsection{Core point}\label{core-point-36}

The three frames do not describe three different theories. They are
three choices of field variables, distinguished by which metric is
called the spacetime metric. Their exact relation in ten dimensions is
\[
\boxed{
G^{(E)}_{MN}=e^{-\Phi/2}G^{(S)}_{MN},
\qquad
g^{(mE)}_{MN}=e^{-\phi/2}G^{(S)}_{MN}
=\sqrt{g_s}\,G^{(E)}_{MN}
},
\] where \[
\Phi=\Phi_0+\phi,
\qquad
g_s=e^{\Phi_0},
\qquad
\phi\xrightarrow{r\to\infty}0.
\] Equivalently, \[
G^{(S)}_{MN}=e^{\Phi/2}G^{(E)}_{MN}
=e^{\phi/2}g^{(mE)}_{MN},
\qquad
G^{(E)}_{MN}=\frac1{\sqrt{g_s}}g^{(mE)}_{MN}.
\]

Thus the ordinary and modified Einstein metrics differ only by a
constant. They have the same local advantages: in both frames the Ricci
scalar has no dilaton prefactor and the graviton and dilaton kinetic
terms are diagonal. The reason for introducing the modified frame is
instead the boundary condition. If the string-frame metric approaches
the coordinate Minkowski metric, then \[
G^{(S)}_{MN}\to\eta_{MN},
\qquad
g^{(mE)}_{MN}\to\eta_{MN},
\qquad
G^{(E)}_{MN}\to\frac{\eta_{MN}}{\sqrt{g_s}}.
\] Consequently, the modified frame is convenient for computing the ADM
mass or tension of the brane solutions without performing a final
constant conversion.

The rest of the answer first writes the action in all three frames, then
derives the Weyl transformation explicitly, and finally explains where
each frame is used in the book.

\subsubsection{String-frame action}\label{string-frame-action}

For the bosonic type-II fields, omitting the Chern--Simons terms for the
moment, the string-frame action can be written schematically as \[
S_S
=\frac1{2\widetilde\kappa_{10}^{\,2}}
\int d^{10}x\,\sqrt{-G^{(S)}}
\left[
e^{-2\Phi}
\left(
R_S+4(\partial\Phi)^2_S-\frac1{2\cdot3!}H_{MNP}H^{MNP}_S
\right)
-\sum_{n\in\mathrm{RR}}
\frac1{2n!}F_n^2{}_S
\right]
+S_{\mathrm{CS}}.
\] For type IIA the non-democratic R--R degrees are \(F_2,F_4\), while
for type IIB they are \(F_1,F_3,F_5\). In the IIB pseudo-action the
\(F_5\) term has one extra factor of \(1/2\), \[
-\frac14|F_5|^2,
\] and the self-duality condition \(F_5=*F_5\) is imposed after varying
the pseudo-action.

This is the structure displayed explicitly in equation (16.134). The
decisive feature is that every tree-level NS--NS term has the common
factor \(e^{-2\Phi}\), whereas the R--R kinetic terms do not. This
difference is not arbitrary. A closed-string sphere has Euler
characteristic \(\chi=2\) and hence contributes \(g_s^{-2}\); the NS--NS
bulk action is therefore naturally organized with \(e^{-2\Phi}\). The
R--R fields used in (16.134) have been normalized so that their kinetic
terms have no such exponential.

It is useful to rewrite the same action in the Chapter 18 normalization.
Split off the asymptotic dilaton and define \[
\kappa_{10}=g_s\widetilde\kappa_{10},
\qquad
\widetilde F_{\mathrm{RR}}=g_sF_{\mathrm{RR}},
\qquad
\widetilde H_3=H_3.
\] Then \[
S_S
=\frac1{2\kappa_{10}^{\,2}}
\int d^{10}x\,\sqrt{-G^{(S)}}
\left[
e^{-2\phi}
\left(
R_S+4(\partial\phi)^2_S-\frac1{2\cdot3!}H_3^2{}_S
\right)
-\sum_{n\in\mathrm{RR}}
\frac1{2n!}\widetilde F_n^2{}_S
\right]
+S_{\mathrm{CS}}.
\] Indeed, \[
\frac{e^{-2\Phi}}{\widetilde\kappa_{10}^{\,2}}
=\frac{e^{-2\phi}}{\kappa_{10}^{\,2}},
\qquad
\frac{F_n^2}{\widetilde\kappa_{10}^{\,2}}
=\frac{\widetilde F_n^2}{\kappa_{10}^{\,2}}.
\] This second form makes the passage to equation (18.26) transparent.

The special properties of string frame are:

\begin{itemize}
\tightlist
\item
  \(G^{(S)}_{MN}\) is the metric that appears directly in the Polyakov
  sigma-model action, \[
  S_{\sigma}= -\frac1{4\pi\alpha'}
  \int d^2\sigma\sqrt{-h}\,h^{\alpha\beta}
  G^{(S)}_{MN}(X)\partial_\alpha X^M\partial_\beta X^N+\cdots.
  \]
\item
  The world-sheet beta functions directly give the string-frame
  equations of motion. This is why equation (14.8), derived from the
  vanishing of the sigma-model beta functions, is in string frame.
\item
  The factors \(e^{-2\Phi}\) for a sphere and \(e^{-\Phi}\) for a disk
  make string-loop counting manifest. In particular, the string-frame
  fundamental-string tension is independent of \(g_s\), a D-brane
  tension scales as \(g_s^{-1}\), and an NS5-brane tension scales as
  \(g_s^{-2}\).
\item
  The Buscher T-duality rules, equations (14.15)--(14.16), take their
  standard form for \(G^{(S)}\), \(B_2\), and \(\Phi\).
\item
  The disadvantage is that \(e^{-2\Phi}R_S\) produces kinetic mixing
  between the trace of the metric fluctuation and the dilaton. Thus the
  graviton and dilaton are not diagonal variables at quadratic order.
\end{itemize}

\subsubsection{Ordinary Einstein-frame
action}\label{ordinary-einstein-frame-action}

Define the ordinary Einstein metric by equation (16.141), \[
\boxed{G^{(E)}_{MN}=e^{-\Phi/2}G^{(S)}_{MN}}.
\] In this frame the type-II bulk action is \[
S_E
=\frac1{2\widetilde\kappa_{10}^{\,2}}
\int d^{10}x\,\sqrt{-G^{(E)}}
\left[
R_E
-\frac12(\partial\Phi)^2_E
-\frac{e^{-\Phi}}{2\cdot3!}H_3^2{}_E
-\sum_{n\in\mathrm{RR}}
\frac{e^{(5-n)\Phi/2}}{2n!}F_n^2{}_E
\right]
+S_{\mathrm{CS}}.
\] For \(n=p+2\), the R--R exponent is \[
\frac{5-n}{2}=\frac{3-p}{2}=a_p.
\] Thus, for example, \[
F_2:\ e^{3\Phi/2},
\qquad
F_3:\ e^{\Phi},
\qquad
F_4:\ e^{\Phi/2},
\qquad
F_5:\ e^0.
\] The NS--NS three-form instead has \(e^{-\Phi}\), corresponding to
\(a_1=-1\). The coincidence \(a_3=0\) for the R--R five-form explains
why the D3-brane can have a constant dilaton.

The defining properties of Einstein frame are now visible directly:

\begin{itemize}
\tightlist
\item
  The gravitational term is the canonical Einstein--Hilbert term
  \(\sqrt{-G^{(E)}}R_E\), with no dilaton-dependent multiplier.
\item
  The scalar kinetic term is canonically normalized as
  \(-\frac12(\partial\Phi)^2\).
\item
  At quadratic order around a constant-dilaton background, the graviton
  and dilaton propagators are diagonal. This is the point made after
  equation (16.115): the vertex operators used there create orthogonal
  graviton and dilaton states, so their low-energy action naturally
  first appears in Einstein frame.
\item
  The ordinary Einstein metric is the natural metric for spacetime
  duality symmetries. In type IIB, the action can be organized using \[
  \tau=C_0+ie^{-\Phi},
  \] and the classical \(SL(2,\mathbb R)\) transformation leaves
  \(G^{(E)}_{MN}\) invariant. This is the form used in equations
  (16.141)--(16.144) and in the S-duality discussion of Section 18.6.
\end{itemize}

The cost is that the direct relation to the world-sheet expansion is
less visible: the various form-field terms now carry different dilaton
exponentials. Also, if one keeps the same coordinates in which
\(G^{(S)}\to\eta\), the standard Einstein metric does not approach
\(\eta\) but rather \(\eta/\sqrt{g_s}\).

\subsubsection{Modified Einstein-frame action and equation
(18.26)}\label{modified-einstein-frame-action-and-equation-18.26}

The modified Einstein frame removes only the fluctuating part of the
dilaton from the string metric: \[
\boxed{
g^{(mE)}_{MN}=e^{-\phi/2}G^{(S)}_{MN}
}.
\] Since \(\Phi=\Phi_0+\phi\), comparison with the ordinary Einstein
metric gives \[
g^{(mE)}_{MN}
=e^{\Phi_0/2}G^{(E)}_{MN}
=\sqrt{g_s}\,G^{(E)}_{MN}.
\] This is exactly the constant rescaling mentioned in the quoted
paragraph.

With the physical gravitational coupling and rescaled form fields
defined above, the action becomes \[
S_{mE}
=\frac1{2\kappa_{10}^{\,2}}
\int d^{10}x\,\sqrt{-g^{(mE)}}
\left[
R_{mE}
-\frac12(\partial\phi)^2_{mE}
-\frac{e^{-\phi}}{2\cdot3!}H_3^2{}_{mE}
-\sum_{n\in\mathrm{RR}}
\frac{e^{(5-n)\phi/2}}{2n!}\widetilde F_n^2{}_{mE}
\right]
+S_{\mathrm{CS}}.
\] If only one \((p+2)\)-form is retained, this reduces precisely to
equation (18.26): \[
\boxed{
S
=\frac1{2\kappa_d^2}
\int d^dx\sqrt{-g}
\left(
R-\frac12g^{MN}\partial_M\phi\partial_N\phi
-\frac1{2(p+2)!}e^{a_p\phi}\widetilde F_{p+2}^{\,2}
\right)
},
\] with \[
a_1=-1\quad\text{for }H_3,
\qquad
a_p=\frac{3-p}{2}\quad\text{for an R--R }F_{p+2}.
\]

The special properties of the modified Einstein frame are:

\begin{itemize}
\tightlist
\item
  The local kinetic terms are just as canonical as in ordinary Einstein
  frame, because the two metrics differ only by a constant.
\item
  The fluctuating field obeys \(\phi\to0\), and the brane metrics are
  chosen to obey \(g^{(mE)}_{MN}\to\eta_{MN}\). ADM masses and tensions
  obtained from the asymptotic falloff can therefore be read in standard
  Minkowski units.
\item
  All constant \(g_s\) dependence has been moved into \(\kappa_{10}\),
  the physical brane tensions, and the rescaled R--R fields. The
  remaining exponentials depend only on the local fluctuation \(\phi\).
\item
  It is especially convenient when comparing supergravity brane
  solutions with the physical tensions previously extracted from string
  amplitudes.
\end{itemize}

\subsubsection{Explicit Weyl-rescaling
derivation}\label{explicit-weyl-rescaling-derivation}

The numerical exponents are worth deriving, because this is where the
distinction between \(\Phi\) and \(\phi\) matters.

For a Weyl rescaling in \(D\) dimensions, \[
G_{MN}=e^{2\omega}g_{MN},
\] one has \[
\sqrt{-G}=e^{D\omega}\sqrt{-g},
\] \[
R[G]=e^{-2\omega}
\left[
R[g]-2(D-1)\nabla^2\omega
-(D-1)(D-2)(\partial\omega)^2
\right],
\] and for an \(n\)-form, \[
F_n^2{}_G=e^{-2n\omega}F_n^2{}_g,
\] because raising each of its \(n\) indices contributes one inverse
metric.

For the string-to-modified-Einstein transformation in ten dimensions, \[
G^{(S)}_{MN}=e^{\phi/2}g^{(mE)}_{MN},
\qquad
\omega=\frac\phi4.
\] The overall exponential multiplying the Ricci scalar is \[
e^{10\omega}e^{-2\phi}e^{-2\omega}
=e^{8\omega-2\phi}
=e^{2\phi-2\phi}=1.
\] This proves that the Ricci scalar has no dilaton prefactor. More
explicitly, \[
\sqrt{-G^{(S)}}e^{-2\phi}R_S
=\sqrt{-g^{(mE)}}
\left[
R_{mE}-\frac92\nabla^2\phi
-\frac92(\partial\phi)^2
\right].
\] The \(\nabla^2\phi\) term is a total derivative and is discarded when
the appropriate boundary term is understood. The original string-frame
dilaton term becomes \[
\sqrt{-G^{(S)}}e^{-2\phi}\,4(\partial\phi)^2_S
=\sqrt{-g^{(mE)}}\,4(\partial\phi)^2_{mE}.
\] Combining it with the \(-\frac92(\partial\phi)^2\) generated by the
curvature gives \[
4-\frac92=-\frac12,
\] which is the canonical dilaton kinetic term in (18.26).

For the NS--NS three-form, the three factors are \[
\sqrt{-G^{(S)}}:\ e^{5\phi/2},
\qquad
e^{-2\phi},
\qquad
H_3^2{}_S:\ e^{-3\phi/2}H_3^2{}_{mE}.
\] Their product is \[
e^{5\phi/2}e^{-2\phi}e^{-3\phi/2}=e^{-\phi},
\] so \(a_1=-1\).

For an R--R \(n\)-form there is no initial \(e^{-2\phi}\) factor.
Therefore \[
\sqrt{-G^{(S)}}F_n^2{}_S
=\sqrt{-g^{(mE)}}
e^{5\phi/2}e^{-n\phi/2}F_n^2{}_{mE}
=\sqrt{-g^{(mE)}}
e^{(5-n)\phi/2}F_n^2{}_{mE}.
\] Putting \(n=p+2\) gives \[
a_p=\frac{5-(p+2)}2=\frac{3-p}{2},
\] exactly as stated below equation (18.26).

Finally, the constant rescaling from ordinary to modified Einstein frame
also reproduces every constant in (18.26). Starting with \[
G^{(E)}_{MN}=g_s^{-1/2}g^{(mE)}_{MN},
\qquad
\Phi=\Phi_0+\phi,
\] the gravitational term supplies \(g_s^{-2}\): \[
\sqrt{-G^{(E)}}R_E
=g_s^{-5/2}g_s^{1/2}\sqrt{-g^{(mE)}}R_{mE}
=g_s^{-2}\sqrt{-g^{(mE)}}R_{mE}.
\] Since \(\kappa_{10}^2=g_s^2\widetilde\kappa_{10}^{\,2}\), \[
\frac{g_s^{-2}}{\widetilde\kappa_{10}^{\,2}}
=\frac1{\kappa_{10}^{\,2}}.
\] For an R--R form the remaining factor is canceled by
\(\widetilde F_n=g_sF_n\). This explains why equation (18.26) contains
\(\kappa_{10}\) and \(\widetilde F\), rather than the Chapter 16
quantities \(\widetilde\kappa_{10}\) and \(F\).

\subsubsection{Brane source actions in the three
frames}\label{brane-source-actions-in-the-three-frames}

The same transformation is visible in the D\(p\)-brane action. In string
frame, \[
S_{Dp}^{(S)}
=-\mu_p\int_{\Sigma_{p+1}}e^{-\Phi}
\sqrt{-\det P[G^{(S)}]}
+\mu_p\int_{\Sigma_{p+1}}P[C_{p+1}]
+\cdots.
\] Using \(G^{(S)}=e^{\Phi/2}G^{(E)}\), the induced
\((p+1)\)-dimensional volume element transforms as \[
\sqrt{-\det P[G^{(S)}]}
=e^{(p+1)\Phi/4}
\sqrt{-\det P[G^{(E)}]}.
\] Hence \[
S_{Dp}^{(E)}
=-\mu_p\int e^{(p-3)\Phi/4}
\sqrt{-\det P[G^{(E)}]}
+\mu_p\int P[C_{p+1}]
+\cdots.
\] Since \(a_p=(3-p)/2\), this exponential is \(e^{-a_p\Phi/2}\).

In modified Einstein frame, define \[
\tau_p=\frac{\mu_p}{g_s},
\qquad
\widetilde C_{p+1}=g_sC_{p+1}.
\] Then \[
S_{Dp}^{(mE)}
=-\tau_p\int e^{-a_p\phi/2}
\sqrt{-\det P[g^{(mE)}]}
+\tau_p\int P[\widetilde C_{p+1}]
+\cdots,
\] which is exactly the source structure written in equation (18.48).
The common coefficient \(\tau_p\) of the DBI and Wess--Zumino terms
makes the BPS equality between tension and charge manifest in the
normalization of Section 18.5.

\subsubsection{Masses and tensions under the constant
rescaling}\label{masses-and-tensions-under-the-constant-rescaling}

The string and modified Einstein metrics have the same asymptotic
normalization, so an asymptotic observer finds \[
m_S^2=m_{mE}^2.
\] Because \[
G^{(E)}_{MN}=g_s^{-1/2}g^{(mE)}_{MN},
\] their inverse metrics obey \[
(G^{(E)})^{MN}=g_s^{1/2}(g^{(mE)})^{MN}.
\] Applying the mass-shell condition in the two metrics gives \[
\boxed{
m_S^2=m_{mE}^2=\frac{m_E^2}{\sqrt{g_s}}
},
\] which is the relation stated on page 690.

For a \(p\)-brane tension, which is mass per spatial \(p\)-volume, the
corresponding constant Weyl factor gives \[
\tau_p^{(E)}
=g_s^{(p+1)/4}\tau_p^{(mE)}.
\] For example, \[
\tau_{F1}^{(E)}=\sqrt{g_s}\,\tau_{F1}^{(S)},
\qquad
\tau_{D1}^{(E)}=\sqrt{g_s}\,\tau_{D1}^{(S)}
=\frac1{\sqrt{g_s}}\frac1{2\pi\alpha'},
\] so the fundamental string and D-string tensions are exchanged by
\(g_s\leftrightarrow1/g_s\) in ordinary Einstein frame, as in equation
(18.81).

\subsubsection{Which frame is used where in the
book}\label{which-frame-is-used-where-in-the-book}

The book changes frame according to the question being studied:

\begin{enumerate}
\def\labelenumi{\arabic{enumi}.}
\item
  \textbf{World-sheet sigma models, beta functions, and T-duality:
  string frame.} The nonlinear sigma-model metric in equation (14.1) is
  \(G^{(S)}\). The beta-function action (14.8) and the Buscher rules
  (14.15)--(14.16) are string-frame formulas.
\item
  \textbf{Ten-dimensional supergravity actions as obtained from string
  perturbation theory: mostly string frame.} Equations (16.118),
  (16.134), and (16.146) display the heterotic/type-II/type-I actions in
  string frame because their dilaton dependence keeps track of sphere
  and disk topology.
\item
  \textbf{Diagonal graviton--dilaton variables and type-IIB \(SL(2)\)
  symmetry: ordinary Einstein frame.} Equation (16.103) is identified as
  Einstein frame because the graviton and dilaton decouple at quadratic
  order. Equations (16.141)--(16.144) use ordinary Einstein frame to
  make the type-IIB \(SL(2,\mathbb R)\) symmetry manifest; this metric
  is invariant under the classical duality.
\item
  \textbf{Flux compactifications: frame stated case by case.} For
  example, the type-IIB flux action in equation (17.48) is in Einstein
  frame, whereas the type-II supersymmetry variations in equation
  (17.64) are explicitly in string frame. One must not combine their
  formulas without converting the metric, gamma matrices, and spinors
  consistently.
\item
  \textbf{The brane solutions of Section 18.5: modified Einstein frame
  unless the book explicitly labels a string-frame expression.} Equation
  (18.26), the equations of motion (18.27), the ansatz (18.28), the
  general solution (18.34), the ADM tension calculation, and the source
  action (18.48) all use \(g^{(mE)}\). The sentence ``from here on, in
  this section, Einstein frame always means modified Einstein frame''
  applies specifically to Section 18.5. The supersymmetry variations
  used in (18.52) are nevertheless imported in string frame, so the book
  explicitly converts the background when applying them. Table 18.3 then
  lists each ten-dimensional solution in both the modified Einstein
  frame and string frame.
\item
  \textbf{The S-duality tension comparison of Section 18.6: ordinary
  Einstein frame.} Equation (18.81) uses the standard Einstein-frame
  scaling \(\tau_{F1}^{(E)}\propto\sqrt{g_s}\) and
  \(\tau_{D1}^{(E)}\propto1/\sqrt{g_s}\). These would not be the
  scalings in the modified frame, whose asymptotic masses agree with
  string frame. Thus the special convention declared inside Section 18.5
  should not be silently extended to all later occurrences of the words
  ``Einstein frame.''
\item
  \textbf{DBI actions and perturbative brane tensions: naturally string
  frame, often converted to modified Einstein frame for supergravity
  matching.} The disk factor \(e^{-\Phi}\) and the scaling
  \(\tau_{Dp}\sim g_s^{-1}\) are most transparent in string frame.
  Equation (18.48) is their modified-Einstein-frame form, chosen to
  match the normalization of the supergravity charge and ADM tension.
\end{enumerate}

\subsubsection{Checks}\label{checks-30}

For a D3-brane, \(p=3\), so \[
a_3=0.
\] Its R--R five-form kinetic term and its modified-Einstein-frame DBI
volume term have no local dilaton factor. This is consistent with the
constant-dilaton D3 solution (18.59) and with the self-duality of the
D3-brane under type-IIB S-duality.

For the fundamental string, the relevant NS--NS field strength is
\(H_3\), so \(a_1=-1\). Its modified-Einstein-frame source factor is \[
e^{-a_1\phi/2}=e^{\phi/2},
\] which is exactly what follows by applying
\(G^{(S)}=e^{\phi/2}g^{(mE)}\) to a two-dimensional Nambu--Goto area
element.

For an R--R D1-brane, \(a_1=(3-1)/2=1\), so its source factor is
\(e^{-\phi/2}\). The opposite signs for F1 and D1 are the local dilaton
version of their exchange under type-IIB S-duality.

\subsubsection{Result}\label{result-36}

The complete chain of metric redefinitions is \[
\boxed{
\Phi=\Phi_0+\phi,
\quad g_s=e^{\Phi_0},
\quad
G^{(E)}=e^{-\Phi/2}G^{(S)},
\quad
g^{(mE)}=e^{-\phi/2}G^{(S)}
=\sqrt{g_s}\,G^{(E)}
}.
\]

The characteristic bulk actions are \[
\boxed{
S_S\sim\int\sqrt{-G^{(S)}}
\left[
e^{-2\Phi}\left(R_S+4(\partial\Phi)^2-\frac12|H_3|^2\right)
-\frac12\sum_{\mathrm{RR}}|F_n|^2
\right]
},
\] \[
\boxed{
S_E\sim\int\sqrt{-G^{(E)}}
\left[
R_E-\frac12(\partial\Phi)^2
-\frac12e^{-\Phi}|H_3|^2
-\frac12\sum_{\mathrm{RR}}e^{(5-n)\Phi/2}|F_n|^2
\right]
},
\] and \[
\boxed{
S_{mE}\sim\int\sqrt{-g^{(mE)}}
\left[
R_{mE}-\frac12(\partial\phi)^2
-\frac12e^{-\phi}|H_3|^2
-\frac12\sum_{\mathrm{RR}}e^{(5-n)\phi/2}|\widetilde F_n|^2
\right]
}.
\] Here the omitted common prefactors, factorials, Chern--Simons terms,
and the type-IIB five-form qualification are those stated above. For a
single field strength \(F_{p+2}\), the last action is equation (18.26),
with \[
\boxed{
a_p=-1\ \text{for NS--NS }H_3,
\qquad
a_p=\frac{3-p}{2}\ \text{for R--R }F_{p+2}
}.
\]

Conceptually: string frame is adapted to the world-sheet and
perturbative \(g_s\) counting; ordinary Einstein frame is adapted to
canonical gravity and S-duality; modified Einstein frame is adapted to
asymptotically Minkowski brane solutions and direct ADM mass/tension
comparisons.

\section{Q2: Deriving the BPS warp-factor
condition}\label{q2-deriving-the-bps-warp-factor-condition}

In page692 eq(18.33) is said to be derived from the requirement that the
p-brane solution preserves half of the supercharges. This equation is
very simple and clear, how to derived it?

\begin{quote}
\textbf{Textbook location and related material.} Chapter 18, Section
18.5 (pp.~688--705): the supergravity \(p\)-brane ansatz,
Killing-spinor/BPS conditions, harmonic functions, and the relation
among warp factors.
\end{quote}

\subsection{Answer}\label{answer-37}

\subsubsection{Core point}\label{core-point-37}

The statement to derive is \[
\boxed{nA+\tilde nB=\mathrm{const}},
\qquad
n=p+1,
\qquad
\tilde n=d-n-2.
\] For an asymptotically Minkowski solution, \(A,B\to0\) as
\(r\to\infty\), so the constant is zero.

This condition is a consequence of the first-order Killing-spinor
equations. A charged brane imposes one algebraic projector on the
supersymmetry parameter. That projector leaves one half of the spinors.
Once it is imposed, the longitudinal and transverse components of the
gravitino variation contain the same radial flux function. Their
relative gamma-matrix coefficients force \[
A'=-\frac{2\tilde n}{\Delta(d-2)}u(r),
\qquad
B'=\frac{2n}{\Delta(d-2)}u(r).
\] Therefore \[
nA'+\tilde nB'
=-\frac{2n\tilde n}{\Delta(d-2)}u
+\frac{2n\tilde n}{\Delta(d-2)}u
=0.
\] Integration gives equation (18.33).

The book does not perform this calculation at page 692. It first imposes
(18.33) while solving the bosonic equations, then later gives the
relevant Killing-spinor equations in (18.52)--(18.55). Footnote 9 on
page 697 says that one could instead insert the unsolved ansatz
(18.28)--(18.29) directly into those supersymmetry variations. The
derivation below supplies that omitted calculation.

\subsubsection{The metric ansatz and the meaning of the
coefficients}\label{the-metric-ansatz-and-the-meaning-of-the-coefficients}

The metric in equation (18.28) is \[
ds^2
=e^{2A(r)}\eta_{\mu\nu}dx^\mu dx^\nu
+e^{2B(r)}\delta_{mn}dy^m dy^n,
\] with \[
\mu=0,\ldots,n-1,
\qquad
m=1,\ldots,q,
\qquad
q=d-n=\tilde n+2.
\] Thus \(\tilde n\) is not the number of transverse dimensions; it is
two less: \[
\tilde n=q-2.
\] This explains why (18.33) is not simply a statement about the
determinant of the metric.

Choose the orthonormal coframe \[
e^{\hat\mu}=e^A dx^\mu,
\qquad
e^{\hat m}=e^Bdy^m.
\] The nonzero spin-connection components are precisely those displayed
in equation (18.53). Schematically, the parts relevant for a radial
Killing spinor are \[
\nabla_{\hat\mu}\epsilon
\ \propto\ e^{-B}A'\Gamma_{\hat\mu\hat r}\epsilon,
\qquad
\nabla_{\hat i}\epsilon
\ \propto\ e^{-B}B'\Gamma_{\hat i\hat r}\epsilon,
\] where \(\hat i\) is tangent to the transverse sphere and \[
\Gamma_{\hat r}=\frac{y^m}{r}\Gamma_{\hat m}.
\] The first term measures the radial variation of the longitudinal
vielbein, while the second measures the radial variation of the
transverse vielbein.

For an electric \(p\)-brane, the field strength has an orthonormal
component \[
F_{\hat r\hat0\cdots\hat p}(r).
\] Its gamma-matrix contraction is proportional to \[
F_{\hat r\hat0\cdots\hat p}
\Gamma^{\hat r\hat0\cdots\hat p}
=F_{\hat r\hat0\cdots\hat p}
\Gamma^{\hat r}\Gamma_\parallel,
\qquad
\Gamma_\parallel=\Gamma^{\hat0\cdots\hat p}.
\] Consequently the form-field term in the supersymmetry variation can
cancel the spin connection only on spinors obeying a brane projector.

\subsubsection{The brane projector and the loss of half the
supercharges}\label{the-brane-projector-and-the-loss-of-half-the-supercharges}

Equation (18.54) has the form \[
\left(1-s\,\mathcal K_p\right)\epsilon=0,
\qquad
\mathcal K_p=(-1)^p\Gamma^{\hat0\cdots\hat p}\mathcal P_{p+2},
\qquad
s=\pm1,
\] up to the book's choice of which sign is called a brane. The
gamma-matrix and type-II chirality factors obey \[
\mathcal K_p^2=1.
\] Hence \[
\Pi_s=\frac12(1+s\mathcal K_p)
\] is a projector. Its two eigenspaces have equal dimension, so \[
\Pi_s\epsilon=\epsilon
\] retains 16 of the original 32 type-II supercharges. Changing \(s\)
reverses the R--R charge and gives the anti-brane projector.

The projector is important for two separate reasons:

\begin{enumerate}
\def\labelenumi{\arabic{enumi}.}
\tightlist
\item
  it is the algebraic statement that the configuration is at most
  \(1/2\) BPS;
\item
  it replaces \(\Gamma_\parallel\mathcal P_{p+2}\) by the number \(s\)
  inside the differential Killing-spinor equations.
\end{enumerate}

After this replacement, all flux terms have the same gamma structures as
the spin-connection terms. The Killing-spinor equations therefore reduce
to ordinary first-order equations for the radial functions.

\subsubsection{General first-order BPS
equations}\label{general-first-order-bps-equations}

For those Einstein-dilaton-form systems in (18.26) that arise as
supersymmetric type-II or eleven-dimensional supergravity truncations,
the longitudinal gravitino equation, the angular transverse gravitino
equation, and the dilatino equation reduce to \[
A'
=-\frac{2\tilde n}{\Delta(d-2)}u(r),
\] \[
B'
=\frac{2n}{\Delta(d-2)}u(r),
\] \[
\phi'
=\frac{2\zeta a_p}{\Delta}u(r),
\qquad
\Delta=a_p^2+\frac{2n\tilde n}{d-2}.
\] Here \(\zeta=+1\) for the electric representative and \(\zeta=-1\)
for its magnetic dual, while the definition of \(u\) can absorb the
orientation sign \(s\). At this stage \(u(r)\) is simply the one common
function built from the orthonormal form-field component and the dilaton
exponential. The radial gravitino equation separately determines the
radial dependence of \(\epsilon\).

The crucial fact is the relative coefficient of the first two equations.
It comes from the gamma-matrix trace subtraction in the form-field
contribution to the gravitino variation: the flux acts with weight
\(\tilde n\) along a world-volume direction and with weight \(n\) along
a transverse direction, with opposite signs. The overall normalization
of \(u\), the value of \(a_p\), and the electric/magnetic sign do not
affect their ratio.

Taking the weighted sum immediately gives \[
\begin{aligned}
\frac{d}{dr}\left(nA+\tilde nB\right)
&=nA'+\tilde nB'\\
&=-\frac{2n\tilde n}{\Delta(d-2)}u
+\frac{2\tilde n n}{\Delta(d-2)}u\\
&=0.
\end{aligned}
\] Therefore \[
\boxed{nA+\tilde nB=C}.
\] The boundary conditions used below equation (18.28) are \[
A(r)\to0,
\qquad
B(r)\to0,
\qquad
\phi(r)\to0
\qquad
(r\to\infty),
\] so \[
\boxed{C=0}.
\]

Solving the remaining Maxwell and dilaton equations identifies \[
u(r)=(\log H(r))'
\] in a convenient normalization. Integration then gives \[
A=-\frac{2\tilde n}{\Delta(d-2)}\log H,
\qquad
B=\frac{2n}{\Delta(d-2)}\log H,
\qquad
\phi=\frac{2\zeta a_p}{\Delta}\log H,
\] which is equation (18.34). Notice that (18.33) follows before one
knows the explicit harmonic function \(H=1+N\alpha/r^{\tilde n}\).

\subsubsection{\texorpdfstring{Explicit derivation for a type-II
D\(p\)-brane}{Explicit derivation for a type-II Dp-brane}}\label{explicit-derivation-for-a-type-ii-dp-brane}

It is useful to see the cancellation without hiding it inside the
general coefficients. Equations (18.52) are written in string frame,
whereas \(A\) and \(B\) in (18.33) are modified-Einstein-frame warp
factors. Write \[
ds_S^2
=e^{2A_S(r)}\eta_{\mu\nu}dx^\mu dx^\nu
+e^{2B_S(r)}\delta_{mn}dy^m dy^n.
\] Because \[
G^{(S)}_{MN}=e^{\phi/2}g^{(mE)}_{MN},
\] the warp factors are related by \[
\boxed{
A_S=A+\frac{\phi}{4},
\qquad
B_S=B+\frac{\phi}{4}
}.
\]

Define the oriented orthonormal flux function \[
\mathcal Q(r)
=s\,e^{\Phi+B_S}
F_{\hat r\hat0\cdots\hat p}(r),
\] where numerical signs associated with \(\mathcal P_{p+2}\) are
absorbed into \(s\). Insert the spin connection (18.53) and the
projector (18.54) into the D\(p\)-brane equations (18.52). The four
independent components become \[
\delta\Psi_\mu=0:
\qquad
A_S'-\frac14\mathcal Q=0,
\] \[
\delta\Psi_i=0:
\qquad
B_S'+\frac14\mathcal Q=0,
\] \[
\delta\lambda=0:
\qquad
\phi'+\frac{3-p}{4}\mathcal Q=0,
\] and \[
\delta\Psi_r=0:
\qquad
\epsilon'-\frac18\mathcal Q\,\epsilon=0.
\] The overall sign of every \(\mathcal Q\) term flips if the opposite
orientation is chosen, but the relations obtained by eliminating
\(\mathcal Q\) are unchanged.

The first three equations imply \[
\boxed{
B_S'=-A_S',
\qquad
\phi'=(p-3)A_S'
}.
\] With the asymptotic conditions \(A_S,B_S,\phi\to0\), they integrate
to \[
B_S=-A_S,
\qquad
\phi=(p-3)A_S.
\] The radial gravitino equation gives \[
\epsilon(r)=e^{A_S(r)/2}\epsilon_0,
\qquad
\Pi_s\epsilon_0=\epsilon_0.
\] Thus the solution has a nonzero Killing spinor satisfying one
rank-half projector.

Now convert back to the modified Einstein frame: \[
\begin{aligned}
A
&=A_S-\frac{\phi}{4}
=\left(1-\frac{p-3}{4}\right)A_S
=\frac{7-p}{4}A_S,\\
B
&=B_S-\frac{\phi}{4}
=-\left(1+\frac{p-3}{4}\right)A_S
=-\frac{p+1}{4}A_S.
\end{aligned}
\] In ten dimensions, \[
n=p+1,
\qquad
\tilde n=10-(p+1)-2=7-p.
\] Therefore \[
\begin{aligned}
nA+\tilde nB
&=(p+1)\frac{7-p}{4}A_S
+(7-p)\left(-\frac{p+1}{4}A_S\right)\\
&=0.
\end{aligned}
\] This is equation (18.33) derived directly from the type-II
Killing-spinor equations quoted later in the same section.

For \(p=3\), the coefficient \(3-p\) in the dilatino equation vanishes.
Then \(\phi'=0\), as expected for a D3-brane. The brane projector is
obtained from the gravitino/five-form equation rather than from the
dilatino equation alone, and the same warp-factor result still follows.

\subsubsection{Why the condition makes the second-order equations
simple}\label{why-the-condition-makes-the-second-order-equations-simple}

There is also a useful geometric explanation of why the authors impose
precisely this combination. For a radial scalar \(f(r)\), \[
\sqrt{-g}
\ \propto\
e^{nA+qB}r^{q-1},
\qquad
g^{rr}=e^{-2B}.
\] Hence \[
\begin{aligned}
\nabla^2 f
&=\frac1{\sqrt{-g}}
\partial_r\left(\sqrt{-g}\,g^{rr}f'\right)\\
&=e^{-2B}
\left[
f''
+\left(
\frac{q-1}{r}
+nA'
+(q-2)B'
\right)f'
\right].
\end{aligned}
\] Since \(q-2=\tilde n\), \[
\nabla^2 f
=e^{-2B}
\left[
f''
+\left(
\frac{q-1}{r}
+nA'
+\tilde nB'
\right)f'
\right].
\] When the BPS condition holds, \[
nA'+\tilde nB'=0,
\] the extra warp-factor derivative disappears. Apart from the overall
\(e^{-2B}\), the radial operator becomes the flat transverse Laplacian,
\[
f''+\frac{q-1}{r}f'.
\] This is why the remaining equations collapse to the statement that
\(H\) is harmonic in the transverse space: \[
\frac1{r^{q-1}}
\frac{d}{dr}\left(r^{q-1}\frac{dH}{dr}\right)=0
\qquad
(r\neq0).
\] Thus (18.33) is often called the harmonic or extremal gauge
condition. In this presentation it is not imposed merely as a convenient
coordinate gauge: supersymmetry selects it for the isotropic BPS ansatz.

\subsubsection{Checks}\label{checks-31}

For a ten-dimensional D\(p\)-brane, \(\Delta=4\), and equation (18.34)
gives \[
A=-\frac{7-p}{16}\log H,
\qquad
B=\frac{p+1}{16}\log H.
\] Therefore \[
(p+1)A+(7-p)B=0.
\]

For the M2-brane, \(d=11\), \(n=3\), and \(\tilde n=6\). Its metric is
\[
ds^2=H^{-2/3}ds_\parallel^2+H^{1/3}ds_\perp^2,
\] so \[
A=-\frac13\log H,
\qquad
B=\frac16\log H,
\] and \[
3A+6B=0.
\]

For the M5-brane, \(n=6\) and \(\tilde n=3\): \[
ds^2=H^{-1/3}ds_\parallel^2+H^{2/3}ds_\perp^2,
\] so \[
A=-\frac16\log H,
\qquad
B=\frac13\log H,
\] and again \[
6A+3B=0.
\]

These checks show that the relation is independent of the presence of a
dilaton and survives electric/magnetic duality.

\subsubsection{Result}\label{result-37}

The brane projector \[
\Pi_s\epsilon=\epsilon,
\qquad
\Pi_s=\frac12(1+s\mathcal K_p),
\qquad
\mathcal K_p^2=1,
\] leaves one half of the supercharges. After imposing it, the
Killing-spinor equations reduce to \[
\boxed{
A'=-\frac{2\tilde n}{\Delta(d-2)}u,
\qquad
B'=\frac{2n}{\Delta(d-2)}u,
\qquad
\phi'=\frac{2\zeta a_p}{\Delta}u
}.
\] The first two equations imply \[
\boxed{
\frac{d}{dr}(nA+\tilde nB)=0
},
\] and hence \[
\boxed{
nA+\tilde nB=\mathrm{const}=0
}
\] for asymptotically Minkowski boundary conditions.

For a D\(p\)-brane, the same result follows explicitly from \[
B_S'=-A_S',
\qquad
\phi'=(p-3)A_S',
\qquad
A=A_S-\frac{\phi}{4},
\qquad
B=B_S-\frac{\phi}{4}.
\] The simplicity of (18.33) is therefore not accidental: it is the
flux-independent combination of the longitudinal and transverse
first-order BPS equations.

\section{\texorpdfstring{Q3: Supersymmetry and the slope rule for
\((p,q)\) five-brane
webs}{Q3: Supersymmetry and the slope rule for (p,q) five-brane webs}}\label{q3-supersymmetry-and-the-slope-rule-for-pq-five-brane-webs}

Consider the type IIB \((p,q)\) 5-brane web. The 5-branes extends in the
01234 directions and one direction in the \(x^5-x^6\) plane. Suppose we
have D5-brane along \(x^5\) and NS5-brane along \(x^6\), we can find the
susy condition \[
\begin{aligned}
\text{D5}&:\epsilon_{L}=\Gamma^0\Gamma^1\Gamma^2\Gamma^3\Gamma^4\Gamma^5\epsilon_{R}\\
\text{NS5}&:\epsilon_{L}=\Gamma^0\Gamma^1\Gamma^2\Gamma^3\Gamma^4\Gamma^6\epsilon_{L} \ , \ \epsilon_{R}=-\Gamma^0\Gamma^1\Gamma^2\Gamma^3\Gamma^4\Gamma^6\epsilon_{R} 
\end{aligned}
\] Show that 8 of the 32 susy are preserved. In the construction of
\((p,q)\) 5-brane web, one crucial fact is that a \((p,q)\) 5-brane
along the direction with slope
\(\delta x^5+i\delta x^6\propto p+\tau q\) in the \(x^5-x^6\) plane,
where \(\tau=C_{0}+ie^{-\Phi}\) can be set to \(\tau=i\). Derive this
fact.

\begin{quote}
\textbf{Textbook location and related material.} Chapter 18, Section
18.6 (pp.~705--716): type-IIB \(SL(2,\mathbb Z)\), \((p,q)\) objects,
and five-brane webs. Section 18.4 supplies the projector and
central-charge logic.
\end{quote}

\subsection{Answer}\label{answer-38}

\subsubsection{Core point}\label{core-point-38}

There are two related statements.

First, the D5-brane along \((012345)\) and the NS5-brane along
\((012346)\) impose two independent but commuting half-BPS projectors.
Therefore their common eigenspace has \[
\frac12\cdot\frac12\cdot32=8
\] real supercharges. This is the minimal \(\mathcal N=1\) supersymmetry
of the common five-dimensional spacetime \((01234)\).

Second, a general \((p,q)_5\) charge has the complex BPS phase \[
e^{i\alpha_{p,q}}
=\frac{p+\tau q}{|p+\tau q|}.
\] Its charge-space projector preserves the same eight spinors as the
D5--NS5 intersection only if its spatial angle \(\theta\) in the
\(x^5\)-\(x^6\) plane equals this phase: \[
\boxed{
\theta=\alpha_{p,q}=\arg(p+\tau q)
}.
\] Equivalently, \[
\boxed{
\delta x^5+i\delta x^6
\ \propto\
p+\tau q
}.
\]

The slope rule and the normalized-coordinate convention \(\tau=i\) are
equations (11) and (16)--(18) of the original five-brane-web
construction by \href{https://arxiv.org/abs/hep-th/9710116}{Aharony,
Hanany, and Kol}. The supersymmetry projectors used below follow the
book's page 698 conventions.

\subsubsection{Writing the two conditions as
projectors}\label{writing-the-two-conditions-as-projectors}

The type-IIB supersymmetry parameter is \[
\epsilon=
\begin{pmatrix}
\epsilon_L\\
\epsilon_R
\end{pmatrix},
\] where \(\epsilon_L\) and \(\epsilon_R\) are Majorana--Weyl spinors of
the same chirality, each with 16 real components.

Define \[
\Gamma_D=\Gamma^{012345},
\qquad
\Gamma_N=\Gamma^{012346}.
\] For a product of six gamma matrices containing one timelike
direction, \[
\Gamma_D^2=\Gamma_N^2=1.
\]

The D5 condition in the question, \[
\epsilon_L=\Gamma_D\epsilon_R,
\] also implies \[
\epsilon_R=\Gamma_D\epsilon_L.
\] Therefore it is equivalent to \[
K_D\epsilon=\epsilon,
\qquad
K_D=\Gamma_D\sigma_1,
\] because \[
\Gamma_D\sigma_1
\begin{pmatrix}
\epsilon_L\\
\epsilon_R
\end{pmatrix}
=
\begin{pmatrix}
\Gamma_D\epsilon_R\\
\Gamma_D\epsilon_L
\end{pmatrix}.
\]

Similarly, the two NS5 conditions are equivalent to \[
K_N\epsilon=\epsilon,
\qquad
K_N=\Gamma_N\sigma_3,
\] since \[
\Gamma_N\sigma_3
\begin{pmatrix}
\epsilon_L\\
\epsilon_R
\end{pmatrix}
=
\begin{pmatrix}
\Gamma_N\epsilon_L\\
-\Gamma_N\epsilon_R
\end{pmatrix}.
\]

Both operators are involutions: \[
K_D^2=K_N^2=1.
\] Thus \[
\Pi_D=\frac12(1+K_D),
\qquad
\Pi_N=\frac12(1+K_N)
\] are rank-half projectors.

\subsubsection{Why the two projectors
commute}\label{why-the-two-projectors-commute}

Let \[
\Gamma_{(5)}=\Gamma^{01234},
\] so \[
\Gamma_D=\Gamma_{(5)}\Gamma^5,
\qquad
\Gamma_N=\Gamma_{(5)}\Gamma^6.
\] Because \(\Gamma_{(5)}\) contains five gamma matrices, both
\(\Gamma^5\) and \(\Gamma^6\) anticommute with it. It follows that \[
\Gamma_D\Gamma_N=-\Gamma_N\Gamma_D.
\] The Pauli matrices also anticommute: \[
\sigma_1\sigma_3=-\sigma_3\sigma_1.
\] The two minus signs cancel in the full operators: \[
\begin{aligned}
K_DK_N
&=(\Gamma_D\Gamma_N)(\sigma_1\sigma_3)\\
&=(-\Gamma_N\Gamma_D)(-\sigma_3\sigma_1)\\
&=K_NK_D.
\end{aligned}
\] Hence \[
[K_D,K_N]=0.
\]

The two projectors are independent: their product contains
\(\Gamma^{56}\sigma_2\) and is not \(\pm1\). They can therefore be
diagonalized simultaneously, and their common \(+1\) eigenspace has
dimension \[
\dim\ker(1-K_D)\cap\ker(1-K_N)
=\frac{32}{2\cdot2}
=8.
\]

\subsubsection{\texorpdfstring{The same count directly in terms of
\(\epsilon_L,\epsilon_R\)}{The same count directly in terms of \textbackslash epsilon\_L,\textbackslash epsilon\_R}}\label{the-same-count-directly-in-terms-of-epsilon_lepsilon_r}

The NS5 condition on \(\epsilon_R\) is \[
\Gamma_N\epsilon_R=-\epsilon_R.
\] Since \(\Gamma_N^2=1\), its \(-1\) eigenspace has dimension eight.
Choose any \(\epsilon_R\) in this eigenspace and define \[
\epsilon_L=\Gamma_D\epsilon_R.
\] Then, using \(\{\Gamma_D,\Gamma_N\}=0\), \[
\Gamma_N\epsilon_L
=\Gamma_N\Gamma_D\epsilon_R
=-\Gamma_D\Gamma_N\epsilon_R
=\Gamma_D\epsilon_R
=\epsilon_L.
\] Thus the NS5 condition on \(\epsilon_L\) is automatic. No further
constraint is produced, so the solution space is parametrized by the
eight allowed components of \(\epsilon_R\).

This also shows why it is essential that the D5 and NS5 occupy different
directions in the \(x^5\)-\(x^6\) plane. The gamma-matrix
anticommutation is exactly what converts the \(-1\) NS5 eigenvalue of
\(\epsilon_R\) into the \(+1\) eigenvalue required for \(\epsilon_L\).

\subsubsection{\texorpdfstring{The projector for a general angled
\((p,q)\)
five-brane}{The projector for a general angled (p,q) five-brane}}\label{the-projector-for-a-general-angled-pq-five-brane}

Let the oriented spatial direction of a five-brane in the plane be \[
\hat t_\theta
=\cos\theta\,\partial_5+\sin\theta\,\partial_6,
\] and define \[
\Gamma_\theta
=\cos\theta\,\Gamma^5+\sin\theta\,\Gamma^6.
\]

In the convention of the question, \[
(1,0)_5=\mathrm{D5},
\qquad
(0,1)_5=\mathrm{NS5}.
\]

\paragraph{The charge-space matrix from the IIB BPS
algebra}\label{the-charge-space-matrix-from-the-iib-bps-algebra}

The two Pauli matrices in the elementary projectors tell us how the two
elementary five-brane charges enter the type-IIB supersymmetry algebra:
\[
\mathrm{D5\ charge}:\ \sigma_1,
\qquad
\mathrm{NS5\ charge}:\ \sigma_3
\] at the canonical point \(\tau=i\). More invariantly, the bare integer
charge vector \((p,q)\) is an \(SL(2)\) doublet, whereas the
supersymmetry spinors transform under the local \(SO(2)\) associated
with the coset \[
\frac{SL(2,\mathbb R)}{SO(2)}.
\] The axio-dilaton supplies the coset vielbein that converts the
integer \(SL(2)\) charge vector into a local orthonormal \(SO(2)\)
vector.

A convenient real coset representative is \[
V(\tau)
=\frac1{\sqrt{\tau_2}}
\begin{pmatrix}
1&\tau_1\\
0&\tau_2
\end{pmatrix}.
\] Thus the local orthonormal charge vector is \[
\begin{pmatrix}
z_1\\
z_2
\end{pmatrix}
=V(\tau)
\begin{pmatrix}
p\\
q
\end{pmatrix}
=\frac1{\sqrt{\tau_2}}
\begin{pmatrix}
p+\tau_1q\\
\tau_2q
\end{pmatrix},
\] or, in complex notation, \[
z_1+iz_2=\frac{p+\tau q}{\sqrt{\tau_2}}.
\] The common positive factor \(1/\sqrt{\tau_2}\) affects the
Einstein-frame magnitude but not the direction of this vector. In the
string-frame normalization convention used for five-brane webs, the
corresponding five-brane central-charge matrix acting on the IIB spinor
doublet is \[
\boxed{
\mathcal Z_5(p,q;\tau)
=T_{\mathrm{D5}}
\left[
(p+\tau_1q)\sigma_1+\tau_2q\,\sigma_3
\right]
}
\] or equivalently \[
\boxed{
\mathcal Z_5
=T_{\mathrm{D5}}
\left[
\operatorname{Re}(p+\tau q)\sigma_1
+\operatorname{Im}(p+\tau q)\sigma_3
\right].
}
\] This is the precise step that was missing from the previous answer.
It is the central-charge, or zero-world-volume-flux limit, of the
\(SL(2,\mathbb Z)\)-covariant kappa-symmetry construction of a type-IIB
\((p,q)\) five-brane; see, for example, the
\href{https://arxiv.org/abs/hep-th/9905019}{covariant super-five-brane
construction}.

Since \[
\sigma_1^2=\sigma_3^2=1,
\qquad
\{\sigma_1,\sigma_3\}=0,
\] the square of this matrix is \[
\begin{aligned}
\mathcal Z_5^2
&=T_{\mathrm{D5}}^2
\left[
(p+\tau_1q)^2+(\tau_2q)^2
\right]\mathbf1\\
&=T_{\mathrm{D5}}^2|p+\tau q|^2\mathbf1.
\end{aligned}
\] The largest eigenvalue of the central-charge matrix gives the BPS
bound. Thus, in the string-frame convention used for the web, \[
T_{(p,q)_5}
=T_{\mathrm{D5}}\,|p+\tau q|.
\] For a saturated brane, divide the central-charge matrix by its
tension: \[
\begin{aligned}
\widehat{\mathcal Z}_5
&=\frac{\mathcal Z_5}{T_{(p,q)_5}}\\
&=\frac{\operatorname{Re}(p+\tau q)}{|p+\tau q|}\sigma_1
+\frac{\operatorname{Im}(p+\tau q)}{|p+\tau q|}\sigma_3.
\end{aligned}
\] Defining \[
e^{i\alpha}
=\frac{p+\tau q}{|p+\tau q|}
\] gives \[
\boxed{
\widehat{\mathcal Z}_5
=\cos\alpha\,\sigma_1+\sin\alpha\,\sigma_3
}.
\] It is an involution, \[
\widehat{\mathcal Z}_5^{\,2}=1,
\] as required for the internal part of a kappa-symmetry or
supersymmetry projector.

At \(\tau=i\) this statement is especially transparent: \[
\mathcal Z_5
=T_{\mathrm{D5}}(p\sigma_1+q\sigma_3),
\] so the Pauli factor is simply the unit vector in the two-dimensional
D5--NS5 charge plane. At a general \(\tau\), the integer charges remain
\((p,q)\), but the scalar background dresses them into local orthonormal
components proportional to \((p+\tau_1q,\tau_2q)\), with the common
factor \(1/\sqrt{\tau_2}\) shown above.

There is an equivalent rotation description. Let \[
U_\alpha
=\exp\left(\frac{i\alpha}{2}\sigma_2\right).
\] Using the Pauli algebra, \[
U_\alpha\sigma_1U_\alpha^{-1}
=\cos\alpha\,\sigma_1+\sin\alpha\,\sigma_3
=\widehat{\mathcal Z}_5.
\] The half-angle in \(U_\alpha\) is present because it acts on the
spinor doublet; conjugation rotates the Pauli-vector by the full angle
\(\alpha\).

\paragraph{Why the normalized central charge enters a supersymmetry
projector}\label{why-the-normalized-central-charge-enters-a-supersymmetry-projector}

The remaining step follows directly from the BPS algebra. For a static
planar five-brane, the energy-density term and the five-brane
central-charge term in the type-IIB supercharge anticommutator can be
organized, after suppressing the positive charge-conjugation matrix, as
\[
\{Q,Q\}_{(p,q)_5}
\ \sim\
T_{(p,q)_5}\mathbf1
-\Gamma_{(5)}\Gamma_\theta\,\mathcal Z_5.
\] The overall sign here is fixed so that the elementary D5 condition is
\(K_D\epsilon=\epsilon\). Positivity of \(\{Q,Q\}\) gives the BPS
inequality \[
T_{(p,q)_5}\geq \|\mathcal Z_5\|.
\] A supersymmetry generated by \(\epsilon\) is preserved precisely when
the corresponding eigenvalue of this anticommutator vanishes. For a
saturated brane, \(T_{(p,q)_5}=\|\mathcal Z_5\|\), so \[
\left[
T_{(p,q)_5}\mathbf1
-\Gamma_{(5)}\Gamma_\theta\mathcal Z_5
\right]\epsilon=0.
\] Dividing by the tension gives \[
\Gamma_{(5)}\Gamma_\theta
\frac{\mathcal Z_5}{T_{(p,q)_5}}\epsilon
=\epsilon.
\] Therefore the operator that must appear in the
preserved-supersymmetry condition is not the unnormalized charge matrix,
but \[
\boxed{
K_{(p,q)}(\theta)
=\Gamma_{(5)}\Gamma_\theta\widehat{\mathcal Z}_5
},
\qquad
\widehat{\mathcal Z}_5
=\frac{\mathcal Z_5}{T_{(p,q)_5}}.
\] This also explains why normalization by the BPS tension is
compulsory: it makes \(K_{(p,q)}^2=1\), so \((1+K_{(p,q)})/2\) is a
genuine rank-half projector.

\paragraph{The spacetime gamma matrix from the geometric
rotation}\label{the-spacetime-gamma-matrix-from-the-geometric-rotation}

The reference D5-brane lies along \(x^5\) and has spacetime factor \[
\Gamma_{(5)}\Gamma^5=\Gamma^{012345}.
\] Rotate its final tangent direction through angle \(\theta\) in the
\(x^5\)-\(x^6\) plane. The spin lift of this rotation is \[
S_\theta
=\exp\left(-\frac{\theta}{2}\Gamma^{56}\right).
\] The Clifford algebra gives \[
S_\theta\Gamma^5S_\theta^{-1}
=\cos\theta\,\Gamma^5+\sin\theta\,\Gamma^6
=\Gamma_\theta.
\] Since \(\Gamma_{(5)}=\Gamma^{01234}\) is invariant under rotations in
the \(56\) plane, \[
S_\theta
\left(\Gamma_{(5)}\Gamma^5\right)
S_\theta^{-1}
=\Gamma_{(5)}\Gamma_\theta.
\]

\paragraph{Combining the independent
rotations}\label{combining-the-independent-rotations}

Start from the reference D5 involution \[
K_D=\Gamma_{(5)}\Gamma^5\sigma_1.
\] The spatial rotation \(S_\theta\) acts only on spacetime spinor
indices, while \(U_\alpha\) acts only on the IIB \((L,R)\) doublet.
Hence they commute. The involution for the rotated, charge-dressed brane
is obtained by conjugation: \[
\begin{aligned}
K_{(p,q)}(\theta)
&=S_\theta U_\alpha K_DU_\alpha^{-1}S_\theta^{-1}\\
&=
\left[
S_\theta(\Gamma_{(5)}\Gamma^5)S_\theta^{-1}
\right]
\left[
U_\alpha\sigma_1U_\alpha^{-1}
\right]\\
&=\Gamma_{(5)}\Gamma_\theta
\left(
\cos\alpha\,\sigma_1+\sin\alpha\,\sigma_3
\right).
\end{aligned}
\] Therefore \[
\boxed{
K_{(p,q)}(\theta)
=\Gamma_{(5)}\Gamma_\theta
\left(
\cos\alpha\,\sigma_1+\sin\alpha\,\sigma_3
\right)
}.
\] This formula is thus the product of two independently normalized
factors: \[
\boxed{
K_{(p,q)}
=
\underbrace{\Gamma_{(5)}\Gamma_\theta}_{\text{oriented world-volume}}
\underbrace{\widehat{\mathcal Z}_5}_{\text{normalized IIB charge matrix}}
}.
\] It is the standard BPS form ``world-volume gamma matrix times
normalized central charge.''

It obeys \[
K_{(p,q)}(\theta)^2=1,
\] so an isolated angled \((p,q)\) five-brane is half-BPS.

At \(\tau=i\) this reproduces the two elementary cases: \[
\begin{array}{c|c|c|c}
(p,q)&\alpha&\theta&K_{(p,q)}\\ \hline
(1,0)&0&0&\Gamma^{012345}\sigma_1=K_D\\
(0,1)&\pi/2&\pi/2&\Gamma^{012346}\sigma_3=K_N.
\end{array}
\]

\subsubsection{Alignment of the charge phase and spatial
angle}\label{alignment-of-the-charge-phase-and-spatial-angle}

Now restrict to a spinor \(\epsilon\) satisfying \[
K_D\epsilon=\epsilon,
\qquad
K_N\epsilon=\epsilon.
\] The two mixed gamma--Pauli structures act as \[
\Gamma_{(5)}\Gamma^5\sigma_3\epsilon
=\Gamma^{56}\epsilon,
\] and \[
\Gamma_{(5)}\Gamma^6\sigma_1\epsilon
=-\Gamma^{56}\epsilon.
\] To see the first identity, for example, use \(K_N\epsilon=\epsilon\)
to replace \[
\sigma_3\epsilon=\Gamma_{(5)}\Gamma^6\epsilon.
\]

Expanding the general projector gives \[
\begin{aligned}
K_{(p,q)}(\theta)\epsilon
={}&
\cos\theta\cos\alpha\,
K_D\epsilon
+\sin\theta\sin\alpha\,
K_N\epsilon\\
&+\cos\theta\sin\alpha\,
\Gamma_{(5)}\Gamma^5\sigma_3\epsilon\\
&+\sin\theta\cos\alpha\,
\Gamma_{(5)}\Gamma^6\sigma_1\epsilon.
\end{aligned}
\] Using the common D5--NS5 conditions, \[
K_{(p,q)}(\theta)\epsilon
=
\left[
\cos(\theta-\alpha)
-\sin(\theta-\alpha)\Gamma^{56}
\right]\epsilon.
\] Since \[
(\Gamma^{56})^2=-1,
\] the operator \[
\left[\cos(\theta-\alpha)-1\right]
-\sin(\theta-\alpha)\Gamma^{56}
\] has a nontrivial kernel only if \[
\theta-\alpha=0\pmod{2\pi}
\] for an oriented leg. Reversing both the geometric orientation and the
charge vector gives the same unoriented line.

Therefore \[
\boxed{
\theta=\alpha=\arg(p+\tau q)
},
\] which is equivalent to \[
\boxed{
\delta x^5+i\delta x^6
=\ell\,(p+\tau q),
\qquad
\ell>0
}.
\] This proves that adding any number of correctly oriented \((p,q)\)
legs does not impose a third independent projector: every such leg
preserves the same eight-dimensional common eigenspace.

\subsubsection{Force balance gives the same slope
rule}\label{force-balance-gives-the-same-slope-rule}

The supersymmetry result has a simple mechanical consequence. Treat the
plane as the complex plane. The tension vector of an oriented leg is \[
\mathcal T_{p,q}
=T_{(p,q)_5}e^{i\theta}.
\] Using the BPS tension and the angle rule, \[
\mathcal T_{p,q}
=T_{\mathrm{D5}}|p+\tau q|
\frac{p+\tau q}{|p+\tau q|}
=T_{\mathrm{D5}}(p+\tau q).
\] At a junction, charge conservation is \[
\sum_i p_i=0,
\qquad
\sum_i q_i=0.
\] It immediately implies \[
\sum_i\mathcal T_{p_i,q_i}
=T_{\mathrm{D5}}
\left(
\sum_i p_i+\tau\sum_iq_i
\right)
=0.
\] Thus the same relation that aligns the supersymmetry projector also
guarantees zero net force at every vertex.

\subsubsection{\texorpdfstring{How \(\tau\) is normalized to
\(i\)}{How \textbackslash tau is normalized to i}}\label{how-tau-is-normalized-to-i}

Write \[
\tau=\tau_1+i\tau_2,
\qquad
\tau_2>0.
\] The physical slope rule says \[
\delta x^5+i\delta x^6
=\ell\left(p+\tau q\right).
\] Separating real and imaginary parts, \[
\delta x^5
=\ell(p+\tau_1q),
\qquad
\delta x^6
=\ell\tau_2q.
\]

Now introduce coordinates in the real basis \((1,\tau)\): \[
\boxed{
x^5=\tilde x^5+\tau_1\tilde x^6,
\qquad
x^6=\tau_2\tilde x^6
}.
\] Equivalently, \[
\boxed{
\tilde x^5=x^5-\frac{\tau_1}{\tau_2}x^6,
\qquad
\tilde x^6=\frac{x^6}{\tau_2}
}.
\] Then every \((p,q)\) leg obeys \[
\delta\tilde x^5=\ell p,
\qquad
\delta\tilde x^6=\ell q,
\] or \[
\boxed{
\delta\tilde x^5+i\delta\tilde x^6
=\ell(p+iq)
}.
\] Thus, in the normalized web diagram, \[
\boxed{\tau_{\mathrm{web}}=i},
\] and the slope is simply \[
\delta\tilde x^5:\delta\tilde x^6=p:q.
\]

For \(\tau_1\neq0\), this transformation shears the plane; for
\(\tau_2\neq1\), it rescales one direction. Indeed, the physical
Euclidean metric becomes \[
(dx^5)^2+(dx^6)^2
=
\left(d\tilde x^5+\tau_1d\tilde x^6\right)^2
+\tau_2^2(d\tilde x^6)^2.
\] Therefore the normalized diagram is not claiming that the physical
axio-dilaton has dynamically changed to \(i\) while every other datum
remains fixed. It is a convenient affine drawing of the web in which the
\((p,q)\) charge lattice itself supplies the coordinate grid. Physical
lengths and tensions must still be interpreted with the transformed
metric and tension scale.

For the decoupled low-energy five-dimensional theory, one can
simultaneously rescale and shear the web while holding its BPS mass
parameters fixed. In that restricted sense \(\tau\) is redundant, which
is the sense intended in the five-brane-web construction.

\subsubsection{\texorpdfstring{Distinction from an \(SL(2,\mathbb Z)\)
duality}{Distinction from an SL(2,\textbackslash mathbb Z) duality}}\label{distinction-from-an-sl2mathbb-z-duality}

At the level of classical type-IIB supergravity, the continuous group
\(SL(2,\mathbb R)\) acts transitively on the upper half-plane. The real
matrix \[
M_\tau
=
\begin{pmatrix}
1/\sqrt{\tau_2}&-\tau_1/\sqrt{\tau_2}\\
0&\sqrt{\tau_2}
\end{pmatrix}
\in SL(2,\mathbb R)
\] sends \[
\tau\longmapsto
\frac{a\tau+b}{c\tau+d}
=i.
\] However, as the book stresses around equation (18.76), the exact
string-theory duality group is \(SL(2,\mathbb Z)\), not all of
\(SL(2,\mathbb R)\). A generic \(M_\tau\) is not integral and does not
preserve the integer \((p,q)\) charge lattice. Consequently, \[
\text{generic }\tau
\quad\not\sim_{SL(2,\mathbb Z)}\quad
i.
\]

The safe statement is therefore:

\begin{itemize}
\tightlist
\item
  one may set \(\tau=i\) in the abstract normalized web coordinates
  without changing the web's low-energy combinatorial data;
\item
  one may also map \(\tau\) to \(i\) at the level of classical
  supergravity by an \(SL(2,\mathbb R)\) transformation;
\item
  one may not generally claim that an arbitrary full type-IIB string
  vacuum is \(SL(2,\mathbb Z)\)-equivalent to the vacuum at \(\tau=i\).
\end{itemize}

\subsubsection{Checks}\label{checks-32}

For a \((1,1)_5\) leg at \(\tau=i\), \[
\alpha_{1,1}=\arg(1+i)=\frac{\pi}{4}.
\] Thus it must run at \(45^\circ\) and has projector \[
K_{(1,1)}
=\Gamma_{(5)}
\frac{\Gamma^5+\Gamma^6}{\sqrt2}
\frac{\sigma_1+\sigma_3}{\sqrt2}.
\] On the common D5--NS5 eigenspace this acts as the identity, so it
introduces no additional supersymmetry breaking.

At the elementary three-leg vertex with incoming charges \[
(1,0)+(0,1)+(-1,-1)=0,
\] the three tension vectors at \(\tau=i\) are proportional to \[
1,\qquad i,\qquad-1-i,
\] which sum to zero.

If \(C_0=0\) but \(g_s\neq1\), then \[
\tau=\frac{i}{g_s}.
\] A \((p,q)\) leg has physical slope \[
\tan\theta=\frac{q}{g_sp}.
\] The normalized rescaling \(\tilde x^6=g_sx^6\) turns this into \[
\tan\tilde\theta=\frac{q}{p},
\] which is precisely the \(\tau=i\) drawing convention.

\subsubsection{Result}\label{result-38}

The elementary projectors are \[
\boxed{
K_D=\Gamma^{012345}\sigma_1,
\qquad
K_N=\Gamma^{012346}\sigma_3
}.
\] They satisfy \[
K_D^2=K_N^2=1,
\qquad
[K_D,K_N]=0,
\] and are independent, so \[
\boxed{32\longrightarrow16\longrightarrow8}.
\]

The scalar-dressed charge matrix, its BPS norm, and its normalized form
are \[
\boxed{
\begin{aligned}
\mathcal Z_5
&=T_{\mathrm{D5}}
\left[
\operatorname{Re}(p+\tau q)\sigma_1
+\operatorname{Im}(p+\tau q)\sigma_3
\right],\\
T_{(p,q)_5}
&=\|\mathcal Z_5\|
=T_{\mathrm{D5}}|p+\tau q|,\\
\widehat{\mathcal Z}_5
&=\frac{\mathcal Z_5}{T_{(p,q)_5}}
=\cos\alpha\,\sigma_1+\sin\alpha\,\sigma_3,
\qquad
e^{i\alpha}=\frac{p+\tau q}{|p+\tau q|}.
\end{aligned}
}
\] BPS saturation gives \[
\left[
T_{(p,q)_5}\mathbf1
-\Gamma^{01234}\Gamma_\theta\mathcal Z_5
\right]\epsilon=0,
\] so the preserved-supersymmetry involution must be the world-volume
gamma matrix multiplied by the normalized charge matrix.

For a general leg, \[
\boxed{
K_{(p,q)}(\theta)
=\Gamma^{01234}
\left(\cos\theta\,\Gamma^5+\sin\theta\,\Gamma^6\right)
\left(\cos\alpha\,\sigma_1+\sin\alpha\,\sigma_3\right),
\qquad
e^{i\alpha}=\frac{p+\tau q}{|p+\tau q|}
}.
\] It preserves the common eight supercharges exactly when \[
\boxed{
\theta=\alpha
\quad\Longleftrightarrow\quad
\delta x^5+i\delta x^6\propto p+\tau q
}.
\]

Finally, for \(\tau=\tau_1+i\tau_2\), \[
\boxed{
x^5=\tilde x^5+\tau_1\tilde x^6,
\qquad
x^6=\tau_2\tilde x^6
}
\] turns the slope rule into \[
\boxed{
\delta\tilde x^5+i\delta\tilde x^6
\propto p+iq
}.
\] This is the precise meaning of setting \(\tau=i\) in a normalized
\((p,q)\) five-brane web.

\section{Q4: Type-I D-string and heterotic-string spectrum
match}\label{q4-type-i-d-string-and-heterotic-string-spectrum-match}

In page710-711 they showed the massless excitations of the D-string in
type I are precisely the fields appearing on the world-sheet of the
heterotic string. What are the massless excitations of a D-brane? Derive
this equivalence in details.

\begin{quote}
\textbf{Textbook location and related material.} Chapter 18, Section
18.6: type-I/heterotic duality and the D-string realization of the
heterotic string. Sections 9.4 and 13.3 provide the type-I projections
and D-brane spectrum.
\end{quote}

\subsection{Answer}\label{answer-39}

\subsubsection{Core point}\label{core-point-39}

For an ordinary oriented type-II D\(p\)-brane, the massless open strings
with both endpoints on the brane form the dimensional reduction of
ten-dimensional \(\mathcal N=1\) super-Yang--Mills to \(p+1\)
dimensions: \[
\boxed{
A_M
\longrightarrow
\left(A_\alpha,\Phi^i\right),
\qquad
\lambda_{10}
\longrightarrow
\lambda_{p+1}
}.
\] Here \(A_\alpha\) is a gauge field tangent to the D\(p\)-brane, the
\(9-p\) fields \(\Phi^i\) describe transverse fluctuations of the brane,
and \(\lambda_{p+1}\) is the dimensionally reduced ten-dimensional
Majorana--Weyl gaugino. For \(N\) coincident oriented branes, all these
fields are in the adjoint of \(U(N)\).

The type-I D-string is more special because type I is the orientifold \[
\text{type I}
=\frac{\text{type IIB}}{\Omega}
+32\ \mathrm{D9\text{-}branes}.
\] Consequently, the massless fields on a single D1 receive
contributions from two kinds of open strings:

\begin{enumerate}
\def\labelenumi{\arabic{enumi}.}
\tightlist
\item
  D1--D1 strings give eight transverse bosons and eight right-moving
  fermions;
\item
  D1--D9 and D9--D1 strings together give 32 real left-moving fermions.
\end{enumerate}

Thus the final physical world-sheet fields are \[
\boxed{
X^i(\sigma,\tau),
\qquad
\psi^A(\sigma^-),
\qquad
\lambda^a(\sigma^+)
},
\] with \[
i=2,\ldots,9,
\qquad
A=1,\ldots,8,
\qquad
a=1,\ldots,32.
\] These are exactly the Green--Schwarz world-sheet fields of the
ten-dimensional \(\operatorname{Spin}(32)/\mathbb Z_2\) heterotic
string: eight transverse non-chiral bosons, eight right-moving
Green--Schwarz fermions, and 32 left-moving real fermions generating the
\(\mathfrak{so}(32)_1\) current algebra.

The derivation has three essential ingredients: the ordinary D-brane
open-string spectrum, the type-I \(\Omega\) projection in the D1--D1
sector, and the change of oscillator moding in the eight mixed D1--D9
directions.

\subsubsection{\texorpdfstring{Massless open strings on a general
D\(p\)-brane}{Massless open strings on a general Dp-brane}}\label{massless-open-strings-on-a-general-dp-brane}

Take a flat D\(p\)-brane along \(x^0,\ldots,x^p\). An open string whose
two endpoints lie on it obeys \[
\left.\partial_\sigma X^\alpha\right|_{\partial\Sigma}=0,
\qquad
\alpha=0,\ldots,p,
\] and \[
\left.\delta X^i\right|_{\partial\Sigma}=0,
\qquad
i=p+1,\ldots,9.
\] The first are Neumann--Neumann boundary conditions and the second are
Dirichlet--Dirichlet boundary conditions. The open-string center-of-mass
momentum therefore has only world-volume components: \[
k^M=(k^\alpha,0),
\qquad
k^2=k_\alpha k^\alpha.
\]

For NN or DD boundary conditions, the world-sheet fermions are
half-integer moded in the NS sector and integer moded in the R sector.
The mass formulas are \[
\alpha'M_{\mathrm{NS}}^2
=N_{\mathrm{osc}}-\frac12,
\qquad
\alpha'M_{\mathrm R}^2
=N_{\mathrm{osc}}.
\]

\paragraph{NS sector: gauge field and transverse
scalars}\label{ns-sector-gauge-field-and-transverse-scalars}

The NS ground state has \[
\alpha'M^2=-\frac12
\] and is removed by the open-string GSO projection. The first GSO-even
states have one fermionic oscillator: \[
\epsilon_M b_{-1/2}^M|0;k\rangle_{\mathrm{NS}}.
\] Their oscillator number is \(N_{\mathrm{osc}}=\frac12\), so they are
massless. Splitting the polarization into NN and DD components gives \[
\epsilon_\alpha b_{-1/2}^\alpha|0;k\rangle_{\mathrm{NS}}
\quad\longleftrightarrow\quad
A_\alpha,
\] and \[
\phi_i b_{-1/2}^i|0;k\rangle_{\mathrm{NS}}
\quad\longleftrightarrow\quad
\Phi^i.
\]

The super-Virasoro constraint \(G_{1/2}|\mathrm{phys}\rangle=0\) gives
\[
k^\alpha\epsilon_\alpha=0.
\] Null states identify \[
\epsilon_\alpha
\sim
\epsilon_\alpha+k_\alpha\Lambda,
\] which is the linearized gauge invariance \[
A_\alpha
\sim
A_\alpha+\partial_\alpha\Lambda.
\] There is no analogous condition on \(\phi_i\), because the open
string has no momentum \(k^i\) in a Dirichlet direction. The \(\Phi^i\)
are therefore world-volume scalars.

Their geometrical meaning follows either directly from the Dirichlet
boundary value or by T-duality: \[
\boxed{
X^i_{\mathrm{brane}}
=2\pi\alpha'\Phi^i
}
\] for a single brane, up to the conventional normalization of the
scalar. A constant expectation value of \(\Phi^i\) moves the D-brane in
the \(x^i\) direction.

\paragraph{R sector: the world-volume
gaugino}\label{r-sector-the-world-volume-gaugino}

In the R sector the zero-point energy vanishes. The ground states are
already massless and the zero modes obey \[
\{b_0^M,b_0^N\}
=\eta^{MN}.
\] Thus they generate the ten-dimensional Clifford algebra and the R
ground states form a ten-dimensional spinor. The GSO projection selects
one Majorana--Weyl chirality, leaving 16 real components before the
massless Dirac equation is imposed: \[
\Gamma_{11}u=+u.
\] The zero-mode super-Virasoro constraint is \[
G_0|u;k\rangle_{\mathrm R}=0
\quad\Longleftrightarrow\quad
k_\alpha\Gamma^\alpha u=0.
\] For a nonzero null world-volume momentum this leaves eight on-shell
fermionic polarizations. Together with \(A_\alpha\) and \(\Phi^i\),
these are the fields obtained by reducing ten-dimensional
\(\mathcal N=1\) super-Yang--Mills to the D\(p\)-brane world-volume.

For a Chan--Paton stack of \(N\) identical oriented branes, the states
carry endpoint labels \(|r,s\rangle\), \(r,s=1,\ldots,N\). Hence \[
A_\alpha,\Phi^i,\lambda
\in\operatorname{End}(\mathbb C^N),
\] and the low-energy gauge group is \(U(N)\). The abelian scalars give
the center-of-mass position, while eigenvalue differences measure
relative brane positions.

\subsubsection{A useful zero-point-energy formula for mixed boundary
conditions}\label{a-useful-zero-point-energy-formula-for-mixed-boundary-conditions}

The D1--D9 sector differs from the D1--D1 sector because it has mixed
boundary conditions. It is useful to derive the relevant zero-point
energy once.

For one real bosonic or fermionic oscillator, zeta-function
regularization gives \[
\sum_{n=1}^\infty n
=-\frac1{12},
\qquad
\sum_{n=0}^\infty\left(n+\frac12\right)
=\frac1{24}.
\] Consequently, \[
\begin{array}{c|cc}
&\text{integer modes}&\text{half-integer modes}\\ \hline
\text{boson}
&-\dfrac1{24}
&+\dfrac1{48}\\[2mm]
\text{fermion}
&+\dfrac1{24}
&-\dfrac1{48}
\end{array}
\] where the opposite sign in the fermionic row comes from fermionic
normal ordering.

For an NN or DD direction, the NS fermion is half-integer moded, so one
transverse boson--fermion pair contributes \[
-\frac1{24}-\frac1{48}
=-\frac1{16}.
\] For an ND or DN direction, the moding is reversed: the boson is
half-integer moded and the NS fermion is integer moded. Such a pair
contributes \[
\frac1{48}+\frac1{24}
=+\frac1{16}.
\] If \(\nu\) of the eight light-cone transverse directions are mixed,
then \[
\boxed{
a_{\mathrm{NS}}(\nu)
=(8-\nu)\left(-\frac1{16}\right)
+\nu\left(\frac1{16}\right)
=-\frac12+\frac{\nu}{8}
}.
\] In the R sector, bosons and fermions have the same moding for every
type of boundary condition and cancel pairwise: \[
\boxed{a_{\mathrm R}=0}.
\]

For a string with both endpoints on the same D-brane, \(\nu=0\) and
\(a_{\mathrm{NS}}=-\frac12\), reproducing the spectrum above. For a
D1--D9 string, \(\nu=8\) and therefore \[
\boxed{
a_{\mathrm{NS}}=+\frac12,
\qquad
a_{\mathrm R}=0
}.
\] This single calculation explains why the D1--D9 sector has massless
fermions but no massless bosons.

\subsubsection{The D1--D1 NS sector after the type-I
projection}\label{the-d1d1-ns-sector-after-the-type-i-projection}

The D1 extends along \(x^0,x^1\). Thus a D1--D1 string has NN boundary
conditions for \(M=\alpha=0,1\) and DD boundary conditions for
\(M=i=2,\ldots,9\).

Before the orientifold projection, its massless NS states are \[
b_{-1/2}^\alpha|0;k\rangle_{\mathrm{NS}}
\quad\longleftrightarrow\quad
A_\alpha,
\] and \[
b_{-1/2}^i|0;k\rangle_{\mathrm{NS}}
\quad\longleftrightarrow\quad
X^i.
\] The notation \(X^i\) is natural here because these scalars are
precisely the transverse embedding coordinates of the D-string.

Equation (9.75) gives the action of world-sheet parity on the fermionic
oscillators: \[
\Omega b_r^M\Omega^{-1}
=\eta_M e^{i\pi r}b_r^M,
\qquad
\eta_\alpha=+1,
\qquad
\eta_i=-1.
\] The sign difference between \(\eta_\alpha\) and \(\eta_i\) is the
difference between NN and DD boundary conditions. With the type-I
convention (9.76), \[
\Omega|0\rangle_{\mathrm{NS}}
=-i|0\rangle_{\mathrm{NS}}.
\] For \(r=-\frac12\), \[
e^{i\pi r}=e^{-i\pi/2}=-i.
\] Therefore \[
\begin{aligned}
\Omega\left(
b_{-1/2}^M|0\rangle_{\mathrm{NS}}
\right)
&=
\eta_M(-i)(-i)
b_{-1/2}^M|0\rangle_{\mathrm{NS}}\\
&=-\eta_M
b_{-1/2}^M|0\rangle_{\mathrm{NS}}.
\end{aligned}
\] The result is \[
\begin{array}{c|c|c}
\text{state}&\Omega\text{ eigenvalue}&\text{type-I projection}\\ \hline
b_{-1/2}^{0,1}|0\rangle_{\mathrm{NS}}
&-1
&\text{removed},\\
b_{-1/2}^{i}|0\rangle_{\mathrm{NS}}
&+1
&\text{kept}.
\end{array}
\] Thus \[
\boxed{
\text{D1--D1 NS sector:}\qquad
X^i(\sigma,\tau),
\quad i=2,\ldots,9
}.
\]

For a single type-I D1 the Chan--Paton group is
\(O(1)\cong\mathbb Z_2\), so there is no continuous world-volume gauge
field. Notice also that \(A_0,A_1\) would carry no local propagating
polarization in two dimensions even before the projection. The eight
surviving scalars are the Goldstone modes for the eight translations
broken by placing the D1 at a definite position in \(\mathbb R^8\).

\subsubsection{The D1--D1 R sector and its
chirality}\label{the-d1d1-r-sector-and-its-chirality}

All ten D1--D1 directions are either NN or DD, so all R-sector fermions
are integer moded. The ten zero modes generate \(\mathrm{Cl}(1,9)\): \[
b_0^M
=\frac1{\sqrt2}\Gamma^M.
\] After the open-string GSO projection, choose the positive-chirality
spinor \[
\Gamma_{11}=+1.
\] Under \[
SO(1,9)
\supset
SO(1,1)\times SO(8),
\] it decomposes in the book's conventions as \[
\boxed{
\mathbf{16}
\longrightarrow
\mathbf 8_s^{+}\oplus\mathbf 8_c^{-}
}.
\]

The superscript is the two-dimensional chirality. To determine which
summand survives, use the weight basis \[
|s\rangle
=|s_0;s_1,s_2,s_3,s_4\rangle,
\qquad
s_I=\pm\frac12,
\quad I=0,\ldots,4.
\] For a D1, equation (9.75) at \(r=0\) says that \(\Omega\) commutes
with the two NN zero modes and anticommutes with the eight DD zero
modes: \[
\Omega b_0^\alpha\Omega^{-1}=+b_0^\alpha,
\qquad
\Omega b_0^i\Omega^{-1}=-b_0^i.
\] Thus its action on the spinor module is, up to the ground-state phase
fixed by (9.76), the transverse chirality operator. The book obtains \[
\Omega|s\rangle
=-e^{i\pi(s_1+s_2+s_3+s_4)}|s\rangle.
\] Let \(n_-\) be the number of minus signs among \(s_1,\ldots,s_4\).
Then \[
s_1+s_2+s_3+s_4
=2-n_-,
\] and hence \[
\Omega|s\rangle
=-(-1)^{n_-}|s\rangle.
\] The invariant states have odd \(n_-\). Those are precisely the
\(\mathbf 8_c\) weights. Positive ten-dimensional chirality then
correlates them with negative \(SO(1,1)\) chirality, so \[
\boxed{
\mathbf{16}
\ \xrightarrow{\ \Omega\ }\
\mathbf 8_c^{-}
}.
\]

The remaining physical-state condition is the two-dimensional massless
Dirac equation \[
k_\alpha\Gamma^\alpha u_-=0.
\] In position space, multiply by \(\Gamma^0\) and use
\(\Gamma^{01}u_-=-u_-\). With mostly-plus signature, \[
\left(
-\partial_\tau+\Gamma^{01}\partial_\sigma
\right)u_-=0
\quad\Longrightarrow\quad
\left(\partial_\tau+\partial_\sigma\right)u_-=0.
\] Therefore \[
u_-=u_-(\tau-\sigma)=u_-(\sigma^-),
\] which is right-moving in the convention of the book. We obtain \[
\boxed{
\text{D1--D1 R sector:}\qquad
\psi^A(\sigma^-),
\quad A=1,\ldots,8
}.
\] These are the eight right-moving Green--Schwarz fermions paired by
supersymmetry with the eight transverse bosons \(X^i\).

\subsubsection{The D1--D9 and D9--D1
sectors}\label{the-d1d9-and-d9d1-sectors}

A D9 fills all ten dimensions. A D1--D9 string therefore has \[
\begin{array}{c|c}
x^{0,1}&\mathrm{NN}\\
x^{2},\ldots,x^9&\mathrm{DN}
\end{array}
\] and hence eight mixed directions: \[
\nu=8.
\]

\paragraph{Absence of massless bosons}\label{absence-of-massless-bosons}

The zero-point formula gives \[
\alpha'M_{\mathrm{NS}}^2
=N_{\mathrm{osc}}+\frac12.
\] Since \(N_{\mathrm{osc}}\geq0\), the entire NS sector is massive at
its lowest level: \[
M_{\mathrm{NS}}^2
\geq\frac1{2\alpha'}.
\] Thus the D1--D9 strings supply no additional massless bosonic fields.

\paragraph{The 32 massless chiral
fermions}\label{the-32-massless-chiral-fermions}

The R-sector zero-point energy is zero, so its ground state is massless.
In an ND or DN direction, however, the R fermions are half-integer
moded. Therefore the eight fields in directions \(2,\ldots,9\) have no R
zero modes. Only the common NN directions \(0,1\) have zero modes: \[
\{b_0^\alpha,b_0^\beta\}
=\eta^{\alpha\beta},
\qquad
\alpha,\beta=0,1.
\] The ground state is consequently only a two-dimensional spinor,
rather than a ten-dimensional spinor: \[
|s_0,a\rangle_{\mathrm R},
\qquad
s_0=\pm\frac12.
\] The label \(a=1,\ldots,32\) is carried by the D9 endpoint because
tadpole cancellation requires 32 D9-branes.

The GSO projection keeps \[
s_0=+\frac12.
\] Equivalently, the mutual-locality condition between D1--D1 and D1--D9
vertex operators selects this chirality. The physical state equation is
again \[
k_\alpha\Gamma^\alpha u_+=0.
\] Now \(\Gamma^{01}u_+=+u_+\), so \[
\left(
-\partial_\tau+\partial_\sigma
\right)u_+=0
\quad\Longrightarrow\quad
\left(\partial_\tau-\partial_\sigma\right)u_+=0.
\] It follows that \[
u_+=u_+(\tau+\sigma)=u_+(\sigma^+),
\] which is left-moving in the book's convention.

An oppositely oriented string belongs to the D9--D1 sector. World-sheet
parity reverses the string orientation and exchanges the two sectors: \[
\Omega:
\quad
\mathcal H_{\mathrm{D1-D9}}
\longleftrightarrow
\mathcal H_{\mathrm{D9-D1}}.
\] The orientifold projection therefore identifies their states rather
than adding a second independent copy. The result is exactly \[
\boxed{
\text{D1--D9 plus D9--D1:}\qquad
\lambda^a(\sigma^+),
\quad a=1,\ldots,32
},
\] where the \(\lambda^a\) are real left-moving Majorana--Weyl fermions
transforming in the vector representation of \(SO(32)\).

\subsubsection{Matching to the heterotic Green--Schwarz
world-sheet}\label{matching-to-the-heterotic-greenschwarz-world-sheet}

Collecting the two sectors gives \[
\begin{array}{c|c|c|c}
\text{D1 sector}
&\text{field}
&SO(8)\times SO(32)
&\text{two-dimensional propagation}\\ \hline
\mathrm{D1-D1,\ NS}
&X^i
&(\mathbf8_v,\mathbf1)
&\text{left and right}\\
\mathrm{D1-D1,\ R}
&\psi^A
&(\mathbf8_c,\mathbf1)
&\text{right-moving}\\
\mathrm{D1-D9/D9-D1,\ R}
&\lambda^a
&(\mathbf1,\mathbf{32})
&\text{left-moving}.
\end{array}
\]

This is exactly the physical Green--Schwarz description of the heterotic
string in static or light-cone gauge:

\begin{itemize}
\tightlist
\item
  \(X^i\) are the eight transverse embedding coordinates;
\item
  \(\psi^A\) are their eight right-moving supersymmetric partners;
\item
  \(\lambda^a\) are the 32 left-moving fermions of the gauge sector.
\end{itemize}

The gauge currents are the bilinears \[
\boxed{
J^{ab}
=i\lambda^a\lambda^b,
\qquad
J^{ab}=-J^{ba}
}.
\] There are \[
\frac{32\cdot31}{2}=496
\] such currents, and their operator products generate the level-one
affine algebra \[
\widehat{\mathfrak{so}}(32)_1.
\] This explains how open strings ending on the 32 D9-branes become the
left-moving gauge degrees of freedom of the dual heterotic closed
string.

The matching of chiral central charges is also exact. In light-cone
gauge, \[
c_L
=8+\frac{32}{2}
=24,
\] while \[
c_R
=8+\frac{8}{2}
=12.
\] These are the heterotic light-cone matter central charges: the
left-moving side is bosonic and the right-moving side is supersymmetric.

\subsubsection{How the heterotic GSO projection arises on the
D-string}\label{how-the-heterotic-gso-projection-arises-on-the-d-string}

The equality of local massless fields is not yet the complete global
statement. Thirty-two free left-moving fermions have four
\(\mathfrak{so}(32)_1\) conjugacy classes, conventionally denoted \[
O,\qquad V,\qquad S,\qquad C.
\] The \(\operatorname{Spin}(32)/\mathbb Z_2\) heterotic string keeps
the root and one spinor class and removes the vector and
conjugate-spinor classes: \[
\boxed{
O\oplus S
\quad\text{kept},
\qquad
V\oplus C
\quad\text{projected out}
}
\] up to the convention that exchanges \(S\) and \(C\). In the fermionic
formulation, this is implemented by summing over spin structures for all
32 fermions.

On one type-I D1 the continuous gauge field was projected out, but its
Chan--Paton group is not trivial: \[
O(1)\cong\mathbb Z_2.
\] Therefore a D1 world-sheet with nontrivial cycles can carry flat
\(\mathbb Z_2\) gauge bundles. On a torus, \[
H^1(T^2,\mathbb Z_2)
\cong
\mathbb Z_2\oplus\mathbb Z_2,
\] so there are four choices of holonomy: \[
(w_A,w_B)
\in
\{(+,+),(+,-),(-,+),(-,-)\}.
\] Because the D1--D9 fermions are charged under this \(O(1)\), a
negative holonomy reverses the sign of \(\lambda^a\) around the
corresponding cycle. Summing over the four flat bundles is therefore
precisely the common spin-structure sum for the 32 left-moving fermions.
It produces the heterotic GSO projection and refines the perturbative
Lie algebra \(\mathfrak{so}(32)\) to the global gauge group \[
\boxed{
\operatorname{Spin}(32)/\mathbb Z_2
}.
\]

Thus both the local fields and their global projection agree.

\subsubsection{Checks}\label{checks-33}

For the D1--D1 strings, \(\nu=0\), so \[
a_{\mathrm{NS}}=-\frac12.
\] One \(b_{-1/2}\) oscillator raises the level by \(\frac12\) and gives
a massless state, as required.

For the D1--D9 strings, \(\nu=8\), so \[
a_{\mathrm{NS}}=+\frac12.
\] There is no possible nonnegative oscillator level that makes the NS
state massless. The R ground state remains massless because
\(a_{\mathrm R}=0\).

The \(\Omega\) projection retains eight D1--D1 bosons and eight D1--D1
fermions, so the supersymmetric transverse sector is balanced: \[
8_B=8_F.
\] The 32 additional fermions are all of the opposite world-sheet
chirality and form the heterotic gauge sector rather than additional
supersymmetric partners.

Finally, \(\Omega\) identifies the D1--D9 and D9--D1 orientations.
Without this identification one would incorrectly obtain 64 rather than
32 real chiral fermions.

\subsubsection{Result}\label{result-39}

For an oriented type-II D\(p\)-brane, \[
\boxed{
\begin{aligned}
\mathrm{NS}:&\quad
b_{-1/2}^{\alpha}|0;k\rangle
\leftrightarrow A_\alpha,
\qquad
b_{-1/2}^{i}|0;k\rangle
\leftrightarrow\Phi^i,\\
\mathrm R:&\quad
|u;k\rangle,
\qquad
\Gamma_{11}u=u,
\qquad
k_\alpha\Gamma^\alpha u=0.
\end{aligned}
}
\] This is ten-dimensional \(\mathcal N=1\) super-Yang--Mills reduced to
\(p+1\) dimensions.

For \(\nu\) mixed ND or DN directions, \[
\boxed{
a_{\mathrm{NS}}
=-\frac12+\frac{\nu}{8},
\qquad
a_{\mathrm R}=0
}.
\]

For the type-I D1, \[
\boxed{
\begin{array}{c|c}
\mathrm{D1-D1,\ NS}
&8\,X^i(\sigma,\tau)\\
\mathrm{D1-D1,\ R}
&8\,\psi^A(\sigma^-)\\
\mathrm{D1-D9/D9-D1,\ NS}
&\text{no massless states}\\
\mathrm{D1-D9/D9-D1,\ R}
&32\,\lambda^a(\sigma^+).
\end{array}
}
\] Therefore \[
\boxed{
\mathcal H_{\mathrm{massless}}^{\mathrm{type\ I\ D1}}
=
\left\{
8X^i,\,
8\psi^A_R,\,
32\lambda^a_L
\right\}
=
\mathcal H_{\mathrm{world\text{-}sheet}}^{
\operatorname{Spin}(32)/\mathbb Z_2\ \mathrm{heterotic}}
}.
\] The four flat \(O(1)\cong\mathbb Z_2\) bundles on a toroidal D1
world-sheet implement the heterotic spin-structure sum and hence the
heterotic GSO projection.

\section{Q5: Deriving the dilaton exponents in the M-theory/type-IIA
metric
ansatz}\label{q5-deriving-the-dilaton-exponents-in-the-m-theorytype-iia-metric-ansatz}

In equation (18.105), the ansatz is \[
ds_{11}^2=e^{-2\Phi/3}g_{\mu \nu}dx^\mu dx^\nu+e^{4\Phi/3}(dx^{11}+C_{\mu}dx^\mu)^2 \ , \ \mu,\nu=0,\cdots,9
\] I want to know why the factors \(e^{-2\Phi/3}\) and \(e^{4\Phi/3}\)
are chosen like these? It seems that from the 11d action (18.104) to 10d
IIA (16.134a), we have the relation \[
\frac{1}{\kappa_{10}^2}\sqrt{ -g }R_{d=10}\sim \frac{2\pi g_{s}\sqrt{ \alpha' } }{\kappa_{11}^2}\sqrt{ -G }R_{d=11} \implies \sqrt{ -G }R_{d=11}=\sqrt{ -g }R_{d=10}
\] When we take the determinant of the 11d metric we find \[
\sqrt{ -G }=e^{2\Phi/3}\cdot e^{-10\Phi/3}\sqrt{ -g }=e^{-8\Phi/3}\sqrt{ -g }
\] What about the Ricci scalar part?

\begin{quote}
\textbf{Textbook location and related material.} Chapter 18, Section
18.7 (pp.~716--724): reduction of eleven-dimensional supergravity on a
circle and the M-theory/type-IIA metric ansatz. Section 16.4 lists the
ten- and eleven-dimensional actions being matched.
\end{quote}

\subsection{Answer}\label{answer-40}

\subsubsection{Core point}\label{core-point-40}

Your determinant is correct: \[
\boxed{
\sqrt{-G_{11}}
=e^{-8\Phi/3}\sqrt{-g_{10}}
}.
\] The missing point is that the eleven-dimensional Ricci scalar is not
equal to the ten-dimensional one. Its term proportional to \(R_{10}\)
carries the inverse Weyl factor \[
R_{11}
\supset
e^{2\Phi/3}R_{10}.
\] Therefore \[
\sqrt{-G_{11}}\,R_{11}
\supset
\sqrt{-g_{10}}\,
e^{-8\Phi/3}e^{2\Phi/3}R_{10}
=
\sqrt{-g_{10}}\,e^{-2\Phi}R_{10}.
\] This is exactly the string-frame gravitational coupling. Equation
(16.134a) does not contain a bare \(\sqrt{-g}\,R\) term; its NS--NS part
is \[
\boxed{
\sqrt{-g}\,e^{-2\Phi}
\left[
R+4(\partial\Phi)^2-\frac12|H_3|^2
\right]
}.
\]

The full metric reduction gives, before integrating by parts, \[
\boxed{
R_{11}
=e^{2\Phi/3}
\left[
R_{10}
+\frac{14}{3}\Box\Phi
-\frac{16}{3}(\partial\Phi)^2
-\frac14e^{2\Phi}F_2^2
\right]
},
\] where \(F_2=dC_1\). Consequently, \[
\boxed{
\sqrt{-G_{11}}R_{11}
=\sqrt{-g}
\left\{
e^{-2\Phi}
\left[
R+\frac{14}{3}\Box\Phi-\frac{16}{3}(\partial\Phi)^2
\right]
-\frac14F_2^2
\right\}.
}
\] After integrating the \(\Box\Phi\) term by parts, its contribution
combines with \(-\frac{16}{3}(\partial\Phi)^2\) to give
\(+4(\partial\Phi)^2\). The rest of the answer derives these statements
and explains why the two exponents in (18.105) are uniquely selected.

\subsubsection{The target is the type-IIA string-frame
action}\label{the-target-is-the-type-iia-string-frame-action}

Ignoring the Chern--Simons term, the relevant bosonic type-IIA action is
\[
S_{\mathrm{IIA}}
=\frac1{2\widetilde\kappa_{10}^{\,2}}
\int d^{10}x\,\sqrt{-g}
\left\{
e^{-2\Phi}
\left[
R+4(\partial\Phi)^2-\frac12|H_3|^2
\right]
-\frac12|F_2|^2
-\frac12|F_4|^2
\right\}.
\] Thus the two classes of fields have different dilaton behavior: \[
\begin{aligned}
\text{NS--NS fields }(g_{\mu\nu},\Phi,B_2):
&\qquad e^{-2\Phi},\\
\text{R--R fields }(C_1,C_3):
&\qquad \text{no string-frame dilaton prefactor}.
\end{aligned}
\] The ten-dimensional metric \(g_{\mu\nu}\) in (18.105) is specifically
the string-frame metric. Choosing another ten-dimensional frame would
change the first exponent in the eleven-dimensional ansatz.

If the asymptotic coupling has already been extracted from the
gravitational constant, write \[
\Phi=\Phi_0+\phi,
\qquad
g_s=e^{\Phi_0}.
\] Then the same local statement becomes \[
\frac{e^{-2\Phi}}{\widetilde\kappa_{10}^{\,2}}
=\frac{e^{-2\phi}}{\kappa_{10}^{\,2}},
\qquad
\kappa_{10}=g_s\widetilde\kappa_{10}.
\] This constant convention does not change any of the following local
reduction formulas.

\subsubsection{Why the exponents are
fixed}\label{why-the-exponents-are-fixed}

Begin with the more general Kaluza--Klein ansatz \[
ds_{11}^2
=e^{2a\Phi}g_{\mu\nu}dx^\mu dx^\nu
+e^{2b\Phi}
\left(dx^{11}+C_\mu dx^\mu\right)^2.
\] There are two unknown constants, \(a\) and \(b\).

The determinant scales as \[
\sqrt{-G_{11}}
=e^{(10a+b)\Phi}\sqrt{-g}.
\] The part of the eleven-dimensional Ricci scalar containing the
ten-dimensional Ricci scalar scales as \[
R_{11}
\supset
e^{-2a\Phi}R_{10}.
\] Hence the reduced Einstein--Hilbert density contains \[
\sqrt{-G_{11}}R_{11}
\supset
\sqrt{-g}\,
e^{(8a+b)\Phi}R_{10}.
\] To reproduce the string-frame factor \(e^{-2\Phi}\), one must impose
\[
\boxed{8a+b=-2}.
\]

The off-diagonal metric component is the Kaluza--Klein gauge field: \[
C_\mu
\equiv
\frac{G_{\mu\,11}}{G_{11\,11}}.
\] In type IIA it is the R--R one-form, so its kinetic term must not
carry the NS--NS factor \(e^{-2\Phi}\). The higher-dimensional curvature
contains \[
R_{11}
\supset
-\frac14e^{(2b-4a)\Phi}F_{\mu\nu}F^{\mu\nu},
\] where the indices on \(F_2\) are raised with \(g_{\mu\nu}\). After
multiplying by the determinant, its dilaton factor is \[
e^{(10a+b)\Phi}e^{(2b-4a)\Phi}
=e^{(6a+3b)\Phi}.
\] Requiring no dilaton prefactor gives \[
6a+3b=0
\quad\Longleftrightarrow\quad
\boxed{2a+b=0}.
\]

Solving the two equations, \[
8a+b=-2,
\qquad
2a+b=0,
\] gives \[
\boxed{
a=-\frac13,
\qquad
b=\frac23
}.
\] Therefore \[
e^{2a\Phi}=e^{-2\Phi/3},
\qquad
e^{2b\Phi}=e^{4\Phi/3},
\] which proves the exponents in equation (18.105).

This argument also shows that the exponents are not chosen merely to
make the determinant simple. They simultaneously arrange the
string-frame NS--NS prefactor and the R--R normalization of the
Kaluza--Klein vector.

\subsubsection{Determinant of the full Kaluza--Klein
metric}\label{determinant-of-the-full-kaluzaklein-metric}

For the actual values \(a=-1/3\), \(b=2/3\), the metric components are
\[
\begin{aligned}
G_{\mu\nu}
&=e^{-2\Phi/3}g_{\mu\nu}
+e^{4\Phi/3}C_\mu C_\nu,\\
G_{\mu\,11}
&=e^{4\Phi/3}C_\mu,\\
G_{11\,11}
&=e^{4\Phi/3}.
\end{aligned}
\] Using the Schur complement, \[
\begin{aligned}
\det G_{11}
&=G_{11\,11}
\det\left(
G_{\mu\nu}
-\frac{G_{\mu\,11}G_{\nu\,11}}{G_{11\,11}}
\right)\\
&=e^{4\Phi/3}
\det\left(e^{-2\Phi/3}g_{\mu\nu}\right)\\
&=e^{4\Phi/3}e^{-20\Phi/3}\det g\\
&=e^{-16\Phi/3}\det g.
\end{aligned}
\] Thus \[
\boxed{
\sqrt{-G_{11}}
=e^{-8\Phi/3}\sqrt{-g}
}.
\] The cancellation of the \(C_\mu C_\nu\) terms is why the
Kaluza--Klein vector does not occur in the determinant.

\subsubsection{Explicit reduction of the Ricci
scalar}\label{explicit-reduction-of-the-ricci-scalar}

The curvature calculation is clearest in two steps. Define \[
h_{\mu\nu}
=e^{-2\Phi/3}g_{\mu\nu},
\qquad
\sigma=\frac{2\Phi}{3},
\qquad
\eta=dx^{11}+C_1.
\] Then \[
ds_{11}^2
=h_{\mu\nu}dx^\mu dx^\nu
+e^{2\sigma}\eta^2.
\]

\paragraph{Step 1: ordinary circle
reduction}\label{step-1-ordinary-circle-reduction}

For a one-dimensional fiber with all fields independent of \(x^{11}\),
direct evaluation of the Christoffel symbols, or equivalently the Cartan
structure equations, gives \[
\boxed{
R_{11}
=R[h]
-\frac14e^{2\sigma}F_{\mu\nu}F^{\mu\nu}_{(h)}
-2\Box_h\sigma
-2(\partial\sigma)^2_h
}.
\] The four terms have a transparent origin:

\begin{itemize}
\tightlist
\item
  \(R[h]\) is the curvature of the ten-dimensional base;
\item
  \(F_2=dC_1\) measures the twisting of the circle over the base;
\item
  the last two terms arise because the local circle radius \(e^\sigma\)
  varies over the base.
\end{itemize}

For \(C_1=0\), this reduces to the familiar scalar-warped-product
formula \[
R_{11}
=R[h]-2e^{-\sigma}\Box_h e^\sigma.
\]

\paragraph{\texorpdfstring{Step 2: Weyl-transform from \(h_{\mu\nu}\) to
the string
metric}{Step 2: Weyl-transform from h\_\{\textbackslash mu\textbackslash nu\} to the string metric}}\label{step-2-weyl-transform-from-h_munu-to-the-string-metric}

Write \[
h_{\mu\nu}=e^{2\omega}g_{\mu\nu},
\qquad
\omega=-\frac{\Phi}{3}.
\] In \(d=10\), the Weyl transformation of the Ricci scalar is \[
R[h]
=e^{-2\omega}
\left[
R[g]
-18\Box\omega
-72(\partial\omega)^2
\right].
\] Substituting \(\omega=-\Phi/3\) gives \[
\boxed{
R[h]
=e^{2\Phi/3}
\left[
R[g]
+6\Box\Phi
-8(\partial\Phi)^2
\right]
}.
\]

For a scalar \(f\), \[
\Box_hf
=e^{-2\omega}
\left[
\Box f+8\,\partial_\mu\omega\,\partial^\mu f
\right].
\] Therefore, with \(\sigma=2\Phi/3\), \[
\boxed{
\Box_h\sigma
=e^{2\Phi/3}
\left[
\frac23\Box\Phi
-\frac{16}{9}(\partial\Phi)^2
\right]
},
\] and \[
\boxed{
(\partial\sigma)^2_h
=e^{2\Phi/3}
\frac49(\partial\Phi)^2
}.
\]

Finally, because a two-form contains two inverse metrics, \[
F_{\mu\nu}F^{\mu\nu}_{(h)}
=e^{4\Phi/3}F_{\mu\nu}F^{\mu\nu}_{(g)}.
\]

\paragraph{Step 3: combine all terms}\label{step-3-combine-all-terms}

Substitution into the circle-reduction formula gives \[
\begin{aligned}
R_{11}
=e^{2\Phi/3}\Bigg[
&R+6\Box\Phi-8(\partial\Phi)^2\\
&-\frac43\Box\Phi
+\frac{32}{9}(\partial\Phi)^2
-\frac89(\partial\Phi)^2\\
&-\frac14e^{2\Phi}F_2^2
\Bigg].
\end{aligned}
\] Combining coefficients, \[
6-\frac43=\frac{14}{3},
\] and \[
-8+\frac{32}{9}-\frac89
=-\frac{16}{3}.
\] Hence \[
\boxed{
R_{11}
=e^{2\Phi/3}
\left[
R
+\frac{14}{3}\Box\Phi
-\frac{16}{3}(\partial\Phi)^2
-\frac14e^{2\Phi}F_2^2
\right]
}.
\] This is the Ricci-scalar factor missing from the determinant-only
argument.

\subsubsection{Combining the determinant and Ricci
scalar}\label{combining-the-determinant-and-ricci-scalar}

Multiplying by \[
\sqrt{-G_{11}}
=e^{-8\Phi/3}\sqrt{-g}
\] gives \[
\begin{aligned}
\sqrt{-G_{11}}R_{11}
=\sqrt{-g}\Bigg\{
&e^{-2\Phi}
\left[
R+\frac{14}{3}\Box\Phi-\frac{16}{3}(\partial\Phi)^2
\right]\\
&-\frac14F_2^2
\Bigg\}.
\end{aligned}
\] The \(F_2\) term has no dilaton factor, exactly as required for an
R--R field in string frame.

Now discard a boundary term and integrate by parts: \[
\begin{aligned}
\int d^{10}x\,\sqrt{-g}\,
e^{-2\Phi}\Box\Phi
&=
-\int d^{10}x\,\sqrt{-g}\,
\partial_\mu(e^{-2\Phi})\partial^\mu\Phi\\
&=
2\int d^{10}x\,\sqrt{-g}\,
e^{-2\Phi}(\partial\Phi)^2.
\end{aligned}
\] It follows that \[
\frac{14}{3}\Box\Phi
-\frac{16}{3}(\partial\Phi)^2
\quad\xrightarrow{\text{inside the action}}\quad
\left(
\frac{28}{3}-\frac{16}{3}
\right)(\partial\Phi)^2
=4(\partial\Phi)^2.
\] Thus, suppressing the constant circle integral for the moment, the
eleven-dimensional Einstein--Hilbert term reduces to \[
\boxed{
\int d^{11}x\,\sqrt{-G}\,R_{11}
\longrightarrow
\int dx^{11}
\int d^{10}x\,\sqrt{-g}
\left[
e^{-2\Phi}
\left(R+4(\partial\Phi)^2\right)
-\frac14F_2^2
\right]
}
\] This formula isolates the local field dependence. The constant
normalization of \(x^{11}\), its asymptotic proper radius, and the
gravitational couplings are treated separately below.

\subsubsection{Reduction of the eleven-dimensional
three-form}\label{reduction-of-the-eleven-dimensional-three-form}

The same exponents also give the correct NS--NS/R--R split of the
three-form. Equation (18.105) writes \[
A_3=C_3+B_2\wedge dx^{11},
\] or, in explicitly Kaluza--Klein-invariant form, \[
dA_3
=F_4+H_3\wedge\eta,
\qquad
\eta=dx^{11}+C_1,
\] with \[
H_3=dB_2,
\qquad
F_4=dC_3-H_3\wedge C_1.
\]

For a four-form entirely along the ten-dimensional base, four inverse
base metrics give \[
|F_4|_{11}^2
=e^{8\Phi/3}|F_4|_{10}^2.
\] The determinant cancels this factor: \[
\sqrt{-G}\,|F_4|_{11}^2
=\sqrt{-g}\,|F_4|_{10}^2.
\] Thus \(F_4\) is R--R and has no string-frame dilaton prefactor.

For \(H_3\wedge\eta\), the three inverse base metrics contribute
\(e^{2\Phi}\), while the inverse fiber metric contributes
\(e^{-4\Phi/3}\). Hence \[
|H_3\wedge\eta|_{11}^2
=e^{2\Phi/3}|H_3|_{10}^2,
\] and therefore \[
\sqrt{-G}\,|H_3\wedge\eta|_{11}^2
=\sqrt{-g}\,e^{-2\Phi}|H_3|_{10}^2.
\] Thus \(H_3\) is NS--NS and acquires precisely the same \(e^{-2\Phi}\)
factor as \(R\) and the dilaton kinetic term.

Combining the metric and form reductions gives \[
\boxed{
\begin{aligned}
S_{10}
=\frac1{2\widetilde\kappa_{10}^{\,2}}
\int d^{10}x\sqrt{-g}
\Bigg\{
&e^{-2\Phi}
\left[
R+4(\partial\Phi)^2-\frac12|H_3|^2
\right]\\
&-\frac12|F_2|^2
-\frac12|F_4|^2
\Bigg\}
+S_{\mathrm{CS}},
\end{aligned}
}
\] which is the bosonic type-IIA string-frame action.

\subsubsection{What the gravitational-coupling relation does and does
not
imply}\label{what-the-gravitational-coupling-relation-does-and-does-not-imply}

Equation (16.134a) uses the \(g_s\)-independent coupling
\(\widetilde\kappa_{10}\) together with the full dilaton \(\Phi\).
Equation (18.106), by contrast, uses the physical coupling \[
\kappa_{10}=g_s\widetilde\kappa_{10}
\] and the asymptotically normalized fluctuation \(\phi=\Phi-\Phi_0\).
Indeed, \[
\frac{e^{-2\Phi}}{\widetilde\kappa_{10}^{\,2}}
=\frac{e^{-2\phi}}{\kappa_{10}^{\,2}}.
\]

After this constant normalization has been made, integrating a field
configuration independent of \(x^{11}\) over the circle produces the
overall factor \[
\frac{2\pi R_{11}}{2\kappa_{11}^2}.
\] Equation (18.106), \[
\boxed{
\kappa_{11}^2
=2\pi R_{11}\kappa_{10}^2
},
\] then gives \[
\frac{2\pi R_{11}}{2\kappa_{11}^2}
=\frac1{2\kappa_{10}^2}.
\] This matches the overall normalization of the integrated
ten-dimensional action.

It does not imply \[
\sqrt{-G}\,R_{11}
=\sqrt{-g}\,R_{10}.
\] That would incorrectly omit the local dilaton, its derivatives, and
the Kaluza--Klein field strength. Written in the full-dilaton convention
of equation (16.134a), the actual local relation, modulo a total
derivative, is \[
\boxed{
\sqrt{-G}\,R_{11}
\doteq
\sqrt{-g}
\left[
e^{-2\Phi}
\left(R+4(\partial\Phi)^2\right)
-\frac14F_2^2
\right]
},
\] where \(\doteq\) denotes equality inside the action after discarding
the boundary term. In the physical-coupling convention used in (18.106),
replace \(\Phi\) in the NS--NS factor by the normalized fluctuation
\(\phi=\Phi-\Phi_0\).

The factor \[
2\pi g_s\sqrt{\alpha'}
=2\pi R_{11}
\] in your calculation is therefore responsible only for the relation
between \(\kappa_{11}\) and \(\kappa_{10}\). It cannot be used to equate
the two curvature densities pointwise.

\subsubsection{The same reduction in ten-dimensional Einstein
frame}\label{the-same-reduction-in-ten-dimensional-einstein-frame}

The relation between string and ordinary Einstein metrics is \[
g_{\mu\nu}^{(E)}
=e^{-\Phi/2}g_{\mu\nu}^{(S)}.
\] Substituting \[
g_{\mu\nu}^{(S)}
=e^{\Phi/2}g_{\mu\nu}^{(E)}
\] into (18.105) gives \[
\boxed{
ds_{11}^2
=e^{-\Phi/6}ds_{10,E}^2
+e^{4\Phi/3}(dx^{11}+C_1)^2
}.
\] In this form, the condition that the ten-dimensional
Einstein--Hilbert term have no dilaton prefactor is manifest: \[
8\left(-\frac1{12}\right)+\frac23=0.
\] The reduced metric sector becomes \[
\boxed{
\sqrt{-G}R_{11}
\doteq
\sqrt{-g_E}
\left[
R_E
-\frac12(\partial\Phi)^2
-\frac14e^{3\Phi/2}F_2^2
\right]
}.
\] Thus equation (18.105) can equivalently be understood as the standard
Einstein-frame Kaluza--Klein ansatz, rewritten using
\(g_S=e^{\Phi/2}g_E\).

\subsubsection{Geometrical meaning of the fiber
exponent}\label{geometrical-meaning-of-the-fiber-exponent}

For a constant asymptotic dilaton \(\Phi_0=\log g_s\) and a
dimensionless coordinate \(x^{11}\) of period \(2\pi\), the physical
eleven-dimensional line element is \[
ds_{\mathrm{phys}}^2
=\ell_{11}^2
\left[
g_s^{-2/3}g_{\mu\nu}dx^\mu dx^\nu
+g_s^{4/3}(dx^{11})^2
\right].
\] Consequently, \[
\alpha'
=\ell_{11}^2g_s^{-2/3},
\qquad
R_{11}
=\ell_{11}g_s^{2/3}.
\] These imply \[
\boxed{
\ell_{11}
=g_s^{1/3}\sqrt{\alpha'},
\qquad
R_{11}
=g_s\sqrt{\alpha'}
}.
\] Thus the factor \(e^{4\Phi/3}\) also has an immediate geometrical
interpretation: its square root, \(e^{2\Phi/3}\), measures the M-theory
circle radius in eleven-dimensional Planck units, \[
\boxed{
\frac{R_{11}}{\ell_{11}}
=g_s^{2/3}
}.
\]

\subsubsection{Checks}\label{checks-34}

If \(\Phi\) is constant and \(C_1=0\), then \[
R_{11}
=e^{2\Phi/3}R_{10},
\] so \[
\sqrt{-G}R_{11}
=\sqrt{-g}\,e^{-2\Phi}R_{10}.
\] This already checks the principal Weyl factor without any derivative
terms.

If \(\Phi\) varies but \(C_1=0\), the integration-by-parts identity
gives \[
\frac{14}{3}\Box\Phi-\frac{16}{3}(\partial\Phi)^2
\ \longrightarrow\
4(\partial\Phi)^2,
\] which is the characteristic string-frame dilaton kinetic coefficient.

If \(\Phi\) is constant but \(F_2\neq0\), then \[
\sqrt{-G}R_{11}
\supset
-\frac14\sqrt{-g}\,F_2^2,
\] with no dilaton prefactor. This confirms that the Kaluza--Klein
vector is the type-IIA R--R one-form rather than an NS--NS gauge field.

\subsubsection{Result}\label{result-40}

For \[
ds_{11}^2
=e^{2a\Phi}ds_{10,S}^2
+e^{2b\Phi}(dx^{11}+C_1)^2,
\] string-frame reduction requires \[
\boxed{
8a+b=-2,
\qquad
2a+b=0
},
\] and hence \[
\boxed{
a=-\frac13,
\qquad
b=\frac23
}.
\] Therefore \[
\boxed{
ds_{11}^2
=e^{-2\Phi/3}ds_{10,S}^2
+e^{4\Phi/3}(dx^{11}+C_1)^2
}.
\]

The determinant and curvature are \[
\boxed{
\sqrt{-G}
=e^{-8\Phi/3}\sqrt{-g}
},
\] and \[
\boxed{
R_{11}
=e^{2\Phi/3}
\left[
R
+\frac{14}{3}\Box\Phi
-\frac{16}{3}(\partial\Phi)^2
-\frac14e^{2\Phi}F_2^2
\right]
}.
\] Thus, modulo a total derivative, \[
\boxed{
\sqrt{-G}R_{11}
\doteq
\sqrt{-g}
\left[
e^{-2\Phi}
\left(R+4(\partial\Phi)^2\right)
-\frac14F_2^2
\right]
}.
\] Finally, \[
\boxed{
\kappa_{11}^2
=2\pi R_{11}\kappa_{10}^2
}
\] matches the overall coefficients after integration over the circle;
it does not assert equality of the eleven- and ten-dimensional local
curvature densities.

\section{Q6: Deriving the M2 two-form charge from the Green--Schwarz
action}\label{q6-deriving-the-m2-two-form-charge-from-the-greenschwarz-action}

In page 720, the M-theory susy algebra is given by \[
\{ Q_{\alpha},Q_{\beta} \}=-(\Gamma^\mu C^{-1})_{\alpha\beta}P_{\mu}-\frac{1}{2}(\Gamma^{\mu \nu}C^{-1})_{\alpha\beta}Z_{\mu \nu}
\] It is said that this can be derived from the world-volume action of
the extended objects in the Green-Schwarz formulation. How to derive
this? Show the derivations.

\begin{quote}
\textbf{Textbook location and related material.} Chapter 18, Sections
18.4 and 18.7: the M-theory supersymmetry algebra, M2/M5 central
charges, and Green--Schwarz branes. The explicit Wess--Zumino descent
derivation supplements the book.
\end{quote}

\subsection{Answer}\label{answer-41}

\subsubsection{Core point}\label{core-point-41}

The statement can be made precise as follows. Start with the covariant
Green--Schwarz action of a single M2-brane in flat eleven-dimensional
superspace. Its Nambu--Goto term is \textbf{strictly} invariant under
rigid spacetime supersymmetry. Its Wess--Zumino term is invariant only
\textbf{up to a total derivative}. Noether's theorem therefore gives a
supercharge containing an improvement term inherited from that total
derivative. When the canonical Dirac bracket of two improved
supercharges is evaluated,

\begin{itemize}
\tightlist
\item
  the ordinary part gives the total momentum \(P_M\);
\item
  the Wess--Zumino improvement gives \[
  Z^{MN}=T_{\mathrm{M2}}
  \int_{\Sigma_2}dX^M\wedge dX^N.
  \]
\end{itemize}

This is the two-form charge in equation (18.117), and the last
expression is equation (18.118) with \(p=2\). Thus the two-form term is
not inserted by hand: it is the Noether two-cocycle caused by the
quasi-invariance of the Green--Schwarz Wess--Zumino term.

There is one point of scope to make explicit. The algebra on page 720
actually also permits the five-form term \[
-\frac1{5!}
(\Gamma^{M_1\cdots M_5}C^{-1})_{\alpha\beta}
Z_{M_1\cdots M_5}.
\] The formula in the question is the momentum-plus-M2 truncation. I
will derive the M2 term explicitly. The M5 term follows from the same
Noether mechanism applied to the Green--Schwarz/PST M5 action, with an
additional contribution involving its chiral world-volume two-form. This
distinction agrees with the direct M2 discussion in
\href{https://arxiv.org/abs/hep-th/9712004}{Townsend's superalgebra
lectures} and the direct M5 Noether calculation of
\href{https://arxiv.org/abs/hep-th/9708003}{Sorokin and Townsend}.

\subsubsection{The flat-superspace Green--Schwarz M2
action}\label{the-flat-superspace-greenschwarz-m2-action}

Let \[
\xi^i=(\tau,\sigma^1,\sigma^2)
\] be world-volume coordinates. The embedding fields are \[
X^M(\xi),
\qquad
\Theta^\alpha(\xi),
\qquad
M=0,\ldots,10,
\qquad
\alpha=1,\ldots,32.
\] The fundamental supersymmetry-invariant one-form is \[
\boxed{
\Pi^M=dX^M-i\bar\Theta\Gamma^M d\Theta
},
\qquad
\Pi_i^M=\partial_iX^M-i\bar\Theta\Gamma^M\partial_i\Theta.
\] The induced metric is \[
g_{ij}=\Pi_i^M\Pi_j^N\eta_{MN}.
\] The Green--Schwarz action is \[
\boxed{
S_{\mathrm{M2}}
=-T_{\mathrm{M2}}
\int_{\mathcal M_3}d^3\xi\,\sqrt{-\det g}
+T_{\mathrm{M2}}
\int_{\mathcal M_3}B_3
}.
\] Here \(B_3\) is a super-three-form potential. Its exterior derivative
is the invariant closed super-four-form \[
H_4=dB_3
=-i\,\bar{d\Theta}\Gamma_{MN}d\Theta
\wedge\Pi^M\wedge\Pi^N,
\] up to an overall sign fixed by the orientation chosen for the
M2-brane. Closure, \[
dH_4=0,
\] uses the eleven-dimensional Fierz identity which symmetrically
contracts one factor of \(C\Gamma^M\) with one factor of
\(C\Gamma_{MN}\). This closed four-form is the superspace cocycle
underlying the membrane Wess--Zumino coupling.

The action also has local \(\kappa\)-symmetry. That gauge symmetry
removes half of the components of \(\Theta\), but it is not the origin
of the extra term in the spacetime algebra. The extra term comes from
rigid supersymmetry and the Wess--Zumino coupling. The role of
\(\kappa\)-symmetry will reappear in the BPS check.

\subsubsection{Rigid supersymmetry and strict versus
quasi-invariance}\label{rigid-supersymmetry-and-strict-versus-quasi-invariance}

Rigid eleven-dimensional supersymmetry acts as \[
\delta_\varepsilon\Theta=\varepsilon,
\qquad
\delta_\varepsilon X^M
=i\bar\varepsilon\Gamma^M\Theta,
\qquad
d\varepsilon=0.
\] It follows directly that \[
\begin{aligned}
\delta_\varepsilon\Pi^M
&=d(\delta_\varepsilon X^M)
-i\,\delta_\varepsilon\bar\Theta\Gamma^M d\Theta
-i\bar\Theta\Gamma^M d(\delta_\varepsilon\Theta)\\
&=i\bar\varepsilon\Gamma^M d\Theta
-i\bar\varepsilon\Gamma^M d\Theta
-0\\
&=0.
\end{aligned}
\] Therefore \[
\delta_\varepsilon g_{ij}=0,
\qquad
\delta_\varepsilon
\left(\sqrt{-\det g}\right)=0.
\] The Nambu--Goto part of the action is exactly invariant.

Because \(H_4\) is constructed only from the invariant forms \(\Pi^M\)
and \(d\Theta\), it is also strictly invariant: \[
\delta_\varepsilon H_4=0.
\] But this does \textbf{not} imply that a chosen potential \(B_3\) is
strictly invariant. Since \(H_4=dB_3\), \[
d(\delta_\varepsilon B_3)
=\delta_\varepsilon H_4
=0.
\] Locally, the Poincare lemma gives \[
\boxed{
\delta_\varepsilon B_3=d\Lambda_2(\varepsilon)
}.
\] Consequently, \[
\delta_\varepsilon S_{\mathrm{M2}}
=T_{\mathrm{M2}}
\int_{\mathcal M_3}d\Lambda_2(\varepsilon)
=T_{\mathrm{M2}}
\int_{\partial\mathcal M_3}\Lambda_2(\varepsilon).
\] For a closed membrane, or for boundary conditions that kill the last
integral, the action is supersymmetric. Nevertheless, the Lagrangian
density has changed by a divergence, and that divergence must be
retained in the Noether current.

At the lowest nontrivial order in \(\Theta\), the descent two-form has
the structure \[
\Lambda_2(\varepsilon)
=-\frac{i}{2}
\bar\varepsilon\Gamma_{MN}\Theta\,
dX^M\wedge dX^N
+O(\Theta^3).
\] The higher-order terms replace the \(dX^M\) by the appropriate
supersymmetric combinations and are needed for the exact current. They
do not change the bosonic topological charge that survives in the final
algebra.

\subsubsection{Noether's theorem for a quasi-invariant
Lagrangian}\label{noethers-theorem-for-a-quasi-invariant-lagrangian}

For fields \(\phi^A\) with a variation \[
\delta L=\partial_iK^i,
\] the Noether current is not merely the canonical momentum times
\(\delta\phi^A\). It is \[
\boxed{
j^i
=\frac{\partial L}{\partial(\partial_i\phi^A)}
\delta\phi^A-K^i
}.
\] Indeed, using the Euler--Lagrange equations, \[
\begin{aligned}
\delta L
&=\frac{\partial L}{\partial\phi^A}\delta\phi^A
+\frac{\partial L}{\partial(\partial_i\phi^A)}
\partial_i\delta\phi^A\\
&=\left[
\frac{\partial L}{\partial\phi^A}
-\partial_i\frac{\partial L}{\partial(\partial_i\phi^A)}
\right]\delta\phi^A
+\partial_i\left[
\frac{\partial L}{\partial(\partial_i\phi^A)}
\delta\phi^A
\right],
\end{aligned}
\] and hence on shell \[
\partial_i j^i=0.
\]

For the membrane supersymmetry, write \[
j^i(\varepsilon)=\bar\varepsilon j^i,
\qquad
Q(\varepsilon)=\int_{\Sigma_2}d^2\sigma\,j^0(\varepsilon)
=\bar\varepsilon Q.
\] If \(P_M(\sigma)\) and \(\mathscr P_\alpha(\sigma)\) denote the
canonical momenta conjugate to \(X^M\) and \(\Theta^\alpha\), the
supercharge has the following relevant terms: \[
\boxed{
Q_\alpha
=\int_{\Sigma_2}d^2\sigma
\left[
\mathscr P_\alpha
-iP_M(C\Gamma^M\Theta)_\alpha
-\frac{iT_{\mathrm{M2}}}{2}
\epsilon^{rs}\partial_rX^M\partial_sX^N
(C\Gamma_{MN}\Theta)_\alpha
+O(\Theta^3)
\right]
}.
\] The first two terms are the expected generator of \[
\delta\Theta=\varepsilon,
\qquad
\delta X^M=i\bar\varepsilon\Gamma^M\Theta.
\] The third term is the crucial Wess--Zumino improvement \(-K^0\). If
the total derivative in \(\delta B_3\) were discarded, this term would
be missed and one would incorrectly recover only the unextended
super-Poincare algebra.

The displayed formula keeps precisely the terms needed to expose the
bosonic momentum and membrane charge. The omitted \(O(\Theta^3)\) terms
are not optional in the exact supercharge: together with the fermionic
constraints they supersymmetrically complete the result. Their brackets
either cancel fermion-dependent pieces, combine into world-volume
diffeomorphism and \(\kappa\)-symmetry constraints, or vanish on the
physical constraint surface.

\subsubsection{Canonical bracket of two
supercharges}\label{canonical-bracket-of-two-supercharges}

The fermionic canonical relation is schematically \[
\{\Theta^\alpha(\sigma),
\mathscr P_\beta(\sigma')\}_{D}
=\delta^\alpha{}_{\beta}\,
\delta^{(2)}(\sigma-\sigma').
\] Strictly speaking, the Green--Schwarz fermionic constraints contain
both first- and second-class parts, so the physical calculation uses a
Dirac bracket. For the two field-independent bosonic terms we seek, it
is enough to display the cross-brackets of \(\mathscr P_\alpha\) with
the terms linear in \(\Theta\).

Split the charge density as \[
Q_\alpha=Q^{(0)}_\alpha+Q^{(P)}_\alpha+Q^{(Z)}_\alpha+\cdots,
\] where \[
\begin{aligned}
Q^{(0)}_\alpha
&=\int d^2\sigma\,\mathscr P_\alpha,\\
Q^{(P)}_\alpha
&=-i\int d^2\sigma\,
P_M(C\Gamma^M\Theta)_\alpha,\\
Q^{(Z)}_\alpha
&=-\frac{iT_{\mathrm{M2}}}{2}
\int d^2\sigma\,
\epsilon^{rs}\partial_rX^M\partial_sX^N
(C\Gamma_{MN}\Theta)_\alpha.
\end{aligned}
\]

First consider the momentum term. Since \(C\Gamma^M\) is symmetric in
eleven dimensions, \[
\begin{aligned}
&\{Q^{(0)}_\alpha,Q^{(P)}_\beta\}_{D}
+\{Q^{(P)}_\alpha,Q^{(0)}_\beta\}_{D}\\
&\qquad
=-i(C\Gamma^M)_{\beta\alpha}
\int d^2\sigma\,P_M(\sigma)
-i(C\Gamma^M)_{\alpha\beta}
\int d^2\sigma\,P_M(\sigma)\\
&\qquad
=-2i(C\Gamma^M)_{\alpha\beta}P_M,
\end{aligned}
\] where \[
P_M=\int_{\Sigma_2}d^2\sigma\,P_M(\sigma).
\]

Now consider the Wess--Zumino improvement. In eleven dimensions
\(C\Gamma^{MN}\) is also symmetric on its spinor indices. Therefore \[
\begin{aligned}
&\{Q^{(0)}_\alpha,Q^{(Z)}_\beta\}_{D}
+\{Q^{(Z)}_\alpha,Q^{(0)}_\beta\}_{D}\\
&\qquad
=-iT_{\mathrm{M2}}
(C\Gamma_{MN})_{\alpha\beta}
\int_{\Sigma_2}d^2\sigma\,
\epsilon^{rs}\partial_rX^M\partial_sX^N.
\end{aligned}
\] Define \[
\boxed{
Z^{MN}
=T_{\mathrm{M2}}
\int_{\Sigma_2}d^2\sigma\,
\epsilon^{rs}\partial_rX^M\partial_sX^N
}.
\] Since \[
dX^M\wedge dX^N
=\epsilon^{rs}\partial_rX^M\partial_sX^N
\,d\sigma^1\wedge d\sigma^2,
\] this is equivalently \[
\boxed{
Z^{MN}
=T_{\mathrm{M2}}
\int_{\Sigma_2}dX^M\wedge dX^N
}.
\] Combining the two cross-brackets gives the physical part of the
classical charge algebra, \[
\boxed{
\{Q_\alpha,Q_\beta\}_{D}
=-2i\left[
(C\Gamma^M)_{\alpha\beta}P_M
+\frac12(C\Gamma^{MN})_{\alpha\beta}Z_{MN}
\right]
+\text{constraints}
}.
\] The factor \(1/2\) has the familiar tensor meaning: with unrestricted
sums over \(M,N\), every antisymmetric pair occurs twice, \[
\frac12\Gamma^{MN}Z_{MN}
=\sum_{M<N}\Gamma^{MN}Z_{MN}.
\] On the physical constraint surface, the constraint terms vanish.
Passing from the classical Grassmann Dirac bracket to the quantum
anticommutator, normalizing \(Q\) as in the book, and moving \(C\)
according to the book's placement of spinor indices gives \[
\boxed{
\{Q_\alpha,Q_\beta\}
=-(\Gamma^M C^{-1})_{\alpha\beta}P_M
-\frac12
(\Gamma^{MN}C^{-1})_{\alpha\beta}Z_{MN}
}.
\] Thus the overall minus signs in the question are convention
dependent, but the relative tensor structure, the factor \(1/2\), and
the topological charge are fixed.

\subsubsection{The same result from the Wess--Zumino descent
cocycle}\label{the-same-result-from-the-wesszumino-descent-cocycle}

The canonical computation has a useful conceptual reformulation. Because
\[
\delta_\varepsilon B_3=d\Lambda_2(\varepsilon),
\] two supersymmetry transformations produce the two-form descent
cocycle \[
\Omega_2(\varepsilon_1,\varepsilon_2)
=\delta_{\varepsilon_1}\Lambda_2(\varepsilon_2)
-\delta_{\varepsilon_2}\Lambda_2(\varepsilon_1)
-\Lambda_2([\varepsilon_1,\varepsilon_2]).
\] The commutator of two rigid supersymmetries is an ordinary
translation, and evaluating the field-independent part of this
expression gives \[
\Omega_2(\varepsilon_1,\varepsilon_2)
\propto
\bar\varepsilon_2\Gamma_{MN}\varepsilon_1
\,dX^M\wedge dX^N.
\] Therefore the algebra of the integrated Noether generators contains
\[
T_{\mathrm{M2}}
\bar\varepsilon_2\Gamma_{MN}\varepsilon_1
\int_{\Sigma_2}dX^M\wedge dX^N.
\] This is exactly the term found in the canonical bracket. It also
explains why the extension can be present even though the supersymmetry
transformations close locally on \(X^M\) and \(\Theta\): the extension
is carried by the boundary/improvement data of the action, not by an
additional local transformation of those fields.

\subsubsection{\texorpdfstring{Why \(Z^{MN}\) is topological and
conserved}{Why Z\^{}\{MN\} is topological and conserved}}\label{why-zmn-is-topological-and-conserved}

The integrand is closed identically: \[
d(dX^M\wedge dX^N)=0.
\] Suppose \(\Sigma_2\) and \(\Sigma_2'\) are spatial membrane slices
bounding a segment \(\mathcal W_3\) of the world-volume. Then \[
\begin{aligned}
Z^{MN}(\Sigma_2')-Z^{MN}(\Sigma_2)
&=T_{\mathrm{M2}}
\int_{\partial\mathcal W_3}dX^M\wedge dX^N\\
&=T_{\mathrm{M2}}
\int_{\mathcal W_3}d(dX^M\wedge dX^N)\\
&=0.
\end{aligned}
\] This is conservation without using the membrane equations of motion.
Equivalently, the world-volume current \[
J^{iMN}
=T_{\mathrm{M2}}\epsilon^{ijk}
\partial_jX^M\partial_kX^N
\] obeys \[
\partial_iJ^{iMN}=0
\] identically, because second derivatives are symmetric whereas
\(\epsilon^{ijk}\) is antisymmetric.

There is an important global qualification. If \(\Sigma_2\) is closed
and the functions \(X^M\) are globally single-valued in noncompact
target-space directions, then \[
dX^M\wedge dX^N=d(X^M dX^N)
\] is exact and its integral vanishes. A nonzero finite charge requires
a nontrivial wrapped two-cycle, winding in compact directions, suitable
asymptotic conditions for an infinite membrane, or an open membrane with
boundary data.

For example, take a rectangular target two-torus and \[
\begin{aligned}
X^1&=R_1(m_1\sigma^1+n_1\sigma^2),\\
X^2&=R_2(m_2\sigma^1+n_2\sigma^2),
\qquad
0\leq\sigma^r<2\pi.
\end{aligned}
\] Then \[
\int_{\Sigma_2}dX^1\wedge dX^2
=(2\pi)^2R_1R_2(m_1n_2-m_2n_1),
\] and hence \[
Z^{12}
=T_{\mathrm{M2}}(2\pi)^2R_1R_2
(m_1n_2-m_2n_1).
\] The integer determinant is the oriented wrapping number. Continuous
deformations of the embedding cannot change it.

In the conventions of the book, \[
T_{\mathrm{M2}}=\frac{2\pi}{\ell_{11}^3},
\] so the world-volume result is precisely \[
\boxed{
Z^{MN}
=\frac{2\pi}{\ell_{11}^3}
\int_{\Sigma_2}dX^M\wedge dX^N
},
\] the \(p=2\) case of equation (18.118).

\subsubsection{Short check: a static planar M2-brane and its
projector}\label{short-check-a-static-planar-m2-brane-and-its-projector}

Put a static M2-brane along the \(x^1,x^2\) directions and choose static
gauge, \[
X^0=\tau,
\qquad
X^1=\sigma^1,
\qquad
X^2=\sigma^2.
\] Per unit spatial area, \[
P_0=T_{\mathrm{M2}},
\qquad
Z_{12}=\pm T_{\mathrm{M2}},
\] where the sign is the membrane orientation. With \(C=\Gamma^0\) in a
Majorana basis, the algebra becomes, up to the common positive
normalization, \[
\{Q,Q\}
=T_{\mathrm{M2}}
\left(1\mp\Gamma^{012}\right).
\] Since \[
(\Gamma^{012})^2=1,
\qquad
\operatorname{tr}\Gamma^{012}=0,
\] the two eigenspaces of \(\Gamma^{012}\) each have dimension 16. The
unbroken supersymmetries satisfy \[
\boxed{
\Gamma^{012}\varepsilon=\pm\varepsilon
}.
\] This is exactly the \(\kappa\)-symmetry projector of a static
M2-brane. Positivity of the supercharge anticommutator gives \[
P_0\geq |Z_{12}|,
\] and the fundamental M2 saturates the bound because its charge density
equals its tension. This reproduces equation (18.119) and checks both
the interpretation and the normalization of the two-form charge.

\subsubsection{Extension to the
M5-brane}\label{extension-to-the-m5-brane}

For the M5-brane one repeats the same logical steps with its
six-dimensional Green--Schwarz/PST action. The Wess--Zumino sector
contains the pullbacks of the six-form potential dual to the
eleven-dimensional three-form and terms involving the chiral
world-volume three-form field strength. Its supersymmetry variation is
again an exact form. The corresponding Noether improvement produces \[
Z^{M_1\cdots M_5}
=T_{\mathrm{M5}}
\int_{\Sigma_5}
dX^{M_1}\wedge\cdots\wedge dX^{M_5},
\] which enters as \[
-\frac1{5!}
(\Gamma^{M_1\cdots M_5}C^{-1})_{\alpha\beta}
Z_{M_1\cdots M_5}.
\] The M5 world-volume chiral two-form also allows dissolved M2 charge,
so the complete M5 Noether algebra can contain both the two-form and
five-form charges. The simple statement ``M2 gives \(Z_2\) and M5 gives
\(Z_5\)'' is correct for pure branes, but the full M5 world-volume
algebra is richer.

\subsubsection{Result}\label{result-41}

The Green--Schwarz M2 action and its rigid supersymmetry are \[
S_{\mathrm{M2}}
=-T_{\mathrm{M2}}\int\sqrt{-\det g}
+T_{\mathrm{M2}}\int B_3,
\qquad
\delta\Theta=\varepsilon,
\qquad
\delta X^M=i\bar\varepsilon\Gamma^M\Theta.
\] The invariant pullback satisfies \[
\Pi^M=dX^M-i\bar\Theta\Gamma^M d\Theta,
\qquad
\delta\Pi^M=0,
\] while the Wess--Zumino potential is only quasi-invariant: \[
\delta B_3=d\Lambda_2(\varepsilon).
\] The resulting Noether supercharge contains the improvement \[
Q_\alpha\supset
-\frac{iT_{\mathrm{M2}}}{2}
\int_{\Sigma_2}d^2\sigma\,
\epsilon^{rs}\partial_rX^M\partial_sX^N
(C\Gamma_{MN}\Theta)_\alpha.
\] Its canonical Dirac bracket gives \[
\boxed{
Z^{MN}
=T_{\mathrm{M2}}
\int_{\Sigma_2}dX^M\wedge dX^N
=\frac{2\pi}{\ell_{11}^3}
\int_{\Sigma_2}dX^M\wedge dX^N
}.
\] After imposing the world-volume constraints and translating to the
book's convention, \[
\boxed{
\{Q_\alpha,Q_\beta\}
=-(\Gamma^M C^{-1})_{\alpha\beta}P_M
-\frac12
(\Gamma^{MN}C^{-1})_{\alpha\beta}Z_{MN}
}.
\] Thus the momentum term comes from the strictly invariant kinetic part
of the Green--Schwarz action, whereas the two-form extension comes from
the exact-form supersymmetry variation of the Wess--Zumino term.

\section{Q7: Total seven-brane monodromy, tadpole cancellation, and the
number
18}\label{q7-total-seven-brane-monodromy-tadpole-cancellation-and-the-number-18}

In page 727 it is said that: \textgreater At most 18 of the 24 7-branes
can be D7-branes, while the remaining ones are more general \((p,q)\)
7-branes. An easy way to see that not all 24 branes can be D7's is that
the total monodromy must be trivial, as a cycle enclosing all 24 branes
can be contracted to a point. This is nothing but the tadpole condition
for the seven-branes with a compact transverse space.

The statements are too brief. Why the requirement of trivial total
monodromy is the tadpole cancellation condition? Why this implies there
are at most 18 D7-branes?

\begin{quote}
\textbf{Textbook location and related material.} Chapter 18, Section
18.8 (pp.~724--733): seven-brane monodromies, elliptic K3, F-theory
tadpoles, and the moduli/rank count associated with the number 18.
\end{quote}

\subsection{Answer}\label{answer-42}

\subsubsection{Core point}\label{core-point-42}

There are two logically distinct statements in the quoted paragraph:

\begin{enumerate}
\def\labelenumi{\arabic{enumi}.}
\item
  \textbf{Global monodromy/tadpole condition.} Because the transverse
  space is the compact sphere \(B=\mathbb P^1\), the ordered product of
  the monodromies around all seven-branes must be the identity: \[
  \boxed{
  K_{\mathrm{tot}}
  =K_{24}K_{23}\cdots K_1
  =\mathbf1
  }.
  \] This is the exact, non-Abelian \(SL(2,\mathbb Z)\) version of the
  statement that the integrated seven-brane charge on a compact space
  must vanish.
\item
  \textbf{The number 18.} The product condition immediately proves that
  the 24 branes cannot all be D7-branes, because their product would be
  \(T^{24}\neq\mathbf1\). It does \textbf{not}, by itself, prove the
  sharper inequality ``number of D7-branes \(\leq18\).'' The number 18
  comes from the dimension of the elliptic-K3 complex-structure moduli
  space and equivalently from the rank of the eight-dimensional
  vector-multiplet sector.
\end{enumerate}

Accordingly, the phrase ``at most 18 D7-branes'' should be read as ``at
most 18 independent mutually local brane gauge fields/relative-position
multiplets can be put in one D7 duality frame.'' If it is read as a
literal bound on the number of D7 constituents counted with
multiplicity, it is too strong. I will derive each statement and return
to this qualification explicitly.

\subsubsection{Seven-branes are monodromy
defects}\label{seven-branes-are-monodromy-defects}

The type-IIB axio-dilaton is \[
\tau=C_0+ie^{-\Phi}.
\] A D7-brane is magnetically charged under the axion \(C_0\). Near a
D7-brane at \(z=z_i\), the local solution has the logarithmic form \[
\tau(z)
=\tau_{\mathrm{reg}}(z)
+\frac{1}{2\pi i}\log(z-z_i).
\] On taking \(z-z_i\) once counterclockwise around the origin, \[
\log(z-z_i)\longmapsto\log(z-z_i)+2\pi i,
\] and therefore \[
\tau\longmapsto\tau+1.
\] This is the \(SL(2,\mathbb Z)\) transformation \[
\boxed{
T=
\begin{pmatrix}
1&1\\
0&1
\end{pmatrix}
},
\qquad
K_{\mathrm{D7}}=T.
\]

A general \((p,q)\) seven-brane is obtained by conjugating a D7-brane by
a duality transformation. In the convention of equation (18.89), its
monodromy is \[
\boxed{
K_{[p,q]}
=M(p,q)
=
\begin{pmatrix}
1+pq&p^2\\
-q^2&1-pq
\end{pmatrix}
}.
\] Although every individual \(K_{[p,q]}\) is conjugate to \(T\),
different such matrices generally do not commute. Consequently, ``total
seven-brane charge'' is not globally an ordinary sum of integers. It is
encoded by an \textbf{ordered product} of monodromies.

\subsubsection{\texorpdfstring{Why the total monodromy on
\(\mathbb P^1\) is
trivial}{Why the total monodromy on \textbackslash mathbb P\^{}1 is trivial}}\label{why-the-total-monodromy-on-mathbb-p1-is-trivial}

Remove the 24 seven-brane positions from the base: \[
B^\circ
=\mathbb P^1\setminus\{z_1,\ldots,z_{24}\}.
\] Choose a common base point \(z_*\) and loops \(\gamma_i\) based at
\(z_*\), with \(\gamma_i\) going once around \(z_i\). The fundamental
group of a sphere with 24 punctures has the presentation \[
\pi_1(B^\circ,z_*)
=
\left\langle
\gamma_1,\ldots,\gamma_{24}
\ \middle|\
\gamma_1\gamma_2\cdots\gamma_{24}=1
\right\rangle.
\] Geometrically, a loop enclosing every puncture can be moved to the
opposite side of the sphere, where it encloses no puncture, and then
contracted.

The one-cycles \((a,b)\) of the elliptic fiber form an integral rank-two
local system over \(B^\circ\). Parallel transport defines a monodromy
representation \[
\rho:\pi_1(B^\circ,z_*)\longrightarrow SL(2,\mathbb Z),
\qquad
\rho(\gamma_i)=K_i.
\] Applying \(\rho\) to the fundamental-group relation gives \[
\rho(\gamma_1\cdots\gamma_{24})
=\rho(1)
=\mathbf1.
\] Depending on whether monodromies act on row or column vectors and on
the convention for composing paths, this is written as either \[
K_1K_2\cdots K_{24}=\mathbf1
\] or \[
K_{24}\cdots K_2K_1=\mathbf1.
\] The invariant content is \[
\boxed{K_{\mathrm{tot}}=\mathbf1}.
\]

This condition says that after transporting an arbitrary fiber cycle
around all singular fibers, it must return to itself. Otherwise the
elliptic fibration, and equivalently the type-IIB fields patched by
\(SL(2,\mathbb Z)\), would not be globally well defined.

\subsubsection{Why this is the seven-brane tadpole
condition}\label{why-this-is-the-seven-brane-tadpole-condition}

First consider a perturbative region in which all sources are mutually
local and a single axion field \(C_0\) can be used globally away from
them. Locally, \[
F_1=dC_0.
\] A D7-brane is a magnetic source, so the Bianchi identity takes the
schematic form \[
dF_1
=2\pi\sum_i Q_i\,
\delta^{(2)}(z-z_i),
\] where \(Q_i=+1\) for a D7-brane and negative sources such as
orientifold planes must be included when present.

Integrating over a compact transverse surface \(B\) gives \[
\int_BdF_1
=2\pi\sum_iQ_i.
\] Because \(B\) has no boundary, \[
\int_BdF_1
=\int_{\partial B}F_1
=0,
\] and hence \[
\boxed{\sum_iQ_i=0}.
\] This is precisely tadpole cancellation: the zero-momentum source for
the field dual to \(C_0\) must vanish. Equivalently, the sourced field
equation has a solution on a compact space only if its integrated source
is zero. It is the same elementary obstruction as the impossibility of
solving \[
\nabla^2\phi=\sum_iQ_i\delta^{(2)}(z-z_i)
\] on a compact manifold unless \(\sum_iQ_i=0\), because integrating the
left-hand side gives zero.

For mutually local D7-type sources, a source of charge \(Q_i\) has
monodromy \[
K_i=T^{Q_i}.
\] Since these matrices commute, \[
K_{\mathrm{tot}}
=\prod_iT^{Q_i}
=T^{\sum_iQ_i}.
\] Thus \[
\sum_iQ_i=0
\qquad\Longleftrightarrow\qquad
K_{\mathrm{tot}}=\mathbf1.
\] This explicitly proves the equivalence between the integrated Bianchi
identity and trivial total monodromy in the perturbative, Abelian case.

In a genuine F-theory background, different \((p,q)\) branes are
mutually nonlocal. There is no global duality frame in which all their
charges are integers sourcing the same \(C_0\), and the matrices \(K_i\)
do not commute. The additive equation \(\sum_iQ_i=0\) must then be
replaced by its exact group-valued form, \[
\boxed{
K_{24}\cdots K_1=\mathbf1
}.
\] This is why the book calls trivial total monodromy ``the tadpole
condition'': it is the nonperturbative \(SL(2,\mathbb Z)\)
generalization of the integrated seven-brane Bianchi identity.

One can also see the global condition geometrically. In an elliptic
Calabi--Yau fibration, \[
f\in H^0(B,K_B^{-4}),
\qquad
g\in H^0(B,K_B^{-6}),
\] so \[
\Delta=4f^3+27g^2
\in H^0(B,K_B^{-12}).
\] For \(B=\mathbb P^1\), \[
K_B^{-12}\cong\mathcal O(24),
\] and therefore \[
\deg\Delta=24.
\] This explains why the compact solution has 24 elementary discriminant
zeros. However, ``24 discriminant zeros'' does not mean ``net D7 charge
\(+24\)'': the zeros can carry different, noncommuting \((p,q)\)
monodromies whose ordered product is trivial.

\subsubsection{Why 24 ordinary D7-branes are
impossible}\label{why-24-ordinary-d7-branes-are-impossible}

If all 24 branes were D7-branes in one fixed duality frame, every
monodromy would equal \(T\). Their total monodromy would be \[
K_{\mathrm{tot}}
=T^{24}
=
\begin{pmatrix}
1&24\\
0&1
\end{pmatrix},
\] which is not the identity. Therefore \[
\boxed{
\text{the 24 elementary seven-branes cannot all be D7-branes
in one global duality frame}
}.
\] Some of the branes must have noncommuting monodromies conjugate to
\(T\) so that the full ordered product can become \(\mathbf1\).

Notice exactly what has been established: \[
K_{\mathrm{tot}}=\mathbf1
\quad\Longrightarrow\quad
N_{\mathrm{D7}}\neq24.
\] No inequality \(N_{\mathrm{D7}}\leq18\) follows from this one
calculation. The integer 18 requires additional global information.

\subsubsection{Where the number 18 actually comes
from}\label{where-the-number-18-actually-comes-from}

The elliptic K3 is written in Weierstrass form as \[
y^2=x^3+f(u)x+g(u),
\] with \[
f\in H^0(\mathbb P^1,\mathcal O(8)),
\qquad
g\in H^0(\mathbb P^1,\mathcal O(12)).
\] A degree-eight homogeneous polynomial has nine coefficients, and a
degree-twelve homogeneous polynomial has thirteen. Thus \((f,g)\)
initially contain \[
9+13=22
\] complex parameters.

They are not all physical:

\begin{itemize}
\tightlist
\item
  \(PGL(2,\mathbb C)\) reparametrizations of the base \(\mathbb P^1\)
  remove three complex parameters;
\item
  the weighted rescaling of the Weierstrass coordinates, equivalently
  the common rescaling \(f\mapsto w^2f\), \(g\mapsto w^3g\) used by the
  book, removes one more.
\end{itemize}

Therefore the number of complex-structure moduli is \[
\boxed{
22-3-1=18
}.
\] The 24 roots of \(\Delta\) are consequently not 24 freely specifiable
points. Their positions obey the relations required for them to arise
from the same pair \((f,g)\). Only 18 complex combinations are
independent.

In eight dimensions a vector multiplet with sixteen supercharges
contains one complex scalar. In a local D7 description this scalar is
the transverse position of the brane. Hence the 18 independent complex
deformations are paired with 18 independent vector multiplets: \[
\boxed{
n_V=18,
\qquad
G_{\mathrm{generic,\,brane}}=U(1)^{18}
}.
\] The eight-dimensional supergravity multiplet contains two additional
graviphotons, so the total generic Abelian rank is \[
18+2=20.
\] This matches the dual heterotic string on \(T^2\), whose
vector-multiplet moduli are organized by the Narain lattice
\(\Gamma^{2,18}\).

The naive 24 world-volume \(U(1)\) fields are therefore not 24
independent eight-dimensional gauge fields. Global \(SL(2,\mathbb Z)\)
identifications, couplings to bulk form fields, and Stückelberg
relations remove or identify six combinations. This is the physical
content behind the book's number 18.

This moduli-and-spectrum argument is the one given in the original
construction of F-theory on K3; see
\href{https://arxiv.org/abs/hep-th/9602022}{Vafa, \emph{Evidence for
F-Theory}}.

\subsubsection{The precise meaning of ``at most 18
D7-branes''}\label{the-precise-meaning-of-at-most-18-d7-branes}

The safest precise statement is \[
\boxed{
\text{the independent seven-brane vector-multiplet rank is at most }18
}.
\] Equivalently, at most 18 independent mutually local brane degrees of
freedom can be described simultaneously as ordinary D7 degrees of
freedom in one perturbative duality frame. At enhanced points these
\(U(1)\) fields become the Cartan generators of a non-Abelian algebra,
whose rank is still at most 18.

This is not the same as a literal bound on the number of mutually local
D7 constituents. For example, an elliptic K3 can have a Kodaira
\(I_{19}\) fiber. Its monodromy is \[
K_{I_{19}}=T^{19},
\] and in the seven-brane description it is a stack of 19 mutually local
D7-type constituents. The associated algebra is \[
A_{18}\cong\mathfrak{su}(19),
\] whose rank is 18. Such elliptic K3 surfaces are known explicitly; see
the classification result on
\href{https://arxiv.org/abs/math/0509540}{maximal singular fibers of
elliptic K3 surfaces}.

Thus:

\begin{itemize}
\tightlist
\item
  the monodromy condition rules out \textbf{24} D7-branes in one frame;
\item
  the moduli/spectrum calculation gives at most \textbf{18 independent
  vector multiplets or gauge rank};
\item
  a coincident stack may contain 19 D7 constituents while still
  contributing only rank 18.
\end{itemize}

The book's wording is therefore a useful shorthand in the discussion of
the generic Abelian spectrum, but it should not be interpreted as a
universal constituent-counting theorem.

\subsubsection{Short check: the perturbative orientifold
limit}\label{short-check-the-perturbative-orientifold-limit}

The weak-coupling orientifold limit gives a concrete Abelian version of
the same cancellation. F-theory on K3 becomes type IIB on \[
T^2/\mathbb Z_2
\] with four \(O7^-\)-planes and sixteen D7-branes. In D7 charge units,
\[
Q_{\mathrm{D7}}=+1,
\qquad
Q_{O7^-}=-4,
\] so \[
16Q_{\mathrm{D7}}+4Q_{O7^-}
=16-16
=0.
\] The corresponding monodromies may be written \[
K_{\mathrm{D7}}=T,
\qquad
K_{O7^-}=-T^{-4}.
\] Since four factors of \(-\mathbf1\) multiply to \(\mathbf1\), \[
T^{16}(-T^{-4})^4
=T^{16}T^{-16}
=\mathbf1.
\] Nonperturbatively, each \(O7^-\) splits into a pair of mutually
nonlocal seven-branes. Thus the F-theory discriminant count is \[
16\ \text{D7 constituents}
+4\times2\ \text{nonlocal constituents}
=24,
\] while the total monodromy and the D7 tadpole both cancel.

This limit also shows why ``24 seven-branes'' and ``24 D7-branes'' are
different statements. The first counts elementary discriminant zeros;
the second assumes that every zero can be placed in the same
perturbative duality frame, which global monodromy forbids.

\subsubsection{Result}\label{result-42}

On the punctured sphere, \[
\pi_1\!\left(
\mathbb P^1\setminus\{z_i\}
\right)
=
\left\langle
\gamma_1,\ldots,\gamma_{24}
\ \middle|\
\gamma_1\cdots\gamma_{24}=1
\right\rangle.
\] Therefore the monodromy representation obeys \[
\boxed{
K_{\mathrm{tot}}
=K_{24}\cdots K_1
=\mathbf1
}.
\] For mutually local sources this reduces to the integrated axion
Bianchi identity, \[
dF_1=2\pi\sum_iQ_i\delta_i,
\qquad
0=\int_{\mathbb P^1}dF_1
=2\pi\sum_iQ_i,
\] so trivial total monodromy is precisely the non-Abelian form of
seven-brane tadpole cancellation.

For 24 D7-branes, \[
K_{\mathrm{tot}}
=T^{24}
=
\begin{pmatrix}
1&24\\
0&1
\end{pmatrix}
\neq\mathbf1,
\] so the 24 branes cannot all be D7-branes in one duality frame.

The number of independent elliptic-K3 deformations is instead \[
\boxed{
(9+13)-3-1=18
},
\] giving 18 vector multiplets and maximal independent brane gauge rank
18. Hence the precise conclusion is \[
\boxed{\begin{gathered}
\text{at most 18 independent D7-type gauge/position multiplets,}\\
\text{not necessarily at most 18 D7 constituents}
\end{gathered}}.
\]

\section{\texorpdfstring{Q8: Four orientifold bulk vectors and heterotic
\(U(1)^{20}\)}{Q8: Four orientifold bulk vectors and heterotic U(1)\^{}\{20\}}}\label{q8-four-orientifold-bulk-vectors-and-heterotic-u120}

In page 731-732, they considered compactifying the type IIB string on a
\(T^2\) and taking the orientifold quotient
\(\Omega I_{2}(-1)^{F_{L}}\). This leads to the gauge group \(SO(8)^4\)
since locally there are four O7-planes and on each O7-plane there are
four dynamical D7-branes. While they said that: \textgreater In the
orientifold description, the remaining four Abelian gauge bosons known
to be present in F-theory\ldots\ldots{}

Why there are four remaining Abelian gauge bosons in the orientifold
description?

In the next page they also discussed heterotic string on \(T^2\) and
they also mentioned \textgreater the gauge group for generic
Wilson-lines of the heterotic string on \(T^2\) is also \(U(1)^{20}\)

How to derive this fact?

Make these two points clear. If they were studied somewhere in this
book, search and find the discussions.

\begin{quote}
\textbf{Textbook location and related material.} Chapter 18, Section
18.8: the \(SO(8)^4\) orientifold point and heterotic dual on \(T^2\).
Sections 10.2--10.4 provide the Narain lattice and toroidal vector
counting.
\end{quote}

\subsection{Answer}\label{answer-43}

\subsubsection{Core point}\label{core-point-43}

At the orientifold point, the complete eight-dimensional gauge group is
more precisely \[
\boxed{
SO(8)^4\times U(1)^4
}
\] up to possible global discrete quotients. The factors have different
origins:

\begin{itemize}
\tightlist
\item
  \(SO(8)^4\) comes from open strings on the four D7/O7 stacks;
\item
  \(U(1)^4\) comes from the closed-string fields \[
  B_{\mu a},
  \qquad
  C_{\mu a},
  \qquad
  a=1,2.
  \]
\end{itemize}

There are two vectors from the NS--NS two-form and two from the R--R
two-form. The combined orientifold projection keeps these components
even, although neither set would survive the geometric reflection in
isolation.

On the heterotic side, the generic gauge group follows from the equally
direct count \[
\boxed{
16\ \text{Cartan vectors}
+2\ \text{metric vectors}
+2\ \text{$B$-field vectors}
=20\ \text{vectors}
}.
\] Generic Wilson lines make all non-Cartan root gauge bosons massive,
leaving \[
\boxed{
G_{\mathrm{het,generic}}
=U(1)^{16}\times U(1)^2_{\mathrm{KK}}
\times U(1)^2_{\mathrm{wind}}
=U(1)^{20}
}.
\] In world-sheet language this is \[
U(1)^{18}_L\times U(1)^2_R.
\]

The book develops the ingredients earlier in Chapter 10:

\begin{itemize}
\tightlist
\item
  Section 10.2, especially equation (10.44), derives the \(2D\) vectors
  \(U(1)^D_L\times U(1)^D_R\) of a \(T^D\) compactification.
\item
  Section 10.4, equations (10.115)--(10.127), introduces heterotic
  Wilson lines and proves that the generic rank is \(16+2D\).
\item
  Section 10.6 and Table 10.3 identify the supersymmetric orientifold
  \(\Omega I_2(-1)^{F_L}\) and its \(SO(8)^4\) Chan--Paton group.
\end{itemize}

I will now derive both counts explicitly.

\subsubsection{Rank supplied by the four D7/O7
stacks}\label{rank-supplied-by-the-four-d7o7-stacks}

The involution \[
I_2:(y^1,y^2)\longmapsto(-y^1,-y^2)
\] has four fixed points on \(T^2\). After quotienting, these are the
four \(O7^-\)-planes on \[
T^2/\mathbb Z_2\simeq\mathbb P^1.
\] Local tadpole cancellation places four dynamical D7-branes on each
orientifold plane. The orientifold projection of the Chan--Paton factors
gives \[
SO(8)
\] at each fixed point, and therefore \[
G_{\mathrm{open}}=SO(8)^4.
\] Since \[
\operatorname{rank}SO(8)=4,
\] the open-string sector supplies \[
\operatorname{rank}G_{\mathrm{open}}
=4\times4
=16.
\] This is not yet the entire rank-20 gauge sector of F-theory on K3.
The missing four vectors are closed-string zero modes.

\subsubsection{Reduction of the type-IIB
two-forms}\label{reduction-of-the-type-iib-two-forms}

Split the ten-dimensional coordinates as \[
x^M=(x^\mu,y^a),
\qquad
\mu=0,\ldots,7,
\qquad
a=1,2.
\] Before imposing the orientifold projection, decompose the NS--NS
two-form: \[
\begin{aligned}
B_2
&=\frac12B_{\mu\nu}\,dx^\mu\wedge dx^\nu
+B_{\mu a}\,dx^\mu\wedge dy^a
+\frac12B_{ab}\,dy^a\wedge dy^b.
\end{aligned}
\] For each internal index \(a\), the component \[
B_{\mu a}(x)
\] has one noncompact index and is therefore an eight-dimensional vector
potential. Since \(a=1,2\), \(B_2\) would produce two vectors.

The R--R two-form decomposes in the same way: \[
\begin{aligned}
C_2
&=\frac12C_{\mu\nu}\,dx^\mu\wedge dx^\nu
+C_{\mu a}\,dx^\mu\wedge dy^a
+\frac12C_{ab}\,dy^a\wedge dy^b,
\end{aligned}
\] and its two components \[
C_{\mu1},
\qquad
C_{\mu2}
\] give two more eight-dimensional vector potentials.

Thus dimensional reduction gives the candidate set \[
\boxed{
B_{\mu1},\ B_{\mu2},\
C_{\mu1},\ C_{\mu2}
}.
\] It remains to verify that the orientifold projection retains them.

\subsubsection{Explicit orientifold parity
calculation}\label{explicit-orientifold-parity-calculation}

The orientifold generator is \[
\mathcal O=\Omega I_2(-1)^{F_L}.
\] The three factors act as follows.

First, world-sheet orientation reversal gives the standard type-IIB
parities \[
\Omega:
\qquad
B_2\longmapsto-B_2,
\qquad
C_2\longmapsto+C_2.
\] The first sign follows because the world-sheet coupling \[
\int_{\Sigma_2}B_2
\] changes sign when the world-sheet orientation is reversed. The second
sign is also familiar from the type-I quotient: the R--R two-form
\(C_2\) is retained by \(\Omega\).

Second, \[
(-1)^{F_L}
\] acts trivially on NS--NS fields and by a minus sign on R--R fields:
\[
(-1)^{F_L}:
\qquad
B_2\longmapsto+B_2,
\qquad
C_2\longmapsto-C_2.
\]

Third, under the reflection \[
I_2:y^a\longmapsto-y^a,
\qquad
dy^a\longmapsto-dy^a.
\] A tensor component with exactly one internal index therefore acquires
a minus sign: \[
I_2:
\qquad
B_{\mu a}\longmapsto-B_{\mu a},
\qquad
C_{\mu a}\longmapsto-C_{\mu a}.
\]

Multiplying the three signs gives

\begin{longtable}[]{@{}lrrrr@{}}
\toprule\noalign{}
component & \(\Omega\) & \((-1)^{F_L}\) & \(I_2\) & \(\mathcal O\) \\
\midrule\noalign{}
\endhead
\bottomrule\noalign{}
\endlastfoot
\(G_{\mu a}\) & \(+\) & \(+\) & \(-\) & \(-\) \\
\(B_{\mu a}\) & \(-\) & \(+\) & \(-\) & \(+\) \\
\(C_{\mu a}\) & \(+\) & \(-\) & \(-\) & \(+\) \\
\end{longtable}

Therefore \[
\boxed{
\mathcal O B_{\mu a}=+B_{\mu a},
\qquad
\mathcal O C_{\mu a}=+C_{\mu a}
},
\] while the metric Kaluza--Klein vectors \(G_{\mu a}\) are projected
out.

Hence exactly four bulk Abelian vectors survive: \[
\boxed{
2\ B_{\mu a}
+2\ C_{\mu a}
=4\ U(1)\ \text{gauge bosons}
}.
\]

There is a small geometric subtlety here. The quotient is
\(\mathbb P^1\), and \[
b_1(\mathbb P^1)=0,
\] so these vectors should not be interpreted as ordinary reduction on
harmonic one-forms of \(\mathbb P^1\). They are most transparently
obtained on the covering torus. The one-form \(dy^a\) is odd under
\(I_2\), but this geometric minus sign is canceled by the intrinsic
\(\Omega\) or \(\Omega(-1)^{F_L}\) parity of the corresponding two-form.
On the quotient, the fields are therefore modes valued in the
orientifold sign local system rather than ordinary one-form zero modes.

Combining open and closed sectors gives \[
\boxed{
G_{\mathrm{8d}}
=SO(8)^4\times U(1)^4,
\qquad
\operatorname{rank}G_{\mathrm{8d}}
=16+4
=20
}.
\]

Of these four Abelian fields, two belong to the eight-dimensional
supergravity multiplet and are the graviphotons; the other two belong to
vector multiplets. Which linear combinations of \(B_{\mu a}\) and
\(C_{\mu a}\) are assigned to these multiplets depends on the duality
frame and the scalar background, but the \(2+2\) multiplet split and the
total number four are invariant.

\subsubsection{Ten-dimensional heterotic gauge fields and generic Wilson
lines}\label{ten-dimensional-heterotic-gauge-fields-and-generic-wilson-lines}

The ten-dimensional heterotic gauge group is either \[
E_8\times E_8
\] or \[
\operatorname{Spin}(32)/\mathbb Z_2.
\] Both have rank 16. Let \(H_I\), \(I=1,\ldots,16\), denote Cartan
generators. Compactification on \(T^2\) allows two commuting Wilson
lines \[
W_a
=P\exp\left(i\oint_{\gamma_a}A\right),
\qquad
a=1,2.
\] Because \[
\pi_1(T^2)=\mathbb Z\oplus\mathbb Z
\] is Abelian, \(W_1\) and \(W_2\) must commute. They can therefore be
conjugated into the maximal torus: \[
W_a
=\exp\left(2\pi i\,A_a^IH_I\right).
\]

Decompose the ten-dimensional adjoint algebra into Cartan and root
generators, \[
\mathfrak g
=\mathfrak h
\oplus
\bigoplus_{\alpha\in\Delta_{\mathfrak g}}\mathfrak g_\alpha.
\] For a root generator \(E_\alpha\), \[
W_aE_\alpha W_a^{-1}
=e^{2\pi i\alpha\cdot A_a}E_\alpha.
\] The corresponding gauge boson is unbroken only if it is invariant
around both cycles: \[
\alpha\cdot A_a\in\mathbb Z,
\qquad
a=1,2.
\] For generic values of the Wilson lines, no nonzero root satisfies
both conditions. All root gauge bosons become massive, while the Cartan
generators always commute with \(W_a\). Therefore \[
\boxed{
E_8\times E_8
\ \text{or}\
\operatorname{Spin}(32)/\mathbb Z_2
\quad\xrightarrow{\text{generic }W_1,W_2}\quad
U(1)^{16}
}.
\]

In Kaluza--Klein language the mass shift can be seen schematically as \[
p_a
=\frac1{R_a}
\left(
m_a+\alpha\cdot A_a
\right),
\qquad
m_a\in\mathbb Z.
\] A ten-dimensional root gauge boson can remain massless only if
integers \(m_a\) can cancel both Wilson-line phases. At a generic point
this is impossible for every \(\alpha\neq0\); at special Wilson lines
some roots satisfy the condition and a non-Abelian group reappears.

\subsubsection{Four additional heterotic vectors from the
torus}\label{four-additional-heterotic-vectors-from-the-torus}

The ten-dimensional heterotic NS--NS sector contains the metric
\(G_{MN}\) and the two-form \(B_{MN}\). Reducing on \(T^2\) gives \[
G_{\mu a},
\qquad
B_{\mu a},
\qquad
a=1,2.
\]

The two \(G_{\mu a}\) are ordinary Kaluza--Klein gauge fields. An
internal diffeomorphism with an \(x\)-dependent translation parameter,
\[
y^a\longmapsto y^a+\xi^a(x),
\] acts on them as \[
\delta G_{\mu a}\sim\partial_\mu\xi_a,
\] which is an Abelian gauge transformation. They therefore generate \[
U(1)^2_{\mathrm{KK}}.
\] States carrying integer momentum \(m_a\) around \(T^2\) are
electrically charged under these vectors.

The two \(B_{\mu a}\) arise from the gauge symmetry \[
\delta B_2=d\Lambda_1.
\] Taking an internal component \(\Lambda_a(x)\) gives \[
\delta B_{\mu a}
=\partial_\mu\Lambda_a,
\] so they generate \[
U(1)^2_{\mathrm{wind}}.
\] String states with winding numbers \(n^a\) are charged under these
vectors.

Consequently the complete generic heterotic group is \[
\begin{aligned}
G_{\mathrm{het,generic}}
&=U(1)^{16}_{\mathrm{Cartan}}
\times U(1)^2_{\mathrm{KK}}
\times U(1)^2_{\mathrm{wind}}\\
&=\boxed{U(1)^{20}}.
\end{aligned}
\]

This field-theory reduction already proves the statement in the book.

\subsubsection{World-sheet and Narain-lattice
derivation}\label{world-sheet-and-narain-lattice-derivation}

The string derivation organizes the same fields more naturally into
left- and right-moving combinations. For a \(T^D\) compactification, the
metric and two-form vectors combine schematically as \[
A^a_{\mu,L}
\sim G_{\mu a}+B_{\mu a},
\qquad
A^a_{\mu,R}
\sim G_{\mu a}-B_{\mu a}.
\] Thus a closed string compactified on \(T^D\) has \[
U(1)^D_L\times U(1)^D_R.
\] This is the result derived from the massless vector states in
equation (10.44).

The heterotic left-moving sector has, in addition, sixteen chiral
internal bosons whose currents generate the rank-16 gauge lattice.
Therefore \[
\boxed{
G_{\mathrm{het,generic}}(T^D)
=U(1)^{16+D}_L\times U(1)^D_R
}.
\] For \(D=2\), \[
\boxed{
G_{\mathrm{het,generic}}(T^2)
=U(1)^{18}_L\times U(1)^2_R
=U(1)^{20}
}.
\]

The charges lie in the even self-dual Narain lattice \[
\Gamma^{16+D,D}.
\] For \(D=2\) this is \[
\Gamma^{18,2}.
\] At a generic point in the Narain moduli space, the only massless
vectors are the 20 Cartan currents just counted. Extra left-moving gauge
bosons appear when the lattice contains vectors satisfying \[
p_R=0,
\qquad
p_L^2=2.
\] Such vectors form a simply laced root system. Hence at special values
of the metric, \(B\)-field, and Wilson lines, part of \[
U(1)^{18}_L
\] is enhanced to an ADE group. The two right-moving \(U(1)^2_R\) are in
the supergravity multiplet and do not undergo this non-Abelian
enhancement.

The Narain moduli space quoted in equation (10.127) becomes, for
\(D=2\), \[
\mathcal M_{\mathrm{Narain}}
=
\frac{O(2,18)}
{O(2)\times O(18)}
\Big/
O(2,18;\mathbb Z).
\] This is the same \(O(2,18)\) structure that appears in the
eight-dimensional F-theory/heterotic duality.

\subsubsection{\texorpdfstring{Matching the \(SO(8)^4\) orientifold
point to the heterotic
theory}{Matching the SO(8)\^{}4 orientifold point to the heterotic theory}}\label{matching-the-so84-orientifold-point-to-the-heterotic-theory}

At the special orientifold point, the rank-20 gauge algebra is \[
\boxed{
\mathfrak{so}(8)^{\oplus4}
\oplus
\mathfrak u(1)^{\oplus4}
}.
\] The four \(\mathfrak{so}(8)\) factors have total Cartan rank \[
4\times4=16.
\] On the heterotic side, this means that special Wilson lines have made
the roots of \[
D_4^{\oplus4}
\] massless inside the left-moving Narain lattice. Sixteen left-moving
Cartan generators are thereby enhanced from \[
U(1)^{16}
\quad\text{to}\quad
SO(8)^4,
\] without changing the rank.

Four Abelian fields remain:

\begin{itemize}
\tightlist
\item
  two unenhanced left-moving \(U(1)\) fields in vector multiplets;
\item
  two right-moving \(U(1)\) graviphotons in the supergravity multiplet.
\end{itemize}

Thus \[
\operatorname{rank}
\left[
SO(8)^4\times U(1)^4
\right]
=16+4
=20,
\] exactly matching both the orientifold reduction and the heterotic
Narain count.

The duality chain quoted by the book, \[
\frac{\text{type IIB on }T^2}
{\Omega I_2(-1)^{F_L}}
\xleftrightarrow{\,T\,}
\text{type I on }T^2
\xleftrightarrow{\,S\,}
\operatorname{Spin}(32)/\mathbb Z_2
\text{ heterotic on }T^2
\xleftrightarrow{\,T\,}
E_8\times E_8
\text{ heterotic on }T^2,
\] requires these rank-20 spectra to agree. The four orientifold bulk
vectors \(B_{\mu a},C_{\mu a}\) are mapped by this sequence of T- and
S-dualities to the four toroidal heterotic vectors constructed from
\(G_{\mu a}\) and \(B_{\mu a}\).

\subsubsection{Checks}\label{checks-35}

For a circle, \(D=1\), the same heterotic formula gives \[
U(1)^{17}_L\times U(1)_R,
\] of total rank 18. This agrees with the general result \(16+2D\) in
equation (10.127).

For \(T^4\), \(D=4\), it gives \[
U(1)^{20}_L\times U(1)^4_R,
\] of total rank 24. This is exactly the group displayed in equations
(18.97) and (18.101) in the book's discussion of heterotic/type-IIA
duality in six dimensions.

Finally, at the \(SO(8)^4\) point, \[
\operatorname{rank}SO(8)^4=16
\] and the four bulk vectors restore the universal total rank \[
16+4=20.
\]

\subsubsection{Result}\label{result-43}

The orientifold projection retains \[
\begin{aligned}
B_{\mu a}:&\quad
(-1)_\Omega(+1)_{(-1)^{F_L}}(-1)_{I_2}=+1,\\
C_{\mu a}:&\quad
(+1)_\Omega(-1)_{(-1)^{F_L}}(-1)_{I_2}=+1,
\end{aligned}
\] for \(a=1,2\). Hence \[
\boxed{
G_{\mathrm{orientifold}}
=SO(8)^4\times U(1)^4,
\qquad
U(1)^4
\leftrightarrow
\{B_{\mu1},B_{\mu2},C_{\mu1},C_{\mu2}\}
}.
\]

For the heterotic string, \[
\boxed{
G_{10}
\xrightarrow{\text{generic Wilson lines}}
U(1)^{16}
}
\] and the torus supplies \[
\boxed{
G_{\mu a}:U(1)^2_{\mathrm{KK}},
\qquad
B_{\mu a}:U(1)^2_{\mathrm{wind}}
}.
\] Therefore \[
\boxed{
G_{\mathrm{het,generic}}(T^2)
=U(1)^{16+2+2}
=U(1)^{20}
=U(1)^{18}_L\times U(1)^2_R
}.
\] At special Narain moduli, lattice vectors with \[
p_R=0,
\qquad
p_L^2=2
\] become additional massless gauge bosons and enhance a subgroup of
\(U(1)^{18}_L\) to an ADE group, while the total rank remains 20.

\section{Q9: Localized matter at intersecting D-branes and F-theory
seven-branes}\label{q9-localized-matter-at-intersecting-d-branes-and-f-theory-seven-branes}

In page 732, they briefly mentioned lower dimensional compactifications
of F-theory: \textgreater The singularity enhances further where these
co-dimension one submanifolds, i.e.~the 7-branes, intersect. Similar to
intersecting D-branes, this is where additional matter fields are
localized.

Why intersecting D-branes lead to additional localized matter fields? Is
this discussed somewhere in this book?

\begin{quote}
\textbf{Textbook location and related material.} Chapter 18, Section
18.8 for F-theory localization, read with Chapter 17, Section 17.2 for
matter at intersecting D-branes and its multiplicity/chirality.
\end{quote}

\subsection{Answer}\label{answer-44}

\subsubsection{Core point}\label{core-point-44}

The reason is very concrete in perturbative string theory. Suppose two
D-brane stacks \(a\) and \(b\) intersect. Besides the \(aa\) and \(bb\)
open strings whose two endpoints lie on the same stack, string theory
necessarily contains an \(ab\) sector whose two endpoints obey different
D-brane boundary conditions: \[
X(\tau,0)\in\Pi_a,
\qquad
X(\tau,\pi)\in\Pi_b.
\] The center of mass of such a string can move freely only in
directions common to both branes. In the relative internal directions,
the boundary conditions force its zero mode to sit at \[
\Pi_a\cap\Pi_b.
\] Consequently, the \(ab\) state is a field on the intersection
world-volume, rather than on either full brane world-volume.

At the intersection, the shortest string connecting the two branes has
zero length. Its stretching contribution to the mass therefore vanishes.
The Ramond sector has a massless fermionic ground state; when the
intersection preserves supersymmetry, an appropriate Neveu--Schwarz
state is also massless and combines with it into a matter multiplet. If
the stacks contain \(N_a\) and \(N_b\) branes, the two endpoints give
Chan--Paton indices and the localized state transforms as \[
\boxed{
(\mathbf N_a,\overline{\mathbf N}_b)
\quad\text{under}\quad
U(N_a)\times U(N_b)
}.
\] Thus these are matter fields charged under two gauge factors, not
additional gauge bosons in either adjoint representation.

This is derived explicitly earlier in the book. The main discussion is
Section 10.6, especially equations (10.188)--(10.197), Figures
10.3--10.4, and Table 10.2. The book first proves that the
center-of-mass position is fixed at an intersection, then derives the
fractionally moded open string, its massless Ramond state, and its
bifundamental Chan--Paton representation. Section 15.3 gives an explicit
orientifold example in Table 15.10, and Section 17.2 and Table 17.2
generalize the multiplicity of charged chiral matter to topological
intersection numbers of D6-brane cycles.

\subsubsection{Same-brane strings give gauge fields; mixed strings give
matter}\label{same-brane-strings-give-gauge-fields-mixed-strings-give-matter}

Consider stacks \(a\) and \(b\) with world-volume gauge groups \[
U(N_a)\times U(N_b).
\] There are four oriented open-string sectors: \[
aa,\qquad bb,\qquad ab,\qquad ba.
\]

An \(aa\) string carries a Chan--Paton matrix \((\lambda_{aa})^i{}_j\)
with both indices belonging to stack \(a\). It transforms as \[
\lambda_{aa}\longmapsto
U_a\lambda_{aa}U_a^{-1},
\] so its massless vector state is in the adjoint of \(U(N_a)\). Because
both endpoints may move together anywhere along stack \(a\), this gauge
field propagates on the entire world-volume of that stack. The same
statement applies to the \(bb\) sector.

An \(ab\) string instead carries one index \(i=1,\ldots,N_a\) at its
first endpoint and one index \(\bar j=1,\ldots,N_b\) at its second
endpoint: \[
(\lambda_{ab})^i{}_{\bar j}.
\] It transforms as \[
\lambda_{ab}\longmapsto
U_a\lambda_{ab}U_b^{-1},
\] and therefore lies in \[
(\mathbf N_a,\overline{\mathbf N}_b).
\] Reversing the orientation gives the \(ba\) sector in \[
(\overline{\mathbf N}_a,\mathbf N_b).
\] These two orientations are related by CPT in the lower-dimensional
theory. They do not in general constitute two independent chiral
multiplets.

The representation is already the first indication that this state is
matter. A gauge boson of \(U(N_a)\) would have to lie in
\((\operatorname{adj}_{a},\mathbf1)\), whereas the \(ab\) state is
charged under both groups.

\subsubsection{Boundary conditions pin the string to the
intersection}\label{boundary-conditions-pin-the-string-to-the-intersection}

Now reproduce the local calculation in equations (10.188)--(10.191).
Take a two-dimensional plane with coordinates \((X,Y)\). Put brane \(a\)
along the \(X\)-axis and let brane \(b\) meet it at a relative angle
\(\Delta\phi\).

At the endpoint \(\sigma=0\) on brane \(a\), \(X\) is Neumann and \(Y\)
is Dirichlet: \[
\sigma=0:
\qquad
\partial_\sigma X=0,
\qquad
\partial_\tau Y=0.
\] At the endpoint \(\sigma=\pi\) on brane \(b\), the tangent coordinate
is Neumann and the normal coordinate is Dirichlet: \[
\sigma=\pi:
\qquad
\cos(\Delta\phi)\,\partial_\sigma X
+\sin(\Delta\phi)\,\partial_\sigma Y=0,
\] \[
-\sin(\Delta\phi)\,\partial_\tau X
+\cos(\Delta\phi)\,\partial_\tau Y=0.
\]

To isolate the center-of-mass zero mode, take the part constant in both
\(\sigma\) and \(\tau\): \[
X(\tau,\sigma)=x_0,
\qquad
Y(\tau,\sigma)=y_0.
\] The Dirichlet condition at the first endpoint says that the point
lies on brane \(a\): \[
y_0=0.
\] The Dirichlet condition at the second endpoint says that it lies on
brane \(b\): \[
-\sin(\Delta\phi)x_0
+\cos(\Delta\phi)y_0=0.
\] For a transverse intersection, \(\sin(\Delta\phi)\neq0\), and hence
\[
y_0=0,
\qquad
x_0=0.
\] The only allowed constant position is therefore the intersection
point. If the intersection is located at a point \(P\) rather than the
origin, the same calculation gives \[
(x_0,y_0)=P.
\]

This is the precise meaning of localization. It does not mean that the
state is localized in the noncompact spacetime. If the two D-branes
share \(d\) noncompact directions, their common Neumann boundary
conditions still give momentum zero modes \(k_\mu\) in those directions.
The state propagates as a \(d\)-dimensional field, but its internal
wavefunction is supported at the brane intersection.

\subsubsection{Fractional moding and the missing translational zero
mode}\label{fractional-moding-and-the-missing-translational-zero-mode}

Introduce the complex coordinate \[
Z=\frac1{\sqrt2}(X+iY)
\] and the fractional angle \[
\nu=\frac{\Delta\phi}{\pi}.
\] Solving the two different endpoint conditions gives the shifted
oscillator modes displayed in equation (10.191): \[
\begin{aligned}
Z(\tau,\sigma)
=i\sqrt{\frac{\alpha'}2}
\sum_{n\in\mathbb Z}
\Bigg[
&\frac{\alpha_{n+\nu}}{n+\nu}
e^{-i(n+\nu)(\tau-\sigma)}\\
&+\frac{\widetilde\alpha_{n-\nu}}{n-\nu}
e^{-i(n-\nu)(\tau+\sigma)}
\Bigg].
\end{aligned}
\] The important structural fact is that the modes are shifted by
\(\pm\nu\). For \(\nu\neq0\), there is no ordinary zero-frequency
oscillator and no term \[
z_0+2\alpha' p\tau
\] describing free translation in this plane. This is the
oscillator-language version of the elementary zero-mode calculation
above.

When \(\nu\to0\), the branes become parallel. The ordinary translational
zero mode reappears, and the string is no longer pinned to an isolated
intersection. Thus localization is not an extra assumption: it is
encoded directly in the open-string boundary conditions.

The world-sheet fermions acquire the corresponding shifted modes \[
r+\nu,
\qquad
r-\nu,
\] where \[
r\in\mathbb Z
\quad\text{in the R sector},
\qquad
r\in\mathbb Z+\frac12
\quad\text{in the NS sector}.
\] This is why the book describes the \(ab\) string as a ``twisted open
string'\,': the relative brane angle plays the same role as a twist
acting on a closed-string coordinate.

\subsubsection{Why some localized states become
massless}\label{why-some-localized-states-become-massless}

For two nonintersecting branes separated by a shortest distance \(L\),
the classical stretching energy is \[
E_{\mathrm{stretch}}
=T_{\mathrm F}L
=\frac{L}{2\pi\alpha'}.
\] It contributes \[
M^2_{\mathrm{stretch}}
=\left(\frac{L}{2\pi\alpha'}\right)^2
\] to the open-string mass. The schematic mass formula is therefore \[
M^2
=\left(\frac{L}{2\pi\alpha'}\right)^2
+\frac1{\alpha'}(N+a).
\] At an intersection, \[
L=0,
\] so the geometrical mass obstruction disappears.

For supersymmetric open strings at angles, the bosonic and fermionic
zero-point energies cancel in the Ramond sector: \[
a_{\mathrm R}=0.
\] The GSO-projected Ramond ground state is therefore massless. For
D6-branes wrapping three-cycles and sharing four-dimensional Minkowski
space, it gives a four-dimensional Weyl fermion. Reversing the
orientation supplies its CPT conjugate, and the orientation of the
internal intersection fixes the four-dimensional chirality.

The NS sector is angle dependent. For relative angles \(0\leq\nu_I<1\)
in three complex internal planes, its vacuum energy is \[
a_{\mathrm{NS}}
=-\frac12
+\frac12\sum_{I=1}^3\nu_I.
\] The four light scalar masses obtained from the lowest NS states can
be organized as \[
\begin{aligned}
\alpha'M_1^2&=\frac12(-\nu_1+\nu_2+\nu_3),\\
\alpha'M_2^2&=\frac12(\nu_1-\nu_2+\nu_3),\\
\alpha'M_3^2&=\frac12(\nu_1+\nu_2-\nu_3),\\
\alpha'M_4^2&=1-\frac12(\nu_1+\nu_2+\nu_3).
\end{aligned}
\] Different choices of oriented representatives for the angles permute
these expressions. The supersymmetry condition in the book, \[
\sum_{I=1}^3\Delta\phi^I=0
\pmod{2\pi},
\] is precisely the condition that one of these scalars be massless. It
then combines with the Ramond fermion into a four-dimensional
\(\mathcal N=1\) chiral multiplet. With more special angle relations,
Table 10.2 shows how the fields instead assemble into \(\mathcal N=2\)
hypermultiplets or \(\mathcal N=4\) vector multiplets.

Thus an intersection always creates a distinct localized open-string
sector. It does not, by itself, guarantee a healthy supersymmetric
matter multiplet: for nonsupersymmetric angles some NS states can be
massive or tachyonic. The statement on page 732 concerns supersymmetric
F-theory compactifications, where the relevant localized states organize
into supersymmetric matter multiplets.

\subsubsection{Intersection number gives the multiplicity and
chirality}\label{intersection-number-gives-the-multiplicity-and-chirality}

On a compact space, two wrapped branes may intersect several times. For
one \(T^2\), let their one-cycles have wrapping numbers \[
\Pi_a=n_a[a]+m_a[b],
\qquad
\Pi_b=n_b[a]+m_b[b].
\] Their oriented intersection number is equation (10.193): \[
\boxed{
I_{ab}
=\Pi_a\circ\Pi_b
=n_am_b-m_an_b
}.
\] For transverse intersections, \[
|I_{ab}|
\] counts the distinct intersection points, so there is one copy of the
localized open-string ground state at each point. The sign of \(I_{ab}\)
records the local orientation and therefore the chirality convention.

For factorized D6-branes on \[
T^6=T^2_1\times T^2_2\times T^2_3,
\] equation (10.195) gives \[
\boxed{
I_{ab}
=\prod_{I=1}^3
\left(n_a^Im_b^I-m_a^In_b^I\right)
}.
\] Hence the net chiral spectrum contains \[
I_{ab}
\] copies, with sign interpreted as chirality, in the bifundamental
representation \[
(\mathbf N_a,\overline{\mathbf N}_b).
\] Section 17.2 replaces this toroidal formula by the general
homological pairing \[
I_{ab}=\Pi_a\circ\Pi_b
\] for special Lagrangian three-cycles in a Calabi--Yau orientifold.
This is why charged chiral multiplicities are topological and remain
invariant under continuous deformations that do not cross a wall where
intersections are created or annihilated in oppositely oriented pairs.

\subsubsection{Translation into F-theory
geometry}\label{translation-into-f-theory-geometry}

Now return to the statement on page 732. In an F-theory
compactification, a seven-brane wraps a component \[
D_a\subset B
\] of the discriminant locus \[
\Delta=4f^3+27g^2=0.
\] Since \(D_a\) has complex codimension one in the base, gauge fields
associated with the singular elliptic fiber over \(D_a\) propagate on
the external spacetime together with the internal divisor \(D_a\).

At a generic point of \(D_a\), a collection of fibral curves collapses
according to the root lattice of a gauge algebra \(\mathfrak g_a\). If
another discriminant component \(D_b\) meets it, then over \[
\Sigma_{ab}=D_a\cap D_b
\] the fiber degenerates further. Additional fibral curves become
vanishing cycles only over this complex-codimension-two locus. This is
the geometric meaning of ``the singularity enhances further.'\,'

Locally, the enhanced algebra often organizes the matter representation
through a decomposition \[
\operatorname{adj}(G_{\mathrm{enh}})
\longrightarrow
(\operatorname{adj}G_a,\mathbf1)
\oplus
(\mathbf1,\operatorname{adj}G_b)
\oplus
\bigoplus_{\mathbf R}
\left(\mathbf R_a,\mathbf R_b\right).
\] The first two terms are the gauge fields already present along
\(D_a\) and \(D_b\). The off-diagonal terms are charged under both
groups and are interpreted as matter localized on \(\Sigma_{ab}\).

For example, a local enhancement \[
SU(N_a)\times SU(N_b)
\subset SU(N_a+N_b)
\] gives schematically \[
\begin{aligned}
\operatorname{adj}SU(N_a+N_b)
\longrightarrow{}&
(\operatorname{adj}_{a},\mathbf1)
\oplus
(\mathbf1,\operatorname{adj}_{b})
\oplus
(\mathbf1,\mathbf1)\\
&\oplus
(\mathbf N_a,\overline{\mathbf N}_b)
\oplus
(\overline{\mathbf N}_a,\mathbf N_b).
\end{aligned}
\] The last two terms are precisely the representations of the
perturbative \(ab\) and \(ba\) open strings.

\subsubsection{M-theory explanation of F-theory
localization}\label{m-theory-explanation-of-f-theory-localization}

The nonperturbative F-theory statement is especially precise in the dual
M-theory description. Resolve the singular elliptic fiber. Gauge bosons
arise from M2-branes wrapping fibral curves corresponding to roots. If
\(C_{\mathbf w}\) is an additional fibral curve associated with a matter
weight \(\mathbf w\), then the wrapped M2-brane has mass \[
M_{\mathbf w}
=T_{\mathrm{M2}}
\operatorname{Vol}(C_{\mathbf w}).
\] At a generic point of \(D_a\), this particular curve is absent as an
independent vanishing curve or has nonzero volume, so the corresponding
state is massive. Over the intersection locus \(\Sigma_{ab}\), the fiber
splits or enhances and \[
\operatorname{Vol}(C_{\mathbf w})\longrightarrow0.
\] The wrapped M2-brane therefore becomes massless only there. Since the
vanishing cycle exists only above \(\Sigma_{ab}\), its wavefunction can
propagate along the external spacetime and along \(\Sigma_{ab}\), but
not over the whole of \(D_a\) or \(D_b\). This is the F-theory version
of the missing open-string translational zero mode.

In a weakly coupled type-IIB limit, these wrapped-M2 states reduce to
ordinary strings stretched between D7-brane stacks. For genuinely
mutually nonlocal \((p,q)\) seven-branes, there need not be a globally
valid perturbative open-string description; the states are more properly
described by \((p,q)\) string junctions or by the vanishing fibral
curves. The localization principle remains the same.

\subsubsection{What ``localized matter'\,' means in different F-theory
dimensions}\label{what-localized-matter-means-in-different-f-theory-dimensions}

For F-theory on an elliptically fibered Calabi--Yau threefold with base
surface \(B_2\), the seven-branes wrap curves in \(B_2\). Two such
curves meet at points. The matter field is localized at those points
internally and propagates in the six noncompact dimensions.

For F-theory on an elliptically fibered Calabi--Yau fourfold with base
threefold \(B_3\), the seven-branes wrap surfaces. Two surfaces
intersect along a complex curve: \[
\Sigma_{ab}=D_a\cap D_b.
\] This is the familiar matter curve of four-dimensional F-theory
compactifications. The geometry determines which charged representations
can occur. A nontrivial \(G_4\) flux determines the net number of
four-dimensional chiral minus antichiral zero modes. This is why the
paragraph immediately following the quoted sentence adds that chiral
matter in the fourfold case requires a nontrivial \(F_4\) or \(G_4\)
background.

It is therefore useful to distinguish three statements:

\begin{enumerate}
\def\labelenumi{\arabic{enumi}.}
\tightlist
\item
  The codimension-two fiber enhancement produces matter representations
  localized on \(\Sigma_{ab}\).
\item
  The internal Dirac equation on \(\Sigma_{ab}\) determines the actual
  lower-dimensional zero modes.
\item
  Gauge flux, represented by \(G_4\) in F-theory, determines their net
  chirality in four dimensions.
\end{enumerate}

\subsubsection{Checks}\label{checks-36}

If the two branes are separated and do not intersect, then \(L>0\) and
the \(ab\) string has \[
M^2\supset
\left(\frac{L}{2\pi\alpha'}\right)^2>0.
\] There is no massless localized matter.

If the two branes are parallel and coincident, then \(\nu=0\) and the
translational zero mode returns. The mixed strings can propagate along
the full common brane world-volume and usually enlarge the adjoint gauge
algebra; this is gauge enhancement, rather than matter confined to an
isolated internal intersection.

If the branes intersect transversely at \(|I_{ab}|\) points, the local
boundary-condition problem is identical at each point. Hence the same
bifundamental state occurs \(|I_{ab}|\) times, exactly as equations
(10.193) and (10.195) require.

\subsubsection{Result}\label{result-44}

For an \(ab\) open string, \[
X(\tau,0)\in\Pi_a,
\qquad
X(\tau,\pi)\in\Pi_b
\] forces its internal zero mode to obey \[
X_0\in\Pi_a\cap\Pi_b.
\] At an intersection, the stretching length is zero: \[
M^2
=\left(\frac{L}{2\pi\alpha'}\right)^2
+\frac{N+a}{\alpha'},
\qquad
L=0.
\] The massless Ramond state, together with its NS partner when
supersymmetry is preserved, gives localized matter transforming as \[
\boxed{
(\mathbf N_a,\overline{\mathbf N}_b)
\quad\text{with net multiplicity}\quad
I_{ab}=\Pi_a\circ\Pi_b
}.
\]

For factorized cycles on \(T^6\), \[
\boxed{
I_{ab}
=\prod_{I=1}^3
\left(n_a^Im_b^I-m_a^In_b^I\right)
}.
\]

In F-theory, \[
\boxed{
D_a\cap D_b=\Sigma_{ab}
\quad\Longrightarrow\quad
\text{codimension-two fiber enhancement}
\quad\Longrightarrow\quad
\text{matter localized on }\Sigma_{ab}
}.
\] Equivalently, an additional fibral curve collapses only over
\(\Sigma_{ab}\), \[
M_{\mathbf w}
=T_{\mathrm{M2}}\operatorname{Vol}(C_{\mathbf w})
\longrightarrow0,
\] which is the nonperturbative geometric counterpart of a zero-length
open string at a D-brane intersection.

\section{\texorpdfstring{Q10: The regimes \(g_s\ll1\), \(g_sN\ll1\), and
\(g_sN\gg1\)}{Q10: The regimes g\_s\textbackslash ll1, g\_sN\textbackslash ll1, and g\_sN\textbackslash gg1}}\label{q10-the-regimes-g_sll1-g_snll1-and-g_sngg1}

In page 734: \textgreater The D-brane picture is valid at weak string
coupling, \(g_{s}\gg 1\), where the string is perturbative and open
strings end on D-branes which are heavy and themselves not in the
perturbative spectrum of the theory. This picture also assumes that the
back-reaction of N coincident branes on Minkowski space is small, which
is the case if \(g_{s}N\gg 1\). Otherwise their distortion of space-time
cannot be neglected. In that situation we have a description as
solutions of the SUGRA equations of motion with fluxes. This also
requires \(g_{s}\) small but \(g_{s}N\) large, because otherwise the
SUGRA description (18.26) which ignores string loop effects and
\(\alpha'\) corrections due to massive string modes is not valid and the
equations of motions (18.27) and therefore also their brane solutions
have to be corrected.

I am confused with these statements: 1. The limit \(g_{s}\gg 1\) is easy
to understand: we can consider the supergravity as an tree-level
effective field theory description without higher string loop
corrections. 2. How the second limit \(g_{s}N\gg 1\) relates to the
back-reaction of N coincident branes on the Minkowski space? 3. Why we
need \(g_{s}N\) large compared to \(g_{s}\) i.e.~\(N\gg 1\)? To ignore
\(\alpha'\) corrections we need to make sure that the curvature of the
SUGRA solution is small compared to \(1 /\alpha'\), how to obtain the
requirement for \(g_{s}N\)?

\begin{quote}
\textbf{Textbook location and related material.} Chapter 18, Section
18.9 around p.~734: open-string, supergravity, and large-\(N\) regimes
for coincident branes; the D3 near-horizon scale and suppression of
\(\alpha'\) corrections.
\end{quote}

\subsection{Answer}\label{answer-45}

\subsubsection{Core point}\label{core-point-45}

The apparent contradiction comes from two reversed inequalities in the
question. The printed text on page 734 says \[
\boxed{
g_s\ll1,
\qquad
g_sN\ll1
\quad
\text{for the perturbative D-brane picture in nearly flat space}
}
\] and \[
\boxed{
g_s\ll1,
\qquad
g_sN\gg1
\quad
\text{for the classical supergravity brane geometry}
}.
\] It does not say \(g_s\gg1\), and it does not say that back-reaction
is small for \(g_sN\gg1\).

The two parameters control different approximations:

\begin{itemize}
\tightlist
\item
  \(g_s\) controls closed-string loops. Tree-level string theory
  requires \(g_s\ll1\).
\item
  \(g_sN\) controls the collective strength of \(N\) D-branes and, for
  D3-branes, the curvature radius in string units.
\end{itemize}

For \(N\) D3-branes the exact relation quoted in equation (18.149) is \[
L^4=4\pi g_sN\alpha'^2.
\] The curvature scale is \(|\mathrm{Riemann}|^{1/2}\sim L^{-2}\), so \[
\alpha'|\mathrm{Riemann}|^{1/2}
\sim\frac{\alpha'}{L^2}
=\frac1{\sqrt{4\pi g_sN}}.
\] Therefore \[
\boxed{
\alpha'\text{ corrections small}
\quad\Longleftrightarrow\quad
L\gg\ell_s
\quad\Longleftrightarrow\quad
g_sN\gg1
}.
\]

Combining this with \(g_s\ll1\) is possible only if \(N\) is large. This
is the essential large-\(N\) window: \[
\boxed{
\frac1N\ll g_s\ll1
}
\] up to the numerical factors suppressed in the parametric
inequalities.

The book derives these facts in two places. Section 18.5, especially
pages 695--696, obtains the D-brane gravitational strength \(g_sN\)
directly from Newton's constant and the D\(p\)-brane tension. Section
18.9 then specializes to D3-branes in equations (18.149), (18.187), and
(18.188), relating \(g_sN\) to the radius \(L\) and to the 't Hooft
coupling.

\subsubsection{\texorpdfstring{Why tree level requires \(g_s\ll1\), not
\(g_s\gg1\)}{Why tree level requires g\_s\textbackslash ll1, not g\_s\textbackslash gg1}}\label{why-tree-level-requires-g_sll1-not-g_sgg1}

A connected closed-string world-sheet of genus \(h\) carries the dilaton
weight \[
g_s^{-\chi}
=g_s^{2h-2},
\qquad
\chi=2-2h.
\] Relative to the sphere contribution, each additional handle
contributes a factor \[
g_s^2.
\] Consequently, higher closed-string loops are suppressed when \[
\boxed{
g_s\ll1
}.
\]

If \(g_s\gg1\), the perturbative genus expansion is not controlled.
Type-IIB string theory may then admit an S-dual weakly coupled
description, but it is not the same weakly coupled D-brane expansion
being discussed on page 734.

It is also important that ``tree-level string theory'' and
``supergravity'' are not identical approximations. Classical
supergravity requires two independent suppressions: \[
\boxed{
\text{string loops:}\quad e^\Phi\ll1,
\qquad
\text{massive-string or derivative corrections:}\quad
\alpha'|\mathrm{Riemann}|\ll1
}.
\] For the D3-brane solution the dilaton is constant, \[
e^\Phi=g_s,
\] so the first condition is simply \(g_s\ll1\). The second condition
will lead to \(g_sN\gg1\).

\subsubsection{\texorpdfstring{The source strength of one
D\(p\)-brane}{The source strength of one Dp-brane}}\label{the-source-strength-of-one-dp-brane}

The ten-dimensional bulk action begins schematically as \[
S_{\mathrm{bulk}}
=\frac1{2\kappa_{10}^2}
\int d^{10}x\,\sqrt{-G}\,R+\cdots,
\] while \(N\) coincident D\(p\)-branes contribute the source action \[
S_{\mathrm{source}}
=-N\tau_p
\int d^{p+1}\xi\,
\sqrt{-\det G_{\mathrm{ind}}}
+\cdots.
\] In the conventions used in the book, \[
2\kappa_{10}^2
=(2\pi)^7g_s^2\alpha'^4,
\] and the physical D\(p\)-brane tension is \[
\tau_p
=\frac1{(2\pi)^pg_s\alpha'^{(p+1)/2}}.
\]

The tension is large at weak coupling, \[
\tau_p\sim\frac1{g_s},
\] which is why a D-brane is not an ordinary perturbative closed-string
state. Nevertheless, its gravitational field is governed not by
\(\tau_p\) alone but by Newton's constant times the tension: \[
\begin{aligned}
\kappa_{10}^2N\tau_p
&=
\frac{(2\pi)^7}{2}
g_s^2\alpha'^4
\frac{N}{(2\pi)^pg_s\alpha'^{(p+1)/2}}\\
&=
\frac{(2\pi)^{7-p}}2
g_sN\,
\alpha'^{(7-p)/2}.
\end{aligned}
\] Thus the factor \(1/g_s\) in the D-brane tension cancels one of the
two powers of \(g_s\) in \(\kappa_{10}^2\), leaving \[
\boxed{
\kappa_{10}^2N\tau_p
\sim
g_sN\ell_s^{\,7-p}
}.
\] This is the elementary origin of the combination \(g_sN\).

\subsubsection{Linearized back-reaction on Minkowski
space}\label{linearized-back-reaction-on-minkowski-space}

Let \[
G_{MN}=\eta_{MN}+h_{MN}
\] and linearize the Einstein equation around flat space. Suppressing
tensor indices and numerical coefficients, a stack localized in the
\(9-p\) transverse directions obeys \[
\nabla_\perp^2h
\sim
\kappa_{10}^2N\tau_p\,
\delta^{(9-p)}(y).
\] For \(p<7\), the Green function in \(9-p\) transverse dimensions
behaves as \[
G_\perp(r)\sim\frac1{r^{7-p}}.
\] Therefore \[
h(r)
\sim
\frac{\kappa_{10}^2N\tau_p}{r^{7-p}}
\sim
g_sN
\left(\frac{\ell_s}{r}\right)^{7-p}.
\]

At distances of order the string length, \[
r\sim\ell_s,
\] the dimensionless metric perturbation is parametrically \[
h(\ell_s)\sim g_sN.
\] Hence \[
\boxed{
g_sN\ll1
\quad\Longrightarrow\quad
\text{weak back-reaction outside the string-scale core}
}.
\] In this regime it is consistent to treat the D-branes as
hypersurfaces embedded in nearly flat Minkowski space and quantize open
strings ending on them.

Conversely, \[
\boxed{
g_sN\gtrsim1
\quad\Longrightarrow\quad
\text{the flat-background approximation fails}
}.
\] The branes must then be replaced by their back-reacted geometry.

There is a small qualification. For every nonzero source, the linearized
field becomes large sufficiently close to the idealized brane position.
The claim \(g_sN\ll1\) means that the breakdown occurs inside the
string-scale region, where the geometric supergravity description was
not trustworthy anyway and the microscopic open-string description takes
over.

\subsubsection{\texorpdfstring{The same result from the nonlinear
D\(p\)-brane
solution}{The same result from the nonlinear Dp-brane solution}}\label{the-same-result-from-the-nonlinear-dp-brane-solution}

The full BPS solution for \(N\) coincident D\(p\)-branes, with \(p<7\),
contains a harmonic function of the form \[
H_p(r)
=1+\frac{R_p^{\,7-p}}{r^{7-p}},
\] where \[
R_p^{\,7-p}
=c_p\,g_sN\alpha'^{(7-p)/2}.
\] The constant \(c_p\) depends only on \(p\) and on normalization
conventions. Thus \[
\left(\frac{R_p}{\ell_s}\right)^{7-p}
=c_p\,g_sN.
\]

If \[
g_sN\ll1,
\] then \(R_p\ll\ell_s\): the would-be strongly warped region is smaller
than the string length and is not a macroscopic supergravity throat.
Outside the string-scale core, \(H_p\simeq1\).

If instead \[
g_sN\gg1,
\] then \(R_p\gg\ell_s\): the region in which \(H_p\) differs
substantially from one is macroscopic in string units. The geometry is
strongly deformed away from Minkowski space and back-reaction cannot be
ignored.

This explains an important distinction: \[
\boxed{
\text{large back-reaction does not mean large curvature}
}.
\] A geometry can differ greatly from flat space but vary over a very
large radius \(R_p\). In that case its curvature is small even though
its warping is strong. This is precisely the regime in which classical
supergravity is useful.

\subsubsection{Explicit D3-brane derivation of the curvature
condition}\label{explicit-d3-brane-derivation-of-the-curvature-condition}

For D3-branes the book gives \[
ds^2
=H_3^{-1/2}\eta_{\mu\nu}dx^\mu dx^\nu
+H_3^{1/2}
\left(
dr^2+r^2d\Omega_5^2
\right),
\] with \[
H_3(r)
=1+\frac{L^4}{r^4},
\qquad
L^4=4\pi g_sN\alpha'^2.
\] The same relation immediately measures the back-reaction. At
\(r\sim\ell_s\), \[
H_3-1
\sim
\frac{L^4}{\ell_s^4}
=4\pi g_sN.
\] Thus the flat-space description requires \(g_sN\ll1\), up to the
inessential factor \(4\pi\).

In the near-horizon region \(r\ll L\), the metric becomes \[
ds^2
=\frac{r^2}{L^2}\eta_{\mu\nu}dx^\mu dx^\nu
+\frac{L^2}{r^2}dr^2
+L^2d\Omega_5^2,
\] which is \[
AdS_5(L)\times S^5(L).
\] Both factors have curvature radius \(L\). Their Riemann tensors scale
as \[
R_{MNPQ}\sim\frac1{L^2}.
\] Although the ten-dimensional scalar curvature cancels between the
\(AdS_5\) and \(S^5\) factors, \[
R_{AdS_5}+R_{S^5}=0,
\] the spacetime is not flat. For example, \[
R_{MNPQ}R^{MNPQ}
=\frac{80}{L^4}.
\] The appropriate dimensionless derivative-expansion parameter is
therefore \[
\alpha'
\sqrt{R_{MNPQ}R^{MNPQ}}
\sim
\frac{\alpha'}{L^2}
=\frac1{\sqrt{4\pi g_sN}}.
\] Thus \[
\boxed{
L\gg\ell_s
\quad\Longleftrightarrow\quad
4\pi g_sN\gg1
}.
\]

The leading higher-derivative term in type-IIB string theory is
schematically \[
\alpha'^3R^4.
\] Relative to the two-derivative Einstein term, its size on a
background of radius \(L\) is controlled by \[
\left(\frac{\alpha'}{L^2}\right)^3
\sim
(4\pi g_sN)^{-3/2}.
\] This makes the suppression of \(\alpha'\) corrections at large
\(g_sN\) explicit.

\subsubsection{\texorpdfstring{Why large \(N\) is
necessary}{Why large N is necessary}}\label{why-large-n-is-necessary}

Classical D3-brane supergravity requires both \[
g_s\ll1
\] and \[
g_sN\gg1.
\] Dividing the second inequality by the first parameter shows that the
number of branes must be large: \[
N\gg\frac1{g_s}\gg1.
\] Equivalently, there must be a parametric window \[
\boxed{
\frac1N\ll g_s\ll1
}.
\]

It is slightly misleading to say that ``\(g_sN\) must be large compared
with \(g_s\).'' The two quantities control different effects:

\begin{itemize}
\tightlist
\item
  each individual D-brane is a weak gravitational source of order
  \(g_s\);
\item
  \(N\) coincident branes source the same fields coherently, giving a
  total strength of order \(g_sN\);
\item
  closed-string loops remain controlled by \(g_s\) itself.
\end{itemize}

Large \(N\) therefore permits a useful separation: \[
\boxed{
\text{each brane is weakly coupled}
\quad\text{but}\quad
\text{the stack has a large classical field}
}.
\] This is analogous to a classical electromagnetic field produced by
many individually weak sources.

\subsubsection{Gauge-theory interpretation and the 't Hooft
coupling}\label{gauge-theory-interpretation-and-the-t-hooft-coupling}

For \(N\) coincident D3-branes, the low-energy open-string theory is
four-dimensional \(\mathcal N=4\) \(U(N)\) super-Yang--Mills. The book
uses \[
g_{\mathrm{YM}}^2=2\pi g_s.
\] Its 't Hooft coupling is therefore \[
\lambda
=g_{\mathrm{YM}}^2N
=2\pi g_sN.
\] Equation (18.188) can be written as \[
L^4
=4\pi g_sN\alpha'^2
=2\lambda\alpha'^2.
\] Hence \[
\frac{\alpha'}{L^2}
=\frac1{\sqrt{2\lambda}}.
\] The two descriptions have complementary perturbative regimes: \[
\boxed{
\lambda\ll1
\quad\Longleftrightarrow\quad
\text{weakly coupled open-string gauge theory}
}
\] whereas \[
\boxed{
\lambda\gg1
\quad\Longleftrightarrow\quad
\text{small-curvature gravitational description}
}.
\] This is the strong/weak character of the AdS/CFT correspondence.

At fixed \(\lambda\), \[
g_s
=\frac{\lambda}{2\pi N}.
\] Thus taking \(N\to\infty\) suppresses string loops. To have both
large \(\lambda\) and small \(g_s\), the more precise parametric window
is \[
\boxed{
1\ll\lambda\ll 2\pi N
}.
\]

One can also see bulk quantum-loop suppression directly.
Compactification on the \(S^5\) of radius \(L\) gives parametrically \[
G_5\sim\frac{G_{10}}{L^5}.
\] The dimensionless five-dimensional loop parameter is \[
\frac{G_5}{L^3}
\sim
\frac{g_s^2\alpha'^4}{L^8}
\sim
\frac1{N^2}.
\] Therefore the classical bulk limit is simultaneously the planar
large-\(N\) limit of the gauge theory.

\subsubsection{The three regimes should not be
conflated}\label{the-three-regimes-should-not-be-conflated}

For D3-branes, the approximations can be summarized as follows:

\begin{longtable}[]{@{}
  >{\raggedright\arraybackslash}p{(\linewidth - 4\tabcolsep) * \real{0.3333}}
  >{\raggedright\arraybackslash}p{(\linewidth - 4\tabcolsep) * \real{0.3333}}
  >{\raggedright\arraybackslash}p{(\linewidth - 4\tabcolsep) * \real{0.3333}}@{}}
\toprule\noalign{}
\begin{minipage}[b]{\linewidth}\raggedright
regime
\end{minipage} & \begin{minipage}[b]{\linewidth}\raggedright
coupling conditions
\end{minipage} & \begin{minipage}[b]{\linewidth}\raggedright
appropriate description
\end{minipage} \\
\midrule\noalign{}
\endhead
\bottomrule\noalign{}
\endlastfoot
Weak flat-space D-brane regime & \(g_s\ll1\), \(g_sN\ll1\) &
Perturbative open strings; weakly coupled world-volume gauge theory;
back-reaction neglected \\
Classical but string-scale curved regime & \(g_s\ll1\), \(g_sN\sim1\) &
String loops are small, but the full \(\alpha'\)-corrected classical
string background is needed \\
Classical supergravity regime & \(g_s\ll1\), \(g_sN\gg1\) & Back-reacted
D3 geometry; both string loops and \(\alpha'\) corrections suppressed \\
Quantum-string regime & \(g_s\not\ll1\) & String loops cannot be
neglected; a dual formulation or full quantum string theory is
required \\
\end{longtable}

The two perturbatively useful descriptions do not overlap: \[
\text{weak open-string/gauge theory:}\quad g_sN\ll1,
\] \[
\text{weakly curved supergravity:}\quad g_sN\gg1.
\] AdS/CFT asserts that they are two descriptions of the same system
and, in its strongest form, extends beyond either perturbative limit.

\subsubsection{\texorpdfstring{General D\(p\)-brane
caveat}{General Dp-brane caveat}}\label{general-dp-brane-caveat}

The particularly clean condition \(g_sN\gg1\) for small curvature
applies to the D3-brane throat because the D3 solution has a constant
dilaton and one constant curvature radius \(L\).

For a general D\(p\)-brane, \[
e^\Phi
=g_sH_p^{(3-p)/4}
\] varies with position, and so do the curvature and the effective gauge
coupling. In the decoupling variables one commonly organizes the local
conditions using \[
g_{\mathrm{eff}}^2(U)
\sim
g_{\mathrm{YM}}^2N\,U^{p-3}.
\] Small curvature requires \(g_{\mathrm{eff}}^2(U)\gg1\), while loop
suppression imposes a separate condition on the local dilaton. Thus for
\(p\neq3\) the trustworthy supergravity region can be only a range of
radial positions. The page-734 discussion immediately specializes to
D3-branes, for which \[
g_{\mathrm{eff}}^2=\lambda
\] is independent of scale and the two conditions reduce to the simple
statements above.

\subsubsection{Checks}\label{checks-37}

For one D3-brane at weak coupling, \[
N=1,
\qquad
g_s\ll1,
\] one necessarily has \(g_sN\ll1\). The perturbative D-brane
description is valid, but \[
L\lesssim\ell_s,
\] so the corresponding classical supergravity throat is string-scale
and cannot be trusted.

For fixed small \(g_s\) and sufficiently large \(N\), \[
g_s\ll1,
\qquad
g_sN\gg1,
\] the throat satisfies \(L\gg\ell_s\). The metric is far from Minkowski
space near the stack, but its curvature is small. There is no
contradiction between strong back-reaction and a valid derivative
expansion.

Finally, in the 't Hooft limit \[
N\to\infty,
\qquad
\lambda\ \text{fixed},
\] one has \[
g_s=\frac{\lambda}{2\pi N}\to0.
\] This suppresses string loops, but supergravity additionally requires
\[
\lambda\gg1.
\]

\subsubsection{Result}\label{result-45}

The two transcription corrections are \[
\boxed{
\begin{array}{ccl}
\text{flat-space D-brane picture}
&:&
g_s\ll1,\quad g_sN\ll1,\\[2mm]
\text{classical SUGRA brane picture}
&:&
g_s\ll1,\quad g_sN\gg1.
\end{array}
}
\]

The collective source strength is \[
\boxed{
\kappa_{10}^2N\tau_p
\sim
g_sN\alpha'^{(7-p)/2}
}
\] and therefore, for \(p<7\), \[
h(r)
\sim
g_sN
\left(\frac{\ell_s}{r}\right)^{7-p}.
\]

For D3-branes, \[
\boxed{
H_3=1+\frac{L^4}{r^4},
\qquad
L^4=4\pi g_sN\alpha'^2
}
\] so the \(\alpha'\) expansion parameter is \[
\boxed{
\frac{\alpha'}{L^2}
=\frac1{\sqrt{4\pi g_sN}}
=\frac1{\sqrt{2\lambda}}
}.
\] Hence \[
\boxed{
g_s\ll1
\ \text{suppresses string loops},
\qquad
g_sN\gg1
\ \text{suppresses curvature corrections}
}.
\] Both can hold simultaneously because \[
\boxed{
N\gg1,
\qquad
\frac1N\ll g_s\ll1
}.
\]

\section{Q11: Minkowski space as a section of the projective null
cone}\label{q11-minkowski-space-as-a-section-of-the-projective-null-cone}

In page 738-739 they discussed the conformally compactified Minkowski
space. They embed \(\mathbb{R}^{d-1,1}\) with coordinates
\((y^0,\cdots,y^{d-1})\) into \(\mathbb{R}^{d,2}\) with coordinates
\((y^0,\cdots,y^{d-1},y^d,y^{d+1})\). They said that the Minkowski space
is the intersection of the light-cone with the hyperplane
\(v=y^{d+1}-y^d=1\). How to see this? Then they intersect the light-cone
with the sphere \((y^0)^2+(y^{d+1})^2+\sum_{i=1}^d(y^i)^2=2\), the
resulting space is topologically \(S^1\times S^{d-1}\). Is this the
conformally compactified Minkowski space? How it is relate to the common
Minkowski space?

\begin{quote}
\textbf{Textbook location and related material.} Chapter 18, Section
18.9 around pp.~738--739: the projective null cone in
\(\mathbb R^{d,2}\), conformal compactification, and the Einstein
cylinder.
\end{quote}

\subsection{Answer}\label{answer-46}

\subsubsection{Core point}\label{core-point-46}

Yes. The construction is the standard projective-null-cone description
of conformally compactified Minkowski space. The logical steps are:

\begin{enumerate}
\def\labelenumi{\arabic{enumi}.}
\tightlist
\item
  The ambient null cone in \(\mathbb R^{d,2}\) has dimension \(d+1\),
  one dimension too many.
\item
  Points differing by a nonzero rescaling represent the same conformal
  point: \[
  Y^M\sim\lambda Y^M.
  \]
\item
  On the open patch \(v\neq0\), one may use this rescaling to impose
  \(v=1\). The cone equation then determines \(u=x^2\), leaving exactly
  the \(d\) arbitrary Minkowski coordinates \(x^\mu\).
\item
  The patch \(v=1\) misses the rays with \(v=0\). These are precisely
  the points at conformal infinity.
\item
  Intersecting the cone with the Euclidean sphere fixes the scale
  without discarding the \(v=0\) rays. The result is the compact space
  \(S^1\times S^{d-1}\), or more precisely its antipodal quotient if one
  identifies \(Y\) and \(-Y\).
\end{enumerate}

Thus the precise global statement is \[
\boxed{
\overline{\mathbb R^{d-1,1}}
=
\frac{\{Y\in\mathbb R^{d,2}\setminus\{0\}\mid Y^2=0\}}
{Y\sim\lambda Y}
\simeq
\frac{S^1\times S^{d-1}}{\mathbb Z_2}
}.
\] The book writes \(S^1\times S^{d-1}\) because it is using the double
cover, or equivalently quotienting only by positive rescalings. Ordinary
Minkowski space is an open dense patch of this compact space.

\subsubsection{Ambient metric and the null
cone}\label{ambient-metric-and-the-null-cone}

To avoid confusing ambient and physical coordinates, denote coordinates
on \(\mathbb R^{d,2}\) by \[
Y^M
=
\left(
Y^\mu,Y^d,Y^{d+1}
\right),
\qquad
\mu=0,\ldots,d-1.
\] The ambient metric is \[
\eta_{MN}
=\operatorname{diag}
\left(
-1,+1,\ldots,+1,-1
\right).
\] Therefore \[
Y^2
=
\eta_{\mu\nu}Y^\mu Y^\nu
+(Y^d)^2-(Y^{d+1})^2.
\] Define \[
u=Y^{d+1}+Y^d,
\qquad
v=Y^{d+1}-Y^d.
\] Since \[
(Y^d)^2-(Y^{d+1})^2=-uv,
\] the null-cone equation \(Y^2=0\) becomes equation (18.163b): \[
\boxed{
\eta_{\mu\nu}Y^\mu Y^\nu=uv
}.
\]

The differential ambient metric is correspondingly \[
dS^2
=
\eta_{\mu\nu}dY^\mu dY^\nu
-du\,dv.
\]

The cone has dimension \[
(d+2)-1=d+1.
\] The extra dimension is the radial scaling direction \[
Y^M\longmapsto\lambda Y^M.
\] Because the cone equation is homogeneous, a whole ray lies on the
cone whenever one nonzero point of that ray does. Conformal
compactification uses the ray, not a particular representative on it.

\subsubsection{\texorpdfstring{Why the section \(v=1\) is exactly
Minkowski
space}{Why the section v=1 is exactly Minkowski space}}\label{why-the-section-v1-is-exactly-minkowski-space}

Restrict first to the open part of the cone on which \[
v\neq0.
\] Using the scaling \(Y\sim\lambda Y\), choose the representative with
\[
v=1.
\] Let \[
x^\mu=Y^\mu
\] on this section and define \[
x^2
=\eta_{\mu\nu}x^\mu x^\nu.
\] The cone equation now gives \[
x^2=u.
\] Consequently every point on the \(v=1\) section is uniquely of the
form \[
\left(
Y^\mu,u,v
\right)
=
\left(
x^\mu,x^2,1
\right),
\] with arbitrary \[
x^\mu\in\mathbb R^{d-1,1}.
\] Conversely, every \(x^\mu\in\mathbb R^{d-1,1}\) produces a null
ambient vector of this form. This is already a one-to-one identification
between the section and Minkowski space.

In terms of the original final two coordinates, \[
Y^{d+1}
=\frac{u+v}{2}
=\frac{1+x^2}{2},
\] \[
Y^d
=\frac{u-v}{2}
=\frac{x^2-1}{2}.
\] The embedding is therefore \[
\boxed{
Y^M(x)
=
\left(
x^\mu,
\frac{x^2-1}{2},
\frac{x^2+1}{2}
\right)
}.
\] It is manifestly null: \[
\begin{aligned}
Y^2
&=
x^2
+\frac{(x^2-1)^2}{4}
-\frac{(x^2+1)^2}{4}\\
&=0.
\end{aligned}
\]

The identification is not merely topological. It also gives the
Minkowski metric. Since \(dv=0\) on the \(v=1\) section, \[
dS^2\big|_{v=1}
=
\eta_{\mu\nu}dx^\mu dx^\nu-du\,dv
=
\eta_{\mu\nu}dx^\mu dx^\nu.
\] Equivalently, \(dY^d=dY^{d+1}=x_\mu dx^\mu\), so the last two
contributions to the ambient line element cancel. Thus \[
\boxed{
\mathcal N\cap\{v=1\}
\cong
\mathbb R^{d-1,1}
}
\] as a metric space.

\subsubsection{Why projectivization is needed for the conformal
group}\label{why-projectivization-is-needed-for-the-conformal-group}

The group \(SO(d,2)\) acts linearly on \(Y^M\) and preserves the null
cone. A generic \(SO(d,2)\) transformation does not preserve the
particular section \(v=1\). If \[
Y^M(x)\longmapsto Y'^M,
\] one must rescale the transformed vector back to the section: \[
Y'^M
\longmapsto
\frac{Y'^M}{v'}.
\] The induced transformation of the physical coordinates is \[
x'^\mu
=\frac{Y'^\mu}{v'}.
\] Because the denominator depends on \(x\), a linear ambient
transformation becomes a nonlinear conformal transformation of Minkowski
space.

For some \(x\), one may obtain \[
v'=0.
\] In the ordinary Minkowski chart, this means that \(x'^\mu\) has been
sent to infinity. This is exactly what happens for a special conformal
transformation at the point where its denominator vanishes. The
projective cone includes these \(v=0\) rays, allowing \(SO(d,2)\) to act
globally.

This explains why ordinary Minkowski space alone is insufficient: it is
not closed under finite conformal transformations. Its conformal
compactification is.

\subsubsection{\texorpdfstring{Why the sphere section is
\(S^1\times S^{d-1}\)}{Why the sphere section is S\^{}1\textbackslash times S\^{}\{d-1\}}}\label{why-the-sphere-section-is-s1times-sd-1}

Now fix the scale globally by intersecting the null cone with the
Euclidean sphere used in the book: \[
(Y^0)^2+(Y^{d+1})^2
+\sum_{i=1}^d(Y^i)^2
=2.
\] Introduce \[
A=(Y^0)^2+(Y^{d+1})^2,
\qquad
B=\sum_{i=1}^d(Y^i)^2.
\] The null-cone equation (18.163a) says \[
A=B.
\] The sphere equation says \[
A+B=2.
\] Therefore \[
A=B=1.
\] That is, \[
(Y^0)^2+(Y^{d+1})^2=1
\] and \[
\sum_{i=1}^d(Y^i)^2=1.
\] The first equation defines \(S^1\), while the second defines
\(S^{d-1}\). Hence \[
\boxed{
\mathcal N
\cap
\left\{
|Y|_{\mathrm E}^2=2
\right\}
\simeq
S^1\times S^{d-1}
}.
\]

This sphere is not being used as a spacetime metric condition. Its role
is only to choose a finite-norm representative of every null ray. Every
positive ray intersects it exactly once, including rays with \(v=0\), so
unlike the affine section \(v=1\), it gives a compact global section.

\subsubsection{\texorpdfstring{The antipodal \(\mathbb Z_2\)
subtlety}{The antipodal \textbackslash mathbb Z\_2 subtlety}}\label{the-antipodal-mathbb-z_2-subtlety}

If projective equivalence allows every nonzero real scaling, \[
Y\sim\lambda Y,
\qquad
\lambda\in\mathbb R^\times,
\] then \[
Y\sim-Y.
\] Both \(Y\) and \(-Y\) lie on the sphere section. Thus the sphere
intersects each unoriented projective ray twice, and the true projective
cone is \[
\boxed{
\mathbb P\mathcal N
\simeq
\frac{S^1\times S^{d-1}}{\mathbb Z_2}
},
\] where the \(\mathbb Z_2\) acts antipodally on both factors.

If instead one quotients only by positive scalings, \[
\lambda>0,
\] then the rays are oriented and the result is the double cover \[
S^1\times S^{d-1}.
\] The book states the topology of this double cover and suppresses the
antipodal quotient. Both conventions are common; locally they give the
same conformal geometry.

For Lorentzian physics one often goes one step further and unwraps the
timelike circle: \[
S^1\times S^{d-1}
\longrightarrow
\mathbb R\times S^{d-1}.
\] This universal cover is the Einstein static cylinder and avoids the
closed timelike curves associated with periodic cylinder time.

\subsubsection{Explicit conformal map to the Einstein
cylinder}\label{explicit-conformal-map-to-the-einstein-cylinder}

Parametrize the sphere section by \[
Y^0=\sin T,
\qquad
Y^{d+1}=\cos T,
\] and write the unit vector on \(S^{d-1}\) as \[
Y^d=-\cos\chi,
\qquad
Y^a=\sin\chi\,\widehat n^a,
\qquad
a=1,\ldots,d-1,
\] where \[
\widehat n\in S^{d-2},
\qquad
0\leq\chi\leq\pi.
\] The induced ambient metric on this section is \[
\boxed{
ds_{\mathrm{cyl}}^2
=
-dT^2+d\chi^2
+\sin^2\chi\,d\Omega_{d-2}^2
}.
\] This is the metric on the Lorentzian Einstein cylinder
\(S^1\times S^{d-1}\).

For this parametrization, \[
v
=Y^{d+1}-Y^d
=\cos T+\cos\chi.
\] On the patch \(v\neq0\), rescale to \(v=1\) and define Minkowski
coordinates: \[
x^\mu=\frac{Y^\mu}{v}.
\] Writing \(t=x^0\) and \[
r^2=\sum_{a=1}^{d-1}(x^a)^2,
\] one obtains \[
\boxed{
t
=\frac{\sin T}{\cos T+\cos\chi},
\qquad
r
=\frac{\sin\chi}{\cos T+\cos\chi}
}.
\] Equivalently, \[
\boxed{
t+r
=\tan\left(\frac{T+\chi}{2}\right),
\qquad
t-r
=\tan\left(\frac{T-\chi}{2}\right)
}.
\]

Substitution gives the conformal relation \[
\begin{aligned}
ds_{\mathrm M}^2
&=
-dt^2+dr^2+r^2d\Omega_{d-2}^2\\
&=
\frac1{(\cos T+\cos\chi)^2}
\left(
-dT^2+d\chi^2
+\sin^2\chi\,d\Omega_{d-2}^2
\right).
\end{aligned}
\] Thus \[
\boxed{
ds_{\mathrm{cyl}}^2
=
\Omega^2ds_{\mathrm M}^2,
\qquad
\Omega=\cos T+\cos\chi
}.
\] The cylinder metric is not equal to the Minkowski metric; it is
conformally related to it. That is why the result is called a conformal
compactification.

\subsubsection{Minkowski space is a diamond inside the
cylinder}\label{minkowski-space-is-a-diamond-inside-the-cylinder}

Choose the representative with \[
\Omega=\cos T+\cos\chi>0.
\] The tangent formulas imply \[
-\pi<T+\chi<\pi,
\qquad
-\pi<T-\chi<\pi.
\] Equivalently, \[
\boxed{
|T|+\chi<\pi
}.
\] This is the familiar diamond-shaped Minkowski region in the Einstein
cylinder.

The ordinary Minkowski metric becomes singular when \[
\Omega=0,
\] but the rescaled cylinder metric remains regular. The new finite loci
are the points at infinity:

\begin{itemize}
\tightlist
\item
  future null infinity \(\mathscr I^+\) lies on \(T+\chi=\pi\);
\item
  past null infinity \(\mathscr I^-\) lies on \(T-\chi=-\pi\);
\item
  future and past timelike infinity are \[
  i^+:(T,\chi)=(\pi,0),
  \qquad
  i^-:(T,\chi)=(-\pi,0);
  \]
\item
  spatial infinity is \[
  i^0:(T,\chi)=(0,\pi).
  \]
\end{itemize}

Accordingly, conformal compactification does not make the physical
Minkowski metric regular at infinity. It replaces it by a conformally
equivalent metric that extends smoothly to a boundary at finite
coordinate distance.

\subsubsection{\texorpdfstring{Why equation (18.165) uses the patch
\(u\neq0\)}{Why equation (18.165) uses the patch u\textbackslash neq0}}\label{why-equation-18.165-uses-the-patch-uneq0}

There is a small change of chart between equations (18.163) and (18.165)
that the book does not emphasize.

The sentence before equation (18.164) uses the \(v\neq0\) chart and
fixes \[
v=1.
\] Equation (18.165), however, uses the other affine chart, \[
u\neq0,
\] and writes \[
\left(
Y^\mu,u,v
\right)
=
\left(
u x^\mu,u,ux^2
\right).
\] Projectively setting \(u=1\) gives \[
\left(
Y^\mu,u,v
\right)
=
\left(
x^\mu,1,x^2
\right),
\] which is another copy of Minkowski space.

Let \(x_v^\mu=Y^\mu/v\) be the coordinates in the \(v\)-chart and
\(x_u^\mu=Y^\mu/u\) those in the \(u\)-chart. On their overlap, the cone
equation gives \[
\frac{u}{v}=x_v^2,
\] and hence \[
\boxed{
x_u^\mu
=\frac{x_v^\mu}{x_v^2}
}.
\] The transition function is conformal inversion. Thus the two affine
sections are not contradictory; they are two conformally related
Minkowski charts on the same projective null cone.

\subsubsection{Checks}\label{checks-38}

The dimension count is \[
\dim\mathbb R^{d,2}
-1_{\text{ cone equation}}
-1_{\text{ projective scaling}}
=d,
\] as required for compactified \(d\)-dimensional Minkowski space.

On the \(v=1\) section, \[
Y^d=\frac{x^2-1}{2},
\qquad
Y^{d+1}=\frac{x^2+1}{2},
\] so \[
Y^{d+1}-Y^d=1
\] and \(Y^2=0\) identically. This directly verifies the book's claim.

At the center of the Minkowski diamond, \[
T=0,
\qquad
\chi=0,
\] the conformal map gives \[
t=0,
\qquad
r=0.
\] As \(T+\chi\to\pi\), one has \[
t+r\to+\infty,
\] showing explicitly how null infinity is moved to a finite cylinder
coordinate.

\subsubsection{Result}\label{result-46}

The ambient cone is \[
\mathcal N
=
\left\{
Y\neq0:
\eta_{\mu\nu}Y^\mu Y^\nu=uv
\right\},
\qquad
u=Y^{d+1}+Y^d,
\quad
v=Y^{d+1}-Y^d.
\] On the \(v=1\) section, \[
\boxed{
u=x^2,
\qquad
Y^M(x)
=
\left(
x^\mu,
\frac{x^2-1}{2},
\frac{x^2+1}{2}
\right),
\qquad
ds_{\mathrm{ind}}^2
=\eta_{\mu\nu}dx^\mu dx^\nu
}.
\]

The global conformal compactification is the projective null cone: \[
\boxed{
\overline{\mathbb R^{d-1,1}}
=\mathbb P\mathcal N
\simeq
\frac{S^1\times S^{d-1}}{\mathbb Z_2}
}.
\] The book's \(S^1\times S^{d-1}\) is its double cover.

With \[
\Omega=\cos T+\cos\chi,
\] ordinary Minkowski space is the open diamond \[
\boxed{
|T|+\chi<\pi
}
\] and its metric is related to the Einstein-cylinder metric by \[
\boxed{
ds_{\mathrm M}^2
=
\Omega^{-2}ds_{\mathrm{cyl}}^2,
\qquad
ds_{\mathrm{cyl}}^2
=
-dT^2+d\chi^2
+\sin^2\chi\,d\Omega_{d-2}^2
}.
\] The added locus \(\Omega=0\) is conformal infinity. Equation (18.165)
uses the complementary \(u\neq0\) chart, whose coordinates are related
to those of the \(v\neq0\) chart by inversion: \[
\boxed{
x_u^\mu
=\frac{x_v^\mu}{x_v^2}
}.
\]

\section{\texorpdfstring{Q12: Finite-temperature \(\mathcal N=4\) SYM
and its thermal stress
tensor}{Q12: Finite-temperature \textbackslash mathcal N=4 SYM and its thermal stress tensor}}\label{q12-finite-temperature-mathcal-n4-sym-and-its-thermal-stress-tensor}

In page759, they interpret the result with \(\mathcal{N}=4\) SYM at
finite temperature on the boundary. They said the energy-momentum tensor
is the one of a relativistic plasma at rest in a thermal state \[
\langle T_{\mu \nu}\rangle=\frac{\pi^2}{8}T^4N^2 \mathrm{diag}(3,1,1,1)
\] Give a brief introduction on \(\mathcal{N}=4\) SYM at finite
temperature and this energy-momentum tensor.

\begin{quote}
\textbf{Textbook location and related material.} Chapter 18, Section
18.9 around p.~759: finite-temperature \(\mathcal N=4\) SYM, the AdS
black brane, holographic thermodynamics, and the thermal stress tensor.
\end{quote}

\subsection{Answer}\label{answer-47}

\subsubsection{Core point}\label{core-point-47}

The formula on page 759 is the equilibrium stress tensor of strongly
coupled, large-\(N\) \(\mathcal N=4\) \(SU(N)\) super-Yang--Mills
theory: \[
\boxed{
\left\langle T_{\mu\nu}\right\rangle
=
\frac{\pi^2}{8}N^2T^4
\operatorname{diag}(3,1,1,1)
}.
\] It contains two logically distinct ingredients.

First, equilibrium symmetries and conformal invariance fix the tensor
structure: \[
\left\langle T_{\mu\nu}\right\rangle
=
\operatorname{diag}(\epsilon,p,p,p),
\qquad
\epsilon=3p.
\] Second, AdS/CFT fixes the strong-coupling pressure: \[
p
=\frac{\pi^2}{8}N^2T^4.
\] Therefore \[
\epsilon
=3p
=\frac{3\pi^2}{8}N^2T^4.
\]

The coefficient is not determined by conformal symmetry alone. It is the
classical-supergravity prediction in the limit \[
N\gg1,
\qquad
\lambda=g_{\mathrm{YM}}^2N\gg1,
\qquad
g_s\ll1.
\] At weak coupling the tensor has the same perfect-fluid and traceless
form but a different coefficient.

\subsubsection{\texorpdfstring{Zero-temperature \(\mathcal N=4\)
SYM}{Zero-temperature \textbackslash mathcal N=4 SYM}}\label{zero-temperature-mathcal-n4-sym}

Four-dimensional \(\mathcal N=4\) SYM with gauge group \(SU(N)\)
contains only adjoint fields:

\begin{itemize}
\tightlist
\item
  one gauge field \(A_\mu\);
\item
  six real scalar fields \(\Phi^I\), \(I=1,\ldots,6\);
\item
  four Weyl fermions \(\lambda^A\), \(A=1,\ldots,4\).
\end{itemize}

All fields transform in the adjoint representation, whose dimension is
\[
\dim SU(N)=N^2-1.
\] The theory has sixteen Poincare supercharges and an
\(SU(4)_R\simeq SO(6)_R\) R-symmetry. Its beta function vanishes: \[
\beta(g_{\mathrm{YM}})=0.
\] Consequently, on flat space and at zero temperature it is an exactly
conformal quantum field theory for every value of the coupling.

For \(N\) coincident D3-branes, the microscopic open-string theory
initially has gauge group \(U(N)\). Its diagonal \(U(1)\) describes the
free center-of-mass motion of the stack and decouples from the
interacting sector. The nontrivial CFT is therefore the \(SU(N)\)
theory. Classical gravity retains only the leading large-\(N\) scaling
\[
N^2-1=N^2\left(1-\frac1{N^2}\right)
\longrightarrow N^2.
\]

\subsubsection{What finite temperature
means}\label{what-finite-temperature-means}

The thermal theory is defined by the canonical density matrix \[
\rho_{\mathrm{th}}
=\frac{e^{-\beta H}}{Z(\beta)},
\qquad
Z(\beta)
=\operatorname{Tr}e^{-\beta H},
\qquad
\beta=\frac1T.
\] Thermal expectation values are \[
\langle\mathcal O\rangle_T
=
\operatorname{Tr}
\left(
\rho_{\mathrm{th}}\mathcal O
\right).
\]

In the Euclidean path integral, this is implemented by compactifying
Euclidean time: \[
\tau_E\sim\tau_E+\beta.
\] Bosons are periodic and fermions are antiperiodic: \[
\Phi(\tau_E+\beta)=\Phi(\tau_E),
\qquad
\lambda(\tau_E+\beta)=-\lambda(\tau_E).
\] The unequal boundary conditions break supersymmetry at every nonzero
temperature. Thus ``finite-temperature \(\mathcal N=4\) SYM'' means the
thermal state of a theory whose zero-temperature Lagrangian has
\(\mathcal N=4\) supersymmetry; the thermal state itself is not
supersymmetric.

The temperature introduces a scale, so the thermal state is not
invariant under dilatations. Nevertheless, the operator identity
expressing conformal invariance still constrains its equation of state.
On flat spacetime there is no curvature contribution to the Weyl
anomaly, and hence \[
\left\langle T^\mu{}_\mu\right\rangle_T=0.
\]

\subsubsection{Perfect-fluid form from equilibrium
symmetries}\label{perfect-fluid-form-from-equilibrium-symmetries}

The thermal equilibrium state at rest is:

\begin{itemize}
\tightlist
\item
  invariant under time translations;
\item
  invariant under spatial translations;
\item
  invariant under spatial rotations;
\item
  homogeneous and isotropic;
\item
  at zero chemical potential in the discussion of the book.
\end{itemize}

These symmetries imply \[
\langle T_{0i}\rangle=0,
\] and rotational invariance implies \[
\langle T_{ij}\rangle=p\,\delta_{ij}.
\] Defining \[
\epsilon=\langle T_{00}\rangle,
\] one obtains \[
\boxed{
\langle T_{\mu\nu}\rangle
=
\operatorname{diag}(\epsilon,p,p,p)
}.
\]

Covariantly, the equilibrium stress tensor of a relativistic perfect
fluid is \[
\boxed{
T_{\mu\nu}
=(\epsilon+p)u_\mu u_\nu
+p\,\eta_{\mu\nu}
},
\] where \[
u^\mu u_\mu=-1.
\] In the rest frame and in the mostly-plus convention, \[
u^\mu=(1,0,0,0),
\qquad
u_\mu=(-1,0,0,0),
\] so this formula gives \[
T_{00}=\epsilon,
\qquad
T_{ij}=p\,\delta_{ij}.
\]

This is also the origin of the boosted expression in equation (18.231).
Once \(\epsilon=3p\), \[
T_{\mu\nu}
=p\left(
\eta_{\mu\nu}+4u_\mu u_\nu
\right).
\]

\subsubsection{\texorpdfstring{Conformal invariance fixes
\(\epsilon=3p\)}{Conformal invariance fixes \textbackslash epsilon=3p}}\label{conformal-invariance-fixes-epsilon3p}

Using \[
\eta_{\mu\nu}
=\operatorname{diag}(-,+,+,+),
\] the trace is \[
\begin{aligned}
\left\langle T^\mu{}_\mu\right\rangle
&=
\eta^{\mu\nu}
\left\langle T_{\mu\nu}\right\rangle\\
&=
-\epsilon+3p.
\end{aligned}
\] Flat-space conformal invariance requires the trace to vanish: \[
-\epsilon+3p=0.
\] Therefore \[
\boxed{
\epsilon=3p
}.
\]

Dimensional analysis supplies the temperature dependence. In four
dimensions, \(T_{\mu\nu}\) has mass dimension four, and a homogeneous
thermal state has no scale other than \(T\). Hence \[
p
=N^2T^4f(\lambda)
+O(N^0),
\] where \(f(\lambda)\) is a dimensionless function of the exactly
marginal 't Hooft coupling. Conformal symmetry fixes the power \(T^4\)
and the relation \(\epsilon=3p\), but it does not fix \(f(\lambda)\).

\subsubsection{The dual AdS black brane}\label{the-dual-ads-black-brane}

At large \(N\) and large \(\lambda\), the thermal state has a classical
gravitational dual. It is not empty \(AdS_5\times S^5\), which
represents the vacuum, but an asymptotically AdS black three-brane.

Restoring the AdS radius \(L\), the metric in the coordinates of
equation (18.227) is \[
ds^2
=L^2
\left[
\frac{dU^2}{U^2f(U)}
-U^2f(U)\,dt^2
+U^2d\mathbf x^2
+d\Omega_5^2
\right],
\] where \[
f(U)=1-\frac{U_0^4}{U^4}.
\] The event horizon is at \[
U=U_0.
\]

After Wick rotation \(t=-i\tau_E\), regularity of the Euclidean metric
at the horizon requires a definite period for \(\tau_E\). This period
gives the Hawking temperature \[
\boxed{
T=T_{\mathrm{BH}}=\frac{U_0}{\pi}
}.
\] AdS/CFT identifies the Hawking temperature of the bulk black brane
with the temperature of the boundary SYM state.

The translationally invariant planar horizon explains why the boundary
state is homogeneous. Its rotational symmetry in the three \(\mathbf x\)
directions explains why all three pressures are equal.

\subsubsection{Derivation from the holographic stress
tensor}\label{derivation-from-the-holographic-stress-tensor}

The book first derives the answer through holographic renormalization.
In its Fefferman--Graham expansion, the normalizable boundary datum is
chosen as \[
g^{(2)}_{\mu\nu}
=c\,
\operatorname{diag}(3,1,1,1),
\qquad
c=\frac1{\rho_0^2}.
\] For a flat four-dimensional boundary, equation (18.220) gives \[
\left\langle T_{\mu\nu}\right\rangle
=
\frac{4L^3}{16\pi G_5}
g^{(2)}_{\mu\nu}
=
\frac{L^3}{4\pi G_5}
g^{(2)}_{\mu\nu}.
\]

Equation (18.225) relates \(\rho_0\) to the Hawking temperature: \[
T
=\frac1\pi\sqrt{\frac2{\rho_0}}.
\] Therefore \[
\frac1{\rho_0^2}
=\frac{\pi^4T^4}{4}.
\] Dimensional reduction of type-IIB supergravity on \(S^5\) gives \[
\boxed{
\frac{L^3}{G_5}
=\frac{2N^2}{\pi}
}
\] at leading order in large \(N\). Substituting these relations yields
\[
\begin{aligned}
\left\langle T_{\mu\nu}\right\rangle
&=
\frac{L^3}{4\pi G_5}
\frac{\pi^4T^4}{4}
\operatorname{diag}(3,1,1,1)\\
&=
\frac{\pi^2}{8}N^2T^4
\operatorname{diag}(3,1,1,1).
\end{aligned}
\] This is equation (18.228).

\subsubsection{Independent thermodynamic derivation from the
horizon}\label{independent-thermodynamic-derivation-from-the-horizon}

The same coefficient follows from the Bekenstein--Hawking entropy. The
spatial horizon area per unit boundary volume is \[
\frac{A_{\mathrm h}}{V_3}
=L^3U_0^3.
\] Therefore the entropy density is \[
s
=\frac{A_{\mathrm h}}{4G_5V_3}
=\frac{L^3U_0^3}{4G_5}.
\] Using \[
U_0=\pi T,
\qquad
\frac{L^3}{G_5}=\frac{2N^2}{\pi},
\] one obtains \[
\boxed{
s
=\frac{\pi^2}{2}N^2T^3
}.
\]

At zero chemical potential, the thermodynamic relations are \[
s=\frac{\partial p}{\partial T},
\qquad
\epsilon=Ts-p.
\] Taking \(p(0)=0\) and integrating gives \[
p(T)
=\int_0^T s(T')\,dT'
=\frac{\pi^2}{8}N^2T^4.
\] Then \[
\epsilon
=Ts-p
=\frac{3\pi^2}{8}N^2T^4.
\] Hence \[
\boxed{
\epsilon=3p,
\qquad
p=\frac{\pi^2}{8}N^2T^4,
\qquad
s=\frac{\pi^2}{2}N^2T^3
}.
\]

\subsubsection{Comparison with the free SYM
plasma}\label{comparison-with-the-free-sym-plasma}

At zero coupling, the energy density can be computed by counting
massless adjoint degrees of freedom.

For each adjoint generator, the bosonic degrees of freedom are \[
n_b
=2_{\text{ gauge polarizations}}
+6_{\text{ scalars}}
=8.
\] The four Weyl fermions contribute \[
n_f
=4\times2
=8
\] fermionic degrees of freedom. The energy density of a free
relativistic gas is \[
\epsilon_{\mathrm{free}}
=
\frac{\pi^2}{30}
\left(
n_b+\frac78n_f
\right)
(N^2-1)T^4.
\] Since \[
n_b+\frac78n_f
=8+7
=15,
\] this gives \[
\boxed{
\epsilon_{\mathrm{free}}
=
\frac{\pi^2}{2}
(N^2-1)T^4
\simeq
\frac{\pi^2}{2}N^2T^4
}.
\] The free pressure is \[
p_{\mathrm{free}}
=\frac13\epsilon_{\mathrm{free}}
\simeq
\frac{\pi^2}{6}N^2T^4.
\]

The strong-coupling holographic result is therefore \[
\boxed{
\frac{\epsilon_{\mathrm{strong}}}
{\epsilon_{\mathrm{free}}}
=
\frac{p_{\mathrm{strong}}}
{p_{\mathrm{free}}}
=\frac34
}
\] at leading order in large \(N\). This is equation (18.229).

Unlike the Weyl-anomaly coefficients of \(\mathcal N=4\) SYM, the
thermal free energy is not protected against coupling dependence. The
factor \(3/4\) does not mean that one quarter of the fields disappear.
It means that interactions change the thermodynamic free energy.
Finite-\(\lambda\) corrections arise from \(\alpha'\) corrections in the
bulk, while finite-\(N\) corrections arise from string loops.

\subsubsection{What ``relativistic plasma at rest'' does and does not
mean}\label{what-relativistic-plasma-at-rest-does-and-does-not-mean}

The diagonal form \[
\operatorname{diag}(\epsilon,p,p,p)
\] means:

\begin{itemize}
\tightlist
\item
  there is no momentum or energy flux in the chosen rest frame, so
  \(T_{0i}=0\);
\item
  there is no shear stress in equilibrium;
\item
  the pressure is isotropic;
\item
  the state is homogeneous and time independent.
\end{itemize}

It does not mean that the strongly coupled plasma is a free ideal gas.
The displayed tensor is only the zeroth-order equilibrium form. If
temperature and velocity vary slowly in spacetime, hydrodynamics adds
derivative corrections: \[
T_{\mu\nu}
=
(\epsilon+p)u_\mu u_\nu
+p\eta_{\mu\nu}
-2\eta_{\mathrm{shear}}\sigma_{\mu\nu}
+\cdots.
\] Those corrections vanish for the static homogeneous state in equation
(18.228). This is why the next part of the book proceeds from the
equilibrium black brane to holographic hydrodynamics.

\subsubsection{Checks}\label{checks-39}

The tensor is traceless: \[
\eta^{\mu\nu}
\left\langle T_{\mu\nu}\right\rangle
=
\frac{\pi^2}{8}N^2T^4
(-3+1+1+1)
=0.
\]

It is conserved because it is constant: \[
\partial^\mu
\left\langle T_{\mu\nu}\right\rangle
=0.
\]

The thermodynamic identity also checks: \[
Ts
=
T\left(
\frac{\pi^2}{2}N^2T^3
\right)
=
\frac{\pi^2}{2}N^2T^4
=
\epsilon+p.
\]

Finally, as \(T\to0\), \[
\epsilon,p,s\longrightarrow0,
\] and the planar black-brane horizon disappears, returning to the
vacuum Poincare patch of \(AdS_5\).

\subsubsection{Result}\label{result-47}

Finite-temperature \(\mathcal N=4\) SYM is defined by \[
\rho_{\mathrm{th}}
=Z^{-1}e^{-H/T},
\] with periodic bosons and antiperiodic fermions on the Euclidean
thermal circle. The thermal boundary conditions break supersymmetry.

Homogeneity, isotropy, and conformal invariance imply \[
\boxed{
\left\langle T_{\mu\nu}\right\rangle
=
(\epsilon+p)u_\mu u_\nu
+p\eta_{\mu\nu},
\qquad
\epsilon=3p
}.
\]

In the large-\(N\), large-\(\lambda\) holographic regime, \[
\boxed{
T=\frac{U_0}{\pi},
\qquad
s=\frac{\pi^2}{2}N^2T^3,
\qquad
p=\frac{\pi^2}{8}N^2T^4,
\qquad
\epsilon=\frac{3\pi^2}{8}N^2T^4
}.
\] Therefore, in the rest frame, \[
\boxed{
\left\langle T_{\mu\nu}\right\rangle
=
\frac{\pi^2}{8}N^2T^4
\operatorname{diag}(3,1,1,1)
}.
\]

For the free theory, \[
\epsilon_{\mathrm{free}}
\simeq
\frac{\pi^2}{2}N^2T^4,
\] so \[
\boxed{
\epsilon_{\mathrm{strong}}
=\frac34\epsilon_{\mathrm{free}}
}
\] at leading order in the classical-supergravity limit.

\appendix

\chapter{Essential Notation and Detailed Roadmap of the
Textbook}\label{essential-notation-and-detailed-roadmap-of-the-textbook}

This appendix is designed as a re-entry map. The notation list retains
only symbols used repeatedly across the book and the conversations. The
roadmap then follows the logic of \emph{Basic Concepts of String Theory}
chapter by chapter and section by section, with printed page numbers
from the book.

\section{Essential notation and
conventions}\label{essential-notation-and-conventions}

\subsection{World-sheet geometry and conformal field
theory}\label{world-sheet-geometry-and-conformal-field-theory}

\begin{itemize}
\tightlist
\item
  \(\sigma^a=(\tau,\sigma)\) are world-sheet coordinates; \(z,\bar z\)
  are Euclidean complex coordinates. The world-sheet metric is
  \(h_{ab}\), while \(G_{MN}\) denotes the target-space metric.
\item
  \(T(z)\) and \(\bar T(\bar z)\) are the holomorphic and
  antiholomorphic stress tensors. Their modes are \(L_n\) and
  \(\bar L_n\), with central charges \(c\) and \(\bar c\).
\item
  A local operator has conformal weights \((h,\bar h)\). A physical
  integrated closed-string matter insertion has \((h,\bar h)=(1,1)\); an
  unintegrated insertion includes \(c\bar c\) ghosts and has total
  weight \((0,0)\).
\item
  On the torus, \(\tau=\tau_1+i\tau_2\) is the complex-structure modulus
  and \(q=e^{2\pi i\tau}\). The modular generators act by
  \(S:\tau\mapsto-1/\tau\) and \(T:\tau\mapsto\tau+1\).
\item
  \(\eta(\tau)\) is the Dedekind eta function and \(\vartheta_i\) are
  Jacobi theta functions. A characteristic is written
  \(\vartheta\bigl[\begin{smallmatrix}\alpha\\\beta\end{smallmatrix}\bigr]\).
\item
  In an RCFT, \(\chi_i(\tau)\) is the character of the representation
  \(i\); \(S\) and \(T\) are its modular matrices. Charge conjugation is
  \(C_{ij}=\delta_{i,j^*}\), and \(S^2=C\).
\end{itemize}

\subsection{Quantization, ghosts, and
superstrings}\label{quantization-ghosts-and-superstrings}

\begin{itemize}
\tightlist
\item
  \(\alpha'\) is the inverse string tension scale, with
  \(T_{\mathrm F1}=1/(2\pi\alpha')\); conventions for \(l_s\) differ by
  harmless factors of \(2\pi\).
\item
  \(N\), \(\bar N\) denote oscillator levels. The normal-ordering
  intercept is \(a=1\) for the bosonic string, \(a_{\mathrm{NS}}=1/2\),
  and \(a_{\mathrm R}=0\) in the RNS string.
\item
  \((b,c)\) are reparametrization ghosts. \((\beta,\gamma)\) are
  superconformal ghosts, often bosonized using \(\phi\);
  \(q_{\mathrm{pic}}\) denotes picture number when confusion with the
  torus nome is possible.
\item
  \(Q_{\mathrm{BRST}}\) is the BRST charge. Physical states are BRST
  cohomology classes: \(Q_{\mathrm{BRST}}\ket{\Psi}=0\) modulo
  \(\ket{\Psi}\sim\ket{\Psi}+Q_{\mathrm{BRST}}\ket{\Lambda}\).
\item
  NS and R denote antiperiodic and periodic cylinder fermions,
  respectively. \((-1)^F\) is world-sheet fermion parity; \((-1)^{F_L}\)
  is left-moving spacetime fermion number.
\item
  \(\Omega\) denotes world-sheet orientation reversal. \(I_n\) is
  reflection of \(n\) target coordinates. \(D^\mu{}_{\nu}\) is the
  D-brane gluing matrix, \(+1\) along Neumann and \(-1\) along Dirichlet
  directions.
\end{itemize}

\subsection{Compactification and
geometry}\label{compactification-and-geometry}

\begin{itemize}
\tightlist
\item
  \(X\) is a compactification manifold, often a Calabi--Yau threefold.
  \(\omega\) is its Kähler form and \(\Omega\) its nowhere-vanishing
  holomorphic volume form.
\item
  \(g_{i\bar j}\) is a Kähler metric; barred indices are kept explicit.
  \(H^{p,q}(X)\) and \(h^{p,q}\) denote Dolbeault cohomology and Hodge
  numbers.
\item
  For a cycle \(\Sigma\subset X\), \(N_\Sigma\) is its normal bundle. A
  special-Lagrangian deformation is represented by a harmonic one-form,
  so its infinitesimal moduli are counted by \(b_1(\Sigma)\).
\item
  \(T\), \(U\), and \(S\) commonly denote Kähler, complex-structure, and
  axiodilaton moduli, with the precise assignment stated locally. \(K\)
  and \(W\) are the four-dimensional Kähler potential and
  superpotential.
\item
  In heterotic compactifications, \(P\) and \(Q\) denote weight and root
  lattices when discussing Lie algebras. For an affine algebra, \(k\) is
  the level, \(\psi\) the highest root, and \(\lambda\) an integrable
  highest weight satisfying \((\psi,\lambda)\le k\).
\end{itemize}

\subsection{Couplings, branes, and
dualities}\label{couplings-branes-and-dualities}

\begin{itemize}
\tightlist
\item
  \(\Phi\) is the ten-dimensional dilaton and \(g_s=e^{\Phi_0}\) its
  asymptotic exponential. The string and Einstein metrics obey
  \(G^{(E)}_{MN}=e^{-\Phi/2}G^{(S)}_{MN}\) in ten dimensions.
\item
  \(\kappa_d\) is the \(d\)-dimensional gravitational coupling.
  \(\tau_p\) or \(T_p\) is a D\(p\)-brane tension; \(\mu_p\) is its R--R
  charge in conventions that distinguish the two.
\item
  \(B_2\) is the NS--NS two-form, \(H_3=dB_2\), and \(C_p\) denotes R--R
  potentials. The gauge-invariant D-brane combination is
  \(P[B_2]+2\pi\alpha'F\).
\item
  \(\Gamma^M\) are spacetime gamma matrices. BPS projectors act on
  supersymmetry parameters and reduce the number of preserved
  supercharges.
\item
  \(g_sN\) measures the collective backreaction of \(N\) coincident
  D-branes. Perturbative branes require \(g_sN\ll1\), whereas a weakly
  curved classical near-horizon geometry requires \(g_sN\gg1\) together
  with \(g_s\ll1\).
\end{itemize}

\section{The book's overall argument}\label{the-books-overall-argument}

The book follows a deliberate chain:

\[
\begin{aligned}
\text{classical string}
&\longrightarrow \text{quantization and CFT}
\longrightarrow \text{BRST and perturbative amplitudes} \\
&\longrightarrow \text{superstrings and compactification} \\
&\longrightarrow \text{D-branes, fluxes, dualities, M/F-theory, and AdS/CFT}.
\end{aligned}
\]

The first six chapters construct the bosonic string and its perturbation
theory. Chapters 7--9 add world-sheet supersymmetry, spacetime fermions,
GSO projection, D-branes, and orientifolds. Chapters 10--13 build the
algebraic technology needed for compactified superstrings: Narain
lattices, affine Lie algebras, superconformal theories, BRST vertices,
and covariant lattices. Chapters 14--17 translate world-sheet
consistency into spacetime geometry and effective field theory, then add
branes and fluxes. Chapter 18 reorganizes the perturbative theories into
a nonperturbative network and ends with holography.

\section{Chapter-by-chapter and section-by-section
roadmap}\label{chapter-by-chapter-and-section-by-section-roadmap}

\subsection{Chapter 1: Introduction
(p.~1)}\label{chapter-1-introduction-p.-1}

The introduction motivates strings as quantum gravity rather than
hadronic models. It emphasizes the unavoidable massless closed-string
spin-two state, the role of supersymmetry and anomaly cancellation, the
five ten-dimensional superstring theories, D-branes, dualities,
M-theory, compactification, and the gauge/gravity correspondence. Read
it as a list of destinations whose mechanisms are supplied later.

\subsection{Chapter 2: The Classical Bosonic String
(pp.~7--34)}\label{chapter-2-the-classical-bosonic-string-pp.-734}

\subsubsection{2.1 The relativistic particle
(p.~7)}\label{the-relativistic-particle-p.-7}

The point-particle action introduces reparametrization invariance, the
einbein formulation, constraints, and gauge fixing. This miniature model
previews the Polyakov formulation: an auxiliary world-volume metric
makes symmetries transparent, and its equation of motion returns the
geometric action.

\subsubsection{2.2 The Nambu--Goto action
(p.~10)}\label{the-nambugoto-action-p.-10}

The world-line is replaced by a world-sheet, and proper length by area.
The induced metric gives the Nambu--Goto action and the string tension
sets the fundamental scale. The square-root form is geometrically direct
but inconvenient for quantization.

\subsubsection{2.3 The Polyakov action and its symmetries
(p.~12)}\label{the-polyakov-action-and-its-symmetries-p.-12}

An independent world-sheet metric produces a quadratic action
classically equivalent to Nambu--Goto. The section develops target-space
Poincaré symmetry, world-sheet diffeomorphisms, Weyl symmetry, equations
of motion, stress-tensor constraints, conformal gauge, residual
conformal symmetry, and open-string Neumann/Dirichlet boundary
conditions. These constraints later become Virasoro conditions.

\subsubsection{2.4 Oscillator expansions
(p.~23)}\label{oscillator-expansions-p.-23}

Classical solutions are decomposed into left- and right-moving modes.
Open and closed boundary conditions determine the allowed oscillators,
zero modes encode center-of-mass position and momentum, and the
constraints become relations among mode coefficients. This is the
classical input for canonical quantization.

\subsubsection{2.5 Examples of classical string solutions
(p.~31)}\label{examples-of-classical-string-solutions-p.-31}

Rotating and stretched strings illustrate how energy, angular momentum,
and tension produce Regge behavior. The examples connect the geometric
world-sheet description with the particle spectrum anticipated after
quantization.

\subsection{Chapter 3: The Quantized Bosonic String
(pp.~35--62)}\label{chapter-3-the-quantized-bosonic-string-pp.-3562}

\subsubsection{3.1 Canonical quantization of the bosonic string
(p.~35)}\label{canonical-quantization-of-the-bosonic-string-p.-35}

Oscillator Poisson brackets become commutators, and the Fock space
initially contains timelike negative-norm excitations. Virasoro
constraints, normal ordering, and the intercept define physical states;
null states encode gauge redundancy.

\subsubsection{3.2 Light-cone quantization of the bosonic string
(p.~42)}\label{light-cone-quantization-of-the-bosonic-string-p.-42}

Light-cone gauge solves the constraints explicitly and leaves only
transverse oscillators. Lorentz-algebra closure fixes the critical
dimension and intercept, giving a manifestly positive physical Hilbert
space at the cost of manifest covariance.

\subsubsection{3.3 Spectrum of the bosonic string
(p.~46)}\label{spectrum-of-the-bosonic-string-p.-46}

The open and closed mass formulae, level matching, tachyon, massless
vector, graviton, antisymmetric tensor, and dilaton are identified. The
little-group decomposition explains the spacetime spins carried by
oscillator states.

\subsubsection{3.4 Covariant path-integral quantization
(p.~53)}\label{covariant-path-integral-quantization-p.-53}

The Polyakov functional integral is gauge fixed, producing determinants
and a sum/integral over inequivalent metrics. This viewpoint prepares
the ghost system, moduli integration, and higher-genus perturbation
theory.

\subsubsection{3.5 Appendix: the Virasoro algebra
(p.~59)}\label{appendix-the-virasoro-algebra-p.-59}

The central extension of the conformal algebra is derived from
oscillator normal ordering. The anomaly connects representation theory,
the critical dimension, and the consistency of the gauge-fixed quantum
theory.

\subsection{Chapter 4: Introduction to Conformal Field Theory
(pp.~63--106)}\label{chapter-4-introduction-to-conformal-field-theory-pp.-63106}

\subsubsection{4.1 General introduction
(p.~63)}\label{general-introduction-p.-63}

The section develops conformal transformations, complex coordinates,
operator-product expansions, primary fields, Ward identities, radial
quantization, the state--operator correspondence, the Virasoro algebra,
and basic representation theory. It supplies the language in which a
string background is a two-dimensional CFT.

\subsubsection{4.2 Application to closed string theory
(p.~85)}\label{application-to-closed-string-theory-p.-85}

Free-boson CFT is matched to the closed-string oscillator construction.
Vertex operators encode string states, conformal weights impose
mass-shell conditions, and correlation functions become the local
ingredients of amplitudes.

\subsubsection{4.3 Boundary conformal field theory
(p.~93)}\label{boundary-conformal-field-theory-p.-93}

A boundary identifies left and right conformal transformations, leaving
a diagonal algebra. Open-string spectra, boundary conditions, the
doubling trick, and open--closed channel duality are formulated in CFT
language.

\subsubsection{4.4 Free-boson boundary states
(p.~101)}\label{free-boson-boundary-states-p.-101}

Neumann and Dirichlet conditions are imposed as closed-string gluing
equations. Their coherent-state solutions represent D-branes in the
closed channel and reproduce cylinder amplitudes through overlaps.

\subsubsection{4.5 Crosscap states for the free boson
(p.~103)}\label{crosscap-states-for-the-free-boson-p.-103}

Orientation-reversing identifications produce crosscap gluing
conditions. Crosscap states are the closed-channel representation of
unoriented world-sheets and become the basic description of orientifold
planes.

\subsection{Chapter 5: Parametrization Ghosts and BRST Quantization
(pp.~107--120)}\label{chapter-5-parametrization-ghosts-and-brst-quantization-pp.-107120}

\subsubsection{5.1 The ghost system as a conformal field theory
(p.~107)}\label{the-ghost-system-as-a-conformal-field-theory-p.-107}

Faddeev--Popov gauge fixing of world-sheet diffeomorphisms introduces
the \((b,c)\) system. Its OPEs, conformal weights, stress tensor,
central charge, zero modes, and ghost-number selection rules compensate
the unphysical metric degrees of freedom.

\subsubsection{5.2 BRST quantization
(p.~110)}\label{brst-quantization-p.-110}

Matter and ghost constraints are combined into a nilpotent BRST charge.
Nilpotency encodes criticality; physical states are cohomology classes,
BRST-exact states decouple, and vertex operators acquire their precise
ghost dressing. This is the covariant definition used in amplitude
calculations.

\subsection{Chapter 6: String Perturbation Theory and One-Loop
Amplitudes
(pp.~121--174)}\label{chapter-6-string-perturbation-theory-and-one-loop-amplitudes-pp.-121174}

\subsubsection{6.1 String perturbation expansion
(p.~121)}\label{string-perturbation-expansion-p.-121}

Splitting and joining strings generates world-sheets classified by
topology. The dilaton coupling weights a surface by its Euler
characteristic, turning \(g_s\) into the genus expansion parameter and
organizing closed, open, and unoriented diagrams.

\subsubsection{6.2 Polyakov path integral for the closed bosonic string
(p.~126)}\label{polyakov-path-integral-for-the-closed-bosonic-string-p.-126}

Gauge fixing separates Weyl/diffeomorphism directions from moduli and
conformal Killing transformations. Ghost zero modes, vertex positions,
and the integration over moduli space are derived rather than inserted
ad hoc.

\subsubsection{6.3 The torus partition function
(p.~146)}\label{the-torus-partition-function-p.-146}

The genus-one vacuum amplitude is evaluated both as a path integral and
as a Hamiltonian trace. Modular invariance identifies equivalent tori
and restricts integration to a fundamental domain with invariant
measure; level matching is enforced by the real part of \(\tau\).

\subsubsection{6.4 Torus partition functions for rational CFTs
(p.~152)}\label{torus-partition-functions-for-rational-cfts-p.-152}

The torus amplitude is expanded in finitely many characters. Modular
matrices \(S,T\), charge conjugation, and integer multiplicity matrices
classify modular invariants and connect geometry of the torus with
fusion and representation data.

\subsubsection{6.5 The cylinder partition function
(p.~157)}\label{the-cylinder-partition-function-p.-157}

The annulus is described as an open-string one-loop trace or a closed
string propagating between boundaries. The modular transformation
between channels exposes tadpoles and long-range forces, and the measure
\(dt/(2t)\) reflects its single real modulus and residual automorphism.

\subsubsection{6.6 Boundary states and cylinder amplitudes for RCFTs
(p.~163)}\label{boundary-states-and-cylinder-amplitudes-for-rcfts-p.-163}

Ishibashi states solve chiral gluing conditions, while Cardy states
impose open--closed consistency. The Cardy condition relates boundary
coefficients to the modular \(S\) matrix and ensures nonnegative integer
open-string spectra.

\subsubsection{6.7 Crosscap states, Klein bottle, and Möbius strip
amplitudes
(p.~166)}\label{crosscap-states-klein-bottle-and-muxf6bius-strip-amplitudes-p.-166}

Unoriented one-loop surfaces are written in direct and transverse
channels. Crosscap coefficients, parity projections, Möbius consistency,
and tadpole factorization provide the general RCFT framework for
orientifolds.

\subsubsection{6.8 Appendix: D-brane tension
(p.~171)}\label{appendix-d-brane-tension-p.-171}

The massless part of the cylinder is compared with field-theory
graviton/dilaton exchange. This fixes the normalization and tension of a
D-brane and illustrates how world-sheet factorization determines
spacetime source couplings.

\subsection{Chapter 7: The Classical Fermionic String
(pp.~175--194)}\label{chapter-7-the-classical-fermionic-string-pp.-175194}

\subsubsection{7.1 Motivation for the fermionic string
(p.~175)}\label{motivation-for-the-fermionic-string-p.-175}

World-sheet fermions are introduced to remove central defects of the
bosonic theory and to build spacetime supersymmetric spectra. The
section motivates the RNS formulation.

\subsubsection{7.2 Superstring action and its symmetries
(p.~176)}\label{superstring-action-and-its-symmetries-p.-176}

The Polyakov action is extended to a locally supersymmetric world-sheet
theory. Global and local supersymmetry, zweibein, gravitino, and their
gauge transformations identify the supergravity multiplet whose
redundancies must be fixed.

\subsubsection{7.3 Superconformal gauge
(p.~180)}\label{superconformal-gauge-p.-180}

Gauge fixing leaves free bosons and Majorana fermions plus
superconformal constraints. The supercurrent joins the stress tensor,
and residual transformations form the super-Virasoro symmetry used in
quantization.

\subsubsection{7.4 Oscillator expansions
(p.~189)}\label{oscillator-expansions-p.-189}

Boundary conditions yield Ramond and Neveu--Schwarz moding. Mode
expansions, canonical brackets, and classical super-Virasoro generators
set up the different sectors of the quantum RNS string.

\subsubsection{7.5 Appendix: spinor algebra in two dimensions
(p.~192)}\label{appendix-spinor-algebra-in-two-dimensions-p.-192}

Two-dimensional gamma matrices, chirality, Majorana conditions, and
spinor bilinears fix conventions underlying the action and supersymmetry
transformations.

\subsection{Chapter 8: The Quantized Fermionic String
(pp.~195--222)}\label{chapter-8-the-quantized-fermionic-string-pp.-195222}

\subsubsection{8.1 Canonical quantization
(p.~195)}\label{canonical-quantization-p.-195}

Bosonic and fermionic modes are quantized and the super-Virasoro algebra
acquires central terms. Physical-state conditions and normal-ordering
constants differ between NS and R sectors.

\subsubsection{8.2 Light-cone quantization
(p.~200)}\label{light-cone-quantization-p.-200}

Superconformal constraints eliminate longitudinal modes, leaving
transverse bosons and fermions. Closure and zero-point energies select
ten dimensions and provide a ghost-free description.

\subsubsection{8.3 Spectrum and GSO projection
(p.~202)}\label{spectrum-and-gso-projection-p.-202}

The unprojected NS tachyon and R spinor are reorganized by the GSO
projection. It removes the tachyon, selects chirality, produces equal
bosonic and fermionic degeneracies, and yields spacetime supersymmetry.

\subsubsection{8.4 Path-integral quantization
(p.~208)}\label{path-integral-quantization-p.-208}

Gauge fixing local world-sheet supersymmetry introduces the
\((\beta,\gamma)\) superghosts in addition to \((b,c)\). Central-charge
cancellation and spin-structure sums anticipate the covariant BRST
formalism.

\subsubsection{\texorpdfstring{8.5 Appendix: Dirac matrices and spinors
in \(d\) dimensions
(p.~210)}{8.5 Appendix: Dirac matrices and spinors in d dimensions (p.~210)}}\label{appendix-dirac-matrices-and-spinors-in-d-dimensions-p.-210}

Clifford algebras, Weyl and Majorana conditions, chirality, and
dimensional dependence are collected for use in spectra, R-sector ground
states, supersymmetry, and brane projectors.

\subsection{Chapter 9: Superstrings
(pp.~223--262)}\label{chapter-9-superstrings-pp.-223262}

\subsubsection{9.1 Spin structures and the superstring partition
function
(p.~223)}\label{spin-structures-and-the-superstring-partition-function-p.-223}

Fermion boundary conditions on the torus are summed with GSO phases.
Theta functions organize the four spin structures; modular invariance
fixes their combination, and the Jacobi identity makes the
supersymmetric vacuum amplitude vanish.

\subsubsection{9.2 Boundary states for fermions
(p.~232)}\label{boundary-states-for-fermions-p.-232}

Fermionic gluing conditions are solved in NS--NS and R--R sectors. Spin
labels, zero modes, channel duality, and GSO projection determine the
consistent boundary-state combinations.

\subsubsection{9.3 D-branes (p.~235)}\label{d-branes-p.-235}

Bosonic and fermionic boundary states combine into BPS D-branes. Their
R--R charges, tensions, preserved supersymmetry, open-string spectra,
and cancellation of NS--NS attraction against R--R repulsion are
derived.

\subsubsection{9.4 The type I string
(p.~241)}\label{the-type-i-string-p.-241}

The \(\Omega\) orientifold of type IIB introduces an O9-plane and
D9-branes. Klein-bottle, annulus, and Möbius tadpoles factorize,
requiring 32 D9-branes and producing the anomaly-free \(SO(32)\) theory.

\subsubsection{9.5 Stable non-BPS branes
(p.~251)}\label{stable-non-bps-branes-p.-251}

Orientifold and orbifold projections can remove tachyons from
brane--antibrane constructions. K-theoretic charges and stability beyond
the BPS sector are motivated through explicit spectra.

\subsubsection{9.6 Appendix: theta functions and twisted fermion
partition functions
(p.~254)}\label{appendix-theta-functions-and-twisted-fermion-partition-functions-p.-254}

Products, characteristics, modular transformations, and identities of
theta functions are compiled and related to fermionic traces with
twisted boundary conditions.

\subsection{Chapter 10: Toroidal Compactifications and the
Ten-Dimensional Heterotic String
(pp.~263--320)}\label{chapter-10-toroidal-compactifications-and-the-ten-dimensional-heterotic-string-pp.-263320}

\subsubsection{10.1 Motivation (p.~263)}\label{motivation-p.-263}

Compactification is introduced as the bridge from critical ten
dimensions to lower-dimensional physics, with moduli and new
momentum/winding quantum numbers.

\subsubsection{10.2 Toroidal compactification of the closed bosonic
string
(p.~264)}\label{toroidal-compactification-of-the-closed-bosonic-string-p.-264}

Momentum and winding form left/right lattices controlled by the metric
and \(B\)-field. Mass formulae, level matching, T-duality, enhanced
gauge symmetry, and the Narain moduli space emerge from lattice
geometry.

\subsubsection{10.3 Toroidal partition functions
(p.~280)}\label{toroidal-partition-functions-p.-280}

The compact zero modes are summed into lattice theta functions. Modular
invariance requires even self-dual Lorentzian lattices, explaining why
consistency is encoded arithmetically.

\subsubsection{\texorpdfstring{10.4 The \(E_8\times E_8\) and \(SO(32)\)
heterotic string theories
(p.~284)}{10.4 The E\_8\textbackslash times E\_8 and SO(32) heterotic string theories (p.~284)}}\label{the-e_8times-e_8-and-so32-heterotic-string-theories-p.-284}

Left-moving bosonic and right-moving supersymmetric sectors are
combined. Even self-dual 16-dimensional Euclidean lattices select the
two heterotic gauge groups, and the massless gauge/gravity spectrum is
read off.

\subsubsection{10.5 Toroidal orbifolds
(p.~294)}\label{toroidal-orbifolds-p.-294}

Discrete quotients introduce untwisted projections and twisted sectors
localized at fixed loci. Modular invariance, level matching, fixed-point
degeneracy, and gauge embeddings determine consistent spectra and
furnish solvable singular compactifications.

\subsubsection{10.6 D-branes on toroidal compactifications
(p.~308)}\label{d-branes-on-toroidal-compactifications-p.-308}

T-duality converts Wilson lines into positions and Neumann into
Dirichlet directions. Wrapped branes, images, orientifold planes,
Chan--Paton projections, and tadpoles are organized geometrically; Table
10.3 summarizes a key T-duality orbit.

\subsection{Chapter 11: CFT II - Lattices and Kac--Moody Algebras
(pp.~321--354)}\label{chapter-11-cft-ii---lattices-and-kacmoody-algebras-pp.-321354}

\subsubsection{11.1 Kac--Moody algebras
(p.~321)}\label{kacmoody-algebras-p.-321}

Currents generate affine Lie algebras through their OPEs. The level,
central extension, Sugawara stress tensor, and conformal weights connect
spacetime gauge symmetry to chiral CFT data.

\subsubsection{11.2 Lattices and Lie algebras
(p.~327)}\label{lattices-and-lie-algebras-p.-327}

Roots, coroots, weights, Weyl reflections, conjugacy classes, and ADE
lattices are reviewed. This arithmetic controls both current-algebra
representations and modular invariant compactification lattices.

\subsubsection{11.3 Frenkel--Kac--Segal construction
(p.~337)}\label{frenkelkacsegal-construction-p.-337}

Simply-laced level-one current algebras are realized by free bosons
compactified on root lattices. Cartan currents arise from derivatives
and step operators from lattice vertex operators with cocycles.

\subsubsection{11.4 Fermionic construction and bosonization
(p.~340)}\label{fermionic-construction-and-bosonization-p.-340}

Free fermion bilinears generate current algebras, while bosonization
translates fermions into exponential bosonic vertices. The four \(D_n\)
conjugacy classes encode NS/R sectors and spinor/vector representations.

\subsubsection{11.5 Unitary representations and characters
(p.~343)}\label{unitary-representations-and-characters-p.-343}

Affine primaries, integrability bounds, finite representation sets at
fixed level, characters, and modular transformations are derived. This
makes Kac--Moody theories rational CFTs.

\subsubsection{\texorpdfstring{11.6 Highest-weight representations of
\(\widehat{\mathfrak{su}}(2)_k\)
(p.~349)}{11.6 Highest-weight representations of \textbackslash widehat\{\textbackslash mathfrak\{su\}\}(2)\_k (p.~349)}}\label{highest-weight-representations-of-widehatmathfraksu2_k-p.-349}

The general theory is made explicit: allowed spins, null vectors,
characters, modular \(S\) matrix, and fusion illustrate how level
truncates representation theory.

\subsection{Chapter 12: CFT III - Superconformal Field Theory
(pp.~355--390)}\label{chapter-12-cft-iii---superconformal-field-theory-pp.-355390}

\subsubsection{\texorpdfstring{12.1 \(N=1\) superconformal symmetry
(p.~355)}{12.1 N=1 superconformal symmetry (p.~355)}}\label{n1-superconformal-symmetry-p.-355}

The stress tensor and supercurrent form the \(N=1\) algebra. NS/R
representations, superprimaries, spectral properties, and minimal models
provide the internal-CFT language needed for spacetime supersymmetric
strings.

\subsubsection{\texorpdfstring{12.2 \(N=2\) superconformal symmetry
(p.~370)}{12.2 N=2 superconformal symmetry (p.~370)}}\label{n2-superconformal-symmetry-p.-370}

Two supercurrents and a \(U(1)\) current enlarge the algebra. Spectral
flow relates NS and R sectors, charges BPS states, and turns internal
\(N=2\) symmetry into spacetime supersymmetry after compactification.

\subsubsection{12.3 Chiral ring and topological conformal field theory
(p.~382)}\label{chiral-ring-and-topological-conformal-field-theory-p.-382}

Chiral primaries saturate a unitarity bound and close under nonsingular
OPEs. Twisting converts a supercharge into a BRST operator, producing
topological sectors whose ring data encode protected geometry.

\subsection{Chapter 13: Covariant Vertex Operators, BRST, and Covariant
Lattices
(pp.~391--426)}\label{chapter-13-covariant-vertex-operators-brst-and-covariant-lattices-pp.-391426}

\subsubsection{13.1 Bosonization and first-order systems
(p.~391)}\label{bosonization-and-first-order-systems-p.-391}

General \((b,c)\)-type systems and the superghosts are bosonized.
Background charge, conformal weights, zero-mode anomalies, spin fields,
and picture number prepare a covariant treatment of RNS vertices.

\subsubsection{13.2 Covariant vertex operators, BRST, and picture
changing
(p.~403)}\label{covariant-vertex-operators-brst-and-picture-changing-p.-403}

Physical NS and R vertices are constructed in different pictures. BRST
closure imposes mass shell, transversality, and Dirac equations;
picture-changing relates equivalent representatives and makes superghost
charge balance possible in amplitudes.

\subsubsection{13.3 D-branes and spacetime supersymmetry
(p.~411)}\label{d-branes-and-spacetime-supersymmetry-p.-411}

Fermion gluing on the upper half-plane is lifted to spin fields and
supercharges. A D-brane preserves a diagonal combination \(Q_L+P Q_R\),
and gamma-matrix chirality selects allowed BPS D\(p\)-branes in type
IIA/IIB.

\subsubsection{13.4 The covariant lattice
(p.~415)}\label{the-covariant-lattice-p.-415}

Bosonized fermions and ghosts are assembled into a lattice description.
GSO locality becomes a lattice condition, conjugacy classes organize
sectors, and one-loop characters acquire a uniform covariant form.

\subsubsection{13.5 Heterotic strings in covariant lattice description
(p.~422)}\label{heterotic-strings-in-covariant-lattice-description-p.-422}

Right-moving superstring and left-moving gauge lattices are combined.
Modular invariance and lattice extensions classify ten-dimensional
heterotic constructions and prepare four-dimensional covariant-lattice
models.

\subsection{Chapter 14: String Compactifications
(pp.~427--520)}\label{chapter-14-string-compactifications-pp.-427520}

\subsubsection{14.1 Conformal invariance and spacetime geometry
(p.~427)}\label{conformal-invariance-and-spacetime-geometry-p.-427}

The nonlinear sigma model treats \(G\), \(B\), and \(\Phi\) as
world-sheet couplings. Vanishing Weyl-anomaly coefficients gives their
spacetime field equations order by order in \(\alpha'\), establishing
the world-sheet origin of gravitational dynamics.

\subsubsection{14.2 T-duality and Buscher rules
(p.~432)}\label{t-duality-and-buscher-rules-p.-432}

Gauging an isometry and integrating out different fields derives the
Buscher transformation of \(G\), \(B\), and \(\Phi\). T-duality is
extended from a circle spectrum to general sigma-model backgrounds.

\subsubsection{14.3 Compactification
(p.~439)}\label{compactification-p.-439}

Kaluza--Klein reduction relates internal geometry to lower-dimensional
fields, masses, gauge symmetries, and supersymmetry. Holonomy and
covariantly constant spinors identify the geometry required by unbroken
supercharges.

\subsubsection{14.4 Mathematical preliminaries
(p.~450)}\label{mathematical-preliminaries-p.-450}

Complex, Hermitian, Kähler, symplectic, and Riemannian geometry are
reviewed together with differential forms, cohomology, characteristic
classes, bundles, and index-theoretic tools. This is the reference
section for the later Calabi--Yau analysis.

\subsubsection{14.5 Calabi--Yau manifolds
(p.~473)}\label{calabiyau-manifolds-p.-473}

Calabi--Yau holonomy, the Kähler form, holomorphic volume form, Hodge
diamond, moduli, and special geometry are developed. Topological data
are connected to harmonic representatives and deformations.

\subsubsection{14.6 Type-II compactification on Calabi--Yau threefolds
(p.~487)}\label{type-ii-compactification-on-calabiyau-threefolds-p.-487}

Ten-dimensional fields are expanded in harmonic forms to obtain
four-dimensional \(\mathcal N=2\) multiplets. Type IIA and IIB exchange
Kähler and complex-structure sectors under mirror symmetry.

\subsubsection{14.7 Heterotic compactification on Calabi--Yau threefolds
(p.~497)}\label{heterotic-compactification-on-calabiyau-threefolds-p.-497}

Supersymmetry and anomaly cancellation constrain holomorphic vector
bundles through the Hermitian Yang--Mills equations and Bianchi
identity. The standard embedding, surviving gauge group, chiral index,
and bundle moduli connect geometry to particle spectra.

\subsubsection{14.8 Appendix: Riemannian geometry and Hodge duals
(p.~508)}\label{appendix-riemannian-geometry-and-hodge-duals-p.-508}

Curvature, spin connection, form conventions, Hodge stars, and
dimensional-reduction identities are collected for calculations in
Chapters 14--18.

\subsection{Chapter 15: CFTs for Type-II and Heterotic String Vacua
(pp.~521--584)}\label{chapter-15-cfts-for-type-ii-and-heterotic-string-vacua-pp.-521584}

\subsubsection{15.1 Motivation (p.~521)}\label{motivation-p.-521}

The chapter shifts from large-radius geometry to exact world-sheet
constructions. Orbifolds, orientifolds, Gepner models, and covariant
lattices provide complementary solvable points in the string-vacuum
moduli space.

\subsubsection{15.2 Supersymmetric orbifold compactifications
(p.~522)}\label{supersymmetric-orbifold-compactifications-p.-522}

Twist vectors preserving supersymmetry act on tori and gauge lattices.
Untwisted and twisted spectra, fixed points, Hodge numbers, standard
embeddings, modular constraints, and blow-up modes connect singular CFT
data with smooth Calabi--Yau geometry.

\subsubsection{15.3 Orientifold compactifications
(p.~534)}\label{orientifold-compactifications-p.-534}

World-sheet parity is combined with geometric involutions and
\((-1)^{F_L}\). Closed projections, O-plane loci, Klein-bottle tadpoles,
fractional D-branes, annulus/Möbius amplitudes, and open spectra are
worked out in concrete toroidal models.

\subsubsection{15.4 World-sheet aspects of supersymmetric
compactifications
(p.~544)}\label{world-sheet-aspects-of-supersymmetric-compactifications-p.-544}

Internal \(N=2\) SCFT, bosonization, spectral flow, GSO conditions, and
covariant vertices are used to derive massless type-II and heterotic
spectra. R--R fields appear through field-strength vertices, while
internal current superfields generate gauge bosons.

\subsubsection{15.5 Gepner models (p.~567)}\label{gepner-models-p.-567}

Tensor products of \(N=2\) minimal models with the correct total central
charge give exactly solvable Calabi--Yau compactifications.
Simple-current projections, spectral flow, massless spectra, and the
relation to geometry are explained.

\subsubsection{15.6 Four-dimensional heterotic strings via covariant
lattices
(p.~578)}\label{four-dimensional-heterotic-strings-via-covariant-lattices-p.-578}

Covariant lattice methods construct modular invariant heterotic vacua
directly in four dimensions. Lattice vectors encode spacetime
supersymmetry, gauge symmetry, chirality, and matter representations.

\subsection{Chapter 16: String Scattering Amplitudes and Low-Energy
Effective Field Theory
(pp.~585--640)}\label{chapter-16-string-scattering-amplitudes-and-low-energy-effective-field-theory-pp.-585640}

\subsubsection{16.1 Generalities (p.~585)}\label{generalities-p.-585}

String amplitudes are defined as moduli integrals of BRST-invariant
vertex correlators, with coupling and topology factors. The section
explains what data are on shell and how factorization constrains
consistency.

\subsubsection{16.2 String scattering amplitudes
(p.~588)}\label{string-scattering-amplitudes-p.-588}

Explicit sphere, disk, and related amplitudes produce
Veneziano/Virasoro--Shapiro structures. Ghost fixing, Koba--Nielsen
factors, upper-half-plane propagators, D-brane scattering, poles,
residues, and soft limits are computed.

\subsubsection{16.3 From amplitudes to the low-energy effective field
theory
(p.~615)}\label{from-amplitudes-to-the-low-energy-effective-field-theory-p.-615}

The \(\alpha'\) expansion separates massless poles from contact terms.
Matching on-shell amplitudes determines effective interactions up to
field redefinitions and equation-of-motion terms, explaining both the
power and limits of S-matrix reconstruction.

\subsubsection{16.4 Effective supergravities in ten and eleven
dimensions
(p.~622)}\label{effective-supergravities-in-ten-and-eleven-dimensions-p.-622}

The bosonic actions and field contents of type IIA, type IIB,
heterotic/type I, and eleven-dimensional supergravity are summarized.
Frame changes, dual field strengths, Chern--Simons terms, and
self-duality conventions prepare the duality chapter.

\subsubsection{16.5 Effective action for D-branes
(p.~630)}\label{effective-action-for-d-branes-p.-630}

The Dirac--Born--Infeld and Chern--Simons actions encode tension, gauge
fields, transverse scalars, NS--NS couplings, and R--R charges.
Pullbacks, the combination \(B+2\pi\alpha'F\), and curvature/derivative
corrections translate disk amplitudes into brane dynamics.

\subsubsection{16.6 Appendix: integrals in Veneziano and
Virasoro--Shapiro amplitudes
(p.~636)}\label{appendix-integrals-in-veneziano-and-virasoroshapiro-amplitudes-p.-636}

Beta-function and complex-plane integrals used in tree amplitudes are
derived and analytically continued, clarifying their pole expansions.

\subsection{Chapter 17: Type-II Compactifications with D-Branes and
Fluxes
(pp.~641--674)}\label{chapter-17-type-ii-compactifications-with-d-branes-and-fluxes-pp.-641674}

\subsubsection{17.1 Brane worlds and fluxes
(p.~641)}\label{brane-worlds-and-fluxes-p.-641}

D-branes localize gauge sectors and fluxes generate potentials for
moduli. The section sets the phenomenological and geometric questions:
supersymmetry, chirality, tadpoles, and stabilization.

\subsubsection{17.2 Supersymmetry, tadpoles, and massless spectra
(p.~644)}\label{supersymmetry-tadpoles-and-massless-spectra-p.-644}

Orientifolded Calabi--Yau compactifications with D-branes are
constrained by calibration conditions and R--R charge cancellation.
Intersection numbers determine chiral matter, same-stack strings give
adjoints, and anomalies are tied to tadpoles and Green--Schwarz
couplings.

\subsubsection{17.3 Flux compactifications
(p.~657)}\label{flux-compactifications-p.-657}

NS--NS and R--R fluxes induce the Gukov--Vafa--Witten superpotential.
Flux quantization, D3 tadpoles, F-term conditions, imaginary
self-duality, and no-scale structure explain how complex-structure
moduli and the axiodilaton are stabilized while Kähler moduli remain
flat at tree level.

\subsubsection{\texorpdfstring{17.4 Fluxes and \(SU(3)\) structure group
(p.~663)}{17.4 Fluxes and SU(3) structure group (p.~663)}}\label{fluxes-and-su3-structure-group-p.-663}

Allowing torsion generalizes Calabi--Yau holonomy to \(SU(3)\)
structure. Intrinsic torsion classes, non-closed \(J\) and \(\Omega\),
flux supersymmetry equations, and warped backgrounds organize solutions
beyond Ricci-flat compactifications.

\subsubsection{17.5 Appendix (p.~670)}\label{appendix-p.-670}

Technical conventions and geometric identities used in the flux and
brane analysis are collected for reference.

\subsection{Chapter 18: String Dualities and M-Theory
(pp.~675--774)}\label{chapter-18-string-dualities-and-m-theory-pp.-675774}

\subsubsection{18.1 General remarks
(p.~675)}\label{general-remarks-p.-675}

Duality is presented as an equivalence between descriptions with
different expansion parameters or geometries. Perturbative theories
become coordinate patches on a larger nonperturbative structure.

\subsubsection{18.2 Simple examples: modular invariance and T-duality
(p.~677)}\label{simple-examples-modular-invariance-and-t-duality-p.-677}

Earlier exact symmetries are reinterpreted as prototypes of duality:
modular transformations identify world-sheet descriptions, while
T-duality exchanges momentum/winding, radius, brane dimension, and type
IIA/IIB.

\subsubsection{18.3 Extended objects: generalities
(p.~680)}\label{extended-objects-generalities-p.-680}

\(p\)-branes couple electrically and magnetically to form fields, with
tensions and charges controlled by \(g_s\) and frame choice. Dirac
quantization and world-volume actions set the vocabulary for
nonperturbative states.

\subsubsection{18.4 Central charges and the BPS bound
(p.~684)}\label{central-charges-and-the-bps-bound-p.-684}

Extended charges appear in supersymmetry algebras. Positivity gives BPS
inequalities; saturation implies shortened multiplets, preserved
projector equations, and protected masses/tensions that can be followed
across strong coupling.

\subsubsection{18.5 Brane solutions in supergravity
(p.~688)}\label{brane-solutions-in-supergravity-p.-688}

Symmetry and supersymmetry determine warped metrics, dilaton profiles,
form fluxes, and harmonic functions. Fundamental strings, NS5-branes,
D\(p\)-branes, and their near-horizon regions connect source actions
with classical backreaction.

\subsubsection{18.6 Nonperturbative dualities
(p.~705)}\label{nonperturbative-dualities-p.-705}

Type-IIB self-duality organizes \((p,q)\) strings and five-branes; type
I and heterotic strings are S-dual; compactified theories exhibit
U-duality. Brane spectra and BPS quantities supply quantitative tests
beyond perturbation theory.

\subsubsection{18.7 M-theory (p.~716)}\label{m-theory-p.-716}

Strongly coupled type IIA opens an eleventh dimension. Reduction of
eleven-dimensional supergravity yields IIA fields, while M2/M5 branes
reduce or wrap to strings and D-branes. The supersymmetry algebra
unifies their charges.

\subsubsection{18.8 F-theory (p.~724)}\label{f-theory-p.-724}

The varying type-IIB axiodilaton becomes the complex structure of an
auxiliary elliptic fiber. Seven-brane monodromies, elliptic K3,
orientifold limits, heterotic duality, singular fibers, gauge
enhancement, and localized matter are encoded geometrically.

\subsubsection{18.9 AdS/CFT correspondence
(p.~733)}\label{adscft-correspondence-p.-733}

The D3-brane near-horizon limit produces \(AdS_5\times S^5\) and leads
to the equivalence with four-dimensional \(\mathcal N=4\) SYM. The
section develops the parameter map, conformal compactification,
bulk/boundary dictionary, correlation functions, Wilson loops, black
branes, and finite-temperature thermodynamics, ending with holography as
the strong-coupling culmination of the book.

\section{Rapid re-entry guide}\label{rapid-re-entry-guide}

When returning after a long break, use the following routes:

\begin{itemize}
\tightlist
\item
  For \textbf{world-sheet foundations}, revisit Sections 2.3, 3.2,
  4.1--4.3, 5.2, and 6.2--6.3.
\item
  For \textbf{superstring spectra and D-branes}, revisit Sections 8.3,
  9.1--9.4, 13.2--13.3, and 16.5.
\item
  For \textbf{modular invariance, lattices, and exact CFT}, revisit
  Sections 6.4--6.7, 10.2--10.5, 11.1--11.6, 12.1--12.3, and 15.2--15.6.
\item
  For \textbf{Calabi--Yau compactification}, revisit Sections 14.1 and
  14.4--14.7, followed by 17.2--17.4.
\item
  For \textbf{effective actions and amplitudes}, revisit Sections
  16.1--16.5 and the beta-function argument of Section 14.1.
\item
  For \textbf{dualities and modern geometry}, read Chapter 18 in order;
  its logic depends especially on BPS bounds (18.4), brane solutions
  (18.5), and the compactification technology of Chapters 10, 14, and
  17.
\end{itemize}

The shortest conceptual reconstruction is: world-sheet gauge symmetry
forces CFT and BRST; modular consistency and supersymmetry restrict the
perturbative theories; compactification converts CFT data into geometry
and lower-dimensional fields; D-branes and fluxes make gauge sectors,
chirality, and potentials; BPS protection lets those objects be followed
to strong coupling; dualities then reveal M-theory, F-theory, and
holography as different geometric organizations of the same underlying
degrees of freedom.

\backmatter
\end{document}
